\documentclass[12pt]{report}

% ---------------------------------------------------------------
% 字型與中文排版
% ---------------------------------------------------------------
% 作者在 macOS 上用 Times New Roman／Songti TC／PingFang TC；
% 其他系統（例如 Linux）找不到這些字型時，改用 TeX Gyre Termes 與
% Noto Serif／Sans CJK TC，同一份原始碼兩邊都能編。
\usepackage{fontspec}
\usepackage{xeCJK}
\IfFontExistsTF{Times New Roman}
    {\setmainfont{Times New Roman}}
    {\setmainfont{TeX Gyre Termes}}
\IfFontExistsTF{Songti TC}
    {\setCJKmainfont{Songti TC}}
    {\setCJKmainfont{Noto Serif CJK TC}}
\IfFontExistsTF{PingFang TC}
    {\setCJKsansfont{PingFang TC}}
    {\setCJKsansfont{Noto Sans CJK TC}}
% 等寬字用 Latin Modern Mono。它不是系統字型，只在 TeX Live 樹裡；
% 用「檔名＋副檔名」讓 XeTeX 透過 kpathsea 找檔案，不寫死 TeX Live
% 的年份路徑。
\setmonofont{lmmono10-regular}[
    Extension = .otf,
    AutoFakeBold = 2.5,
]

% ---------------------------------------------------------------
% 本書排版套件（packages/ 目錄）
% ---------------------------------------------------------------
\usepackage{sleek}
\usepackage{sleek-code}
\usepackage{sleek-diagrams}
\usepackage{sleek-recap}
\usepackage{sleek-aki}
% 注意：sleek-theorems.sty 不載入——本書沒有用到定理類環境，而它依賴的 thmtools
% 和 2026-11 以後的 LaTeX 核心衝突（\newtheorem{lem}[thm] 會報「\c@lem already
% defined」），載入的話整本書編不過。
% 注意：sleek-listings.sty 不載入——它是舊版 lstlistings/inconsolata 排版套件，
% 本書實際用的是 sleek-code.sty 的 minted 版 code/example 環境，兩者同時載入會讓
% minted 抓不到 \setmonofont 設定的等寬字（回退成 Latin Modern Roman），
% 且會讓程式碼裡的直引號被轉成全形彎引號。

% 章節標題中文化：「第 N 章」
\usepackage{titlesec}
\titleformat{\chapter}{\normalfont\huge\bfseries}{第~\thechapter~章}{1em}{}
\titlespacing*{\chapter}{0pt}{-20pt}{20pt}
% 頁首與目錄標題中文化（預設是 CHAPTER 1. 和 CONTENTS）
\renewcommand{\contentsname}{目錄}
\renewcommand{\chaptermark}[1]{\markboth{第~\thechapter~章\quad #1}{}}

% 行內程式碼（\texttt）不能斷行，段落裡出現較長的識別字時容易凸出右邊界；
% 允許 TeX 多拉開一點字距來排版。
\setlength{\emergencystretch}{3em}

% PDF 的書名資訊另外寫，不從 \title 抓（\title 裡有換行和系列名稱）。
\hypersetup{
    pdftitle={C++ 演算法解題 2：動態規劃與圖論},
    pdfauthor={林子祐},
    pdfsubject={註解 Annotated 演算法系列 第 2 本},
}

\title{C++ 演算法解題 2\\[0.3em]{\LARGE 動態規劃與圖論}\\[1.2em]{\large 註解 Annotated 演算法系列}}
\author{林子祐}
\date{\today}

\begin{document}

\maketitle
\romantableofcontents

% ---------------------------------------------------------------
% 各章內容放在 "chapters/<單元>/chapter.tex"。
% 章號一開始就用最後的編號：還沒寫的章先註解掉，已寫的章前面用
% \setcounter{chapter}{N - 1} 把章號定好，之後補上新章時不用改號。
% （Makefile 編單章時會自己從目錄名稱算章號，這幾行不影響單章 PDF。）
% ---------------------------------------------------------------

\chapter{背包 DP}
\label{ch:knapsack}

出國前一晚，手機只剩 $8$ GB 的空間。 想下載的東西有四樣：離線地圖 $2$ GB、
一組播客 $3$ GB、一款遊戲 $4$ GB、一季影集 $5$ GB。 你替它們打了「想要的程度」，
分別是 $3$、$4$、$5$、$7$ 分。 每一樣都只能整個下載，怎麼挑才能讓分數加起
來最高？ 往下讀之前先自己挑一次。 再想一個問題：如果候選有一百樣、空間是
$64$ GB，你挑的方法還用得上嗎？

《C++ 演算法解題 1：APCS 實作題核心》（以下簡稱第 1 本）第~7~章 7.4 節試過「分數除以大小，比值大的先拿」。 這裡地圖的比值 $1.5$ 最
高，影集 $1.4$ 第二，先下載這兩樣，用掉 $7$ GB、得到 $10$ 分，剩下的 $1$ GB
什麼都放不下。 可是播客加影集剛好 $8$ GB，有 $11$ 分。 第 1 本第~9~章 9.7 節接著說
明了只記「前 $i$ 樣的最佳值」不夠，還要記用掉多少空間，狀態變成兩維。 這一
章把這件事做完。 這類「容量有限、每樣東西有大小和價值、挑一些放進去」的問題
叫\textbf{背包問題}（knapsack problem）。

第~\ref{sec:knap-01}~節先用兩維的表解最基本的 0/1 背包，第~\ref{sec:knap-1d}~節
把表壓成一維，說清楚迴圈為什麼要倒著跑。 之後的幾節都是同一張表的小改
動：改初始值，就能要求「剛好裝滿」並標出湊不出來的容量；改迴圈方向，每樣東西
就能拿很多次；把取最大值換成加法，就能數方法數；把有數量上限的東西拆成幾包，
就回到 0/1 背包。 最後一節整理幾種常見的變形。

這一章用到第 1 本第~9~章的動態規劃四件事（狀態、轉移、初始值、計算順序）和第 1 本第~10~章
10.3 節的滾動陣列，不再重講。 語法只用到 \texttt{vector}（包括
\texttt{vector<vector<int>>} 的兩維表）、\texttt{long long} 和 \texttt{sort}
加 lambda。 《C++ 程式設計：給學過 C 的人》（以下簡稱銜接書）第~10~章的練習 3
寫過一維的 0/1 背包，當時只簡短說明了迴圈要由大到小，第~\ref{sec:knap-1d}~節會把理由講完整。

\section{兩維的表}
\label{sec:knap-01}

先把問題寫清楚。 有 $n$ 件物品，第 $i$ 件的重量（大小）是 $w_i$、價值是
$v_i$，背包最多裝重量 $W$。 每件物品只有一件，可以放或不放，求放進去的總價值
最大是多少。 每件只有「放」和「不放」兩種，所以叫 \textbf{0/1 背包}（0/1
knapsack）。 章首的手機就是 $n = 4$、$W = 8$，四件物品依序是 $(2, 3)$、
$(3, 4)$、$(4, 5)$、$(5, 7)$，括號裡是（重量，價值）。

最直接的做法是第 1 本第~4~章 4.2 節的子集合枚舉：每件放或不放，$n$ 件共 $2^n$ 種，
每種算出總重量和總價值。 $n = 4$ 只有 $16$ 種，手算都行；$n = 100$ 時
$2^{100}$ 種，永遠跑不完。 要像第 1 本第~9~章那樣，找出會重複出現的小問題。

\subsection{狀態多記一維}

照第 1 本第~9~章的四件事來設計。 一件一件決定，決定到第 $i$ 件時，之後還能放什麼，
只取決於「已經用掉多少容量」，和前面挑了哪幾件無關。 所以狀態要有兩個數：
看到第幾件，以及容量。

\begin{itemize}
\item \textbf{狀態}：$\mathrm{best}(i, c)$ ＝ 只考慮前 $i$ 件物品、總重量不超
過 $c$ 時，能得到的最大總價值。
\item \textbf{轉移}：看第 $i$ 件放不放。 不放，前 $i - 1$ 件仍然可以用滿 $c$
的容量，得到 $\mathrm{best}(i - 1, c)$；放（需要 $c \ge w_i$），前 $i - 1$
件只剩 $c - w_i$ 的容量，得到 $\mathrm{best}(i - 1, c - w_i) + v_i$。 兩者取
大的。
\item \textbf{初始值}：$\mathrm{best}(0, c) = 0$，一件都沒有時什麼都放不了，對
每個 $c$ 都一樣。
\item \textbf{計算順序}：第 $i$ 列只用到第 $i - 1$ 列，所以 $i$ 由小到大，一
列一列填。 同一列裡的 $c$ 誰先誰後都可以。
\end{itemize}

答案是 $\mathrm{best}(n, W)$。 寫成式子：
\[
\mathrm{best}(i, c) = \max\bigl(\mathrm{best}(i - 1, c),\ \mathrm{best}(i - 1, c - w_i) + v_i\bigr),
\quad \text{第二項只在 } c \ge w_i \text{ 時才有。}
\]

下圖是手機的例子填完的表。 第 $i$ 列的每一格，只看正上方和左上方某一格：正
上方是不放第 $i$ 件，往左 $w_i$ 格的那一格是放。

\begin{center}
\begin{tikzpicture}[x=0.8cm, y=0.62cm,
    t/.style={draw=black!45, minimum width=0.8cm, minimum height=0.62cm,
              inner sep=0pt, font=\small\ttfamily},
    src/.style={t, draw=coreStroke, thick, fill=coreFill},
    dst/.style={t, draw=memStroke, very thick, fill=memFill}]
\foreach \c in {0,...,8} { \node[font=\small] at (\c,0.9) {\c}; }
\node[font=\small, left] at (-0.6,0.9) {$c$};
\foreach \y/\lab in {0/{$i=0$},1/{地圖 (2,3)},2/{播客 (3,4)},3/{遊戲 (4,5)},5.2/{影集 (5,7)}} {
    \node[font=\small, left] at (-0.6,-\y) {\lab};
}
\foreach \c/\v in {0/0,1/0,2/0,3/0,4/0,5/0,6/0,7/0,8/0} { \node[t] at (\c,0) {\v}; }
\foreach \c/\v in {0/0,1/0,2/3,3/3,4/3,5/3,6/3,7/3,8/3} { \node[t] at (\c,-1) {\v}; }
\foreach \c/\v in {0/0,1/0,2/3,3/4,4/4,5/7,6/7,7/7,8/7} { \node[t] at (\c,-2) {\v}; }
\foreach \c/\v in {0/0,1/0,2/3,3/4,4/5,5/7,6/8,7/9,8/9} { \node[t] at (\c,-3) {\v}; }
\foreach \c/\v in {0/0,1/0,2/3,3/4,4/5,5/7,6/8,7/10} { \node[t] at (\c,-5.2) {\v}; }
\node[src] (up) at (8,-3) {9};
\node[src] (ul) at (3,-3) {4};
\node[dst] (me) at (8,-5.2) {11};
\draw[diagarrow] (up) -- node[right, font=\scriptsize, xshift=4pt] {不放：$9$} (me);
\draw[diagarrow] (ul.south east) -- (me.north west);
\node[font=\scriptsize] at (4.6,-4.3) {放：$4 + 7$};
\end{tikzpicture}
\diagcaption{手機例子的 $\mathrm{best}(i, c)$。 右下角的 $11$ 從兩格算出：不下
載影集是正上方的 $9$；下載影集要先留 $5$ GB，前三樣只能用 $3$ GB，得到左上方
的 $4$，再加影集的 $7$。}
\label{fig:knap-table}
\end{center}

讀幾格確認意思。 第 $2$ 列（地圖和播客）的 $c = 5$ 是 $7$：兩樣都下載剛好
$5$ GB。 第 $3$ 列的 $c = 8$ 是 $9$：前三樣全下載要 $9$ GB，放不下，最好是播
客加遊戲。 右下角就是答案 $11$。

\subsection{找出選了哪幾件}

表格只存了最大價值。 要知道是哪幾件，從右下角往回走：如果
$\mathrm{best}(i, c)$ 和正上方的 $\mathrm{best}(i - 1, c)$ 相同，不放第 $i$
件也能得到這個值，就當作沒放，往上走一列；不同，就一定放了第 $i$ 件，往上走一
列並把 $c$ 減掉 $w_i$。

\begin{example}[label=ex:knap-2d]{cpp}{(phone\_2d.cpp) 兩維的表，並找出放了哪幾件}
#include <algorithm>
#include <iostream>
#include <vector>
using namespace std;

int main() {
    int n, cap;
    cin >> n >> cap;
    vector<int> w(n + 1), v(n + 1);         // 物品編號 1 到 n
    for (int i = 1; i <= n; i++) {
        cin >> w[i] >> v[i];
    }
    // best[i][c] = 只看前 i 件、總重量不超過 c 的最大價值
    vector<vector<int>> best(n + 1, vector<int>(cap + 1, 0));
    for (int i = 1; i <= n; i++) {
        for (int c = 0; c <= cap; c++) {
            best[i][c] = best[i - 1][c];             // 不放
            if (c >= w[i]) {                         // 放得下才能放
                best[i][c] = max(best[i][c], best[i - 1][c - w[i]] + v[i]);
            }
        }
    }
    cout << best[n][cap] << '\n';

    // 從右下角往回走，找出放了哪幾件
    vector<int> picked;
    for (int i = n, c = cap; i >= 1; i--) {
        if (best[i][c] != best[i - 1][c]) {          // 放了第 i 件
            picked.push_back(i);
            c -= w[i];
        }
    }
    reverse(picked.begin(), picked.end());       // 改成由小到大
    for (int k = 0; k < (int)picked.size(); k++) {
        cout << (k > 0 ? " " : "") << picked[k];
    }
    cout << '\n';
    return 0;
}
\end{example}

\iosample{
4 8 \\
2 3 \\
3 4 \\
4 5 \\
5 7 \\
}{
11 \\
2 4 \\
}

往回走的過程：$\mathrm{best}(4, 8) = 11 \ne \mathrm{best}(3, 8) = 9$，放了第
$4$ 件，$c$ 變成 $3$；$\mathrm{best}(3, 3) = \mathrm{best}(2, 3) = 4$，第 $3$
件沒放；$\mathrm{best}(2, 3) = 4 \ne \mathrm{best}(1, 3) = 3$，放了第 $2$ 件，
$c$ 變成 $0$；最後 $\mathrm{best}(1, 0) = \mathrm{best}(0, 0)$。 放的是播客和
影集。

時間是表格的格數 $O(nW)$，每格 $O(1)$。 要注意這個 $W$ 指的是容量的數值大小，和
輸入裡有幾個數無關：$n = 100$、$W = 10^5$ 是 $10^7$ 格，很快；同樣 $100$ 件物品，
如果 $W$ 到 $10^9$，表格根本開不出來。 所以看到背包題，先看容量（或其他要當
成表格一維的量）最大是多少。

\section{一維陣列：為什麼要倒著填}
\label{sec:knap-1d}

$n = 200$、$W = 10^5$ 時，\texttt{long long} 的兩維表要 $200 \times 10^5
\times 8$ 位元組，大約 $160$ MB，已經接近很多題目的記憶體限制。 轉移只用到上一
列，可以照第 1 本第~10~章 10.3 節的滾動陣列，只留一列 \texttt{best[0..W]}：處理第 $i$ 件之前，它
存的是第 $i - 1$ 列；處理完，變成第 $i$ 列。

第 1 本第~10~章的格子是由左往右填，因為轉移需要的是「這一列」的左邊那格。 背包的轉
移需要的是「上一列」的 \texttt{best[c - w]}，在 $c$ 的左邊。 如果由左往右填，
填到 \texttt{best[c]} 時，左邊的 \texttt{best[c - w]} 已經被這一列的新值蓋掉
了。 用手機的第一件物品（地圖，$w = 2$、$v = 3$）看兩種順序，一開始整列都是
$0$：

\begin{center}
\begin{tabular}{l|ccccccccc}
\toprule
$c$ & 0 & 1 & 2 & 3 & 4 & 5 & 6 & 7 & 8 \\
\midrule
處理前（第 $0$ 列） & 0 & 0 & 0 & 0 & 0 & 0 & 0 & 0 & 0 \\
$c$ 由大到小 & 0 & 0 & 3 & 3 & 3 & 3 & 3 & 3 & 3 \\
$c$ 由小到大 & 0 & 0 & 3 & 3 & 6 & 6 & 9 & 9 & 12 \\
\bottomrule
\end{tabular}
\end{center}

由小到大時，$c = 2$ 先算出 $3$（放一份地圖）；到了 $c = 4$，讀到的
\texttt{best[2]} 已經是 $3$，再加 $3$ 得到 $6$，等於放了兩份地圖。 $c = 8$
的 $12$ 是四份地圖。 地圖只有一份，這是錯的。

由大到小時，填 \texttt{best[c]} 用到的 \texttt{best[c - w]} 在左邊，還沒輪到，
存的一定是上一列的值：

\begin{center}
\begin{tikzpicture}[x=0.95cm, y=1cm]
\foreach \c/\v in {0/0,1/0,2/0,3/0,4/0,5/0} { \node[cell, fill=black!8] (b\c) at (\c,0) {\v}; }
\node[cell, draw=optStroke, very thick] (b6) at (6,0) {0};
\foreach \c/\v in {7/3,8/3} { \node[hlcell] (b\c) at (\c,0) {\v}; }
\foreach \c in {0,...,8} { \node[idxlbl] at (\c,-0.65) {\c}; }
\draw[decorate, decoration={brace, amplitude=5pt, mirror}] (-0.45,-0.95) -- (5.45,-0.95)
    node[midway, below=5pt, font=\small] {還是第 $0$ 列（舊值）};
\draw[decorate, decoration={brace, amplitude=5pt, mirror}] (6.55,-0.95) -- (8.45,-0.95)
    node[midway, below=5pt, font=\small] {已經是第 $1$ 列};
\draw[diagarrow] (b4.north) to[bend left=35] node[above, font=\small] {$\texttt{best[4]} + 3$} (b6.north);
\end{tikzpicture}
\diagcaption{由大到小處理地圖，正要填 \texttt{best[6]}。 它需要的
\texttt{best[4]} 在灰色的區域，還是處理地圖之前的值，所以地圖最多只會被放一次。}
\label{fig:knap-reverse}
\end{center}

\begin{example}[label=ex:knap-1d]{cpp}{(phone\_1d.cpp) 0/1 背包的一維寫法}
#include <algorithm>
#include <iostream>
#include <vector>
using namespace std;

int main() {
    int n, cap;
    cin >> n >> cap;
    // best[c]：目前看過的物品裡，重量不超過 c 的最大價值
    vector<int> best(cap + 1, 0);
    for (int i = 0; i < n; i++) {
        int w, v;
        cin >> w >> v;              // 邊讀邊算，不用存下所有物品
        for (int c = cap; c >= w; c--) {          // 由大到小
            best[c] = max(best[c], best[c - w] + v);
        }
    }
    cout << best[cap] << '\n';
    return 0;
}
\end{example}

\iosample{
4 8 \\
2 3 \\
3 4 \\
4 5 \\
5 7 \\
}{
11 \\
}

\begin{itemize}
\item 內層迴圈停在 \texttt{c >= w}。 $c < w$ 的格子放不下這件，維持上一列的值，
本來就不用動，也就不會讀到 \texttt{best} 的負索引。
\item 外層一定是物品、內層才是容量。 每件物品處理一次，處理完它，陣列就整個前
進一列。
\item 把內層改成 \texttt{for (int c = w; c <= cap; c++)}，範例的輸入會印出
\texttt{12}，就是四份地圖。 範例的答案很容易剛好對，這種錯要用小輸入手算或
對拍才抓得到。
\end{itemize}

壓成一維之後，用第~\ref{sec:knap-01}~節的方法往回找選了哪幾件就做不到了，因為
前面幾列已經被蓋掉。 題目要輸出選法時，用兩維的表（或另外記錄每一格是怎麼來
的）。

\subsection{例題：總務股長的挑選法差多少}
\label{sec:ex-treasurer}

\textbf{題意}（\akiref{how-far-off-is-the-treasurers-rule}{總務股長的挑選法差多少}）

\begin{itemize}
\item 班費 $B$ 元，有 $n$ 種材料，第 $i$ 種價格 $p_i$、效益 $e_i$，每種最多買一
份。
\item 總務股長的方法：依 $e_i / p_i$ 由大到小排；比值相同時價格低的在前，再相
同時編號小的在前。 依序看，剩下的錢買得起就買，買不起就跳過、繼續看下一種。
\item 輸出兩個數：照這個方法得到的效益總和，以及總價不超過 $B$ 時效益總和的最
大值。
\item $n \le 200$，$B \le 10^5$，$p_i \le 10^5$，$e_i \le 10^9$。
\end{itemize}

把章首的手機換成材料：$B = 8$，四種材料的（價格，效益）是 $(2, 3)$、$(3, 4)$、
$(4, 5)$、$(5, 7)$。 比值由大到小是第 $1$ 種（$1.5$）、第 $4$ 種（$1.4$）、第
$2$ 種、第 $3$ 種。 總務股長買第 $1$、$4$ 種，剩 $1$ 元，效益 $10$；最大值是買
第 $2$、$4$ 種的 $11$。 輸出 \texttt{10 11}。

\textbf{想法}

題目要兩個數，各自是一個已經學過的東西。 第一個數是第 1 本第~7~章那條比值規則，照題
目規定模擬就好；第二個數是 0/1 背包，容量是班費，重量是價格。 暴力解枚舉所有
買法要 $2^{200}$ 種，不可能；背包的表是 $n \times B = 2 \times 10^7$ 格，可以，
但兩維的 \texttt{long long} 表要 $160$ MB，題目又不用輸出買了哪些，所以用一維。

模擬的部分要小心比較比值。 $e_a / p_a > e_b / p_b$ 等於
$e_a \cdot p_b > e_b \cdot p_a$（價格都是正的，兩邊同乘不會變號），改成交叉相
乘就不用浮點數，也不會有「看起來相等、其實差一點點」的問題。 乘積最大是
$10^9 \times 10^5 = 10^{14}$，要用 \texttt{long long}。

這題也把第 1 本第~7~章的結論變成數字：比值規則在可以切開的時候是對的，物品不能切開時
可能差一截。 同樣的模型出現在很多地方：預算內挑幾個專案、硬碟空間有限時挑哪
些檔案備份、一天的時間裡挑哪幾件工作做。

\textbf{程式}

\begin{example}[label=ex:treasurer]{cpp}{(treasurer.cpp) 總務股長的挑選法差多少}
#include <algorithm>
#include <iostream>
#include <vector>
using namespace std;

int main() {
    ios::sync_with_stdio(false);
    cin.tie(nullptr);
    int n, budget;
    cin >> n >> budget;
    vector<long long> price(n), gain(n);
    for (int i = 0; i < n; i++) {
        cin >> price[i] >> gain[i];
    }

    // 總務股長的方法：比值大的在前，比值相同時便宜的在前，再相同時編號小的在前
    vector<int> order(n);
    for (int i = 0; i < n; i++) {
        order[i] = i;
    }
    sort(order.begin(), order.end(), [&](int a, int b) {
        long long left = gain[a] * price[b], right = gain[b] * price[a];
        if (left != right) {
            return left > right;          // gain[a]/price[a] > gain[b]/price[b]
        }
        if (price[a] != price[b]) {
            return price[a] < price[b];
        }
        return a < b;
    });
    long long money = budget, greedy = 0;
    for (int i : order) {
        if (price[i] <= money) {          // 買不起就跳過，繼續看下一種
            money -= price[i];
            greedy += gain[i];
        }
    }

    // 0/1 背包：best[c] = 目前看過的材料裡，總價不超過 c 的最大效益
    vector<long long> best(budget + 1, 0);
    for (int i = 0; i < n; i++) {
        for (int c = budget; c >= price[i]; c--) {   // 由大到小
            best[c] = max(best[c], best[c - price[i]] + gain[i]);
        }
    }
    cout << greedy << ' ' << best[budget] << '\n';
    return 0;
}
\end{example}

\iosample{
4 8 \\
2 3 \\
3 4 \\
4 5 \\
5 7 \\
}{
10 11 \\
}

\begin{itemize}
\item \texttt{sort} 排的是編號 \texttt{order}，比較函式用 \texttt{[\&]} 抓
\texttt{price} 和 \texttt{gain} 來比（銜接書第~9~章的 lambda）。 最後比編號，
任兩種材料都分得出先後，排序的結果只有一種。
\item 價格可能大於 $B$。 這時內層迴圈的條件 \texttt{c >= price[i]} 一開始就不
成立，這種材料自然被跳過，不用另外判斷。
\item 效益總和最大大約 $200 \times 10^9 = 2 \times 10^{11}$，\texttt{greedy}
和 \texttt{best} 都要用 \texttt{long long}。
\end{itemize}

\textbf{複雜度}

排序 $O(n \log n)$，背包 $O(nB)$，$n = 200$、$B = 10^5$ 時是 $2 \times 10^7$
次更新。 最大的測試資料在本書測試的環境下約 $0.01$ 秒。

\subsection{常見錯誤}

\begin{itemize}
\item \textbf{買不起就停下來}：模擬時寫成 \texttt{if (price[i] > money)
break;}。 輸入 \texttt{3 6}、\texttt{5 10}、\texttt{4 6}、\texttt{1 1}（每個
是一行）的正確輸出是 \texttt{11 11}：買了第一種之後剩 $1$ 元，第二種買不起，
跳過，第三種買得起。 停下來的寫法印出 \texttt{10 11}。
\item \textbf{交叉相乘用 \texttt{int}}：把 \texttt{price}、\texttt{gain} 和
\texttt{left}、\texttt{right} 都改成 \texttt{int}。 輸入
\begin{code}{text}{乘積超過 int 的輸入}
2 10
7 1000000000
7 1
\end{code}
正確輸出是 \texttt{1000000000 1000000000}。 $10^9 \times 7$ 超過
\texttt{int}，在本書測試的環境下比較的結果顛倒，先買了效益 $1$ 的那一種，印出
\texttt{1 1000000000}。
\item \textbf{背包的迴圈由小到大}：範例的輸入印出 \texttt{10 12}，同一種材料買
了好幾份。 0/1 背包的一維寫法，容量一定由大到小。
\end{itemize}

\section{剛好裝滿與湊不出來的容量}
\label{sec:knap-exact}

到目前為止，$\mathrm{best}(c)$ 的意思都是「總重量\textbf{不超過} $c$」。 有些
題目要求\textbf{剛好}：廣播時段要剛好播滿、展示櫃要剛好排滿、郵資要剛好貼滿。
這時有些容量根本湊不出來，表格裡要有一個值代表「這個容量不可能」。

只要改初始值。 狀態改成「總重量剛好是 $c$ 的最大價值」，一件物品都沒有時，只
有 $c = 0$ 湊得出來（什麼都不放）：
\[
\mathrm{best}(0) = 0,\qquad \mathrm{best}(c) = \text{湊不出來}\quad (c \ge 1).
\]
轉移不變，只多一條規則：湊不出來的格子不能拿來轉移，因為「在湊不出來的組合上再
放一件」也還是湊不出來。 手機的例子改成「剛好用完空間」，四件物品處理完是：

\begin{center}
\begin{tabular}{l|ccccccccc}
\toprule
$c$ & 0 & 1 & 2 & 3 & 4 & 5 & 6 & 7 & 8 \\
\midrule
不超過 $c$（初始全部 $0$） & 0 & 0 & 3 & 4 & 5 & 7 & 8 & 10 & 11 \\
剛好 $c$（初始只有 $0$ 是 $0$） & 0 & -- & 3 & 4 & 5 & 7 & 8 & 10 & 11 \\
\bottomrule
\end{tabular}
\end{center}

「--」是湊不出來。 這四件物品的重量最小是 $2$，$1$ GB 怎樣都湊不出。 其他容
量兩列剛好一樣，但只看前三件時就不同了：前三件的「剛好 $8$」湊不出來（$2$、$3$、
$4$ 的組合沒有一種加起來是 $8$），「不超過 $8$」是 $9$。 兩種定義的關係是：
「不超過 $c$」等於「剛好 $0$、剛好 $1$、……、剛好 $c$」裡最大的那個。

「湊不出來」在程式裡有兩種常見的寫法：

\begin{itemize}
\item 價值都不是負數時，用 \texttt{-1}。 轉移前檢查
\texttt{best[c - w] != -1}，最後直接輸出，\texttt{-1} 剛好就是很多題目要的
「做不到」。 忘了檢查的話，\texttt{-1 + v} 會被當成一個真的價值。
\item 用一個很小的數 \texttt{NEG}，例如 \texttt{LLONG\_MIN / 4}。 它加上任何
可能的價值總和仍然非常小，取 \texttt{max} 時永遠輸，可以不檢查；不能直接用
\texttt{LLONG\_MIN}，減一點點就溢位。 輸出前判斷 \texttt{best[c]} 是不是還在
\texttt{NEG} 附近（例如小於 \texttt{NEG / 2}），是的話就是湊不出來。
\end{itemize}

求最小值的題目反過來，用一個很大的 \texttt{INF}，第 1 本第~9~章練習第~2~題的郵票就
是這樣。 不管哪一種，先用一句話寫清楚狀態是「不超過」還是「剛好」，初始值跟
著它走。

「不超過」的表還有一個好用的性質：$c$ 越大能放的東西只會越多，所以
$\mathrm{best}(c)$ 隨 $c$ 不會變小。 最大值第一次出現的位置，就是「要拿到最大
價值，最少要用多少容量」。 第~\ref{sec:knap-bounded}~節的例題會用到。

\section{完全背包}
\label{sec:knap-unbounded}

如果每種物品都可以拿任意多個，叫\textbf{完全背包}（unbounded knapsack）。
第~\ref{sec:knap-1d}~節那個錯誤的寫法，在這裡就是對的：容量由小到大，填
\texttt{best[c]} 時讀到的 \texttt{best[c - w]} 可能已經放了這種物品，再放一個
就是放第二個、第三個。
\[
\mathrm{best}(c) = \max\bigl(\mathrm{best}(c),\ \mathrm{best}(c - w_i) + v_i\bigr),
\quad c \text{ 由小到大}.
\]

用兩維的表想也一樣清楚。 完全背包的 $\mathrm{best}(i, c)$ 中，「放第 $i$ 種」
之後還可以再放第 $i$ 種，所以要從同一列的 $\mathrm{best}(i, c - w_i)$ 轉
移。 同一列的左邊要先算好，所以由小到大。 0/1 和完全背包的一維程式只
差在內層迴圈的方向：

\begin{center}
\begin{tabular}{lll}
\toprule
 & 放了第 $i$ 種之後 & 一維的內層迴圈 \\
\midrule
0/1 背包 & 只能看前 $i - 1$ 種（上一列） & \texttt{c} 由大到小 \\
完全背包 & 還能再放第 $i$ 種（同一列） & \texttt{c} 由小到大 \\
\bottomrule
\end{tabular}
\end{center}

第 1 本第~9~章的郵票其實就是完全背包：每種郵票張數不限，價值都是「一張」，要剛好貼滿
並讓張數最少。 第 1 本第~9~章的迴圈是金額在外層、面額在內層，這裡是物品在外層、容量在
內層。 求最大或最小值時兩種都對：算 \texttt{best[x]} 的時候，比它小的金額都已
經是最終的答案，「最後一張是哪一種」全部試過一次。 求方法數時兩種順序的答案不
同，第~\ref{sec:knap-count}~節會看到。

\subsection{例題：剛好排滿的展示櫃}
\label{sec:ex-shelf}

\textbf{題意}（\akiref{a-perfectly-filled-display-shelf}{剛好排滿的展示櫃}）

\begin{itemize}
\item 展示櫃有 $q$ 層，第 $j$ 層寬 $L_j$。 有 $n$ 種模型盒，第 $i$ 種寬 $w_i$、
展示分數 $v_i$，每種數量不限。
\item 每一層都要剛好排滿：盒子寬度加起來剛好等於這一層的寬度。 各層互不影響。
\item 對每一層輸出最大的分數總和，排不滿的層輸出 \texttt{-1}。
\item $n \le 100$，$q \le 10^4$，$w_i, L_j \le 10^4$，$v_i \le 10^6$。
\end{itemize}

例如三種盒子（寬，分數）是 $(4, 5)$、$(6, 8)$、$(9, 12)$，四層寬 $10$、$11$、
$12$、$13$。 寬 $10$ 是 $4 + 6$，$13$ 分；寬 $11$ 怎麼組都湊不出來，輸出
\texttt{-1}；寬 $12$ 有 $6 + 6$（$16$ 分）和 $4 + 4 + 4$（$15$ 分），取 $16$；
寬 $13$ 是 $4 + 9$，$17$ 分。

\textbf{想法}

一層就是一次「剛好裝滿」的完全背包，容量是這一層的寬度。 最直接的做法是每一層
各做一次：一次 $O(nL)$，最多 $100 \times 10^4 = 10^6$；$q = 10^4$ 層要
$10^{10}$，來不及。

能改進，是因為狀態的意思：$\mathrm{best}(c)$ 是「剛好排滿寬度 $c$」的最大分數，填到最
寬的那一層時，比它窄的每一個寬度都順便算好了。 所以只要做一次，容量開到所有層
裡最寬的 $\max L$，每一層直接查 \texttt{best[L[j]]}。

這題要剛好排滿，所以用第~\ref{sec:knap-exact}~節的初始值；盒子數量不限，所以
容量由小到大。 分數都是正的，「湊不出來」用 \texttt{-1}。 下面是範例的表：

\begin{center}
\begin{tabular}{l|cccccccccccccc}
\toprule
寬 $c$ & 0 & 1 & 2 & 3 & 4 & 5 & 6 & 7 & 8 & 9 & 10 & 11 & 12 & 13 \\
\midrule
\texttt{best[c]} & 0 & -1 & -1 & -1 & 5 & -1 & 8 & -1 & 10 & 12 & 13 & -1 & 16 & 17 \\
\bottomrule
\end{tabular}
\end{center}

\texttt{best[8]} $= 10$ 是兩個寬 $4$ 的盒子，同一種用了兩次；
\texttt{best[12]} 在處理寬 $6$ 的盒子時，從 \texttt{best[6]} $= 8$ 再加一個寬
$6$ 的盒子，得到 $16$，比處理寬 $4$ 時從 \texttt{best[8]} 得到的 $15$ 大。

同樣的模型也出現在別的地方：把一捲布剪成固定幾種長度的布條、一點都不能剩；用
幾種規格的木板鋪滿一段走道。

\textbf{程式}

\begin{example}[label=ex:shelf]{cpp}{(shelf.cpp) 剛好排滿的展示櫃}
#include <algorithm>
#include <iostream>
#include <vector>
using namespace std;

int main() {
    ios::sync_with_stdio(false);
    cin.tie(nullptr);
    int n, q;
    cin >> n >> q;
    vector<int> width(n);
    vector<long long> score(n);
    for (int i = 0; i < n; i++) {
        cin >> width[i] >> score[i];
    }
    vector<int> len(q);
    int maxLen = 0;
    for (int j = 0; j < q; j++) {
        cin >> len[j];
        maxLen = max(maxLen, len[j]);
    }

    const long long NONE = -1;                // 湊不出這個寬度
    vector<long long> best(maxLen + 1, NONE); // best[c] = 剛好排滿寬度 c 的最大分數
    best[0] = 0;
    for (int i = 0; i < n; i++) {
        for (int c = width[i]; c <= maxLen; c++) {   // 由小到大：同一種可以重複放
            if (best[c - width[i]] != NONE) {
                best[c] = max(best[c], best[c - width[i]] + score[i]);
            }
        }
    }
    for (int j = 0; j < q; j++) {
        cout << best[len[j]] << '\n';
    }
    return 0;
}
\end{example}

\iosample{
3 4 \\
4 5 \\
6 8 \\
9 12 \\
10 11 12 13 \\
}{
13 \\
-1 \\
16 \\
17 \\
}

\begin{itemize}
\item 盒子比最寬的層還寬時，\texttt{c} 從 \texttt{width[i]} 開始就已經超過
\texttt{maxLen}，迴圈一次都不跑。
\item \texttt{best[c]} 本來是 \texttt{-1}，第一次找到排法時，\texttt{max} 一
定選新的值，因為分數至少是 $1$。
\item 一層最多 $10^4$ 個寬 $1$ 的盒子，每個 $10^6$ 分，總和 $10^{10}$，超過
\texttt{int}。
\end{itemize}

\textbf{複雜度}

一次完全背包是 $O(n \cdot \max L) = 10^6$，加上 $q$ 次查表。 最大的測試資料在
本書測試的環境下不到 $0.01$ 秒。

\subsection{常見錯誤}

以下都用範例的輸入，正確輸出是 \texttt{13 -1 16 17}（每個一行）。

\begin{itemize}
\item \textbf{初始值全部設成 $0$}：算出的是「不超過」，寬 $11$ 的那層印出
\texttt{13}（放 $4 + 6$，留 $1$ 的空隙）。 而且永遠不會印出 \texttt{-1}。
\item \textbf{轉移前沒有檢查 \texttt{-1}}：寬 $11$ 印出 \texttt{12}，是從湊
不出來的格子加上分數得到的假答案。
\item \textbf{容量由大到小}：變成每種只能用一個，寬 $12$ 印出 \texttt{-1}，
因為 $4$、$6$、$9$ 各用一次湊不出 $12$。
\item \textbf{每一層重做一次背包}：答案都對，但最大的測試資料要做 $10^{10}$
次更新，超過時間限制。
\end{itemize}

\section{子集合和與方法數}
\label{sec:knap-count}

背包的表格不一定存「最大價值」。 換成別的東西，就能回答別的問題。

\subsection{做不做得到：子集合和}

只問「能不能剛好湊出 $c$」時，每一格存 \texttt{true} 或 \texttt{false}，叫
\textbf{子集合和}（subset sum）。 例如一箱東西要分給兩個人，重量要剛好一樣，就
是問「能不能從裡面挑出總和剛好一半的東西」。 每件東西只有一件，所以倒著填：

\begin{code}{cpp}{子集合和（0/1）}
vector<char> can(total + 1, 0);   // can[c]：挑出總和剛好 c 的方法存不存在
can[0] = 1;                       // 什麼都不挑
for (int x : a) {
    for (int c = total; c >= x; c--) {
        if (can[c - x]) {
            can[c] = 1;
        }
    }
}
\end{code}

用 \texttt{vector<char>} 而不用 \texttt{vector<bool>}，是因為後者為了省空間把
八格擠在一個位元組裡，存取比較慢，在本書測試的環境下要花大約三倍的時間。 處理完之
後，\texttt{can} 一次回答了所有的 $c$：哪些總和湊得出來、最接近一半的是多少、
最大的湊不出來的數是多少，都是掃一遍 \texttt{can} 的事。

\subsection{有幾種方法}

把 \texttt{max} 換成加法，每一格存「剛好湊出 $c$ 的方法數」：
\[
\mathrm{ways}(c) \mathrel{+}= \mathrm{ways}(c - w_i),\qquad \mathrm{ways}(0) = 1.
\]
$\mathrm{ways}(0) = 1$ 是「什麼都不放」這一種方法，和第 1 本第~9~章計數題的起點一樣。
0/1 還是完全，一樣看迴圈方向。

計數題要多想一件事：怎樣算「不同的方法」。 例子：集點卡有 $10$ 點，可以換三種
贈品，分別要 $2$、$3$、$5$ 點，每種想換幾個都可以，點數要剛好用完。

\begin{itemize}
\item 只看每種換幾個：$5 + 5$、$5 + 3 + 2$、$3 + 3 + 2 + 2$、
$2 + 2 + 2 + 2 + 2$，共 $4$ 種。 這叫\textbf{組合}。
\item 一次換一個，換的先後不同就算不同（$2 + 3 + 5$ 和 $5 + 3 + 2$ 是兩種）：
共 $14$ 種。 這是第 1 本第~9~章跳格子的算法，叫\textbf{排列}。
\end{itemize}

兩種答案用的是同一個轉移，只差在兩層迴圈誰在外面：

\begin{example}[label=ex:count-ways]{cpp}{(count\_ways.cpp) 組合數和排列數}
#include <iostream>
#include <vector>
using namespace std;

int main() {
    const long long MOD = 1000000007;
    int n, target;
    cin >> n >> target;
    vector<int> cost(n);
    for (int &x : cost) {
        cin >> x;
    }

    // 只看每種換幾個：物品在外層
    vector<long long> combo(target + 1, 0);
    combo[0] = 1;
    for (int i = 0; i < n; i++) {
        for (int c = cost[i]; c <= target; c++) {
            combo[c] = (combo[c] + combo[c - cost[i]]) % MOD;
        }
    }

    // 先後不同算不同：點數在外層
    vector<long long> seq(target + 1, 0);
    seq[0] = 1;
    for (int c = 1; c <= target; c++) {
        for (int i = 0; i < n; i++) {
            if (cost[i] <= c) {
                seq[c] = (seq[c] + seq[c - cost[i]]) % MOD;
            }
        }
    }
    cout << combo[target] << ' ' << seq[target] << '\n';
    return 0;
}
\end{example}

\iosample{
3 10 \\
2 3 5 \\
}{
4 14 \\
}

物品在外層時，下表是處理完每一種贈品後的 \texttt{combo}：

\begin{center}
\begin{tabular}{l|ccccccccccc}
\toprule
點數 $c$ & 0 & 1 & 2 & 3 & 4 & 5 & 6 & 7 & 8 & 9 & 10 \\
\midrule
只有 $2$ 點的 & 1 & 0 & 1 & 0 & 1 & 0 & 1 & 0 & 1 & 0 & 1 \\
加上 $3$ 點的 & 1 & 0 & 1 & 1 & 1 & 1 & 2 & 1 & 2 & 2 & 2 \\
加上 $5$ 點的 & 1 & 0 & 1 & 1 & 1 & 2 & 2 & 2 & 3 & 3 & 4 \\
\bottomrule
\end{tabular}
\end{center}

為什麼物品在外層就不會重複？ 處理第 $i$ 種贈品時，表格裡的每一種方法，都是
「先決定前 $i - 1$ 種各換幾個，再決定第 $i$ 種換幾個」。 每一種組合只會用這個
固定的順序被數到一次：$5 + 3 + 2$ 只會被當成「$2$ 點的一個、$3$ 點的一個、$5$
點的一個」。 點數在外層時，算 $\mathrm{seq}(c)$ 是在問「最後一個換的是哪一
種」，每一種都試，$2 + 3 + 5$ 和 $5 + 3 + 2$ 的最後一個不同，就被分開數。

第 1 本第~9~章練習第~3~題的郵票也是這樣：$1$、$5$、$6$ 元貼滿 $12$ 元，依序貼的貼法
有 $29$ 種；把同一支程式的輸入改成 \texttt{3 12} 和 \texttt{1 5 6}，第一個數
是 $6$，就是不管先後的貼法數。

\begin{apcsnote}
實作題由電腦自動評分，只看輸出。 計數題的範例通常很小，組合數和排列數有時剛好
一樣（例如只有一種物品，或目標小到只有一種湊法），範例過了不代表數對了東西。
寫之前先照題目對「不同」的規定，手算一組有兩三種物品的小例子，再決定哪一層迴
圈在外面。 答案要取餘數時，每做一次加法就取一次。
\end{apcsnote}

\section{多重背包：二進位拆分}
\label{sec:knap-bounded}

第三種情況介在中間：第 $i$ 種物品有 $k_i$ 個，最多拿 $k_i$ 個，叫\textbf{多重
背包}（bounded knapsack）。

最直接的做法是把它當成 $k_i$ 件一模一樣的物品，做 0/1 背包。 這是對的，但慢：
物品總數變成 $\sum k_i$，時間是 $O(W \sum k_i)$。 $n = 200$、$k_i = 1000$、
$W = 10^5$ 時是 $2 \times 10^{10}$。

\textbf{二進位拆分}（binary splitting）的想法：把 $k$ 個拆成幾包，每包當成一件
0/1 物品，只要「挑某幾包」能湊出 $0$ 到 $k$ 的每一種個數就夠了。 包的大小取
$1, 2, 4, \ldots$，最後一包放剩下的。 例如 $k = 13$：
\[
13 = 1 + 2 + 4 + 6.
\]
$0$ 到 $7$ 用前三包就湊得出來（就是二進位）；$8$ 到 $13$ 先拿 $6$ 那一包，剩
下的 $2$ 到 $7$ 再用前三包。 而四包全拿是 $13$，不會超過庫存。 一種物品原本
要當成 $13$ 件，現在只要 $4$ 件；$k = 1000$ 只要 $10$ 件。 每一包的重量和價值
是單個的倍數，例如 $4$ 個一包，重量 $4w$、價值 $4v$。

\begin{whybox}[拆出來的包剛好湊得出 $0$ 到 $k$]
照程式的拆法，包的大小是 $1, 2, \ldots, 2^{m-1}$，加上最後一包 $r$，其中
$1 \le r \le 2^m$，而且 $1 + 2 + \cdots + 2^{m-1} + r = (2^m - 1) + r = k$。

\begin{itemize}
\item 湊得出來：$0 \le x \le 2^m - 1$ 時，$x$ 的二進位寫法就告訴我們拿前 $m$ 包
的哪幾包。 $2^m \le x \le k$ 時，先拿 $r$ 那一包，剩下 $x - r$。 因為
$x \le k$，$x - r \le k - r = 2^m - 1$；因為 $x \ge 2^m \ge r$，$x - r \ge 0$。
所以剩下的也能用前 $m$ 包湊出。
\item 不會超過：所有包加起來剛好 $k$，任挑幾包都不超過 $k$。
\end{itemize}

0/1 背包會考慮「每包拿或不拿」的所有組合，等於考慮了 $0$ 到 $k$ 的每一種個數，
不多也不少，所以答案和多重背包一樣。
\end{whybox}

\subsection{例題：有限庫存的兌換獎品}
\label{sec:ex-prizes}

\textbf{題意}（\akiref{prizes-in-limited-stock}{有限庫存的兌換獎品}）

\begin{itemize}
\item 手上有 $T$ 點。 有 $n$ 種獎品，第 $i$ 種每個 $p_i$ 點、滿意度 $v_i$，只剩
$k_i$ 個。
\item 每種換 $0$ 到 $k_i$ 個，花的點數總和不超過 $T$。
\item 先讓總滿意度最大；最大值有好幾種換法時，取花最少點數的。 輸出最大總滿意
度，以及達到它最少要花幾點。
\item $n \le 200$，$T \le 10^5$，$v_i \le 10^5$，$k_i \le 1000$。
\end{itemize}

例如 $T = 10$，鑰匙圈每個 $2$ 點、滿意度 $5$，剩 $3$ 個；便條本 $3$ 點、滿意
度 $6$，剩 $1$ 個；抱枕 $5$ 點、滿意度 $10$，剩 $1$ 個。 換三個鑰匙圈加便條本，
花 $9$ 點、滿意度 $21$；換一個鑰匙圈、便條本和抱枕，花 $10$ 點，也是 $21$。
沒有更好的換法，輸出 \texttt{21 9}。 如果忘了庫存，鑰匙圈換五個就有 $25$。

\textbf{想法}

暴力解是枚舉每種換幾個，$1001^{200}$ 種，不可能。 每種獎品當成 $k_i$ 件物品做
0/1 背包，前面算過要 $2 \times 10^{10}$ 次，太慢。 用二進位拆分，每種最多 $10$
包，$200$ 種是 $2000$ 包，0/1 背包做 $2000 \times 10^5 = 2 \times 10^8$ 次簡單
的更新，可以。

第二個數用第~\ref{sec:knap-exact}~節最後的性質。 \texttt{best[c]} 是「花不超
過 $c$ 點」的最大滿意度，隨 $c$ 不減。 設 $c^*$ 是 \texttt{best[c]} 第一次
等於 \texttt{best[T]} 的位置：

\begin{itemize}
\item 有一種換法花不超過 $c^*$ 點就拿到最大值，所以最少的花費 $\le c^*$。
\item 如果某種換法只花 $c' < c^*$ 點就拿到最大值，\texttt{best[c']} 就至少是
最大值，第一次出現的位置會在 $c'$ 或更前面，和 $c^*$ 的定義矛盾。
\end{itemize}

所以 $c^*$ 就是答案。 範例的 \texttt{best[9]} $=$ \texttt{best[10]} $= 21$，
\texttt{best[8]} $= 16$，$c^* = 9$。

同樣的模型也出現在別的地方：倉庫裡每種貨只剩幾箱，貨車載重有限；餐廳的每道菜
限量供應，預算內點菜。

\textbf{程式}

\begin{example}[label=ex:prizes]{cpp}{(prizes.cpp) 有限庫存的兌換獎品}
#include <algorithm>
#include <iostream>
#include <vector>
using namespace std;

int main() {
    ios::sync_with_stdio(false);
    cin.tie(nullptr);
    int n, points;
    cin >> n >> points;
    // best[c] = 花不超過 c 點的最大總滿意度
    vector<long long> best(points + 1, 0);
    for (int i = 0; i < n; i++) {
        long long cost, joy;
        int stock;
        cin >> cost >> joy >> stock;
        // 把 stock 個拆成 1, 2, 4, ... 和剩下的，每一包當成一件 0/1 物品
        for (int size = 1; stock > 0; size *= 2) {
            int take = min(size, stock);
            stock -= take;
            long long w = take * cost, v = take * joy;
            for (long long c = points; c >= w; c--) {
                best[c] = max(best[c], best[c - w] + v);
            }
        }
    }
    int least = 0;                          // 最大值第一次出現的位置
    while (best[least] != best[points]) {
        least++;
    }
    cout << best[points] << ' ' << least << '\n';
    return 0;
}
\end{example}

\iosample{
3 10 \\
2 5 3 \\
3 6 1 \\
5 10 1 \\
}{
21 9 \\
}

\begin{itemize}
\item 拆包的迴圈裡，\texttt{size} 每次加倍，\texttt{take} 是這一包實際的個數。
最後一輪 \texttt{stock} 比 \texttt{size} 小，\texttt{take} 就是剩下的。 例如
$3$ 個拆成 $1$、$2$，$13$ 個拆成 $1$、$2$、$4$、$6$。
\item 一包的重量 \texttt{w} 可能超過 \texttt{points}，內層迴圈不跑，這一包等於
不存在，不影響答案：拿得到的個數本來就受點數限制。
\item 總滿意度最大約 $200 \times 1000 \times 10^5 = 2 \times 10^{10}$，要用
\texttt{long long}。
\end{itemize}

\textbf{複雜度}

包的總數不超過 $\sum \lceil \log_2(k_i + 1) \rceil \le 200 \times 10$，每包掃一次
容量，共約 $2 \times 10^8$ 次。 最大的測試資料在本書測試的環境下約 $0.07$ 秒，
這題的時間限制是 $2$ 秒。

\subsection{常見錯誤}

\begin{itemize}
\item \textbf{忽略庫存}：當成完全背包，範例的輸入印出 \texttt{25 10}，換了五
個鑰匙圈。
\item \textbf{拆包時漏掉剩下的，或拆過頭}：$13$ 個只拆成 $1$、$2$、$4$，最多
只能拿到 $7$ 個；拆成 $1$、$2$、$4$、$8$，可以拿到 $15$ 個，超過庫存。 範例
的庫存 $3$、$1$ 剛好拆得乾淨，看不出來。
\item \textbf{第二個數直接輸出 $T$}：範例印出 \texttt{21 10}。 最佳的換法不一
定把點數用完。
\item \textbf{每一包的迴圈由小到大}：一包被拿了好幾次，範例同樣印出
\texttt{25 10}。 拆出來的包是 0/1 物品，一定由大到小。
\end{itemize}

\section{背包的變形}
\label{sec:knap-variants}

很多題目讀起來不像背包，寫出狀態才發現是。 下面幾種變形都只是在前面的表格上
改一兩個地方。 認題目的時候問自己三件事：哪個量有上限（它就是容量）、每樣東西
能拿幾次（0/1、完全、多重）、要「不超過」還是「剛好」。

\textbf{每組恰好選一個。} 東西分成好幾組，每組一定要選一個、也只能選一個，例
如每門課從開出的幾個時段裡選一個。 每處理一組，就開一個新的陣列，初始全部是「湊
不出來」，只從「選了這組的某一個」轉移過來：

\begin{code}{cpp}{每組恰好選一個}
vector<long long> next(cap + 1, NONE);         // 這一組一定要選一個
for (auto [w, v] : group) {
    for (int c = w; c <= cap; c++) {
        if (best[c - w] != NONE) {
            next[c] = max(next[c], best[c - w] + v);
        }
    }
}
best = next;
\end{code}

\texttt{next} 只從舊的 \texttt{best} 轉移，所以同一組不會選兩個（迴圈方向也就
不重要了）；\texttt{next} 不從「不選」轉移，所以每組一定選一個。 如果規則是
「每組最多選一個」，\texttt{next} 改成從 \texttt{best} 複製一份再更新。

\textbf{容量會是負的。} 例如每個動作讓室溫上升或下降，最後要落在某個範
圍。 把所有可能的總和加上一個固定的偏移量 \texttt{OFF}，平移到從 $0$ 開始的索
引。 0/1 的掃描方向要看轉移從哪邊來：$w > 0$ 時從左邊來，由大到小；$w < 0$ 時
從右邊來，由小到大。 原則都一樣：讀到的格子必須還是舊值。

\textbf{兩種限制同時存在。} 例如同時限制重量和體積，狀態變成
$\mathrm{best}(a, b)$，表格多一維，0/1 時兩層容量迴圈都由大到小。 時間和空間
都乘上第二個容量，先估一下放不放得下。

\textbf{需要「至少」。} 例如要買到至少 $X$ 度電、總價最低。 超過 $X$ 的部分都
一樣好，所以容量只開到 $X$，所有超過的都算在第 $X$ 格：完全背包可以寫成
$\mathrm{cost}(c) = \min_i \bigl(\mathrm{cost}(\max(0, c - a_i)) + p_i\bigr)$，
$c$ 由小到大。 如果不截在 $X$，一包裡的數量很大時，表格要開到很遠。

\textbf{容量太大。} 容量到 $10^9$，表格開不出來。 先看別的量是不是很小：
如果價值總和很小，就把價值當成表格的一維，存「拿到這麼多價值最少要多少重量」，
最後找重量不超過 $W$ 的最大價值。 如果物品只有二十件左右，用第 1 本第~4~章的子集
合枚舉。 兩個都很大，通常是題目還有別的性質沒看出來。

\subsection{什麼時候用背包}

\begin{itemize}
\item 有一個「總和不能超過」或「總和要剛好」的量（重量、預算、時間、點數），每
樣東西各佔一些，要挑一些出來。
\item 這個量的上限不大，大約 $10^5$ 以內，乘上物品數在 $10^8$ 左右以內。
\item 東西不能切開。 能切開時用第 1 本第~7~章練習第~2(b)~題的比值貪心，快很多。
\end{itemize}

看起來像、但不能直接用的情況：

\begin{itemize}
\item 容量和價值都到 $10^9$、物品又多：狀態開不出來，要找題目的其他性質，或者
用貪心並證明它。
\item 物品的價值會受到別的物品影響，例如「同時買 A 和 B 打八折」。 單純的背包假
設每樣東西的價值固定；要先想辦法把會互相影響的東西綁成一組，變成「每組選一
個」。
\end{itemize}

\section{本章常見錯誤}
\label{sec:knap-pitfalls}

各節的錯誤，歸納起來是四類。 寫背包題時，交出去之前逐項檢查。

\begin{center}
\begin{tabular}{p{0.2\linewidth}p{0.46\linewidth}>{\raggedright\arraybackslash}p{0.22\linewidth}}
\toprule
錯誤 & 例子 & 見 \\
\midrule
迴圈方向錯 & 0/1 背包或拆出來的包由小到大，同一件被拿好幾次；完全背包由大到
小，每種只能拿一個 & 第~\ref{sec:knap-1d}、\ref{sec:ex-shelf}、
\ref{sec:ex-prizes}~節 \\
初始值錯 & 要「剛好」卻全部設成 $0$；湊不出來的 \texttt{-1} 拿去轉移 &
第~\ref{sec:knap-exact}、\ref{sec:ex-shelf}~節 \\
數錯東西 & 數組合時把容量放在外層；忽略庫存；第二個答案直接輸出容量 &
第~\ref{sec:knap-count}、\ref{sec:ex-prizes}~節 \\
該重用卻重算 & 每個詢問重做一次背包；多重背包展開成 $k$ 件 &
第~\ref{sec:ex-shelf}、\ref{sec:knap-bounded}~節 \\
\bottomrule
\end{tabular}
\end{center}

另外，背包的價值總和很容易超過 \texttt{int}：動手前先算「物品數 $\times$ 最大
價值 $\times$ 最多拿幾個」。 還有一個只有小輸入抓得到的錯：範例的答案常常剛好
不受迴圈方向影響，寫完用兩三件物品手算一次，或用第 1 本第~4~章的子集合枚舉對拍。

\section{本章練習}
\label{sec:knap-exercises}

\begin{drills}[label=drills:knap]
\item 四件物品（重量，價值）依序是 $(3, 5)$、$(1, 2)$、$(4, 7)$、$(2, 3)$，
容量 $6$。
\begin{enumerate}
\item 用一維陣列、容量由大到小，寫出處理完每一件之後的 \texttt{best[0..6]}。
\item 如果容量寫成由小到大，最後的 \texttt{best[6]} 是多少？ 它對應哪一種「拿
法」？
\end{enumerate}

\item 某人要算「$2$、$3$、$5$ 點的贈品換剛好 $10$ 點，只看每種換幾個的方法數」，
寫成：
\begin{code}{cpp}{練習 2}
ways[0] = 1;
for (int c = 1; c <= 10; c++) {
    for (int x : {2, 3, 5}) {
        if (x <= c) {
            ways[c] += ways[c - x];
        }
    }
}
\end{code}
它會印出多少？ 錯在哪裡？ 只看 \texttt{ways[5]}，多數了哪些？

\item 用二進位拆分。
\begin{enumerate}
\item $k = 100$ 要拆成哪幾包？
\item 要拿 $77$ 個時，拿哪幾包？
\item 有人把 $100$ 拆成 $1, 2, 4, 8, 16, 32, 64$。 錯在哪裡？
\end{enumerate}

\item 估計下列情況的時間或記憶體。
\begin{enumerate}
\item 範例~\ref{ex:treasurer} 如果改用兩維的 \texttt{long long} 表，大約要多少
記憶體？
\item 範例~\ref{ex:prizes} 如果把每種獎品展開成 $k_i$ 件 0/1 物品，最壞情況要
幾次更新？ 用二進位拆分呢？
\end{enumerate}

\item 0/1 背包要求剛好裝滿，有人用 \texttt{-1} 代表湊不出來，寫成
\begin{code}{cpp}{練習 5}
vector<long long> best(cap + 1, -1);
best[0] = 0;
for (int i = 0; i < n; i++) {
    for (int c = cap; c >= w[i]; c--) {
        best[c] = max(best[c], best[c - w[i]] + v[i]);
    }
}
\end{code}
物品是 $(3, 4)$、$(2, 1)$，容量 $4$。 \texttt{best[4]} 應該是多少？ 這段程式算
出多少？ 怎麼改？

\item 完全背包求最大價值。 有人把兩層迴圈對調，容量在外層由小到大、物品在內層：
\begin{code}{cpp}{練習 6}
for (int c = 1; c <= cap; c++) {
    for (int i = 0; i < n; i++) {
        if (w[i] <= c) {
            best[c] = max(best[c], best[c - w[i]] + v[i]);
        }
    }
}
\end{code}
答案還對嗎？ 如果改成 0/1 背包，能不能只把外層改成由大到小？ 各說明理由。
\end{drills}

\section{實作練習}
\label{sec:knap-practice}

下面是 Aki-Judge 上和本章相關的題目，依小節分組。 難度分四階：暖身是直接套用
本章的寫法；基本要自己認出容量和物品，或處理一個細節；實戰的題敘比較長，要先把
題目轉換成本章的形狀；挑戰會和其他章的技巧合在一起。

\begin{center}
\begin{tabular}{llc}
\toprule
小節 & 題目 & 難度 \\
\midrule
0/1 背包與剛好裝滿 & \akiref{filling-the-lunch-broadcast-exactly}{剛好填滿的午休廣播時段} & 暖身 \\
 & \akiref{fewest-stamps-for-exact-postage}{剛好貼滿郵資的最少郵票張數} & 暖身 \\
 & \akiref{how-big-a-backpack-is-enough}{要多大的背包才夠} & 基本 \\
 & \akiref{club-shopping-with-one-half-price-coupon}{一張半價券的福利社採購} & 實戰 \\
 & \akiref{how-many-best-packings}{最佳裝法有幾種} & 實戰 \\
 & \akiref{gear-that-must-come-along}{哪些裝備一定得帶} & 挑戰 \\
 & \akiref{the-cheapest-one-is-free}{最便宜那件免費} & 挑戰 \\
\midrule
完全背包 & \akiref{sticker-counts-you-cannot-buy}{買不到的貼紙張數} & 暖身 \\
 & \akiref{the-no-change-pay-station}{不找零的繳費機} & 基本 \\
 & \akiref{listing-the-fewest-stamps}{列出最少郵票的貼法} & 基本 \\
 & \akiref{cheer-routine-without-repeats}{不重複招式的啦啦隊編排} & 實戰 \\
 & \akiref{cutting-a-plank-with-saw-kerf}{考慮鋸縫的木條裁切} & 實戰 \\
 & \akiref{supply-bundles-for-freshman-orientation}{迎新物資的補給包} & 挑戰 \\
\midrule
子集合和 & \akiref{two-luggage-holds}{兩個行李艙裝得下嗎} & 基本 \\
 & \akiref{how-light-can-the-heaviest-box-be}{最重那箱可以多輕} & 實戰 \\
 & \akiref{fewest-moves-to-balance-two-teams}{最少移動幾人讓兩隊力量相等} & 實戰 \\
 & \akiref{rivals-and-partners-in-the-tug-of-war}{有恩怨也有搭檔的分組} & 挑戰 \\
\midrule
變形 & \akiref{one-plan-per-club}{每個社團恰選一個方案} & 暖身 \\
 & \akiref{camp-lunches-with-one-free-voucher}{營隊午餐與一張兌換券} & 基本 \\
 & \akiref{toppings-within-the-sweetness-range}{甜度剛好落在區間的配料} & 實戰 \\
\bottomrule
\end{tabular}
\end{center}

本章的三題例題也在 Aki-Judge 上，讀完各節之後可以先自己寫一次再對照。

\section{讀完你應該會}
\label{sec:knap-checklist}

\begin{itemize}
\item 看到「容量有限、挑一些東西」的題目，能說出容量是哪個量、每樣東西能拿幾
次，並估計 $n \times W$ 來不來得及。
\item 能寫出 0/1 背包的兩維表格，並從表格往回找出選了哪幾件。
\item 能寫出 0/1 和完全背包的一維寫法，並用一件物品的小例子說明迴圈方向為什
麼不能反。
\item 能依「不超過」或「剛好」設定初始值，用 \texttt{-1} 或 \texttt{NEG} 標出
湊不出來的容量，並在轉移時避開它們。
\item 數方法數時，能依題目對「不同」的定義，決定物品和容量哪一層在外面；遇到
有數量上限的物品，能用二進位拆分改成 0/1 背包。
\end{itemize}

% 頁面剩不到五行時先換頁，避免「延伸閱讀」單獨留在頁尾（效果同 needspace）。
\par\begingroup\dimen0=5\baselineskip \dimen2=\pagegoal \advance\dimen2 by -\pagetotal
\ifdim\dimen0>\dimen2 \ifdim\dimen2>0pt \vfil\fi \break\fi\endgroup
\textbf{延伸閱讀}（書末會整理完整的參考資料）

\begin{itemize}
\item Antti Laaksonen，\textit{Competitive Programmer's Handbook}，第 7 章的
knapsack 一節：用子集合和介紹背包，篇幅很短。
\item 劉汝佳，《算法競賽入門經典》第 2 版，9.3.2 節：0/1 背包和滾動陣列，並
討論壓成一維之後為什麼不好印出選法。
\item Steven Halim 等，\textit{Competitive Programming 4}，Book 1，3.5.2 節：
0/1 背包和找零錢的記憶化寫法，附有大量分類好的練習題。
\end{itemize}

\addcontentsline{toc}{section}{練習解答}
\section*{練習解答}

\begin{enumerate}
\item
\begin{enumerate}
\item 一開始全是 $0$。
\begin{center}
\begin{tabular}{l|ccccccc}
\toprule
$c$ & 0 & 1 & 2 & 3 & 4 & 5 & 6 \\
\midrule
$(3, 5)$ 之後 & 0 & 0 & 0 & 5 & 5 & 5 & 5 \\
$(1, 2)$ 之後 & 0 & 2 & 2 & 5 & 7 & 7 & 7 \\
$(4, 7)$ 之後 & 0 & 2 & 2 & 5 & 7 & 9 & 9 \\
$(2, 3)$ 之後 & 0 & 2 & 3 & 5 & 7 & 9 & 10 \\
\bottomrule
\end{tabular}
\end{center}
答案 $10$：$(4, 7)$ 加 $(2, 3)$，或 $(3, 5)$、$(1, 2)$、$(2, 3)$ 三件，都是
重量 $6$、價值 $10$。
\item $12$，等於六份重量 $1$、價值 $2$ 的物品。 由小到大時每件都可以重複拿，
算出的是完全背包的答案；在完全背包裡，價值除以重量最高的是 $(1, 2)$，裝滿六份
最好。
\end{enumerate}

\item 印出 $14$。 容量在外層、物品在內層，算的是「最後一個換的是哪一種」，先後
不同的換法被分開數，得到的是排列數；組合數是 $4$。 \texttt{ways[5]} 會是 $3$：
$5$、$2 + 3$、$3 + 2$，其中 $2 + 3$ 和 $3 + 2$ 是同一種組合，多數了一次。 把兩
層迴圈對調，物品在外層，就是範例~\ref{ex:count-ways} 的 \texttt{combo}。

\item
\begin{enumerate}
\item $1, 2, 4, 8, 16, 32, 37$。 前六包加起來是 $63$，剩下 $100 - 63 = 37$。
\item $77 \ge 64$，先拿 $37$ 那一包，剩下 $40 = 32 + 8$。 拿 $37$、$32$、$8$
三包。
\item 七包加起來是 $127$，可以拿到 $101$ 到 $127$ 個，超過庫存。 最後一包要放
剩下的數量，不能繼續加倍。
\end{enumerate}

\item
\begin{enumerate}
\item $(200 + 1) \times (10^5 + 1)$ 格，每格 $8$ 位元組，大約 $160$ MB。 一維只
要 $800$ KB。
\item 展開後物品數是 $\sum k_i \le 200 \times 1000 = 2 \times 10^5$，每件掃一次
$10^5$ 的容量，最壞 $2 \times 10^{10}$ 次。 二進位拆分後每種最多 $10$ 包
（$1 + 2 + \cdots + 256 = 511$，第十包是剩下的 $489$），共 $2000$ 包，
$2 \times 10^8$ 次。
\end{enumerate}

\item 應該是 $-1$：$3 + 2 = 5$，兩件都湊不出 $4$。 這段程式算出 $3$。 處理
$(3, 4)$ 時 \texttt{best[3] = 4}，\texttt{best[4]} 從 \texttt{best[1] + 4}
得到 $-1 + 4 = 3$；處理 $(2, 1)$ 時 \texttt{best[2]} 還是 $-1$，\texttt{best[4]} 從
\texttt{best[2] + 1} $= 0$ 和自己的 $3$ 比，留下 $3$。 這個 $3$ 是從湊不出來的 \texttt{best[1]}
來的假答案。 轉移前加上 \texttt{if (best[c - w[i]] != -1)}。

\item 第一個問題：還是對的。 算 \texttt{best[c]} 時，比 $c$ 小的格子都已經考慮
過所有物品、是最終的答案，內層對每一種「最後放的是哪一種」都試一次，正好是完全
背包的轉移。 求最大值時同一種放法被看到幾次都沒關係，所以和物品在外層的答案一
樣；求方法數時才會多數（練習第~2~題）。

第二個問題：不行。 外層由大到小時，\texttt{best[c - w[i]]} 還沒算到，存的都是
初始的 $0$，算出的只是「放一件」的最好結果。 例如物品 $(1, 1)$、$(1, 1)$、容量
$2$，正確答案是 $2$，這樣寫得到 $1$。 0/1 背包的一維寫法需要「每處理完一件，
整列前進一次」，物品一定要在外層。
\end{enumerate}

\chapter{子序列 DP}
\label{ch:subseq}

週末你和朋友各自排了逛夜市的順序。 你想照這個順序吃：雞排、珍奶、臭豆腐、
蚵仔煎、烤玉米、地瓜球、糖葫蘆。 朋友的順序是：臭豆腐、雞排、蚵仔煎、珍奶、
地瓜球、烤玉米、糖葫蘆。 兩個人想一起逛一段，一起吃的攤位要同時符合兩人各自
的先後順序；自己清單上的攤位可以跳過，跳過的那幾攤之後再自己去。 最多能一起
吃幾攤？ 往下讀之前先試一試：先挑雞排的話，接下來還能一起吃哪幾攤？

從一串東西裡刪掉一些（也可以都不刪），剩下的保持原來的先後，叫做原來那串的
\textbf{子序列}（subsequence）。 子序列可以跳著挑；如果規定一定要挑相鄰的一
段，那叫\textbf{子字串}（substring），是《C++ 演算法解題 1：APCS 實作題核心》（以下簡稱第 1 本）第~5~章前綴和、滑動視窗處理的對象。
上面的問題是在找兩份清單共同的子序列裡最長的一個。 這一章處理三種跟子序列有
關的問題：

\begin{itemize}
\item 一串數字裡，最長的遞增子序列。 長度怎麼算，銜接書已經教過；這一章補上
「把子序列本身印出來」和「每個位置各自的答案」。
\item 兩串東西的共同子序列，以及把一串改成另一串最少要改幾次。
\item 一串東西和它自己倒過來的那一串：最長的回文子序列。
\end{itemize}

三題例題分別練最長遞增子序列的延伸、把共同子序列的「相等」換成別的條件並且
只留兩列表格，以及在一長串文字裡找最像某句話的一段。

算上第 1 本的兩章，這一章是動態規劃的第四章。 第 1 本第~9~章的四件事（狀態、轉移、初始值、計算順序）、
第 1 本第~10~章的二維表格和滾動陣列都直接拿來用，不再重講。 最長遞增子序列的兩種寫
法在《C++ 程式設計：給學過 C 的人》（以下簡稱銜接書）第~10~章 10.3 節，本章
只用一個框回顧；\texttt{lower\_bound} 在銜接書第~9~章。 沒讀過銜接書的讀者，
只要知道 \texttt{lower\_bound(v.begin(), v.end(), x)} 會在排好的
\texttt{v} 裡找到第一個 $\ge x$ 的位置，就能讀這一章。

\section{子序列問題的兩種狀態}
\label{sec:ss-states}

先想最直接的做法。 長度 $n$ 的一串，每個元素可以選或不選，一共有 $2^n$ 個子
序列。 要找最長的遞增子序列，可以一個一個檢查；要找兩串的最長共同子序列，可
以把第一串的每個子序列拿去第二串裡找找看。 $n = 20$ 時大約一百萬個，還來得及；
$n = 40$ 就超過一兆。 這一章的題目，$n$ 常常到幾千甚至 $2 \times 10^5$，一定
要換成動態規劃。 不過小範圍的暴力解仍然要寫，用來和動態規劃的程式比對。

子序列的動態規劃，狀態幾乎都是下面兩種之一：

\begin{center}
\begin{tabular}{p{0.2\linewidth}p{0.42\linewidth}p{0.28\linewidth}}
\toprule
題目的形狀 & 狀態 & 例子 \\
\midrule
一串，挑出的子序列要滿足某種前後關係 & \texttt{f[i]}：最後一個挑的是第 $i$
個時，最好能多好 & 最長遞增子序列 \\
兩串，兩邊各挑一個子序列互相對應 & \texttt{f[i][j]}：只看第一串的前 $i$ 個和
第二串的前 $j$ 個時，最好能多好 & 最長共同子序列、編輯距離、回文子序列 \\
\bottomrule
\end{tabular}
\end{center}

兩種狀態都在問「最後一步」：第一種問最後一個挑的是誰、它前面接的是誰；第二種
看兩串各自的最後一個元素，它們配在一起，還是其中一個沒用到。 這一章的每個轉移
式都是這樣想出來的。

\section{最長遞增子序列：從長度到序列本身}
\label{sec:ss-lis}

\begin{bridgerecap}[銜接書 10.3 節：最長遞增子序列的兩種寫法]
\textbf{最長遞增子序列}（longest increasing subsequence，LIS）：挑出的數字保持
原來的先後，而且嚴格遞增，最長能挑幾個。
\begin{itemize}
\item $O(n^2)$：\texttt{dp[i]} 是以 \texttt{a[i]} 結尾的最長長度，等於所有
$j < i$、\texttt{a[j] < a[i]} 的 \texttt{dp[j] + 1} 裡最大的，沒有就是 $1$。
答案是所有 \texttt{dp[i]} 的最大值。
\item $O(n \log n)$：\texttt{tail[k]} 是「長度 $k + 1$ 的遞增子序列，結尾最小
可以是多少」。 \texttt{tail} 一定遞增；每讀一個 $x$，用
\texttt{lower\_bound} 找第一個 $\ge x$ 的結尾，找到就換成 $x$，找不到就把 $x$
接在最後。 \texttt{tail} 的長度就是答案。
\item 允許相等（不下降）時，把 \texttt{lower\_bound} 換成 \texttt{upper\_bound}。
\item \texttt{tail} 只保證長度對，它本身不一定是一個子序列。
\end{itemize}
\end{bridgerecap}

這一節補上兩件事：把一個最長的子序列真的印出來，以及算出「以每個位置結尾」的
最長長度。 第二件事是這一節例題的關鍵。

\subsection{印出一個最長的}
\label{sec:ss-lis-path}

$O(n^2)$ 的寫法很好改：算 \texttt{dp[i]} 時，記下是哪一個 $j$ 給出最大值，存
在 \texttt{from[i]}。 最後從 \texttt{dp} 最大的那個位置出發，沿著
\texttt{from} 往回走，走到 $-1$ 為止，經過的數字倒過來就是答案。

\begin{code}{cpp}{O(n²) 的寫法多記一個 from}
vector<int> dp(n, 1), from(n, -1);
for (int i = 0; i < n; i++) {
    for (int j = 0; j < i; j++) {
        if (a[j] < a[i] && dp[j] + 1 > dp[i]) {
            dp[i] = dp[j] + 1;
            from[i] = j;        // a[i] 接在 a[j] 後面
        }
    }
}
\end{code}

$O(n \log n)$ 的寫法也可以記 \texttt{from}，只是要多記一樣東西。
\texttt{a[i]} 放進 \texttt{tail[k]} 時，它的意思是「\texttt{a[i]} 可以接在某
個長度 $k$ 的遞增子序列後面」。 那個長度 $k$ 的子序列，結尾就是現在的
\texttt{tail[k - 1]}；但 \texttt{tail} 只存了值，不知道它在哪個位置。 所以另
外開一個 \texttt{tailAt}，和 \texttt{tail} 同步更新，記下每個結尾在原數列的
位置。

\begin{example}[label=ex:ss-lis-path]{cpp}{(lis\_path.cpp) 印出一個最長遞增子序列}
#include <algorithm>
#include <iostream>
#include <vector>
using namespace std;

int main() {
    int n;
    cin >> n;
    vector<int> a(n);
    for (int &x : a) {
        cin >> x;
    }
    vector<int> tail;        // tail[k]：長度 k + 1 的最小結尾（值）
    vector<int> tailAt;      // tailAt[k]：這個結尾在 a 的哪個位置
    vector<int> from(n, -1); // from[i]：a[i] 接在誰後面，-1 是開頭
    for (int i = 0; i < n; i++) {
        auto it = lower_bound(tail.begin(), tail.end(), a[i]);
        int k = it - tail.begin();
        if (k == (int)tail.size()) {
            tail.push_back(a[i]);
            tailAt.push_back(i);
        } else {
            tail[k] = a[i];
            tailAt[k] = i;
        }
        if (k > 0) {
            from[i] = tailAt[k - 1];   // 接在長度 k 的最小結尾後面
        }
    }
    vector<int> seq;
    for (int i = tailAt.back(); i != -1; i = from[i]) {
        seq.push_back(a[i]);
    }
    reverse(seq.begin(), seq.end());
    cout << seq.size() << '\n';
    for (int i = 0; i < (int)seq.size(); i++) {
        cout << seq[i] << (i + 1 < (int)seq.size() ? ' ' : '\n');
    }
    return 0;
}
\end{example}

\iosample{
8 \\
7 2 8 4 9 5 3 6 \\
}{
4 \\
2 4 5 6 \\
}

範例每一輪結束時的狀態如下（位置從 $0$ 開始數）：

\begin{center}
\begin{tabular}{cccllc}
\toprule
$i$ & \texttt{a[i]} & $k$ & \texttt{tail} & \texttt{tailAt} & \texttt{from[i]} \\
\midrule
0 & 7 & 0 & 7 & 0 & $-1$ \\
1 & 2 & 0 & 2 & 1 & $-1$ \\
2 & 8 & 1 & 2 8 & 1 2 & 1 \\
3 & 4 & 1 & 2 4 & 1 3 & 1 \\
4 & 9 & 2 & 2 4 9 & 1 3 4 & 3 \\
5 & 5 & 2 & 2 4 5 & 1 3 5 & 3 \\
6 & 3 & 1 & 2 3 5 & 1 6 5 & 1 \\
7 & 6 & 3 & 2 3 5 6 & 1 6 5 7 & 5 \\
\bottomrule
\end{tabular}
\end{center}

最後 \texttt{tailAt} 的最後一格是 $7$，往回走：位置 $7$（值 $6$）、
\texttt{from[7]} $= 5$（值 $5$）、\texttt{from[5]} $= 3$（值 $4$）、
\texttt{from[3]} $= 1$（值 $2$）、\texttt{from[1]} $= -1$ 停下。 倒過來是
\texttt{2 4 5 6}。 對照最後的 \texttt{tail} $=$ \texttt{2 3 5 6}：其中的 $3$
在位置 $6$，排在 $5$（位置 $5$）後面，不是合法的子序列，所以真的不能直接印
\texttt{tail}。

\begin{itemize}
\item \texttt{from[i]} 在第 $i$ 輪就決定，之後不會再改。 後來
\texttt{tailAt[k - 1]} 換成別的位置，也不影響已經記下的 \texttt{from[i]}：當
時那條長度 $k$ 的子序列仍然存在，只是 \texttt{tail} 不再記它。
\item 往回走的起點是 \texttt{tailAt.back()}，也就是目前最長的那一種長度的最小
結尾。 任何一個長度等於答案的結尾都可以當起點，挑這一個最方便。
\item \texttt{(int)tail.size()}：\texttt{size()} 的型別是不帶正負號的
\texttt{size\_t}，和 \texttt{int} 比較時 \texttt{-Wall -Wextra} 會警告，先轉
成 \texttt{int}。
\end{itemize}

這個範例的最長遞增子序列只有 \texttt{2 4 5 6} 一個，但一般來說可能有好幾個，
這支程式只印其中一個。 題目如果要求印出特定的一個（例如字典序最小），回溯的方
式就要照題目的規定重新設計，不能直接用這支程式。

\subsection{以每個位置結尾的長度}
\label{sec:ss-lis-each}

範例~\ref{ex:ss-lis-path} 還藏著一個副產品：\texttt{a[i]} 放進 \texttt{tail[k]}
時，$k + 1$ 剛好就是「以 \texttt{a[i]} 結尾的最長遞增子序列」的長度，也就是
$O(n^2)$ 寫法裡的 \texttt{dp[i]}。 上表 $k$ 那一欄加一，得到
$1, 1, 2, 2, 3, 3, 2, 4$，和用 $O(n^2)$ 算出來的完全一樣。 所以 $n$ 到
$2 \times 10^5$、又需要每個位置各自的長度時，不用退回 $O(n^2)$。

\begin{whybox}[$k + 1$ 為什麼剛好是以 a[i] 結尾的最長長度]
\textbf{至少 $k + 1$。} $k > 0$ 時，\texttt{lower\_bound} 回傳 $k$，表示
\texttt{tail[k - 1]} $<$ \texttt{a[i]}；而 \texttt{tail[k - 1]} 是某條長度
$k$ 的遞增子序列的結尾，位置在 $i$ 前面。 把 \texttt{a[i]} 接上去，就是長度
$k + 1$。

\textbf{不會更多。} 假設有一條長度 $k + 2$ 的遞增子序列以 \texttt{a[i]} 結尾，
它的倒數第二個數 \texttt{a[j]}（$j < i$，\texttt{a[j]} $<$ \texttt{a[i]}）是
一條長度 $k + 1$ 的結尾。 處理完 \texttt{a[j]} 之後，\texttt{tail[k]} 就存在，
而且不超過 \texttt{a[j]}；\texttt{tail} 的每一格只會變小，所以處理
\texttt{a[i]} 時 \texttt{tail[k]} $\le$ \texttt{a[j]} $<$ \texttt{a[i]}。
這樣 \texttt{lower\_bound} 不會停在 $k$，矛盾。
\end{whybox}

\subsection{例題：指定要唱的那首歌}
\label{sec:ss-ex-sing}

\textbf{題意}（\akiref{the-song-i-must-sing}{指定要唱的那首歌}）

\begin{itemize}
\item 歌單上依序有 $n$ 首歌，第 $i$ 首的難度是 $a_i$。
\item 參賽者從歌單挑一些歌照順序唱，難度要\textbf{嚴格}遞增，至少唱一首。
\item 有 $q$ 位參賽者，第 $j$ 位指定第 $p_j$ 首一定要唱。 對每一位，輸出在包
含第 $p_j$ 首的前提下最多能唱幾首。
\item $n, q \le 2 \times 10^5$，難度在 $1$ 到 $10^9$ 之間。
\end{itemize}

例如難度依序是 $3, 7, 2, 5, 6, 1, 8, 4$。 指定第 $1$ 首（難度 $3$），可以唱
$3, 5, 6, 8$，共 $4$ 首；指定第 $2$ 首（難度 $7$），前面只能接 $3$，後面只能
接 $8$，共 $3$ 首，雖然不指定的話能唱到 $4$ 首。

\textbf{想法}

最直接的做法：每一位參賽者都重算一次。 把「一定要唱第 $p$ 首」加進
$O(n^2)$ 的動態規劃不難，但 $q$ 次就是 $O(qn^2)$；就算每次用 $O(n \log n)$，
$q \cdot n \log n$ 也大約 $7 \times 10^{11}$，來不及。 要找一個先算好、每個
詢問直接查表的做法。

在第 $p$ 首把挑出來的歌切成兩段：第 $p$ 首和它前面的歌，是一條\textbf{以第
$p$ 首結尾}的遞增子序列；第 $p$ 首和它後面的歌，是一條\textbf{以第 $p$ 首開
頭}的遞增子序列。 兩段只在第 $p$ 首相接，前一段怎麼挑完全不影響後一段能怎麼
挑，因為後一段只要求比 $a_p$ 大、位置在 $p$ 後面。 所以兩段可以各自取最長的：
設 $L_p$ 是以第 $p$ 首結尾的最長長度、$R_p$ 是以第 $p$ 首開頭的最長長度，答
案是 $L_p + R_p - 1$（第 $p$ 首算了兩次）。

$L_p$ 就是第~\ref{sec:ss-lis-each}~小節的 $k + 1$，對整個陣列跑一次就有全部的
$L$。 $R_p$ 反過來想：從右往左讀，「以第 $p$ 首開頭、往右嚴格遞增」變成「以第
$p$ 首結尾、往左嚴格\textbf{遞減}」。 把難度都加上負號，遞減又變回遞增，可以
直接用同一個函式。 範例的兩個方向：

\begin{center}
\begin{tabular}{l*{8}{c}}
\toprule
第幾首 & 1 & 2 & 3 & 4 & 5 & 6 & 7 & 8 \\
難度 & 3 & 7 & 2 & 5 & 6 & 1 & 8 & 4 \\
\midrule
$L$（以它結尾） & 1 & 2 & 1 & 2 & 3 & 1 & 4 & 2 \\
$R$（以它開頭） & 4 & 2 & 4 & 3 & 2 & 2 & 1 & 1 \\
$L + R - 1$ & 4 & 3 & 4 & 4 & 4 & 2 & 4 & 2 \\
\bottomrule
\end{tabular}
\end{center}

同樣的模型也出現在別的地方：規劃一條一定要經過某個景點、而且一路往上爬的登山
路線；在一年的營收紀錄裡，找包含某一個月的最長成長紀錄。

\textbf{程式}

\begin{example}[label=ex:ss-sing]{cpp}{(must\_sing.cpp) 指定要唱的那首歌}
#include <algorithm>
#include <iostream>
#include <vector>
using namespace std;

// 回傳 len：len[i] 是以 v[i] 結尾的最長嚴格遞增子序列長度
vector<int> endingLengths(const vector<int> &v) {
    vector<int> tail, len(v.size());
    for (int i = 0; i < (int)v.size(); i++) {
        auto it = lower_bound(tail.begin(), tail.end(), v[i]);
        int k = it - tail.begin();
        len[i] = k + 1;                 // v[i] 接在長度 k 的後面
        if (k == (int)tail.size()) {
            tail.push_back(v[i]);
        } else {
            tail[k] = v[i];
        }
    }
    return len;
}

int main() {
    ios::sync_with_stdio(false);
    cin.tie(nullptr);
    int n;
    cin >> n;
    vector<int> a(n);
    for (int &x : a) {
        cin >> x;
    }
    vector<int> endAt = endingLengths(a);
    // 從右往左讀、難度取負號：「往右嚴格遞增」變成「往左嚴格遞增」
    vector<int> rev(n);
    for (int i = 0; i < n; i++) {
        rev[i] = -a[n - 1 - i];
    }
    vector<int> revLen = endingLengths(rev);
    int q;
    cin >> q;
    while (q--) {
        int p;
        cin >> p;
        p--;                                    // 改成從 0 開始
        int startAt = revLen[n - 1 - p];   // 以 a[p] 開頭的長度
        cout << endAt[p] + startAt - 1 << '\n';   // a[p] 算了兩次
    }
    return 0;
}
\end{example}

\iosample{
8 \\
3 7 2 5 6 1 8 4 \\
4 \\
1 2 6 8 \\
}{
4 \\
3 \\
2 \\
2 \\
}

\begin{itemize}
\item \texttt{rev[i]} 是原本的第 $n - 1 - i$ 首，所以原本第 $p$ 首在
\texttt{rev} 裡是第 $n - 1 - p$ 個。 兩個方向的位置換算，寫錯一格就全部錯開，
寫完先用範例對一次表格。
\item 難度最小是 $1$，加負號之後在 $-10^9$ 到 $-1$ 之間，\texttt{int} 裝得下。
如果難度可能是 $-2^{31}$，取負號會溢位，那時改用 \texttt{long long}，或把
\texttt{lower\_bound} 的比較改成 \texttt{greater<int>()}。
\item 一位參賽者只要查兩個數，所以詢問再多都不怕。
\end{itemize}

\textbf{複雜度}

兩次 \texttt{endingLengths} 各 $O(n \log n)$，每個詢問 $O(1)$，總共
$O(n \log n + q)$。 $n = q = 2 \times 10^5$ 時，二分搜尋大約
$2 \times 2 \times 10^5 \times 18 \approx 7 \times 10^6$ 次比較。 最大的測試資
料在本書測試的環境下約 $0.1$ 秒。

\subsection{常見錯誤}

\begin{itemize}
\item \textbf{反方向用了 \texttt{upper\_bound}}：從右往左那一次用
\texttt{upper\_bound}，等於允許後面的歌和第 $p$ 首一樣難。 輸入
\begin{code}{text}{反方向用 upper\_bound 會錯的輸入}
6
2 5 5 3 5 6
4
2 3 4 6
\end{code}
正確的輸出是 \texttt{3 3 4 4}（每行一個），這樣寫會印出 \texttt{5 4 4 4}。 指
定第 $2$ 首（難度 $5$）時，它把後面兩首難度 $5$ 的歌也接上去，唱成
$2, 5, 5, 5, 6$。 加負號之後，「後面要比較難」仍然是嚴格的，兩個方向都要用
\texttt{lower\_bound}。
\item \textbf{忘了減一}：寫成 \texttt{endAt[p] + startAt}，第 $p$ 首被算了兩次，
範例印出 \texttt{5 4 3 3}，每個都多 $1$。
\item \textbf{把 \texttt{tail} 的長度當成以它結尾的長度}：先更新
\texttt{tail}，再寫 \texttt{len[i] = tail.size()}。 \texttt{tail.size()} 是
「到目前為止的最長長度」，不一定是以 \texttt{v[i]} 結尾的；範例的第 $2$ 首
（難度 $7$）算對，第 $3$ 首（難度 $2$）就被算成 $2$，程式印出
\texttt{4 5 4 4}。
\end{itemize}

\section{最長共同子序列}
\label{sec:ss-lcs}

兩串 $A$、$B$ 的\textbf{共同子序列}，是同時是兩串子序列的序列；最長的那一個叫
\textbf{最長共同子序列}（longest common subsequence，LCS）。 章首的夜市就是在
找兩份清單的最長共同子序列。 把它想成配對也可以：在 $A$ 和 $B$ 之間畫線，每
條線連起兩個相同的元素，每個元素最多連一條，而且線不能交叉。 線最多能畫幾條，
就是最長共同子序列的長度。

\subsection{表格怎麼填}
\label{sec:ss-lcs-table}

照第~\ref{sec:ss-states}~節的第二種狀態：

\begin{itemize}
\item \textbf{狀態}：\texttt{f[i][j]} 是 $A$ 的前 $i$ 個和 $B$ 的前 $j$ 個的最
長共同子序列長度。 $A$ 的第 $i$ 個寫成 $a_i$，在程式裡是 \texttt{a[i - 1]}。
\item \textbf{轉移}：看兩串最後的 $a_i$ 和 $b_j$。
  \begin{itemize}
  \item $a_i = b_j$：讓它們配在一起，前面剩下 $A$ 的前 $i - 1$ 個和 $B$ 的前
  $j - 1$ 個，\texttt{f[i][j] = f[i - 1][j - 1] + 1}。
  \item $a_i \ne b_j$：它們不可能互相配對，所以至少有一個完全沒用到。 $a_i$
  沒用到，答案是 \texttt{f[i - 1][j]}；$b_j$ 沒用到，答案是
  \texttt{f[i][j - 1]}。 兩者取大的。
  \end{itemize}
\item \textbf{初始值}：\texttt{f[0][j] = f[i][0] = 0}，有一串是空的，就沒有東
西可以配。
\item \textbf{計算順序}：每一格用到上、左、左上三格，照列由上而下、每列由左而
右填，和第 1 本第~10~章的格子一樣。
\end{itemize}

$a_i \ne b_j$ 的情況為什麼「至少有一個沒用到」？ 如果兩個都有用到，$a_i$ 配到
$B$ 的某個 $b_{j'}$（$j' < j$），$b_j$ 配到 $A$ 的某個 $a_{i'}$（$i' < i$），
這兩條線一條從 $A$ 的最右邊往左下連、一條從 $B$ 的最右邊往左上連，一定交叉。

$a_i = b_j$ 時直接配，會不會錯過更好的？ 不會。 任取一個最長的配法：如果
$a_i$ 和 $b_j$ 都沒配，加上這條線就更長了，不可能；兩個都配給了別人，那兩條線
會交叉，也不可能；剩下的情況是其中一個配給了別人，例如 $a_i$ 配了 $b_{j'}$，
就把那條線改成連到 $b_j$。 $b_j$ 在 $b_{j'}$ 右邊，
而且是 $B$ 的最後一個，新的線不會和其他線交叉，條數不變。 所以一定有一個最長的
配法讓 $a_i$ 和 $b_j$ 配在一起。

下圖是 $A = $ \texttt{STREAM}、$B = $ \texttt{MASTER} 填好的表。 兩串用的字母
完全相同，只是順序不同，最長共同子序列卻只有 $3$。

\begin{center}
\begin{tikzpicture}[x=0.66cm, y=0.66cm,
    g/.style={draw=black!55, minimum size=0.66cm, inner sep=0pt, font=\small\ttfamily},
    pth/.style={g, fill=black!10},
    hit/.style={g, draw=memStroke, thick, fill=memFill},
    lbl/.style={font=\small\ttfamily}]
\foreach \c [count=\j from 1] in {M,A,S,T,E,R} { \node[lbl] at (\j,1) {\c}; }
\foreach \c [count=\i from 1] in {S,T,R,E,A,M} { \node[lbl] at (-1,-\i) {\c}; }
\foreach \row [count=\r from 0] in {{0,0,0,0,0,0,0},{0,0,0,1,1,1,1},{0,0,0,1,2,2,2},{0,0,0,1,2,2,3},{0,0,0,1,2,3,3},{0,0,1,1,2,3,3},{0,1,1,1,2,3,3}} {
  \foreach \v [count=\c from 0] in \row {
    \node[g] at (\c,-\r) {\v};
  }
}
\foreach \p in {(6,-6),(6,-5),(6,-4),(5,-2),(2,0)} { \node[pth] at \p {}; }
\node[pth] at (6,-6) {3}; \node[pth] at (6,-5) {3}; \node[pth] at (6,-4) {3};
\node[pth] at (5,-2) {2}; \node[pth] at (2,0) {0};
\node[hit] at (6,-3) {3}; \node[hit] at (4,-2) {2}; \node[hit] at (3,-1) {1};
\node[idxlbl] at (0,1) {$j = 0$};
\node[idxlbl, left] at (-1.3,0) {$i = 0$};
\end{tikzpicture}
\diagcaption{$A = $ \texttt{STREAM}（列）、$B = $ \texttt{MASTER}（行）的
\texttt{f[i][j]}。 灰色和橘色是從右下角往回走經過的格子，橘色是配對的地方，
讀出 \texttt{S}、\texttt{T}、\texttt{R}。}
\label{fig:ss-lcs-table}
\end{center}

\subsection{從表格還原}
\label{sec:ss-lcs-path}

要印出共同子序列本身，從右下角 \texttt{f[n][m]} 往回走，每一步照填表時的理由
倒推：

\begin{itemize}
\item $a_i = b_j$：這一對配在一起，把這個字記下來，走到左上
$(i - 1, j - 1)$。
\item $a_i \ne b_j$：答案來自上面或左邊比較大的那一格，就往那裡走。 一樣大時
走哪一邊都可以，但要固定一種規則。
\end{itemize}

走到某一串用完（$i = 0$ 或 $j = 0$）就停。 記下來的字是倒著的，最後反轉。

\begin{example}[label=ex:ss-lcs-path]{cpp}{(lcs\_path.cpp) 最長共同子序列的長度和一個例子}
#include <algorithm>
#include <iostream>
#include <string>
#include <vector>
using namespace std;

int main() {
    string a, b;
    cin >> a >> b;
    int n = a.size(), m = b.size();
    // f[i][j]：a 的前 i 個字和 b 的前 j 個字，最長共同子序列的長度
    vector<vector<int>> f(n + 1, vector<int>(m + 1, 0));
    for (int i = 1; i <= n; i++) {
        for (int j = 1; j <= m; j++) {
            if (a[i - 1] == b[j - 1]) {
                f[i][j] = f[i - 1][j - 1] + 1;
            } else {
                f[i][j] = max(f[i - 1][j], f[i][j - 1]);
            }
        }
    }
    // 從右下角往回走，倒著收集配對的字
    string common;
    int i = n, j = m;
    while (i > 0 && j > 0) {
        if (a[i - 1] == b[j - 1]) {
            common += a[i - 1];
            i--;
            j--;
        } else if (f[i - 1][j] >= f[i][j - 1]) {
            i--;                        // a 的第 i 個字沒用到
        } else {
            j--;                        // b 的第 j 個字沒用到
        }
    }
    reverse(common.begin(), common.end());
    cout << f[n][m] << '\n' << common << '\n';
    return 0;
}
\end{example}

\iosample{
STREAM MASTER \\
}{
3 \\
STR \\
}

照圖~\ref{fig:ss-lcs-table} 走一次：從 $(6, 6)$ 開始，\texttt{M} 和
\texttt{R} 不同，上面和左邊都是 $3$，規則是一樣大時往上，走到 $(5, 6)$、再到
$(4, 6)$；$(3, 6)$ 的 \texttt{R} 和 \texttt{R} 相同，記下 \texttt{R}，走到
$(2, 5)$；\texttt{T} 和 \texttt{E} 不同，上面 $1$、左邊 $2$，往左到 $(2, 4)$，
記下 \texttt{T}；再到 $(1, 3)$ 記下 \texttt{S}，走到 $(0, 2)$ 停下。

\texttt{STE} 也是長度 $3$ 的共同子序列。 一樣大時改成往左走，就會印出它（本章
練習第~2~題）。 題目只要長度時，這些都不用管；要印出子序列時，先看題目規定印
哪一個。

\begin{itemize}
\item 還原要用到整張表，所以這支程式沒有省記憶體。 下一節的滾動陣列只能求長
度。
\item 往回走最多 $n + m$ 步，比填表的 $nm$ 快得多。
\end{itemize}

\subsection{其中一串沒有重複時}
\label{sec:ss-lcs-lis}

回到章首的夜市。 把你的清單依序編號：雞排 $1$、珍奶 $2$、臭豆腐 $3$、蚵仔煎
$4$、烤玉米 $5$、地瓜球 $6$、糖葫蘆 $7$。 朋友的清單換成編號是
\[
3,\ 1,\ 4,\ 2,\ 6,\ 5,\ 7.
\]
一起吃的攤位，在朋友的清單裡照順序出現，在你的清單裡也要照順序，也就是編號要
遞增。 所以答案是這串編號的最長遞增子序列：$1, 4, 6, 7$（雞排、蚵仔煎、地瓜
球、糖葫蘆），共 $4$ 攤。 $1, 2, 5, 7$ 也可以。

這個轉換只要求\textbf{其中一串沒有重複}：每個元素在那一串裡只有一個位置，另
一串的元素換成這個位置（不在那一串裡的直接丟掉），共同子序列就變成位置遞增的
子序列。 一般的最長共同子序列要 $O(nm)$，兩串都到 $2 \times 10^5$ 時有
$4 \times 10^{10}$ 格，做不到；換成最長遞增子序列只要 $O(n \log n)$。 兩串都
有重複時，同一個元素有好幾個位置可以選，這個轉換就不成立了。

\section{滾動陣列：只留兩列}
\label{sec:ss-rolling}

兩串各 $5000$ 個時，最長共同子序列的表有 $5001 \times 5001$ 格，計算大約
$2500$ 萬次，很快；但一格 \texttt{int} 是 $4$ 位元組，整張表大約 $100$ MB。
記憶體限制比這小的時候，就要用第 1 本第~10~章 10.3 節的滾動陣列。

看轉移式，第 $i$ 列只用到第 $i - 1$ 列和第 $i$ 列自己，所以留兩列就夠：
\texttt{prev} 是第 $i - 1$ 列，\texttt{cur} 是正在填的第 $i$ 列，填完一列就
\texttt{swap(prev, cur)}。 記憶體從 $O(nm)$ 變成 $O(m)$。

想只留一列，要多注意一件事。 最長共同子序列用到左上角 \texttt{f[i-1][j-1]}，
由左往右填到第 $j$ 格時，\texttt{f[j - 1]} 已經被換成第 $i$ 列的值，左上角的
舊值不見了。 第 1 本第~10~章 10.3 節最後提過這種情況：用一個變數在被蓋掉之前先存起
來。

\begin{code}{cpp}{只留一列：用 diag 存住左上角}
vector<int> f(m + 1, 0);
for (int i = 1; i <= n; i++) {
    int diag = 0;                 // 上一列的 f[j - 1]，也就是左上角
    for (int j = 1; j <= m; j++) {
        int up = f[j];            // 上一列的 f[j]，等一下會被蓋掉
        if (a[i - 1] == b[j - 1]) {
            f[j] = diag + 1;
        } else {
            f[j] = max(up, f[j - 1]);
        }
        diag = up;                // 下一格的左上角，就是這一格的上面
    }
}
\end{code}

兩列的寫法多用一列記憶體，但每一格「讀的是哪一列」一目了然；一列的寫法省一
半，要小心變數更新的先後。 記憶體夠的時候，建議用兩列。

滾動陣列丟掉了前面的列，就沒辦法照第~\ref{sec:ss-lcs-path}~小節的方法還原。
題目同時要求印出子序列、記憶體又放不下整張表時，需要更進一步的技巧，超出本書
的範圍。

\subsection{例題：兩座測站的容許誤差比對}
\label{sec:ss-ex-stations}

\textbf{題意}（\akiref{cross-checking-two-weather-stations}{兩座測站的容許誤差比對}）

\begin{itemize}
\item 測站 A 依時間記錄了 $n$ 筆讀數 $a_1, \ldots, a_n$，測站 B 記錄了 $m$ 筆
$b_1, \ldots, b_m$。
\item 給定容許誤差 $D$，$|a_i - b_j| \le D$ 時說兩筆讀數「吻合」。
\item 挑出一些配對 $(i, j)$，每對都要吻合，兩邊的編號都要嚴格遞增（配對不交
叉，每筆讀數最多用一次）。 求最多能挑幾對。
\item $n, m \le 5000$，讀數和 $D$ 在 $0$ 到 $10^9$ 之間，記憶體限制 $32$ MB。
\end{itemize}

例如 $D = 2$，$a = 12, 15, 9, 20, 14$，$b = 10, 17, 13, 21, 8, 15$。 可以挑
$12$ 配 $10$、$15$ 配 $13$、$20$ 配 $21$、$14$ 配 $15$，共 $4$ 對。

\textbf{想法}

把配對畫成線，這就是最長共同子序列，只是「兩個元素相同」換成「兩筆讀數吻
合」。 回頭檢查第~\ref{sec:ss-lcs-table}~小節的兩段理由，用到的只有「$a_i$ 和
$b_j$ 能不能配」：

\begin{itemize}
\item 不能配時，至少一個沒用到：交叉的理由和「能不能配」怎麼定義無關。
\item 能配時直接配：把 $a_i$ 原本的線改連到 $b_j$，需要的是 $a_i$ 和 $b_j$ 能
配，這正是這個情況的前提。
\end{itemize}

所以轉移式一個字都不用改，只把 \texttt{a[i - 1] == b[j - 1]} 換成
\texttt{abs(a[i - 1] - b[j - 1]) <= D}。

要注意的是，吻合\textbf{沒有遞移性}：$D = 3$ 時，$2$ 和 $5$ 吻合、$5$ 和 $8$
吻合，$2$ 和 $8$ 卻不吻合。 所以不能先把讀數分組（例如除以 $D + 1$ 取商），
再當成「同一組就相同」去做一般的最長共同子序列。 每一對都要直接算差。

最後是記憶體。 整張表 $5001 \times 5001$ 個 \texttt{int}，大約 $100$ MB，超
過 $32$ MB，所以用兩列的滾動陣列，記憶體只剩幾十 KB。

同樣的模型也出現在別的地方：兩支手環各自記錄的心跳、兩台收銀機各自的交易時
間，取樣時間不同步、偶爾漏記，要對齊兩串紀錄時，都可以用「誤差在範圍內就算
吻合」的最長共同子序列。

\textbf{程式}

\begin{example}[label=ex:ss-stations]{cpp}{(stations.cpp) 兩座測站的容許誤差比對}
#include <algorithm>
#include <cstdlib>
#include <iostream>
#include <vector>
using namespace std;

int main() {
    ios::sync_with_stdio(false);
    cin.tie(nullptr);
    int n, m, d;
    cin >> n >> m >> d;
    vector<int> a(n), b(m);
    for (int &x : a) {
        cin >> x;
    }
    for (int &x : b) {
        cin >> x;
    }
    // prev：表格的第 i - 1 列；cur：正在填的第 i 列
    vector<int> prev(m + 1, 0), cur(m + 1, 0);
    for (int i = 1; i <= n; i++) {
        cur[0] = 0;
        for (int j = 1; j <= m; j++) {
            cur[j] = max(prev[j], cur[j - 1]);       // 有一筆沒用到
            if (abs(a[i - 1] - b[j - 1]) <= d) {
                cur[j] = max(cur[j], prev[j - 1] + 1);   // 配成一對
            }
        }
        swap(prev, cur);
    }
    cout << prev[m] << '\n';
    return 0;
}
\end{example}

\iosample{
5 6 2 \\
12 15 9 20 14 \\
10 17 13 21 8 15 \\
}{
4 \\
}

\begin{itemize}
\item 迴圈結束時剛好又 \texttt{swap} 了一次，最後一列在 \texttt{prev} 裡，所以
輸出 \texttt{prev[m]}。 印 \texttt{cur[m]} 會印到倒數第二列。
\item 能配的時候，這裡仍然和上面、左邊比一次大小。 上一節說過能配時直接配不會
吃虧，所以 \texttt{prev[j - 1] + 1} 一定不比另外兩個小，多比一次只是不用另外
想，結果一樣。
\item 讀數到 $10^9$，兩個相減在 $-10^9$ 到 $10^9$ 之間，\texttt{int} 不會溢
位。 \texttt{abs} 對 \texttt{int} 要 \texttt{\#include <cstdlib>}。
\end{itemize}

\textbf{複雜度}

$O(nm)$ 時間，$n = m = 5000$ 時 $2500$ 萬格，每格只做幾個比較；記憶體
$O(n + m)$。 最大的測試資料在本書測試的環境下約 $0.1$ 秒，整支程式的記憶體
大約 $4$ MB，離 $32$ MB 的限制很遠。

\subsection{常見錯誤}

\begin{itemize}
\item \textbf{只留一列，左上角被蓋掉}：把 \texttt{prev}、\texttt{cur} 合成一
個陣列 \texttt{f}，能配時寫 \texttt{f[j - 1] + 1}。 範例會印出 \texttt{5}，正
確是 $4$。 \texttt{f[j - 1]} 已經是這一列的值，裡面可能已經用過 $a_i$，再加一
等於讓 $a_i$ 配了兩次。 要只留一列，就照上一節用 \texttt{diag} 存住左上角。
\item \textbf{開整張表}：\texttt{vector<vector<int>> f(n + 1, vector<int>(m + 1))}
在 $n = m = 5000$ 時大約 $100$ MB，超過記憶體限制。 動手前先用格子數乘上每格的
位元組數估一次。
\item \textbf{先分組再比相等}：讀數除以 $D + 1$ 取商，同一組才算吻合。 題目
的第二組範例（$D = 3$，$a = 2, 9, 5, 6$，$b = 8, 4, 1, 3$）正確答案是 $3$，
這樣做只得到 $2$：$9$ 和 $8$ 吻合，卻被分到不同組。
\item \textbf{$\le$ 寫成 $<$}：範例會印出 \texttt{3}。 $D = 0$ 時更明顯，只
有完全相等才吻合，寫成 $<$ 就一對都配不到。
\end{itemize}

\section{編輯距離}
\label{sec:ss-edit}

打字打錯了，要改幾個字才會對？ 對一個字串可以做三種\textbf{編輯}：插入一個字、
刪掉一個字、把一個字換成另一個字。 把字串 $A$ 變成 $B$ 最少要幾次編輯，叫做
$A$ 到 $B$ 的\textbf{編輯距離}（edit distance），也叫 Levenshtein 距離。 拼字
檢查、比對兩份文件的差異、比對兩段 DNA，都會用到它。

\subsection{三種操作，三個方向}
\label{sec:ss-edit-table}

狀態和最長共同子序列一樣是兩串的前綴：\texttt{f[i][j]} 是把 $A$ 的前 $i$ 個
字改成 $B$ 的前 $j$ 個字，最少要幾次編輯。 看兩串的最後一個字，最佳的改法只有
三種結局：

\begin{itemize}
\item $a_i$ 最後變成 $b_j$：一樣就不用改，不一樣就換一次。 前面是 $A$ 的前
$i - 1$ 個改成 $B$ 的前 $j - 1$ 個，從左上角 \texttt{f[i - 1][j - 1]} 來。
\item $a_i$ 被刪掉：花一次，前面是 $A$ 的前 $i - 1$ 個改成 $B$ 的前 $j$ 個，
從上面 \texttt{f[i - 1][j]} 來。
\item $b_j$ 是插入的：花一次，前面是 $A$ 的前 $i$ 個改成 $B$ 的前 $j - 1$ 個，
從左邊 \texttt{f[i][j - 1]} 來。
\end{itemize}

三者取最小。 初始值也要想清楚：$A$ 的前 $i$ 個改成空字串，只能全部刪掉，
\texttt{f[i][0] = i}；空字串改成 $B$ 的前 $j$ 個，只能全部插入，
\texttt{f[0][j] = j}。 這兩行和最長共同子序列的「全部是 $0$」不同，是最常寫錯
的地方。

\begin{example}[label=ex:ss-edit]{cpp}{(edit.cpp) 編輯距離}
#include <algorithm>
#include <iostream>
#include <string>
#include <vector>
using namespace std;

int main() {
    string a, b;
    cin >> a >> b;
    int n = a.size(), m = b.size();
    // f[i][j]：把 a 的前 i 個字改成 b 的前 j 個字，最少要幾次編輯
    vector<vector<int>> f(n + 1, vector<int>(m + 1));
    for (int i = 0; i <= n; i++) {
        f[i][0] = i;                    // 全部刪掉
    }
    for (int j = 0; j <= m; j++) {
        f[0][j] = j;                    // 全部插入
    }
    for (int i = 1; i <= n; i++) {
        for (int j = 1; j <= m; j++) {
            int same = (a[i - 1] == b[j - 1]) ? 0 : 1;
            f[i][j] = min({f[i - 1][j - 1] + same,   // 不改或替換
                           f[i - 1][j] + 1,          // 刪掉 a[i]
                           f[i][j - 1] + 1});        // 插入 b[j]
        }
    }
    cout << f[n][m] << '\n';
    return 0;
}
\end{example}

\iosample{
CAFE CHAI \\
}{
3 \\
}

\begin{center}
\begin{tikzpicture}[x=0.7cm, y=0.7cm,
    g/.style={draw=black!55, minimum size=0.7cm, inner sep=0pt, font=\small\ttfamily},
    pth/.style={g, draw=memStroke, thick, fill=memFill},
    lbl/.style={font=\small\ttfamily}]
\foreach \c [count=\j from 1] in {C,H,A,I} { \node[lbl] at (\j,1) {\c}; }
\foreach \c [count=\i from 1] in {C,A,F,E} { \node[lbl] at (-1,-\i) {\c}; }
\foreach \row [count=\r from 0] in {{0,1,2,3,4},{1,0,1,2,3},{2,1,1,1,2},{3,2,2,2,2},{4,3,3,3,3}} {
  \foreach \v [count=\c from 0] in \row {
    \node[g] at (\c,-\r) {\v};
  }
}
\node[pth] at (0,0) {0}; \node[pth] at (1,-1) {0}; \node[pth] at (2,-1) {1};
\node[pth] at (3,-2) {1}; \node[pth] at (4,-3) {2}; \node[pth] at (4,-4) {3};
\node[font=\scriptsize, right] at (5.2,0) {$(0,0) \to (1,1)$：\texttt{C} 不用改};
\node[font=\scriptsize, right] at (5.2,-1) {$(1,1) \to (1,2)$：插入 \texttt{H}};
\node[font=\scriptsize, right] at (5.2,-2) {$(1,2) \to (2,3)$：\texttt{A} 不用改};
\node[font=\scriptsize, right] at (5.2,-3) {$(2,3) \to (3,4)$：\texttt{F} 換成 \texttt{I}};
\node[font=\scriptsize, right] at (5.2,-4) {$(3,4) \to (4,4)$：刪掉 \texttt{E}};
\end{tikzpicture}
\diagcaption{\texttt{CAFE}（列）改成 \texttt{CHAI}（行）的表，答案 $3$。 橘色是
一條最佳的路徑：往右是插入，往下是刪除，往右下是不用改或替換。}
\label{fig:ss-edit-table}
\end{center}

路徑上的每一步都對應一次操作，所以還原的方法和最長共同子序列一樣：從右下角往
回走，每一步找一個「算得出這一格」的來源。 這張表裡最佳的改法不只一種：把
\texttt{A}、\texttt{F}、\texttt{E} 依序換成 \texttt{H}、\texttt{A}、
\texttt{I}，也是 $3$ 次，對應表裡一路往右下的對角線。

編輯距離和最長共同子序列是近親。 如果不准替換，只能插入和刪除，那麼沒被刪掉的
字和沒被插入的字要一一對上，它們就是一個共同子序列；保留得越多，要刪、要插的
就越少。 所以這時的最少次數是 $n + m - 2 \times$ 最長共同子序列的長度（本章練
習第~5~題）。

\subsection{例題：在聊天紀錄裡找打錯的句子}
\label{sec:ss-ex-typo}

\textbf{題意}（\akiref{finding-a-mistyped-sentence-in-the-chat-log}{在聊天紀錄裡找打錯的句子}）

\begin{itemize}
\item 聊天紀錄是一個小寫字母字串 $T$，要找的句子是 $P$，都去掉了空白和標點。
\item $T$ 的一段連續子字串 $S = t_l \cdots t_r$，代價是 $P$ 到 $S$ 的編輯距離。
\item 求最小的代價 $c$。 代價都是 $c$ 的子字串裡，選最短的；一樣短就選開始位
置 $l$ 最小的。 輸出 $c$、$l$、$r$（位置從 $1$ 開始）。
\item $|T| \le 5000$，$|P| \le 1000$。
\end{itemize}

例如 $T = $ \texttt{wemeetinlibaryatfive}、$P = $ \texttt{library}。 第 $9$ 到
第 $14$ 個字是 \texttt{libary}，從 \texttt{library} 刪掉一個 \texttt{r} 就得到，
代價 $1$；$T$ 裡沒有一段和 \texttt{library} 完全相同，所以答案是
\texttt{1 9 14}。 前面的 \texttt{wemeetin} 不在 $S$ 裡，不算代價。

\textbf{想法}

最直接的做法：枚舉所有 $(l, r)$，各算一次編輯距離。 子字串有大約
$1.25 \times 10^7$ 個，每個再乘上 $|P| \times |S|$，完全不行。 稍微好一點：固定
起點 $l$，對 $P$ 和 $t_l t_{l+1} \cdots t_n$ 做一次編輯距離，表格的最後一列就同
時給出所有結束位置 $r$ 的代價。 這樣是 $|T|$ 次、每次 $|P| \times |T|$，總共
$2.5 \times 10^{10}$，還是太慢。

\textbf{觀察一：改初始值，起點就自由了。} 照編輯距離的狀態，讓
\texttt{f[i][j]} 表示「$P$ 的前 $i$ 個字改成某個\textbf{結束在 $t_j$} 的子字
串」的最少次數，子字串從哪裡開始都可以。 和上一小節只差在第 $0$ 列：$P$ 的前
$0$ 個字（空字串）要變成一個結束在 $t_j$ 的子字串，選空的那一段就好，不用任何
操作，所以 \texttt{f[0][j] = 0}，而不是 $j$。 $t_1 \cdots t_j$ 這些字不在
$S$ 裡，本來就不該算錢。 其他的轉移完全不變：第 $0$ 行
\texttt{f[i][0] = i} 也不變，$P$ 的字不能免費丟掉。 最後一列
\texttt{f[|P|][j]} 是結束在 $t_j$ 的最小代價，對所有 $j$ 取最小，結束位置也自
由了。 這樣只要做一次 $|P| \times |T| = 5 \times 10^6$ 格的表。

\textbf{觀察二：每格多記一個起點。} 題目還要最短的子字串。 同樣結束在 $t_j$、
代價一樣，起點越晚越短，所以每格除了代價，再記「同代價下最晚的起點」。 第 $0$
列的空字串結束在 $t_j$，起點記成 $j + 1$。 三種轉移都只是把某一格的代價加上
$0$ 或 $1$，起點原封不動帶過來。 所以把一格寫成
$(\text{代價}, -\text{起點})$ 這一對，先比代價、再比 $-$起點，取最小的那一對。
為什麼每格只記最好的一對就夠？ 從同一個來源格出發的所有改法，代價都加上同一個
數，原本誰比較好，加完還是誰比較好；所以來源格只要交出它最好的那一對，不會漏
掉更好的。 \texttt{pair} 的 \texttt{<} 剛好就是先比第一個、再比第二個（銜接
書第~8~章）。

\textbf{一個邊界：空字串。} 最後一列的某一格，可能是「把 $P$ 全刪光、對到空字
串」，代價 $|P|$，起點 $j + 1$ 比結束位置還大，要跳過。 被它擋住的情況只有一
種：最小代價剛好是 $|P|$。 這時任何一個字的子字串代價都不超過 $|P|$（換掉一
個字、刪掉其他 $|P| - 1$ 個），所以答案就是 $|P|$，取最短、最左的 $t_1$。

同樣的模型也出現在別的地方：在一長段 DNA 裡找最像某個片段的區域；在搜尋引擎裡
輸入打錯的關鍵字，仍然在文章中標出最接近的那幾個字。

\textbf{程式}

\begin{example}[label=ex:ss-typo]{cpp}{(typo.cpp) 在聊天紀錄裡找打錯的句子}
#include <algorithm>
#include <iostream>
#include <string>
#include <utility>
#include <vector>
using namespace std;

int main() {
    ios::sync_with_stdio(false);
    cin.tie(nullptr);
    string t, p;
    cin >> t >> p;
    int n = t.size(), m = p.size();
    // 一格存 {代價, -起點}。
    // 第 i 列是 P 的前 i 個字，第 j 行是「結束在 t_j」。
    // 第 0 列：空字串對到結束在 t_j 的空字串，起點 j + 1。
    vector<pair<int, int>> prev(n + 1), cur(n + 1);
    for (int j = 0; j <= n; j++) {
        prev[j] = {0, -(j + 1)};
    }
    for (int i = 1; i <= m; i++) {
        cur[0] = {i, -1};                    // P 的前 i 個字全刪
        for (int j = 1; j <= n; j++) {
            auto [c1, s1] = prev[j - 1];     // P 的第 i 個字變成 t_j
            auto [c2, s2] = prev[j];         // 刪掉 P 的第 i 個字
            auto [c3, s3] = cur[j - 1];      // t_j 是插入的
            int diff = (p[i - 1] != t[j - 1]) ? 1 : 0;
            cur[j] = min({make_pair(c1 + diff, s1),
                          make_pair(c2 + 1, s2),
                          make_pair(c3 + 1, s3)});
        }
        swap(prev, cur);
    }
    int bestC = m + 1, bestL = 0, bestR = 0;
    for (int j = 1; j <= n; j++) {
        int c = prev[j].first, l = -prev[j].second;
        if (l > j) {
            continue;                        // 空字串，不算
        }
        bool better = c < bestC ||
            (c == bestC && (j - l < bestR - bestL ||
                            (j - l == bestR - bestL && l < bestL)));
        if (better) {
            bestC = c;
            bestL = l;
            bestR = j;
        }
    }
    if (bestC >= m) {                        // 沒有比 m 好的：取 t_1
        bestC = m;
        bestL = 1;
        bestR = 1;
    }
    cout << bestC << ' ' << bestL << ' ' << bestR << '\n';
    return 0;
}
\end{example}

\iosample{
wemeetinlibaryatfive \\
library \\
}{
1 9 14 \\
}

\begin{itemize}
\item 這裡的列是 $P$、行是 $T$，滾動的是 $P$ 的列，所以兩列各有 $|T| + 1$ 格。
填完之後 \texttt{prev} 是最後一列，裡面是每個結束位置的答案。 一整張表是
$1001 \times 5001$ 個 \texttt{pair}，大約 $40$ MB，這題的記憶體也放得下，但
用兩列一樣好寫。
\item \texttt{auto [c1, s1] = prev[j - 1];} 是結構化綁定（銜接書第~8~章），把
\texttt{pair} 的兩個成員分別取出來。 沒讀過銜接書的話，把它看成
\texttt{c1 = prev[j - 1].first; s1 = prev[j - 1].second;} 即可。
\item \texttt{min} 收一串用大括號包起來的值，型別要完全相同，所以三個候選都用
\texttt{make\_pair} 做成 \texttt{pair<int, int>}。
\item 選答案時，長度用 $j - l$ 比較，少加的 $1$ 兩邊都少，不影響大小。
\end{itemize}

\textbf{複雜度}

一張 $|P| \times |T| = 5 \times 10^6$ 格的表，每格比三個 \texttt{pair}，記憶體
$O(|T|)$。 最大的測試資料在本書測試的環境下約 $0.1$ 秒。

\subsection{常見錯誤}

\begin{itemize}
\item \textbf{照抄編輯距離的初始值}：第 $0$ 列寫成 \texttt{\{j, -1\}}，等於規
定 $S$ 一定從 $t_1$ 開始，前面的字全都要算插入。 範例會印出 \texttt{7 1 1}：
最便宜的對齊也要把 \texttt{wemeetin} 算進去，代價都超過 $7$，最後落到「沒有比
$|P|$ 更好」的那一段。
\item \textbf{同代價時保留最早結束的}：選答案時只比 \texttt{c < bestC}。 輸入
\begin{code}{text}{只保留最早結束會錯的輸入}
seeyauthenseyouandseyou
seeyou
\end{code}
正確答案是 \texttt{1 11 15}（\texttt{seyou}，刪掉一個 \texttt{e}），這樣寫會
印出 \texttt{1 1 6}：\texttt{seeyau} 結束得比較早，代價也是 $1$，但比較長。
\item \textbf{沒處理空字串}：拿掉 \texttt{l > j} 的判斷和最後的
\texttt{if}，輸入 $T = $ \texttt{xyz}、$P = $ \texttt{ab}（兩者沒有共同的字）
時印出 \texttt{2 2 1}，開始位置比結束位置還大。 正確答案是 \texttt{2 1 1}。 只
拿掉最後的 \texttt{if}，會印出 \texttt{3 0 0}。
\end{itemize}

\section{回文子序列}
\label{sec:ss-palin}

正著讀和倒著讀一樣的字串叫\textbf{回文}（palindrome），例如
\texttt{LEVEL}、\texttt{NOON}。 一個字串裡最長的回文子序列有多長？
\texttt{BANANAS} 的答案是 $5$（\texttt{ANANA}）。

一個字串和它自己倒過來的那一串，是兩串。 回文子序列倒過來還是它自己，所以它
同時是 $s$ 和「$s$ 倒過來」的子序列，也就是兩者的共同子序列。 反過來也成立：
兩者的最長共同子序列的長度，剛好等於最長回文子序列的長度（理由見下面的框）。
所以直接套用最長共同子序列，再用兩列的滾動陣列：

\begin{example}[label=ex:ss-palin]{cpp}{(palin.cpp) 最長回文子序列與最少插入次數}
#include <algorithm>
#include <iostream>
#include <string>
#include <vector>
using namespace std;

int main() {
    string s;
    cin >> s;
    string r(s.rbegin(), s.rend());     // s 倒過來
    int n = s.size();
    vector<int> prev(n + 1, 0), cur(n + 1, 0);
    for (int i = 1; i <= n; i++) {
        for (int j = 1; j <= n; j++) {
            if (s[i - 1] == r[j - 1]) {
                cur[j] = prev[j - 1] + 1;
            } else {
                cur[j] = max(prev[j], cur[j - 1]);
            }
        }
        swap(prev, cur);
    }
    int longest = prev[n];
    cout << longest << ' ' << n - longest << '\n';
    return 0;
}
\end{example}

\iosample{
BANANAS \\
}{
5 2 \\
}

倒過來的字串是用反向迭代器 \texttt{rbegin()}、\texttt{rend()} 建出來的（銜接
書第~9~章）。

第二個數字是「最少插入幾個字，才能讓整個字串變成回文」。 留下最長回文子序列，
其他每個字都在對稱的位置插入一個一樣的字和它配對：\texttt{BANANAS} 留下
\texttt{ANANA}，在前面補 \texttt{S}、在後面補 \texttt{B}，變成
\texttt{SBANANABS}，插入 $2$ 個。 不可能更少：插入 $k$ 個字變成回文之後，原本
的字裡，對稱位置上也是原本的字的那些，本身就組成 $s$ 的一個回文子序列；對稱位
置是插入的字的，最多 $k$ 個。 所以 $k \ge n - $ 最長回文子序列。 改成「最少刪
幾個字」，答案也是同一個數。

\begin{whybox}[共同子序列夠長，就有一樣長的回文子序列]
設 $s$ 和倒過來的 $r$ 有一個長度 $L$ 的共同子序列。 $r$ 的第 $j$ 個字是 $s$
的第 $n + 1 - j$ 個字，所以這個共同子序列可以寫成 $L$ 對位置
$(x_1, y_1), \ldots, (x_L, y_L)$：$s[x_k] = s[y_k]$，$x$ 由小到大，$y$ 由大到
小（$y$ 是換回 $s$ 的位置）。

$x$ 越來越大、$y$ 越來越小，所以前面若干對有 $x_k < y_k$，之後都是
$x_k \ge y_k$。 設前面有 $t$ 對。
\begin{itemize}
\item 前 $t$ 對：$x_1 < \cdots < x_t < y_t < \cdots < y_1$，依序讀
$s[x_1], \ldots, s[x_t], s[y_t], \ldots, s[y_1]$，是長度 $2t$ 的回文子序列。
\item 後 $L - t$ 對：$y_L < \cdots < y_{t+1} \le x_{t+1} < \cdots < x_L$，同樣
讀出一個回文子序列，長度 $2(L - t)$，$y_{t+1} = x_{t+1}$ 時中間共用一個字，
長度 $2(L - t) - 1$。
\end{itemize}
兩個裡面至少有一個不短於 $L$：$2t \ge L$ 的話第一個就是；否則 $L - t > L / 2$，
$2(L - t) - 1 \ge L$。
\end{whybox}

這個框也說明了一件事：共同子序列本身不一定是回文。 \texttt{ACBAC} 倒過來是
\texttt{CABCA}，最長共同子序列長度 $3$，\texttt{ABC} 就是其中一個，但它不是回
文。 要印出回文子序列，可以用框裡的方法從配對拼出來，或改用第~3~章的區間動態
規劃：直接以「$s$ 的第 $l$ 到第 $r$ 個字」為狀態，比較兩端的字。 那種寫法在回
文問題上更直接，加權、限制條件也比較好加。

\section{什麼時候用子序列 DP}
\label{sec:ss-when}

\begin{itemize}
\item 題目說「照原本的順序」「可以跳過」「不必相鄰」，答案是一些位置挑出來的
結果，就是子序列。 一定要相鄰的是子字串，回到第 1 本第~5~章的前綴和、雙指標。
\item 只有一串，挑出來的東西前後要滿足一種關係（變大、變小、不超過某個差距）：
用「以第 $i$ 個結尾」的狀態。 關係是單純的大小比較、$n$ 到 $2 \times 10^5$
時，想辦法變成最長遞增子序列，用 $O(n \log n)$；關係比較複雜、$n$ 在幾千以內，
$O(n^2)$ 就好。
\item 兩串要互相對應，配對不能交叉：用「前 $i$ 個和前 $j$ 個」的表格。 先算
$nm$ 夠不夠快，再算 $nm$ 乘上每格的位元組數放不放得下，放不下就滾動。
\item 「最少改幾次才相同」：編輯距離，先確認三種操作各花多少、哪些不准用，初始
值跟著改。
\item 回文：和自己倒過來的那串做共同子序列，或用第~3~章的區間動態規劃。
\end{itemize}

看起來像、其實不是的情況：兩串都很長（各 $2 \times 10^5$）又都有重複時，一般
的最長共同子序列做不到，要找題目的特殊結構，例如第~\ref{sec:ss-lcs-lis}~小節
「其中一串沒有重複」。 子序列要求「總和最大」而不是「最長」時，最長遞增子序列
的 \texttt{tail} 就不能用了（\texttt{tail} 靠的是「長度相同時只要記最小的結
尾」），要回到 $O(n^2)$，或用之後才會學到的資料結構。

\begin{apcsnote}
APCS 的實作題由程式自動評分，只比對輸出。 要印出子序列的題目，答案不只一個時
印哪一個，要照題目的規定；寫回溯之前先在題目裡找到這句話，「一樣大時往哪邊走」
照它決定。 只要長度的題目，就不要多寫回溯，也不用留整張表。
\end{apcsnote}

\section{本章常見錯誤}
\label{sec:ss-pitfalls}

\begin{center}
\begin{tabular}{p{0.22\linewidth}p{0.46\linewidth}p{0.2\linewidth}}
\toprule
錯誤 & 例子 & 見 \\
\midrule
嚴格與不嚴格搞混 & 嚴格遞增卻用 \texttt{upper\_bound}；反方向加負號後換了函式 &
第~\ref{sec:ss-ex-sing}~節 \\
把 \texttt{tail} 當成答案 & 直接印 \texttt{tail}；用 \texttt{tail.size()} 當成
以它結尾的長度 & 第~\ref{sec:ss-lis-path}、\ref{sec:ss-ex-sing}~節 \\
初始值照抄 & 編輯距離的第 $0$ 列、第 $0$ 行寫成 $0$；找子字串時第 $0$ 列沒改
成 $0$ & 第~\ref{sec:ss-edit-table}、\ref{sec:ss-ex-typo}~節 \\
滾動陣列讀到新值 & 只留一列卻用 \texttt{f[j - 1]} 當左上角；最後印了
\texttt{cur} & 第~\ref{sec:ss-rolling}、\ref{sec:ss-ex-stations}~節 \\
沒估記憶體 & $5000 \times 5000$ 的 \texttt{int} 表要 $100$ MB &
第~\ref{sec:ss-ex-stations}~節 \\
同分的處理 & 回溯時一樣大往哪邊走；同代價時要最短還是最早 &
第~\ref{sec:ss-lcs-path}、\ref{sec:ss-ex-typo}~節 \\
\bottomrule
\end{tabular}
\end{center}

另外，這一章每一個動態規劃都可以用 $2^n$ 的暴力解檢查：列出所有子序列，逐一
判斷。 $n \le 12$ 的隨機資料上比對幾千組，上面這些錯誤大多會現形。

\section{本章練習}
\label{sec:ss-exercises}

\begin{drills}[label=drills:ss]
\item 數列 $5, 1, 6, 2, 7, 3, 4$。
\begin{enumerate}
\item 照範例~\ref{ex:ss-lis-path} 的方式，寫出每一輪的 $k$、\texttt{tail}、
\texttt{tailAt}、\texttt{from[i]}，以及程式印出的子序列。
\item 寫出以每個位置結尾的最長長度，和以每個位置開頭的最長長度。 包含 $6$ 的
遞增子序列最長多長？
\end{enumerate}

\item 把範例~\ref{ex:ss-lcs-path} 的 \texttt{>=} 改成 \texttt{>}（一樣大時往左
走），輸入 \texttt{STREAM MASTER} 會印出什麼？ 照圖~\ref{fig:ss-lcs-table} 寫
出經過的格子。

\item 有人把最長共同子序列寫成只留一列：
\begin{code}{cpp}{練習 3}
vector<int> f(m + 1, 0);
for (int i = 1; i <= n; i++) {
    for (int j = 1; j <= m; j++) {
        if (a[i - 1] == b[j - 1]) {
            f[j] = f[j - 1] + 1;
        } else {
            f[j] = max(f[j], f[j - 1]);
        }
    }
}
\end{code}
輸入 $A = $ \texttt{A}、$B = $ \texttt{AA} 時 \texttt{f[m]} 是多少？ 正確答案
是多少？ 錯在哪裡？

\item 估計記憶體，判斷要不要用滾動陣列。
\begin{enumerate}
\item 兩串各 $3000$ 個，求最長共同子序列的長度，表格用 \texttt{int}，記憶體限
制 $64$ MB。
\item 同上，但要印出一個最長共同子序列。
\item 編輯距離，兩串各 $2 \times 10^4$ 個，表格用 \texttt{int}，記憶體限制
$256$ MB。 時間來得及嗎？
\end{enumerate}

\item 只准插入和刪除、不准替換時，把 $A$ 改成 $B$ 最少要幾次？
\begin{enumerate}
\item 說明答案是 $n + m - 2L$，$L$ 是最長共同子序列的長度。
\item 用這個式子算 \texttt{STREAM} 改成 \texttt{MASTER}，和准替換時的編輯距
離比較。
\end{enumerate}

\item 範例~\ref{ex:ss-sing} 改成「難度可以相等」（不下降），兩次
\texttt{lower\_bound} 各要怎麼改？ 用 $2, 5, 5, 3, 5, 6$ 檢查指定第 $2$ 首時
的答案。

\item 字串 \texttt{ABCBDA}。
\begin{enumerate}
\item 最長回文子序列是多少？ 寫出一個。
\item 最少插入幾個字可以變成回文？ 寫出插入後的字串。
\end{enumerate}
\end{drills}

\section{實作練習}
\label{sec:ss-practice}

下面是 Aki-Judge 上和本章相關的題目，依小節分組。 難度分四階：暖身是直接套用
本章的寫法；基本要自己改出狀態或轉移；實戰的題敘比較長，要先把題目轉換成本章
的形狀；挑戰會和其他章的技巧合在一起，或需要更進一步的觀察。

\begin{center}
\begin{tabular}{llc}
\toprule
小節 & 題目 & 難度 \\
\midrule
最長遞增子序列 & \akiref{longest-improving-streak-ending-each-day}{每天的進步紀錄長度} & 暖身 \\
 & \akiref{daily-best-on-the-sticker-card}{集點卡每日最長紀錄} & 暖身 \\
 & \akiref{relabeling-the-lockers}{重編置物櫃號碼} & 基本 \\
 & \akiref{rising-record-without-long-gaps}{不能斷太久的成長紀錄} & 基本 \\
 & \akiref{card-binders}{分冊收藏} & 實戰 \\
 & \akiref{freshmen-joining-the-club-lineage}{新生接上社團傳承線} & 挑戰 \\
\midrule
最長共同子序列 & \akiref{prefix-matching-of-dance-logs}{舞步紀錄的前綴比對} & 暖身 \\
 & \akiref{recital-check-letter-score}{背誦檢查：字母計分} & 暖身 \\
 & \akiref{the-two-most-similar-submissions}{最相似的兩份作業} & 基本 \\
 & \akiref{where-two-judges-agree}{兩位評審的共識名單} & 實戰 \\
 & \akiref{songs-that-might-be-kept}{哪些歌曲可能被保留} & 實戰 \\
 & \akiref{subtitle-sync-with-a-drift-limit}{字幕同步的漂移上限} & 實戰 \\
 & \akiref{cheapest-shared-stop-sequence}{兩人共搭的最省站序} & 實戰 \\
 & \akiref{the-looping-melody}{循環播放的旋律} & 挑戰 \\
\midrule
編輯距離與回文 & \akiref{flags-take-down-or-recolor}{拆旗換色排回文} & 基本 \\
 & \akiref{aligning-a-scanned-handout-with-missing-lines}{掃描講義缺行的對齊成本} & 挑戰 \\
\bottomrule
\end{tabular}
\end{center}

本章的三題例題也在 Aki-Judge 上，讀完各節之後可以先自己寫一次再對照。 這一章
的題單還會再增加。

\section{讀完你應該會}
\label{sec:ss-checklist}

\begin{itemize}
\item 用 $O(n \log n)$ 的寫法印出一個最長遞增子序列，並算出以每個位置結尾、以
每個位置開頭的最長長度。
\item 寫出最長共同子序列和編輯距離的狀態、三個方向的轉移和初始值，手動填一張
小表，從右下角還原出一組答案。
\item 估計表格的記憶體，放不下時改成兩列的滾動陣列；只留一列時，用一個變數保
住左上角。
\item 遇到「相等」換成別的條件、起點或終點可以自由選的變形，說出轉移式和初始值
各要改哪裡。
\item 用「和倒過來的自己求共同子序列」算最長回文子序列和最少插入次數，並知道
共同子序列不一定是回文。
\end{itemize}

\textbf{延伸閱讀}（書末會整理完整的參考資料）

\begin{itemize}
\item Thomas H. Cormen 等，\textit{Introduction to Algorithms}，第 4 版，
14.4 節：最長共同子序列的完整證明，以及用另一張表記錄方向的還原方法。
\item Steven S. Skiena，\textit{The Algorithm Design Manual}，第 3 版，10.2
節：編輯距離，以及改初始值、改成本就能變出的幾種變形（找子字串、最長共同子序
列）。
\item Steven Halim 等，\textit{Competitive Programming 4}，Book 1 的 3.5.2 節
（最長遞增子序列的 $O(n \log k)$ 寫法與還原）、Book 2 的 6.3 節（字串的對齊與
最長共同子序列）。
\item Antti Laaksonen，\textit{Competitive Programmer's Handbook}，第 7 章：最長
遞增子序列和編輯距離各一頁，適合快速複習。
\end{itemize}

\addcontentsline{toc}{section}{練習解答}
\section*{練習解答}

\begin{enumerate}
\item
\begin{enumerate}
\item 位置從 $0$ 開始：
\begin{center}
\begin{tabular}{cccllc}
\toprule
$i$ & \texttt{a[i]} & $k$ & \texttt{tail} & \texttt{tailAt} & \texttt{from[i]} \\
\midrule
0 & 5 & 0 & 5 & 0 & $-1$ \\
1 & 1 & 0 & 1 & 1 & $-1$ \\
2 & 6 & 1 & 1 6 & 1 2 & 1 \\
3 & 2 & 1 & 1 2 & 1 3 & 1 \\
4 & 7 & 2 & 1 2 7 & 1 3 4 & 3 \\
5 & 3 & 2 & 1 2 3 & 1 3 5 & 3 \\
6 & 4 & 3 & 1 2 3 4 & 1 3 5 6 & 5 \\
\bottomrule
\end{tabular}
\end{center}
從 \texttt{tailAt} 的最後一格 $6$ 往回走：$6 \to 5 \to 3 \to 1$，印出
\texttt{1 2 3 4}，長度 $4$。
\item 以它結尾（$k + 1$）：$1, 1, 2, 2, 3, 3, 4$。 以它開頭：$5$ 後面可以接
$6, 7$，是 $3$；依序是 $3, 4, 2, 3, 1, 2, 1$。 包含 $6$ 的最長是
$2 + 2 - 1 = 3$，例如 $5, 6, 7$ 或 $1, 6, 7$。
\end{enumerate}

\item 印出 \texttt{3} 和 \texttt{STE}。 從 $(6, 6)$ 開始，上面和左邊都是 $3$，
改成往左：$(6, 5)$ 的 \texttt{M} 和 \texttt{E} 不同，上面 $3$、左邊 $2$，往上
到 $(5, 5)$、$(4, 5)$；$(4, 5)$ 是 \texttt{E} 和 \texttt{E}，記下 \texttt{E}，
到 $(3, 4)$；\texttt{R} 和 \texttt{T} 不同，上面 $2$、左邊 $1$，往上到
$(2, 4)$，記下 \texttt{T}；到 $(1, 3)$ 記下 \texttt{S}。 倒過來是
\texttt{STE}。

\item \texttt{f[m]} 是 $2$，正確答案是 $1$。 $i = 1$、$j = 1$ 時兩個
\texttt{A} 配對，\texttt{f[1]} 變成 $1$；$j = 2$ 時又相等，程式用
\texttt{f[1] + 1}，可是 \texttt{f[1]} 已經是這一列的值，裡面用過 $A$ 唯一的
\texttt{A}，等於同一個 \texttt{A} 配了兩次。 左上角要用上一列的
\texttt{f[1]}，也就是 $0$，見第~\ref{sec:ss-rolling}~節的 \texttt{diag}。

\item
\begin{enumerate}
\item $3001 \times 3001 \times 4$ 位元組大約 $36$ MB，放得下，不一定要滾動；
用兩列也只要幾十 KB。
\item 要還原就要整張表，$36$ MB 放得下。 想再省，可以只存方向，一格一個
\texttt{char}，大約 $9$ MB；但長度的部分仍然可以用兩列。
\item 整張表 $4 \times 10^8$ 格，要 $1.6$ GB，一定要滾動。 時間是
$4 \times 10^8$ 格，每格幾個比較，照銜接書 10.1 節的估計，時間限制一秒時很可能
來不及，要先確認題目給的時間限制。
\end{enumerate}

\item
\begin{enumerate}
\item 只准插入、刪除時，$A$ 裡沒被刪掉的字，依序就是 $B$ 裡沒被插入的字，所以
留下來的是一個共同子序列，設長度 $\ell$。 要刪 $n - \ell$ 個、插入 $m - \ell$
個，共 $n + m - 2\ell$ 次；$\ell$ 越大越少，取 $\ell = L$ 最少。
\item $6 + 6 - 2 \times 3 = 6$ 次，例如刪掉 \texttt{R}、\texttt{A}、
\texttt{M}，在前面插入 \texttt{M}、\texttt{A}，在最後插入 \texttt{R}。 准替換
時，用範例~\ref{ex:ss-edit} 跑出來是 $5$：最後一步不刪 \texttt{M} 再插
\texttt{R}，改成把 \texttt{M} 換成 \texttt{R}，省一次。 一次替換抵得過一次刪除
加一次插入，所以准替換時不會更多。
\end{enumerate}

\item 兩次都改成 \texttt{upper\_bound}。 正方向是不下降，用
\texttt{upper\_bound}；反方向加負號之後，「後面不比較簡單」變成「負號之後不下
降」，也是 \texttt{upper\_bound}。 $2, 5, 5, 3, 5, 6$ 指定第 $2$ 首：前面
$2, 5$，後面 $5, 5, 6$（第 $3$、$5$、$6$ 首），共 $5$ 首。

\item
\begin{enumerate}
\item 長度 $5$，\texttt{ABCBA}（第 $1$、$2$、$3$、$4$、$6$ 個字）。
\item $6 - 5 = 1$ 個。 \texttt{D} 不在回文子序列裡，在它的對稱位置（第 $1$ 個
字 \texttt{A} 後面）補一個 \texttt{D}，得到 \texttt{ADBCBDA}。
\end{enumerate}
\end{enumerate}

\chapter{區間 DP}
\label{ch:interval-dp}

老師的桌上有一排四疊作業簿，由左到右是 $6$、$6$、$5$、$7$ 本。 要把它們合成
一疊交回教務處，可是每次只能把\textbf{相鄰}的兩疊搬到一起，搬一次花的力氣
是這兩疊加起來的本數。 合成一疊要搬三次，怎麼搬最省力？ 往下讀之前，先想一
條規則，例如「每次合併加起來最少的相鄰兩疊」，照你的規則算一次總共花多少
力氣。

《C++ 演算法解題 1：APCS 實作題核心》（以下簡稱第 1 本）第~9~章 9.6 節整理過一維動態規劃常見的四種狀態：前 $i$ 個、前 $i$ 個再加上
第 $i$ 個的情況、以第 $i$ 個結尾、把前 $i$ 個切成好幾段。 這四種都從序列的
左端長出來，子問題永遠是「從頭開始的一段」。 作業簿的問題不是這樣：先合併中
間的兩疊之後，剩下要處理的東西夾在中間，左右兩端都可能還沒動。 這一章的狀態
要用兩個數字描述一段：左端 $l$ 和右端 $r$。 這種動態規劃叫\textbf{區間
DP}（interval DP）。

這一章先用作業簿把區間 DP 的四件事走一遍，重點是狀態的意思和計算順序，並估
計 $O(n^3)$ 到底能處理多大的 $n$。 接著三題例題練三種常見的轉移：在區間裡挑
一個切點把它分成兩半、看最左邊的人和誰配成一對、從兩端往內縮。 切點那一節
之後，另外用一節講演算法課本最常拿來介紹區間 DP 的矩陣鏈乘，並示範怎麼從表格
印出最好的做法。 最後處理排成一圈的情況，以及什麼時候不該用區間 DP。

動態規劃的四件事、記憶化和填表的關係，第 1 本第~9~章都講過，這一章直接用。 前綴和
在第 1 本第~5~章。 用到的語法只有二維的 \texttt{vector}（銜接書第~3~章 3.6 節）。
沒讀過銜接書的讀者，只要知道 \texttt{vector<vector<long long>> dp(n, vector<long long>(n, 0))}
會建出一張 $n \times n$、全部是 $0$ 的表格，就能讀這一章。

\section{一段一段的子問題}
\label{sec:ivdp-basic}

\subsection{先試貪心，再寫暴力解}

章首的規則「每次合併加起來最少的相鄰兩疊」，在 $6, 6, 5, 7$ 上是這樣走的：相
鄰兩疊的和是 $12$、$11$、$12$，先合併 $6$ 和 $5$；接著剩下 $6, 11, 7$，和是
$17$、$18$，合併 $6$ 和 $11$；最後 $17$ 和 $7$。 下表把它和另一種搬法並排：

\begin{center}
\begin{tabular}{lccc|c}
\toprule
搬法 & 第 1 次 & 第 2 次 & 第 3 次 & 總力氣 \\
\midrule
每次合併最少的 & $6 + 5 = 11$ & $6 + 11 = 17$ & $17 + 7 = 24$ & $52$ \\
左右各合一半 & $6 + 6 = 12$ & $5 + 7 = 12$ & $12 + 12 = 24$ & $48$ \\
\bottomrule
\end{tabular}
\end{center}

貪心先合併的 $5$ 本跑進了 $11$ 本那一疊，之後每次搬都要把它們再搬一次。 最
後一次一定是把全部 $24$ 本搬在一起，哪種搬法都一樣；差別在前面兩次。

最直接的做法是試遍所有搬法。 有 $n$ 疊時，第一次有 $n - 1$ 對相鄰的可以選，
第二次有 $n - 2$ 對，一直到最後只剩一對，共 $(n - 1)!$ 種搬法。 $n = 20$ 時
大約是 $1.2 \times 10^{17}$ 種，一定來不及。 和第 1 本第~9~章一樣，要找「大問題可
以用小問題的答案算出來」的關係。

\subsection{看最後一次合併}
\label{sec:ivdp-last}

不管怎麼搬，最後一次合併之前桌上剩兩疊。 合併只發生在相鄰的兩疊之間，所以
任何時刻桌上的每一疊，都是原本\textbf{連續的一段}作業簿：左邊那疊是第 $1$ 到
第 $k$ 疊合起來的，右邊那疊是第 $k + 1$ 到第 $n$ 疊合起來的，$k$ 是
$1$ 到 $n - 1$ 之間的某個數。

\begin{center}
\begin{tikzpicture}[x=1cm, y=1cm]
\foreach \v [count=\i] in {6,6,5,7} {
  \node[cell] at (\i,0) {\v};
  \node[idxlbl] at (\i,0.7) {\i};
}
\draw[decorate, decoration={brace, amplitude=5pt, mirror}, thick, draw=memStroke]
  (0.6,-0.6) -- (2.4,-0.6) node[midway, below=6pt, font=\small] {第 $1$～$2$ 疊};
\draw[decorate, decoration={brace, amplitude=5pt, mirror}, thick, draw=optStroke]
  (2.6,-0.6) -- (4.4,-0.6) node[midway, below=6pt, font=\small] {第 $3$～$4$ 疊};
\draw[dashed, black!50] (2.5,0.55) -- (2.5,-1.4);
\node[font=\small, right] at (4.8,0) {最後一次合併：左右兩疊，$k = 2$};
\end{tikzpicture}
\diagcaption{最後一次合併把一排作業簿分成左右兩段。 左段裡怎麼搬，和右段完全
無關。}
\label{fig:ivdp-last}
\end{center}

把 $k$ 固定下來看：

\begin{itemize}
\item 最後一次合併花的力氣是全部的本數，和 $k$ 無關。
\item 在那之前，左段第 $1$ 到 $k$ 疊自己合成一疊，右段第 $k + 1$ 到 $n$ 疊自
己合成一疊。 兩段之間從來沒有合併過，左段怎麼搬不影響右段的花費，反過來也
一樣。 所以兩段都應該用各自最省力的搬法。
\end{itemize}

左段是「第 $1$ 到第 $k$ 疊合成一疊」，右段是「第 $k + 1$ 到第 $n$ 疊合成一疊」。
右段不是從第 $1$ 疊開始的，而它自己再拆一次，又會出現兩端都在中間的一段。 所
以子問題要能描述任意一段：

\begin{center}
\fbox{\parbox{0.86\linewidth}{\centering
$\mathrm{dp}(l, r)$：把第 $l$ 到第 $r$ 疊（連續的一段）合成一疊，最少要花多少力氣。}}
\end{center}

照第 1 本第~9~章的四件事寫下來：

\begin{enumerate}
\item \textbf{狀態}：如上。 $l \le r$，共 $n(n + 1)/2$ 個。
\item \textbf{轉移}：最後一次合併的分界 $k$ 有 $r - l$ 種可能，每一種都算一
次，取最小，再加上這一段的總本數 $S(l, r)$：
\[
\mathrm{dp}(l, r) = \min_{l \le k < r} \bigl(\mathrm{dp}(l, k) + \mathrm{dp}(k + 1, r)\bigr) + S(l, r).
\]
$S(l, r)$ 用第 1 本第~5~章的前綴和算，一次減法就有。
\item \textbf{初始值}：只有一疊時不用搬，$\mathrm{dp}(i, i) = 0$。
\item \textbf{計算順序}：下一小節講。
\end{enumerate}

答案是 $\mathrm{dp}(1, n)$。 這裡的 $k$ 叫做\textbf{切點}（split point），區間
DP 很多題目的轉移都是「枚舉切點，左右兩段各自用最好的做法」。

\subsection{狀態的第五種寫法}
\label{sec:ivdp-shapes}

把第 1 本第~9~章 9.6 節的表加上一列：

\begin{center}
\begin{tabular}{p{0.25\linewidth}p{0.36\linewidth}p{0.29\linewidth}}
\toprule
狀態 & 轉移在看什麼 & 例子 \\
\midrule
前 $i$ 個（或金額 $x$） & 最後一步是哪一種 & 郵票、跳格子（第 1 本第~9~章） \\
前 $i$ 個，而且第 $i$ 個是某種情況 & 第 $i$ 個和第 $i - 1$ 個能不能相鄰 & 招牌燈、大夜班（第 1 本第~9~章） \\
以第 $i$ 個結尾的連續一段 & 接在前一段後面，或重新開始 & 最大區間和（第 1 本第~9~章） \\
前 $i$ 個切成好幾段 & 最後一段從哪裡開始 & 分段（第 1 本第~9~章） \\
\textbf{第 $l$ 到第 $r$ 個這一段} & \textbf{切點在哪裡，或兩端怎麼處理} & \textbf{本章} \\
\bottomrule
\end{tabular}
\end{center}

前四種的子問題都是「前 $i$ 個」或「到第 $i$ 個為止」，只要一個數字。 什麼時
候需要第五種？ 判斷的方法是：把問題拆一次之後，看剩下的子問題長什麼樣子。
如果剩下的總是從頭開始的一段，用前四種；如果剩下的會是夾在中間的一段（像作業
簿的右段），或是做法本來就從兩端往內縮（後面的回文），就要用區間。

第~2~章的 \texttt{f[i][j]} 也有兩個數字，但意思不同：$i$、$j$ 是\textbf{兩
個}序列各取多長的前綴。 區間 DP 的 $l$、$r$ 是\textbf{同一個}序列上一段的兩端。

和第 1 本第~9~章一樣，狀態也可以再多記一點資訊。 例如要沿著一條街走過一段攤位，除了
「已經走過的是哪一段」，還要知道「人現在站在這一段的左端還是右端」，狀態就變
成 $(l, r, \text{哪一端})$。 本章的實作練習有這樣的題目。

\subsection{計算順序：由短到長}
\label{sec:ivdp-order}

第 1 本第~9~章的表格是一排，從左填到右就好。 區間 DP 的表格是一張 $n \times n$ 的表，
只用到 $l \le r$ 的那一半。 先看 $\mathrm{dp}(l, r)$ 會用到哪些格子：$k$ 從
$l$ 到 $r - 1$，用到的 $\mathrm{dp}(l, k)$ 在同一列的左邊，$\mathrm{dp}(k + 1, r)$
在同一行的下面。

\begin{center}
\begin{tikzpicture}[x=1.05cm, y=1.05cm]
\foreach \c in {1,...,4} {\node[idxlbl] at (\c,0.65) {$r = \c$};}
\foreach \rr in {1,...,4} {\node[idxlbl] at (0.2,-\rr+1) {$l = \rr$};}
\foreach \l/\r/\v in {2/2/0, 3/3/0, 2/3/11} {
  \node[cell] at (\r,-\l+1) {\v};
}
\foreach \l/\r/\v in {1/1/0, 1/2/12, 1/3/28, 2/4/29, 3/4/12, 4/4/0} {
  \node[cell, fill=coreFill, draw=coreStroke] at (\r,-\l+1) {\v};
}
\node[hlcell] at (4,0) {48};
\node[font=\small, right, align=left] at (4.8,-0.4)
  {$\mathrm{dp}(1, 4)$ 只用到\\同一列左邊、同一行下面\\的格子（藍色），\\它們都比較短};
\node[font=\small, right, align=left] at (4.8,-2.6)
  {同一條斜線上的格子\\長度相同};
\end{tikzpicture}
\diagcaption{作業簿 $6, 6, 5, 7$ 的表格。 第 $l$ 列第 $r$ 行是 $\mathrm{dp}(l, r)$，
左下半部用不到。 主對角線是長度 $1$，往右上每一條斜線長度加 $1$。}
\label{fig:ivdp-table}
\end{center}

$\mathrm{dp}(l, k)$ 和 $\mathrm{dp}(k + 1, r)$ 都比 $[l, r]$ 短。 所以只要
\textbf{依區間長度由短到長}填，算到某一格時，它需要的格子一定都算好了。
圖上同一條斜線的格子長度相同，這個順序等於把斜線依序填滿，右上角的
$\mathrm{dp}(1, n)$ 最後才填。

最自然的寫法「$l$ 由小到大、$r$ 由小到大」是錯的。 它先把第 $1$ 列整列填完，
可是 $\mathrm{dp}(1, 4)$ 要用 $\mathrm{dp}(2, 4)$，那在第 $2$ 列，還沒填。 沒
填的格子還是 $0$，程式會把「右段不用花力氣」當真，算出太小的答案（見本節的常
見錯誤）。 第 1 本第~10~章的格子一列一列填沒有問題，是因為那裡每一格只用到上面和左
邊的格子；這裡用到的是\textbf{下面}的格子。

由短到長以外，還有別的順序也對，例如 $l$ 由大到小、$r$ 由小到大：第 $l$ 列
需要的下面那幾列早就填好，同一列左邊的格子也先填了。 只要每一格用到的格子都
先算好，順序就是對的，本章練習第~2~題會再看幾種。 本書統一用「由短到長」，
因為它不用想，對每一種區間 DP 都成立。

不想安排順序時，也可以寫成第 1 本第~9~章的記憶化遞迴：要 $\mathrm{dp}(l, r)$ 就呼叫
它，它再去呼叫比較短的區間。 每呼叫一層區間至少短 $1$，遞迴最多大約 $n$ 層，
不用擔心太深。

\begin{example}[label=ex:ivdp-merge]{cpp}{(merge.cpp) 合併一排作業簿的最少力氣}
#include <climits>
#include <iostream>
#include <vector>
using namespace std;

int main() {
    int n;
    cin >> n;
    vector<long long> pre(n + 1, 0);          // pre[i]：前 i 疊的本數總和
    for (int i = 1; i <= n; i++) {
        long long x;
        cin >> x;
        pre[i] = pre[i - 1] + x;
    }
    // dp[l][r]：把第 l 到第 r 疊合成一疊的最少力氣
    vector<vector<long long>> dp(n + 1, vector<long long>(n + 1, 0));
    for (int len = 2; len <= n; len++) {      // 區間長度由短到長
        for (int l = 1; l + len - 1 <= n; l++) {
            int r = l + len - 1;
            long long best = LLONG_MAX;
            for (int k = l; k < r; k++) {     // 最後一次合併：[l, k] 和 [k + 1, r]
                best = min(best, dp[l][k] + dp[k + 1][r]);
            }
            dp[l][r] = best + (pre[r] - pre[l - 1]);
        }
    }
    cout << dp[1][n] << '\n';
    return 0;
}
\end{example}

\iosample{
4 \\
6 6 5 7 \\
}{
48 \\
}

\begin{itemize}
\item 疊的編號從 $1$ 開始，\texttt{pre[0] = 0}，第 $l$ 到第 $r$ 疊的總本數是
\texttt{pre[r] - pre[l - 1]}，$l = 1$ 時也不用特別處理。
\item 三層迴圈的外層是長度，中間是左端，右端由這兩個算出來。 這三行是區間 DP
最常見的骨架，後面的例題都用它。
\item \texttt{dp[i][i]} 開表時就是 $0$，所以長度從 $2$ 開始。 長度至少 $2$ 時
$k$ 至少有一個可以選，\texttt{best} 一定會被更新，不會拿 \texttt{LLONG\_MAX}
去加總本數而溢位。
\end{itemize}

填表的過程（圖~\ref{fig:ivdp-table} 就是填完的樣子）：

\begin{center}
\begin{tabular}{cccl}
\toprule
長度 & 區間 & 本數 & 各個切點 $k$ 的 $\mathrm{dp}(l, k) + \mathrm{dp}(k + 1, r)$ \\
\midrule
2 & $[1, 2]$ & 12 & $k = 1$：$0 + 0$ \quad 所以 $\mathrm{dp} = 12$ \\
2 & $[2, 3]$ & 11 & $k = 2$：$0 + 0$ \quad 所以 $\mathrm{dp} = 11$ \\
2 & $[3, 4]$ & 12 & $k = 3$：$0 + 0$ \quad 所以 $\mathrm{dp} = 12$ \\
3 & $[1, 3]$ & 17 & $k = 1$：$0 + 11$；$k = 2$：$12 + 0$ \quad 所以 $\mathrm{dp} = 28$ \\
3 & $[2, 4]$ & 18 & $k = 2$：$0 + 12$；$k = 3$：$11 + 0$ \quad 所以 $\mathrm{dp} = 29$ \\
4 & $[1, 4]$ & 24 & $k = 1$：$0 + 29$；$k = 2$：$12 + 12$；$k = 3$：$28 + 0$ \quad 所以 $\mathrm{dp} = 48$ \\
\bottomrule
\end{tabular}
\end{center}

$\mathrm{dp}(1, 4)$ 在 $k = 2$ 取到最小，也就是表格上「左右各合一半」的搬法。

\subsection{$O(n^3)$ 能處理多大的 $n$}
\label{sec:ivdp-cost}

照第 1 本第~9~章的估法：狀態數乘上每個狀態要看的選擇數。 長度為 $len$ 的區間有
$n - len + 1$ 個，每個要試 $len - 1$ 個切點，所以最內層那一行一共執行
\[
\sum_{len = 2}^{n} (n - len + 1)(len - 1) = \frac{n^3 - n}{6}
\]
次，大約是 $n^3 / 6$。 狀態有 $n^2/2$ 個左右，每個平均看 $n/3$ 個切點，乘起
來也是這個量級。 這裡的 $1/6$ 不能丟：它讓 $O(n^3)$ 的程式實際上比「$n^3$ 次」
快好幾倍。

下表是範例~\ref{ex:ivdp-merge} 在隨機輸入上的實測，表格用 \texttt{long long}：

\begin{center}
\begin{tabular}{rrrr}
\toprule
$n$ & 最內層執行次數 & 表格大小 & 在本書測試的環境下 \\
\midrule
$200$ & $1.3 \times 10^{6}$ & $0.3$ MB & 不到 $0.01$ 秒 \\
$500$ & $2.1 \times 10^{7}$ & $2$ MB & 約 $0.01$ 秒 \\
$1000$ & $1.7 \times 10^{8}$ & $8$ MB & 約 $0.1$ 秒 \\
$2000$ & $1.3 \times 10^{9}$ & $32$ MB & 約 $0.9$ 秒 \\
\bottomrule
\end{tabular}
\end{center}

$n$ 變成兩倍，時間變成八、九倍，這是 $n^3$ 的樣子。 記憶體是 $n^2$ 格，
$n = 5000$ 時光是一張 \texttt{long long} 的表就要 $200$ MB。 所以看到資料範
圍時可以這樣判斷：

\begin{itemize}
\item $n$ 在幾百以內：每個狀態枚舉一個切點的 $O(n^3)$ 很安全。
\item $n$ 到一千左右：$O(n^3)$ 還來得及，但轉移只要多做幾件事就會變慢，先估
一次最內層的次數。
\item $n$ 到兩千以上：只能讓每個狀態用 $O(1)$ 個選擇，總共 $O(n^2)$，例如第
~\ref{sec:ivdp-pal}~節的回文。 也要先算表格的記憶體。
\item $n$ 到 $10^5$：連 $n^2$ 個狀態都存不下，通常不是區間 DP，要換想法。
\end{itemize}

\begin{apcsnote}
題目會寫出資料範圍和記憶體限制，動手前先照上面的方法算一次：最內層大約執行
幾次、表格要幾 MB。 如果題目把測試資料分成幾個子題，$n$ 很小的子題可以先用
暴力解（試遍所有順序）拿下來，這個暴力解之後也可以拿來和區間 DP 對拍。
\end{apcsnote}

\subsection{常見錯誤}

\begin{itemize}
\item \textbf{$l$、$r$ 都由小到大}：把外兩層改成 \texttt{for (l = 1; l <= n; l++)}、
\texttt{for (r = l + 1; r <= n; r++)}。 範例的輸入印出 \texttt{24}，正確答案
是 $48$。 算 $\mathrm{dp}(1, 4)$ 時第 $2$、$3$ 列還沒填，全是 $0$，程式以
為「$k = 1$：左段 $0$、右段 $0$」，只算了最後一次合併的 $24$。
\item \textbf{忘了加這一段的總本數}：轉移只寫 \texttt{dp[l][k] + dp[k + 1][r]}，
每一格都是 $0$，印出 \texttt{0}。 每個狀態「這一步」的花費要記得加。
\item \textbf{取最小值的起點設成 $0$}：把 \texttt{best} 的初始值寫成
\texttt{0}，\texttt{min} 永遠是 $0$，每一格都只剩這一段的總本數，範例的輸
入印出 \texttt{24}。 取最小值時起點要夠大，例如 \texttt{LLONG\_MAX}，或直接
用第一個切點的值當起點。
\end{itemize}

\section{切點枚舉：先做哪一件}
\label{sec:ivdp-split}

作業簿看的是\textbf{最後}一次合併：它把一段分成左右兩段，兩段各自獨立。 有
些題目反過來，看\textbf{第一}個動作：第一刀切下去，或第一個東西拿掉，一整段
就斷成左右兩段，之後兩邊互不影響。 兩種看法的目的相同：挑一個動作，讓剩下的
問題變成兩段連續、互相獨立的區間。

\subsection{例題：拆除一排相連的攤位}
\label{sec:ex-ivdp-demolish}

\textbf{題意}（\akiref{demolishing-a-row-of-stalls}{拆除一排相連的攤位}）

\begin{itemize}
\item 一排 $n$ 間相連的組合屋攤位，第 $i$ 間的震動係數是 $a_i$。 還沒拆的攤
位裡，連續相鄰、中間沒有空位的一段叫一個區塊，一開始整排是一個區塊。
\item 每次選一間還沒拆的攤位拆掉，花費是 $a_k \times L$，$L$ 是拆的當下它所
在區塊的攤位數。 拆完後區塊在這裡斷開，變成左右兩個區塊（一邊可能是空的），
不會再接回去。
\item 求拆完全部攤位的最小總花費。 $n \le 150$，$1 \le a_i \le 1000$。
\end{itemize}

例如 $a = 3, 1, 4, 2$：先拆第 $2$ 間，整排 $4$ 間是一個區塊，花 $1 \times 4$；
剩下 $\{1\}$ 和 $\{3, 4\}$。 再拆第 $1$ 間花 $3 \times 1$，拆第 $4$ 間花
$2 \times 2$（區塊 $\{3, 4\}$ 有兩間），拆第 $3$ 間花 $4 \times 1$。 總共
$4 + 3 + 4 + 4 = 15$，這是最小的。

\textbf{想法}

最直接的做法是試遍拆的順序，有 $n!$ 種，$n = 150$ 完全不可能。 但它仍然有
用：$n \le 8$ 時用 \texttt{next\_permutation} 枚舉順序、照規則模擬，就是對
拍用的暴力解。

兩條直覺的貪心都會錯。 「每次先拆係數最大的」錯得最明顯：係數大的攤位在區塊
還很長的時候拆，乘上的 $L$ 最大。 上面的例子先拆 $4$ 要付 $4 \times 4 = 16$，
總共 $25$。 「每次先拆係數最小的」聽起來比較合理，但拆的位置同時決定了區塊
怎麼斷開。 $a = 1, 2, 3, 4$ 時依序拆第 $1$、$2$、$3$、
$4$ 間，花費 $1 \times 4 + 2 \times 3 + 3 \times 2 + 4 \times 1 = 20$；先拆第
$2$ 間，整排斷成 $\{1\}$ 和 $\{3, 4\}$，總共只要
$2 \times 4 + 1 + 3 \times 2 + 4 = 19$。

\textbf{看最後拆的那間行不行？} 學過作業簿，第一個念頭可能是看「最後一個拆的
攤位」。 可是知道最後拆誰，什麼都沒切開：在它之前，區塊怎麼斷、每次斷成多
長，全部要看前面的順序。 子問題沒辦法用一段區間描述。

\textbf{看第一個拆的。} 整排 $[1, n]$ 第一個拆的是第 $k$ 間：

\begin{itemize}
\item 這一下的花費是 $a_k \times n$，因為當下整排是同一個區塊。
\item 拆完之後，$[1, k - 1]$ 和 $[k + 1, n]$ 是兩個區塊，以後永遠不會接回去。
左邊的攤位怎麼拆，只影響左邊區塊的長度，和右邊無關。 兩邊各自用最好的拆法。
\end{itemize}

右邊區塊 $[k + 1, n]$ 再拆一次，又會出現兩端都在中間的區塊，所以狀態是：

\begin{center}
\fbox{\parbox{0.86\linewidth}{\centering
$\mathrm{dp}(l, r)$：第 $l$ 到第 $r$ 間剛好是一個完整的區塊時，把它拆完的最小花費。}}
\end{center}

區塊 $[l, r]$ 裡第一個拆的是第 $k$ 間（$l \le k \le r$），當下區塊有
$r - l + 1$ 間：
\[
\mathrm{dp}(l, r) = \min_{l \le k \le r} \bigl(a_k \times (r - l + 1) + \mathrm{dp}(l, k - 1) + \mathrm{dp}(k + 1, r)\bigr).
\]
和作業簿相比有兩點不同。 第一，切點 $k$ 本身被拆掉，左右兩段是 $[l, k - 1]$
和 $[k + 1, r]$，不含 $k$。 第二，$k = l$ 或 $k = r$ 時其中一段是空的，空區
塊不用花錢，$\mathrm{dp}(l, l - 1) = \mathrm{dp}(r + 1, r) = 0$。 一間的區塊
$\mathrm{dp}(i, i) = a_i$，套進式子也對（$k$ 只能是 $i$，兩邊都空），不用另
外寫。

\begin{center}
\begin{tikzpicture}[x=1cm, y=1cm]
\foreach \v [count=\i] in {3,1,4,2} {
  \node[idxlbl] at (\i,0.7) {\i};
}
\node[cell] at (1,0) {3};
\node[hlcell] at (2,0) {1};
\node[cell] at (3,0) {4};
\node[cell] at (4,0) {2};
\node[font=\small] at (2,-0.8) {先拆，$1 \times 4$};
\draw[decorate, decoration={brace, amplitude=4pt, mirror}, thick, draw=coreStroke]
  (0.6,-1.3) -- (1.4,-1.3) node[midway, below=5pt, font=\small] {$\mathrm{dp}(1, 1)$};
\draw[decorate, decoration={brace, amplitude=5pt, mirror}, thick, draw=coreStroke]
  (2.6,-1.3) -- (4.4,-1.3) node[midway, below=6pt, font=\small] {$\mathrm{dp}(3, 4)$};
\node[font=\small, right] at (5,0) {$\mathrm{dp}(1, 4)$ 在 $k = 2$ 的候選：};
\node[font=\small, right] at (5,-0.7) {$4 + \mathrm{dp}(1, 1) + \mathrm{dp}(3, 4) = 4 + 3 + 8 = 15$};
\end{tikzpicture}
\diagcaption{區塊 $[1, 4]$ 先拆第 $2$ 間，剩下 $[1, 1]$ 和 $[3, 4]$ 兩個獨立
的區塊。 $\mathrm{dp}(3, 4) = 8$：先拆第 $4$ 間（$2 \times 2$）再拆第 $3$ 間
（$4$）。}
\label{fig:ivdp-demolish}
\end{center}

計算順序仍然是長度由短到長：$[l, k - 1]$ 和 $[k + 1, r]$ 都比 $[l, r]$ 短。

同樣的模型也出現在別的地方：一條生產線上相連的機台要逐台停機保養，停一台會
影響同一段還在運轉的機台；一根木料要在指定的位置鋸開，每一刀的成本和當下那
段木料的長度有關。 只要「第一個動作把一段切成兩段、兩段從此無關」，就可以
枚舉第一個動作的位置。

\textbf{程式}

\begin{example}[label=ex:ivdp-demolish]{cpp}{(demolition.cpp) 拆除一排相連的攤位}
#include <climits>
#include <iostream>
#include <vector>
using namespace std;

int main() {
    int n;
    cin >> n;
    vector<int> a(n + 2, 0);
    for (int i = 1; i <= n; i++) {
        cin >> a[i];
    }
    // dp[l][r]：第 l 到第 r 間是一個完整的區塊時，把它拆完的最小花費
    // l > r 表示空的區塊，花費 0。開 n + 2 列，dp[l][l - 1]、dp[r + 1][r] 都在表內
    vector<vector<int>> dp(n + 2, vector<int>(n + 2, 0));
    for (int len = 1; len <= n; len++) {
        for (int l = 1; l + len - 1 <= n; l++) {
            int r = l + len - 1;
            int best = INT_MAX;
            for (int k = l; k <= r; k++) {    // 這個區塊第一個拆第 k 間
                int cost = a[k] * len + dp[l][k - 1] + dp[k + 1][r];
                best = min(best, cost);
            }
            dp[l][r] = best;
        }
    }
    cout << dp[1][n] << '\n';
    return 0;
}
\end{example}

\iosample{
4 \\
3 1 4 2 \\
}{
15 \\
}

\begin{itemize}
\item 表格開 $n + 2$ 列、$n + 2$ 行。 $k = r = n$ 時會讀 \texttt{dp[n + 1][n]}，
$k = l = 1$ 時會讀 \texttt{dp[1][0]}，這些空區塊的格子都在表內，而且開表時就
是 $0$。 空區間是區間 DP 最常越界的地方，多開一圈最省事。
\item 這題長度從 $1$ 開始，因為一間的區塊也要花錢。
\item 答案最大大約是 $150$ 間 $\times$ 每間最多 $1000 \times 150$，約
$2.3 \times 10^7$，\texttt{int} 裝得下。 動手前先估一次，就知道要不要用
\texttt{long long}。
\end{itemize}

\textbf{複雜度}

最內層執行 $\sum_{len=1}^{n} (n - len + 1) \cdot len = n(n + 1)(n + 2)/6$ 次，
$n = 150$ 時是 $573800$ 次。 最大的測試資料在本書測試的環境下不到 $0.01$ 秒。

\subsection{常見錯誤}

\begin{itemize}
\item \textbf{用貪心}：每次先拆係數最小的，輸入 \texttt{4}、\texttt{1 2 3 4}
印出 \texttt{20}，正確答案是 $19$；每次先拆係數最大的，範例的輸入印出
\texttt{25}。 兩條規則都只看係數，沒有考慮拆的位置會怎麼切開區塊。
\item \textbf{區塊長度用整排的 $n$}：把 \texttt{a[k] * len} 寫成
\texttt{a[k] * n}，範例的輸入印出 \texttt{40}（每間都乘 $4$）。 斷開的兩段
不會接回去，拆的時候看的是當下所在區塊的長度。
\item \textbf{計算順序錯}：外兩層改成 $l$、$r$ 都由小到大，範例的輸入印出
\texttt{7}。 算 $\mathrm{dp}(1, 4)$ 時 \texttt{dp[3][4]}、\texttt{dp[2][4]}
都還是 $0$。
\item \textbf{表格開太小}：開成 $n + 1$ 列，\texttt{dp[n + 1][n]} 就在表外。
在本書測試的環境下，範例的輸入直接當掉（segmentation fault）；加上
\texttt{-fsanitize=address} 編譯，會指出讀到了陣列外面。 這是未定義行為，換
一台機器或換一組輸入，也可能不當掉、印出奇怪的數字。
\end{itemize}

\section{矩陣鏈乘}
\label{sec:ivdp-matrix}

演算法課本講區間 DP，幾乎都從\textbf{矩陣鏈乘}（matrix-chain multiplication）
開始。 它和作業簿一樣看最後一次合併，差別在最後一次的花費會隨切點改變。 這一
節用它把切點的寫法再走一遍，並補上前面沒做的一件事：除了最小的花費，也把最好
的做法印出來。

\subsection{乘的順序會影響花費}

沒學過矩陣也能讀這一節，只要知道兩條規則。 一個 $p \times q$ 的\textbf{矩陣}
是排成 $p$ 列、$q$ 行的一張數字表。

\begin{itemize}
\item $p \times q$ 的矩陣只能乘上 $q \times r$ 的矩陣（前一個的行數等於後一個的
列數），結果是 $p \times r$。
\item 結果的每一格，是拿前一個的一列和後一個的一行逐項相乘再加起來，要做 $q$
次乘法。 結果有 $p \times r$ 格，所以一共做 $p \cdot q \cdot r$ 次乘法。
\end{itemize}

矩陣乘法可以任意加括號，$(AB)C$ 和 $A(BC)$ 算出來的矩陣相同，但做的乘法次數
可能差很多。 設 $A$ 是 $2 \times 20$、$B$ 是 $20 \times 3$、$C$ 是
$3 \times 15$：

\begin{center}
\begin{tabular}{lllc}
\toprule
算法 & 第一次 & 第二次 & 總次數 \\
\midrule
$(AB)C$ & $AB$：$2 \cdot 20 \cdot 3 = 120$，得到 $2 \times 3$ & $2 \cdot 3 \cdot 15 = 90$ & $210$ \\
$A(BC)$ & $BC$：$20 \cdot 3 \cdot 15 = 900$，得到 $20 \times 15$ & $2 \cdot 20 \cdot 15 = 600$ & $1500$ \\
\bottomrule
\end{tabular}
\end{center}

先乘 $AB$，中間產物只有 $2 \times 3$；先乘 $BC$，會先做出一張 $20 \times 15$
的大表，下一次乘法也跟著貴。 結果一樣，次數差了七倍多。 電腦繪圖和科學計算
常要把一長串矩陣乘起來，矩陣一大，順序選得好不好，計算時間會差好幾倍。

問題是這樣：$n$ 個矩陣 $A_1, A_2, \ldots, A_n$ 要依序相乘，第 $i$ 個是
$d_{i-1} \times d_i$。 相鄰兩個的維度本來就要接得上，所以 $n + 1$ 個數字
$d_0, d_1, \ldots, d_n$ 就描述了全部。 矩陣的先後不能換，只能決定括號怎麼加。
求最少的乘法次數，並印出一種達到它的加括號方式。

最直接的做法是試遍所有加括號的方式。 $3$ 個矩陣有 $2$ 種，$4$ 個有 $5$ 種，
$10$ 個有 $4862$ 種，$20$ 個超過 $17$ 億種（這串數叫\textbf{卡特蘭數}，Catalan
number），只能拿來對拍。

貪心也不行。 以 $d = 5, 8, 6, 2, 7$ 為例，四個矩陣是 $5 \times 8$、
$8 \times 6$、$6 \times 2$、$2 \times 7$。 「每次先做目前最便宜的那一次相鄰
乘法」會先乘 $A_3 A_4$（$6 \cdot 2 \cdot 7 = 84$），接著 $A_1 A_2$（$240$），
最後 $5 \cdot 6 \cdot 7 = 210$，總共 $534$，比老老實實從左乘到右的 $370$ 還
多。 最好的做法是 $(A_1 (A_2 A_3)) A_4$，只要 $96 + 80 + 70 = 246$。 第一步
貪便宜，把最小的維度 $2$ 消掉了；留著它，後面每一次乘法都會乘上這個 $2$，反
而便宜。

\subsection{狀態與轉移}

和作業簿一樣看\textbf{最後一次}乘法。 把 $A_l$ 到 $A_r$ 乘起來，最後一次一定
是「左邊一串的乘積」乘上「右邊一串的乘積」：$(A_l \cdots A_k)(A_{k+1} \cdots A_r)$。
左邊的乘積是 $d_{l-1} \times d_k$，右邊是 $d_k \times d_r$，這一次乘法花
$d_{l-1} \cdot d_k \cdot d_r$。 左右兩串各自怎麼加括號互不影響，都用最好的。

\begin{center}
\fbox{\parbox{0.86\linewidth}{\centering
$\mathrm{dp}(l, r)$：把 $A_l$ 到 $A_r$ 乘起來，最少要做幾次乘法。}}
\end{center}

\[
\mathrm{dp}(l, r) = \min_{l \le k < r} \bigl(\mathrm{dp}(l, k) + \mathrm{dp}(k + 1, r) + d_{l-1} \cdot d_k \cdot d_r\bigr),
\qquad \mathrm{dp}(i, i) = 0.
\]

\begin{center}
\begin{tikzpicture}[x=1.5cm, y=1cm]
\foreach \v [count=\i] in {5 \times 8, 8 \times 6, 6 \times 2, 2 \times 7} {
  \node[cell, minimum width=1.3cm] at (\i,0) {$\v$};
  \node[idxlbl] at (\i,0.7) {$A_{\i}$};
}
\draw[decorate, decoration={brace, amplitude=5pt, mirror}, thick, draw=memStroke]
  (0.55,-0.6) -- (3.45,-0.6) node[midway, below=6pt, font=\small] {$A_1 A_2 A_3$：$5 \times 2$};
\draw[decorate, decoration={brace, amplitude=5pt, mirror}, thick, draw=optStroke]
  (3.55,-0.6) -- (4.45,-0.6) node[midway, below=6pt, font=\small] {$2 \times 7$};
\draw[dashed, black!50] (3.5,0.55) -- (3.5,-1.4);
\node[font=\small, right] at (4.7,0) {$k = 3$：最後一次花 $5 \cdot 2 \cdot 7 = 70$};
\end{tikzpicture}
\diagcaption{最後一次乘法切在 $k = 3$。 左邊三個的乘積是 $d_0 \times d_3 = 5 \times 2$，
右邊是 $d_3 \times d_4 = 2 \times 7$。}
\label{fig:ivdp-matrix}
\end{center}

和作業簿比，差別只有一處：作業簿最後一次的花費是整段的總本數，和 $k$ 無關，
可以放在 $\min$ 外面；這裡的 $d_{l-1} \cdot d_k \cdot d_r$ 跟著 $k$ 變，要放
在裡面一起比。 計算順序仍然是長度由短到長，迴圈和範例~\ref{ex:ivdp-merge}
一模一樣，時間 $O(n^3)$（最內層大約 $n^3/6$ 次），表格 $O(n^2)$。 $n$ 到幾百
都很快。

以 $d = 5, 8, 6, 2, 7$ 填表：

\begin{center}
\begin{tabular}{ccl}
\toprule
長度 & 區間 & 各個切點 $k$ 的 $\mathrm{dp}(l, k) + \mathrm{dp}(k + 1, r) + d_{l-1} d_k d_r$ \\
\midrule
2 & $[1, 2]$ & $k = 1$：$0 + 0 + 240$ \quad 所以 $\mathrm{dp} = 240$ \\
2 & $[2, 3]$ & $k = 2$：$0 + 0 + 96$ \quad 所以 $\mathrm{dp} = 96$ \\
2 & $[3, 4]$ & $k = 3$：$0 + 0 + 84$ \quad 所以 $\mathrm{dp} = 84$ \\
3 & $[1, 3]$ & $k = 1$：$0 + 96 + 80 = 176$；$k = 2$：$240 + 0 + 60 = 300$ \quad 取 $k = 1$ \\
3 & $[2, 4]$ & $k = 2$：$0 + 84 + 336 = 420$；$k = 3$：$96 + 0 + 112 = 208$ \quad 取 $k = 3$ \\
4 & $[1, 4]$ & $k = 1$：$488$；$k = 2$：$240 + 84 + 210 = 534$；$k = 3$：$176 + 0 + 70 = 246$ \\
\bottomrule
\end{tabular}
\end{center}

$k = 1$ 那一格是 $0 + 208 + 5 \cdot 8 \cdot 7 = 488$。 答案 $\mathrm{dp}(1, 4) = 246$，
在 $k = 3$ 取到。

\subsection{還原加括號的方式}

表格只算出最少 $246$ 次，沒有說括號怎麼加。 做法和第~2~章從表格還原一樣：
填表時多開一張表 $\mathrm{cut}(l, r)$，記下 $\mathrm{dp}(l, r)$ 是在哪個 $k$
取到最小的。 印的時候從整段 $[1, n]$ 開始：

\begin{itemize}
\item 只有一個矩陣（$l = r$），印 $A_l$。
\item 否則 $k = \mathrm{cut}(l, r)$，印「左括號、$[l, k]$ 的寫法、$[k + 1, r]$
的寫法、右括號」，左右兩段再用同樣的方法印。
\end{itemize}

範例的表裡 $\mathrm{cut}(1, 4) = 3$，所以整段是 $(\,[1, 3]\ A_4)$；
$\mathrm{cut}(1, 3) = 1$，$[1, 3]$ 是 $(A_1\ [2, 3])$；$\mathrm{cut}(2, 3) = 2$，
$[2, 3]$ 是 $(A_2\ A_3)$。 合起來就是 $((A_1 (A_2 A_3)) A_4)$。 遞迴每一層區間至
少短 $1$，最多 $n$ 層。

不開 $\mathrm{cut}$ 也可以：印 $[l, r]$ 時重新試每個 $k$，找一個讓
$\mathrm{dp}(l, k) + \mathrm{dp}(k + 1, r) + d_{l-1} d_k d_r$ 剛好等於
$\mathrm{dp}(l, r)$ 的。 多花一點時間，省一張表。

\begin{example}[label=ex:ivdp-matrix]{cpp}{(matrix\_chain.cpp) 矩陣鏈乘的最少乘法次數與加括號方式}
#include <climits>
#include <iostream>
#include <string>
#include <vector>
using namespace std;

vector<vector<int>> cut;   // cut[l][r]：dp[l][r] 取到最小時，最後一次乘法的切點

// 照 cut 表寫出 A_l 到 A_r 的加括號方式
string paren(int l, int r) {
    if (l == r) {
        return "A" + to_string(l);
    }
    int k = cut[l][r];
    return "(" + paren(l, k) + " " + paren(k + 1, r) + ")";
}

int main() {
    int n;
    cin >> n;
    vector<long long> d(n + 1);   // 第 i 個矩陣是 d[i - 1] x d[i]
    for (int i = 0; i <= n; i++) {
        cin >> d[i];
    }
    // dp[l][r]：把 A_l 到 A_r 乘起來的最少乘法次數
    vector<vector<long long>> dp(n + 1, vector<long long>(n + 1, 0));
    cut.assign(n + 1, vector<int>(n + 1, 0));
    for (int len = 2; len <= n; len++) {
        for (int l = 1; l + len - 1 <= n; l++) {
            int r = l + len - 1;
            dp[l][r] = LLONG_MAX;
            for (int k = l; k < r; k++) {     // 最後一次：(A_l..A_k)(A_k+1..A_r)
                long long cost = dp[l][k] + dp[k + 1][r] + d[l - 1] * d[k] * d[r];
                if (cost < dp[l][r]) {
                    dp[l][r] = cost;
                    cut[l][r] = k;
                }
            }
        }
    }
    cout << dp[1][n] << '\n';
    cout << paren(1, n) << '\n';
    return 0;
}
\end{example}

\iosample{
4 \\
5 8 6 2 7 \\
}{
246 \\
((A1 (A2 A3)) A4) \\
}

\begin{itemize}
\item 輸入先給 $n$，再給 $n + 1$ 個維度。 \texttt{d} 從 $0$ 開始編號、矩陣從
$1$ 開始，第 $l$ 個矩陣的列數是 \texttt{d[l - 1]}，這個 $-1$ 是最常寫錯的地方。
\item 乘法次數很容易超過 \texttt{int}：維度到 $1000$ 時，一次乘法就可能是
$10^9$ 次。 \texttt{d} 直接開 \texttt{long long}，三個數相乘時就不會先在
\texttt{int} 裡溢位。
\item 用嚴格的 \texttt{<} 更新，幾個切點一樣好時留下最小的 $k$。 最好的加括號
方式可能不只一種，題目要求特定的那一種時，要照題目的規定挑。
\end{itemize}

同樣的模型也出現在別的地方：把一個凸多邊形切成三角形，每個三角形的花費由三
個頂點決定；資料庫要把好幾張表依序合併，每次合併的成本看兩邊的大小。 只要
是「一排東西只能相鄰合併，合併的花費看左右兩塊的邊界」，就可以照這一節的寫法。

\subsection{常見錯誤}

\begin{itemize}
\item \textbf{維度差一}：把 \texttt{d[l - 1]} 寫成 \texttt{d[l]}，範例的輸入印出
\texttt{312}，正確答案是 $246$。 程式把第 $l$ 個矩陣的行數當成列數，算的是另
一組矩陣。 手算一個 $n = 2$ 的例子（答案就是 $d_0 d_1 d_2$）就能抓到。
\item \textbf{最後一次的花費放在 $\min$ 外面}：照作業簿的寫法，先找出
\texttt{dp[l][k] + dp[k + 1][r]} 最小的 $k$，再加上那個 $k$ 的
$d_{l-1} d_k d_r$。 輸入 \texttt{3}、\texttt{2 1 3 4} 印出 \texttt{30}，正確答
案是 $20$：$(A_1 A_2)$ 只花 $6$，所以選了 $k = 2$，但最後一次要花
$2 \cdot 3 \cdot 4 = 24$；$k = 1$ 的 $12 + 2 \cdot 1 \cdot 4 = 20$ 才是最好的。
範例的輸入碰巧印出正確的 $246$，只拿範例測不出這個錯。
\end{itemize}

\section{括號型：最左邊的和誰一組}
\label{sec:ivdp-pair}

有一類題目，最後的結果是「每個東西和另一個東西配成一對」，而且配對像括號一樣
不交叉：$($ 和 $)$ 可以一層包一層，也可以並排，但不會出現 $( [ ) ]$ 這種
交錯。 這類題目的轉移看區間最左邊的那一個：它和誰配成一對，那一對就把區間分
成「被包在裡面的」和「在右邊的」兩段。

\subsection{例題：兩兩相鄰離隊搭車}
\label{sec:ex-ivdp-pairs}

\textbf{題意}（\akiref{leaving-the-line-in-pairs}{兩兩相鄰離隊搭車}）

\begin{itemize}
\item $n$ 位同學排成一排（$n$ 是偶數），第 $i$ 位的身高是 $h_i$。
\item 每次從目前的隊伍裡叫兩位\textbf{相鄰}的同學一起離隊，這一車的不協調
度是 $|h_i - h_j|$。 兩人離開後隊伍往中間靠攏，原本在他們兩邊的人變成相鄰。
\item 重複到所有人都離隊，求不協調度總和的最小值。 $n \le 400$，
$h_i \le 10^9$。
\end{itemize}

例如身高 $3, 8, 2, 2, 9, 4$：先叫中間兩位 $2$（差 $0$），隊伍變成
$3, 8, 9, 4$；再叫 $8$ 和 $9$（差 $1$）；最後 $3$ 和 $4$（差 $1$），總和 $2$。

\textbf{想法}

最直接的是模擬每一種叫人的順序：一開始有 $n - 1$ 對相鄰的可以叫，之後
$n - 3$ 對……全部乘起來，$n = 20$ 時已經超過 $6 \times 10^8$ 種，只能拿來在
很小的 $n$ 對拍。 「每次叫目前身高差最小的相鄰兩人」也會錯：身高
$1, 5, 6, 10$ 時先叫 $5$ 和 $6$（差 $1$），剩下 $1$ 和 $10$ 只能一起走（差
$9$），總和 $10$；先叫 $1, 5$ 再叫 $6, 10$ 只要 $8$。

\textbf{順序不重要，重要的是誰和誰一組。} 總和只看最後的配對，和叫人的順序
無關。 那麼哪些配對做得到？ 把每一對畫成一條弧，連起兩人原本的位置：

\begin{center}
\begin{tikzpicture}[x=1cm, y=1cm]
\foreach \v [count=\i] in {3,8,2,2,9,4} {
  \node[cell] (p\i) at (\i,0) {\v};
  \node[idxlbl] at (\i,-0.7) {\i};
}
\draw[thick, memStroke] (p1.north) to[out=70, in=110] (p6.north);
\draw[thick, coreStroke] (p2.north) to[out=70, in=110] (p5.north);
\draw[thick, optStroke] (p3.north) to[out=70, in=110] (p4.north);
\node[font=\small] at (8.6,0.9) {$(\,(\,(\,)\,)\,)$};
\node[font=\small] at (8.6,0.3) {左端寫 (，右端寫 )};
\end{tikzpicture}
\diagcaption{範例的配對：$(1, 6)$、$(2, 5)$、$(3, 4)$。 三條弧一層包一層，
寫成括號就是 \texttt{((()))}。}
\label{fig:ivdp-arcs}
\end{center}

兩條弧不會交叉。 假設 $i < p < j < q$，$i$ 和 $j$ 一組、$p$ 和 $q$ 一組。 $i$、
$j$ 離隊時兩人必須相鄰，夾在中間的 $p$ 已經先走了；$p$、$q$ 離隊時，夾在中
間的 $j$ 也已經先走了。 兩組都比對方先走，不可能。 反過來，任何不交叉的配
對都做得到：最內層的那一對中間沒有人，先叫它們；拿掉之後，剩下的配對仍然不
交叉，又有新的最內層。 所以問題變成：在所有\textbf{不交叉}的配對裡，找總差距
最小的那一種。

\textbf{看最左邊的人。} 區間 $[l, r]$ 裡，第 $l$ 位和第 $k$ 位一組。 弧不交
叉，所以：

\begin{itemize}
\item 夾在中間的 $[l + 1, k - 1]$ 只能彼此配對，這一段的人數必須是偶數，也就
是 $k - l$ 是奇數。
\item 右邊的 $[k + 1, r]$ 也只能彼此配對，各自獨立。
\end{itemize}

\begin{center}
\fbox{\parbox{0.86\linewidth}{\centering
$\mathrm{dp}(l, r)$：第 $l$ 到第 $r$ 位（偶數個人）全部配成不交叉的對，差距總和的最小值。}}
\end{center}

\[
\mathrm{dp}(l, r) = \min_{k = l + 1,\, l + 3,\, \ldots,\, r} \bigl(|h_l - h_k| + \mathrm{dp}(l + 1, k - 1) + \mathrm{dp}(k + 1, r)\bigr),
\]
空區間是 $0$，只有偶數長度的區間有意義。 這和括號字串的題目是同一個形狀：最
左邊的「(」配到哪一個「)」，就把字串分成括號裡面和括號右邊兩段。

同樣的模型也出現在別的地方：一排待焊接的接點要兩兩用導線連起來，導線不能交
叉；一條 RNA 的鹼基兩兩配對，配對不能交錯。 看到「兩兩配對、不交叉」，就可
以看最左邊的那一個和誰一組。

\textbf{程式}

\begin{example}[label=ex:ivdp-pairs]{cpp}{(pairs.cpp) 兩兩相鄰離隊搭車}
#include <climits>
#include <cstdlib>
#include <iostream>
#include <vector>
using namespace std;

int main() {
    int n;
    cin >> n;
    vector<long long> h(n + 2, 0);
    for (int i = 1; i <= n; i++) {
        cin >> h[i];
    }
    // dp[l][r]：第 l 到第 r 位（偶數個人）全部離隊的最小總和；l > r 時是 0
    vector<vector<long long>> dp(n + 2, vector<long long>(n + 2, 0));
    for (int len = 2; len <= n; len += 2) {   // 只有偶數長度有意義
        for (int l = 1; l + len - 1 <= n; l++) {
            int r = l + len - 1;
            long long best = LLONG_MAX;
            for (int k = l + 1; k <= r; k += 2) {   // l 的搭檔是 k
                long long cost = llabs(h[l] - h[k]) + dp[l + 1][k - 1] + dp[k + 1][r];
                best = min(best, cost);
            }
            dp[l][r] = best;
        }
    }
    cout << dp[1][n] << '\n';
    return 0;
}
\end{example}

\iosample{
6 \\
3 8 2 2 9 4 \\
}{
2 \\
}

\begin{itemize}
\item \texttt{len += 2} 只算偶數長度；\texttt{k += 2} 讓 $k - l$ 保持奇數。
奇數長度的格子從來沒有被填過，也從來不會被讀到。
\item $k = l + 1$ 時中間是空的 \texttt{dp[l + 1][l]}；$k = r$ 時右邊是空的
\texttt{dp[r + 1][r]}。 和上一題一樣，表格多開一圈。
\item 身高到 $10^9$，一對的差距最大約 $10^9$，$200$ 對加起來約
$2 \times 10^{11}$，要用 \texttt{long long}。 \texttt{llabs} 是
\texttt{long long} 版的絕對值，在 \texttt{<cstdlib>}。
\end{itemize}

\textbf{複雜度}

偶數長度的區間大約 $n^2/4$ 個，每個試大約 $len/2$ 個搭檔，最內層大約執行
$n^3/24$ 次，$n = 400$ 時約 $2.7 \times 10^6$ 次。 最大的測試資料在本書測試
的環境下不到 $0.01$ 秒。

\subsection{常見錯誤}

\begin{itemize}
\item \textbf{每次叫差距最小的}：輸入 \texttt{4}、\texttt{1 5 6 10} 印出
\texttt{10}，正確答案是 $8$。
\item \textbf{只配原本就相鄰的}：只考慮 $(1, 2), (3, 4), \ldots$ 這種配法，範例
的輸入得到 $|3 - 8| + |2 - 2| + |9 - 4| = 10$。 漏掉了一層包一層的配對。
\item \textbf{沒有限制搭檔的奇偶}：把 \texttt{k += 2} 寫成 \texttt{k++}，範例
的輸入印出 \texttt{1}。 $k - l$ 是偶數時，夾在中間的是奇數個人，配不完；程式
卻去讀奇數長度的格子，那些格子從沒填過，還是 $0$，等於假裝那一段的人不用
花任何代價就走光了。
\item \textbf{總和用 \texttt{int}}：輸入 $n = 400$，身高 $1$ 和 $10^9$ 交錯，
正確答案是 \texttt{199999999800}。 把 \texttt{long long} 都改成
\texttt{int}，在本書測試的環境下印出 \texttt{-1863463112}。
\end{itemize}

\section{回文：從兩端往內縮}
\label{sec:ivdp-pal}

\textbf{回文}在第~2~章 2.6 節出現過：正著讀和倒著讀都一樣的字串，例如
\texttt{abba}、\texttt{level}。 回文的定義本身就是區間的形狀：一段字串是回
文，剛好是在「兩端的字元相同，而且拿掉兩端之後中間那段也是回文」的時候。 這
一節的轉移都只看兩端，每個狀態 $O(1)$，所以 $n$ 可以到幾千。

\subsection{哪些區間是回文}

先建一張表，記下每一段是不是回文：

\begin{center}
\fbox{\parbox{0.86\linewidth}{\centering
$\mathrm{pal}(l, r)$：$s[l..r]$ 是不是回文。}}
\end{center}

長度 $1$ 一定是；長度 $2$ 看兩個字元是否相同；長度至少 $3$ 時，
$\mathrm{pal}(l, r)$ 是「$s[l] = s[r]$ 而且 $\mathrm{pal}(l + 1, r - 1)$」。
$[l + 1, r - 1]$ 比 $[l, r]$ 短，所以仍然依長度由短到長填：

\begin{code}{cpp}{回文表}
// pal[l][r]：s[l..r] 是不是回文（0 起算）
vector<vector<char>> pal(n, vector<char>(n, 0));
for (int len = 1; len <= n; len++) {
    for (int l = 0; l + len - 1 < n; l++) {
        int r = l + len - 1;
        pal[l][r] = (s[l] == s[r]) && (len <= 2 || pal[l + 1][r - 1]);
    }
}
\end{code}

表格用 \texttt{char} 存 $0$ 和 $1$，每格一個位元組，$n = 3000$ 時約 $9$ MB。
建好之後，「$s$ 的某一段是不是回文」只要查一格。 很多題目先建這張表，再接別
的 DP：例如要把字串切成好幾段、每段都是回文，就是第 1 本第~9~章「前 $i$ 個切成好幾
段」的狀態，轉移時用這張表檢查最後一段。

第~2~章 2.6 節把\textbf{最長回文子序列}當成「$s$ 和倒過來的 $s$ 的最長共同子
序列」來算。 用區間的狀態更直接：狀態是 $s[l..r]$ 裡最長的回文子序列，看兩端，
相同就兩個都留下、加上中間的答案；不同就至少要丟掉一端。 本章練習第~4~題會
手算一次。

\subsection{例題：拆旗換色排回文}
\label{sec:ex-ivdp-flags}

\textbf{題意}（\akiref{flags-take-down-or-recolor}{拆旗換色排回文}）

\begin{itemize}
\item 一排 $n$ 面彩旗，顏色寫成字串 $s$。 要讓整排從左看和從右看的顏色順序
相同（回文），不能加掛新的旗子。
\item 拆掉第 $i$ 面旗花 $w_i$ 秒，拆下後兩旁的旗子靠攏；把任一面旗換成任意顏
色，每次花 $C$ 秒。 只剩一面旗也算回文。
\item 求最少要花幾秒。 $n \le 2000$，$C$ 和 $w_i$ 都到 $10^9$。
\end{itemize}

例如 $C = 3$、旗子 \texttt{rgbyr}、拆旗秒數 $2, 4, 1, 5, 6$：兩端都是
\texttt{r}；第 $2$ 面 \texttt{g} 和第 $4$ 面 \texttt{y} 對不上，把第 $4$ 面換
成 \texttt{g}，花 $3$ 秒，得到 \texttt{rgbgr}。 只拆不換的話，至少要拆
\texttt{g} 和 \texttt{b}（$4 + 1 = 5$ 秒）或其他組合，都比 $3$ 秒慢。

\textbf{想法}

最直接的做法是枚舉留下哪些旗子，有 $2^n$ 種；對每一種，留下的旗子從兩端往內
兩兩配對，對不上的那一對換一次色。 $n = 2000$ 不可能，但 $n \le 15$ 時可以當
對拍用的暴力解。

\textbf{看最外面的兩面旗。} 只考慮第 $l$ 到第 $r$ 面這一段，它的處理方式只有
三種：

\begin{itemize}
\item 拆掉第 $l$ 面：花 $w_l$，剩下 $[l + 1, r]$ 要修成回文。
\item 拆掉第 $r$ 面：花 $w_r$，剩下 $[l, r - 1]$ 要修成回文。
\item 兩面都留下：它們是修好之後的頭和尾，顏色必須相同。 顏色本來就一樣不用
花錢，不一樣就把其中一面換色，花 $C$。 中間的 $[l + 1, r - 1]$ 再自己修成回
文。
\end{itemize}

這三種涵蓋了所有可能：左端不是被拆掉、就是留下；右端也一樣；兩端都留下時，
它們就是最後那串回文的第一面和最後一面。 每一種都留下一段連續的區間：

\begin{center}
\fbox{\parbox{0.86\linewidth}{\centering
$f(l, r)$：只看第 $l$ 到第 $r$ 面旗時，把它修成回文的最少秒數。}}
\end{center}

\[
f(l, r) = \min\bigl(w_l + f(l + 1, r),\ \ w_r + f(l, r - 1),\ \ f(l + 1, r - 1) + [s_l \ne s_r] \cdot C\bigr),
\]
其中 $[s_l \ne s_r]$ 在兩端顏色不同時是 $1$，相同時是 $0$。 長度 $0$ 或
$1$ 的區間已經是回文，$f = 0$。

\textbf{兩端顏色相同時，也要比三種。} 單看「兩端同色就直接配對」很合理，但拆
旗的秒數每面不同。 看 $C = 6$、旗子 \texttt{rrgr}、拆旗秒數 $1, 4, 5, 3$：
兩端都是 \texttt{r}，配成一對之後中間剩 \texttt{rg}，最便宜也要 $4$ 秒（拆掉
第 $2$ 面）。 改成拆掉第 $1$ 面（$1$ 秒），剩下的 \texttt{rgr} 已經是回文，
總共只要 $1$ 秒。 下表是這組輸入填完的 $f$（從 $0$ 起算）：

\begin{center}
\begin{tabular}{cc|cccc}
\toprule
 & & \multicolumn{4}{c}{$r$} \\
 & & 0 \texttt{r} & 1 \texttt{r} & 2 \texttt{g} & 3 \texttt{r} \\
\midrule
\multirow{4}{*}{$l$} & 0 \texttt{r} & 0 & 0 & 5 & \textbf{1} \\
 & 1 \texttt{r} & & 0 & 4 & 0 \\
 & 2 \texttt{g} & & & 0 & 3 \\
 & 3 \texttt{r} & & & & 0 \\
\bottomrule
\end{tabular}
\end{center}

$f(0, 3)$ 的三種選擇：拆左端 $1 + f(1, 3) = 1 + 0 = 1$；拆右端
$3 + f(0, 2) = 3 + 5 = 8$；兩端配對 $f(1, 2) + 0 = 4$。 最小的是 $1$。

同樣的模型也出現在別的地方：一排展示櫃或座位要調整成左右對稱，每個位置撤掉
的成本不同，也可以花固定的成本換掉其中一個。

\textbf{程式}

\begin{example}[label=ex:ivdp-flags]{cpp}{(flags.cpp) 拆旗換色排回文}
#include <algorithm>
#include <iostream>
#include <string>
#include <vector>
using namespace std;

int main() {
    ios::sync_with_stdio(false);
    cin.tie(nullptr);
    int n;
    long long c;
    string s;
    cin >> n >> c >> s;
    vector<long long> w(n);
    for (long long &x : w) {
        cin >> x;
    }
    // f[l][r]：只留下 s[l..r] 這一段時，把它修成回文的最少秒數
    // 長度 0 或 1 已經是回文，f 是 0
    vector<vector<long long>> f(n, vector<long long>(n, 0));
    for (int len = 2; len <= n; len++) {
        for (int l = 0; l + len - 1 < n; l++) {
            int r = l + len - 1;
            long long inner = (len == 2) ? 0 : f[l + 1][r - 1];
            long long pairUp = inner + (s[l] == s[r] ? 0 : c);   // 兩端配成一對
            long long dropL = w[l] + f[l + 1][r];                 // 拆掉左端
            long long dropR = w[r] + f[l][r - 1];                 // 拆掉右端
            f[l][r] = min({pairUp, dropL, dropR});
        }
    }
    cout << f[0][n - 1] << '\n';
    return 0;
}
\end{example}

\iosample{
4 6 \\
rrgr \\
1 4 5 3 \\
}{
1 \\
}

\begin{itemize}
\item 這題的字串從 $0$ 起算，表格開 $n \times n$。 長度 $2$ 時中間是空的，
\texttt{f[l + 1][r - 1]} 會變成 \texttt{f[l + 1][l]}，左下半部的格子；它開表
時是 $0$，讀到也沒錯，但程式直接用 \texttt{len == 2} 判斷，意思比較清楚。
\item \texttt{min(\{a, b, c\})} 用大括號一次比三個數，第 1 本第~10~章用過，需要
\texttt{<algorithm>}。
\item 拆旗秒數到 $10^9$、最多 $2000$ 面，答案可能到 $10^{12}$ 左右，用
\texttt{long long}。
\end{itemize}

\textbf{複雜度}

狀態約 $n^2/2 = 2 \times 10^6$ 個，每個只比三種，時間 $O(n^2)$。 表格是
$2000 \times 2000$ 個 \texttt{long long}，$32$ MB。 最大的測試資料在本書測試
的環境下約 $0.01$ 秒。 這題的每一格只用到下一列和同一列，改成 $l$ 由大到小填
之後，也可以像第 1 本第~10~章一樣只留兩列；但 $32$ MB 離常見的限制還很遠，整
張表比較好寫、好檢查。

\begin{apcsnote}
區間 DP 的狀態、轉移寫對了，最常錯的是計算順序和邊界（空區間、長度 $1$、長
度 $2$）。 寫完之後，用 $n = 1, 2, 3$ 的輸入各跑一次，再拿範例裡的小例子手
算幾格表格對照；這幾組小輸入幾乎能抓出所有順序和邊界的錯。
\end{apcsnote}

\subsection{常見錯誤}

\begin{itemize}
\item \textbf{兩端同色就直接配對}：寫成「$s_l = s_r$ 時 $f(l, r) = f(l + 1, r - 1)$，
不再比較拆旗」。 範例的輸入印出 \texttt{4}，正確答案是 $1$。 每面旗的拆除
成本一樣時這個直覺是對的，加權之後就不一定。
\item \textbf{從兩端往內一對一對比}：用兩個指標從兩端往中間走，對不上的那一對
付 $C$ 或拆掉便宜的那面。 拆掉一面之後，後面所有的配對都會錯開一格，指標
法沒辦法處理。 範例的輸入兩端已經對上、中間 \texttt{rg} 付 $4$，同樣印出
\texttt{4}。
\item \textbf{忘了換色}：只寫拆旗的兩種轉移，兩端不同色時不考慮配對。 題目一
開始的例子 \texttt{rgbyr} 會印出 \texttt{5}，正確答案是 $3$。
\item \textbf{用 \texttt{int}}：拆旗秒數和 $C$ 都到 $10^9$，拆三面加起來就超過
\texttt{int} 的上限（約 $2.1 \times 10^9$）。
\end{itemize}

\section{環狀：把陣列接一倍}
\label{sec:ivdp-ring}

章首的作業簿如果放在一張圓桌上，第 $n$ 疊和第 $1$ 疊也相鄰，可以合併。 例如
圓桌上依序是 $2, 8, 3, 8, 2$ 本，排成一排時最少要 $56$；在圓桌上，頭尾兩疊
$2$ 和 $2$ 可以先合在一起，最少只要 $50$。 區間 DP 的狀態是一段，圓上的「一
段」要怎麼表示？

\textbf{一定有一個地方沒被跨過。} $n$ 疊排成一圈，相鄰兩疊之間有一個縫，一
共 $n$ 個縫。 每次合併讓一個縫消失（兩疊變成一疊），合成一疊要合併 $n - 1$ 次，
所以到最後剛好有一個縫\textbf{從來沒有被合併跨過}。 從那個縫剪開，圓就變成
一排，那一排用前面的區間 DP 算。

\begin{center}
\begin{tikzpicture}[x=1cm, y=1cm]
\begin{scope}
\foreach \v [count=\i] in {2,8,3,8,2} {
  \node[cell] (q\i) at ({90 - 72*(\i-1)}:1.45) {\v};
}
\draw[very thick, memStroke] ({90 - 72*1.5}:0.7) -- ({90 - 72*1.5}:2.2);
\node[font=\small] at (0,-2.3) {從第 $2$、$3$ 疊之間剪開};
\end{scope}
\begin{scope}[xshift=3.6cm]
\foreach \v [count=\i] in {2,8,3,8,2,2,8,3,8,2} {
  \node[cell, minimum width=0.7cm, minimum height=0.7cm] (d\i) at (0.75*\i,0) {\v};
  \node[idxlbl] at (0.75*\i,0.6) {\i};
}
\foreach \i in {6,...,10} {\node[cell, minimum width=0.7cm, minimum height=0.7cm, fill=black!6] at (0.75*\i,0) {};}
\foreach \i/\v in {6/2,7/8,8/3,9/8,10/2} {\node[font=\ttfamily] at (0.75*\i,0) {\v};}
\draw[decorate, decoration={brace, amplitude=5pt, mirror}, thick, draw=memStroke]
  (0.75*3-0.35,-0.5) -- (0.75*7+0.35,-0.5) node[midway, below=6pt, font=\small] {第 $3$ 到第 $7$ 疊};
\node[font=\small] at (4.1,-1.9) {接一倍之後，任何剪法都是長度 $5$ 的一段};
\end{scope}
\end{tikzpicture}
\diagcaption{左：圓桌上的五疊，從第 $2$、$3$ 疊之間剪開，得到 $3, 8, 2, 2, 8$。
右：把陣列接一倍（灰色是複製的部分），這一排剛好是第 $3$ 到第 $7$ 疊。}
\label{fig:ivdp-ring}
\end{center}

不知道是哪一個縫，就每個縫都試。 從第 $l$ 疊前面剪開，得到的一排是第
$l, l + 1, \ldots, n, 1, \ldots, l - 1$ 疊。 五種剪法的結果：

\begin{center}
\begin{tabular}{clc}
\toprule
從哪一疊開始 & 剪開後的一排 & 排成一排的最少力氣 \\
\midrule
$1$ & $2, 8, 3, 8, 2$ & $56$ \\
$2$ & $8, 3, 8, 2, 2$ & $50$ \\
$3$ & $3, 8, 2, 2, 8$ & $50$ \\
$4$ & $8, 2, 2, 8, 3$ & $50$ \\
$5$ & $2, 2, 8, 3, 8$ & $50$ \\
\bottomrule
\end{tabular}
\end{center}

圓桌的答案是其中最小的 $50$。 每一種剪法各做一次 $O(n^3)$，總共 $O(n^4)$，
太慢。 注意到這些排法都是同一個圈的不同起點：把陣列\textbf{接一倍}，寫成
$a_1, \ldots, a_n, a_1, \ldots, a_n$，那麼「從第 $l$ 疊開始剪開」的那一排，
就是新陣列的第 $l$ 到第 $l + n - 1$ 個。 在長度 $2n$ 的陣列上做一次區間
DP，只算長度不超過 $n$ 的區間，所有剪法的答案都在表裡：
\[
\text{圓桌的答案} = \min_{1 \le l \le n} \mathrm{dp}(l,\ l + n - 1).
\]
不同的剪法共用了大部分的區間，例如 $[3, 5]$ 同時出現在第 $1$、$2$、$3$ 種剪
法裡，只算一次。

\begin{example}[label=ex:ivdp-ring]{cpp}{(ring.cpp) 圓桌上合併作業簿}
#include <climits>
#include <iostream>
#include <vector>
using namespace std;

int main() {
    int n;
    cin >> n;
    vector<long long> a(2 * n + 1, 0);
    for (int i = 1; i <= n; i++) {
        cin >> a[i];
        a[i + n] = a[i];                       // 接一倍：第 n + i 疊就是第 i 疊
    }
    vector<long long> pre(2 * n + 1, 0);
    for (int i = 1; i <= 2 * n; i++) {
        pre[i] = pre[i - 1] + a[i];
    }
    vector<vector<long long>> dp(2 * n + 1, vector<long long>(2 * n + 1, 0));
    for (int len = 2; len <= n; len++) {       // 長度只要算到 n
        for (int l = 1; l + len - 1 <= 2 * n; l++) {
            int r = l + len - 1;
            long long best = LLONG_MAX;
            for (int k = l; k < r; k++) {
                best = min(best, dp[l][k] + dp[k + 1][r]);
            }
            dp[l][r] = best + (pre[r] - pre[l - 1]);
        }
    }
    long long ans = LLONG_MAX;
    for (int l = 1; l <= n; l++) {             // 從第 l 疊開始、長度 n 的一段
        ans = min(ans, dp[l][l + n - 1]);
    }
    cout << ans << '\n';
    return 0;
}
\end{example}

\iosample{
5 \\
2 8 3 8 2 \\
}{
50 \\
}

\begin{itemize}
\item 和範例~\ref{ex:ivdp-merge} 比，只改了三處：陣列和表格開 $2n + 1$；長
度只算到 $n$，但左端走遍整個長度 $2n$ 的陣列；答案取 $n$ 個長度 $n$ 的區間裡
最小的。
\item 長度超過 $n$ 的區間會把同一疊算兩次，沒有意義，所以不算。
\item 狀態約 $3n^2/2$ 個，是一排時的三倍左右；長的區間也比較多，最內層一共
大約執行 $2n^3/3$ 次，是一排時的四倍。 複雜度仍然是 $O(n^3)$。
\end{itemize}

這個做法不限於合併。 只要是「排成一圈，最後的結構一定有一個地方沒被跨過」
的區間 DP，都可以接一倍，再取所有長度 $n$ 的區間的最佳值。 一圈的題目想不到
該從哪裡剪時，先問自己：最後的結果裡，有沒有一個一定沒被跨過的縫？ 也有不用
接一倍的一圈：圓桌旁的人兩兩配對、配對不能交叉，「交不交叉」和從哪裡剪開無
關，從任何一個縫剪開，答案都一樣，直接當成一排做就好。

\subsection{常見錯誤}

\begin{itemize}
\item \textbf{忘了是一圈}：直接用範例~\ref{ex:ivdp-merge}，範例的輸入印出
\texttt{56}，正確答案是 $50$。 只要最好的剪法不是從原本的頭尾之間剪開，就會
出錯。
\item \textbf{答案取 \texttt{dp[1][2n]}}：那是把接一倍之後的 $2n$ 疊全部合
成一疊，每一疊都算了兩次，不是題目要的。 長度 $n$ 的區間才對應一種剪法。
\item \textbf{陣列只開 $n$}：讀第 $n + 1$ 到 $2n$ 疊時超出陣列。 接一倍的時
候，前綴和、表格都要跟著開兩倍。
\end{itemize}

\section{什麼時候用區間 DP}
\label{sec:ivdp-when}

這一章的問題，轉移有三種形狀：

\begin{center}
\begin{tabular}{p{0.2\linewidth}p{0.4\linewidth}p{0.3\linewidth}}
\toprule
形狀 & 轉移在看什麼 & 本章的例子 \\
\midrule
切點 & 最後一次合併的分界，或第一個動作的位置；切開後左右兩段獨立 & 作業簿、拆攤位、矩陣鏈乘 \\
配對 & 最左邊的那一個和誰一組；把區間分成裡面和右邊 & 兩兩離隊 \\
兩端 & 最外面兩個怎麼處理：拆左、拆右、或兩端配對 & 回文表、拆旗換色 \\
\bottomrule
\end{tabular}
\end{center}

適合想到區間 DP 的訊號：

\begin{itemize}
\item 題目的操作只發生在\textbf{相鄰}的東西之間（合併相鄰兩堆、拿掉一個後兩邊
靠攏），或答案和回文、括號、不交叉的配對有關。
\item 拆一次之後，剩下的是一段或兩段\textbf{連續}的區間，而且彼此獨立。
\item $n$ 在一千以內（每個狀態枚舉切點）或幾千（每個狀態只看兩端）。
\end{itemize}

看起來像、其實不是的情況：

\begin{itemize}
\item \textbf{子問題永遠是前綴。} 例如只能從左邊一個一個拿，那是第 1 本第~9~章的
一維狀態，用區間反而多一個維度。
\item \textbf{任意兩堆都可以合併。} 把作業簿換成「桌上任兩疊都能搬在一起」，
就不再有「連續一段」的限制。 這時「每次合併最少的兩疊」反而是對的，用第 1 本第~6~章
的優先佇列每次取出最小的兩疊即可，$O(n \log n)$。 章首的貪心錯，是因為「相
鄰」這個限制；限制拿掉，貪心就對了。
\item \textbf{挑出來的東西不必連續。} 從一堆物品挑一些、總重不超過限制，那是
第~1~章的背包，狀態是「前 $i$ 件、用了多少重量」，不是區間。
\end{itemize}

\section{本章常見錯誤}
\label{sec:ivdp-pitfalls}

各節的錯誤，歸納起來是四類。 寫區間 DP 時，交出去之前逐項檢查。

\begin{center}
\begin{tabular}{p{0.2\linewidth}p{0.46\linewidth}p{0.22\linewidth}}
\toprule
錯誤 & 例子 & 見 \\
\midrule
計算順序 & $l$、$r$ 都由小到大，讀到還沒填的格子（印出 $24$、$7$） &
第~\ref{sec:ivdp-order}、\ref{sec:ex-ivdp-demolish}~節 \\
邊界 & 空區間 \texttt{dp[r + 1][r]} 越界；長度 $2$ 的中間；奇數長度的格子被當成
$0$ & 第~\ref{sec:ex-ivdp-demolish}、\ref{sec:ex-ivdp-pairs}、\ref{sec:ex-ivdp-flags}~節 \\
轉移漏了選擇 & 用貪心；只配原本相鄰的；兩端同色就直接配對；忘了換色 &
第~\ref{sec:ex-ivdp-demolish}、\ref{sec:ex-ivdp-pairs}、\ref{sec:ex-ivdp-flags}~節 \\
這一步的花費 & 忘了加總本數；區塊長度用整排的 $n$；矩陣的維度差一，或把隨切點變的花費放在 $\min$ 外面 & 第~\ref{sec:ivdp-basic}、\ref{sec:ex-ivdp-demolish}、\ref{sec:ivdp-matrix}~節 \\
\bottomrule
\end{tabular}
\end{center}

另外兩個每一題都可能碰到的錯誤：總和超過 \texttt{int}，以及環狀的題目忘了接
一倍。 前者在動手前估答案最大多少；後者讀題時先確認頭尾是否相鄰。

\section{本章練習}
\label{sec:ivdp-exercises}

\begin{drills}[label=drills:ivdp]
\item 一排四疊作業簿是 $3, 1, 4, 2$ 本，規則和第~\ref{sec:ivdp-basic}~節相同。
照長度由短到長，手算 $\mathrm{dp}(l, r)$ 的每一格，寫出最少的力氣，以及最後
一次合併的分界 $k$。

\item 下面四種迴圈順序，用來填範例~\ref{ex:ivdp-merge} 的表格，哪幾種是對的？
錯的那幾種，範例的輸入會印出什麼？
\begin{enumerate}
\item 外層 \texttt{len} 由 $2$ 到 $n$，內層 \texttt{l} 由小到大。
\item 外層 \texttt{l} 由 $n$ 到 $1$，內層 \texttt{r} 由 $l + 1$ 到 $n$。
\item 外層 \texttt{l} 由 $1$ 到 $n$，內層 \texttt{r} 由 $l + 1$ 到 $n$。
\item 外層 \texttt{r} 由 $1$ 到 $n$，內層 \texttt{l} 由 $r - 1$ 到 $1$。
\end{enumerate}

\item 估計下列次數。
\begin{enumerate}
\item 範例~\ref{ex:ivdp-merge} 在 $n = 10$ 時，\texttt{best = min(...)} 那一行
執行幾次？
\item 範例~\ref{ex:ivdp-demolish} 在 $n = 10$ 時，\texttt{int cost = ...} 那一行
執行幾次？ 為什麼比 (a) 多？
\item 範例~\ref{ex:ivdp-merge} 如果 $n = 3000$，最內層大約執行幾次？ 表格要幾
MB？ 依第~\ref{sec:ivdp-cost}~節的實測，大約要跑多久？
\end{enumerate}

\item 字串 \texttt{bdacab} 的最長回文子序列。 設 $g(l, r)$ 是 $s[l..r]$ 裡最長
的回文子序列長度。
\begin{enumerate}
\item 寫出 $g(l, r)$ 的轉移：$s_l = s_r$ 時和 $s_l \ne s_r$ 時各是什麼。 初始值
是什麼？
\item 手算出 $g(1, 6)$，並寫出一個達到這個長度的回文子序列。
\end{enumerate}

\item 有人想數字串裡有幾段連續的子字串是回文，寫了下面的程式片段。
\begin{code}{cpp}{練習 5}
int count = 0;
for (int l = 0; l < n; l++) {
    for (int r = l; r < n; r++) {
        pal[l][r] = (s[l] == s[r]) && (r - l < 2 || pal[l + 1][r - 1]);
        if (pal[l][r]) {
            count++;
        }
    }
}
\end{code}
\texttt{pal} 一開始全是 $0$。 輸入 \texttt{abba} 時它算出多少？ 正確答案是多
少？ 錯在哪裡，怎麼改？

\item 圓桌上的作業簿（第~\ref{sec:ivdp-ring}~節）。
\begin{enumerate}
\item 有人說：圓桌上的答案就是 $\mathrm{dp}(1, n)$，因為接一倍只是為了「讓第
$n$ 疊和第 $1$ 疊相鄰」。 這樣對嗎？
\item 為什麼「一定有一個縫從來沒被跨過」？ 如果題目改成每次可以合併\textbf{三
疊}相鄰的作業簿（$n$ 是奇數），這個說法還成立嗎？
\end{enumerate}

\item 四個矩陣的維度是 $d = 3, 2, 6, 8, 7$，也就是 $3 \times 2$、$2 \times 6$、
$6 \times 8$、$8 \times 7$。
\begin{enumerate}
\item 照長度由短到長，手算每一格 $\mathrm{dp}(l, r)$ 和 $\mathrm{cut}(l, r)$。
\item 照 $\mathrm{cut}$ 寫出加括號的方式，並驗算它的乘法次數。
\end{enumerate}

\item 矩陣鏈乘的暴力解與貪心。
\begin{enumerate}
\item 四個矩陣 $A_1 A_2 A_3 A_4$ 有幾種加括號的方式？ 全部列出來。
\item 有人提議：「每次把目前最大的中間維度消掉，也就是先做中間維度最大的那
一次乘法。」 三個矩陣 $d = 2, 5, 9, 4$ 時，這條規則做幾次乘法？ 最少是幾次？
規則錯在哪裡？
\end{enumerate}
\end{drills}

\section{實作練習}
\label{sec:ivdp-practice}

下面是 Aki-Judge 上和本章相關的題目，依轉移的形狀分組，本章的三題例題也列在
裡面。 難度分四階：暖身是直接套用本章的寫法；基本要自己定出狀態和轉移；實戰
的題敘比較長，要先把題目轉換成本章的形狀；挑戰會和其他章的技巧合在一起。

\begin{center}
\begin{tabular}{llc}
\toprule
形狀 & 題目 & 難度 \\
\midrule
切點 & \akiref{demolishing-a-row-of-stalls}{拆除一排相連的攤位}（例題） & 暖身 \\
\midrule
配對 & \akiref{leaving-the-line-in-pairs}{兩兩相鄰離隊搭車}（例題） & 基本 \\
\midrule
兩端與回文 & \akiref{longest-palindrome-from-each-position}{從每個位置開始的最長回文} & 暖身 \\
 & \akiref{flags-take-down-or-recolor}{拆旗換色排回文}（例題） & 基本 \\
 & \akiref{squared-scores-for-palindromic-pieces}{回文切段的平方得分} & 實戰 \\
\midrule
區間加一點資訊 & \akiref{how-late-can-you-arrive}{最晚幾點到夜市} & 挑戰 \\
\bottomrule
\end{tabular}
\end{center}

例題讀完之後，可以先關上書自己寫一次，再上 Aki-Judge 提交。

\section{讀完你應該會}
\label{sec:ivdp-checklist}

\begin{itemize}
\item 拿到一個題目，能判斷拆一次之後剩下的是前綴還是中間的一段，決定要用第 1 本第~9~
章的一維狀態還是區間狀態，並用一句中文寫出 $\mathrm{dp}(l, r)$ 的意思。
\item 能用「長度由短到長」的三層迴圈填區間 DP 的表，並說得出 $l$、$r$ 都由小
到大為什麼錯。
\item 能從「最後一次合併」「第一個動作」「最左邊的和誰一組」「最外面兩個怎麼處
理」中挑出合適的一種，寫出轉移，並處理空區間和奇偶；題目要求印出做法
時，能另外記下每一段的切點，再用遞迴還原（例如矩陣鏈乘的加括號方式）。
\item 能算出區間 DP 最內層的執行次數（例如 $(n^3 - n)/6$）和表格的記憶體，
判斷給定的 $n$ 來不來得及。
\item 遇到排成一圈的題目，能把陣列接一倍，取所有長度 $n$ 的區間的最佳值。
\end{itemize}

\textbf{延伸閱讀}（書末會整理完整的參考資料）

\begin{itemize}
\item Thomas H. Cormen 等，\textit{Introduction to Algorithms}，第 4 版，14.2
節：矩陣連乘的順序，是枚舉切點最經典的例子，有完整的正確性說明。
\item 劉汝佳，《算法競賽入門經典》第 2 版，9.4.1 節：矩陣連乘、多邊形切割、
切木棍、括號序列幾種區間 DP，可以看同一個想法在不同題目上的樣子。
\item 吳邦一，AP325，6.2.4 節：切棍子和矩陣乘法鏈，題目多，適合要準備 APCS 的
人多練。
\end{itemize}

\addcontentsline{toc}{section}{練習解答}
\section*{練習解答}

\begin{enumerate}
\item 長度 $2$：$\mathrm{dp}(1, 2) = 4$、$\mathrm{dp}(2, 3) = 5$、
$\mathrm{dp}(3, 4) = 6$。 長度 $3$：$\mathrm{dp}(1, 3) = 8 + \min(0 + 5,\ 4 + 0) = 12$；
$\mathrm{dp}(2, 4) = 7 + \min(0 + 6,\ 5 + 0) = 12$。 長度 $4$：
$\mathrm{dp}(1, 4) = 10 + \min(0 + 12,\ 4 + 6,\ 12 + 0) = 20$，在 $k = 2$ 取到。
也就是先合 $3 + 1$、再合 $4 + 2$，最後合成一疊：$4 + 6 + 10 = 20$。

\item (a)、(b)、(d) 是對的，(c) 錯。
\begin{itemize}
\item (b)：算第 $l$ 列時，比它大的列（下面的列）都填完了，所以
$\mathrm{dp}(k + 1, r)$ 已經有值；同一列左邊的 $\mathrm{dp}(l, k)$ 因為
$r$ 由小到大，也先填了。
\item (d)：算 $\mathrm{dp}(l, r)$ 時，$\mathrm{dp}(l, k)$ 的右端 $k < r$，在
前面幾輪外層迴圈就填好了；$\mathrm{dp}(k + 1, r)$ 和它同一個 $r$、左端比較
大，因為內層 $l$ 由大到小，也先填了。
\item (c) 就是第~\ref{sec:ivdp-order}~節的錯誤，印出 \texttt{24}。
\end{itemize}

\item
\begin{enumerate}
\item $(n^3 - n)/6 = (1000 - 10)/6 = 165$ 次。
\item $n(n + 1)(n + 2)/6 = 10 \cdot 11 \cdot 12 / 6 = 220$ 次。 拆攤位的長度從
$1$ 開始，而且切點 $k$ 可以是 $l$ 到 $r$ 的每一個（$len$ 個），作業簿的切點只有
$len - 1$ 個。
\item 大約 $3000^3 / 6 = 4.5 \times 10^9$ 次。 表格 $3001 \times 3001$ 個
\texttt{long long} 約 $72$ MB。 從 $n = 2000$ 約 $0.9$ 秒推算，時間和 $n^3$
成正比，$1.5^3 \approx 3.4$ 倍，大約 $3$ 秒，來不及。 在本書測試的環境下實測
約 $4$ 秒，比推算的還慢：表格大了，讀取記憶體的時間也跟著變長。
\end{enumerate}

\item
\begin{enumerate}
\item $s_l = s_r$ 時 $g(l, r) = g(l + 1, r - 1) + 2$；$s_l \ne s_r$ 時
$g(l, r) = \max\bigl(g(l + 1, r),\ g(l, r - 1)\bigr)$。 初始值 $g(i, i) = 1$，
空區間是 $0$（長度 $2$ 而兩字元相同時，中間是空的，得到 $2$）。
\item $g(1, 6) = 5$。 $s_1 = s_6 = \texttt{b}$，所以 $g(1, 6) = g(2, 5) + 2$；
$s_2 = \texttt{d}$、$s_5 = \texttt{a}$ 不同，$g(2, 5) = \max(g(3, 5), g(2, 4))$；
$g(3, 5)$ 是 \texttt{aca}，等於 $3$。 所以 $g(1, 6) = 3 + 2 = 5$，回文子序列是
\texttt{bacab}（刪掉 \texttt{d}）。
\end{enumerate}

\item 算出 $5$，正確答案是 $6$（\texttt{a}、\texttt{b}、\texttt{b}、\texttt{a}、
\texttt{bb}、\texttt{abba}）。 外層 $l = 0$ 時要用 \texttt{pal[1][2]} 判斷
\texttt{abba}，可是第 $1$ 列要到 $l = 1$ 才填，當時還是 $0$，所以
\texttt{abba} 被當成不是回文。 長度 $1$、$2$ 不需要中間那格，所以這種錯只在
長度至少 $3$ 的回文出現。 改成外層 \texttt{l} 由 $n - 1$ 到 $0$，或照本章用
長度由短到長。

\item
\begin{enumerate}
\item 不對。 $\mathrm{dp}(1, n)$ 是新陣列的第 $1$ 到 $n$ 個，也就是從第 $1$
疊前面剪開的那一種，等於沒把它當成一圈，$2, 8, 3, 8, 2$ 會得到 $56$。 接一倍的
用意是讓「從每一個縫剪開」都變成表裡的一段，所以要取 $n$ 個長度 $n$ 的區間
裡最小的。
\item 一圈有 $n$ 個縫，每次合併相鄰兩疊讓一個縫消失，總共合併 $n - 1$ 次，所
以剩下一個縫從來沒被跨過。 每次合併三疊時，一次讓兩個縫消失，從 $n$ 疊合成
一疊要合併 $(n - 1)/2$ 次，消失 $n - 1$ 個縫，仍然剩下一個，說法照樣成立：
從那個縫剪開，同樣可以接一倍處理。
\end{enumerate}

\item
\begin{enumerate}
\item 長度 $2$：$\mathrm{dp}(1, 2) = 3 \cdot 2 \cdot 6 = 36$、
$\mathrm{dp}(2, 3) = 2 \cdot 6 \cdot 8 = 96$、$\mathrm{dp}(3, 4) = 6 \cdot 8 \cdot 7 = 336$，
切點就是左端。 長度 $3$：$\mathrm{dp}(1, 3) = \min(0 + 96 + 48,\ 36 + 0 + 144) = 144$，
$\mathrm{cut} = 1$；$\mathrm{dp}(2, 4) = \min(0 + 336 + 84,\ 96 + 0 + 112) = 208$，
$\mathrm{cut} = 3$。 長度 $4$：$k = 1$ 是 $0 + 208 + 42 = 250$，$k = 2$ 是
$36 + 336 + 126 = 498$，$k = 3$ 是 $144 + 0 + 168 = 312$，所以
$\mathrm{dp}(1, 4) = 250$，$\mathrm{cut} = 1$。
\item $\mathrm{cut}(1, 4) = 1$，$\mathrm{cut}(2, 4) = 3$，$\mathrm{cut}(2, 3) = 2$，
寫成 $(A_1 ((A_2 A_3) A_4))$。 驗算：$A_2 A_3$ 花 $96$，得到 $2 \times 8$；再乘
$A_4$ 花 $2 \cdot 8 \cdot 7 = 112$，得到 $2 \times 7$；最後 $A_1$ 乘上它花
$3 \cdot 2 \cdot 7 = 42$。 總共 $96 + 112 + 42 = 250$。 這一組的最小維度 $2$ 在
$A_1$ 和 $A_2$ 之間，最好的做法把它留到最後一次乘法。
\end{enumerate}

\item
\begin{enumerate}
\item $5$ 種：
\begin{gather*}
((A_1 A_2) A_3) A_4,\qquad (A_1 (A_2 A_3)) A_4,\qquad (A_1 A_2)(A_3 A_4),\\
A_1 ((A_2 A_3) A_4),\qquad A_1 (A_2 (A_3 A_4)).
\end{gather*}
依最後一次乘法的切點分：
$k = 1$ 時右邊三個有 $2$ 種，$k = 2$ 時左右各 $1$ 種，$k = 3$ 時左邊三個有
$2$ 種，$2 + 1 + 2 = 5$。
\item 中間維度是 $5$ 和 $9$，規則先消掉 $9$，也就是先算 $A_2 A_3$：
$5 \cdot 9 \cdot 4 = 180$，再 $2 \cdot 5 \cdot 4 = 40$，共 $220$。 先算
$A_1 A_2$ 是 $2 \cdot 5 \cdot 9 = 90$，再 $2 \cdot 9 \cdot 4 = 72$，共 $162$，
這才是最少的。 消掉一個維度的那次乘法，花費還要乘上它左右兩個維度；規則只看
要消掉的那個數，沒看另外兩個有多大。 和第~\ref{sec:ivdp-matrix}~節的「先做最
便宜的」一樣，只看眼前一步的規則都可能錯，要靠區間 DP 比過每個切點。
\end{enumerate}
\end{enumerate}

\chapter{位元 DP 與旅行推銷員}
\label{ch:bitmask-dp}

\tikzset{
    bmnode/.style={draw=black!55, rounded corners=2pt, fill=white, font=\small,
        inner sep=3pt, minimum height=0.6cm},
    bmpath/.style={draw=memStroke, thick, fill=memFill},
    bmedge/.style={draw=black!25},
    bmpathedge/.style={draw=memStroke, very thick},
    bmtile/.style={draw=coreStroke, thick, fill=coreFill, rounded corners=2pt},
    bmpillar/.style={fill=black!70},
}

營隊有四個關卡，四位關主 A、B、C、D。 每位關主顧每一關的熟練程度不同，下
表是事前評估的分數，越高越好。 一人顧一關，每一關恰好一人，怎麼分配總分最
高？ 往下讀之前，先想一條規則，例如「從第 $0$ 關開始，每一關挑剩下的人裡最
熟練的」，照你的規則算一次總分。

\begin{center}
\begin{tabular}{c|cccc}
\toprule
關主 \textbackslash{} 關卡 & 0 & 1 & 2 & 3 \\
\midrule
A & 6 & 4 & 9 & 4 \\
B & 6 & 3 & 6 & 8 \\
C & 7 & 3 & 3 & 3 \\
D & 8 & 6 & 6 & 4 \\
\bottomrule
\end{tabular}
\end{center}

《C++ 演算法解題 1：APCS 實作題核心》（以下簡稱第 1 本）第~4~章 4.2 節用一個
整數的二進位表示「挑了哪些東西」，把 $2^n$ 種挑法一個一個試過。 這一章把同
一個整數拿來當動態規劃的\textbf{狀態}：前幾章 $\mathrm{dp}$ 的索引是「前 $i$
個」或「第 $l$ 到第 $r$ 個」，這一章的索引是一個集合，例如「已經排好的是哪幾
位關主」「已經去過哪幾個地方」。 集合用位元遮罩存，所以這種動態規劃叫\textbf{位元
DP}（bitmask DP），也常被叫做狀態壓縮 DP。

第~\ref{sec:bm-set}~節用關主的問題說明集合為什麼可以當狀態，把 $n!$ 種排法降
到 $2^n \cdot n$。 第~\ref{sec:bm-tsp}~節是最有名的位元 DP：從家出發，每個地
方去一次再回家，怎麼走最短，也就是\textbf{旅行推銷員問題}（traveling salesman
problem，常簡稱 TSP），這時還要記「現在停在哪裡」。 第~\ref{sec:bm-submask}
~節列舉子集合的子集合，第~\ref{sec:bm-tiling}~節一列一列鋪地磚，中間也會算
清楚這種表格要多少記憶體。

動態規劃的四件事（狀態、轉移、初始值、計算順序）和往前推、往回拉兩種寫法，
第 1 本第~9~章都講過；第 1 本第~4~章 4.2 節講過位元遮罩，第~\ref{sec:bm-set}
~節會用一個框回顧。 語法只用到 \texttt{vector} 和 GCC 內建的
\texttt{\_\_builtin\_popcount}、\texttt{\_\_builtin\_ctz}，用到時再說明。

\section{把集合當成狀態}
\label{sec:bm-set}

\subsection{先試貪心，再寫暴力解}

章首的規則「每一關挑剩下的人裡最熟練的」是這樣走的：第 $0$ 關四個人的分數
是 $6, 6, 7, 8$，挑 D；第 $1$ 關剩 A、B、C，分數 $4, 3, 3$，挑 A；第 $2$ 關
剩 B、C，挑 B；第 $3$ 關只剩 C。 總分 $8 + 4 + 6 + 3 = 21$。 另一種排法是
C 顧第 $0$ 關、D 顧第 $1$ 關、A 顧第 $2$ 關、B 顧第 $3$ 關，總分
$7 + 6 + 9 + 8 = 30$。

貪心在第 $0$ 關拿走 D，只比 C 多 $1$ 分；可是第 $1$ 關除了 D 以外，每個人都
只有 $3$、$4$ 分。 前面的選擇把後面需要的人用掉了，這是貪心常見的失敗方式。

最直接的做法是試遍所有排法：第 $0$ 關有 $n$ 個人可以選，第 $1$ 關剩 $n - 1$
個，一共 $n!$ 種。 第 1 本第~4~章 4.3 節用回溯列出排列，\texttt{used[x]} 記錄哪
些數已經用過；這裡照做就好。 $n = 10$ 時 $10! \approx 3.6 \times 10^6$ 種，
來得及；$n = 16$ 時 $16! \approx 2.1 \times 10^{13}$ 種，跑不完。

\subsection{看已經用掉哪些人}
\label{sec:bm-state}

照關卡的順序一關一關排。 假設第 $0$、$1$ 關已經排好，用掉的是 A 和 B，有兩
種排法：

\begin{itemize}
\item A 顧第 $0$ 關、B 顧第 $1$ 關：$6 + 3 = 9$ 分。
\item B 顧第 $0$ 關、A 顧第 $1$ 關：$6 + 4 = 10$ 分。
\end{itemize}

接下來第 $2$、$3$ 關只能由 C、D 來顧，不管前面是哪一種，後面的選擇和得分都
完全一樣。 所以 $9$ 分那種不可能出現在最佳解裡，只要記住 $10$ 就好。 後面的
事只看「哪些人已經用掉」，不看他們用掉的順序；回溯法的 \texttt{used} 陣列，
正好就是這個資訊。

用掉幾個人，就表示排好了幾關，所以「下一關是第幾關」也由這個集合決定。 狀態
定成：

\begin{center}
\fbox{\parbox{0.86\linewidth}{\centering
$\mathrm{dp}(S)$：集合 $S$ 裡的關主剛好顧滿第 $0$ 到第 $|S| - 1$ 關時，總分的最大值。}}
\end{center}

$|S|$ 是 $S$ 的元素個數。 照四件事寫下來：

\begin{enumerate}
\item \textbf{狀態}：$n$ 個人的每一個子集合都是一個狀態，共 $2^n$ 個。
\item \textbf{轉移}：從 $S$ 往前推。 下一關是第 $k = |S|$ 關，挑一個不在 $S$
裡的人 $v$ 來顧：
\[
\mathrm{dp}(S \cup \{v\}) \leftarrow \mathrm{dp}(S) + a_{v,k},
\]
箭頭表示「如果比較大就更新」。 也可以往回拉：$\mathrm{dp}(S)$ 看最後一關
（第 $|S| - 1$ 關）是 $S$ 裡的誰顧的，取
$\max_{v \in S} \bigl(\mathrm{dp}(S \setminus \{v\}) + a_{v,|S|-1}\bigr)$。
\item \textbf{初始值}：$\mathrm{dp}(\varnothing) = 0$，一關都還沒排。
\item \textbf{計算順序}：每個轉移都從小的集合走到多一個元素的集合，所以只要
元素少的先算。 下面會看到，用整數存集合時，直接照整數由小到大就對了。
\end{enumerate}

答案是 $\mathrm{dp}(\{A, B, C, D\})$。 圖~\ref{fig:bm-lattice} 把十六個狀態依元素個數排成五層，
每條線是「多排一個人」，框裡是 $\mathrm{dp}$ 的值。


走一個例子：$\mathrm{dp}(\{A, B\}) = 10$ 有兩條線進來，從 $\{A\}$ 來是 B 顧
第 $1$ 關，$6 + 3 = 9$；從 $\{B\}$ 來是 A 顧第 $1$ 關，$6 + 4 = 10$。 這就是
上面那兩種排法，表格只留下比較好的那個。 一條從頂端走到底部的路就是一種排法，
共 $4! = 24$ 條，可是框只有 $16$ 個：很多條路在中間走進同一個框，動態規劃在那
裡只留最好的一條。

這張圖就是第 1 本第~9~章 9.2 節「把動態規劃畫成圖」的圖：點是狀態，線是轉移。
位元 DP 的點特別多（$2^n$ 個），但每個點只有 $n$ 條線出去。

\begin{center}
\begin{tikzpicture}[x=1.05cm, y=0.78cm]
\node[bmnode, bmpath] (nE) at (0.00,0.00) {$\varnothing$\,:\,0};
\node[bmnode] (nA) at (-4.50,-1.45) {A\,:\,6};
\node[bmnode] (nB) at (-1.50,-1.45) {B\,:\,6};
\node[bmnode, bmpath] (nC) at (1.50,-1.45) {C\,:\,7};
\node[bmnode] (nD) at (4.50,-1.45) {D\,:\,8};
\node[bmnode] (nAB) at (-5.75,-2.90) {AB\,:\,10};
\node[bmnode] (nAC) at (-3.45,-2.90) {AC\,:\,11};
\node[bmnode] (nBC) at (-1.15,-2.90) {BC\,:\,10};
\node[bmnode] (nAD) at (1.15,-2.90) {AD\,:\,12};
\node[bmnode] (nBD) at (3.45,-2.90) {BD\,:\,12};
\node[bmnode, bmpath] (nCD) at (5.75,-2.90) {CD\,:\,13};
\node[bmnode] (nABC) at (-4.50,-4.35) {ABC\,:\,19};
\node[bmnode] (nABD) at (-1.50,-4.35) {ABD\,:\,21};
\node[bmnode, bmpath] (nACD) at (1.50,-4.35) {ACD\,:\,22};
\node[bmnode] (nBCD) at (4.50,-4.35) {BCD\,:\,19};
\node[bmnode, bmpath] (nABCD) at (0.00,-5.80) {ABCD\,:\,30};
\begin{scope}[on background layer]
\draw[bmedge] (nE) -- (nA);
\draw[bmedge] (nE) -- (nB);
\draw[bmedge, bmpathedge] (nE) -- (nC);
\draw[bmedge] (nE) -- (nD);
\draw[bmedge] (nA) -- (nAB);
\draw[bmedge] (nA) -- (nAC);
\draw[bmedge] (nA) -- (nAD);
\draw[bmedge] (nB) -- (nAB);
\draw[bmedge] (nB) -- (nBC);
\draw[bmedge] (nB) -- (nBD);
\draw[bmedge] (nAB) -- (nABC);
\draw[bmedge] (nAB) -- (nABD);
\draw[bmedge] (nC) -- (nAC);
\draw[bmedge] (nC) -- (nBC);
\draw[bmedge, bmpathedge] (nC) -- (nCD);
\draw[bmedge] (nAC) -- (nABC);
\draw[bmedge] (nAC) -- (nACD);
\draw[bmedge] (nBC) -- (nABC);
\draw[bmedge] (nBC) -- (nBCD);
\draw[bmedge] (nABC) -- (nABCD);
\draw[bmedge] (nD) -- (nAD);
\draw[bmedge] (nD) -- (nBD);
\draw[bmedge] (nD) -- (nCD);
\draw[bmedge] (nAD) -- (nABD);
\draw[bmedge] (nAD) -- (nACD);
\draw[bmedge] (nBD) -- (nABD);
\draw[bmedge] (nBD) -- (nBCD);
\draw[bmedge] (nABD) -- (nABCD);
\draw[bmedge, bmpathedge] (nCD) -- (nACD);
\draw[bmedge] (nCD) -- (nBCD);
\draw[bmedge, bmpathedge] (nACD) -- (nABCD);
\draw[bmedge] (nBCD) -- (nABCD);
\end{scope}
\end{tikzpicture}
\diagcaption{關主問題的十六個狀態。 第 $k$ 層是用掉 $k$ 個人、排好前 $k$ 關。
橘色是最佳解的一條路：C、D、A、B 依序顧第 $0$ 到第 $3$ 關。}
\label{fig:bm-lattice}
\end{center}

\subsection{用整數存集合}
\label{sec:bm-bits}

\begin{prevbookrecap}[第~4~章 4.2 節：位元遮罩]
$n$ 個東西編號 $0$ 到 $n - 1$，一個子集合寫成 $n$ 位的二進位數，第 $i$ 位是
$1$ 表示第 $i$ 個在集合裡。 例如 $n = 4$ 時 $\{A, C, D\}$ 是二進位
\texttt{1101}，也就是整數 $13$（最右邊是第 $0$ 位，代表 A）。
\begin{itemize}
\item \texttt{1 << n} 是 $2^n$，\texttt{for (int mask = 0; mask < (1 << n); mask++)}
走過每一個子集合一次。
\item \texttt{(mask >> i) \& 1} 是第 $i$ 位，$1$ 表示 $i$ 在集合裡。
\item \texttt{mask | (1 << i)} 把 $i$ 加進集合（第 1 本第~4~章 4.6 節用過）。
\item \texttt{1 << 31} 已經超出 \texttt{int}；有 \texttt{\&} 又有比較時要加括
號，\texttt{mask \& 1 << i == 0} 會先算 \texttt{==}。
\end{itemize}
\end{prevbookrecap}

位元 DP 還會用到幾個新的寫法。 下表的例子都用 $x = 44$，二進位
\texttt{101100}：

\begin{center}
\begin{tabular}{p{0.3\linewidth}p{0.36\linewidth}l}
\toprule
寫法 & 意思 & $x = 44$ 時 \\
\midrule
\texttt{x \& \textasciitilde(1 << i)} & 把 $i$ 拿掉 & $i = 3$：\texttt{100100} \\
\texttt{x \^{} (1 << i)} & 第 $i$ 位反過來（確定在集合裡時，就是拿掉） & $i = 2$：\texttt{101000} \\
\texttt{(a \& b) == a} & $a$ 是不是 $b$ 的子集合 & \\
\texttt{\_\_builtin\_popcount(x)} & 有幾個 $1$，也就是 $|S|$ & $3$ \\
\texttt{x \& -x} & 只留下最低的那個 $1$ & \texttt{000100}，$4$ \\
\texttt{\_\_builtin\_ctz(x)} & 最低的 $1$ 在第幾位（$x \ne 0$） & $2$ \\
\texttt{x \& (x - 1)} & 把最低的那個 $1$ 拿掉 & \texttt{101000}，$40$ \\
\bottomrule
\end{tabular}
\end{center}

\begin{itemize}
\item \texttt{\_\_builtin\_popcount}、\texttt{\_\_builtin\_ctz} 是 GCC 和 Clang
內建的函式，不用 \texttt{include}。 \texttt{ctz} 是 count trailing zeros，
數最右邊有幾個 $0$；$x = 0$ 時結果沒有定義，不能用。
\item \texttt{x - 1} 會把最低的 $1$ 變成 $0$、它右邊的 $0$ 全變成 $1$，其他位
不動：\texttt{101100} 減 $1$ 是 \texttt{101011}。 和 $x$ 做 AND，最低的 $1$
和它右邊都變成 $0$，就是 \texttt{101000}。
\item \texttt{-x} 在電腦裡是「每一位反過來再加 $1$」。 \texttt{101100} 反過來
是 \texttt{010011}，加 $1$ 是 \texttt{010100}；和 $x$ 只有最低的 $1$ 那一位
相同。
\end{itemize}

用整數存集合還有一個好處：把元素 $v$ 加進集合，整數只會變大，因為
\texttt{mask | (1 << v)} 多了一個 $1$。 所以「元素少的先算」不用另外安排，
\texttt{mask} 從 $0$ 數到 $2^n - 1$ 時，每個狀態被用來往前推之前，所有能走到
它的狀態都已經處理過了。 拿掉一個元素也一樣，\texttt{mask \^{} (1 << v)} 一定
比 \texttt{mask} 小，往回拉也照整數由小到大算。

\subsection{\texorpdfstring{從 $n!$ 到 $2^n \cdot n$}{從 n! 到 2ⁿ·n}}
\label{sec:bm-cost}

關主問題有 $2^n$ 個狀態，每個狀態試 $n$ 個人，總共 $O(2^n \cdot n)$。 和
$n!$ 比一比：

\begin{center}
\begin{tabular}{rrrrrr}
\toprule
$n$ & $n!$ & $2^n$ & $2^n \cdot n$ & $2^n \cdot n^2$ & $3^n$ \\
\midrule
$10$ & $3.6 \times 10^{6}$ & $1024$ & $1.0 \times 10^{4}$ & $1.0 \times 10^{5}$ & $5.9 \times 10^{4}$ \\
$16$ & $2.1 \times 10^{13}$ & $65536$ & $1.0 \times 10^{6}$ & $1.7 \times 10^{7}$ & $4.3 \times 10^{7}$ \\
$18$ & $6.4 \times 10^{15}$ & $2.6 \times 10^{5}$ & $4.7 \times 10^{6}$ & $8.5 \times 10^{7}$ & $3.9 \times 10^{8}$ \\
$20$ & $2.4 \times 10^{18}$ & $1.0 \times 10^{6}$ & $2.1 \times 10^{7}$ & $4.2 \times 10^{8}$ & $3.5 \times 10^{9}$ \\
$25$ & $1.6 \times 10^{25}$ & $3.4 \times 10^{7}$ & $8.4 \times 10^{8}$ & $2.1 \times 10^{10}$ & $8.5 \times 10^{11}$ \\
\bottomrule
\end{tabular}
\end{center}

$2^n$ 仍然是指數，只是長得比 $n!$ 慢很多。 所以位元 DP 只能處理很小的 $n$，
這也是辨認它的訊號：題目裡某個量只到 $15$、$16$、$20$ 左右，其他量（例如時
間、距離、詢問數）卻很大。 這時多半是要把那十幾個東西的「有或沒有」當成狀態。
後面兩欄是本章另外兩種轉移的量：旅行推銷員每個狀態要多看一層「停在哪裡」，是
$2^n \cdot n^2$；第~\ref{sec:bm-submask}~節列舉子集合的子集合，是 $3^n$。

\section{例題：成果發表會的出場順序}
\label{sec:ex-bm-show}

\textbf{題意}（\akiref{running-order-for-the-club-showcase}{成果發表會的出場順序}）

\begin{itemize}
\item 社團成果發表會有 $n$ 個節目，舞台時間切成 $n$ 個時段，一個時段排一個
節目。 節目 $i$ 排在第 $s$ 個時段，評分是 $a_{i,s}$；總評分是每個節目在它所
排時段的評分總和。
\item 另外有 $m$ 條前後限制「節目 $u$ 的時段要比節目 $v$ 早」，兩者不必相
鄰。
\item 求滿足所有限制時總評分的最大值；限制互相矛盾、排不出來時輸出 $-1$。
$n \le 16$，$a_{i,s} \le 10^6$。
\end{itemize}

例如 $n = 4$，評分是下表，限制是「$1$ 比 $2$ 早」「$4$ 比 $3$ 早」：

\begin{center}
\begin{tabular}{c|cccc}
\toprule
節目 \textbackslash{} 時段 & 1 & 2 & 3 & 4 \\
\midrule
1 & 3 & 3 & 8 & 7 \\
2 & 5 & 0 & 0 & 4 \\
3 & 7 & 4 & 3 & 9 \\
4 & 5 & 7 & 5 & 5 \\
\bottomrule
\end{tabular}
\end{center}

不管限制的話，排成 $2, 4, 1, 3$ 可以拿 $5 + 7 + 8 + 9 = 29$，但節目 $2$ 跑到
節目 $1$ 前面了。 滿足限制的排法裡最好的是 $4, 3, 1, 2$，總評分
$5 + 4 + 8 + 4 = 21$。

\textbf{想法}

拿掉限制，這就是上一節的關主問題：節目是關主，時段是關卡。 暴力解仍然是試
$n!$ 種順序、跳過違反限制的，只能拿來在小的 $n$ 對拍。

\textbf{限制也只看集合。} 照時段的順序一個一個排。 要把節目 $v$ 排進下一個時
段，所有必須比 $v$ 早的節目都要已經排好，也就是都在 $S$ 裡。 這個條件只看
「哪些節目已經排了」，和它們的順序無關，所以狀態仍然是 $S$。 把「必須比 $v$
早的節目」也存成一個遮罩 $\mathit{need}_v$，讀到限制 $u, v$ 時把 $u$ 那一位
加進 $\mathit{need}_v$。 檢查 $v$ 能不能排，只要看 $\mathit{need}_v$ 是不是
$S$ 的子集合：
\[
(\mathit{need}_v \mathbin{\&} S) = \mathit{need}_v .
\]

\textbf{有些集合排不出來。} 例如 $\{2\}$：節目 $2$ 要等節目 $1$，它不可能是
第一個時段的節目。 這種集合在前面幾個時段永遠不會出現，$\mathrm{dp}$ 要記成
「不可能」，不能是 $0$。 如果寫成 $0$，程式會假裝「節目 $2$ 先上台、拿 $0$
分」是一種合法的開頭，從它接下去，算出不存在的排法（見本節的常見錯誤）。
評分都不是負的，這裡用 $-1$ 表示不可能；限制矛盾時，全集也走不到，
$\mathrm{dp}$ 的值剛好就是題目要的 $-1$。

\begin{center}
\fbox{\parbox{0.86\linewidth}{\centering
$\mathrm{dp}(S)$：$S$ 裡的節目剛好排滿前 $|S|$ 個時段、而且沒有違反限制時，總評
分的最大值；排不出來時是 $-1$。}}
\end{center}

同樣的模型也出現在別的地方：工廠把 $n$ 道工序排進 $n$ 個時段，每個時段的電
費不同，有些工序要等前一道做完。 看到「把 $n$ 個東西排成一列，分數和位置有
關，$n$ 很小」，就可以用排好的集合當狀態。

\textbf{程式}

\begin{example}[label=ex:bm-show]{cpp}{(show.cpp) 成果發表會的出場順序}
#include <algorithm>
#include <iostream>
#include <vector>
using namespace std;

int main() {
    int n, m;
    cin >> n >> m;
    vector<vector<int>> a(n, vector<int>(n));    // a[i][s]：節目 i 排在第 s 個時段
    for (int i = 0; i < n; i++) {
        for (int s = 0; s < n; s++) {
            cin >> a[i][s];
        }
    }
    vector<int> need(n, 0);                       // need[v]：必須比 v 早的節目
    for (int k = 0; k < m; k++) {
        int u, v;
        cin >> u >> v;
        need[v - 1] |= 1 << (u - 1);
    }
    int full = (1 << n) - 1;
    // dp[mask]：mask 裡的節目剛好排滿前 popcount(mask) 個時段的最大總評分
    // -1 表示這個集合不可能是「前幾個時段」排的節目
    vector<int> dp(1 << n, -1);
    dp[0] = 0;
    for (int mask = 0; mask < full; mask++) {
        if (dp[mask] < 0) {
            continue;
        }
        int slot = __builtin_popcount(mask);      // 下一個要排的時段
        for (int v = 0; v < n; v++) {
            if ((mask >> v) & 1) {
                continue;                         // v 已經排過了
            }
            if ((need[v] & mask) != need[v]) {
                continue;                         // v 的前置還沒全部排好
            }
            int next = mask | (1 << v);
            dp[next] = max(dp[next], dp[mask] + a[v][slot]);
        }
    }
    cout << dp[full] << '\n';
    return 0;
}
\end{example}

\iosample{
4 2 \\
3 3 8 7 \\
5 0 0 4 \\
7 4 3 9 \\
5 7 5 5 \\
1 2 \\
4 3 \\
}{
21 \\
}

\begin{itemize}
\item 節目和時段在程式裡都從 $0$ 編號，讀限制時減 $1$。 \texttt{a[v][slot]}
的 \texttt{slot} 就是 \texttt{mask} 裡有幾個節目，不用另外記現在排到第幾個
時段。
\item 外層只跑到 \texttt{full - 1}：全集已經排滿，不用再往前推。
\item 把兩行 \texttt{need} 的檢查拿掉，就是第~\ref{sec:bm-set}~節關主問題的
程式。
\item 總評分最多 $16 \times 10^6$，\texttt{int} 裝得下。
\end{itemize}

\textbf{複雜度}

$2^n$ 個狀態，每個試 $n$ 個節目，$n = 16$ 時約 $10^6$ 次。 表格是 $2^{16}$
個 \texttt{int}，$256$ KB。 最大的測試資料在本書測試的環境下約 $0.01$ 秒。

\subsection{常見錯誤}

\begin{itemize}
\item \textbf{排不出來的集合記成 $0$}：把 \texttt{dp} 一開始全設成 \texttt{0}，
範例的輸入印出 \texttt{24}，比正確答案 $21$ 還大。 程式從 $\{2\}$ 這種不合
法的開頭接下去，得到「節目 $2$ 先上、拿 $0$ 分，再排 $4, 1, 3$」這種違反限制
的排法。 限制矛盾的輸入（例如 $1$ 比 $2$ 早、$2$ 比 $3$ 早、$3$ 比 $1$ 早）也
不會印出 $-1$。
\item \textbf{限制的方向存反}：寫成 \texttt{need[u - 1] |= 1 << (v - 1)}，等
於要求 $2$ 比 $1$ 早、$3$ 比 $4$ 早，範例的輸入印出 \texttt{22}。
\item \textbf{子集合的檢查少了括號}：寫成 \texttt{need[v] \& mask != need[v]}。
\texttt{!=} 比 \texttt{\&} 先算，這行其實是
\texttt{need[v] \& (mask != need[v])}，只剩 \texttt{need[v]} 的最低位有作用。
範例的輸入印出 \texttt{17}；\texttt{-Wall} 會警告「suggest parentheses」。
\item \textbf{試 $n!$ 種順序}：$n = 16$ 時跑不完。 先檢查限制有沒有矛盾、再沿
著合法的前綴搜尋，也沒有比較快：沒有限制時合法的前綴就是全部排列。
\end{itemize}

\section{旅行推銷員：集合加上最後一站}
\label{sec:bm-tsp}

週末要從家出發，去郵局、五金行、銀行、藥局各一次，最後回家。 下表是兩地之
間走路要幾分鐘（兩個方向一樣）。 怎麼排順序，總時間最短？

\begin{center}
\begin{tabular}{c|ccccc}
\toprule
 & 0 家 & 1 郵局 & 2 五金行 & 3 銀行 & 4 藥局 \\
\midrule
0 家 & -- & 5 & 9 & 3 & 4 \\
1 郵局 & 5 & -- & 9 & 4 & 8 \\
2 五金行 & 9 & 9 & -- & 8 & 9 \\
3 銀行 & 3 & 4 & 8 & -- & 6 \\
4 藥局 & 4 & 8 & 9 & 6 & -- \\
\bottomrule
\end{tabular}
\end{center}

「每次走到最近的、還沒去過的地方」是很自然的規則：家 $\to$ 銀行（$3$）$\to$
郵局（$4$）$\to$ 藥局（$8$）$\to$ 五金行（$9$）$\to$ 家（$9$），共 $33$ 分鐘。
最好的走法是家 $\to 3 \to 1 \to 2 \to 4 \to$ 家，
$3 + 4 + 9 + 9 + 4 = 29$ 分鐘。 五金行離每個地方都遠，最近鄰的規則把它拖到最
後，結果去程、回程都付了遠的代價。

\textbf{只記集合不夠。} 試著照上一節的方法，用「去過哪些地方」當狀態。 下一
段路要花多少時間，取決於\textbf{現在在哪裡}，集合 $\{1, 3\}$ 說不出人在郵局還
是在銀行。 所以狀態要再加一個數：

\begin{center}
\fbox{\parbox{0.86\linewidth}{\centering
$\mathrm{dp}(S, i)$：從家出發、去過的地方剛好是 $S$、現在停在 $i$（$i \in S$）時，
已經走的最短時間。}}
\end{center}

為什麼這樣就夠了？ 從 $(S, i)$ 之後還要去 $S$ 以外的每個地方再回家，這段路
怎麼走、要多久，只和「人在 $i$」「還剩哪些地方」有關，和前面怎麼走到 $i$ 無關。
所以走進同一個 $(S, i)$ 的各條路線，只要留下最短的。 例如下圖兩條路線都去過
$\{1, 2, 3\}$、停在五金行：

\begin{center}
\begin{tikzpicture}[x=1.45cm, y=1cm,
  pl/.style={draw=black!55, circle, minimum size=7mm, inner sep=0pt, font=\small}]
\node[pl] (a0) at (0,0.6) {家};
\node[pl] (a1) at (1,0.6) {1};
\node[pl] (a3) at (2,0.6) {3};
\node[pl] (a2) at (3,0.6) {2};
\draw[diagarrow] (a0) -- node[above, font=\scriptsize] {5} (a1);
\draw[diagarrow] (a1) -- node[above, font=\scriptsize] {4} (a3);
\draw[diagarrow] (a3) -- node[above, font=\scriptsize] {8} (a2);
\node[font=\small, right] at (3.4,0.6) {共 $17$};
\node[pl] (b0) at (0,-0.6) {家};
\node[pl] (b3) at (1,-0.6) {3};
\node[pl] (b1) at (2,-0.6) {1};
\node[pl, draw=memStroke, thick, fill=memFill] (b2) at (3,-0.6) {2};
\draw[diagarrow] (b0) -- node[above, font=\scriptsize] {3} (b3);
\draw[diagarrow] (b3) -- node[above, font=\scriptsize] {4} (b1);
\draw[diagarrow] (b1) -- node[above, font=\scriptsize] {9} (b2);
\node[font=\small, right] at (3.4,-0.6) {共 $16$};
\node[draw=memStroke, thick, fill=memFill, rounded corners=2pt, font=\small, align=left]
  at (6.3,0) {狀態 $(\{1, 2, 3\}, 2)$\\只留 $16$};
\end{tikzpicture}
\diagcaption{兩條路線走進同一個狀態。 之後去藥局再回家的路完全相同，所以 $17$ 那
條不可能比 $16$ 那條好。}
\label{fig:bm-merge}
\end{center}

轉移是「從 $i$ 走到一個還沒去過的 $j$」：
\[
\mathrm{dp}(S \cup \{j\}, j) \leftarrow \mathrm{dp}(S, i) + d_{i,j}, \qquad j \notin S.
\]
初始值是從家直接走到第一站：$\mathrm{dp}(\{i\}, i) = d_{0,i}$。 每個地方都
去過之後還要回家，答案是
\[
\min_{i} \bigl(\mathrm{dp}(\text{全集}, i) + d_{i,0}\bigr).
\]
家不放進集合裡：它是固定的起點，集合只記 $n$ 個要去的地方。 計算順序和上一節
一樣，$S \cup \{j\}$ 比 $S$ 大，\texttt{mask} 由小到大就對。 這個做法叫
\textbf{Held–Karp 演算法}，1962 年提出。 旅行推銷員問題到現在都沒有人找到多
項式時間的解法，$2^n \cdot n^2$ 已經是實用上很好的精確解。

\subsection{填表}
\label{sec:bm-tsp-table}

把四個地方的每個狀態都算出來（橫排是停在哪裡，「--」表示 $i$ 不在 $S$ 裡）：

\begin{center}
\begin{tabular}{cc|cccc}
\toprule
$S$ & \texttt{mask} & $i = 1$ & $i = 2$ & $i = 3$ & $i = 4$ \\
\midrule
$\{1\}$ & \texttt{0001} & 5 & -- & -- & -- \\
$\{2\}$ & \texttt{0010} & -- & 9 & -- & -- \\
$\{3\}$ & \texttt{0100} & -- & -- & 3 & -- \\
$\{4\}$ & \texttt{1000} & -- & -- & -- & 4 \\
\midrule
$\{1, 2\}$ & \texttt{0011} & 18 & 14 & -- & -- \\
$\{1, 3\}$ & \texttt{0101} & 7 & -- & 9 & -- \\
$\{2, 3\}$ & \texttt{0110} & -- & 11 & 17 & -- \\
$\{1, 4\}$ & \texttt{1001} & 12 & -- & -- & 13 \\
$\{2, 4\}$ & \texttt{1010} & -- & 13 & -- & 18 \\
$\{3, 4\}$ & \texttt{1100} & -- & -- & 10 & 9 \\
\midrule
$\{1, 2, 3\}$ & \texttt{0111} & 20 & \textbf{16} & 22 & -- \\
$\{1, 2, 4\}$ & \texttt{1011} & 22 & 21 & -- & 23 \\
$\{1, 3, 4\}$ & \texttt{1101} & 14 & -- & 16 & 15 \\
$\{2, 3, 4\}$ & \texttt{1110} & -- & 18 & 21 & 20 \\
\midrule
$\{1, 2, 3, 4\}$ & \texttt{1111} & 25 & 23 & 26 & 25 \\
加上回家 & & 30 & 32 & \textbf{29} & \textbf{29} \\
\bottomrule
\end{tabular}
\end{center}

地點 $i$ 對應第 $i - 1$ 位，\texttt{mask} 照習慣由右往左寫。 粗體的
$\mathrm{dp}(\{1, 2, 3\}, 2) = 16$ 就是圖~\ref{fig:bm-merge} 的那個狀態。 最
後一列的 $29$ 出現兩次：停在銀行的那條是家 $\to 4 \to 2 \to 1 \to 3 \to$ 家，
停在藥局的是它倒過來走。 距離是對稱的，同一個圈正反走都一樣長。

\subsection{把路線找回來}
\label{sec:bm-route}

表格只告訴我們最短是 $29$。 要知道怎麼走，可以像第~1~章 1.1 節「找出選了哪幾
件」那樣從答案往回推；這裡更簡單的做法是在填表時就記下來：每次更新 $\mathrm{dp}(S \cup \{j\}, j)$ 時，記下是從哪一站來的，存在
$\mathit{from}(S \cup \{j\}, j)$。 最後從答案的那個狀態開始，沿著
$\mathit{from}$ 往回走：知道現在在 $j$、前一站是 $i$，就把 $j$ 從集合拿掉，
回到狀態 $(S, i)$，一直走到集合只剩一個地方為止。 這樣得到的是倒過來的路線，
最後反轉一次。

\begin{example}[label=ex:bm-route]{cpp}{(route.cpp) 從家出發走遍每個地方再回家}
#include <algorithm>
#include <climits>
#include <iostream>
#include <vector>
using namespace std;

int main() {
    int n;                                        // 地點 0 是家，1 到 n 是要去的地方
    cin >> n;
    vector<vector<int>> d(n + 1, vector<int>(n + 1));
    for (auto &row : d) {
        for (int &x : row) {
            cin >> x;
        }
    }
    const int INF = INT_MAX / 2;
    int full = (1 << n) - 1;
    // 第 i 位代表地點 i + 1
    // dp[mask][last]：從家出發、去過的地方剛好是 mask、停在 last 的最短路程
    // from[mask][last]：達到這個最短路程時，前一站是哪裡（-1 表示前一站是家）
    vector<vector<int>> dp(1 << n, vector<int>(n, INF));
    vector<vector<int>> from(1 << n, vector<int>(n, -1));
    for (int i = 0; i < n; i++) {
        dp[1 << i][i] = d[0][i + 1];
    }
    for (int mask = 1; mask <= full; mask++) {
        for (int last = 0; last < n; last++) {
            if (dp[mask][last] == INF) {
                continue;
            }
            for (int j = 0; j < n; j++) {
                if ((mask >> j) & 1) {
                    continue;
                }
                int next = mask | (1 << j);
                int cost = dp[mask][last] + d[last + 1][j + 1];
                if (cost < dp[next][j]) {
                    dp[next][j] = cost;
                    from[next][j] = last;
                }
            }
        }
    }
    int best = INF, end = -1;
    for (int last = 0; last < n; last++) {        // 最後一站回家
        int total = dp[full][last] + d[last + 1][0];
        if (total < best) {
            best = total;
            end = last;
        }
    }
    // 從終點沿著 from 往回走，得到倒過來的路線
    vector<int> route;
    int mask = full, cur = end;
    while (cur != -1) {
        route.push_back(cur + 1);
        int prev = from[mask][cur];
        mask ^= 1 << cur;                         // 拿掉 cur，回到前一個狀態
        cur = prev;
    }
    reverse(route.begin(), route.end());
    cout << best << '\n';
    cout << 0;
    for (int p : route) {
        cout << ' ' << p;
    }
    cout << " 0\n";
    return 0;
}
\end{example}

\iosample{
4 \\
0 5 9 3 4 \\
5 0 9 4 8 \\
9 9 0 8 9 \\
3 4 8 0 6 \\
4 8 9 6 0 \\
}{
29 \\
0 4 2 1 3 0 \\
}

\begin{itemize}
\item 地點有 $n + 1$ 個，遮罩只有 $n$ 位：第 $i$ 位代表地點 $i + 1$。 讀表時
\texttt{d[last + 1][j + 1]} 的 \texttt{+ 1} 都是在換回地點編號。
\item \texttt{INF} 用 \texttt{INT\_MAX / 2}，兩個 \texttt{INF} 相加也不會溢
位。 程式在 \texttt{dp[mask][last] == INF} 時就跳過，不會拿 \texttt{INF} 去
加，但留一點空間比較保險。
\item \texttt{from} 是 $-1$ 的只有第一站（前一站是家），往回走到這裡就停。
\item 印出的是停在銀行那一條。 把找終點那個迴圈的 \texttt{total < best} 改成
\texttt{<=}，會印出停在藥局的 \texttt{0 3 1 2 4 0}。 題目如果要求特定的一條
（例如字典序最小），平手時怎麼選要另外想。
\end{itemize}

\subsection{路徑和迴路}
\label{sec:bm-path}

經過每個點恰好一次的路線，叫\textbf{漢米爾頓路徑}（Hamiltonian path）；最後
再回到起點的，叫\textbf{漢米爾頓迴路}（Hamiltonian cycle）。 旅行推銷員問題找
的是最短的漢米爾頓迴路。 題目常常只改起點、終點的規定，表格不用變，改的是初
始值和答案：

\begin{center}
\begin{tabular}{p{0.32\linewidth}p{0.3\linewidth}p{0.28\linewidth}}
\toprule
規定 & 初始值 & 答案 \\
\midrule
從家出發，最後回家 & $\mathrm{dp}(\{i\}, i) = d_{0,i}$ & $\min_i \mathrm{dp}(\text{全集}, i) + d_{i,0}$ \\
從家出發，停在哪裡都可以 & 同上 & $\min_i \mathrm{dp}(\text{全集}, i)$ \\
從 $s$ 出發，停在 $t$ & $\mathrm{dp}(\{s\}, s) = 0$ & $\mathrm{dp}(\text{全集}, t)$ \\
起點、終點都自己選 & 每個 $i$ 都 $\mathrm{dp}(\{i\}, i) = 0$ & $\min_i \mathrm{dp}(\text{全集}, i)$ \\
\bottomrule
\end{tabular}
\end{center}

第三列的起點 $s$ 本身在集合裡，因為它和其他點一樣是要經過一次的地方。 只問
「有沒有」漢米爾頓路徑，把 $\mathrm{dp}$ 換成 \texttt{bool} 就好；不能直接走
的兩點，轉移時跳過。

\textbf{迴路可以固定起點。} 一個圈從哪一點開始算都是同一個圈，所以要走遍
$n$ 個點再回到原點時，可以直接規定從第 $0$ 個點出發。 包含第 $0$ 個點的集合
只有 $2^{n - 1}$ 個，表格可以少一半。

\subsection{表格要多大}
\label{sec:bm-memory}

$\mathrm{dp}(S, i)$ 有 $2^n \cdot n$ 格。 $n = 20$ 時是 $2^{20} \times 20 \approx 2.1 \times 10^7$
格，每格一個 \texttt{int}（$4$ 位元組）：
\[
2^{20} \times 20 \times 4 = 83{,}886{,}080 \text{ 位元組} = 80 \text{ MB}.
\]
這裡的 MB 是 $2^{20}$ 位元組。 換成 \texttt{long long} 要 $160$ MB，再加一張
$\mathit{from}$ 又多 $80$ MB。 很多題目的記憶體限制是 $256$ MB，甚至
$64$ MB，這種表格寫下去之前一定要先算。 下表是 $n = 20$ 的旅行推銷員（固定
從第 $0$ 個點出發）在本書測試的環境下實際用掉的記憶體：

\begin{center}
\begin{tabular}{lrr}
\toprule
表格的寫法 & 記憶體 & 時間 \\
\midrule
全域陣列 \texttt{int dp[1 << 20][20]} & $84$ MB & 約 $0.2$ 秒 \\
\texttt{vector<vector<int>> dp(1 << n, vector<int>(n))} & $123$ MB & 約 $0.2$ 秒 \\
\texttt{vector<vector<int>> dp(n, vector<int>(1 << n))} & $87$ MB & 約 $0.2$ 秒 \\
全域陣列 \texttt{long long dp[1 << 20][20]} & $164$ MB & 約 $0.2$ 秒 \\
\bottomrule
\end{tabular}
\end{center}

\begin{itemize}
\item 第二種多了將近 $40$ MB：$2^{20}$ 個小 \texttt{vector} 各有自己的管理資料
和配置的額外空間。 維度對調（第三種）只開 $20$ 個大 \texttt{vector}，幾乎沒有
額外負擔。 $n \le 16$ 時都只有幾 MB，本章的例題用第二種，順序和
$\mathrm{dp}(S, i)$ 的寫法相同。
\item 這麼大的陣列不能宣告在 \texttt{main} 裡面。 函式裡的陣列放在堆疊上，堆
疊通常只有幾 MB（本書測試的環境是 $8$ MB），\texttt{int dp[1 << 20][20];} 寫
在 \texttt{main} 裡，程式一執行就當掉（segmentation fault）。 宣告成全域變數，
或用 \texttt{vector}。
\item 能固定起點時（上一小節），表格可以只留 $2^{n - 1}$ 個集合，再省一半。
\end{itemize}

\begin{apcsnote}
題目會寫出記憶體限制。 寫位元 DP 之前，先算「狀態數 $\times$ 每格位元組數」：
$2^{16} \times 16$ 個 \texttt{long long} 只要 $8$ MB，$2^{20} \times 20$ 個
\texttt{int} 就是 $80$ MB。 如果題目把測試資料分成幾個子題，$n \le 8$ 的子題
可以先用 \texttt{next\_permutation} 試遍所有順序拿下來，這個暴力解之後也能和
位元 DP 對拍。
\end{apcsnote}

\subsection{常見錯誤}

\begin{itemize}
\item \textbf{兩層迴圈的順序對調}：把外層改成 \texttt{last}、內層改成
\texttt{mask}。 範例的輸入印出 \texttt{32}，正確答案是 $29$。 外層
\texttt{last = 0} 跑完之後，停在地點 $1$ 的狀態就再也不會被拿來往前推；後來才
從別的地方走到地點 $1$ 的更新，全都沒有繼續傳下去。 往前推的寫法裡，每個狀態
被拿來推的時候必須已經算完整，所以集合要放在最外層。
\item \textbf{忘了回家}：答案取 $\min_i \mathrm{dp}(\text{全集}, i)$，範例的輸
入印出 \texttt{23}，那是「停在哪裡都可以」的答案。 讀題時先確定起點、終點的
規定，再對照上一小節的表。
\end{itemize}

\section{例題：越推越重的回收推車}
\label{sec:ex-bm-cart}

\textbf{題意}（\akiref{the-ever-heavier-recycling-cart}{越推越重的回收推車}）

\begin{itemize}
\item 環保股長推著自重 $C$ 的回收推車，從回收站（地點 $0$）出發，到 $n$ 個班
級（地點 $1$ 到 $n$）各收一次回收袋，最後回到回收站。 第 $i$ 班的袋子重
$w_i$，收了就一直留在車上。
\item 從地點 $u$ 直接走到 $v$ 的費力值是
$d_{u,v} \times (C + \text{離開 } u \text{ 時車上袋子的總重})$。 距離對稱，滿
足三角不等式。
\item 求總費力值的最小值。 $n \le 16$，$C$、$w_i$、$d_{u,v}$ 都不超過
$10^4$。
\end{itemize}

例如 $n = 3$、$C = 2$、袋子重 $3, 1, 2$，距離如下表。 照
$0 \to 3 \to 1 \to 2 \to 0$ 走：$5 \times 2 + 3 \times (2 + 2) + 2 \times (2 + 5) + 3 \times (2 + 6) = 10 + 12 + 14 + 24 = 60$，
這是最小的。

\begin{center}
\begin{tabular}{c|cccc}
\toprule
 & 0 & 1 & 2 & 3 \\
\midrule
0 & 0 & 4 & 3 & 5 \\
1 & 4 & 0 & 2 & 3 \\
2 & 3 & 2 & 0 & 4 \\
3 & 5 & 3 & 4 & 0 \\
\bottomrule
\end{tabular}
\end{center}

\textbf{想法}

最直接的是試 $n!$ 種順序，$16!$ 太多。 幾條直覺的規則都會錯：每次走最近的班
級（$0 \to 2 \to 1 \to 3 \to 0$）總共 $70$；袋子由輕到重收，把最重的留到最後，
在題目的範例 2 也不是最好的。 最好的走法常常是趁車子還輕時先走遠路，但「遠」
和「重」要一起權衡，沒有簡單的規則。

\textbf{這題看起來不像上一節。} 一段路的費力值會乘上當時的重量，而重量看起來
和前面怎麼走有關。 第一個反應可能是在狀態裡多記一維「目前的重量」，可是重量
可以到 $16 \times 10^4$，表格開不下。

\textbf{重量只看集合。} 停下來想一想：走到狀態 $(S, i)$ 時，車上有哪些袋子？
答案是 $S$ 裡每一班的袋子，一個不多、一個不少，和收的順序無關。 所以重量是
$S$ 的函數，不用另外記：
\[
W(S) = \sum_{j \in S} w_j .
\]
從 $(S, i)$ 走到 $j$ 的費力值是 $d_{i,j} \times (C + W(S))$，只用到 $S$、$i$、$j$，
仍然是上一節的 $\mathrm{dp}(S, i)$：

\begin{center}
\fbox{\parbox{0.86\linewidth}{\centering
$\mathrm{dp}(S, i)$：收過的班級剛好是 $S$、人停在班級 $i$ 時，到目前為止的最小
總費力值。}}
\end{center}

\[
\mathrm{dp}(S \cup \{j\}, j) \leftarrow \mathrm{dp}(S, i) + d_{i,j} \times \bigl(C + W(S)\bigr), \qquad j \notin S.
\]

乘的是 $W(S)$，不是 $W(S \cup \{j\})$：費力值看的是\textbf{離開} $i$ 時的重
量，$j$ 的袋子要到了 $j$ 才放上車。 第一段車上是空的，
$\mathrm{dp}(\{i\}, i) = d_{0,i} \times C$；最後回程時所有袋子都在車上，答案
是 $\min_i \bigl(\mathrm{dp}(\text{全集}, i) + d_{i,0} \times (C + W(\text{全集}))\bigr)$。

\textbf{一次算好每個集合的重量。} $W(S)$ 每次用迴圈加，會多一個 $n$ 倍。 可以
先把 $2^n$ 個集合的重量都算好：拿掉 $S$ 最低的那個元素，剩下的集合
\texttt{S \& (S - 1)} 比 $S$ 小、已經算好，再加上那個元素的重量就是 $W(S)$。
最低的元素是第幾個，用 \texttt{\_\_builtin\_ctz(S)}。 這種「從比較小的集合遞推」
的預先計算，位元 DP 很常用。

同樣的模型也出現在別的地方：貨車沿路裝貨、油耗隨載重增加；無人機逐點投遞包
裹、耗電和剩下的包裹重量有關。 轉移的成本只取決於「已經做過哪些事」時，就還
是 $\mathrm{dp}(S, i)$；成本和做過的順序有關時就不夠了，本章練習第~6~題會看
這種情況。

\textbf{程式}

\begin{example}[label=ex:bm-cart]{cpp}{(cart.cpp) 越推越重的回收推車}
#include <algorithm>
#include <climits>
#include <iostream>
#include <vector>
using namespace std;

int main() {
    int n;
    long long c;
    cin >> n >> c;
    vector<long long> w(n);
    for (long long &x : w) {
        cin >> x;
    }
    vector<vector<long long>> d(n + 1, vector<long long>(n + 1));
    for (auto &row : d) {
        for (long long &x : row) {
            cin >> x;
        }
    }
    int full = (1 << n) - 1;
    // load[mask]：收過 mask 這些班級之後，車上回收袋的總重
    vector<long long> load(1 << n, 0);
    for (int mask = 1; mask <= full; mask++) {
        int low = __builtin_ctz(mask);                    // 最低的 1 是第幾位
        load[mask] = load[mask & (mask - 1)] + w[low];
    }
    const long long INF = LLONG_MAX / 2;
    // dp[mask][last]：收過的班級剛好是 mask、人停在班級 last 的最小費力值
    // 班級在程式裡編號 0 到 n - 1，地點編號是班級編號加 1
    vector<vector<long long>> dp(1 << n, vector<long long>(n, INF));
    for (int i = 0; i < n; i++) {
        dp[1 << i][i] = d[0][i + 1] * c;                  // 第一段車上還沒有袋子
    }
    for (int mask = 1; mask <= full; mask++) {
        long long weight = c + load[mask];                // 離開 last 時車子的重量
        for (int last = 0; last < n; last++) {
            if (dp[mask][last] == INF) {
                continue;
            }
            for (int j = 0; j < n; j++) {
                if ((mask >> j) & 1) {
                    continue;
                }
                long long cost = dp[mask][last] + d[last + 1][j + 1] * weight;
                long long &target = dp[mask | (1 << j)][j];
                target = min(target, cost);
            }
        }
    }
    long long ans = INF;
    for (int last = 0; last < n; last++) {                // 最後回到回收站，車上全滿
        ans = min(ans, dp[full][last] + d[last + 1][0] * (c + load[full]));
    }
    cout << ans << '\n';
    return 0;
}
\end{example}

\iosample{
3 2 \\
3 1 2 \\
0 4 3 5 \\
4 0 2 3 \\
3 2 0 4 \\
5 3 4 0 \\
}{
60 \\
}

\begin{itemize}
\item 車子的重量 \texttt{weight} 只和 \texttt{mask} 有關，放在 \texttt{last} 那
層迴圈外面算一次。
\item \texttt{long long \&target} 是參考，讓 \texttt{dp[mask | (1 << j)][j]}
不用寫兩次。 如果你沒讀過銜接書，只要知道對 \texttt{target} 賦值就是對那一格
賦值即可。
\item 數值範圍：單段最大約 $10^4 \times (10^4 + 16 \times 10^4) = 1.7 \times 10^9$，
已經超過 \texttt{int}；整趟最大約 $1.5 \times 10^{10}$。 距離和重量都存成
\texttt{long long}，乘法也在 \texttt{long long} 裡算。
\end{itemize}

\textbf{複雜度}

$2^n \cdot n$ 個狀態，每個試 $n$ 個下一站，$O(2^n \cdot n^2)$，$n = 16$ 時約
$1.7 \times 10^7$ 次。 表格 $2^{16} \times 16$ 個 \texttt{long long} 是
$8$ MB。 最大的測試資料在本書測試的環境下約 $0.02$ 秒。

\subsection{常見錯誤}

\begin{itemize}
\item \textbf{回程只乘 $C$}：最後一段寫成 \texttt{d[last + 1][0] * c}，範例的
輸入印出 \texttt{40}。 回程是車子最重的一段，不能漏掉袋子的重量。
\item \textbf{把下一班的袋子也算進去}：乘上 \texttt{c + load[mask | (1 << j)]}，
範例的輸入印出 \texttt{71}。 費力值看的是離開時的重量。
\item \textbf{用 \texttt{int}}：$n = 16$、所有數值都是 $10^4$ 時，正確答案是
\texttt{15300000000}。 把程式裡的 \texttt{long long} 全改成 \texttt{int}，在
本書測試的環境下印出 \texttt{-1521225473}。
\item \textbf{以為正反走一樣}：距離對稱，但重量讓方向有差。 範例的
$0 \to 3 \to 1 \to 2 \to 0$ 是 $60$，倒過來的 $0 \to 2 \to 1 \to 3 \to 0$ 是
$70$。 不能只算一個方向。
\end{itemize}

\section{子集合的子集合}
\label{sec:bm-submask}

前面兩題的轉移都是「多加一個元素」。 有些題目一次要決定一整組：把人分成好幾
組、把工作分給幾台機器，轉移是「從剩下的人裡挑出一組」，要列舉某個集合的所有
子集合。

\subsection{列舉一個集合的子集合}

最直接的寫法是從 $0$ 數到 $2^n - 1$，再用 \texttt{(sub \& mask) == sub} 挑出
是 \texttt{mask} 子集合的。 \texttt{mask} 只有三個 $1$ 時，$2^n$ 個數裡只有
$8$ 個用得到，大部分時間都浪費了。 下面這個迴圈只走 \texttt{mask} 的子集合，
由大到小，一個不漏：

\begin{code}{cpp}{列舉 mask 的所有子集合（含空集合）}
for (int sub = mask; ; sub = (sub - 1) & mask) {
    // 處理 sub
    if (sub == 0) {
        break;
    }
}
\end{code}

以 \texttt{mask} $=$ \texttt{10110} 為例：

\begin{center}
\begin{tabular}{ccc}
\toprule
\texttt{sub} & \texttt{sub - 1} & \texttt{(sub - 1) \& mask} \\
\midrule
\texttt{10110} & \texttt{10101} & \texttt{10100} \\
\texttt{10100} & \texttt{10011} & \texttt{10010} \\
\texttt{10010} & \texttt{10001} & \texttt{10000} \\
\texttt{10000} & \texttt{01111} & \texttt{00110} \\
\texttt{00110} & \texttt{00101} & \texttt{00100} \\
\texttt{00100} & \texttt{00011} & \texttt{00010} \\
\texttt{00010} & \texttt{00001} & \texttt{00000} \\
\texttt{00000} & & 結束 \\
\bottomrule
\end{tabular}
\end{center}

只看 \texttt{mask} 的三個 $1$（第 $1$、$2$、$4$ 位），\texttt{sub} 在這三位
上依序是 \texttt{111}、\texttt{110}、\texttt{101}、\texttt{100}、\texttt{011}、
\texttt{010}、\texttt{001}、\texttt{000}，就是一個三位的二進位數由 $7$ 往下數到 $0$。 減 $1$ 是在做「往下數一個」；和 \texttt{mask} 做 AND，把借位時跑
到 \texttt{mask} 以外的 $1$ 清掉，等於跳過不在 \texttt{mask} 裡的位。 所以
$k$ 個 $1$ 的 \texttt{mask} 剛好走 $2^k$ 次。

空集合要不要處理，看題目。 迴圈寫在 \texttt{if (sub == 0) break;} 之前的部分
會處理到空集合。 兩種常見的寫錯法：寫成 \texttt{for (sub = mask; sub > 0; ...)}
會漏掉空集合；寫成 \texttt{sub >= 0} 則永遠停不下來，因為
\texttt{(0 - 1) \& mask} 又回到 \texttt{mask}。

\subsection{\texorpdfstring{為什麼是 $3^n$}{為什麼是 3ⁿ}}
\label{sec:bm-3n}

如果每個 \texttt{mask} 都把它的子集合列一遍，總共要走幾次？ 換個方式數：每一
對「\texttt{mask} 和它的一個子集合 \texttt{sub}」，每個元素都剛好是下面三種
之一：

\begin{itemize}
\item 不在 \texttt{mask} 裡；
\item 在 \texttt{mask} 裡，但不在 \texttt{sub} 裡；
\item 在 \texttt{sub} 裡（當然也在 \texttt{mask} 裡）。
\end{itemize}

$n$ 個元素各自三選一，共 $3^n$ 對，所以總次數是 $3^n$。 $n = 16$ 時是
$4.3 \times 10^7$，來得及；$n = 20$ 時是 $3.5 \times 10^9$，就太慢了。 如果每
個 \texttt{mask} 都從 $0$ 數到 $2^n - 1$ 再挑子集合，總共是
$2^n \times 2^n = 4^n$，$n = 16$ 時 $4.3 \times 10^9$，差了一百倍。

\subsection{例題：湊齊器材的專題小組}
\label{sec:ex-bm-kit}

\textbf{題意}（\akiref{teams-with-every-piece-of-gear}{湊齊器材的專題小組}）

\begin{itemize}
\item 攝影社 $n$ 位同學要組外拍小組，外拍需要 $k$ 種器材。 每位同學帶哪些器
材，用長度 $k$ 的 01 字串表示。
\item 一個小組至少一人，組內所有人的器材合起來每一種都至少有一件，才算合格。
每位同學最多加入一組，也可以不加入。
\item 最多能組成幾個合格的小組？ $n \le 16$，$k \le 10$。
\end{itemize}

例如 $n = 4$、$k = 3$，四位同學的器材是 \texttt{110}、\texttt{001}、
\texttt{101}、\texttt{010}。 第 $0$、$1$ 位合起來是 \texttt{111}，第 $2$、$3$
位合起來也是 \texttt{111}，可以組成兩組。

\textbf{想法}

\textbf{先試貪心。} 「每次挑人數最少的合格小組」聽起來省人，但會出錯。 看
$k = 5$、六位同學 \texttt{11001}、\texttt{00111}、\texttt{10000}、
\texttt{01000}、\texttt{00100}、\texttt{00010}：兩個人就湊齊的只有第 $0$、
$1$ 位，貪心先讓他們一組；可是第 $5$ 種器材只有這兩個人有，剩下的四個人湊不
出第二組，只有 $1$ 組。 讓第 $0$ 位和第 $4$、$5$ 位一組，第 $1$ 位和第
$2$、$3$ 位一組，就有 $2$ 組。 兩個「各自都很有用」的人被綁在同一組，浪費了。

暴力解是替每個人決定去哪一組或不參加，$n = 16$ 時分法超過 $8 \times 10^{10}$
種，只能在很小的 $n$ 對拍。

\textbf{狀態是「還能用的人」。} 一群人能組幾組，只和是哪些人有關。 令

\begin{center}
\fbox{\parbox{0.86\linewidth}{\centering
$\mathrm{dp}(S)$：只用集合 $S$ 裡的同學，最多能組成幾個合格的小組。}}
\end{center}

\textbf{固定編號最小的那個人。} 如果寫成「從 $S$ 裡挑任意一組合格的 $G$，
$\mathrm{dp}(S) = \max \mathrm{dp}(S \setminus G) + 1$」，同一種分法會因為「先
挑哪一組」被算很多次，也沒有處理「有人不參加」。 好的寫法是只看 $S$ 裡編號最
小的那位同學 $p$，他只有兩種情況：

\begin{itemize}
\item 他不參加：$\mathrm{dp}(S \setminus \{p\})$。
\item 他在某個合格小組 $G$ 裡，$p \in G \subseteq S$：
$\mathrm{dp}(S \setminus G) + 1$。
\end{itemize}

任何一種分法裡，$p$ 不是沒參加，就是剛好在一組，所以這兩種情況不重複也不遺
漏。 $G$ 是 $p$ 再加上 $S \setminus \{p\}$ 的一個子集合（可以是空的，表示他一
個人就湊齊了），用上一小節的迴圈列舉。 $S \setminus G$ 比 $S$ 小，所以
\texttt{mask} 由小到大算。

\textbf{一次算好每個集合合不合格。} 和回收推車的重量一樣，集合 $S$ 的器材聯
集可以從 $S \setminus \{\text{最低的人}\}$ 遞推：$\mathit{cover}(S) = \mathit{cover}(\texttt{S \& (S - 1)})$
OR 最低那個人的器材，$O(2^n)$ 就把每個集合的聯集算好。

同樣的模型也出現在別的地方：把一群工程師分成值班小組，每組都要涵蓋所有系統
的專長；把一批零件分裝成好幾套，每套都要配齊。 看到「把 $n$ 個東西分成好幾
組，每組要滿足某個條件，$n$ 很小」，就列舉子集合的子集合。

\textbf{程式}

\begin{example}[label=ex:bm-kit]{cpp}{(kit.cpp) 湊齊器材的專題小組}
#include <algorithm>
#include <iostream>
#include <string>
#include <vector>
using namespace std;

int main() {
    int n, k;
    cin >> n >> k;
    vector<int> gear(n, 0);                     // gear[i]：第 i 位同學的器材，k 位元
    for (int i = 0; i < n; i++) {
        string s;
        cin >> s;
        for (int j = 0; j < k; j++) {
            if (s[j] == '1') {
                gear[i] |= 1 << j;
            }
        }
    }
    int everyone = (1 << n) - 1, allGear = (1 << k) - 1;
    // cover[S]：同學集合 S 合起來有哪些器材
    vector<int> cover(1 << n, 0);
    for (int S = 1; S <= everyone; S++) {
        cover[S] = cover[S & (S - 1)] | gear[__builtin_ctz(S)];
    }
    // dp[S]：只用 S 裡的同學，最多能組成幾個合格的小組
    vector<int> dp(1 << n, 0);
    for (int S = 1; S <= everyone; S++) {
        int p = S & -S;                         // 編號最小的同學（只有他那一位是 1）
        int others = S ^ p;
        int best = dp[others];                  // 他不參加
        // 他所在的小組 = p 加上 others 的子集合 sub
        for (int sub = others; ; sub = (sub - 1) & others) {
            if (cover[sub | p] == allGear) {
                best = max(best, dp[others ^ sub] + 1);
            }
            if (sub == 0) {
                break;
            }
        }
        dp[S] = best;
    }
    cout << dp[everyone] << '\n';
    return 0;
}
\end{example}

\iosample{
6 5 \\
11001 \\
00111 \\
10000 \\
01000 \\
00100 \\
00010 \\
}{
2 \\
}

\begin{itemize}
\item \texttt{p} 是一個只有一位是 $1$ 的遮罩（代表那位同學），不是編號；所以
小組寫成 \texttt{sub | p}，剩下的人是 \texttt{others \^{} sub}。
\item \texttt{sub} 從 \texttt{others} 本身開始，最後一個是 $0$，表示他一個人
自成一組。 這一輪要先處理再跳出。
\end{itemize}

\textbf{複雜度}

固定最低位之後，每個 $S$ 只列舉 $S \setminus \{p\}$ 的子集合，次數是
$2^{|S| - 1}$。 全部加起來是 $(3^n - 1)/2$，$n = 16$ 時約
$2.2 \times 10^7$，比不固定時少一半。 最大的測試資料在本書測試的環境下約
$0.02$ 秒。

\subsection{常見錯誤}

\begin{itemize}
\item \textbf{漏掉空的 \texttt{sub}}：迴圈寫成
\texttt{for (int sub = others; sub > 0; ...)}。 輸入 \texttt{2 2}、
\texttt{11}、\texttt{11} 印出 \texttt{1}，正確答案是 $2$：每個人自己就湊齊
了，但「一個人一組」從來沒有被試過。
\item \textbf{忘了可以不參加}：把 \texttt{best} 的起點寫成一個很小的負數，
只考慮 $p$ 加入某一組。 只要有人湊得出一組，多的人塞進那組也還合格，答案不會
錯；但一組都湊不出來時，例如輸入 \texttt{2 2}、\texttt{10}、\texttt{10}，就
印出那個負數，正確答案是 $0$。
\item \textbf{一次加一個人}：仿照前兩節，狀態記「已經組好幾組、目前這一組有哪
些器材」，每次加一個人。 兩個未完成的小組，例如一組有相機和腳架、另一組只有
麥克風，沒辦法比較哪個比較好，只留一個就可能丟掉後來需要的那個。 這種做法在
隨機的輸入上常常碰巧答對，所以更要小心。
\end{itemize}

\section{一列一列鋪地磚}
\label{sec:bm-tiling}

前面的遮罩記的都是「$n$ 個東西裡用過哪些」。 這一節的遮罩記的是一列格子的
樣子。

\textbf{題目。} 一塊 $h$ 列、$w$ 行的地板，有些格子是柱子，其餘的要用
$1 \times 2$ 的地磚（可以直放或橫放）剛好鋪滿，地磚不能蓋到柱子。 共有幾種鋪
法？ 答案很大時取除以 $10^9 + 7$ 的餘數。 $w \le 10$，$h \le 1000$。

例如 $h = 4$、$w = 3$，左上角和右下角是柱子，共有 $5$ 種鋪法。

\textbf{一列的狀態。} 從上往下一列一列鋪。 第 $r$ 列鋪好之後，下面的列只受到
一件事影響：第 $r$ 列有哪些直磚伸進了第 $r + 1$ 列。 那些格子已經被蓋住，第
$r + 1$ 列只能鋪剩下的格子；第 $r$ 列以上其他的事情，下面完全看不到。 所以狀
態是「第 $r$ 列的哪幾格已經被上一列的直磚蓋住」，用一個 $w$ 位的遮罩存。

\begin{center}
\begin{tikzpicture}[x=0.6cm, y=0.6cm]
\begin{scope}
\foreach \r in {0,...,3} { \foreach \c in {0,...,2} {
  \draw[black!35] (\c,-\r) rectangle (\c+1,-\r-1); } }
\fill[bmpillar] (0,0) rectangle (1,-1);
\fill[bmpillar] (2,-3) rectangle (3,-4);
\draw[bmtile] (1.1,-0.1) rectangle (2.9,-0.9);
\draw[bmtile] (1.1,-1.1) rectangle (2.9,-1.9);
\draw[bmtile, draw=memStroke, fill=memFill] (0.1,-1.1) rectangle (0.9,-2.9);
\fill[black!8] (1,-2) rectangle (3,-3);
\foreach \r in {0,...,3} {\node[idxlbl, left] at (-0.1,-\r-0.5) {\r};}
\node[font=\small] at (1.5,-4.7) {(a)};
\end{scope}
\begin{scope}[xshift=3.6cm]
\foreach \r in {0,...,3} { \foreach \c in {0,...,2} {
  \draw[black!35] (\c,-\r) rectangle (\c+1,-\r-1); } }
\fill[bmpillar] (0,0) rectangle (1,-1);
\fill[bmpillar] (2,-3) rectangle (3,-4);
\draw[bmtile] (1.1,-0.1) rectangle (1.9,-1.9);
\draw[bmtile] (2.1,-0.1) rectangle (2.9,-1.9);
\draw[bmtile, draw=memStroke, fill=memFill] (0.1,-1.1) rectangle (0.9,-2.9);
\fill[black!8] (1,-2) rectangle (3,-3);
\node[font=\small] at (1.5,-4.7) {(b)};
\end{scope}
\node[font=\small, right, align=left] at (10.2,-2)
  {(a)、(b) 鋪好第 $0$、$1$ 列的方式不同，\\
   但都只有一塊直磚伸進第 $2$ 列的第 $0$ 格（橘色）。\\
   第 $2$ 列以下怎麼鋪，兩種情況完全一樣，\\
   所以合在同一個狀態「第 $2$ 列被蓋住的是 $\{0\}$」，\\
   方法數是 $2$。 灰色是第 $2$ 列還要鋪的格子。};
\end{tikzpicture}
\diagcaption{$4 \times 3$ 的地板，黑色是柱子。 兩種不同的上半部，留給下半部的
是同一個樣子。}
\label{fig:bm-profile}
\end{center}

\begin{center}
\fbox{\parbox{0.86\linewidth}{\centering
$f_r(M)$：第 $0$ 到第 $r - 1$ 列都鋪好，而且伸進第 $r$ 列的直磚剛好蓋住 $M$
這幾格的方法數。}}
\end{center}

一開始沒有東西伸進第 $0$ 列，$f_0(\varnothing) = 1$。 鋪完最後一列之後不能有
磚伸出地板，答案是 $f_h(\varnothing)$。

\textbf{一列怎麼轉移。} 已知第 $r$ 列被蓋住的格子 $M$（再加上這一列的柱子），
從左到右一格一格決定還空著的格子：

\begin{itemize}
\item 已經被蓋住或是柱子：跳過。
\item 放一塊直磚往下：這一格蓋住了，下一列的同一格也被蓋住，記進下一列的遮罩。
\item 放一塊橫磚：右邊那格也要空著，兩格一起蓋住。
\end{itemize}

走到這一列的最右邊，得到下一列的遮罩 $M'$，就把 $f_r(M)$ 加到
$f_{r+1}(M')$。 如果上一列伸下來的直磚剛好戳到這一列的柱子（$M$ 和柱子有交
集），這個狀態不合法，直接跳過。 「一格一格決定」用遞迴寫最直接：

\begin{example}[label=ex:bm-tiling]{cpp}{(tiling.cpp) 有柱子的地板鋪地磚}
#include <iostream>
#include <string>
#include <vector>
using namespace std;

const long long MOD = 1000000007;
int w;
vector<long long> nxt;   // 下一列各種「被伸進來的格子」的鋪法數

// 這一列處理到第 c 格；filled：這一列已經被蓋住（或是柱子）的格子
// out：往下一列伸出去的直磚蓋住了下一列的哪些格子；ways：到這裡的鋪法數
void go(int c, int filled, int out, long long ways) {
    if (c == w) {
        nxt[out] = (nxt[out] + ways) % MOD;
        return;
    }
    if ((filled >> c) & 1) {                      // 這格已經有東西，換下一格
        go(c + 1, filled, out, ways);
        return;
    }
    go(c + 1, filled, out | (1 << c), ways);      // 放一塊直磚，蓋住下一列的第 c 格
    if (c + 1 < w && !((filled >> (c + 1)) & 1)) {
        go(c + 2, filled, out, ways);             // 放一塊橫磚，蓋住第 c、c + 1 格
    }
}

int main() {
    int h;
    cin >> h >> w;
    vector<int> pillar(h, 0);                     // pillar[r]：第 r 列哪幾格是柱子
    for (int r = 0; r < h; r++) {
        string s;
        cin >> s;
        for (int c = 0; c < w; c++) {
            if (s[c] == '#') {
                pillar[r] |= 1 << c;
            }
        }
    }
    // cur[mask]：前面幾列都鋪好，而且上一列的直磚剛好伸進這一列 mask 這幾格的鋪法數
    vector<long long> cur(1 << w, 0);
    cur[0] = 1;                                   // 第 0 列上面沒有東西伸下來
    for (int r = 0; r < h; r++) {
        nxt.assign(1 << w, 0);
        for (int mask = 0; mask < (1 << w); mask++) {
            if (cur[mask] == 0 || (mask & pillar[r]) != 0) {
                continue;                         // 沒有這種情況，或直磚戳到柱子
            }
            go(0, mask | pillar[r], 0, cur[mask]);
        }
        cur = nxt;
    }
    cout << cur[0] << '\n';                       // 最後一列不能有磚伸出地板
    return 0;
}
\end{example}

\iosample{
4 3 \\
\#.. \\
... \\
... \\
..\# \\
}{
5 \\
}

範例每一列處理完之後的狀態（集合寫的是被蓋住的行號）：

\begin{center}
\begin{tabular}{cl}
\toprule
鋪完第幾列 & 伸進下一列的格子：鋪法數 \\
\midrule
開始 & $\varnothing$：$1$ \\
第 $0$ 列 & $\varnothing$：$1$，\quad $\{1, 2\}$：$1$ \\
第 $1$ 列 & $\{0\}$：$2$，\quad $\{2\}$：$1$，\quad $\{0, 1, 2\}$：$1$ \\
第 $2$ 列 & $\varnothing$：$4$，\quad $\{0, 1\}$：$1$，\quad $\{1, 2\}$：$2$ \\
第 $3$ 列 & $\varnothing$：$5$，\quad $\{0, 1\}$：$4$ \\
\bottomrule
\end{tabular}
\end{center}

第 $1$ 列之後的 $\{0\}$：$2$ 就是圖~\ref{fig:bm-profile} 的兩種情況。 處理第
$3$ 列時，$\{1, 2\}$ 的 $2$ 種被丟掉了：第 $3$ 列的第 $2$ 格是柱子，直磚戳不
進去。 第 $3$ 列之後的 $\{0, 1\}$ 表示有磚伸出地板，不算；答案是
$\varnothing$ 的 $5$。

\begin{itemize}
\item 直磚伸進下一列時沒有檢查下一列那格是不是柱子，留到處理下一列時用
\texttt{mask \& pillar[r]} 一起擋掉；最後一列伸出去的直磚，則由「答案只取
\texttt{cur[0]}」擋掉。
\item 每一列的計算只用到上一列的 \texttt{cur}，和第 1 本第~10~章的滾動陣列一
樣，只留兩排。
\end{itemize}

\textbf{複雜度。} 遞迴展開到最後的每一種結果對應一對 $(M, M')$：每一格不是
被上面蓋住（在 $M$ 裡）、往下放直磚（在 $M'$ 裡），就是橫磚的一半，而且剩下
的空格只有一種橫磚鋪法。 所以每一格三選一，一列最多 $3^w$ 種。 $w = 10$ 時每列最多約 $6 \times 10^4$ 種，
$h = 1000$ 列最多約 $6 \times 10^7$；實際上很多組合湊不出來，沒有柱子的
$1000 \times 10$ 在本書測試的環境下不到 $0.01$ 秒。 寬和高差很多時，讓比較窄
的那一邊當 $w$。

\textbf{地板很長的時候。} 實作練習的\akiref{tiling-a-pillared-hallway}{有柱子的長廊鋪磚}
長度到 $10^{18}$，一列一列做來不及。 那題的柱子每 $p$ 列重複一次，從一列到
下一列的轉移可以寫成一張 $2^w \times 2^w$ 的表（矩陣），整段長廊就是很多張表
連乘，再用第 1 本第~8~章 8.6 節快速冪的想法加速。 矩陣乘法第 3 本會正式介紹；
想先挑戰的讀者可以從「把 \texttt{cur} 到 \texttt{nxt} 的轉移寫成矩陣」開始。

\section{什麼時候用位元 DP}
\label{sec:bm-when}

第~3~章 3.1 節的狀態表，再加上這一章的三種：

\begin{center}
\begin{tabular}{p{0.27\linewidth}p{0.36\linewidth}p{0.27\linewidth}}
\toprule
狀態 & 轉移在看什麼 & 本章的例子 \\
\midrule
一個集合 $S$ & 下一個加進來的是誰；或一次挑出一整組（子集合的子集合） & 關主、出場順序、專題小組 \\
集合 $S$ 加上最後一個 $i$ & 從 $i$ 走到哪一個還沒去過的 $j$ & 跑腿路線、回收推車 \\
一列格子的樣子 & 這一列怎麼鋪，伸進下一列的是哪些格子 & 鋪地磚 \\
\bottomrule
\end{tabular}
\end{center}

適合想到位元 DP 的訊號：

\begin{itemize}
\item 某個量特別小：$n \le 15$、$16$、$20$，或格子的寬度不超過十幾，而其他
量很大。
\item 題目要的是排列或分組（排順序、走遍每個點、分成幾組），暴力解是 $n!$ 或
更多。
\item 列出狀態後，「之後怎麼做」只看「哪些已經做過」，不看做的順序。 回收推
車的重量就是這樣。
\end{itemize}

看起來像、其實不是的情況：

\begin{itemize}
\item \textbf{只在乎總和，不在乎是哪些。} 從 $n$ 件物品挑一些、總重不超過限
制，物品之間沒有交互作用，那是第~1~章的背包，狀態是「前 $i$ 件、用了多少重
量」，$n$ 可以很大。 物品之間有交互作用（某兩件不能同時拿、某幾件湊齊有加分），
$n$ 又很小，才需要整個集合。
\item \textbf{順序本來就固定。} 東西只能照輸入的順序處理，例如一排攤位由左到
右，狀態是「前 $i$ 個」或第~3~章的區間，不需要集合。
\item \textbf{$n$ 到幾十以上。} $2^{30}$ 就開不出表了。 這時要找別的結構：圖
是一棵樹、一條線，或者其實不用經過每個點。
\end{itemize}

第 1 本第~12~章的狀態 BFS 也常用遮罩，例如迷宮裡「拿到哪幾把鑰匙」：每一步代
價都是 $1$ 時用 BFS，代價不同、要計數或取最大值時用動態規劃。

\textbf{一個位置不只兩種情況時。} 遮罩的每一位只有 $0$、$1$。 如果每種東西
要記 $0$ 到 $3$ 個，可以每種用兩個位元，等於一個四進位數，$k$ 種東西就有
$4^k$ 個狀態。 這時遮罩的位數是 $2k$，$k = 16$ 就要 $32$ 位，\texttt{1 << i}
要改成 \texttt{1LL << i}，遮罩本身也要用 \texttt{long long}。 實作練習的
\akiref{orientation-prop-bundles}{迎新道具組合包} 是這種題目。

\section{本章常見錯誤}
\label{sec:bm-pitfalls}

各節的錯誤歸納成五類。 寫位元 DP 時，交出去之前逐項檢查。

\begin{center}
\begin{tabular}{p{0.18\linewidth}p{0.5\linewidth}p{0.2\linewidth}}
\toprule
錯誤 & 例子 & 見 \\
\midrule
位元運算 & \texttt{need[v] \& mask != need[v]} 少括號；\texttt{1 << 31} 超出
\texttt{int}；兩個位元一組時忘了 \texttt{1LL} & 第~\ref{sec:ex-bm-show}、\ref{sec:bm-when}~節 \\
不可能的狀態 & 排不出來的集合記成 $0$，從不合法的開頭接下去（印出 $24$） &
第~\ref{sec:ex-bm-show}~節 \\
計算順序 & 旅行推銷員的外層放 \texttt{last}（印出 $32$）；子集合的迴圈漏掉空集
合或停不下來 & 第~\ref{sec:bm-tsp}、\ref{sec:bm-submask}~節 \\
轉移的成本 & 回程忘了加、重量算錯時間點、題目要不要回到起點 &
第~\ref{sec:bm-tsp}、\ref{sec:ex-bm-cart}~節 \\
記憶體與數值 & $2^{20} \times 20$ 的表 $80$ MB、宣告在 \texttt{main} 裡當掉；費
力值超過 \texttt{int} & 第~\ref{sec:bm-memory}、\ref{sec:ex-bm-cart}~節 \\
\bottomrule
\end{tabular}
\end{center}

每一題都值得做的檢查：用 $n \le 8$ 的輸入，拿試遍所有排列
（\texttt{next\_permutation}）的暴力解對拍。

\section{本章練習}
\label{sec:bm-exercises}

\begin{drills}[label=drills:bm]
\item 用第~\ref{sec:bm-set}~節關主的分數表，改用往回拉的寫法：
$\mathrm{dp}(S) = \max_{v \in S} \bigl(\mathrm{dp}(S \setminus \{v\}) + a_{v,|S|-1}\bigr)$。
\begin{enumerate}
\item \texttt{mask} $=$ \texttt{1011} 代表哪幾位關主？ 他們顧的是哪幾關？
\item 照圖~\ref{fig:bm-lattice} 的值，算出 $\mathrm{dp}(\texttt{1011})$，寫出三個候
選和取到最大值的是哪一個。
\end{enumerate}

\item 設 \texttt{x} $= 52$（二進位 \texttt{110100}）。
\begin{enumerate}
\item 求 \texttt{x \& -x}、\texttt{x \& (x - 1)}、\texttt{\_\_builtin\_ctz(x)}、
\texttt{\_\_builtin\_popcount(x)}。
\item 有人想判斷第 $j$ 個人不在集合裡，寫成 \texttt{if (mask \& 1 << j == 0)}。
這個條件什麼時候成立？
\end{enumerate}

\item 估計下列表格的大小與執行次數。
\begin{enumerate}
\item 旅行推銷員 $n = 18$，$\mathrm{dp}(S, i)$ 用 \texttt{long long}：表格幾
MB？ 最內層大約執行幾次？
\item 第~\ref{sec:ex-bm-kit}~節的子集合 DP 如果 $n = 18$，固定最低位之後最內層
大約執行幾次？
\end{enumerate}

\item 照第~\ref{sec:bm-submask}~節的迴圈，列出 \texttt{mask} $=$ \texttt{11010}
的所有子集合，依程式走過的順序。 一共幾個？

\item 第~\ref{sec:bm-tsp-table}~節的跑腿表。
\begin{enumerate}
\item 如果辦完事不用回家，最後停在藥局（地點 $4$），最短要幾分鐘？
\item 如果不用回家、停在哪裡都可以呢？ 停在哪裡？
\end{enumerate}

\item 回收推車的費力值如果改成「距離 $\times$ 出發以來已經走過的總距離」，
$\mathrm{dp}(S, i)$ 還夠嗎？ 說明理由。

\item 用第~\ref{sec:bm-tiling}~節的方法，處理 $3$ 列 $2$ 行、沒有柱子的地板。
列出每一列處理完的狀態和方法數，最後共有幾種鋪法？
\end{drills}

\section{實作練習}
\label{sec:bm-practice}

下面是 Aki-Judge 上和本章相關的題目，依狀態的形狀分組，三題例題也列在裡面。
難度分四階：暖身直接套用本章的寫法；基本要自己定出狀態和轉移；實戰要先把題目
轉換成本章的形狀；挑戰會和其他章的技巧合在一起。

\begin{center}
\begin{tabular}{p{0.2\linewidth}p{0.56\linewidth}c}
\toprule
形狀 & 題目 & 難度 \\
\midrule
集合 & \akiref{running-order-for-the-club-showcase}{成果發表會的出場順序}（例題） & 暖身 \\
 & \akiref{color-chain-card-duel}{接色牌對戰} & 實戰 \\
\midrule
集合加最後一站 & \akiref{fair-patrol-route}{園遊會糾察的巡視路線} & 暖身 \\
 & \akiref{the-ever-heavier-recycling-cart}{越推越重的回收推車}（例題） & 基本 \\
 & \akiref{booth-hopping-at-the-club-fair}{社團博覽會的體驗時間} & 實戰 \\
 & \akiref{restocking-change-at-the-school-fair}{園遊會零錢補給} & 挑戰 \\
 & \akiref{errands-before-school-starts}{開學前的採買路線} & 挑戰 \\
\midrule
子集合的子集合 & \akiref{teams-with-every-piece-of-gear}{湊齊器材的專題小組}（例題） & 基本 \\
\midrule
其他遮罩 & \akiref{orientation-prop-bundles}{迎新道具組合包} & 挑戰 \\
 & \akiref{tiling-a-pillared-hallway}{有柱子的長廊鋪磚} & 挑戰 \\
\bottomrule
\end{tabular}
\end{center}

幾點提示：接色牌對戰是兩人輪流的遊戲，狀態除了剩下的牌，還要記牌鏈兩端露出的
顏色，值是「輪到的人會不會贏」。 園遊會零錢補給和開學前的採買路線的地圖很大，
要先用第~8~章的 Dijkstra 算出幾個重要地點之間的最短距離，再做位元 DP。 例題
讀完之後，可以先關上書自己寫一次，再上 Aki-Judge 提交。

\section{讀完你應該會}
\label{sec:bm-checklist}

\begin{itemize}
\item 看到某個量只有十幾、題目要排順序或分組時，能判斷要不要用集合當狀態，並
用一句中文寫出 $\mathrm{dp}(S)$ 或 $\mathrm{dp}(S, i)$ 的意思。
\item 能用 \texttt{|}、\texttt{\&}、\texttt{\^{}}、\texttt{x \& -x}、
\texttt{x \& (x - 1)}、\texttt{\_\_builtin\_popcount} 寫出加入、拿掉、檢查
子集合、找最低位，並在有比較時加上括號。
\item 能寫出旅行推銷員的 $\mathrm{dp}(S, i)$，依題目改初始值和答案（要不要回
起點、起點終點是否固定），並用 $\mathit{from}$ 找回路線。
\item 能用 \texttt{sub = (sub - 1) \& mask} 列舉子集合，說出總次數為什麼是 $3^n$。
\item 動手前能算出位元 DP 表格要幾 MB，知道 $2^{20} \times 20$ 個 \texttt{int}
是 $80$ MB。
\end{itemize}

\textbf{延伸閱讀}（書末會整理完整的參考資料）

\begin{itemize}
\item Antti Laaksonen，\textit{Guide to Competitive Programming}，第 3 版，6.2.4、
6.2.5 節：從排列改成子集合的動態規劃，以及一列一列數鋪磚方法的寫法，篇幅短。
\item Steven Halim 等，\textit{Competitive Programming 4}，Book 1 第 3.5.2 節的
旅行推銷員，和 Book 2 第 8 章更多用遮罩當狀態的題型。
\item Steven S. Skiena，\textit{The Algorithm Design Manual}，第 3 版，10.9 節：
用最長簡單路徑說明為什麼「沒有順序」的問題，動態規劃的狀態會變成指數個。
\end{itemize}

\addcontentsline{toc}{section}{練習解答}
\section*{練習解答}

\begin{enumerate}
\item
\begin{enumerate}
\item 第 $0$、$1$、$3$ 位是 $1$，代表 A、B、D，他們顧第 $0$ 到第 $2$ 關。
\item 最後一關是第 $2$ 關。 拿掉 A：$\mathrm{dp}(\{B, D\}) + a_{A,2} = 12 + 9 = 21$；
拿掉 B：$\mathrm{dp}(\{A, D\}) + a_{B,2} = 12 + 6 = 18$；拿掉 D：
$\mathrm{dp}(\{A, B\}) + a_{D,2} = 10 + 6 = 16$。 最大是 $21$，A 顧第 $2$ 關，
和圖上的值相同。
\end{enumerate}

\item
\begin{enumerate}
\item \texttt{x \& -x} $= 4$（\texttt{000100}）；\texttt{x \& (x - 1)} $= 48$
（\texttt{110000}）；\texttt{\_\_builtin\_ctz(x)} $= 2$；
\texttt{\_\_builtin\_popcount(x)} $= 3$。
\item \texttt{==} 比 \texttt{<<} 晚、比 \texttt{\&} 早算，條件是
\texttt{mask \& ((1 << j) == 0)}。 \texttt{1 << j} 永遠不是 $0$，括號裡永遠是
$0$，整個條件永遠不成立。 要寫成 \texttt{((mask >> j) \& 1) == 0}。
\end{enumerate}

\item
\begin{enumerate}
\item $2^{18} \times 18 \times 8 = 37{,}748{,}736$ 位元組，$36$ MB。 最內層約
$2^{18} \times 18 \times 18 \approx 8.5 \times 10^7$ 次。
\item $(3^{18} - 1)/2 \approx 1.9 \times 10^8$ 次，是 $n = 16$ 的九倍。 每次只做
一次查表和比較，在本書測試的環境下 $n = 18$ 的隨機輸入約 $0.13$ 秒，還來得及；
$n = 20$ 再多九倍，就很緊了。
\end{enumerate}

\item \texttt{11010}、\texttt{11000}、\texttt{10010}、\texttt{10000}、
\texttt{01010}、\texttt{01000}、\texttt{00010}、\texttt{00000}，共 $8$ 個，
$= 2^3$。

\item
\begin{enumerate}
\item $\mathrm{dp}(\{1, 2, 3, 4\}, 4) = 25$ 分鐘（家 $\to 3 \to 1 \to 2 \to 4$，
$3 + 4 + 9 + 9$）。
\item 表的最後一列取最小，$\mathrm{dp}(\{1, 2, 3, 4\}, 2) = 23$，停在五金行。 例
如家 $\to 4 \to 3 \to 1 \to 2$：$4 + 6 + 4 + 9 = 23$。
\end{enumerate}

\item 不夠。 已經走過的總距離和順序有關，不是 $S$ 的函數。 走進同一個
$(S, i)$ 的兩條路線，可能一條費力值比較小、但走過的距離比較長，另一條相反；
之後每一段都要乘上那個距離，哪一條好要看後面，沒辦法只留一條。 回收推車能用
$\mathrm{dp}(S, i)$，正是因為重量只看 $S$。

\item 開始 $\varnothing$：$1$。 第 $0$ 列（兩塊直磚或一塊橫磚）：
$\{0, 1\}$：$1$，$\varnothing$：$1$。 第 $1$ 列：$\varnothing$ 又分出兩種，
$\{0, 1\}$ 整列已被蓋住、只得到 $\varnothing$，合計 $\varnothing$：$2$，
$\{0, 1\}$：$1$。 第 $2$ 列：$\varnothing$：$3$，$\{0, 1\}$：$2$。 答案是
$\varnothing$ 的 $3$ 種。
\end{enumerate}

\chapter{數位 DP}
\label{ch:digitdp}

有些大樓的樓層、醫院的病床號碼會跳過含有 $4$ 的數字。 假設一家連鎖咖啡店的會
員卡號也這樣發：從 $1$ 號開始，遇到 $4$、$14$、$40$ 到 $49$、$104$ 這類含有
$4$ 的號碼就跳過。 有兩個問題：卡號 $3517$ 以內實際發出了幾張卡？ 第 $2526$
位會員拿到的是幾號？ 往下讀之前先想一個做法。 再想一下，如果卡號是
$10^{18}$，你的做法還來得及嗎？

從 $1$ 數到 $3517$，一個一個檢查，電腦一下子就做完了。 數到 $10^{18}$ 就不行：
照第 1 本第~1~章的估計，一秒大約 $10^8$ 次簡單運算，$10^{18}$ 次要三百多年。
這一章的做法不看一個一個的數，改看數字的每一\textbf{位}：一次決定一位，同一
類的數整批算。 這種做法叫\textbf{數位 DP}（digit DP）。 $10^{18}$ 只有 $19$
位，每一位最多試 $10$ 個數字，所以運算次數和位數有關，和數值的大小無關。

《C++ 演算法解題 1：APCS 實作題核心》（以下簡稱第 1 本）第~8~章 8.3 節的記憶化、
第 1 本第~9~章的動態規劃四件事（狀態、轉移、初始值、計算順序），這一章都直接
拿來用。 第~\ref{sec:dg-prefix}~節先不寫動態規劃，用手算解決上面的第一個問題，
看出「比 $N$ 小的數長什麼樣子」。 第~\ref{sec:dg-skeleton}~節把它寫成數位 DP
的固定骨架，並說清楚哪些狀態可以記住。 之後幾節依序加上前導零、多記一點資訊、
「不含某一串數字」的狀態，以及反過來找第 $K$ 個數。 最後兩節處理兩個數一起填，
以及很大的數和溢位。

\section{比 \texorpdfstring{$N$}{N} 小的數長什麼樣子}
\label{sec:dg-prefix}

\subsection{先寫暴力解，再換成兩個前綴}

會員卡的第一個問題，最直接的寫法是從 $1$ 跑到 $N$，把每個數轉成字串，看裡面有
沒有字元 \texttt{'4'}。 這個暴力解很短，而且一定對，之後拿來對拍（用暴力解在
小範圍比對）正好。 它的時間是 $O(N \log N)$：$N$ 個數，每個數有大約
$\log_{10} N$ 位。

題目常常問一段區間 $[L, R]$ 裡有幾個數符合條件。 這時先定義
\[
f(N) = 0 \text{ 到 } N \text{ 之中符合條件的個數},
\]
區間的答案就是 $f(R) - f(L - 1)$，和第 1 本第~5~章的前綴和是同一個想法。 這
樣只要會算「$0$ 到 $N$」一種問題。 $L = 0$ 時 $L - 1 = -1$，規定 $f(-1) = 0$，
程式裡要先判斷，不要把 $-1$ 轉成字串去算。

\subsection{第一個不一樣的位置}

把 $0$ 到 $N$ 的每個數都在左邊補 $0$，補到和 $N$ 一樣長。 以 $N = 3517$ 為
例，$7$ 寫成 \texttt{0007}，$352$ 寫成 \texttt{0352}。 任何一個比 $N$ 小的數
$x$，從最高位往右比，一定有一個\textbf{第一個不一樣的位置}：在這之前，$x$ 的每
一位都和 $N$ 一樣；在這一位，$x$ 的數字比 $N$ 的小；在這之後，$x$ 的每一位都
可以是 $0$ 到 $9$ 的任何數字，因為前面已經比較小了，後面再大也不會超過 $N$。

所以比 $N$ 小的數可以照「第一個不一樣的位置」和「那一位填多少」分成很多塊，
每一塊的形狀都是：一段和 $N$ 相同的開頭、一個比較小的數字、後面全部自由。
$N$ 本身另外算一個。 以 $3517$ 數不含 $4$ 的個數，後面自由的每一位有 $9$ 種
選法（不能是 $4$）：

\begin{center}
\begin{tabular}{llll}
\toprule
第一個不一樣的位置 & 這一位可以填 & 這一塊的樣子 & 不含 $4$ 的個數 \\
\midrule
千位（$N$ 是 $3$） & $0, 1, 2$ & \texttt{0???}、\texttt{1???}、\texttt{2???} & $3 \times 9^3 = 2187$ \\
百位（$N$ 是 $5$） & $0, 1, 2, 3$（跳過 $4$） & \texttt{30??} 到 \texttt{33??} & $4 \times 9^2 = 324$ \\
十位（$N$ 是 $1$） & $0$ & \texttt{350?} & $1 \times 9 = 9$ \\
個位（$N$ 是 $7$） & $0, 1, 2, 3, 5, 6$ & \texttt{3510} 到 \texttt{3516} & $6$ \\
沒有（就是 $N$） & & \texttt{3517} & $1$ \\
\bottomrule
\end{tabular}
\end{center}

加起來是 $2527$，這是 $0$ 到 $3517$ 的個數；$0$ 不是卡號，扣掉之後是 $2526$
張。 每一塊互不重疊，因為一個數的「第一個不一樣的位置」只有一個；所有比 $N$ 小
的數也都落在某一塊裡。 所以這樣分塊，每個數剛好算一次。

如果 $N$ 的某一位就是 $4$，例如 $N = 3457$，走到百位的 $4$ 就要停：開頭是
\texttt{34} 的數都含 $4$，再往下走的塊全部不能要，$N$ 本身也不能算。

\begin{example}[label=ex:dg-no4]{cpp}{(no4.cpp) 1 到 N 之中不含 4 的數}
#include <iostream>
#include <string>
using namespace std;

int main() {
    long long n;
    cin >> n;
    string s = to_string(n);
    int len = s.size();
    long long pw[20];                 // pw[r] = 9 的 r 次方
    pw[0] = 1;
    for (int r = 1; r < 20; r++) {
        pw[r] = pw[r - 1] * 9;
    }
    long long count = 0;              // 0 到 n 之中不含 4 的個數
    bool stopped = false;
    for (int i = 0; i < len; i++) {
        int c = s[i] - '0';
        for (int d = 0; d < c; d++) {
            if (d != 4) {
                // 前面和 n 一樣，這一位填 d，後面任意
                count += pw[len - 1 - i];
            }
        }
        if (c == 4) {                 // 再貼著 n 走，一定含 4
            stopped = true;
            break;
        }
    }
    if (!stopped) {
        count++;                      // n 本身也不含 4
    }
    cout << count - 1 << '\n';        // 扣掉 0
    return 0;
}
\end{example}

\iosample{
3517 \\
}{
2526 \\
}

\begin{itemize}
\item 迴圈的第 $i$ 輪處理「第一個不一樣的位置是第 $i$ 位」的那幾塊，後面還有
\texttt{len - 1 - i} 位自由。 一輪做完，等於確定第 $i$ 位和 $N$ 一樣，往下一
位走。
\item 這支程式也數了「有前導零」的數，例如 \texttt{0007}。 這裡沒有問題，因為
補上去的 $0$ 不是 $4$，不會改變答案。 第~\ref{sec:dg-started}~節會看到前導零
會影響答案的題目。
\item 時間是 $O(10 \times \text{位數})$，$N = 10^{18}$ 也只有 $19$ 位。
\end{itemize}

這一節的做法有一個前提：後面自由的那幾位有幾種填法，算得出一個公式（這裡是
$9^r$）。 條件一複雜，例如「數字和是 $7$ 的倍數」，後面幾位的填法數就和前面
填了什麼有關，沒有簡單的公式。 下一節用動態規劃處理這件事。

\section{數位 DP 的骨架}
\label{sec:dg-skeleton}

換一個問題：$L$ 到 $R$ 之中，各位數字的和是 $k$ 的倍數的數有幾個？
$R \le 10^{18}$，$k \le 100$。 例如 $k = 5$ 時，$32$（和 $5$）、$136$（和
$10$）都算，$36$ 不算。

照上一節的分塊，第一個不一樣的位置之後的每一位仍然自由，可是「自由的那幾位要
填出多少的和」取決於前面的數字和。 前面加起來除以 $k$ 餘 $2$，後面就要補一個
除以 $k$ 餘 $k - 2$ 的和。 所以除了「填到第幾位」，還要記住目前的數字和除以
$k$ 的餘數。

\subsection{四件事}

把 $N$ 寫成字串 \texttt{s}，長度 $n$，從最高位開始一位一位填。

\begin{itemize}
\item \textbf{狀態}：$\mathrm{dfs}(\mathit{pos}, \mathit{rem}, \mathit{tight})$
＝ 已經填好前 $\mathit{pos}$ 位、這幾位的和除以 $k$ 餘 $\mathit{rem}$ 時，把剩
下的位數填完、讓整個數不超過 $N$ 而且數字和是 $k$ 的倍數，有幾種填法。
$\mathit{tight}$ 是一個真假值，表示前 $\mathit{pos}$ 位是不是每一位都和 $N$
一樣，也就是還\textbf{貼著上界}。
\item \textbf{轉移}：這一位填 $d$。 貼著上界時 $d$ 只能從 $0$ 填到 $N$ 的這一
位 \texttt{s[pos]}，否則 $0$ 到 $9$ 都可以。 填完之後餘數變成
$(\mathit{rem} + d) \bmod k$；原本貼著、而且 $d$ 剛好等於 \texttt{s[pos]}，才
繼續貼著。 把每個 $d$ 的結果加起來。
\item \textbf{初始值}：$\mathit{pos} = n$ 時已經填完，是一個數。 餘數是 $0$ 就
算一個，否則不算。
\item \textbf{計算順序}：$\mathit{pos}$ 只會變大，用記憶化的遞迴最自然。 答案
是 $\mathrm{dfs}(0, 0, \text{true})$：什麼都還沒填時，當然和 $N$ 一樣（空的開
頭），算是貼著上界。
\end{itemize}

$\mathit{tight}$ 就是上一節的「第一個不一樣的位置」換個說法：$\mathit{tight}$
從真變成假的那一位，就是第一個比 $N$ 小的位置。 貼著上界的路只有一條，沿著
$N$ 的數字走；每一步往旁邊岔出去，就是一整塊後面自由的數。 下圖是 $N = 347$
的情形，三層岔出去的塊加上 $N$ 本身，就是 $0$ 到 $347$ 的全部：

\begin{center}
\begin{tikzpicture}[x=1cm, y=1cm,
    blk/.style={draw=black!55, fill=black!5, font=\small, align=center, inner sep=4pt}]
\node[cell, minimum width=1.2cm] (r) at (0,0) {開頭};
\node[hlcell] (a) at (3,0) {3};
\node[hlcell] (b) at (6,0) {4};
\node[hlcell] (c) at (9,0) {7};
\draw[diagarrow, draw=memStroke] (r) -- node[above, font=\small] {填 $3$} (a);
\draw[diagarrow, draw=memStroke] (a) -- node[above, font=\small] {填 $4$} (b);
\draw[diagarrow, draw=memStroke] (b) -- node[above, font=\small] {填 $7$} (c);
\node[blk] (x) at (0,-2) {\texttt{0??}、\texttt{1??}、\texttt{2??}\\後面 $2$ 位自由};
\node[blk] (y) at (3.3,-2) {\texttt{30?} 到 \texttt{33?}\\後面 $1$ 位自由};
\node[blk] (z) at (6.4,-2) {\texttt{340} 到 \texttt{346}\\沒有自由的位};
\node[blk] (w) at (9.4,-2) {\texttt{347}\\$N$ 本身};
\draw[diagarrow] (r) -- node[left, font=\small] {填 $0$ 到 $2$} (x);
\draw[diagarrow] (a) -- node[left, font=\small] {填 $0$ 到 $3$} (y);
\draw[diagarrow] (b) -- node[left, font=\small] {填 $0$ 到 $6$} (z);
\draw[diagarrow] (c) -- (w);
\end{tikzpicture}
\diagcaption{$N = 347$。 上面一排是貼著上界的狀態，沿著 $N$ 的數字走；每往下
岔出去一次，就離開了上界，後面的位數都自由。}
\label{fig:dg-tight-path}
\end{center}

$k = 5$ 時每一塊各有幾個：

\begin{center}
\begin{tabular}{lccc}
\toprule
塊 & 前面的和 $\bmod 5$ & 後面要補的和 $\bmod 5$ & 個數 \\
\midrule
\texttt{0??}、\texttt{1??}、\texttt{2??} & $0, 1, 2$ & $0, 4, 3$ & 各 $20$ \\
\texttt{30?} 到 \texttt{33?} & $3, 4, 0, 1$ & $2, 1, 0, 4$ & 各 $2$ \\
\texttt{340} 到 \texttt{346} & $(7 + d) \bmod 5$ & 沒有自由的位 & $1$（\texttt{343}） \\
\texttt{347} & $4$ & & $0$ \\
\bottomrule
\end{tabular}
\end{center}

兩位數字 \texttt{00} 到 \texttt{99} 的和除以 $5$，五種餘數各有 $20$ 個；一位
數字 $0$ 到 $9$，五種餘數各 $2$ 個。 \texttt{340} 到 \texttt{346} 那一塊沒有
自由的位，前面 $3 + 4$ 加上個位 $d$，只有 $d = 3$ 時和是 $10$。 \texttt{347}
的數字和是 $14$，不算。 合計 $60 + 8 + 1 = 69$，這是 $f(347)$，其中包含 $0$。
問 $1$ 到 $347$ 時，再減掉 $f(0) = 1$，得到 $68$。

\begin{example}[label=ex:dg-digitsum]{cpp}{(digitsum.cpp) 數字和是 k 的倍數：記憶化的寫法}
#include <cstring>
#include <iostream>
#include <string>
using namespace std;

int k;
string s;                    // 上界 N 的十進位寫法
// memo[剩下幾位][數字和 mod k]，-1 表示還沒算
long long memo[20][100];

// 從第 pos 位填到最後一位。 前面填的數字和除以 k 餘 rem；
// tight 表示前面每一位都和 N 一樣。
long long dfs(int pos, int rem, bool tight) {
    int n = s.size();
    if (pos == n) {                       // 填完了，這是一個數
        return rem == 0 ? 1 : 0;
    }
    if (!tight && memo[n - pos][rem] != -1) {
        return memo[n - pos][rem];
    }
    int up = tight ? s[pos] - '0' : 9;    // 這一位最多填多少
    long long total = 0;
    for (int d = 0; d <= up; d++) {
        total += dfs(pos + 1, (rem + d) % k, tight && d == up);
    }
    if (!tight) {
        memo[n - pos][rem] = total;
    }
    return total;
}

long long f(long long N) {       // 0 到 N 之中符合的個數
    if (N < 0) {
        return 0;
    }
    s = to_string(N);
    return dfs(0, 0, true);
}

int main() {
    long long L, R;
    cin >> L >> R >> k;
    memset(memo, -1, sizeof(memo));
    cout << f(R) - f(L - 1) << '\n';
    return 0;
}
\end{example}

\iosample{
1 347 5 \\
}{
68 \\
}

\begin{itemize}
\item \texttt{up} 是這一位能填的最大數字。 新的 \texttt{tight} 寫成
\texttt{tight \&\& d == up}：不貼著的時候 \texttt{up} 是 $9$，填 $9$ 也不會變
回貼著，靠的是前面的 \texttt{tight} 已經是假。
\item \texttt{memset} 把每個位元組設成 \texttt{0xFF}，\texttt{long long} 的
八個位元組都是 \texttt{0xFF} 剛好是 $-1$，所以可以用來把整張表設成「還沒算」。
這個技巧只對 $0$ 和 $-1$ 有用。
\item $10^{18}$ 有 $19$ 位，\texttt{memo} 的第一維開 $20$。
\end{itemize}

\subsection{哪些狀態可以記住}
\label{sec:dg-memo}

程式只在 \texttt{tight} 是假的時候讀寫 \texttt{memo}，原因有兩個。

第一，貼著上界的狀態不用記。 從圖~\ref{fig:dg-tight-path} 看，貼著上界的路只
有一條，每一層只有一個貼著的狀態，每個只會走到一次，記了也不會再用到。

第二，貼著上界的狀態\textbf{不能}和不貼的共用同一格。 不貼的時候，答案只和「剩
下幾位」「目前餘數」有關，和 $N$ 長什麼樣子無關：剩下兩位、餘數 $1$，不管 $N$
是 $347$ 還是 $999$，答案都是「兩位數字的和補成餘數 $0$ 的填法數」。 貼著的
時候，剩下的位數還受 $N$ 的數字限制，答案不一樣。 如果把兩種狀態存進同一格，
就會把「受限」的答案拿去給「自由」的狀態用，或者反過來。

不貼的狀態和 $N$ 無關，還有一個好處：\texttt{memo} 用「剩下幾位」
\texttt{n - pos} 當索引，算 $f(R)$ 存下來的值，算 $f(L - 1)$ 時可以直接拿來
用，不用清空。 如果改用「目前是第幾位」\texttt{pos} 當索引，$R$ 和 $L - 1$ 的
位數不同時，同一個 \texttt{pos} 代表的剩餘位數就不同，又沒有清空的話會讀到錯的
值。 把範例~\ref{ex:dg-digitsum} 的 \texttt{memo[n - pos]} 都改成
\texttt{memo[pos]}，輸入 \texttt{100 347 5} 的正確答案是 $49$，程式卻印出
\texttt{-59}：$f(99)$ 在 \texttt{pos = 1}（只剩一位）讀到的是 $f(347)$ 存的
「剩兩位」的值。

不過有一種情況還是要清空：狀態的意義依賴題目給的其他數。 這一題的 \texttt{memo}
依賴 $k$。 如果一個輸入有好幾組測試資料、每組的 $k$ 不一樣，每組開始時都要重新
\texttt{memset}。

狀態數：剩下的位數最多 $19$ 種、餘數 $k$ 種，每個狀態試 $10$ 個數字，總共大約
$19 \times 100 \times 10 = 1.9 \times 10^4$ 次，再加上貼著上界的那 $19$ 個狀
態。 和 $N$ 是 $10^{18}$ 相比，小得可以忽略。

\subsection{改成填表}
\label{sec:dg-table}

記憶化的遞迴是從答案往回問。 也可以像第 1 本第~9~章那樣先把表填好。 不貼上界
的部分和 $N$ 無關，所以先算一張「自由的位數」的表：
\[
\mathrm{ways}(r, m) = r \text{ 個數字（每個都可以是 } 0 \text{ 到 } 9 \text{）的和除以 } k \text{ 餘 } m \text{ 的填法數}.
\]
初始值 $\mathrm{ways}(0, 0) = 1$（零個數字，和是 $0$），轉移是看最前面那個數字
填多少。 表填好之後，沿著 $N$ 的數字走一遍，就是第~\ref{sec:dg-prefix}~節的分
塊：每一位試所有比 \texttt{s[i]} 小的 $d$，那一塊的個數直接查表。

\begin{example}[label=ex:dg-digitsum-table]{cpp}{(digitsum\_table.cpp) 數字和是 k 的倍數：填表的寫法}
#include <iostream>
#include <string>
using namespace std;

int k;
// ways[r][m]：r 個數字（可以是 0）的和除以 k 餘 m 的填法數
long long ways[19][100];

long long f(long long N) {   // 0 到 N 之中符合的個數
    if (N < 0) {
        return 0;
    }
    string s = to_string(N);
    int n = s.size();
    long long total = 0;
    int pre = 0;             // 貼著 N 走到這裡，前面的數字和 mod k
    for (int i = 0; i < n; i++) {
        int c = s[i] - '0';
        for (int d = 0; d < c; d++) {   // 填比 N 小的 d，後面自由
            int m = (pre + d) % k;
            // 後面的和要補成 k 的倍數
            total += ways[n - 1 - i][(k - m) % k];
        }
        pre = (pre + c) % k;
    }
    if (pre == 0) {
        total++;             // N 本身
    }
    return total;
}

int main() {
    long long L, R;
    cin >> L >> R >> k;
    ways[0][0] = 1;          // 0 個數字，和是 0
    for (int r = 1; r <= 18; r++) {
        for (int m = 0; m < k; m++) {
            for (int d = 0; d <= 9; d++) {   // 最前面那一位填 d
                ways[r][(m + d) % k] += ways[r - 1][m];
            }
        }
    }
    cout << f(R) - f(L - 1) << '\n';
    return 0;
}
\end{example}

\iosample{
1 347 5 \\
}{
68 \\
}

\begin{itemize}
\item 前面的和餘 $m$ 時，後面要補的是 $(k - m) \bmod k$。 最外層再取一次
\texttt{\% k}，是為了 $m = 0$ 時得到 $0$ 而不是 $k$。
\item 表只開到 $r = 18$。 $N \le 10^{18}$ 最多 $19$ 位，第一位之後最多剩 $18$
位自由。 $\mathrm{ways}(r, \cdot)$ 加起來是 $10^r$，$r = 19$ 時 $k = 1$ 的那一
格就是 $10^{19}$，超過 \texttt{long long}。 把迴圈改成 $r \le 19$，用
\texttt{-fsanitize=undefined} 編譯後輸入 \texttt{1 347 1}，會看到
\texttt{signed integer overflow} 的錯誤訊息；沒有開這個選項的話，程式默默算出
一個沒有用到的錯誤值。 不用的格子就不要算。
\item 「$N$ 本身」要另外加，因為迴圈只算了比 $N$ 小的數。
\end{itemize}

兩種寫法算的是同一個東西。 記憶化的寫法比較接近「一位一位決定」的想法，加狀
態時只要改參數；填表的寫法把和 $N$ 無關的部分抽出來，同一張表可以拿來回答很多
個不同的 $N$，第~\ref{sec:dg-kth}~節找第 $K$ 個數時就要用到這一點。 這一章的
例題兩種都會用。

\section{前導零}
\label{sec:dg-started}

前兩節都把短的數補成和 $N$ 一樣長，因為補上去的 $0$ 不會改變答案：$0$ 不是
$4$，加進數字和也是加 $0$。 很多條件沒有這麼幸運。 數「含幾個 $0$」時，$7$
寫成 \texttt{007} 會多出兩個 $0$；看「相鄰兩位」時，$5$ 寫成 \texttt{05} 會多
出一組 \texttt{0} 和 \texttt{5}。

解法是再加一個狀態 $\mathit{started}$：前面是不是已經寫下第一個非零的數字。
還沒開始的時候，這一位填 $0$ 代表「這個數比較短，這一位不存在」，不算任何東西；
填一個非零數字，數就從這一位開始。 開始之後，填的 $0$ 才是真的 $0$。 一路都
沒有開始，表示每一位都是 $0$，這是整數 $0$，要依題目另外判斷它算不算。

\subsection{例題：恰好 \texorpdfstring{$k$}{k} 個零的整數}
\label{sec:dg-ex-zeros}

\textbf{題意}（\akiref{integers-with-exactly-k-zeros}{恰好 k 個零的整數}）

\begin{itemize}
\item 給 $L$、$R$、$k$，數 $L$ 到 $R$ 的整數之中，十進位寫法裡恰好有 $k$ 個
數字 $0$ 的有幾個。
\item 寫法照一般習慣：正整數沒有前導零，所以 $7$ 沒有 $0$，$100$ 有兩個；整數
$0$ 寫成 \texttt{0}，有一個 $0$。
\item $0 \le L \le R \le 10^{18}$，$0 \le k \le 18$。
\end{itemize}

例如 $L = 0$、$R = 120$、$k = 1$：整數 $0$ 算一個；兩位數的 $10, 20, \ldots,
90$ 九個；三位數的 $101$ 到 $109$ 九個（$100$ 有兩個 $0$），以及 $110$、
$120$。 共 $21$ 個。

\textbf{想法}

暴力解從 $L$ 跑到 $R$，每個數數一次 $0$。 $R - L$ 可以到 $10^{18}$，來不及。
照前兩節，答案寫成 $f(R) - f(L - 1)$，$f(N)$ 是 $0$ 到 $N$ 之中恰有 $k$ 個 $0$
的個數，$f(-1) = 0$。

狀態要記：填到第幾位、目前寫了幾個 $0$、是否貼著上界、是否已經開始。

\begin{itemize}
\item 這一位填 $d$。 已經開始而且 $d = 0$，$0$ 的個數加一；還沒開始時填 $0$ 不
算。 $d \ne 0$ 時，從這一位開始（已經開始的維持開始）。
\item 填完時，如果一直沒有開始，這是整數 $0$，只有 $k = 1$ 時算；否則看 $0$ 的
個數是不是剛好 $k$。
\item $0$ 的個數超過 $k$ 就可以直接回傳 $0$，後面再填也不會變少。
\end{itemize}

記憶化只記「不貼上界、已經開始」的狀態，用 \texttt{memo[剩下幾位][0 的個數]}。
還沒開始的狀態和已經開始的不能共用：同樣剩兩位、目前沒有 $0$，已經開始的話
\texttt{?0} 的 $0$ 要算，還沒開始的話 \texttt{0?} 的 $0$ 不算，填法數不同。
還沒開始的狀態每一層也只有一個（前面全是 $0$、不貼上界），不記也不慢。

$N = 120$、$k = 1$ 時，照前兩位分開來看：

\begin{center}
\begin{tabular}{llll}
\toprule
前兩位 & 涵蓋的數 & 走到第三位時的狀態 & 符合的個數 \\
\midrule
\texttt{00} & $0$ 到 $9$ & 還沒開始 & $1$（只有整數 $0$） \\
\texttt{01} 到 \texttt{09} & $10$ 到 $99$ & 已開始、$0$ 個 $0$、不貼 & 各 $1$，共 $9$ \\
\texttt{10} & $100$ 到 $109$ & 已開始、$1$ 個 $0$、不貼 & $9$ \\
\texttt{11} & $110$ 到 $119$ & 已開始、$0$ 個 $0$、不貼 & $1$ \\
\texttt{12} & $120$ & 貼著上界 & $1$ \\
\bottomrule
\end{tabular}
\end{center}

第二列的九種開頭走到同一個狀態（剩一位、$0$ 個 $0$），只有第一次真的算，後面
八次都從 \texttt{memo} 拿；第四列 \texttt{11} 也是這個狀態。 合計 $21$。

同樣的模型也出現在別的地方：發票號碼、車牌號碼、訂單編號裡某個數字出現幾次，
都是「數字的某個性質在一個大範圍裡出現幾次」。

\textbf{程式}

\begin{example}[label=ex:dg-zeros]{cpp}{(zeros.cpp) 恰好 k 個零的整數}
#include <cstring>
#include <iostream>
#include <string>
using namespace std;

int k;
string s;
// memo[剩下幾位][已經寫了幾個 0]，只存 !tight 且 started 的狀態
long long memo[20][20];

// started：前面是否已經寫下第一個非零數字；
// 還沒開始時填的 0 是前導零
long long dfs(int pos, int zeros, bool tight, bool started) {
    if (zeros > k) {
        return 0;                          // 0 已經太多了
    }
    int n = s.size();
    if (pos == n) {
        if (!started) {          // 每一位都是 0：整數 0，寫作 "0"
            return k == 1 ? 1 : 0;
        }
        return zeros == k ? 1 : 0;
    }
    bool useMemo = !tight && started;
    if (useMemo && memo[n - pos][zeros] != -1) {
        return memo[n - pos][zeros];
    }
    int up = tight ? s[pos] - '0' : 9;
    long long total = 0;
    for (int d = 0; d <= up; d++) {
        // 已經開始之後的 0 才算，前導零不算
        int nz = zeros + ((started && d == 0) ? 1 : 0);
        bool ns = started || d != 0;
        total += dfs(pos + 1, nz, tight && d == up, ns);
    }
    if (useMemo) {
        memo[n - pos][zeros] = total;
    }
    return total;
}

long long f(long long N) {    // 0 到 N 之中恰有 k 個 0 的個數
    if (N < 0) {
        return 0;
    }
    s = to_string(N);
    return dfs(0, 0, true, false);
}

int main() {
    long long L, R;
    cin >> L >> R >> k;
    memset(memo, -1, sizeof(memo));
    cout << f(R) - f(L - 1) << '\n';
    return 0;
}
\end{example}

\iosample{
0 120 1 \\
}{
21 \\
}

\begin{itemize}
\item \texttt{started || d != 0} 是新的 $\mathit{started}$：已經開始的維持開始，
還沒開始的遇到非零數字才開始。
\item \texttt{memo} 依賴 $k$（結尾比的是 \texttt{zeros == k}），這題一個輸入
只有一個 $k$，所以只在 \texttt{main} 開頭清一次。
\item 答案最大大約 $3.2 \times 10^{17}$（$L = 0$、$R = 10^{18}$、$k = 1$ 時），
要用 \texttt{long long}。
\end{itemize}

\textbf{複雜度}

會記憶化的狀態最多 $19 \times 19$ 個，每個試 $10$ 個數字；貼著上界和還沒開始
的狀態各只有 $19$ 個左右。 算兩次 $f$，總共幾千次運算，最大的測試資料在本書測
試的環境下不到 $0.01$ 秒。

\subsection{常見錯誤}

\begin{itemize}
\item \textbf{前導零也算成 $0$}：把 \texttt{nz} 寫成 \texttt{zeros + (d == 0 ?
1 : 0)}，等於把每個數都補成和 $N$ 一樣長再數。 輸入 \texttt{95 1005 2} 的正
確答案是 $14$，程式印出 \texttt{248}：$95$ 被看成 \texttt{0095}，$7$ 被看成
\texttt{0007}，全都多了 $0$。
\item \textbf{忘了整數 $0$}：結尾寫成只有 \texttt{return zeros == k ? 1 : 0;}，
整數 $0$ 的 \texttt{zeros} 是 $0$，被當成「沒有 $0$」。 範例的輸入印出
\texttt{20}，少了 $0$ 本身；$k = 0$ 時則會多算一個。
\item \textbf{還沒開始的狀態也記}：\texttt{useMemo} 寫成只看 \texttt{!tight}。
範例的輸入剛好還是 $21$，輸入 \texttt{95 1005 2} 卻印出 \texttt{5}。 算
$f(1005)$ 時，還沒開始、剩兩位、$0$ 個 $0$ 的狀態（\texttt{00??} 的後兩位）先存
進表裡；接著已經開始的 \texttt{01??}、\texttt{02??} 等等走到同一格，讀到的是「後
兩位的 $0$ 是前導零、不算」的答案。 $100, 200, \ldots, 900$ 這九個就這樣漏掉
了。
\item \textbf{貼著上界的狀態也記}：\texttt{useMemo} 寫成只看 \texttt{started}。
輸入 \texttt{1 20 0} 的正確答案是 $18$（$1$ 到 $9$、$11$ 到 $19$），程式印出
\texttt{27}：貼著上界的 \texttt{2?} 讀到了不貼的 \texttt{1?} 存的值，把 $21$
到 $29$ 也算了進去。 兩個範例的輸入都剛好答對，這種錯要靠對拍。
\item \textbf{$L - 1$ 用無號整數}：把 $L$、$R$ 宣告成 \texttt{unsigned long
long}，並拿掉 \texttt{N < 0} 的判斷。 $L = 0$ 時 \texttt{L - 1} 變成
$2^{64} - 1$，範例的輸入在本書測試的環境下印出 \texttt{-5394818876798596471}。
\end{itemize}

\section{狀態多帶一點資訊}
\label{sec:dg-extra}

到目前為止，除了位置、$\mathit{tight}$、$\mathit{started}$ 之外，狀態只多記了
一個數：數字和的餘數，或是 $0$ 的個數。 設計數位 DP 時，要問的問題都一樣：一
位一位由左往右填，填到一半時，\textbf{要知道前面的什麼，才能判斷後面怎麼填會
符合條件}？ 常見的有這幾種：

\begin{center}
\begin{tabular}{p{0.38\linewidth}p{0.3\linewidth}p{0.2\linewidth}}
\toprule
條件 & 要多記的東西 & 可能的值 \\
\midrule
數字和、整個數是 $k$ 的倍數 & 數字和或整個數除以 $k$ 的餘數 & $k$ 種 \\
相鄰兩位有關（奇偶交錯、遞增） & 前一位數字 & $10$ 種 \\
某個數字恰好、最多出現幾次 & 已經出現幾次 & 最多 $19$ 種 \\
同一個數字不能連續太多次 & 前一位數字、目前連續幾次 & $10 \times$ 上限 \\
數字裡不能出現某一串 & 目前對上那一串的幾個字 & 那一串的長度 \\
\bottomrule
\end{tabular}
\end{center}

整個數除以 $k$ 的餘數也可以一位一位算：前面的數是 $x$、這一位填 $d$，新的數是
$10x + d$，餘數是 $(10 \cdot \mathit{rem} + d) \bmod k$。 最後一列在第
~\ref{sec:dg-automaton}~節講。

估計狀態數的方法和第 1 本第~9~章一樣，把每一個量可能的值相乘，再乘上每個狀態試
的 $10$ 個數字。 前一位和連續次數一起記，是 $19 \times 10 \times 5 \times 10$
左右，都很小。 數位 DP 很少因為狀態太多而太慢，比較常見的問題是狀態少記了一
樣，或者 $\mathit{started}$、$\mathit{tight}$ 沒處理好。

\subsection{要的是總和，不是個數}

有些題目問的是「每個數都有一個量，全部加起來是多少」，
例如 $1$ 到 $N$ 所有數的數字和。 這時每個狀態要回傳兩個數：

\begin{itemize}
\item $\mathit{cnt}$：從這個狀態往後，有幾種填法；
\item $\mathit{sum}$：這些填法在後面產生的量，全部加起來。
\end{itemize}

關鍵在於：這一位造成的量，會出現在後面\textbf{每一種}填法裡。 以數字和為例，
這一位填 $d$，後面有 $c$ 種填法，這個 $d$ 就在 $c$ 個數裡各加了一次，總共是
$d \times c$。 所以
\[
\mathit{sum} \mathrel{+}= \mathit{sum}' + d \cdot \mathit{cnt}', \qquad
\mathit{cnt} \mathrel{+}= \mathit{cnt}',
\]
$(\mathit{cnt}', \mathit{sum}')$ 是填了 $d$ 之後那個狀態的回傳值。 只加 $d$
而不是 $d \cdot \mathit{cnt}'$，是這類題目最常見的錯。

\subsection{例題：所有號碼的上坡次數總和}
\label{sec:dg-ex-ascent}

\textbf{題意}（\akiref{total-climbs-over-all-ticket-numbers}{所有號碼的上坡次數總和}）

\begin{itemize}
\item 把一個正整數的十進位寫法（沒有前導零）由左往右看，相鄰兩位如果左邊的數
字比右邊的小，就是一次\textbf{上坡}。 例如 $1324$ 有兩次（$1 \to 3$、
$2 \to 4$），$5577$ 有一次，$5$ 沒有。
\item 給 $L$、$R$，求 $L$ 到 $R$ 每個整數的上坡次數加起來是多少。
\item $1 \le L \le R \le 10^{17}$，答案在 \texttt{long long} 的範圍內。
\end{itemize}

例如 $L = 1$、$R = 20$：一位數都沒有上坡；$10$、$11$、$20$ 也沒有；$12$ 到
$19$ 各一次。 答案是 $8$。

\textbf{想法}

暴力解對每個數數一次上坡，$R$ 到 $10^{17}$，來不及。 令 $g(x)$ 為 $1$ 到 $x$
的上坡次數總和，答案是 $g(R) - g(L - 1)$。

判斷這一位會不會造成上坡，要知道前一位的數字，所以狀態是（位置、前一位數字、
是否貼著上界、是否已經開始）。 前導零在這題一定要處理：$5$ 不能看成
\texttt{05}，不然 $0 \to 5$ 會被當成一次上坡。 還沒開始時沒有「前一位」，第一
個非零數字也不會和前面形成上坡。

這題要的是總和，所以每個狀態回傳 $(\mathit{cnt}, \mathit{sum})$。 這一位填
$d$，如果已經開始而且前一位比 $d$ 小，這一次上坡會出現在後面每一種填法裡，要加
$\mathit{cnt}'$。

用不貼上界的狀態檢查一次。 剩一位、前一位是 $p$：後面有 $10$ 種填法，其中
$9 - p$ 種比 $p$ 大，所以是 $(10, 9 - p)$。 剩兩位、前一位是 $1$：這兩位寫成
$ab$，$a$ 有 $8$ 種比 $1$ 大，每一種後面接 $10$ 個 $b$，這一步貢獻
$8 \times 10 = 80$；$a$ 到 $b$ 的上坡，是十個 $(10, 9 - a)$ 的 $\mathit{sum}$
加起來，$9 + 8 + \cdots + 0 = 45$。 所以是 $(100, 125)$，也就是 $100$ 到
$199$ 的上坡次數總和是 $125$。 如果第一步只加 $8$，就變成 $53$，少算了很多。

同樣的模型也出現在別的地方：統計一大批流水號裡總共用了幾個某種數字、全部號碼
的數字和、每個號碼的某種分數加總，都是「每個數一個量，求全部的總和」。

\textbf{程式}

\begin{example}[label=ex:dg-ascent]{cpp}{(ascent.cpp) 所有號碼的上坡次數總和}
#include <iostream>
#include <string>
using namespace std;

struct Info {
    long long cnt;   // 從這裡往後有幾種填法
    long long sum;   // 這些填法在後面產生的上坡次數加起來
};

string s;
// memo[剩下幾位][前一位數字]，只存 !tight 且 started 的狀態
Info memo[20][10];
bool done[20][10];

Info dfs(int pos, int prev, bool tight, bool started) {
    int n = s.size();
    if (pos == n) {
        return {1, 0};             // 填完了：一種填法，沒有上坡
    }
    bool useMemo = !tight && started;
    if (useMemo && done[n - pos][prev]) {
        return memo[n - pos][prev];
    }
    int up = tight ? s[pos] - '0' : 9;
    Info res = {0, 0};
    for (int d = 0; d <= up; d++) {
        bool nt = tight && d == up;
        Info sub;
        if (!started && d == 0) {      // 還是前導零：沒有寫下數字
            sub = dfs(pos + 1, 0, nt, false);
        } else {
            sub = dfs(pos + 1, d, nt, true);
            if (started && prev < d) {
                res.sum += sub.cnt;    // 出現在後面每一種填法裡
            }
        }
        res.cnt += sub.cnt;
        res.sum += sub.sum;
    }
    if (useMemo) {
        done[n - pos][prev] = true;
        memo[n - pos][prev] = res;
    }
    return res;
}

long long climbs(long long x) {   // 1 到 x 所有整數的上坡次數總和
    if (x <= 0) {
        return 0;
    }
    s = to_string(x);
    return dfs(0, 0, true, false).sum;
}

int main() {
    long long L, R;
    cin >> L >> R;
    cout << climbs(R) - climbs(L - 1) << '\n';
    return 0;
}
\end{example}

\iosample{
1 20 \\
}{
8 \\
}

\begin{itemize}
\item 回傳兩個數，用一個小 \texttt{struct}。 \texttt{\{1, 0\}} 這種寫法直接依
成員的順序設初值。 也可以用 \texttt{pair<long long, long long>}，只是
\texttt{.first}、\texttt{.second} 比較看不出意思。
\item \texttt{memo} 的初值不容易用一個「不可能的值」表示，所以另外開一個
\texttt{done} 記哪些格子算過。 全域的 \texttt{bool} 陣列一開始都是
\texttt{false}。
\item \texttt{memo} 和題目的其他數無關（不像前一題依賴 $k$），算 $g(R)$ 存下
的值，算 $g(L - 1)$ 時直接用。
\item $x = 0$ 時 \texttt{climbs} 直接回傳 $0$，所以 $L = 1$ 不用另外判斷。 整數
$0$ 的上坡次數本來就是 $0$，算不算進去都一樣，這題不用特別處理它。
\end{itemize}

\textbf{複雜度}

不貼上界的狀態最多 $18 \times 10$ 個，每個試 $10$ 個數字，算兩次。
$g(10^{17})$ 大約是 $7.1 \times 10^{17}$，在 \texttt{long long} 的範圍內。 最大的測試資料在本書測試的環境下不到 $0.01$ 秒。

\subsection{常見錯誤}

\begin{itemize}
\item \textbf{沒有 $\mathit{started}$}：把每個數補成一樣長，第一位以後一律和前
一位比。 範例的輸入印出 \texttt{17}：$1$ 到 $9$ 被看成 \texttt{01} 到
\texttt{09}，多了九次 $0$ 往上的上坡。
\item \textbf{上坡只加一}：把 \texttt{res.sum += sub.cnt;} 寫成
\texttt{res.sum += 1;}。 範例的輸入剛好是 $8$（兩位數時後面只剩一種填法），
輸入 \texttt{1 1000} 的正確答案是 $801$，程式印出 \texttt{477}。
\item \textbf{區間少一個端點}：寫成 \texttt{climbs(R) - climbs(L)}，$L$ 本身的
上坡被扣掉了。 輸入 \texttt{12 12} 應該是 $1$，這樣寫印出 \texttt{0}。
\end{itemize}

\section{不含某一串數字：小型自動機}
\label{sec:dg-automaton}

回到章首的會員卡，把「不含 $4$」換成「不含某一串數字 $P$」，例如 $P =$
\texttt{13} 或 \texttt{112}：號碼的十進位寫法裡，$P$ 不能以連續的一段出現。
$|P| = 1$ 就是章首的問題。

一位一位填的時候，要記住前面的什麼？ 直覺的答案是「目前已經連續對上 $P$ 的幾
個字」，記成 $j$：下一個數字等於 $P$ 的第 $j + 1$ 個字就 $j$ 加一，到 $|P|$
表示 $P$ 出現了；對不上就歸零（如果剛好等於 $P$ 的第一個字，就記成 $1$）。

這個做法會漏。 $P =$ \texttt{112}，號碼 \texttt{1112}：讀完 \texttt{11}，
$j = 2$；第三個字 \texttt{1} 對不上 \texttt{2}，歸零後因為它等於 \texttt{1}，記
成 $j = 1$；第四個字 \texttt{2} 對不上 $P$ 的第二個字 \texttt{1}，又歸零。 結
論是「不含」，可是 \texttt{1112} 的後三個字就是 \texttt{112}。 問題出在第三個字
對不上時：剛讀過的 \texttt{11} 本身又是 \texttt{112} 的開頭，應該回到 $j = 2$，
不是 $1$。

\subsection{狀態該記什麼}

正確的狀態是：\textbf{目前讀到的字串，最長的一段結尾，同時是 $P$ 的開頭，長度
是多少}。 記成 $j$，$0 \le j \le |P|$；$j = |P|$ 表示 $P$ 已經出現。 剛才的例
子：\texttt{1}、\texttt{11}、\texttt{111}、\texttt{1112} 的 $j$ 依序是
$1, 2, 2, 3$，在最後一個字發現 \texttt{112}。

狀態從 $j$ 讀到數字 $c$ 會變成多少，記成 $\mathrm{nxt}(j, c)$。 它只和 $P$ 的
前 $j$ 個字、以及 $c$ 有關，和更前面讀過什麼無關（理由見下面的框）。 所以可以
事先把 $|P| \times 10$ 個值算好：令 $t$ 是「$P$ 的前 $j$ 個字接上 $c$」，從長
到短試 $t$ 的每一段結尾，第一個等於 $P$ 開頭的長度就是答案。 $P$ 很短時，這樣
直接比對就夠快，例如 $|P| \le 5$ 時全部只要幾百次字串比較。 $P =$ \texttt{112}
的表：

\begin{center}
\begin{tabular}{cccc}
\toprule
$j$（已對上） & 讀到 \texttt{1} & 讀到 \texttt{2} & 讀到其他數字 \\
\midrule
$0$（空） & $1$ & $0$ & $0$ \\
$1$（\texttt{1}） & $2$ & $0$ & $0$ \\
$2$（\texttt{11}） & $2$ & $3$（出現 $P$） & $0$ \\
\bottomrule
\end{tabular}
\end{center}

$j = 2$ 讀到 \texttt{1}，$t = $ \texttt{111}，長度 $3$ 的結尾 \texttt{111} 不
是 \texttt{112}，長度 $2$ 的結尾 \texttt{11} 是 \texttt{112} 的開頭，所以是
$2$。 這一格就是剛才那個錯誤做法填成 $1$ 的地方。

這樣一張表，加上「從哪個狀態開始、到哪個狀態算失敗」，叫做\textbf{自動機}
（automaton）：一個有限的狀態集合，每讀一個字元，照表換到下一個狀態。 數位 DP
的狀態只要多記一個 $j$，轉移時查 $\mathrm{nxt}$ 表，走到 $|P|$ 的就不要了。

\begin{whybox}[為什麼只看 $P$ 的前 $j$ 個字接上 $c$ 就夠]
設讀過的字串是 $u$，它最長的「同時是 $P$ 開頭的結尾」長度是 $j$。 現在讀到 $c$，
要找 $uc$ 最長的同樣性質的結尾，設長度是 $\ell \ge 1$。 這一段去掉最後的 $c$，
是 $u$ 的一段結尾，長度 $\ell - 1$，而且是 $P$ 的開頭。 照 $j$ 的定義，
$\ell - 1 \le j$，所以這一段整個落在「$u$ 的最後 $j$ 個字接上 $c$」裡面。
$u$ 的最後 $j$ 個字就是 $P$ 的前 $j$ 個字，因此只要在「$P$ 的前 $j$ 個字接上
$c$」裡找，不會漏掉。
\end{whybox}

這張表有更快的建法，叫做 KMP 的失敗函數，可以在 $O(|P|)$ 的時間內算出每個 $j$
退回到哪裡，$P$ 很長時才需要；《C++ 演算法解題 3》講字串時會完整介紹。 這一章
的題目 $P$ 都很短，直接比對就好，重點是想清楚狀態的意思。

同一個想法也可以處理「必須含有 $P$」：含有 $P$ 的個數，等於全部的個數減掉不含
$P$ 的個數；或者讓走到 $|P|$ 的狀態留在 $|P|$，最後只數停在 $|P|$ 的。

\section{第 \texorpdfstring{$K$}{K} 個數：沿著位數往下走}
\label{sec:dg-kth}

章首的第二個問題：第 $2526$ 位會員拿到幾號？ 可用的號碼由小到大是
$1, 2, 3, 5, \ldots$，要找第 $K$ 個。

一個方法是第 1 本第~3~章的對答案二分搜：「$1$ 到 $x$ 有幾個可用號碼」隨 $x$
不會變小，找最小的 $x$ 讓個數 $\ge K$。 每次判斷做一次第~\ref{sec:dg-prefix}~
節的計數，大約 $60$ 次就找到了。 只問一次時這樣很好；題目一次問 $10^5$ 個 $K$
時，就多了 $60$ 倍的工作。

更直接的方法是把第~\ref{sec:dg-prefix}~節的分塊倒過來用。 同樣位數的號碼，數
值的大小和字串的大小順序相同，所以：

\begin{enumerate}
\item \textbf{先決定位數。} 一位數的可用號碼有 $8$ 個（$1$ 到 $9$，去掉 $4$），
兩位數有 $8 \times 9 = 72$ 個，三位數 $8 \times 81 = 648$ 個，四位數
$8 \times 729 = 5832$ 個。 $K = 2526$ 先扣掉前三種，剩
$2526 - 8 - 72 - 648 = 1798$，答案是第 $1798$ 個四位數。
\item \textbf{再從最高位往下，每一位由小到大試。} 第一位填 $1$ 的四位數有
$9^3 = 729$ 個可用的。 $1798 > 729$，表示答案不在這一塊，跳過整塊，剩
$1069$；第一位填 $2$ 也是 $729$ 個，跳過，剩 $340$；第一位填 $3$ 有 $729$ 個，
$340 \le 729$，答案在這一塊裡，第一位是 $3$。 接著在這一塊裡用同樣的方法決定
第二位。
\end{enumerate}

\begin{center}
\begin{tabular}{cllc}
\toprule
決定第幾位 & 依序試的數字和那一塊的個數 & 停在 & 剩下的 $K$ \\
\midrule
第 1 位 & $1$：$729$、$2$：$729$、$3$：$729$ & $3$ & $1798 \to 340$ \\
第 2 位 & $0$ 到 $3$ 各 $81$、$4$ 跳過、$5$：$81$ & $5$ & $340 \to 16$ \\
第 3 位 & $0$：$9$、$1$：$9$ & $1$ & $16 \to 7$ \\
第 4 位 & $0$ 到 $3$ 各 $1$、$4$ 跳過、$5$、$6$ 各 $1$、$7$：$1$ & $7$ & $7 \to 1$ \\
\bottomrule
\end{tabular}
\end{center}

答案是 $3517$，和第~\ref{sec:dg-prefix}~節「$3517$ 以內有 $2526$ 張」對得起
來。 每一步需要的都是「從某個狀態出發、後面還有 $r$ 位自由，有幾種填法」，正
好是第~\ref{sec:dg-table}~小節那種和 $N$ 無關的表。 表建一次，每個 $K$ 只要
$O(10 \times \text{位數})$。

決定第一位時從 $1$ 開始試，不從 $0$ 開始：位數已經固定，第一位是 $0$ 的其實是
比較短的數，已經在第一步扣掉了。

\subsection{例題：跳過不吉利字串的第 \texorpdfstring{$K$}{K} 個號碼}
\label{sec:dg-ex-skip}

\textbf{題意}（\akiref{the-k-th-ticket-without-the-unlucky-string}{跳過不吉利字串的第 K 個號碼}）

\begin{itemize}
\item 飲料店的號碼牌不使用十進位寫法（沒有前導零）裡出現字串 $P$ 的號碼，$P$ 以
連續的一段出現就不行。 $P$ 由 $1$ 到 $5$ 個數字組成，可能以 \texttt{0} 開頭。
\item 可用的號碼由小到大發給客人。 有 $Q$ 筆詢問，每筆給 $K$，回答第 $K$ 小的
可用正整數。
\item $Q \le 10^5$，$K \le 10^{15}$，答案一定小於 $10^{17}$。
\end{itemize}

例如 $P =$ \texttt{13}：$1$ 到 $12$ 都可以用，$13$ 跳過，第 $13$ 位拿到 $14$；
$1$ 到 $101$ 之中只跳過了 $13$，所以第 $100$ 位拿到 $101$。 $P =$ \texttt{05}
時，$5$ 的寫法是 \texttt{5}，不是 \texttt{05}，可以用；$105$ 不能用。

\textbf{想法}

暴力解從 $1$ 往上數，每個號碼檢查一次。 $K$ 到 $10^{15}$，數不完。 前兩節的東
西合起來就是這題：狀態是第~\ref{sec:dg-automaton}~節的 $j$，找第 $K$ 個用這一
節的走法。

\begin{itemize}
\item 先建 $\mathrm{nxt}(j, c)$，$0 \le j < |P|$。
\item 再建 $\mathrm{ways}(r, j)$ ＝ 目前對上 $j$ 個字，後面再接 $r$ 個數字（每個
$0$ 到 $9$），一路都沒有讓 $P$ 出現的接法數。 $\mathrm{ways}(0, j) = 1$；
$\mathrm{ways}(r, j)$ 是所有 $\mathrm{nxt}(j, c) < |P|$ 的
$\mathrm{ways}(r - 1, \mathrm{nxt}(j, c))$ 加起來。
\item 每筆詢問：位數 $\ell$ 的可用號碼有 $\sum_{d=1}^{9}
\mathrm{ways}(\ell - 1, \mathrm{nxt}(0, d))$ 個（只加 $\mathrm{nxt}(0, d) <
|P|$ 的），從 $\ell = 1$ 開始扣，決定位數；再一位一位試，填 $d$ 那一塊的大小是
$\mathrm{ways}(\text{剩下的位數}, \mathrm{nxt}(j, d))$。
\end{itemize}

前導零在這題不會造成麻煩：位數先固定，第一位從 $1$ 開始，自動機只讀到真的寫出
來的數字。 如果把號碼補成固定長度再讓自動機讀，$P =$ \texttt{05} 時 $5$ 會被
看成 \texttt{\ldots 05} 而跳過，就錯了。

答案小於 $10^{17}$，最多 $17$ 位，所以 $\mathrm{ways}$ 只要到 $r = 16$，最大
不超過 $10^{16}$。

如果改用二分搜，答案小於 $10^{17}$，每筆詢問要做大約 $57$ 次計數，每次沿著 $N$
走大約 $17 \times 10$ 步，$Q = 10^5$ 時大約 $10^9$ 步，比直接走多了幾十倍。

同樣的模型也出現在別的地方：電話號碼、車牌、票券流水號跳過特定字串，以及把一
組有限制的編號「依序排好」之後找第幾個。

\textbf{程式}

\begin{example}[label=ex:dg-skip]{cpp}{(skip.cpp) 跳過不吉利字串的第 K 個號碼}
#include <iostream>
#include <string>
using namespace std;

int main() {
    ios::sync_with_stdio(false);
    cin.tie(nullptr);
    string P;
    int q;
    cin >> P >> q;
    int m = P.size();

    // nxt[j][c]：已經對上 P 的前 j 個字，再讀到數字 c 之後
    // 對上幾個。 做法：t = P 的前 j 個字接上 c，
    // 找 t 最長的、同時是 P 開頭的結尾。
    int nxt[6][10];
    for (int j = 0; j < m; j++) {
        for (int c = 0; c <= 9; c++) {
            string t = P.substr(0, j) + char('0' + c);
            int len = t.size();
            while (len > 0 &&
                   t.substr(t.size() - len) != P.substr(0, len)) {
                len--;
            }
            nxt[j][c] = len;           // len == m 表示 P 出現了
        }
    }

    // ways[r][j]：目前對上 j 個字，後面再接 r 個數字
    // （0 到 9 都可以），一路都沒有讓 P 出現的接法數
    long long ways[17][6];
    for (int j = 0; j < m; j++) {
        ways[0][j] = 1;
    }
    for (int r = 1; r <= 16; r++) {
        for (int j = 0; j < m; j++) {
            ways[r][j] = 0;
            for (int c = 0; c <= 9; c++) {
                if (nxt[j][c] < m) {
                    ways[r][j] += ways[r - 1][nxt[j][c]];
                }
            }
        }
    }

    while (q--) {
        long long K;
        cin >> K;
        // 第一步：決定位數。 len 位的可用號碼，第一位是 1 到 9。
        int len = 1;
        while (true) {
            long long block = 0;
            for (int d = 1; d <= 9; d++) {
                if (nxt[0][d] < m) {
                    block += ways[len - 1][nxt[0][d]];
                }
            }
            if (K <= block) {
                break;
            }
            K -= block;                // 整個 len 位都跳過
            len++;
        }
        // 第二步：從最高位開始，每一位由小到大試
        string ans;
        int j = 0;
        for (int pos = 0; pos < len; pos++) {
            for (int d = (pos == 0 ? 1 : 0); d <= 9; d++) {
                if (nxt[j][d] == m) {
                    continue;          // 填 d 就出現 P，整塊不能用
                }
                long long block = ways[len - 1 - pos][nxt[j][d]];
                if (K <= block) {      // 答案在這一塊裡
                    ans += char('0' + d);
                    j = nxt[j][d];
                    break;
                }
                K -= block;            // 跳過這一塊
            }
        }
        cout << ans << '\n';
    }
    return 0;
}
\end{example}

\iosample{
112 \\
3 \\
111 \\
112 \\
1111 \\
}{
111 \\
113 \\
1113 \\
}

範例的 $P$ 是 \texttt{112}。 $1$ 到 $111$ 都可以用；$112$ 跳過，所以第 $112$
個是 $113$；$1$ 到 $1113$ 之中跳過了 $112$ 和 $1112$，第 $1111$ 個是 $1113$。
\texttt{1112} 正是第~\ref{sec:dg-automaton}~節那個例子，表裡 $\mathrm{nxt}(2,
\texttt{1}) = 2$ 這一格讓程式跳過了它。

\begin{itemize}
\item \texttt{t.substr(t.size() - len)} 是 $t$ 的最後 \texttt{len} 個字，
\texttt{P.substr(0, len)} 是 $P$ 的前 \texttt{len} 個字。 \texttt{len} 從
$t$ 的長度往下試，第一個相等的就是最長的。
\item 這題輸出很多行，用 \texttt{'\textbackslash n'} 而不是 \texttt{endl}，
並關掉和 C 的同步（第 1 本第~1~章），不然 $10^5$ 行的輸出會慢很多。
\item 決定位數的迴圈裡，\texttt{block} 是「剛好 \texttt{len} 位的可用號碼」。
題目保證答案小於 $10^{17}$，所以 \texttt{len} 不會超過 $17$，
\texttt{ways[len - 1]} 不會越界。
\end{itemize}

\textbf{複雜度}

建 $\mathrm{nxt}$ 是 $|P| \times 10$ 格，每格最多 $|P| + 1$ 次長度 $\le |P| + 1$
的字串比較；建 $\mathrm{ways}$ 是 $16 \times |P| \times 10$。 每筆詢問決定位數
最多 $17 \times 9$ 步，逐位最多 $17 \times 10$ 步，$Q = 10^5$ 時總共大約
$3 \times 10^7$ 步。 最大的測試資料在本書測試的環境下約 $0.02$ 秒。

\subsection{常見錯誤}

\begin{itemize}
\item \textbf{對不上就歸零}：$\mathrm{nxt}$ 寫成「等於 $P$ 的下一個字就加一，否
則看是不是 $P$ 的第一個字」。 範例的輸入第三行會印出 \texttt{1112}，它含有
\texttt{112}。 $P$ 只有一個字、或 $P$ 的開頭和結尾沒有重疊（例如
\texttt{13}）時這樣寫不會錯，所以只用 \texttt{13} 測試會以為是對的。
\item \textbf{第一位從 $0$ 開始試}：決定位數之後，\texttt{d} 一律從 $0$ 開始。
第一位填 $0$ 的那一塊其實是比較短的號碼，已經扣過一次，又被算了一次。 範例的
輸入印出 \texttt{011}、\texttt{012}、\texttt{0113}。
\item \textbf{用 \texttt{int}}：$K$ 到 $10^{15}$，\texttt{ways} 到 $10^{16}$，都
要 \texttt{long long}。
\end{itemize}

\section{兩個數一起填}
\label{sec:dg-pairs}

有些題目數的是數對：$0 \le x \le A$、$0 \le y \le B$，$x$ 和 $y$ 對齊之後每一
位要滿足某個關係，例如「直式加法每一位都不進位」：每一位上 $x$ 的數字加 $y$ 的
數字都不超過 $9$。

做法是把 $x$、$y$ 都補到同樣的長度（$A$、$B$ 之中比較長的那個），每一位同時決
定兩個數字 $(a, b)$，一共最多 $100$ 種組合，只留下滿足關係的。 上界有兩個，
\textbf{貼著上界的旗標也要兩個}：$x$ 是否還貼著 $A$、$y$ 是否還貼著 $B$，各自
照第~\ref{sec:dg-skeleton}~節的規則更新，合起來 $4$ 種狀態。

只用一個旗標會錯。 假設只記「兩個數是不是都還貼著」，其中一個先變小就當作兩個
都自由了。 $A = B = 25$ 時，十位填 $(1, 2)$：$x$ 已經比 $25$ 小，$y$ 還貼著
$25$。 只有一個旗標的話，個位就讓 $y$ 填到 $9$，$x = 10$、$y = 29$ 也被算進
去（$1 + 2$、$0 + 9$ 都不進位），可是 $29 > B$。

前導零要不要處理，看補上去的 $0$ 會不會影響關係。 不進位的條件裡，兩邊都補
$0$ 時是 $0 + 0 \le 9$，一定成立；一邊補 $0$ 時也只是 $0$ 加上另一邊的數字，
和真的直式加法一樣。 所以補 $0$ 剛好就是直式加法的對齊方式，不需要
$\mathit{started}$。 換成別的關係時，要像第~\ref{sec:dg-started}~節那樣先問
一次：補上去的 $0$ 會不會被當成真的數字？

還有一類題目，條件是從右往左傳的，例如加法的進位：個位進位會影響十位，十位又影
響百位。 這和「從最高位往右填、用 $\mathit{tight}$ 比大小」的方向相反，直接套
用骨架寫不出轉移。 遇到這種題目，先想清楚兩個方向要怎麼接起來，再決定狀態。

\section{溢位與很長的數}
\label{sec:dg-overflow}

數位 DP 的題目，上界常常就是 \texttt{long long} 附近，數字的大小要先估好：

\begin{center}
\begin{tabular}{ll}
\toprule
量 & 大小 \\
\midrule
\texttt{long long} 的最大值 & 大約 $9.2 \times 10^{18}$ \\
\texttt{unsigned long long} 的最大值 & 大約 $1.8 \times 10^{19}$ \\
$10^{18}$ 的位數 & $19$（陣列開 $20$） \\
$0$ 到 $10^{18}$ 符合條件的「個數」 & 最多 $10^{18} + 1$，\texttt{long long} 裝得下 \\
$1$ 到 $10^{18}$ 的數字和加總 & 大約 $8.1 \times 10^{19}$，兩種都裝不下 \\
\bottomrule
\end{tabular}
\end{center}

\begin{itemize}
\item \textbf{個數}：只要上界不超過 $10^{18}$，個數就不會超過 $10^{18} + 1$，用
\texttt{long long} 就好。 要小心的是中間的表：像第~\ref{sec:dg-table}~小節，
多算一列 $r = 19$，那一列的值就可能到 $10^{19}$。
\item \textbf{總和}：每個數的量乘上 $10^{18}$ 個數，很容易超過。 先估最大值：
$1$ 到 $10^{18}$ 的數字和，平均每個數大約 $81$，總和就超過了兩種型別。 這時題目通常會要求輸出除以某個數的餘數，或者要用
\texttt{\_\_int128}（g++ 提供的 $128$ 位元整數，大約到 $1.7 \times 10^{38}$）。
\texttt{\_\_int128} 不能直接用 \texttt{cin}、\texttt{cout}，輸出要自己一位一位
印：
\begin{code}{cpp}{印出非負的 \_\_int128}
void print128(__int128 x) {     // x >= 0
    if (x >= 10) {
        print128(x / 10);
    }
    cout << char('0' + x % 10);
}
\end{code}
\item \textbf{無號整數和 $L - 1$}：用 \texttt{unsigned long long} 存 $L$，
$L = 0$ 時 $L - 1$ 會繞回 $2^{64} - 1$，見第
~\ref{sec:dg-ex-zeros}~小節的常見錯誤。 先判斷 $L = 0$，或者全部用有號的
\texttt{long long}。
\end{itemize}

上界大到任何整數型別都存不下時，例如幾萬位的數字，就直接把它當成字串，DP 照
著字串一位一位走，程式本來就只用到 \texttt{s[pos]}。 要注意三件事：

\begin{itemize}
\item 位數一多，遞迴的深度就是位數，記憶化的表也要開到位數那麼大。 這時改用
第~\ref{sec:dg-table}~小節的填表寫法比較穩：先建「剩 $r$ 位自由」的表，再沿著
字串走一次。
\item $L - 1$ 要在字串上做減法：從最後一位往前，遇到 $0$ 就變成 $9$ 繼續借，
遇到第一個非零的減一。 $1000$ 減一是 $999$，最前面變成 $0$ 時要拿掉。
\begin{code}{cpp}{字串減一}
// s 是正整數的十進位寫法，沒有前導零
string minusOne(string s) {
    int i = s.size() - 1;
    while (s[i] == '0') {        // 借位：結尾的 0 都變成 9
        s[i] = '9';
        i--;
    }
    s[i]--;
    if (s.size() > 1 && s[0] == '0') {
        s.erase(0, 1);           // 1000 減 1 是 999，少一位
    }
    return s;
}
\end{code}
也可以不做減法，改算 $f(R) - f(L)$，再單獨檢查 $L$ 本身符不符合，符合就加一。
\item 答案要取餘數時，$f(R) - f(L - 1)$ 可能是負的（兩個餘數相減），要先加上模
數再取一次餘數。
\end{itemize}

\section{什麼時候用數位 DP}
\label{sec:dg-when}

\begin{itemize}
\item 題目問「$L$ 到 $R$ 之中有幾個數符合某個條件」，或是「所有數的某個量加起來」，
而 $R$ 到 $10^{18}$，甚至是一個很長的字串。 這麼大的範圍不可能一個一個看。
\item 條件可以\textbf{一位一位檢查}：由左往右讀數字時，只要記住一點點東西（餘
數、前一位、次數、對上 $P$ 的長度），讀完就能判斷。 先把這「一點點東西」用一句
話寫出來，就是狀態。
\item 問「第 $K$ 個符合條件的數」：建一張和上界無關的表，沿著位數往下走；只問
一次的話，二分搜加計數也可以。
\end{itemize}

看起來像、其實不適合的情況：

\begin{itemize}
\item 條件要看整個數的「值」而不是數字，例如「$x$ 是質數」「$x$ 是完全平方數」。
這種條件一位一位讀的時候記不住，數位 DP 用不上。
\item 範圍不大（例如 $R - L \le 10^6$），每個數直接檢查就好，不用冒險。 這個暴
力解也正好拿來對拍。
\item 條件和「每一位」有關，但要記的東西會隨位數暴增，例如「所有數字都不相同」
要記用過哪些數字，是 $2^{10}$ 種集合。 這還做得到，狀態就是第~4~章那種「一個小
集合」；集合再大就不行了。
\end{itemize}

\begin{apcsnote}
APCS 實作題由程式自動評分，只比對輸出。 數位 DP 的錯誤大多在邊界：前導零、整
數 $0$、$L = 0$、$N$ 本身。 範例的輸入常常剛好避開這些情況，本章好幾個錯誤寫
法都能通過範例。 寫完之後，用「一個一個檢查」的暴力解，在 $R \le 2000$ 的所有
$(L, R)$ 上比對，比只測範例可靠得多。
\end{apcsnote}

\section{本章常見錯誤}
\label{sec:dg-pitfalls}

\begin{center}
\begin{tabular}{p{0.2\linewidth}p{0.46\linewidth}p{0.22\linewidth}}
\toprule
錯誤 & 例子 & 見 \\
\midrule
前導零當成數字 & 數 $0$ 時把 $7$ 看成 \texttt{007}；上坡時把 $5$ 看成 \texttt{05}；
補成固定長度再讓自動機讀 &
第~\ref{sec:dg-ex-zeros}、\ref{sec:dg-ex-ascent}、\ref{sec:dg-ex-skip}~節 \\
整數 $0$ 和 $L = 0$ & 整數 $0$ 寫作 \texttt{0}；$f(-1)$ 沒有判斷；無號整數的
$L - 1$ & 第~\ref{sec:dg-ex-zeros}~節 \\
記錯狀態 & 貼著上界或還沒開始的狀態也記憶化；記憶化用 \texttt{pos} 當索引，位數
變了沒清空 & 第~\ref{sec:dg-memo}、\ref{sec:dg-ex-zeros}~節 \\
狀態少記一樣 & 「連續對上幾個」對不上就歸零；兩個上界共用一個旗標 &
第~\ref{sec:dg-automaton}、\ref{sec:dg-pairs}~節 \\
總和只加一 & 新產生的量沒有乘上後面的填法數 & 第~\ref{sec:dg-ex-ascent}~節 \\
區間端點 & $f(R) - f(L)$ 漏了 $L$；只算比 $N$ 小的，忘了 $N$ 本身 &
第~\ref{sec:dg-prefix}、\ref{sec:dg-ex-ascent}~節 \\
溢位 & 表多算一列到 $10^{19}$；總和超過 \texttt{long long}；取餘數相減變負數 &
第~\ref{sec:dg-table}、\ref{sec:dg-overflow}~節 \\
\bottomrule
\end{tabular}
\end{center}

\section{本章練習}
\label{sec:dg-exercises}

\begin{drills}[label=drills:dg]
\item 不寫程式，照第~\ref{sec:dg-prefix}~節的分塊，算出 $1$ 到 $2025$ 之中不含
$4$ 的數有幾個。 列出每一位加了多少。

\item 範例~\ref{ex:dg-digitsum} 輸入 \texttt{1 347 5} 時：
\begin{enumerate}
\item 算 $f(347)$ 的過程中，\texttt{dfs} 被呼叫時 \texttt{tight} 是真的有幾
次？
\item 接著算 $f(0)$ 時，有沒有讀到 $f(347)$ 存下來的 \texttt{memo}？ 為什麼？
\end{enumerate}

\item 有人把範例~\ref{ex:dg-zeros} 的記憶化改成不看 \texttt{tight}：
\begin{code}{cpp}{練習 3}
bool useMemo = started;
\end{code}
輸入 \texttt{1 20 0} 會印出什麼？ 用圖~\ref{fig:dg-tight-path} 的方式說明錯
在哪裡。

\item 下面每個條件，寫出數位 DP 除了位置、$\mathit{tight}$ 之外要記的狀態，並
說明要不要 $\mathit{started}$。
\begin{enumerate}
\item 相鄰兩位數字的差都不超過 $2$，例如 $1312$、$5$、$886$。
\item 至少有一組相鄰的兩位數字相同，例如 $1223$、$100$。
\item 數字 $7$ 出現的次數是奇數。
\item 整個數是 $7$ 的倍數，而且數字裡沒有 $7$。
\end{enumerate}

\item $P =$ \texttt{1213}。
\begin{enumerate}
\item 寫出 $\mathrm{nxt}(3, c)$，也就是已經對上 \texttt{121} 時，讀到每個數字之
後的狀態。
\item 號碼 \texttt{121213} 含不含 $P$？ 寫出正確的狀態序列，以及「對不上就歸
零」的寫法會得到什麼結論。
\end{enumerate}

\item 章首的會員卡（跳過含 $4$ 的號碼），照第~\ref{sec:dg-kth}~節的方法，手算第
$100$ 位會員拿到的號碼。

\item 計算 $1$ 到 $N$ 的數字和加總，$N \le 10^{18}$。
\begin{enumerate}
\item 每個狀態要回傳哪兩個數？ 這一位填 $d$ 時，兩個數各怎麼更新？
\item 答案可能超過 \texttt{long long} 嗎？ 如果會，有哪些處理方法？
\end{enumerate}
\end{drills}

\section{實作練習}
\label{sec:dg-practice}

下面是 Aki-Judge 上和本章相關的題目，依小節分組。 難度分四階：暖身是直接套用
本章的骨架；基本要自己想出多記什麼狀態；實戰的題敘比較長，要先把題目轉換成本章
的形狀；挑戰會和其他技巧合在一起，或需要更進一步的觀察。

\begin{center}
\begin{tabular}{llc}
\toprule
小節 & 題目 & 難度 \\
\midrule
骨架與前導零 & \akiref{alternating-parity-ticket-numbers}{奇偶交錯的號碼} & 暖身 \\
 & \akiref{integers-with-exactly-k-zeros}{恰好 k 個零的整數} & 暖身 \\
\midrule
多帶一點資訊 & \akiref{total-climbs-over-all-ticket-numbers}{所有號碼的上坡次數總和} & 基本 \\
 & \akiref{the-code-hidden-in-the-numbers}{號碼裡藏著的暗號} & 基本 \\
 & \akiref{run-limits-on-very-long-serial-numbers}{超長流水號的連號限制} & 基本 \\
 & \akiref{multiples-both-ways}{正讀倒讀都是 K 的倍數} & 實戰 \\
 & \akiref{how-far-the-digit-stickers-go}{數字貼紙夠貼到幾號} & 實戰 \\
\midrule
兩個數一起填 & \akiref{subtraction-without-borrowing}{直式減法不借位的數對} & 實戰 \\
 & \akiref{carries-when-adding-a-fixed-number}{加上固定數時的進位次數} & 挑戰 \\
\midrule
自動機與第 $K$ 個 & \akiref{the-k-th-ticket-without-the-unlucky-string}{跳過不吉利字串的第 K 個號碼} & 挑戰 \\
\bottomrule
\end{tabular}
\end{center}

本章的三題例題也在 Aki-Judge 上，讀完各節之後可以先自己寫一次再對照。 這一章
的題單還會再增加。

\section{讀完你應該會}
\label{sec:dg-checklist}

\begin{itemize}
\item 把「$L$ 到 $R$ 之中有幾個」寫成 $f(R) - f(L - 1)$，並先處理 $L = 0$。
\item 寫出數位 DP 的骨架：從最高位往下填，用 $\mathit{tight}$ 限制這一位的上限，
只記憶化不貼上界的狀態，並說得出為什麼。
\item 看到條件會被前導零影響時，加上 $\mathit{started}$，並單獨判斷整數 $0$。
\item 替一個條件設計要多記的狀態（餘數、前一位、次數、對上 $P$ 的長度），求總
和時讓每個狀態回傳（個數，總和）兩個數。
\item 建一張和上界無關的「剩 $r$ 位自由」的表，用它沿著位數往下走，找出第 $K$
個符合條件的數。
\end{itemize}

\textbf{延伸閱讀}（書末會整理完整的參考資料）

\begin{itemize}
\item Steven Halim 等，\textit{Competitive Programming 4}，Book 2 的 6.3.3 節：
用一題例題示範怎麼一步一步把數位 DP 的狀態補齊，並附一組練習題。
\item 劉汝佳，《算法競賽入門經典》，第 2 版，10.4 節例題 10-22：統計一段區間裡每
個數字各出現幾次，用的也是把 $0$ 到 $n$ 分成幾塊、每塊後面自由的想法。
\end{itemize}

\addcontentsline{toc}{section}{練習解答}
\section*{練習解答}

\begin{enumerate}
\item $1480$ 個。 \texttt{2025} 由左往右：千位 $2$，可以填 $0$、$1$，各
$9^3 = 729$，共 $1458$；百位 $0$，沒有比它小的數字，不加；十位 $2$，可以填
$0$、$1$，各 $9$，共 $18$；個位 $5$，可以填 $0$ 到 $3$（$4$ 跳過），共 $4$；
$2025$ 本身不含 $4$，加 $1$。 $1458 + 18 + 4 + 1 = 1481$ 是 $0$ 到 $2025$ 的
個數，扣掉 $0$ 是 $1480$。

\item
\begin{enumerate}
\item $4$ 次：\texttt{pos} 是 $0$、$1$、$2$、$3$ 時各一次，沿著
\texttt{3}、\texttt{4}、\texttt{7} 走到底。 貼著上界的路只有一條。
\item 沒有。 $f(0)$ 的字串是 \texttt{"0"}，只有一位，\texttt{dfs(0, 0, true)}
的 \texttt{up} 是 $0$，只填 $0$ 而且仍然貼著，直接走到結尾，從頭到尾沒有不貼上
界的狀態，所以不會讀 \texttt{memo}。 即使讀了，存的值也和 $N$ 無關，一樣是對
的。
\end{enumerate}

\item 印出 \texttt{27}，正確答案是 $18$。 $f(20)$ 的第一位試 $0$、$1$、$2$：
填 $1$ 時走到「剩一位、已開始、$0$ 個 $0$、不貼」，算出 $9$（$11$ 到 $19$），
存進 \texttt{memo[1][0]}。 填 $2$ 時走到「剩一位、已開始、$0$ 個 $0$、貼著」，
它的第二位只能填 $0$，答案應該是 $0$（$20$ 有一個 $0$）；可是改過的程式也去讀
\texttt{memo[1][0]}，拿到 $9$。 多算的 $9$ 就是 $21$ 到 $29$。 在圖
~\ref{fig:dg-tight-path} 上，貼著上界的那排節點只能往下走到 $N$ 的數字為止，不
能借用旁邊整塊自由的答案。

\item
\begin{enumerate}
\item 前一位數字（$10$ 種）。 需要 $\mathit{started}$：不處理的話，$5$ 被看成
\texttt{05}，$0$ 和 $5$ 差 $5$，會被誤判成不符合。 開始之前沒有「前一位」，第
一個非零數字不用和任何數字比。
\item 前一位數字，以及「是否已經出現過相鄰相同」（真假）。 需要
$\mathit{started}$：不處理的話，\texttt{005} 的兩個前導零會被當成相鄰相同。
\item 目前 $7$ 出現的次數的奇偶（$2$ 種）。 不需要 $\mathit{started}$，前導零
不是 $7$。
\item 整個數除以 $7$ 的餘數（$7$ 種），轉移時跳過數字 $7$。 不需要
$\mathit{started}$：前導零不是 $7$，而且 $10 \cdot 0 + 0 = 0$，不改變餘數。
\end{enumerate}

\item
\begin{enumerate}
\item $t$ 是 \texttt{121} 接上 $c$。 $c = 3$：\texttt{1213}，就是 $P$，是
$4$。 $c = 2$：\texttt{1212} 的結尾 \texttt{12} 是 $P$ 的開頭，是 $2$。
$c = 1$：\texttt{1211} 的結尾 \texttt{1}，是 $1$。 其他數字都是 $0$。
\item 含有，最後四個字就是 \texttt{1213}。 正確的狀態依序是 $1, 2, 3, 2, 3, 4$，
讀到最後一個字時到達 $4$。 「對不上就歸零」的寫法在第四個字 \texttt{2}（對上
\texttt{121} 之後，期待 \texttt{3}）把狀態變成 $0$，之後是 $1, 0$，結論是不含。
它漏掉的是 \texttt{121} 後面的 \texttt{12} 本身又是 $P$ 的開頭。
\end{enumerate}

\item $121$。 一位數 $8$ 個、兩位數 $72$ 個，扣掉之後是第 $20$ 個三位數。 第一
位填 $1$ 的那一塊有 $81$ 個，$20 \le 81$，第一位是 $1$。 第二位填 $0$ 有 $9$
個，跳過，剩 $11$；填 $1$ 有 $9$ 個，跳過，剩 $2$；填 $2$ 有 $9$ 個，停下。 第
三位填 $0$ 有 $1$ 個，跳過，剩 $1$；填 $1$，停下。 可以用範例~\ref{ex:dg-skip}
輸入 $P =$ \texttt{4}、$K = 100$ 檢查。

\item
\begin{enumerate}
\item 回傳（後面有幾種填法 $\mathit{cnt}$，這些填法在後面的數字和加總
$\mathit{sum}$）。 填 $d$ 之後的狀態回傳 $(\mathit{cnt}', \mathit{sum}')$，則
$\mathit{cnt} \mathrel{+}= \mathit{cnt}'$，$\mathit{sum} \mathrel{+}=
\mathit{sum}' + d \cdot \mathit{cnt}'$。 前導零加的是 $0$，不需要
$\mathit{started}$。
\item 會。 $N = 10^{18}$ 時，$0$ 到 $10^{18} - 1$ 每一位上 $0$ 到 $9$ 各出現
$10^{17}$ 次，$18$ 位共 $18 \times 45 \times 10^{17} = 8.1 \times 10^{19}$，
超過 \texttt{long long} 也超過 \texttt{unsigned long long}。 可以用
\texttt{\_\_int128} 存並自己寫輸出，或者題目要求輸出餘數時，每一步都取餘數（
$d \cdot \mathit{cnt}'$ 先把 $\mathit{cnt}'$ 取餘數再乘）。
\end{enumerate}
\end{enumerate}

\chapter{DAG 與拓樸排序}
\label{ch:dag}

% 本章的圖共用的樣式（工作、強調的工作、箭頭）
\tikzset{
    tk/.style={rounded corners=2pt, draw=black!60, thick, minimum height=0.7cm,
               inner sep=3pt, font=\small},
    tkhl/.style={tk, draw=memStroke, fill=memFill},
    tv/.style={circle, draw=black!60, thick, minimum size=0.7cm, inner sep=0pt, font=\small},
    tvhl/.style={tv, draw=memStroke, fill=memFill},
    tvbad/.style={tv, draw=optStroke, fill=optFill},
    ta/.style={-{Stealth[length=2.2mm]}, thick, draw=black!60, shorten >=1pt},
    tahl/.style={-{Stealth[length=2.2mm]}, very thick, draw=memStroke, shorten >=1pt},
}

週末在家吃火鍋，要做的事有七件，畫在下圖。 箭頭表示先後：買完菜才能洗菜切
菜、調沾醬、煮湯底；湯底要等電磁爐借回來才能煮；最後四件事都好了才開火上
桌。 一次只做一件事，請排出一個做事的順序。 往下讀之前，先在紙上寫一個，並
留意自己是怎麼挑「下一件」的。

\begin{center}
\begin{tikzpicture}[x=2.6cm, y=1.05cm]
\node[tk] (c) at (0,2) {3 買菜};
\node[tk] (f) at (0,0.5) {6 借電磁爐};
\node[tk] (d) at (0,-0.6) {4 擺碗筷};
\node[tk] (b) at (1.15,2.6) {2 洗菜切菜};
\node[tk] (g) at (1.15,1.6) {7 調沾醬};
\node[tk] (e) at (1.15,0.5) {5 煮湯底};
\node[tk] (a) at (2.3,1) {1 開火上桌};
\draw[ta] (c) -- (b);
\draw[ta] (c) -- (g);
\draw[ta] (c) -- (e);
\draw[ta] (f) -- (e);
\draw[ta] (b) -- (a);
\draw[ta] (g) -- (a);
\draw[ta] (e) -- (a);
\draw[ta] (d) -- (a);
\end{tikzpicture}
\diagcaption{吃火鍋前的七件事。 前面的數字是編號，之後的程式都用編號。}
\label{fig:hotpot}
\end{center}

大部分人的做法是：先找「不用等任何事」的（買菜、借電磁爐、擺碗筷），挑一件
做掉；做完之後，看看有沒有哪件事因此可以開始了，再從可以開始的裡面挑一件。
例如 買菜、擺碗筷、借電磁爐、洗菜切菜、調沾醬、煮湯底、開火上桌。 順序不只一
種，但不管怎麼排，每個箭頭都要從前面指到後面。

這種「一堆事情，加上誰要在誰之前」的資料，畫成有向圖，箭頭不能繞成圈，否則
就沒有人能先開始。 《C++ 演算法解題 1：APCS 實作題核心》（以下簡稱第 1 本）第~9~章把這種圖叫做有向無環圖，簡稱 DAG。 那一章說過，
每個動態規劃背後都有一張 DAG，計算時每個點要排在所有指向它的點後面；當時刻
意沒有替這種順序取名字，這一章正式介紹它：它叫\textbf{拓樸順序}。 上面那個
「挑不用等的、做完再看誰可以開始」的做法，第 1 本第~6~章「層層彙整的報告」已經用
過一次，只是那裡的圖是一棵樹。

這一章先把拓樸順序講清楚，接著是兩種排出它的方法：延續第 1 本第~6~章的 Kahn 演算
法，以及用第 1 本第~11~章的深度優先搜尋。 兩種方法都能順便發現圖上有沒有環。 然後
回到第 1 本第~9~章，在輸入給的 DAG 上做動態規劃：最長路、路徑數，以及專案排程裡的
關鍵路徑，並對照資料結構課本裡 AOV 網路、AOE 網路的說法。 最後一節求編號字典序最小的拓樸順序，要把佇列換成優先佇列。

用到的語法：\texttt{queue} 和 \texttt{priority\_queue} 在銜接書第~8~章，自
訂成「小的在上」的寫法也在那裡；鄰接串列的存法在第 1 本第~11~章。 沒讀過銜接書
的讀者，只要會用 \texttt{vector<vector<int>>}、\texttt{queue}，並知道
\texttt{priority\_queue} 可以設定成每次拿出最小的，就能讀這一章。

\section{先後關係與拓樸順序}
\label{sec:topo-order}

把每件事畫成一個點，「$a$ 要在 $b$ 之前」畫成一條有向邊 $a \to b$。 有向圖
上，一個點有幾條邊指進來，叫它的\textbf{入度}（in-degree）；有幾條邊指出去，
叫\textbf{出度}（out-degree）。 圖~\ref{fig:hotpot} 裡，開火上桌（點 $1$）的
入度是 $4$，買菜（點 $3$）的出度是 $3$。 入度是 $0$ 的點不用等任何事，出度是
$0$ 的點沒有任何事在等它。

把所有點排成一列，如果每一條邊 $u \to v$ 都是 $u$ 排在 $v$ 前面，這一列就是
一個\textbf{拓樸順序}（topological order）。 求出一個拓樸順序的過程叫
\textbf{拓樸排序}（topological sort）。 名字裡有「排序」，但它和第 1 本第~2~章的排
序不一樣：這裡沒有大小可以比，只有「誰在誰前面」的規定，沒有規定的兩點誰先
都可以。

用圖~\ref{fig:hotpot} 的編號：

\begin{itemize}
\item $3, 4, 6, 2, 7, 5, 1$ 是拓樸順序。 逐條檢查八條邊，箭尾都排在箭頭前面。
\item $6, 4, 3, 5, 7, 2, 1$ 也是。 同一張圖通常有很多個拓樸順序。
\item $3, 2, 5, 6, 7, 4, 1$ 不是：邊 $6 \to 5$ 的 $6$ 排在 $5$ 後面，湯底還沒有
電磁爐就開始煮了。
\end{itemize}

\subsection{有環就排不出來}

如果圖上有環，例如 $a \to b \to c \to a$，環上三個點誰都不能排第一：排第一的
那個點，一定有環上的另一個點指向它，而那個點排在它後面。 所以有環的圖沒有
拓樸順序。

反過來，沒有環的圖一定排得出來。 理由是：DAG 上一定找得到入度 $0$ 的點。 隨
便挑一個點，如果有邊指向它，就沿著這條邊倒退到前一個點；再有邊指進來就再倒
退。 點只有 $n$ 個，倒退 $n$ 次還沒停，就一定重複經過某個點，也就是找到了一
個環，和沒有環矛盾。 所以一定會停在一個入度 $0$ 的點。 把它排在第一個、從圖
上拿掉，剩下的圖仍然沒有環，再找下一個入度 $0$ 的點，一直做到全部排完。

這段說明本身就是一個演算法，下一節把它寫成程式。 它也說明了一件之後常用的
事：一張有向圖排得出拓樸順序，剛好就是它沒有環的時候。 所以「有沒有環」和
「排不排得出來」是同一個問題，兩種拓樸排序的方法都會順便回答它。

\subsection{回到第 1 本第 9 章}

第 1 本第~9~章說，動態規劃的計算順序要讓「每個點都排在所有指向它的點後面」。 現在
可以直接說：計算順序就是狀態圖的一個拓樸順序。 那一章和第 1 本第~10~章的例子，拓樸
順序都很好找：郵票的箭頭都從小金額指向大金額，由小到大就是一個拓樸順序；格
子圖的箭頭往右、往下，一列一列由左到右也是。 第 1 本第~13~章把樹的 BFS 順序倒過來
算子樹大小，也是在用一個拓樸順序。

這一章的圖直接寫在輸入裡，點的編號和先後沒有關係，看不出哪個順序是對的，要
用演算法算出來。

\section{Kahn 演算法}
\label{sec:kahn}

最直接的做法是照上一節的說明：每一輪把所有點掃一遍，找一個還沒排、入度已經
是 $0$ 的點，排進去，再把它指向的點入度減一。 每一輪掃 $n$ 個點，共 $n$ 輪，
是 $O(n^2)$。 $n = 2 \times 10^5$ 時是 $4 \times 10^{10}$，太慢。

慢的地方在「每一輪都從頭找」。 第 1 本第~6~章的「層層彙整的報告」已經示範過怎麼不
用找：用 \texttt{waiting[i]} 記員工 $i$ 還有幾個部屬沒交，減到 $0$ 就放進佇
列。 那裡的 \texttt{waiting} 就是這一章的入度。 一般的 DAG 和那題只差一點：
員工只要通知一位主管，一般的點可能有好幾條邊指出去，所以要用鄰接串列記下每
個點指向誰。 整理成步驟：

\begin{enumerate}
\item 讀邊時，把 $v$ 加進 $u$ 的串列，並把 $v$ 的入度加一。
\item 把所有入度 $0$ 的點放進佇列。
\item 佇列不空時，從前面拿出一個點 $u$，接在答案後面。 $u$ 的串列裡每個點
$v$ 入度減一；減到 $0$，表示 $v$ 的最後一個前置剛完成，把 $v$ 放進佇列。
\end{enumerate}

這個做法叫 \textbf{Kahn 演算法}（Kahn's algorithm），以提出者的名字命名。 程
式裡多做一件事：結束時檢查排進答案的點夠不夠 $n$ 個，下一小節說明原因。

\begin{example}[label=ex:kahn]{cpp}{(kahn.cpp) Kahn 演算法}
#include <iostream>
#include <queue>
#include <vector>
using namespace std;

int main() {
    int n, m;
    cin >> n >> m;
    vector<vector<int>> adj(n + 1);
    vector<int> indeg(n + 1, 0);           // 入度：還有幾個前置沒完成
    for (int i = 0; i < m; i++) {
        int a, b;
        cin >> a >> b;                     // a 要在 b 之前
        adj[a].push_back(b);
        indeg[b]++;
    }
    queue<int> q;
    for (int v = 1; v <= n; v++) {
        if (indeg[v] == 0) {
            q.push(v);
        }
    }
    vector<int> order;
    while (!q.empty()) {
        int u = q.front();
        q.pop();
        order.push_back(u);
        for (int v : adj[u]) {
            indeg[v]--;
            if (indeg[v] == 0) {           // 最後一個前置剛完成
                q.push(v);
            }
        }
    }
    if ((int)order.size() < n) {           // 有點永遠等不到：圖上有環
        cout << "cycle:";
        for (int v = 1; v <= n; v++) {
            if (indeg[v] > 0) {
                cout << ' ' << v;
            }
        }
        cout << '\n';
        return 0;
    }
    for (int i = 0; i < n; i++) {
        cout << order[i] << (i + 1 < n ? ' ' : '\n');
    }
    return 0;
}
\end{example}

\iosample{
7 8 \\
3 2 \\
3 7 \\
3 5 \\
6 5 \\
2 1 \\
5 1 \\
7 1 \\
4 1 \\
}{
3 4 6 2 7 5 1 \\
}

輸入是圖~\ref{fig:hotpot} 的八條邊。 一開始入度 $0$ 的是 $3$、$4$、$6$，之後
每一步：

\begin{center}
\begin{tabular}{clll}
\toprule
拿出 & 入度減一 & 新放入 & 之後的佇列（前 $\to$ 後） \\
\midrule
 & & & 3 4 6 \\
3 & 2：$1 \to 0$，7：$1 \to 0$，5：$2 \to 1$ & 2、7 & 4 6 2 7 \\
4 & 1：$4 \to 3$ & & 6 2 7 \\
6 & 5：$1 \to 0$ & 5 & 2 7 5 \\
2 & 1：$3 \to 2$ & & 7 5 \\
7 & 1：$2 \to 1$ & & 5 \\
5 & 1：$1 \to 0$ & 1 & 1 \\
1 & & & （空）\\
\bottomrule
\end{tabular}
\end{center}

\begin{itemize}
\item 一個點只在入度\textbf{減到} $0$ 的那一刻放入佇列，所以每個點最多放入一
次。 點 $1$ 的入度被減了四次，前三次都不放。
\item 為什麼拿出來的順序一定合法：$v$ 放入佇列時入度是 $0$，表示每一條指向它
的邊 $u \to v$ 都已經被「減過」了，也就是 $u$ 已經拿出來、排在答案裡。 所以
$v$ 排在每個 $u$ 後面。
\item 複雜度：每個點放入、拿出各一次；每個點拿出時看一遍自己的串列，所有串列
加起來是 $m$ 條邊。 總共 $O(n + m)$。
\end{itemize}

\subsection{找環}
\label{sec:kahn-cycle}

如果圖上有環，環上的點入度永遠不會減到 $0$：環上每個點都在等環上的前一個點，
誰都不會先被拿出來。 所以 Kahn 演算法跑完時，答案裡的點不到 $n$ 個，就表示有
環。 這個檢查只要一行，題目沒有保證無環時一定要做。

例如資料表寫錯，多了一條「煮湯底要在買菜之前」（$5 \to 3$），買菜和煮湯底就
互相等待。 輸入多加一行 \texttt{5 3}（第一行改成 \texttt{7 9}），程式印出
\texttt{cycle: 1 2 3 5 7}。 拿出來的只有 $4$ 和 $6$。 剩下的五個點裡，真正在環上的只有 $3$ 和 $5$；$2$、$7$
要等買菜，$1$ 要等前面幾件事，它們是被環擋住的。 所以 Kahn 剩下的點是「在環
上，或是在某個環的下游」，不一定每個都在環上。

要找出一個真正的環，可以從任一個剩下的點出發往回走：剩下的點入度還不是 $0$，
表示還有一個沒拿出來的點指向它，而那個點也是剩下的點。 一直往回走，走過的點
不會超過 $n$ 個，一定會重複，重複的那一段就是一個環。 往回走需要反向圖（第 1 本第
~11~章）。 第~\ref{sec:dfs-topo}~節的 DFS 可以直接把環印出來，通常比較方便。

\subsection{佇列可以換成別的容器}

Kahn 演算法正確的理由，只用到「放進去的點入度都是 $0$」，和容器裡誰先拿出來
無關。 所以佇列換成堆疊、換成優先佇列，排出來的都是拓樸順序，只是順序不同。
在圖~\ref{fig:hotpot} 上：

\begin{itemize}
\item 佇列（先進先出）：$3, 4, 6, 2, 7, 5, 1$，接近「一層一層」往下。
\item 堆疊（後進先出）：一開始依序放入 $3$、$4$、$6$，先拿出最後放的 $6$，
排出來是 $6, 4, 3, 5, 7, 2, 1$（本章練習第~2~題）。
\item 優先佇列，每次拿編號最小的：$3, 2, 4, 6, 5, 7, 1$。 這是所有拓樸順序裡字
典序最小的一個，第~\ref{sec:lex-topo}~節會說明。
\end{itemize}

題目只要「任一個拓樸順序」時，用哪一種都行。 題目把順序規定死的時候，就要照
題目選容器，下面的例題就是這樣。

\subsection{例題：成果展布置的工作清單}
\label{sec:ex-checklist}

\textbf{題意}（\akiref{setup-checklist-for-the-exhibition}{成果展布置的工作清單}）

\begin{itemize}
\item 有 $n$ 項布置工作和 $m$ 條規定，規定 $a\ b$ 表示工作 $a$ 要在工作 $b$
之前完成。 保證沒有矛盾。
\item 助教用一個待辦佇列決定順序：一開始把沒有前置工作的項目依編號由小到大放
進佇列；每次從最前面拿一項做完；做完後，這一次\textbf{新變成}可以開始的工作
依編號由小到大接到佇列尾端。
\item 輸出完成工作的順序。 $n, m \le 2 \times 10^5$。
\end{itemize}

例如 $n = 4$，規定是 $1 \to 4$、$2 \to 3$。 一開始佇列是 $1, 2$；做完 $1$，
$4$ 可以開始，佇列變成 $2, 4$；做完 $2$，$3$ 可以開始，佇列變成 $4, 3$。 答
案是 \texttt{1 2 4 3}。 $3$ 的編號比 $4$ 小，但 $4$ 比較早進佇列，所以先做。

\textbf{想法}

題目描述的規則就是 Kahn 演算法，只多了兩個要求：一開始的點依編號放入（迴圈
從 $1$ 掃到 $n$ 本來就是這樣），以及同一次新解鎖的點要先排序再放入。 第二個
要求的做法是：拿出 $u$ 之後，先把入度減到 $0$ 的點收集在一個暫存的
\texttt{vector} 裡，排序後再依序放進佇列。

最直接的寫法是照字面模擬：每做完一項，就把所有工作掃一遍，找出新變成可以開
始的。 這是本節開頭的 $O(n^2)$ 做法。 測試資料裡有 $2 \times 10^5$ 項工作排成
一長串的情況，在本書測試的環境下跑了一分鐘還沒結束。

排序會不會太慢？ 每次排序的是「這一次新解鎖的點」，每個點一生只會被解鎖一次，
所以所有批次加起來最多 $n$ 個點，排序的總時間不超過 $O(n \log n)$。

同樣的模型也出現在別的地方：課程有先修規定，每學期只修已經可以修的課；工廠
的生產線上，每道工序要等前面的工序完成。 規定「已經在排隊的不插隊」時，就是
佇列；規定「每次挑編號最小的」時，就是第~\ref{sec:lex-topo}~節的優先佇列。

\textbf{程式}

\begin{example}[label=ex:checklist]{cpp}{(checklist.cpp) 成果展布置的工作清單}
#include <algorithm>
#include <iostream>
#include <queue>
#include <vector>
using namespace std;

int main() {
    ios::sync_with_stdio(false);
    cin.tie(nullptr);
    int n, m;
    cin >> n >> m;
    vector<vector<int>> adj(n + 1);
    vector<int> indeg(n + 1, 0);
    for (int i = 0; i < m; i++) {
        int a, b;
        cin >> a >> b;
        adj[a].push_back(b);
        indeg[b]++;
    }
    queue<int> q;
    for (int v = 1; v <= n; v++) {               // 編號由小到大放入
        if (indeg[v] == 0) {
            q.push(v);
        }
    }
    vector<int> order, fresh;
    while (!q.empty()) {
        int u = q.front();
        q.pop();
        order.push_back(u);
        fresh.clear();                           // 這一次新解鎖的工作
        for (int v : adj[u]) {
            indeg[v]--;
            if (indeg[v] == 0) {
                fresh.push_back(v);
            }
        }
        sort(fresh.begin(), fresh.end());
        for (int v : fresh) {
            q.push(v);
        }
    }
    for (int i = 0; i < n; i++) {
        cout << order[i] << (i + 1 < n ? ' ' : '\n');
    }
    return 0;
}
\end{example}

\iosample{
6 6 \\
3 6 \\
1 4 \\
3 2 \\
4 5 \\
2 5 \\
1 6 \\
}{
1 3 4 2 6 5 \\
}

做完 $3$ 時，$3$ 的串列依輸入是 $6, 2$，兩個都減到 $0$。 \texttt{fresh} 排
序成 $2, 6$ 再放入，佇列從 $4$ 變成 $4, 2, 6$。

\begin{itemize}
\item 題目保證沒有矛盾，所以不檢查 \texttt{order} 的長度。 沒有這個保證的題
目，要加上前面 \texttt{kahn.cpp} 的檢查。
\end{itemize}

\textbf{複雜度}

Kahn 本身 $O(n + m)$，加上各批排序共 $O(n \log n)$。 最大的測試資料在本書測
試的環境下約 $0.15$ 秒。

\subsection{常見錯誤}

\begin{itemize}
\item \textbf{用優先佇列}：把「依編號由小到大」讀成「每次拿可以做的裡面編號最
小的」。 例子裡的輸入 \texttt{4 2}、\texttt{1 4}、\texttt{2 3}（三行）會印出
\texttt{1 2 3 4}，正確答案是 \texttt{1 2 4 3}。 兩種規則在範例~1 碰巧也不同：
優先佇列印出 \texttt{1 3 2 4 5 6}。
\item \textbf{新解鎖的沒有排序}：拿掉 \texttt{sort} 那一行，範例的輸入會照
串列的順序放入 $6$、$2$，印出 \texttt{1 3 4 6 2 5}。 範例~2 的每一批都只有一
個點，這個錯誤在範例~2 看不出來。
\item \textbf{入度一減就放入}：少了 \texttt{if (indeg[v] == 0)}，每減一次就
放一次。 點 $6$ 有兩個前置，會被放進佇列兩次，範例的輸入印出
\texttt{1 3 4 6 2 6}（只印前 $n$ 個），$6$ 出現兩次，$5$ 沒有出現。
\end{itemize}

\section{用 DFS 做拓樸排序}
\label{sec:dfs-topo}

另一種方法用第 1 本第~11~章的深度優先搜尋。 關鍵是 DFS 裡一個點「結束」的時機：
\texttt{dfs(u)} 回傳時，從 $u$ 走得到的每一個點都已經結束了。 換句話說，在有
向無環圖上，每條邊 $u \to v$ 的 $v$ 都比 $u$ 早結束。 把點照結束的先後記下來，
最後\textbf{倒過來}，每條邊的 $u$ 就都排在 $v$ 前面，正好是拓樸順序。

為什麼 $v$ 一定比 $u$ 早結束？ DFS 在 $u$ 的串列看到 $v$ 的時候，$v$ 有三種
可能：

\begin{itemize}
\item $v$ 還沒去過：DFS 會從 $u$ 走進 $v$，等 $v$ 結束才回到 $u$，所以 $v$ 先
結束。
\item $v$ 已經結束：不用說，$v$ 先結束。
\item $v$ 去過了但還沒結束：$v$ 還在呼叫堆疊上，表示 DFS 是從 $v$ 一路走到 $u$
的，再加上 $u \to v$ 這條邊就繞回 $v$，是一個環。 DAG 上不會發生。
\end{itemize}

所以每個點要記三種狀態：還沒去過、進行中（還在堆疊上）、已經結束。 習慣上用
顏色稱呼，白、灰、黑，程式裡用 $0$、$1$、$2$。 第三種情況「看到灰色的點」就
是找到環，而且環就在堆疊上：從那個灰色的點到堆疊頂端，再走回它。

在圖~\ref{fig:hotpot} 上，從編號小的點開始做 DFS，串列的順序照輸入：

\begin{center}
\begin{tabular}{cll}
\toprule
從哪裡開始 & 經過 & 依序結束 \\
\midrule
1 & 1 沒有出去的邊 & 1 \\
2 & $2 \to 1$，1 已結束 & 2 \\
3 & $3 \to 2$ 已結束；$3 \to 7$，$7 \to 1$ 已結束；$3 \to 5$，$5 \to 1$ 已結束
& 7、5、3 \\
4 & $4 \to 1$ 已結束 & 4 \\
5 & 已經是黑色，跳過 & \\
6 & $6 \to 5$ 已結束 & 6 \\
7 & 已經是黑色，跳過 & \\
\bottomrule
\end{tabular}
\end{center}

結束的先後是 $1, 2, 7, 5, 3, 4, 6$，倒過來是 $6, 4, 3, 5, 7, 2, 1$，和 Kahn
排出來的不一樣，但也是拓樸順序。 注意 DFS 不需要從入度 $0$ 的點開始：從點 $1$
開始時，$1$ 就是最先結束的那個，倒過來剛好在最後。

\subsection{寫成迴圈}

第 1 本第~11~章量過，遞迴的 DFS 在一條長鏈上大約十幾萬層就會用光堆疊。 這一章的題
目 $n$ 常常到 $2 \times 10^5$，所以直接用第 1 本第~11~章的「改法二」：自己用一個
\texttt{vector} 當堆疊，每一格記住點和「下一個要看第幾個鄰居」。 鄰居看完、
彈出的那一刻，就是這個點結束的時刻。 堆疊上的點剛好是所有灰色的點，找到環時
要印的那一段就在堆疊裡。

\begin{example}[label=ex:dfs-topo]{cpp}{(dfs\_topo.cpp) 用 DFS 排拓樸順序或找出一個環}
#include <algorithm>
#include <iostream>
#include <utility>
#include <vector>
using namespace std;

int main() {
    int n, m;
    cin >> n >> m;
    vector<vector<int>> adj(n + 1);
    for (int i = 0; i < m; i++) {
        int a, b;
        cin >> a >> b;
        adj[a].push_back(b);
    }
    // 0：還沒去過；1：還在堆疊上（沒結束）；2：已經結束
    vector<int> color(n + 1, 0);
    vector<int> finished;                       // 結束的先後
    for (int s = 1; s <= n; s++) {
        if (color[s] != 0) {
            continue;
        }
        vector<pair<int, int>> st = {{s, 0}};   // （點, 下一個要看的鄰居位置）
        color[s] = 1;
        while (!st.empty()) {
            int u = st.back().first;
            int i = st.back().second;
            if (i == (int)adj[u].size()) {      // 鄰居都看完：u 結束
                color[u] = 2;
                finished.push_back(u);
                st.pop_back();
                continue;
            }
            st.back().second++;
            int v = adj[u][i];
            if (color[v] == 0) {
                color[v] = 1;
                st.push_back({v, 0});
            } else if (color[v] == 1) {         // 走回還沒結束的點：有環
                int k = st.size() - 1;
                while (st[k].first != v) {
                    k--;
                }
                cout << "cycle:";
                for (int j = k; j < (int)st.size(); j++) {
                    cout << ' ' << st[j].first;
                }
                cout << ' ' << v << '\n';
                return 0;
            }
        }
    }
    reverse(finished.begin(), finished.end());
    for (int i = 0; i < n; i++) {
        cout << finished[i] << (i + 1 < n ? ' ' : '\n');
    }
    return 0;
}
\end{example}

\iosample{
7 8 \\
3 2 \\
3 7 \\
3 5 \\
6 5 \\
2 1 \\
5 1 \\
7 1 \\
4 1 \\
}{
6 4 3 5 7 2 1 \\
}

同樣多加一條 \texttt{5 3}，程式印出 \texttt{cycle: 3 5 3}：從 $3$ 走到 $5$，
$5$ 的串列裡先看到 $1$（已結束），再看到 $3$，$3$ 是灰色。 堆疊上從 $3$ 到頂
端是 $3, 5$，接回 $3$ 就是環。 和 Kahn 不同，這裡印的每一個點都在環上。

\begin{itemize}
\item 外層 \texttt{for} 從每個還沒去過的點出發，和第 1 本第~11~章數連通元件一樣。 有
些點從別的點走不到（例如入度 $0$ 的點），不這樣做就會漏掉。
\item 每個點變灰、變黑各一次，每條邊看一次，總共 $O(n + m)$。
\end{itemize}

\subsection{常見錯誤：只用兩種狀態}

很多人第一次寫找環，只用第 1 本第~11~章的 \texttt{seen}：走到已經看過的點就說有環。
在有向圖上這是錯的；無向圖找環的方法又不一樣，第~7~章再談。 四條邊 $1 \to 2$、$1 \to 3$、$2 \to 4$、$3 \to 4$
沒有環，但從 $1$ 走 $1 \to 2 \to 4$，退回後再走 $1 \to 3$，$3$ 的鄰居 $4$ 已
經看過，就被誤判成有環。 $4$ 這時已經結束，是黑色；只有走到灰色的點才是環。
兩條路在 $4$ 會合，在 DAG 上很常見，例如圖~\ref{fig:hotpot} 的點 $1$ 有四條
路可以走到。

\subsection{Kahn 還是 DFS}

\begin{center}
\begin{tabular}{p{0.2\linewidth}p{0.34\linewidth}p{0.34\linewidth}}
\toprule
 & Kahn & DFS \\
\midrule
需要的資料 & 入度、鄰接串列 & 鄰接串列、三種顏色 \\
有環時 & 排出的點不到 $n$ 個；剩下的點在環上或環的下游 & 走到灰色的點；堆疊上就是一個環 \\
順序 & 容器決定，可以換成優先佇列控制 & 由 DFS 的走法決定，很難控制 \\
長鏈 & 本來就是迴圈，不怕 & 遞迴會太深，要改成自己的堆疊 \\
\bottomrule
\end{tabular}
\end{center}

只要排出順序、判斷有沒有環時，兩種都可以，本書的例題多半用 Kahn：它是迴圈，
不用擔心堆疊，順序也比較好控制。 需要印出一個具體的環時，DFS 比較方便。

\begin{apcsnote}
實作題的圖常常到 $2 \times 10^5$ 個點，遞迴的 DFS 有用光堆疊的風險（第 1 本第~11~
章）。 只需要拓樸順序或判斷有沒有環時，直接寫 Kahn。 題目沒有說「保證沒有
矛盾」時，記得檢查排出來的點夠不夠 $n$ 個。
\end{apcsnote}

\section{在 DAG 上做動態規劃}
\label{sec:dag-dp}

第 1 本第~9~章把「最少、最多、有幾種、做不做得到」對應到狀態圖上的「最短路、最長
路、路徑數、走不走得到」，但那時的圖只畫在紙上，計算順序一眼就看得出來。 這一
節的圖寫在輸入裡，順序用 Kahn 算出來，其他和第 1 本第~9~章一樣：照拓樸順序一個點一
個點處理，處理到一個點的時候，所有指向它的點都已經處理完了。

轉移有兩種寫法。 \textbf{往回拉}：處理 $v$ 時，看每一條指向 $v$ 的邊
$u \to v$，用 $u$ 的答案算 $v$，需要反向圖。 \textbf{往前推}：處理 $u$ 時，
$u$ 的答案已經確定，沿著每一條 $u \to v$ 去更新 $v$。 推的寫法直接用鄰接串
列，不用另外建反向圖，本節都用推的。 第 1 本第~9~章的跳格子例題也比較過這兩種寫法。

以「從 $s$ 出發」的題目來說，還有一件第 1 本第~9~章不常碰到的事：從 $s$ 走不到的點。
它們在拓樸順序裡照樣會被處理，答案卻必須是「不存在」。 如果把它們的答案設成
$0$，求最大值時就會被當成「總和 $0$ 的一條路」，算錯。 下面的例題會看到實際的
錯誤輸出。

\subsection{例題：夜市路線的最高滿意度與方案數}
\label{sec:ex-nightmarket}

\textbf{題意}（\akiref{best-night-market-route-and-how-many-ways}{夜市路線的最高滿意度與方案數}）

\begin{itemize}
\item 夜市有 $n$ 個攤位和 $m$ 條單向的巷子，巷子不會繞成圈。 攤位 $i$ 的滿意
度是 $a_i$，可能是負數。
\item 從攤位 $s$ 出發，每一步沿一條巷子走到下一個攤位，經過的每個攤位（含起
點和終點）都要買，得到它的滿意度。 只能在沒有往外巷子的攤位結束；還有往外的
巷子就一定要繼續走。
\item 求所有路線中最大的總滿意度，以及達到最大值的路線有幾條（除以
$10^9 + 7$ 的餘數）。
\item $n \le 2 \times 10^5$，$m \le 4 \times 10^5$，$|a_i| \le 10^9$。
\end{itemize}

例如下圖，從攤位 $1$ 出發，可以停的是 $5$ 和 $6$。 五條路線裡，
$1 \to 2 \to 4 \to 5$ 和 $1 \to 3 \to 4 \to 5$ 都是 $2 + 3 - 1 + 4 = 8$，最
大，答案是 \texttt{8 2}。

\begin{center}
\begin{tikzpicture}[x=1.6cm, y=1.1cm]
\node[tvhl] (1) at (0,0) {1};
\node[tv] (2) at (1,0.8) {2};
\node[tv] (3) at (1,-0.8) {3};
\node[tv] (4) at (2,0) {4};
\node[tv] (5) at (3,0.8) {5};
\node[tv] (6) at (3,-0.8) {6};
\node[idxlbl, below=1pt] at (1.south) {$a = 2$};
\node[idxlbl, above=1pt] at (2.north) {$a = 3$};
\node[idxlbl, below=1pt] at (3.south) {$a = 3$};
\node[idxlbl, below=1pt] at (4.south) {$a = -1$};
\node[idxlbl, above=1pt] at (5.north) {$a = 4$};
\node[idxlbl, below=1pt] at (6.south) {$a = 2$};
\draw[tahl] (1) -- (2);
\draw[tahl] (1) -- (3);
\draw[tahl] (2) -- (4);
\draw[tahl] (3) -- (4);
\draw[tahl] (4) -- (5);
\draw[ta] (4) -- (6);
\draw[ta] (3) -- (6);
\end{tikzpicture}
\diagcaption{範例 1 的夜市。 橘色的箭頭是兩條最好的路線用到的巷子。}
\label{fig:nightmarket}
\end{center}

\textbf{想法}

最直接的做法是從 $s$ 開始 DFS，把每一條路線都走一遍，邊走邊加滿意度。 路線
的條數可能非常多：把攤位排成一層兩個、每一層的兩個都連到下一層的兩個，$k$
層就有 $2^k$ 條路線。 $n = 2 \times 10^5$ 時這是天文數字，一條一條走一定來不
及。 不過這個暴力解在 $n \le 15$ 的小範圍很好寫，可以拿來對拍。

照第 1 本第~9~章的四件事來設計：

\begin{itemize}
\item \textbf{狀態}：\texttt{best[v]} 是「從 $s$ 走到 $v$ 的所有走法裡，經過的
攤位滿意度總和的最大值」，\texttt{ways[v]} 是達到 \texttt{best[v]} 的走法數。
另外用 \texttt{reach[v]} 記 $v$ 從 $s$ 走不走得到。
\item \textbf{轉移}：處理 $u$ 時，對每條 $u \to v$ 算候選值
$\texttt{best[u]} + a_v$。 $v$ 第一次被走到，或候選值比 \texttt{best[v]}
大，就換掉 \texttt{best[v]}，\texttt{ways[v]} 改成 \texttt{ways[u]}；和
\texttt{best[v]} 一樣大，就把 \texttt{ways[u]} 加進 \texttt{ways[v]}。
\item \textbf{初始值}：$\texttt{reach[s]} = \text{true}$，
$\texttt{best[s]} = a_s$，$\texttt{ways[s]} = 1$；其他點都還走不到。
\item \textbf{計算順序}：Kahn 排出的拓樸順序。 走不到的點跳過。
\end{itemize}

最後在「走得到、出度 $0$」的攤位裡找最大的 \texttt{best}，把同樣大的那幾個攤
位的 \texttt{ways} 加起來。

\texttt{ways} 的轉移為什麼對：達到 \texttt{best[v]} 的每一條走法，最後一步都
是某一條 $u \to v$，而且前面那一段一定是走到 $u$ 的最好走法之一（前面更差，
加上 $a_v$ 也會更差）。 依最後一步從哪個 $u$ 來分類，每一類有
\texttt{ways[u]} 條，不重複也不遺漏。 這和第 1 本第~9~章數走法的想法相同，只是多了
「只數最好的」這個條件。

用範例跑一次。 拓樸順序是 $1, 2, 3, 4, 5, 6$：

\begin{center}
\begin{tabular}{cl}
\toprule
處理 & 往前推的結果（\texttt{best}，\texttt{ways}） \\
\midrule
1 & 2：第一次走到，$(5, 1)$；3：第一次走到，$(5, 1)$ \\
2 & 4：第一次走到，$5 - 1 = 4$，$(4, 1)$ \\
3 & 4：候選 $4$ 一樣大，$(4, 2)$；6：第一次走到，$5 + 2 = 7$，$(7, 1)$ \\
4 & 5：第一次走到，$4 + 4 = 8$，$(8, 2)$；6：候選 $4 + 2 = 6 < 7$，不變 \\
5、6 & 沒有出去的巷子 \\
\bottomrule
\end{tabular}
\end{center}

可以停的 $5$、$6$ 分別是 $(8, 2)$ 和 $(7, 1)$，答案 \texttt{8 2}。

同樣的模型也出現在別的地方：工作流程上每個步驟有收益或成本，找最划算的流程，
並知道有幾種同樣好的做法；遊戲的劇情分支圖上，找分數最高的結局路線。

\textbf{程式}

\begin{example}[label=ex:nightmarket]{cpp}{(nightmarket.cpp) 夜市路線的最高滿意度與方案數}
#include <iostream>
#include <vector>
using namespace std;

const long long MOD = 1000000007;

int main() {
    ios::sync_with_stdio(false);
    cin.tie(nullptr);
    int n, m, s;
    cin >> n >> m >> s;
    vector<long long> a(n + 1);
    for (int i = 1; i <= n; i++) {
        cin >> a[i];
    }
    vector<vector<int>> adj(n + 1);
    vector<int> indeg(n + 1, 0);
    for (int i = 0; i < m; i++) {
        int u, v;
        cin >> u >> v;
        adj[u].push_back(v);
        indeg[v]++;
    }

    // 第一步：Kahn 排出拓樸順序（vector 當佇列，order 本身就是結果）
    vector<int> order;
    for (int v = 1; v <= n; v++) {
        if (indeg[v] == 0) {
            order.push_back(v);
        }
    }
    for (int k = 0; k < (int)order.size(); k++) {
        int u = order[k];
        for (int v : adj[u]) {
            indeg[v]--;
            if (indeg[v] == 0) {
                order.push_back(v);
            }
        }
    }

    // 第二步：照拓樸順序往後推
    vector<bool> reach(n + 1, false);   // 從 s 走得到嗎
    vector<long long> best(n + 1, 0);   // 從 s 走到 v 的最大總和
    vector<long long> ways(n + 1, 0);   // 達到 best[v] 的路線數（取餘數）
    reach[s] = true;
    best[s] = a[s];
    ways[s] = 1;
    for (int u : order) {
        if (!reach[u]) {
            continue;
        }
        for (int v : adj[u]) {
            long long cand = best[u] + a[v];
            if (!reach[v] || cand > best[v]) {      // 第一次走到，或更好
                reach[v] = true;
                best[v] = cand;
                ways[v] = ways[u];
            } else if (cand == best[v]) {           // 一樣好：路線數相加
                ways[v] = (ways[v] + ways[u]) % MOD;
            }
        }
    }

    // 第三步：只看走得到、而且沒有往外巷子的攤位
    bool found = false;
    long long ans = 0, cnt = 0;
    for (int v = 1; v <= n; v++) {
        if (!reach[v] || !adj[v].empty()) {
            continue;
        }
        if (!found || best[v] > ans) {
            found = true;
            ans = best[v];
            cnt = ways[v];
        } else if (best[v] == ans) {
            cnt = (cnt + ways[v]) % MOD;
        }
    }
    cout << ans << ' ' << cnt << '\n';
    return 0;
}
\end{example}

\iosample{
6 7 1 \\
2 3 3 -1 4 2 \\
1 2 \\
1 3 \\
2 4 \\
3 4 \\
4 5 \\
4 6 \\
3 6 \\
}{
8 2 \\
}

\begin{itemize}
\item 這裡的 Kahn 用 \texttt{order} 本身當佇列：\texttt{k} 是佇列的前端，
\texttt{push\_back} 是放到尾端。 拿出來的順序就是 \texttt{order} 裡的順序，
不用另外存。 第 1 本第~6~章的 \texttt{queue} 寫法也可以。
\item 比大小用的是 \texttt{best}，它沒有取餘數；取餘數的只有 \texttt{ways}。
路線數取了餘數之後就不能拿來比大小，這題剛好只比 \texttt{best}。
\item \texttt{best} 最大約 $2 \times 10^5 \times 10^9 = 2 \times 10^{14}$，
最小約 $-2 \times 10^{14}$，要用 \texttt{long long}。 \texttt{ways[v] + ways[u]}
兩個都小於 $10^9 + 7$，相加不會超過 \texttt{long long}。
\end{itemize}

\textbf{複雜度}

Kahn 和 DP 各看每條邊一次，$O(n + m)$。 最大的測試資料在本書測試的環境下約
$0.2$ 秒，大部分花在讀入 $4 \times 10^5$ 條邊。

\subsection{常見錯誤}

\begin{itemize}
\item \textbf{沒有記走不走得到}：拿掉 \texttt{reach}，把每個點的 \texttt{best}
當成一開始是 $0$。 範例~2 有五個攤位，滿意度 $5, -3, -4, 6, -1$，巷子
$1 \to 3$、$2 \to 3$、$3 \to 5$、$2 \to 5$、$4 \to 5$。 從攤位 $2$ 出發，路線只有 $2 \to 3 \to 5$（$-8$）和 $2 \to 5$（$-4$），答案是
\texttt{-4 1}。 這樣寫會印出 \texttt{0 0}：攤位 $3$ 的 \texttt{best} 本來就
是 $0$，比從 $2$ 走來的 $-7$ 大，不會被更新，\texttt{ways} 一直是 $0$；攤位
$5$ 也一樣。 只要所有路線的總和都是負的，就會出事。
\item \textbf{在任何攤位都能停}：最後一步沒有檢查出度，在所有走得到的攤位裡
找最大。 範例~2 會印出 \texttt{-3 1}，也就是只買起點攤位就走。
\item \textbf{總和用 \texttt{int}}：輸入一條 $2 \times 10^5$ 個攤位的長鏈、每
個攤位都是 $10^9$，正確答案是 \texttt{200000000000000 1}。 把 \texttt{a}、
\texttt{best}、\texttt{cand}、\texttt{ans} 改成 \texttt{int}，在本書測試的環
境下印出 \texttt{552894464 1}。
\item \textbf{照編號順序算}：用 \texttt{for (int u = 1; u <= n; u++)} 代替拓
樸順序。 範例~1 的編號剛好是拓樸順序，所以會過。 輸入 \texttt{3 2 3}、
\texttt{1 1 1}、\texttt{3 2}、\texttt{2 1}（四行）只有一條路線
$3 \to 2 \to 1$，答案是 \texttt{3 1}，這樣寫會印出 \texttt{0 0}：處理攤位
$2$ 時它還沒被走到，等到處理 $3$ 才把 $2$ 標成走得到，但 $2$ 已經處理過了，
攤位 $1$ 永遠走不到。
\end{itemize}

\subsection{關鍵路徑}
\label{sec:critical-path}

回到章首的火鍋。 現在每件事都有一個時間（分鐘），而且家裡有好幾個人，沒有先
後關係的事可以同時做：

\begin{center}
\begin{tabular}{lccccccc}
\toprule
編號 & 1 & 2 & 3 & 4 & 5 & 6 & 7 \\
工作 & 開火上桌 & 洗菜切菜 & 買菜 & 擺碗筷 & 煮湯底 & 借電磁爐 & 調沾醬 \\
時間 & 5 & 30 & 40 & 5 & 50 & 20 & 10 \\
\bottomrule
\end{tabular}
\end{center}

最快多久能開飯？ 一件事最早什麼時候開始，就是它所有前置工作裡最晚完成的時
刻；沒有前置的從 $0$ 開始。 這正是第 1 本第~6~章「層層彙整的報告」的式子，只是那裡
每人只有一位主管。 照拓樸順序往前推：買菜 $0$ 到 $40$，借電磁爐 $0$ 到 $20$，
煮湯底要等兩者，從 $40$ 開始、$90$ 完成，開火上桌從 $90$ 開始、$95$ 完成。
所以最快 $95$ 分鐘。

再問：哪些事一拖延，開飯就會延後？ 這要倒過來算每件事\textbf{最晚}什麼時候
要完成。 開火上桌最晚 $95$ 完成，所以最晚 $90$ 開始；它的四個前置都必須在
$90$ 以前完成。 煮湯底最晚 $90$ 完成、$40$ 開始，所以買菜和借電磁爐最晚 $40$
要完成。 一般地說，一件事最晚的完成時刻，是它所有後續工作最晚開始時刻的最小
值；沒有後續工作的，就是全部完成的時刻 $95$。 這要照拓樸順序\textbf{倒過來}
處理，因為要先知道後面的事。

最晚和最早的完成時刻之差，叫\textbf{寬裕時間}（slack）：這件事最多可以拖延
這麼久，不影響開飯。 寬裕時間是 $0$ 的事一點都不能拖，它們連成的一條路叫\textbf{關
鍵路徑}（critical path）。 這裡是 買菜 $\to$ 煮湯底 $\to$ 開火上桌，$40 + 50
+ 5 = 95$，剛好是全部的時間。 想提早開飯，只有縮短關鍵路徑上的事才有用；洗菜
切菜再快也沒有幫助。

\begin{example}[label=ex:critical]{cpp}{(critical.cpp) 最早、最晚開始時刻與寬裕時間}
#include <algorithm>
#include <iostream>
#include <vector>
using namespace std;

int main() {
    int n, m;
    cin >> n >> m;
    vector<long long> t(n + 1);
    for (int i = 1; i <= n; i++) {
        cin >> t[i];
    }
    vector<vector<int>> adj(n + 1);
    vector<int> indeg(n + 1, 0);
    for (int i = 0; i < m; i++) {
        int a, b;
        cin >> a >> b;
        adj[a].push_back(b);
        indeg[b]++;
    }
    vector<int> order;                       // Kahn，寫法同夜市的例題
    for (int v = 1; v <= n; v++) {
        if (indeg[v] == 0) {
            order.push_back(v);
        }
    }
    for (int k = 0; k < (int)order.size(); k++) {
        for (int v : adj[order[k]]) {
            if (--indeg[v] == 0) {
                order.push_back(v);
            }
        }
    }

    // 往前推：最早開始 es，最早完成 es + t
    vector<long long> es(n + 1, 0);
    long long total = 0;
    for (int u : order) {
        total = max(total, es[u] + t[u]);
        for (int v : adj[u]) {
            es[v] = max(es[v], es[u] + t[u]);
        }
    }
    // 倒著拉：最晚完成 lf，沒有後續工作的就是 total
    vector<long long> lf(n + 1, total);
    for (int k = n - 1; k >= 0; k--) {
        int u = order[k];
        for (int v : adj[u]) {
            lf[u] = min(lf[u], lf[v] - t[v]);   // v 最晚開始之前要做完
        }
    }

    cout << total << '\n';
    for (int v = 1; v <= n; v++) {
        long long slack = lf[v] - (es[v] + t[v]);
        cout << v << ' ' << es[v] << ' ' << lf[v] - t[v] << ' ' << slack << '\n';
    }
    return 0;
}
\end{example}

\iosample{
7 8 \\
5 30 40 5 50 20 10 \\
3 2 \\
3 7 \\
3 5 \\
6 5 \\
2 1 \\
5 1 \\
7 1 \\
4 1 \\
}{
95 \\
1 90 90 0 \\
2 40 60 20 \\
3 0 0 0 \\
4 0 85 85 \\
5 40 40 0 \\
6 0 20 20 \\
7 40 80 40 \\
}

輸出的每一行是「編號、最早開始、最晚開始、寬裕時間」。 寬裕時間是 $0$ 的是
$1$、$3$、$5$，就是上面找到的關鍵路徑。

\begin{itemize}
\item 往前推和倒著拉用同一個拓樸順序，只是方向相反。 倒著處理時，$u$ 的所有後
續工作都在 $u$ 後面，已經算好了，所以這裡用「拉」的寫法，直接用鄰接串列，不
需要反向圖。
\item 程式假設沒有環。 有環時 \texttt{order} 不到 $n$ 個，\texttt{order[k]}
會讀到陣列外面；題目沒有保證時，要先檢查。
\item 最早完成時刻就是「從某個起點走到這個點、把經過的時間加起來」的最大值，
所以全部的時間就是這張圖上的\textbf{最長路}，關鍵路徑就是最長的那一條。 這
和夜市的例題是同一件事：點上有權重，沿著拓樸順序取最大值。
\end{itemize}

\subsection{資料結構課本的說法：AOV 網路與 AOE 網路}
\label{sec:aov-aoe}

大學的資料結構課本（研究所考試的參考書也是）講這一節的內容時，用的是另一套名
字，圖的畫法也不太一樣。 算的東西和上面完全相同，這裡把兩套說法對照一次，在
課堂上遇到時就認得出來。

\textbf{AOV 網路。} 像圖~\ref{fig:hotpot} 那樣，每件工作畫成一個點、先後關係
畫成邊，課本叫它\textbf{AOV 網路}（activity on vertex network，頂點工作網路）。
名字裡的 activity 就是一件工作，課本譯成\textbf{活動}，下面也跟著這樣叫。
課本在這種圖上做的事就是拓樸排序：排出一個可行的做事順序；排不出來表示圖上
有環，工作的先後規定互相矛盾。 本章第~\ref{sec:kahn}、\ref{sec:dfs-topo}~節
做的正是這件事。

\textbf{AOE 網路。} 另一種畫法反過來：每個活動畫成一條\textbf{邊}，活動的時間
寫在邊上；點代表\textbf{事件}（event），也就是「某幾個活動都做完了，接下來的
活動可以開始」的那一刻。 這種圖叫\textbf{AOE 網路}（activity on edge network，
邊工作網路），通常有一個代表專案開始的起點、一個代表專案完成的終點。 火鍋畫
成 AOE 網路是這樣：

\begin{center}
\begin{tikzpicture}[x=1.3cm, y=1cm]
\node[tv] (v0) at (0,0) {$v_0$};
\node[tv] (v1) at (2.6,1.5) {$v_1$};
\node[tv] (v2) at (2.6,-1.5) {$v_2$};
\node[tv] (v3) at (4.8,-0.5) {$v_3$};
\node[tv] (v4) at (7.6,0) {$v_4$};
\node[tv] (v5) at (10,0) {$v_5$};
\draw[tahl] (v0) -- node[sloped, above, font=\small] {買菜 40} (v1);
\draw[ta] (v0) -- node[sloped, below, font=\small] {借電磁爐 20} (v2);
\draw[ta, rounded corners=6pt] (v0.south) -- (0,-2.4) -- node[above, font=\small] {擺碗筷 5} (7.6,-2.4) -- (v4.south);
\draw[ta] (v1) to[bend left=18] node[above, font=\small] {洗菜切菜 30} (v4);
\draw[ta] (v1) to[bend right=4] node[above right, pos=0.35, font=\small] {調沾醬 10} (v4);
\draw[tahl, dashed] (v1) -- node[left, pos=0.4, font=\small] {0} (v3);
\draw[ta, dashed] (v2) -- node[below right, pos=0.3, font=\small] {0} (v3);
\draw[tahl] (v3) -- node[below right, pos=0.4, font=\small] {煮湯底 50} (v4);
\draw[tahl] (v4) -- node[above, font=\small] {開火上桌 5} (v5);
\end{tikzpicture}
\diagcaption{火鍋的 AOE 網路。 $v_0$ 開始，$v_1$ 菜買回來了，$v_2$ 電磁爐借到
了，$v_3$ 可以煮湯底，$v_4$ 可以開火，$v_5$ 開飯。 虛線是時間 $0$ 的虛擬活動，
橘色是關鍵路徑。}
\label{fig:hotpot-aoe}
\end{center}

圖裡的兩條虛線是\textbf{虛擬活動}（dummy activity），時間是 $0$。 需要它們，
是因為煮湯底要等菜和電磁爐都到，洗菜切菜卻只等菜。 如果把「菜買回來」和「電
磁爐借到」畫成同一個事件，洗菜切菜就被迫等電磁爐，意思變了。 一般地說，把
AOV 網路的每個點拆成「開始」和「完成」兩個事件，中間連一條邊放工作時間，原
本的先後關係改成時間 $0$ 的邊，就得到一張 AOE 網路，只是虛擬活動多一點。 所
以兩種畫法描述的是同一件事。

在 AOE 網路上，課本先算每個事件的兩個時刻，再算每個活動的兩個時刻：

\begin{itemize}
\item 事件 $v$ 的\textbf{最早發生時間} $ve(v)$：照拓樸順序往前推，取所有指進
來的邊 $u \to v$ 的 $ve(u) + \text{時間}$ 的最大值，起點是 $0$。
\item 事件 $v$ 的\textbf{最晚發生時間} $vl(v)$：照拓樸順序倒過來，取所有指出
去的邊 $v \to w$ 的 $vl(w) - \text{時間}$ 的最小值，終點等於 $ve(\text{終點})$。
\item 活動 $u \to v$ 的\textbf{最早開始時間} $e = ve(u)$，\textbf{最晚開始時間}
$l = vl(v) - \text{時間}$。 $l - e$ 就是寬裕時間；$l = e$ 的活動叫\textbf{關
鍵活動}（critical activity），全由關鍵活動組成、從起點到終點的路徑就是關鍵路
徑。
\end{itemize}

火鍋的六個事件，$ve$ 依序是 $0, 40, 20, 40, 90, 95$，$vl$ 依序是
$0, 40, 40, 40, 90, 95$。 由它們算出每個活動的 $e$、$l$：

\begin{center}
\begin{tabular}{lcccc}
\toprule
活動（邊） & $e = ve(\text{起})$ & $l = vl(\text{迄}) - \text{時間}$ & $l - e$ & 範例~\ref{ex:critical} 的輸出 \\
\midrule
買菜 $v_0 \to v_1$ & $0$ & $40 - 40 = 0$ & $0$ & \texttt{3 0 0 0} \\
借電磁爐 $v_0 \to v_2$ & $0$ & $40 - 20 = 20$ & $20$ & \texttt{6 0 20 20} \\
擺碗筷 $v_0 \to v_4$ & $0$ & $90 - 5 = 85$ & $85$ & \texttt{4 0 85 85} \\
洗菜切菜 $v_1 \to v_4$ & $40$ & $90 - 30 = 60$ & $20$ & \texttt{2 40 60 20} \\
調沾醬 $v_1 \to v_4$ & $40$ & $90 - 10 = 80$ & $40$ & \texttt{7 40 80 40} \\
煮湯底 $v_3 \to v_4$ & $40$ & $90 - 50 = 40$ & $0$ & \texttt{5 40 40 0} \\
開火上桌 $v_4 \to v_5$ & $90$ & $95 - 5 = 90$ & $0$ & \texttt{1 90 90 0} \\
\bottomrule
\end{tabular}
\end{center}

每一列都和範例~\ref{ex:critical} 印出的「最早開始、最晚開始、寬裕時間」相同。
程式裡的 \texttt{es} 就是「這件工作可以開始」那個事件的 $ve$，\texttt{lf}
就是「這件工作做完」那個事件的 $vl$。 換句話說，本章的程式是直接在 AOV 網路上
算，省掉了拆事件和虛擬活動。

兩套名字的對照：

\begin{center}
\begin{tabular}{p{0.38\linewidth}p{0.52\linewidth}}
\toprule
本章的說法 & 資料結構課本常見的說法 \\
\midrule
工作畫成點的 DAG & AOV 網路（頂點工作網路） \\
工作畫成邊、點是事件的 DAG & AOE 網路（邊工作網路） \\
一件工作 & 活動（activity） \\
最早開始、最晚開始 & 活動的最早開始時間 $e$、最晚開始時間 $l$；事件的最早、最晚發生時間 $ve$、$vl$ \\
寬裕時間 & slack，$l - e$；專案管理裡叫 float \\
寬裕時間 $0$ 的工作 & 關鍵活動（critical activity） \\
關鍵路徑 & critical path，有的譯本寫「臨界路徑」 \\
\bottomrule
\end{tabular}
\end{center}

課本常接著問：要提早完工，該縮短哪些活動？ 只有縮短關鍵活動才有用，但不一定
夠。 如果有兩條關鍵路徑，只縮其中一條上的活動，另一條還是一樣長，完工時間不
變；要縮兩條共用的活動，或兩條都縮。 縮短之後關鍵路徑也可能換成另一條，要重
新算一次。 工業工程和專案管理把這整套方法叫\textbf{要徑法}（critical path
method，CPM）。

\subsection{什麼時候用、什麼時候不能用}

\begin{itemize}
\item 題目給的是有向圖，而且保證（或可以檢查）沒有環，問最長、最短、幾種、走
不走得到：先 Kahn，再照順序推。 權重是負數也沒關係，這是 DAG 比一般圖好做的
地方。
\item 一般的有向圖（可能有環）上的最長路沒有好的做法，最短路則要用
Dijkstra 等演算法（本書第~8~章）；沒有權重的最短路用第 1 本第~11~章的 BFS。
\item 圖有環、但環上的點可以合併成一個點時（例如一群互相走得到的點可以一次全
部拿走），合併之後就是 DAG。 實作練習裡「不往下走的登山收集」是這種題目。
\end{itemize}

\section{字典序最小的拓樸順序}
\label{sec:lex-topo}

拓樸順序通常不只一個，有些題目要求「字典序最小」的那一個：第一個數盡量小；
第一個數相同時，第二個數盡量小；依此類推，和字典裡的單字比先後的方法一樣，
叫\textbf{字典序}（lexicographic order）。 例如 $3, 1, 5$ 比 $3, 2, 1$ 小，
因為第一個數相同，第二個數 $1 < 2$。

做法是把 Kahn 的佇列換成優先佇列，每次拿出目前入度 $0$ 的點裡編號最小的。 第
~\ref{sec:kahn}~節的圖~\ref{fig:hotpot} 上，這樣排出 $3, 2, 4, 6, 5, 7, 1$。
為什麼對：

\begin{itemize}
\item 排第一個時，能排的只有入度 $0$ 的點。 這些點每一個都真的可以排第一：排
了它之後，剩下的圖仍然沒有環，一定排得完。 所以第一個數最小可以是「入度 $0$
的點裡最小的」，而要字典序最小，就一定要選它。
\item 第一個數決定之後，問題變成在剩下的圖上找字典序最小的順序，同樣的理由再
用一次。 每一步都是同一個理由，所以每一步都選最小的就對了。
\end{itemize}

這是一個第 1 本第~7~章的貪心：每一步選當下最小的，選了不回頭。 能這樣做，是因為
「每個能排的點都不會讓後面排不完」。

用佇列的 Kahn 不保證字典序最小。 第~\ref{sec:ex-checklist}~小節的例題就是反
例：佇列排出 $1, 2, 4, 3$，字典序最小的是 $1, 2, 3, 4$。 DFS 的順序更難控
制，求字典序最小時不要用。

\textbf{另一種「編號小的優先」。} 有些題目要求的是：「$1$ 號排得越前面越好；在這個前提下，$2$ 號
排得越前面越好……」。 兩者聽起來很像，答案卻可能不同。 例如三個點，只有一條
邊 $3 \to 1$：字典序最小是 $2, 3, 1$；要 $1$ 號盡量前面，答案是 $3, 1, 2$，
$1$ 號排在第二，比前一個答案早。 第二種要求的做法是：把每條邊反過來，用每次
拿編號\textbf{最大}的優先佇列排一次，再把結果倒過來（本章練習第~5~題）。 讀
題時看清楚題目要的是哪一種。

\subsection{例題：合併多份名次表}
\label{sec:ex-rankings}

\textbf{題意}（\akiref{merging-partial-rankings}{合併多份名次表}）

\begin{itemize}
\item $n$ 位選手、$q$ 份部分名次表。 每份表是一串不重複的選手編號，表裡排在前
面的人，在總排名中也必須排在後面的人之前。 表的長度 $k_i$ 可以很長，也可以只
有一個人。
\item 排出一個符合所有表的總排名（$1$ 到 $n$ 的排列）。 存在的話輸出字典序最
小的一個；表之間互相矛盾時輸出 \texttt{-1}。
\item $n, q \le 2 \times 10^5$，所有表的長度總和 $L \le 4 \times 10^5$，單份
表可能接近 $n$ 人。
\end{itemize}

例如 $n = 6$，三份表是 $(5, 2, 4)$、$(6, 2)$、$(3, 1)$。 第一名只能從 $3$、
$5$、$6$ 裡選（$1$ 要等 $3$，$2$ 要等 $5$ 和 $6$，$4$ 要等 $2$），選 $3$；接
著 $1$ 可以排了，選 $1$；依此類推，答案是 \texttt{3 1 5 6 2 4}。

\textbf{想法}

先把名次表變成圖。 最直接的想法是：一份表裡每一對「前面的人、後面的人」都建
一條邊。 長度 $k$ 的表有 $k(k - 1)/2$ 對，$k = 2 \times 10^5$ 時大約
$2 \times 10^{10}$ 條邊，記憶體和時間都不夠。 在本書測試的環境下，這樣寫跑最
大的幾筆測試資料，邊還沒建完，記憶體就超過 $4$ GB，程式中止。

觀察：一份表 $(x_1, x_2, \ldots, x_k)$ 只要建相鄰兩人的邊
$x_1 \to x_2 \to \cdots \to x_k$ 就夠了。 如果每條相鄰的邊都遵守，$x_1$ 在
$x_2$ 前、$x_2$ 在 $x_3$ 前，$x_1$ 自然在 $x_3$ 前，表裡任兩人的先後都會遵守。
所以兩種建法允許的總排名完全相同，邊數卻只有 $k - 1$ 條。 全部的邊加起來不超
過 $L$，是 $4 \times 10^5$。

建好圖，剩下的就是本節的優先佇列 Kahn：

\begin{itemize}
\item 字典序最小：每次拿入度 $0$ 的人裡編號最小的。
\item 矛盾：表之間矛盾，圖上就有環，排出來的人不到 $n$ 個，輸出 \texttt{-1}。
\item 沒有出現在任何表裡的人：入度是 $0$，一開始就在優先佇列裡，輪到他最小時
就會排進去，不用特別處理。
\item 同一對人在好幾份表裡都相鄰：會建出重複的邊，入度也加了好幾次。 拿出前面
那個人時，每條重複的邊各減一次，入度剛好減完，不用去掉重複的邊。
\end{itemize}

用範例跑一次。 相鄰的邊是 $5 \to 2$、$2 \to 4$、$6 \to 2$、$3 \to 1$：

\begin{center}
\begin{tabular}{clc}
\toprule
拿出 & 入度減一 & 之後的優先佇列 \\
\midrule
 & & $\{3, 5, 6\}$ \\
3 & 1：$1 \to 0$，放入 & $\{1, 5, 6\}$ \\
1 & & $\{5, 6\}$ \\
5 & 2：$2 \to 1$ & $\{6\}$ \\
6 & 2：$1 \to 0$，放入 & $\{2\}$ \\
2 & 4：$1 \to 0$，放入 & $\{4\}$ \\
4 & & $\{\}$ \\
\bottomrule
\end{tabular}
\end{center}

第三步拿出 $5$ 時，$2$ 還在等 $6$，不能排；所以 $5$、$6$ 都排完才輪到 $2$。

同樣的模型也出現在別的地方：好幾個部門各自交了一份「先做這個、再做那個」的
清單，要合成一份總排程；幾份資料各自記錄了部分的時間先後，要拼出一條完整的
時間線，並發現互相矛盾的紀錄。

\textbf{程式}

\begin{example}[label=ex:rankings]{cpp}{(rankings.cpp) 合併多份名次表}
#include <functional>
#include <iostream>
#include <queue>
#include <vector>
using namespace std;

int main() {
    ios::sync_with_stdio(false);
    cin.tie(nullptr);
    int n, q;
    cin >> n >> q;
    vector<vector<int>> adj(n + 1);
    vector<int> indeg(n + 1, 0);
    for (int i = 0; i < q; i++) {
        int k;
        cin >> k;
        int prev = 0;                       // 這份表的上一位，0 表示還沒有
        for (int j = 0; j < k; j++) {
            int x;
            cin >> x;
            if (prev != 0) {                // 只建相鄰兩人的邊
                adj[prev].push_back(x);
                indeg[x]++;
            }
            prev = x;
        }
    }
    priority_queue<int, vector<int>, greater<int>> pq;   // 小的在上
    for (int v = 1; v <= n; v++) {
        if (indeg[v] == 0) {
            pq.push(v);
        }
    }
    vector<int> order;
    while (!pq.empty()) {
        int u = pq.top();                   // 能排的人裡編號最小的
        pq.pop();
        order.push_back(u);
        for (int v : adj[u]) {
            indeg[v]--;
            if (indeg[v] == 0) {
                pq.push(v);
            }
        }
    }
    if ((int)order.size() < n) {            // 有人永遠排不進去：矛盾
        cout << -1 << '\n';
        return 0;
    }
    for (int i = 0; i < n; i++) {
        cout << order[i] << (i + 1 < n ? ' ' : '\n');
    }
    return 0;
}
\end{example}

\iosample{
6 3 \\
3 5 2 4 \\
2 6 2 \\
2 3 1 \\
}{
3 1 5 6 2 4 \\
}

\begin{itemize}
\item 選手編號從 $1$ 開始，所以用 \texttt{prev = 0} 表示「這份表還沒讀到人」。
每份表開始時都要重設，否則上一份表的最後一個人會和這一份的第一個人連起來。
\item 和第~\ref{sec:ex-checklist}~小節的例題相比，只差在容器：那裡是
\texttt{queue} 加上每批排序，這裡是 \texttt{priority\_queue}。
\end{itemize}

\textbf{複雜度}

建圖 $O(n + L)$。 每個人進出優先佇列各一次，每次 $O(\log n)$；每條邊讓入度減
一次。 總共 $O((n + L) \log n)$。 最大的測試資料在本書測試的環境下約 $0.2$ 秒。

\subsection{常見錯誤}

\begin{itemize}
\item \textbf{每一對都建邊}：答案是對的，但長的表會讓邊數變成 $k^2$ 的量級，
見上面「想法」的實測。
\item \textbf{用一般佇列}：\texttt{priority\_queue} 換成 \texttt{queue}，範例
的輸入印出 \texttt{3 5 6 1 2 4}。 它是合法的總排名，但不是字典序最小：一開始
佇列裡的 $5$、$6$ 排在 $3$ 放進來的 $1$ 前面。
\item \textbf{只對串列去掉重複的邊}：把重複的邊從 \texttt{adj} 去掉，入度卻照
樣每次加一。 輸入 \texttt{3 2}、\texttt{2 1 2}、\texttt{2 1 2}（兩份表都是
$1, 2$），正確答案是 \texttt{1 2 3}，這樣寫會印出 \texttt{-1}：$2$ 的入度是
$2$，卻只會被減一次。 要去重就兩邊一起去，不然就都不要去。
\item \textbf{沒有檢查人數}：少了最後的 \texttt{order.size()} 檢查，直接印出
\texttt{order}。 範例~2（第一行 \texttt{4 3}，三份表 $(2, 3, 4)$、$(2, 3)$、
$(4, 2)$）應該印 \texttt{-1}，這樣寫只印出 \texttt{1}：$2$、$3$、$4$ 在環上，
一個都排不進去。
\end{itemize}

\section{本章常見錯誤}
\label{sec:dag-pitfalls}

\begin{center}
\begin{tabular}{p{0.22\linewidth}p{0.44\linewidth}p{0.22\linewidth}}
\toprule
錯誤 & 例子 & 見 \\
\midrule
以為編號就是順序 & 照編號 $1$ 到 $n$ 做 DP；照編號倒著彙整（第 1 本第~6~章） &
第~\ref{sec:ex-nightmarket}~節 \\
容器選錯 & 題目要先進先出卻用優先佇列；要字典序最小卻用佇列 &
第~\ref{sec:ex-checklist}、\ref{sec:ex-rankings}~節 \\
放入佇列的時機 & 入度一減就放入，同一個點放好幾次 &
第~\ref{sec:ex-checklist}~節 \\
沒有檢查環 & 排出的點不到 $n$ 個卻照樣輸出；DFS 只用兩種狀態，把會合當成環 &
第~\ref{sec:kahn-cycle}、\ref{sec:dfs-topo}、\ref{sec:ex-rankings}~節 \\
走不到的點 & 走不到的點答案設成 $0$，被當成一條總和 $0$ 的路 &
第~\ref{sec:ex-nightmarket}~節 \\
重複的邊 & 串列去重、入度沒有去重 & 第~\ref{sec:ex-rankings}~節 \\
\bottomrule
\end{tabular}
\end{center}

另外兩個每一章都會遇到的錯誤，這一章也有：總和超過 \texttt{int}（夜市的總和
到 $2 \times 10^{14}$），以及遞迴太深（DFS 在長鏈上）。

\section{本章練習}
\label{sec:ch6-exercises}

\begin{drills}[label=drills:ch6]
\item 下面兩張圖各有幾個拓樸順序？ 全部列出來，或說明怎麼算。
\begin{enumerate}
\item 五個點，邊是 $1 \to 3$、$2 \to 3$、$3 \to 4$、$3 \to 5$。
\item 四個點，邊是 $1 \to 2$、$2 \to 3$、$3 \to 4$、$1 \to 4$。 拿掉
$1 \to 4$ 之後呢？
\end{enumerate}

\item 在圖~\ref{fig:hotpot} 上跑 Kahn 演算法，但把佇列換成堆疊（後進先出）：
一開始依編號放入 $3$、$4$、$6$，每次拿最後放進去的，新放入的點照串列順序放。
寫出每一步拿出的點和堆疊的內容，驗證排出來是 $6, 4, 3, 5, 7, 2, 1$。

\item 下面的程式想判斷有向圖有沒有環。 它錯在哪裡？ 給一組沒有環、它卻印出
\texttt{yes} 的輸入，並說明怎麼改。
\begin{code}{cpp}{練習 3}
bool hasCycle = false;
void dfs(int u) {
    seen[u] = true;
    for (int v : adj[u]) {
        if (seen[v]) {
            hasCycle = true;
        } else {
            dfs(v);
        }
    }
}
// main 裡：每個還沒看過的點都 dfs 一次，最後印出 yes 或 no
\end{code}

\item 估計下列做法的時間，$n = m = 2 \times 10^5$。
\begin{enumerate}
\item 第~\ref{sec:kahn}~節開頭的做法：每一輪掃過所有點，找一個入度 $0$ 而且還
沒排的點。 最壞情況大約要檢查幾個點？ 什麼樣的圖會達到？
\item 範例~\ref{ex:checklist} 裡 \texttt{sort} 的總時間，最壞是什麼量級？
\item 範例~\ref{ex:rankings} 如果一份表有 $k$ 人、每一對都建邊，會建幾條邊？
$k = 2 \times 10^5$ 時大約多少？
\end{enumerate}

\item 圖有三個點、一條邊 $3 \to 1$。
\begin{enumerate}
\item 用優先佇列（每次拿最小的）的 Kahn 排出什麼？
\item 把邊反過來變成 $1 \to 3$，用每次拿\textbf{最大}的優先佇列排一次，再把結
果倒過來，得到什麼？ 這個順序有什麼性質？
\item 再加一條邊 $2 \to 1$（原圖是 $3 \to 1$、$2 \to 1$），兩種做法的答案各是
什麼？ 會不會一樣？
\end{enumerate}

\item 第~\ref{sec:critical-path}~小節的火鍋：
\begin{enumerate}
\item 洗菜切菜從 $30$ 分鐘變成 $45$ 分鐘，開飯時間會變嗎？ 變成 $55$ 分鐘呢？
關鍵路徑變成哪一條？
\item 家裡有人可以幫忙，讓某一件事縮短 $10$ 分鐘（開火上桌本身只要 $5$ 分鐘，
不算在內）。 縮短哪一件，開飯時間會提早？ 提早多少？ 如果能縮短 $30$ 分鐘呢？
\end{enumerate}
\end{drills}

\section{實作練習}
\label{sec:ch6-practice}

下面是 Aki-Judge 上和本章相關的題目，依小節分組，本章的三題例題也列在裡面。
難度分四階：暖身是直接套用本章的寫法；基本要自己把題目轉成圖，或在 Kahn 上加
一點變化；實戰的題敘比較長，要先把題目轉換成本章的形狀；挑戰會和其他章的技
巧合在一起。 這一章的題目還在增加中，之後會再補。

\begin{center}
\begin{tabular}{llc}
\toprule
小節 & 題目 & 難度 \\
\midrule
Kahn 與找環 & \akiref{setup-checklist-for-the-exhibition}{成果展布置的工作清單} & 暖身（例題） \\
 & \akiref{achievements-that-never-unlock}{永遠解不開的成就} & 暖身 \\
 & \akiref{rebuild-rounds-after-an-edit}{改檔後的重新編譯輪次} & 基本 \\
 & \akiref{threshold-skill-tree}{門檻式技能樹} & 實戰 \\
 & \akiref{filing-ballots-in-batches}{投票分批整理} & 挑戰 \\
\midrule
DFS 與環 & \akiref{the-first-contradictory-rule}{第一條自相矛盾的排程規則} & 實戰 \\
\midrule
DAG 上的動態規劃 & \akiref{counting-routes-through-the-school-fair}{園遊會攤位的走法數} & 暖身 \\
 & \akiref{best-night-market-route-and-how-many-ways}{夜市路線的最高滿意度與方案數} & 基本（例題） \\
 & \akiref{orienteering-passphrase}{定向越野的通關密語} & 實戰 \\
 & \akiref{stamp-collecting-without-going-downhill}{不往下走的登山收集} & 實戰 \\
\midrule
字典序最小 & \akiref{merging-partial-rankings}{合併多份名次表} & 基本（例題） \\
 & \akiref{back-to-back-presentations}{連著上台的報告順序} & 實戰 \\
\bottomrule
\end{tabular}
\end{center}

\section{讀完你應該會}
\label{sec:ch6-checklist}

\begin{itemize}
\item 看到「誰要在誰之前」「要等前面的做完」的題目，能把它畫成有向圖，說出點和
邊各代表什麼，並檢查一個給定的順序是不是拓樸順序。
\item 能寫出 Kahn 演算法，用排出的點數判斷有沒有環，並依題目的要求選佇列、每批
排序或優先佇列。
\item 能用三種狀態的 DFS 排拓樸順序或印出一個環，並在 $n$ 很大時改成自己的堆
疊。
\item 能照拓樸順序做 DP，求最長路、最佳路線的數量和關鍵路徑，並正確處理走不到
的點。
\item 能用優先佇列求字典序最小的拓樸順序，說得出為什麼每一步選最小的不會錯。
\end{itemize}

\textbf{延伸閱讀}（書末會整理完整的參考資料）

\begin{itemize}
\item Antti Laaksonen，\textit{Competitive Programmer's Handbook}，第 16 章前
半：用 DFS 做拓樸排序與找環，以及在 DAG 上數路徑，篇幅很短。
\item Steven Halim 等，\textit{Competitive Programming 4}，Book 1，4.2.6 節：
DFS 和 Kahn 兩種寫法並列，並提到把佇列換成優先佇列。
\item Thomas H. Cormen 等，\textit{Introduction to Algorithms}，第 4 版，20.4
節（DFS 拓樸排序的證明）與 22.2 節（DAG 上的最短路，以及專案排程的關鍵路徑）。
\item 劉汝佳，《算法競賽入門經典》第 2 版，6.4.3 節（拓撲排序）與 9.2 節（DAG
上的動態規劃）。
\end{itemize}

\addcontentsline{toc}{section}{練習解答}
\section*{練習解答}

\begin{enumerate}
\item
\begin{enumerate}
\item $4$ 個。 $3$ 一定在 $1$、$2$ 之後、$4$、$5$ 之前，所以順序是「$1$、$2$
某種排法，$3$，$4$、$5$ 某種排法」，$2 \times 2 = 4$：$1\,2\,3\,4\,5$、
$1\,2\,3\,5\,4$、$2\,1\,3\,4\,5$、$2\,1\,3\,5\,4$。
\item 只有 $1$ 個，$1, 2, 3, 4$：$1 \to 2 \to 3 \to 4$ 已經把四個點的先後全部
定死。 拿掉 $1 \to 4$ 之後仍然只有這一個，因為 $1$ 在 $4$ 之前可以從另外三
條邊推出來。 這也是範例~\ref{ex:rankings} 只建相鄰邊的理由。
\end{enumerate}

\item 堆疊的頂端寫在右邊。
\begin{center}
\begin{tabular}{cll}
\toprule
拿出 & 入度減一 & 之後的堆疊 \\
\midrule
 & & 3 4 6 \\
6 & 5：$2 \to 1$ & 3 4 \\
4 & 1：$4 \to 3$ & 3 \\
3 & 2：$1 \to 0$，7：$1 \to 0$，5：$1 \to 0$ & 2 7 5 \\
5 & 1：$3 \to 2$ & 2 7 \\
7 & 1：$2 \to 1$ & 2 \\
2 & 1：$1 \to 0$ & 1 \\
1 & & （空）\\
\bottomrule
\end{tabular}
\end{center}
順序是 $6, 4, 3, 5, 7, 2, 1$。 第 3 步拿出 $3$ 時，三個點依串列順序 $2, 7, 5$
放入，最後放的 $5$ 最先拿出。 這個順序剛好和第~\ref{sec:dfs-topo}~節 DFS 排出
來的相同，但這只是巧合，換一張圖通常不同。

\item 它把「看過」都當成環，沒有分辨「還在堆疊上」和「已經結束」。 輸入
$1 \to 2$、$1 \to 3$、$2 \to 4$、$3 \to 4$ 沒有環，但第二次走到 $4$ 時
\texttt{seen[4]} 已經是 \texttt{true}，印出 \texttt{yes}。 改法：用三種狀態，
進入時設成 $1$，\texttt{for} 迴圈結束後設成 $2$，只有 \texttt{color[v] == 1}
才是環。

\item
\begin{enumerate}
\item 共 $n$ 輪，每輪最多掃 $n$ 個點，大約 $n^2 = 4 \times 10^{10}$ 次。 一條
長鏈 $n \to n - 1 \to \cdots \to 1$ 就會達到：每一輪能排的點只有一個，而且是
還沒排的點裡編號最大的，要掃到最後才找到。
\item 每個點只會在某一批裡被排序一次，所以所有批次的大小加起來不超過 $n$。 一
批大小 $b$ 的排序是 $O(b \log b) \le O(b \log n)$，總和是 $O(n \log n)$。
\item $k(k - 1)/2$ 條。 $k = 2 \times 10^5$ 時約 $2 \times 10^{10}$ 條，每條
邊至少要存一個 \texttt{int}，光是記憶體就要 $80$ GB。
\end{enumerate}

\item
\begin{enumerate}
\item 一開始入度 $0$ 的是 $2$、$3$，拿 $2$；再拿 $3$，$1$ 變成可以排。 答案
$2, 3, 1$。
\item 反過來的圖只有 $1 \to 3$。 一開始入度 $0$ 的是 $1$、$2$，拿最大的 $2$；
再拿 $1$，$3$ 可以排了，拿 $3$。 得到 $2, 1, 3$，倒過來是 $3, 1, 2$。 它是原
圖的拓樸順序，而且 $1$ 號排在所有拓樸順序裡最前面的位置（$1$ 號一定要在 $3$
後面，最早只能排第二）。
\item 原圖 $3 \to 1$、$2 \to 1$。 (a) 的做法：$2, 3, 1$。 (b) 的做法：反圖
$1 \to 3$、$1 \to 2$，先拿 $1$，再拿 $3$、$2$，得到 $1, 3, 2$，倒過來是
$2, 3, 1$。 這次一樣：$1$ 號一定是最後一個，兩種做法都把剩下的 $2$、$3$ 小的排
前面。 只有在「讓 $1$ 號提早」和「讓第一個數最小」互相衝突時，兩種答案才不同，
(a)、(b) 就是這種情況。
\end{enumerate}

\item
\begin{enumerate}
\item 洗菜切菜的寬裕時間是 $20$ 分鐘。 變成 $45$ 分鐘只多 $15$，還在寬裕時間
內，開飯仍是 $95$ 分鐘。 變成 $55$ 分鐘時，它在 $40 + 55 = 95$ 分鐘才完成，
開火上桌 $95$ 開始、$100$ 完成，開飯延到 $100$ 分鐘。 這時最長的是 買菜
$\to$ 洗菜切菜 $\to$ 開火上桌（$40 + 55 + 5$），它變成新的關鍵路徑。
\item 從起點到終點的路徑有五條：經過煮湯底的兩條是 $95$（買菜起頭）和 $75$
（借電磁爐起頭），買菜 $\to$ 洗菜切菜 $\to$ 開火上桌是 $75$，買菜 $\to$
調沾醬 $\to$ 開火上桌是 $55$，擺碗筷 $\to$ 開火上桌是 $10$。 開飯時間是最長
那條的長度。 縮短買菜或煮湯底 $10$ 分鐘，$95$ 那條變成 $85$，其他路徑都不超
過 $85$，開飯提早 $10$ 分鐘。 縮短其他事都不在 $95$ 那條上，沒有幫助。 如果
縮短 $30$ 分鐘，$95$ 那條變成 $65$，但另外有一條 $75$ 的路徑沒有變短（縮短買
菜時是借電磁爐那條，縮短煮湯底時是洗菜切菜那條），開飯只提早到 $75$ 分鐘，也
就是只提早 $20$ 分鐘。 關鍵路徑縮短到一個程度，別的路徑就會變成關鍵路徑。
\end{enumerate}
\end{enumerate}

\chapter{並查集、環與二分圖}
\label{ch:dsu}

% 本章的圖共用的樣式（點、強調的點、第二焦點、無向邊、有向邊）
\tikzset{
    dsv/.style={circle, draw=black!60, thick, minimum size=0.72cm, inner sep=0pt, font=\small},
    dsvhl/.style={dsv, draw=memStroke, fill=memFill},
    dsvalt/.style={dsv, draw=optStroke, fill=optFill},
    dsvblue/.style={dsv, draw=coreStroke, fill=coreFill},
    dse/.style={thick, draw=black!60},
    dsehl/.style={very thick, draw=memStroke},
    dsebad/.style={very thick, draw=optStroke, dashed},
    dsd/.style={-{Stealth[length=2.2mm]}, thick, draw=black!60, shorten >=1pt},
    dsdhl/.style={-{Stealth[length=2.2mm]}, very thick, draw=memStroke, shorten >=1pt},
    dslbl/.style={font=\scriptsize, text=black!60},
}

颱風過後，山區六個村子之間的道路全部中斷。 工程隊每天修好一條路，修的順序
畫在下圖，邊上的數字是第幾天修好。 每修好一條，鄉公所都想立刻知道：現在全區
分成幾塊彼此不通的區域？ 往下讀之前，先照順序一天一天算，並且留意第 5 天和
第 7 天修好的路有什麼不一樣。

\begin{center}
\begin{tikzpicture}[x=1.6cm, y=1.2cm]
\node[dsv] (1) at (0,1) {1};
\node[dsv] (2) at (1,1) {2};
\node[dsv] (3) at (2,1) {3};
\node[dsv] (4) at (1,0) {4};
\node[dsv] (5) at (2,0) {5};
\node[dsv] (6) at (3,0.5) {6};
\draw[dse] (1) -- node[above, dslbl] {1} (2);
\draw[dse] (3) -- node[right, dslbl, pos=0.4] {2} (4);
\draw[dse] (2) -- node[above, dslbl] {3} (3);
\draw[dse] (5) -- node[below right, dslbl] {4} (6);
\draw[dsehl] (1) -- node[below left, dslbl] {5} (4);
\draw[dse] (4) -- node[below, dslbl] {6} (5);
\draw[dsehl] (6) -- node[above right, dslbl] {7} (3);
\end{tikzpicture}
\diagcaption{六個村子和七條道路，邊上的數字是修好的日子。 橘色的兩條修好時，
兩頭的村子早就可以繞路互通。}
\label{fig:dsu-villages}
\end{center}

前四天每修好一條，區域就少一塊：$6$ 塊變 $5$、$4$、$3$、$2$。 第 5 天修好 1--4
時，1 和 4 早就可以經過 2、3 繞過去，區域數不變，這條路只是讓圖上多了一個繞
圈的路線。 第 6 天剩下一塊，第 7 天的 6--3 也是繞圈的路。

六個村子用手算很快。 換成 $2 \times 10^5$ 個村子、$2 \times 10^5$ 條路，每修
一條就用《C++ 演算法解題 1：APCS 實作題核心》（以下簡稱第 1 本）第~11~章的 BFS 重新數一次連通元件，每次 $O(n + m)$，總共大約
$4 \times 10^{10}$ 步，來不及。 這一章的第一個工具\textbf{並查集}就是為這種
「邊一條一條加進來，隨時要問誰和誰連通」的情況設計的，每次操作幾乎是常數時間。

上面的問題還藏著第二件事：第 5 天和第 7 天的路讓圖出現了\textbf{環}。 這一章
接著講怎麼在無向圖和有向圖上找環，然後看兩種形狀特別的圖：能不能把所有點分成
兩邊、每條邊都跨過兩邊（二分圖）；以及每個點剛好只有一條往外的箭頭（函數圖），
它和第 1 本第~4~章 4.6 節的週期偵測是同一件事。

這一章用到第 1 本第~11~章的 BFS、DFS 和連通元件，以及第 1 本 11.5 節遞迴太深時改用自己的堆
疊；第 1 本第~13~章「樹剛好有 $n - 1$ 條邊」；第 1 本第~4~章 4.6 節的 $\rho$ 形狀與週期。
這些不再重講，用到時會指回去。 語法上只需要 \texttt{vector}、\texttt{pair}
和 \texttt{queue}。

\section{並查集}
\label{sec:dsu}

先把問題講清楚。 一開始有 $n$ 個點，每個點自己一群。 接下來有一連串的操作，
每個操作是下面兩種之一：

\begin{itemize}
\item \textbf{合併}：把 $a$ 所在的群和 $b$ 所在的群併成一群。
\item \textbf{查詢}：問 $a$ 和 $b$ 是不是在同一群。
\end{itemize}

群只會合併、不會拆開。 能做這兩件事的資料結構叫\textbf{並查集}（disjoint set
union，也叫 union-find）。 「並」是合併，「查」是查詢在哪一群。

\subsection{每群選一個代表}

做法是替每一群選一個\textbf{代表}（representative）。 兩個點在同一群，剛好就
是它們的代表相同。 要讓「找代表」和「合併」都很快，並查集把每一群存成一棵有
根樹，根就是代表。 每個點只記一個數 \texttt{parent[x]}：它在樹上的上一層。 根
的上一層是自己。

\begin{itemize}
\item \textbf{找代表}（\texttt{find(x)}）：從 $x$ 沿著 \texttt{parent} 一直往上
走，走到 \texttt{parent[r] == r} 的 $r$，就是代表。
\item \textbf{合併}（\texttt{unite(a, b)}）：先找出兩邊的代表；相同就表示早已在
同一群，什麼都不做；不同就把其中一個根接到另一個根底下。
\end{itemize}

這裡的樹和第 1 本第~13~章的有根樹一樣只記父節點，但它只是程式內部的記錄方式，和題目
裡的邊沒有關係。 例如圖~\ref{fig:dsu-villages} 裡 1 和 4 有一條路，在並查集的
樹裡 1 卻不一定接在 4 底下。 唯一要維持的是：同一群的點都在同一棵樹上。

\begin{code}{cpp}{最簡單的並查集}
vector<int> parent(n + 1);
for (int i = 1; i <= n; i++) {
    parent[i] = i;                     // 一開始每個點自己一群
}

int find(int x) {
    while (parent[x] != x) {           // 往上走到根
        x = parent[x];
    }
    return x;
}

void unite(int a, int b) {
    a = find(a);
    b = find(b);
    if (a != b) {
        parent[a] = b;                 // a 那棵樹接到 b 底下
    }
}
\end{code}

\subsection{樹會變高}

這個版本是對的，但可能很慢。 照 \texttt{unite(1, 2)}、\texttt{unite(2, 3)}、
\ldots、\texttt{unite(n - 1, n)} 的順序合併：每次 \texttt{a} 本身就是根，
\texttt{parent[a] = b} 讓舊的根接到新的點底下，最後變成一條
$1 \to 2 \to 3 \to \cdots \to n$ 的長鏈。 之後每次 \texttt{find(1)} 都要走
$n - 1$ 步，$n$ 次查詢就是 $O(n^2)$。

問題出在樹的高度。 有兩個互相獨立的改法，一般兩個一起用。

\textbf{改法一：按大小合併}（union by size）。 每個根多記一個數，這一群有幾個
點。 合併時把\textbf{小的}那群的根接到大的那群的根底下。 一個點的深度要加一，
只會發生在它所在的群被接到另一個至少一樣大的群底下，所以每加一層，它所在的
群至少變成兩倍大。 $n$ 個點最多翻倍 $\log_2 n$ 次，任何點的深度都不超過
$\log_2 n$，$n = 2 \times 10^5$ 時是 $17$。

\textbf{改法二：路徑壓縮}（path compression）。 \texttt{find(x)} 找到根之後，
把剛才路上經過的每個點都直接改接到根底下。 樹裡誰接在誰下面本來就不重要，只要
同一群在同一棵樹，所以可以隨便改。 這一次多做了一點事，之後這些點再查代表，
一步就到。 下圖是第 5 天的 \texttt{find(1)}：1 原本要經過 2 才到根 4，壓縮之
後直接接在 4 底下。

\begin{center}
\begin{tikzpicture}[x=1cm, y=1cm]
\begin{scope}
\node[font=\small] at (1,2.7) {壓縮前};
\node[dsv] (4) at (1,2) {4};
\node[dsvhl] (2) at (0.4,1) {2};
\node[dsv] (3) at (1.6,1) {3};
\node[dsvhl] (1) at (0.4,0) {1};
\draw[dsd] (2) -- (4);
\draw[dsd] (3) -- (4);
\draw[dsd] (1) -- (2);
\node[dsv] (6) at (3,2) {6};
\node[dsv] (5) at (3,1) {5};
\draw[dsd] (5) -- (6);
\end{scope}
\begin{scope}[xshift=6cm]
\node[font=\small] at (1,2.7) {壓縮後};
\node[dsv] (4) at (1,2) {4};
\node[dsv] (2) at (0.2,1) {2};
\node[dsv] (3) at (1.8,1) {3};
\node[dsvhl] (1) at (1,1) {1};
\draw[dsd] (2) -- (4);
\draw[dsd] (3) -- (4);
\draw[dsdhl] (1) -- (4);
\node[dsv] (6) at (3,2) {6};
\node[dsv] (5) at (3,1) {5};
\draw[dsd] (5) -- (6);
\end{scope}
\end{tikzpicture}
\diagcaption{第 4 天結束時並查集裡的兩棵樹（箭頭指向 \texttt{parent}），以及第
5 天 \texttt{find(1)} 之後。 兩群是 $\{1, 2, 3, 4\}$ 和 $\{5, 6\}$，代表是 4
和 6。}
\label{fig:dsu-compress}
\end{center}

兩個改法各自能省多少？ 用 $8192$ 個點實際數一次，比較兩種合併順序：一種是上面的
長鏈；另一種是兩兩合併，$1$ 和 $2$、$3$ 和 $4$、\ldots，再把兩個兩人組併成四人組，
像淘汰賽一樣。 合併完之後，每個點各查 $5$ 次代表。 兩種順序都呼叫了五萬七千多次
\texttt{find}，下表是總共沿著 \texttt{parent} 走了幾步：

\begin{center}
\begin{tabular}{lrrrr}
\toprule
合併順序 & 不優化 & 只按大小合併 & 只做路徑壓縮 & 兩個都用 \\
\midrule
長鏈 & 167751680 & 49144 & 49145 & 49144 \\
兩兩合併 & 282596 & 282596 & 61406 & 61406 \\
\bottomrule
\end{tabular}
\end{center}

\begin{itemize}
\item \textbf{長鏈}：不優化要走一億六千多萬步。 只要用上任何一個改法，就降到
五萬步左右，平均每次 \texttt{find} 不到一步。
\item \textbf{兩兩合併}：每次合併的兩群一樣大，按大小合併幫不上忙，樹高是
$\log_2 8192 = 13$，平均每次 \texttt{find} 走大約五步。 加上路徑壓縮，第一輪
查詢就把樹壓平了，後面四輪每次只走一步。
\end{itemize}

兩個改法一起用時，可以證明 $m$ 次操作的總時間是 $O(m \, \alpha(n))$。
$\alpha$ 叫\textbf{反阿克曼函數}（inverse Ackermann function），長得非常慢：
$n$ 就算大到宇宙裡原子的個數，$\alpha(n)$ 也不超過 $4$。 證明很長，這本書不寫；
寫程式時把每次操作當成常數時間就好。 直覺上可以這樣想：按大小合併讓樹一開始
就不高，路徑壓縮讓走過的路以後不必再走第二次，貴的那幾次查詢會被之後大量便宜
的查詢平均掉，和第 1 本第~6~章的攤銷分析是同一種數法。

\subsection{用並查集維護連通元件}

回到章首的修路問題。 下面的程式就是本章之後都會用的並查集寫法：

\begin{itemize}
\item \texttt{find} 寫成遞迴，\texttt{return parent[x] = find(parent[x]);} 先找
到根，把根存進 \texttt{parent[x]} 再回傳。 遞迴一層一層回來時，路上每個點都被
改接到根，這就是路徑壓縮。 因為有按大小合併，樹高不超過 $17$，遞迴不會太深。
\item \texttt{unite} 回傳有沒有真的合併。 回傳 \texttt{true}：兩群併成一群，
區域數減一。 回傳 \texttt{false}：兩頭早就連通，區域數不變；這條路和原本就有的
某條路線合起來，繞成一個環，下一節會再用到。
\end{itemize}

\begin{example}[label=ex:repair-roads]{cpp}{(roads.cpp) 修好每條路之後還有幾塊區域}
#include <iostream>
#include <utility>
#include <vector>
using namespace std;

vector<int> parent, sz;

int find(int x) {
    if (parent[x] == x) {
        return x;
    }
    return parent[x] = find(parent[x]);
}

// 合併 a、b 所在的兩群；本來就在同一群時回傳 false
bool unite(int a, int b) {
    a = find(a);
    b = find(b);
    if (a == b) {
        return false;
    }
    if (sz[a] > sz[b]) {
        swap(a, b);
    }
    parent[a] = b;
    sz[b] += sz[a];
    return true;
}

int main() {
    int n, m;
    cin >> n >> m;
    parent.resize(n + 1);
    sz.assign(n + 1, 1);
    for (int i = 1; i <= n; i++) {
        parent[i] = i;
    }
    int regions = n;                       // 一開始每一區自成一塊
    for (int i = 0; i < m; i++) {
        int a, b;
        cin >> a >> b;
        if (unite(a, b)) {
            regions--;
            cout << regions << '\n';
        } else {
            cout << regions << " loop\n";  // 兩端早就連通：這條路繞成一圈
        }
    }
    return 0;
}
\end{example}

\iosample{
6 7 \\
1 2 \\
3 4 \\
2 3 \\
5 6 \\
1 4 \\
4 5 \\
6 3 \\
}{
5 \\
4 \\
3 \\
2 \\
2 loop \\
1 \\
1 loop \\
}

逐天看 \texttt{parent} 的變化。 兩群一樣大時，程式把 \texttt{a} 那一群接到
\texttt{b} 底下。

\begin{center}
\begin{tabular}{clll}
\toprule
天 & 路 & 發生的事 & \texttt{parent[1..6]} \\
\midrule
0 & --- & --- & 1 2 3 4 5 6 \\
1 & 1--2 & 根 1、2 一樣大，1 接到 2 & 2 2 3 4 5 6 \\
2 & 3--4 & 3 接到 4 & 2 2 4 4 5 6 \\
3 & 2--3 & 根是 2 和 4，一樣大，2 接到 4 & 2 4 4 4 5 6 \\
4 & 5--6 & 5 接到 6 & 2 4 4 4 6 6 \\
5 & 1--4 & \texttt{find(1)} 走 1、2、4，順便把 1 接到 4；根相同 & 4 4 4 4 6 6 \\
6 & 4--5 & 根是 4（4 人）和 6（2 人），6 接到 4 & 4 4 4 4 6 4 \\
7 & 6--3 & 兩邊的根都是 4 & 4 4 4 4 6 4 \\
\bottomrule
\end{tabular}
\end{center}

\begin{apcsnote}
並查集的程式很短，不需要背類別或模板。 考試時直接寫範例~\ref{ex:repair-roads} 的
\texttt{find}、\texttt{unite} 和初始化的迴圈，大約二十行。 只做其中一個改法，
大多數題目也夠快；兩個都寫最保險。
\end{apcsnote}

很多題目不只問「同不同群」，還要知道每一群的人數、總和、最大值這類資料。 做法
都一樣：這些資料只存在根上，合併時把小的那群的資料加到大的那群的根，查詢時先
\texttt{find} 找到根再讀。 \texttt{sz} 就是這樣維護的，下面的例題再多存一個
總和。

\subsection{什麼時候用、什麼時候不能用}

\begin{itemize}
\item 適合：邊一條一條加進來，中間夾著「同不同群」「這群多大」的詢問；或者只
需要知道連通元件，不必建鄰接串列，邊讀邊 \texttt{unite} 就好（第 1 本第~11~章 11.2
節說過，並查集只需要一條一條讀邊）。
\item 不能直接用：會\textbf{刪邊}的題目。 並查集只會合併，拆不開。 如果所有操
作事先都知道，有時可以倒過來處理，把「刪邊」變成「加邊」，這需要另外的觀察，
本章練習第~2~題會用到這個想法。
\item 不能用：要知道兩點之間的路怎麼走、最短要幾步。 並查集只知道在不在同一
群，不知道路徑，這時用第 1 本第~11~章的 BFS。
\end{itemize}

\subsection{例題：社團合併與總積分}
\label{sec:ex-club}

\textbf{題意}（\akiref{club-mergers-and-total-points}{社團合併與總積分}）

\begin{itemize}
\item 有 $n$ 位學生，學生 $i$ 有積分 $p_i$。 一開始每人自己一個社團。
\item 依序有 $q$ 個操作。 \texttt{1 a b}：把 $a$ 和 $b$ 所在的社團合併；已經
同社團（包含 $a = b$）時輸出 \texttt{same}，否則輸出合併後的人數。
\texttt{2 a}：輸出 $a$ 所在社團的人數和積分總和。
\item $n, q \le 2 \times 10^5$，$p_i \le 10^9$。
\end{itemize}

例如積分是 $3, 1, 4, 1, 5$。 合併 1、2 後輸出 \texttt{2}；查學生 2 輸出
\texttt{2 4}；合併 3、4 輸出 \texttt{2}；再合併 2、1，兩人已經同社團，輸出
\texttt{same}；合併 4、1 得到 $\{1, 2, 3, 4\}$，輸出 \texttt{4}；查學生 5 輸出
\texttt{1 5}；查學生 3 輸出 \texttt{4 9}。

\textbf{想法}

最直接的做法是開一個陣列記每個人在哪個社團。 合併時掃過所有學生，把其中一個
社團的編號全部改成另一個；查詢時掃過所有學生，數出同社團的人數和積分。 每個操
作 $O(n)$，總共 $O(nq) = 4 \times 10^{10}$，來不及。

這正是並查集的問題：社團只合併、不拆開。 人數和積分總和照上一小節的說法，只存
在根上：

\begin{itemize}
\item \texttt{members[r]}、\texttt{total[r]}：以 $r$ 為根的社團的人數與積分總和。
一開始 \texttt{members[i] = 1}、\texttt{total[i] = p\_i}。
\item 合併時找出兩個根，相同就輸出 \texttt{same}；不同就把小的接到大的底下，人
數和積分都加到大的那個根上。
\item 查詢時先 \texttt{find(a)}，讀根上的兩個數。
\end{itemize}

積分總和最大是 $2 \times 10^5 \times 10^9 = 2 \times 10^{14}$，要用
\texttt{long long}。

同樣的模型也出現在別的地方：公司合併時，要隨時知道每個集團旗下有幾家公司、總
營收多少；倉庫把可以互相替代的料號合成一類，要知道每一類的總庫存。

\textbf{程式}

\begin{example}[label=ex:club]{cpp}{(club.cpp) 社團合併與總積分}
#include <iostream>
#include <utility>
#include <vector>
using namespace std;

vector<int> parent, members;          // members、total 只在根上有意義
vector<long long> total;

int find(int x) {
    if (parent[x] == x) {
        return x;
    }
    return parent[x] = find(parent[x]);   // 路徑壓縮
}

int main() {
    ios::sync_with_stdio(false);
    cin.tie(nullptr);
    int n, q;
    cin >> n >> q;
    parent.resize(n + 1);
    members.assign(n + 1, 1);
    total.assign(n + 1, 0);
    for (int i = 1; i <= n; i++) {
        parent[i] = i;
        cin >> total[i];
    }
    while (q--) {
        int type;
        cin >> type;
        if (type == 1) {
            int a, b;
            cin >> a >> b;
            a = find(a);
            b = find(b);
            if (a == b) {
                cout << "same\n";
                continue;
            }
            if (members[a] > members[b]) {    // 小的社團接到大的底下
                swap(a, b);
            }
            parent[a] = b;
            members[b] += members[a];
            total[b] += total[a];
            cout << members[b] << '\n';
        } else {
            int a;
            cin >> a;
            a = find(a);
            cout << members[a] << ' ' << total[a] << '\n';
        }
    }
    return 0;
}
\end{example}

\iosample{
5 7 \\
3 1 4 1 5 \\
1 1 2 \\
2 2 \\
1 3 4 \\
1 2 1 \\
1 4 1 \\
2 5 \\
2 3 \\
}{
2 \\
2 4 \\
2 \\
same \\
4 \\
1 5 \\
4 9 \\
}

\texttt{members} 同時是按大小合併用的大小，不用另外開 \texttt{sz}。 範例最後查學
生 3：\texttt{find(3)} 走 $3 \to 4 \to 2$，讀到根 2 上的 \texttt{4 9}。

\textbf{複雜度}

每個操作幾乎是常數時間，總共 $O((n + q)\,\alpha(n))$。 最大的測試資料在本書測
試的環境下約 $0.15$ 秒。

\subsection{常見錯誤}

\begin{itemize}
\item \textbf{積分總和用 \texttt{int}}：把 \texttt{total} 改成
\texttt{vector<int>}。 輸入三位積分都是 $10^9$ 的學生，全部合併後查詢，正確答
案是 \texttt{3 3000000000}，在本書測試的環境下印出 \texttt{3 -1294967296}。
\item \textbf{沒有先判斷已經同社團}：拿掉 \texttt{if (a == b)} 那一段。 同一群
合併時 \texttt{parent[a] = b} 只是把根接到自己，看不出問題，但人數和積分會再
加一次。 上面那組輸入裡，第一個操作 \texttt{1 2 2} 應該印 \texttt{same}，這樣
寫會印 \texttt{2}；最後的查詢印出 \texttt{8 8000000000}。
\item \textbf{查詢時忘了 \texttt{find}}：直接讀 \texttt{members[a]}。 範例最
後一行應該是 \texttt{4 9}，這樣寫印出 \texttt{1 4}，因為學生 3 已經不是根，它
身上的數字停在合併前。
\end{itemize}

\subsection{預告：帶權並查集與奇偶性}
\label{sec:dsu-parity}

有一類題目只有兩隊，給的關係有兩種：「$a$、$b$ 同隊」和「$a$、$b$ 不同隊」。 例如：「1 和 2 不同隊」「2 和 3 不同隊」，可以推出 1 和
3 同隊；如果接著又說「1 和 3 不同隊」，就自相矛盾了。

並查集可以處理這種關係：每個點除了 \texttt{parent}，再多記一個 $0$ 或 $1$，表
示「我和 \texttt{parent} 同隊（$0$）還是不同隊（$1$）」。 沿著樹往上走到根時，
把路上的這些數用互斥或（\texttt{\^}）累加起來，就知道自己和根同隊還是不同隊。
兩個點在同一群時，比較它們各自和根的關係，就能回答「同隊、不同隊」或發現矛盾。
這種在 \texttt{parent} 旁邊多記一個「和上一層的關係」的寫法叫\textbf{帶權並查
集}（weighted union-find），路徑壓縮時這個數也要跟著更新，細節比一般的並查集
多一些。

完整的寫法留到《C++ 演算法解題 3》。 如果所有關係一次給完、不必邊加邊問，這類題目也可以用第
~\ref{sec:bipartite}~節的二分圖上色來做，實作練習裡有兩題可以用兩種方法比較。

\section{無向圖找環}
\label{sec:cycle-undirected}

無向圖上的\textbf{環}（cycle）是一條從某個點出發、每條邊最多走一次、最後回到
出發點的路線。 沒有重邊和自環的圖，環至少有三個點，例如圖~\ref{fig:dsu-villages}
的 $1 - 2 - 3 - 4 - 1$。 有些題目允許重邊，這時兩點之間的兩條邊就能繞成一個環；
自環本身也是一個環。 讀題時先確認題目怎麼算。

下面三種方法都能判斷一張無向圖有沒有環，各有適合的時候。

\textbf{方法一：並查集。} 依序加邊，\texttt{unite} 回傳 \texttt{false} 的那條
邊，兩端早就連通，加上去就形成環。 範例~\ref{ex:repair-roads} 的 \texttt{loop} 就是這
樣標出來的。 這個方法順便回答了「加到第幾條邊時第一次出現環」，邊是一條一條給
的題目最好用。 重邊也會被正確抓到：第二條重邊的兩端已經連通。 自環
\texttt{unite(a, a)} 也會回傳 \texttt{false}。

\textbf{方法二：數點和邊。} 第 1 本第~13~章說過，$c$ 個點的連通圖至少要 $c - 1$ 條
邊，剛好 $c - 1$ 條時是一棵樹，沒有環。 所以先找出每個連通元件，數它有幾個點
$c$、幾條邊 $e$：$e \ge c$ 就有環，$e = c - 1$ 就沒有。 不用真的找出環在哪裡，
寫起來最短。 整張圖要拿掉幾條邊才會沒有環，也能這樣算：每個元件要拿掉
$e - (c - 1)$ 條。

\textbf{方法三：DFS。} 從一個點出發做 DFS，走到 $u$ 時，如果 $u$ 有一個鄰居 $v$
已經走過，而且 $v$ 不是剛才走來 $u$ 的那個點，就找到了環：從 $v$ 沿著 DFS 走到
$u$，再走 $u$--$v$ 這條邊回到 $v$。 要把環上的點印出來，就記下每個點是從哪個點
走來的（\texttt{from}），找到時從 $u$ 沿著 \texttt{from} 往回走到 $v$。

「不是剛才走來的那個點」這個條件要小心。 如果圖有重邊，而題目把兩條重邊算成
一個環，就要改成「不是剛才走來的那\textbf{條邊}」，所以要替每條邊編號：

\begin{code}{cpp}{DFS 找無向圖的環（判斷用）}
// adj[u] 存 (鄰居, 邊的編號)；cameBy[u] 是走到 u 用的那條邊，起點設成 -1
bool foundCycle = false;

void dfs(int u) {
    seen[u] = true;
    for (auto [v, id] : adj[u]) {
        if (id == cameBy[u]) {
            continue;                  // 剛才走來的那條邊，不算
        }
        if (seen[v]) {
            foundCycle = true;         // 從別的邊又碰到走過的點
        } else {
            cameBy[v] = id;
            dfs(v);
        }
    }
}
\end{code}

為什麼碰到走過的點就一定有環？ $u$ 和 $v$ 都已經走過，表示 DFS 早就沿著別的邊
把它們連起來了，再加上 $u$--$v$ 這條沒用過的邊，就繞成一圈。 而且第一次發生這
種情況時，$v$ 一定還沒結束，是 $u$ 在呼叫堆疊裡的祖先：一條無向邊兩端都看得
到，如果 $v$ 先結束，它在結束前就會從自己那一端看到 $u$--$v$。 所以在第一次發
現時停下來，從 $u$ 沿著 \texttt{from} 往回走，一定走得到 $v$。 下一節會看到，有
向圖沒有這個性質，碰到走過的點不代表有環。

和第 1 本第~11~章一樣，點數到 $10^5$ 以上時，遞迴可能太深，改用第 1 本 11.5 節的改法二。

三種方法怎麼選：邊一條一條加、要知道從第幾條邊開始有環，用並查集；只要知道每個
元件有沒有環，數點和邊；要印出環上的點，用 DFS。

\section{有向圖找環}
\label{sec:cycle-directed}

有向圖的環要順著箭頭走：從某個點出發，沿著箭頭走回出發點。 第 1 本第~9~章和本書第~6~
章的有向無環圖，就是沒有這種環的有向圖。 工作排程、課程先修這類「$a$ 要在 $b$
之前」的關係，出現環就表示互相等待，永遠排不出順序。

上一節的 DFS 搬到有向圖上會出錯。 看下圖左邊：$1 \to 2$、$1 \to 3$、
$3 \to 2$，順著箭頭走不出任何環。 從 1 出發的 DFS 先走 2，2 沒有往外的箭頭，結
束；回到 1 再走 3，3 的鄰居 2 已經走過了。 照無向圖的規則，這時會誤判成有環。

\begin{center}
\begin{tikzpicture}[x=1.2cm, y=1.1cm]
\begin{scope}
\node[font=\small] at (1,2) {沒有環};
\node[dsv] (1) at (0,1) {1};
\node[dsv] (2) at (2,1) {2};
\node[dsv] (3) at (1,0) {3};
\draw[dsd] (1) -- (2);
\draw[dsd] (1) -- (3);
\draw[dsdhl] (3) -- (2);
\end{scope}
\begin{scope}[xshift=4.5cm]
\node[font=\small] at (1.5,2) {有環};
\node[dsv] (1) at (0,1) {1};
\node[dsvhl] (2) at (1,1) {2};
\node[dsvhl] (4) at (2,1) {4};
\node[dsvhl] (5) at (2,0) {5};
\node[dsvhl] (3) at (1,0) {3};
\draw[dsd] (1) -- (2);
\draw[dsd] (1) -- (3);
\draw[dsdhl] (3) -- (2);
\draw[dsdhl] (2) -- (4);
\draw[dsdhl] (4) -- (5);
\draw[dsdhl] (5) -- (3);
\end{scope}
\end{tikzpicture}
\diagcaption{左：DFS 從 3 碰到已經結束的 2，但沒有環。 右：多了 $2 \to 4$、
$4 \to 5$、$5 \to 3$，橘色的四條邊繞成環 $2 \to 4 \to 5 \to 3 \to 2$。}
\label{fig:directed-cycle}
\end{center}

差別在於：左圖的 2 已經\textbf{結束}了，DFS 從 2 出發能走的都走完了，也沒有回
到 1 或 3。 碰到一個已經結束的點，只表示兩條路在那裡會合，不表示繞了回來。 真
正的環，是碰到一個\textbf{還沒結束}的點，也就是目前這條路上的祖先：從它走到
$u$，再走 $u \to v$ 回到它。

所以每個點要分三種狀態，不能只記「去過沒有」：

\begin{itemize}
\item $0$：還沒去過。
\item $1$：去過了，但還沒結束，還在目前這條路上（第 1 本第~11~章說的呼叫堆疊裡）。
\item $2$：已經結束，從它出發能走的都走完了。
\end{itemize}

DFS 進入一個點時把它設成 $1$，離開時設成 $2$。 看到一條 $u \to v$：$v$ 是 $0$
就走過去；$v$ 是 $1$ 就找到環；$v$ 是 $2$ 就跳過。

要在離開一個點的那一刻做事，就需要第 1 本第~11~章 11.5 節的改法二：堆疊裡每一格記住
點和「下一個要看的鄰居」，鄰居看完了才彈出。 這個堆疊剛好就是目前這條路，所以
找到環的時候，堆疊裡從 $v$ 到頂端的那一段就是環上的點，不用另外記
\texttt{from}。 \texttt{pos[v]} 記 $v$ 在堆疊裡的位置。

\begin{example}[label=ex:directed-cycle]{cpp}{(directed\_cycle.cpp) 找出有向圖的一個環}
#include <iostream>
#include <utility>
#include <vector>
using namespace std;

int main() {
    int n, m;
    cin >> n >> m;
    vector<vector<int>> adj(n + 1);
    for (int i = 0; i < m; i++) {
        int u, v;
        cin >> u >> v;
        adj[u].push_back(v);
    }
    // state：0 還沒去過，1 在目前的路上，2 已經結束
    vector<int> state(n + 1, 0), pos(n + 1, -1);
    vector<pair<int, int>> st;             // （點, 下一個要看的鄰居位置）
    for (int s = 1; s <= n; s++) {
        if (state[s] != 0) {
            continue;
        }
        state[s] = 1;
        pos[s] = 0;
        st.push_back({s, 0});
        while (!st.empty()) {
            int u = st.back().first;
            int i = st.back().second;
            if (i == (int)adj[u].size()) {     // u 的鄰居都看完了
                state[u] = 2;
                st.pop_back();
                continue;
            }
            st.back().second++;
            int v = adj[u][i];
            if (state[v] == 0) {
                state[v] = 1;
                pos[v] = st.size();            // v 在堆疊裡的位置
                st.push_back({v, 0});
            } else if (state[v] == 1) {        // v 還在路上：找到環
                cout << "cycle:";
                for (int k = pos[v]; k < (int)st.size(); k++) {
                    cout << ' ' << st[k].first;
                }
                cout << ' ' << v << '\n';
                return 0;
            }
        }
    }
    cout << "no cycle\n";
    return 0;
}
\end{example}

\iosample{
5 6 \\
1 2 \\
1 3 \\
3 2 \\
2 4 \\
4 5 \\
5 3 \\
}{
cycle: 2 4 5 3 2 \\
}

輸入是圖~\ref{fig:directed-cycle} 右邊。 下表每一列是「進入一個新的點」或「發現
環」的那一刻：

\begin{center}
\begin{tabular}{clll}
\toprule
步 & 看到的邊 & 發生的事 & 堆疊（底在左） \\
\midrule
1 & --- & 進入 1 & 1 \\
2 & $1 \to 2$ & 2 是 $0$，進入 & 1 2 \\
3 & $2 \to 4$ & 4 是 $0$，進入 & 1 2 4 \\
4 & $4 \to 5$ & 5 是 $0$，進入 & 1 2 4 5 \\
5 & $5 \to 3$ & 3 是 $0$，進入 & 1 2 4 5 3 \\
6 & $3 \to 2$ & 2 是 $1$：找到環 & 1 2 4 5 3 \\
\bottomrule
\end{tabular}
\end{center}

2 在堆疊的第 1 格（從 $0$ 算），印出第 1 格到頂端的 2、4、5、3，再補上 2。 如果
拿掉 $5 \to 3$，DFS 會把每個點都設成 $2$，印出 \texttt{no cycle}。

\begin{itemize}
\item 外層迴圈從每個還沒去過的點出發，因為從 1 出發不一定走得到所有點，環可能
在別的地方。
\item 自環 $u \to u$ 也會被抓到：看到這條邊時 $u$ 自己的狀態就是 $1$。
\item 每個點進出堆疊各一次，每條邊看一次，總共 $O(n + m)$。
\end{itemize}

另一個判斷有向圖有沒有環的方法，是第~6~章的拓樸排序：每次拿掉一個沒有箭頭指
進來的點，如果拿到最後還剩下點，剩下的點就卡在某個環裡（或排在環的後面）。 只
問「有沒有環」時，兩種都可以；要印出環上的點，用三種狀態的 DFS 比較直接。

\section{二分圖}
\label{sec:bipartite}

期末考要把 $n$ 個科目排在星期一和星期二兩天考完。 有些科目有共同的考生，不能
排在同一天。 能不能排得出來？ 把科目畫成點，有共同考生的兩科之間連一條邊，問題
變成：能不能把每個點塗成兩種顏色之一，讓每條邊的兩端顏色都不同？

能這樣塗的無向圖叫\textbf{二分圖}（bipartite graph）：點可以分成兩邊，每條邊都
跨過兩邊，同一邊的點之間沒有邊。 分兩隊、分兩班、兩個櫃子、兩種角色，題目只要
問「能不能分成兩組、某些兩兩不能同組」，就是在問二分圖。

\subsection{用 BFS 上色}

一個點塗了某種顏色之後，它的鄰居全都只能塗另一種，沒有選擇。 所以從任何一個
點開始，先隨便塗一種顏色，接下來整個連通元件的顏色都被決定了。 用 BFS 一層一層
塗：

\begin{itemize}
\item 起點塗 $0$，放進佇列。
\item 從佇列拿出 $u$，看它的每個鄰居 $v$：$v$ 還沒塗色，就塗成和 $u$ 相反的顏
色、放進佇列；$v$ 已經塗了，而且和 $u$ 同色，就是\textbf{衝突}，這個元件不是
二分圖。
\end{itemize}

BFS 從起點出發，距離 $d$ 的點塗的是 $d \bmod 2$：偶數層一種顏色，奇數層另一種。
所以 \texttt{color[v]} 其實就是第 1 本第~11~章的 \texttt{dist[v]} 只留奇偶。 用 DFS 塗
也可以，每走一步換一次顏色；點數大時一樣要注意遞迴深度。

下圖左邊從 1 開始塗。 1 塗藍；2、3 塗粉紅；4 塗藍；5 塗粉紅；6 塗藍。 每條邊兩
端都不同色，它是二分圖。 右邊多了一條 3--5。 BFS 處理 3 的時候，5 還沒上色，就從
3 塗成藍色，距離變成 $2$，和 4 同一層。 接著處理 4，看到鄰居 5 也是藍色，衝突。

\begin{center}
\begin{tikzpicture}[x=1.2cm, y=1.1cm]
\begin{scope}
\node[dsvblue] (1) at (0,1) {1};
\node[dsvalt] (2) at (1,2) {2};
\node[dsvalt] (3) at (1,0) {3};
\node[dsvblue] (4) at (2,1) {4};
\node[dsvalt] (5) at (3,1) {5};
\node[dsvblue] (6) at (4,1) {6};
\draw[dse] (1) -- (2);
\draw[dse] (1) -- (3);
\draw[dse] (2) -- (4);
\draw[dse] (3) -- (4);
\draw[dse] (4) -- (5);
\draw[dse] (5) -- (6);
\foreach \x/\d in {0/0, 1/1, 2/2, 3/3, 4/4} {\node[dslbl] at (\x,-0.7) {$d = \d$};}
\end{scope}
\begin{scope}[xshift=6.3cm]
\node[dsvblue] (1) at (0,1) {1};
\node[dsvalt] (2) at (1,2) {2};
\node[dsvalt] (3) at (1,0) {3};
\node[dsvblue] (4) at (2,1) {4};
\node[dsvblue] (5) at (3,1) {5};
\node[dsvalt] (6) at (4,1) {6};
\draw[dse] (1) -- (2);
\draw[dse] (1) -- (3);
\draw[dse] (2) -- (4);
\draw[dse] (3) -- (4);
\draw[dsebad] (4) -- (5);
\draw[dse] (5) -- (6);
\draw[dse] (3) to[bend right=25] (5);
\end{scope}
\end{tikzpicture}
\diagcaption{BFS 上色，藍色是 $0$、粉紅是 $1$，$d$ 是左圖到點 1 的距離。 右圖多
了 3--5，虛線 4--5 兩端同色；$3 - 4 - 5 - 3$ 是一個三條邊的環。}
\label{fig:bipartite}
\end{center}

右圖的衝突不是上色的順序不好。 3--4--5--3 是一個三條邊的環，繞一圈要換三次顏
色，換奇數次一定回不到原來的顏色，不管怎麼塗都不行。 一般的結論是：

\begin{center}
一張無向圖是二分圖，剛好就是它沒有長度為奇數的環（\textbf{奇環}，odd cycle）。
\end{center}

有奇環就不可能塗，理由就是上面那一句。 反過來，BFS 上色時一旦出現衝突，一定找
得到一個奇環，所以「BFS 沒衝突」和「沒有奇環」是同一件事。 這一半的理由放在下
面的框裡。

\begin{whybox}[BFS 的衝突一定來自奇環]
設 BFS 從 $s$ 出發，看到一條邊 $u$--$v$ 兩端同色。 同色表示 $d(u)$ 和 $d(v)$
同奇偶。 從 $u$ 和 $v$ 各自沿著「是從哪個點走來的」一路往回走，最後都會回到
$s$；設兩條路第一次會合在點 $w$。 走 $w \to \cdots \to u$，經過 $u$--$v$，再走
$v \to \cdots \to w$，是一個環，長度是
\[
\bigl(d(u) - d(w)\bigr) + 1 + \bigl(d(v) - d(w)\bigr)
= d(u) + d(v) - 2d(w) + 1 .
\]
$d(u) + d(v)$ 是偶數，減掉偶數 $2d(w)$ 還是偶數，再加 $1$ 是奇數。 這兩條往回
的路在 $w$ 之前沒有共同的點，所以這真的是一個環。 圖~\ref{fig:bipartite} 右邊：
$u = 4$、$v = 5$，往回走是 $4 \to 2 \to 1$ 和 $5 \to 3 \to 1$，會合在 $1$，得到
長度 $2 + 1 + 2 = 5$ 的奇環 $1 - 2 - 4 - 5 - 3 - 1$。 它和 $3 - 4 - 5 - 3$ 不
同，但一樣是奇數長。
\end{whybox}

\subsection{好幾個連通元件}

圖不一定連通。 每個連通元件各自上色，互不影響，所以外層用第 1 本第~11~章數連通元件的
迴圈，從每個還沒塗色的點開始一次 BFS。 整張圖是二分圖，就是每個元件都是。

一個元件只要能塗，剛好有兩種塗法：起點塗 $0$ 或塗 $1$，其他點跟著決定，兩種塗
法互為顏色對調。 所以有 $c$ 個元件的二分圖一共有 $2^c$ 種塗法。 題目要求答案唯
一時，通常會規定「每個元件編號最小的點塗某一種」，從編號小的點開始掃，每個元
件第一個被掃到的就是它最小的點，正好符合。

\subsection{例題：不相容藥品分櫃}
\label{sec:ex-cabinets}

\textbf{題意}（\akiref{cabinets-for-incompatible-chemicals}{不相容藥品分櫃}）

\begin{itemize}
\item 實驗室有 $n$ 種藥品，有 $m$ 組不相容的配對，配對的兩種不能放在同一個櫃
子。 能經由配對連起來的藥品屬於同一個群組。
\item 每個群組分開處理：能只用 A、B 兩櫃放好，就這樣放，而且群組裡編號最小的
放 A；不管怎麼放都做不到的群組，整組放進隔離櫃 C。
\item 輸出長度 $n$ 的字串，第 $i$ 個字元是藥品 $i$ 的櫃子。 $n, m \le 2 \times
10^5$，配對可能重複出現。
\end{itemize}

例如 $n = 6$，配對是 1--2、3--4、4--5、5--3、2--6。 群組 $\{3, 4, 5\}$ 三種兩兩
不相容，兩個櫃子放不下，全部放 C；群組 $\{1, 2, 6\}$ 放成 1 在 A、2 在 B、6 在 A。
答案是 \texttt{ABCCCA}。

\textbf{想法}

藥品是點、不相容配對是邊，群組就是連通元件，「只用 A、B 放好」就是這個元件是
二分圖。 最直接的做法是每個群組試遍所有放法，群組有 $k$ 種藥品就有 $2^k$ 種，
$k$ 到幾十就來不及。 上一小節已經說過，一個元件只要起點的顏色定了，其他都跟著
決定，不用試。

做法照上一小節：編號由小到大掃，遇到還沒上色的 $s$，它就是新群組裡編號最小的，
從它開始 BFS，$s$ 塗 A。 BFS 時順便把走到的點記在 \texttt{group} 裡，同時記下
有沒有出現衝突。 整個群組走完之後，有衝突就全部標 C，沒有就照顏色標 A、B。

衝突出現時\textbf{不能停下來}。 如果停在一半，群組裡還沒走到的藥品，之後會被
外層迴圈當成新群組的起點，照常放 A、B。 題目要求整組放 C，所以一定要把整個群
組走完。

同樣的模型也出現在別的地方：值班表要把同事分成早晚兩班，有些人不能同班；一場
辯論賽要把報名者分成正反兩方，同一個研究室的人不能同一方。 分不成兩班的那一群
另外處理。

\textbf{程式}

\begin{example}[label=ex:cabinets]{cpp}{(cabinets.cpp) 不相容藥品分櫃}
#include <iostream>
#include <string>
#include <vector>
using namespace std;

int main() {
    ios::sync_with_stdio(false);
    cin.tie(nullptr);
    int n, m;
    cin >> n >> m;
    vector<vector<int>> adj(n + 1);
    for (int i = 0; i < m; i++) {
        int u, v;
        cin >> u >> v;
        adj[u].push_back(v);
        adj[v].push_back(u);
    }
    vector<int> color(n + 1, -1);          // -1：還沒上色；0：A 櫃；1：B 櫃
    string ans(n, '?');
    vector<int> group;                     // 這一群的成員，兼當 BFS 的佇列
    for (int s = 1; s <= n; s++) {
        if (color[s] != -1) {
            continue;
        }
        group.assign(1, s);                // s 是這一群編號最小的藥品
        color[s] = 0;
        bool conflict = false;
        for (int head = 0; head < (int)group.size(); head++) {
            int u = group[head];
            for (int v : adj[u]) {
                if (color[v] == -1) {
                    color[v] = 1 - color[u];
                    group.push_back(v);
                } else if (color[v] == color[u]) {
                    conflict = true;       // 記下來，但這一群還是要走完
                }
            }
        }
        for (int u : group) {
            ans[u - 1] = conflict ? 'C' : (color[u] == 0 ? 'A' : 'B');
        }
    }
    cout << ans << '\n';
    return 0;
}
\end{example}

\iosample{
6 5 \\
1 2 \\
3 4 \\
4 5 \\
5 3 \\
2 6 \\
}{
ABCCCA \\
}

\begin{itemize}
\item \texttt{group} 同時是佇列和成員名單：\texttt{head} 往後走就是從佇列拿出
來，BFS 結束時 \texttt{group} 裡剛好是整個群組。 用 \texttt{queue} 的話，拿出
來的點就不見了，還要另外存一份名單。
\item 範例的 $s = 3$ 時，3 塗 A，4、5 塗 B；處理 4 時看到鄰居 5 也是 B，衝突。
程式繼續把群組走完，三個都標 C。
\end{itemize}

\textbf{複雜度}

每個點進出 \texttt{group} 一次，每條邊從兩端各看一次，$O(n + m)$。 最大的測試
資料在本書測試的環境下約 $0.15$ 秒，大部分花在讀輸入。

\subsection{常見錯誤}

\begin{itemize}
\item \textbf{衝突就停下 BFS}：在 \texttt{for (int head ...)} 的條件加上
\texttt{\&\& !conflict}。 輸入
\begin{code}{text}{停在一半會錯的輸入}
4 4
1 2
2 3
3 1
3 4
\end{code}
正確答案是 \texttt{CCCC}，這樣寫會印 \texttt{CCCA}：處理 2 時就發現 2、3 同色，
停下來時 4 還沒走到，之後被當成新群組放進 A。
\item \textbf{起點不是群組裡編號最小的}：例如從輸入第一條配對的端點開始 BFS。
輸入 \texttt{7 5}、\texttt{2 1}、\texttt{3 2}、\texttt{6 5}、\texttt{5 4}、
\texttt{1 2} 時，正確答案是 \texttt{ABAABAA}，從 2 開始塗會印
\texttt{BABABAA}。
\item \textbf{整張圖一起判斷}：只要有一個衝突就全部放 C。 上面範例的群組
$\{1, 2, 6\}$ 會被拖累，印出 \texttt{CCCCCC}。
\item \textbf{遞迴的 DFS 上色}：測試資料裡有一條將近 $2 \times 10^5$ 個點的長
鏈，在本書測試的環境下 \texttt{Segmentation fault}。
\end{itemize}

\subsection{什麼時候用、什麼時候不能用}

\begin{itemize}
\item 適合：分成兩組，某些兩兩不能同組（或必須不同組）；問「能不能分」、「一共
有幾種分法」、「每個點分到哪邊」。
\item 不能用：分成\textbf{三組}以上。 三種顏色的上色問題，目前沒有人知道有效率
的演算法，題目的點數通常很小，只能用第 1 本第~4~章的回溯。 看到「三組」時先看資料範圍。
\end{itemize}

\section{函數圖}
\label{sec:functional}

第 1 本第~4~章 4.6 節的週期偵測裡，每個狀態的下一步只由它自己決定。 把每個狀態畫成一
個點、從它畫一個箭頭指到下一個狀態，得到的有向圖有一個特色：每個點\textbf{剛好
只有一條}往外的箭頭。 這種圖叫\textbf{函數圖}（functional graph），因為它就是
一個函數 $f$：從 $x$ 出發的箭頭指向 $f(x)$。

第 1 本 4.6 節只從一個起點出發，一步一步走，畫出來是一條尾巴接上一個圈的 $\rho$ 形狀。
這一節換個角度，看整張函數圖長什麼樣子，以及所有點一起處理的方法。 題目常常直
接給 $f$：每個人只把消息傳給一個固定的人、每個房間的傳送門只通到一個房間、每
個格子上的箭頭只指向一格。

\subsection{每一塊剛好一個圈}

下圖是 $n = 10$、$f = (5, 3, 4, 2, 4, 7, 6, 7, 8, 10)$ 的函數圖，也就是
$f(1) = 5$、$f(2) = 3$、\ldots。

\begin{center}
\begin{tikzpicture}[x=1.15cm, y=1.1cm]
\node[dsv] (1) at (0,1) {1};
\node[dsv] (5) at (1,1) {5};
\node[dsvhl] (4) at (2,1) {4};
\node[dsvhl] (2) at (3,1.7) {2};
\node[dsvhl] (3) at (3,0.3) {3};
\draw[dsd] (1) -- (5);
\draw[dsd] (5) -- (4);
\draw[dsdhl] (4) -- (2);
\draw[dsdhl] (2) -- (3);
\draw[dsdhl] (3) -- (4);
\node[dsv] (9) at (4.7,1) {9};
\node[dsv] (8) at (5.7,1) {8};
\node[dsvhl] (7) at (6.7,1) {7};
\node[dsvhl] (6) at (7.9,1) {6};
\draw[dsd] (9) -- (8);
\draw[dsd] (8) -- (7);
\draw[dsdhl] (7) to[bend left=30] (6);
\draw[dsdhl] (6) to[bend left=30] (7);
\node[dsvhl] (10) at (9.3,1) {10};
\draw[dsdhl] (10) to[out=60, in=120, looseness=6] (10);
\end{tikzpicture}
\diagcaption{一張函數圖。 橘色的點和箭頭是圈，其他點都在通往圈的尾巴上。 10 指
向自己，是長度 $1$ 的圈。}
\label{fig:functional}
\end{center}

這張圖分成三塊，每一塊剛好有一個圈，其他點都順著箭頭流進那個圈。 這不是巧合：

\begin{itemize}
\item \textbf{每一塊至少一個圈}：從任何點出發一直走，點只有有限個，總會走到走
過的點，從那裡開始就繞圈。 這就是第 1 本 4.6 節的 $\rho$。
\item \textbf{每一塊至多一個圈}：不看箭頭方向，把一塊當成無向圖。 它有 $c$ 個
點，每個點剛好一條往外的邊，所以剛好 $c$ 條邊。 照第~\ref{sec:cycle-undirected}~
節的方法二，$c$ 個點的連通圖只要 $c - 1$ 條邊就是一棵樹，多出來的那一條邊只能
造出一個環。
\end{itemize}

所以函數圖的每一塊都像一個 $\rho$ 的放大版：一個圈，加上幾棵樹，樹上的箭頭全部
指向圈。 從任何點出發走，先走完一段\textbf{尾巴}，進入它那一塊的圈，之後就一直
繞。 圖~\ref{fig:functional} 裡，從 1 出發走 $1 \to 5 \to 4$，尾巴長 $2$，接著在
長度 $3$ 的圈 $4 \to 2 \to 3 \to 4$ 上繞。

\subsection{一次算出每個點的尾巴和圈}

題目常常要每個點的資料：從 $x$ 出發，走幾步進入圈（尾巴長 \texttt{tail[x]}），
那個圈多長（\texttt{cyc[x]}）。 每個點各自照第 1 本 4.6 節的方法走一次，最壞是
$O(n^2)$：$n$ 個點排成一條長尾巴時，每個點都要走過整條尾巴。

走過的點不必再走。 這和第~\ref{sec:cycle-directed}~節的三種狀態是同一個想法，只
是每個點只有一條出邊，DFS 就簡化成一直往前走。 從還沒走過的 $s$ 出發，沿路把點
設成 $1$ 並記在 \texttt{path} 裡，直到碰見狀態不是 $0$ 的點 $x$。 $x$ 是 $1$，表
示撞上這一趟自己的路，\texttt{path} 從 $x$ 開始的那一段是新的圈；$x$ 是 $2$，表
示接上以前算好的點。 最後把還沒算好的點從靠近圈的一端往回算：
\texttt{tail[u] = tail[f[u]] + 1}、\texttt{cyc[u] = cyc[f[u]]}。

\begin{example}[label=ex:functional]{cpp}{(functional.cpp) 每個點的尾巴長和圈長}
#include <iostream>
#include <vector>
using namespace std;

int main() {
    int n;
    cin >> n;
    vector<int> f(n + 1);
    for (int i = 1; i <= n; i++) {
        cin >> f[i];
    }
    // state：0 還沒走過，1 在這一趟的路上，2 已經算好
    vector<int> state(n + 1, 0), tail(n + 1), cyc(n + 1), pos(n + 1);
    vector<int> path;
    for (int s = 1; s <= n; s++) {
        if (state[s] != 0) {
            continue;
        }
        path.clear();
        int x = s;
        while (state[x] == 0) {            // 一直走到碰見走過的點
            state[x] = 1;
            pos[x] = path.size();
            path.push_back(x);
            x = f[x];
        }
        int end = path.size();             // path[0 .. end-1] 還沒算好
        if (state[x] == 1) {               // 撞上這一趟自己：新的圈
            int len = end - pos[x];
            for (int i = pos[x]; i < end; i++) {
                tail[path[i]] = 0;
                cyc[path[i]] = len;
                state[path[i]] = 2;
            }
            end = pos[x];
        }
        for (int i = end - 1; i >= 0; i--) {   // 尾巴上的點，從靠近圈的往回算
            int u = path[i];
            tail[u] = tail[f[u]] + 1;
            cyc[u] = cyc[f[u]];
            state[u] = 2;
        }
    }
    cout << "tail:";
    for (int i = 1; i <= n; i++) {
        cout << ' ' << tail[i];
    }
    cout << "\ncycle:";
    for (int i = 1; i <= n; i++) {
        cout << ' ' << cyc[i];
    }
    cout << '\n';
    return 0;
}
\end{example}

\iosample{
10 \\
5 3 4 2 4 7 6 7 8 10 \\
}{
tail: 2 0 0 0 1 0 0 1 2 0 \\
cycle: 3 3 3 3 3 2 2 2 2 1 \\
}

外層迴圈每一輪的情形：

\begin{center}
\begin{tabular}{cllll}
\toprule
$s$ & \texttt{path} & 停在 & 狀態 & 結果 \\
\midrule
1 & 1 5 4 2 3 & 4 & $1$ & 圈 4、2、3，長 $3$；5 的尾巴 $1$，1 的尾巴 $2$ \\
6 & 6 7 & 6 & $1$ & 圈 6、7，長 $2$ \\
8 & 8 & 7 & $2$ & 8 的尾巴 $1$ \\
9 & 9 & 8 & $2$ & 9 的尾巴 $2$ \\
10 & 10 & 10 & $1$ & 圈 10，長 $1$ \\
\bottomrule
\end{tabular}
\end{center}

$s = 2, 3, 4, 5, 7$ 時已經算好，直接跳過。 每個點只進入 \texttt{path} 一次，總共 $O(n)$。

\subsection{第 $k$ 步會走到哪裡}

從 $x$ 出發走 $k$ 步，停在哪個點？ 只有一個起點時，第 1 本 4.6 節已經講過：記下走過的
路，找到尾巴長 $p$ 和圈長 $L$，$k \ge p$ 時第 $k$ 步和第
$p + (k - p) \bmod L$ 步相同。 圖~\ref{fig:functional} 從 9 出發，尾巴長 $2$、圈
長 $2$，第 $10^{18}$ 步和第 $2 + (10^{18} - 2) \bmod 2 = 2$ 步相同，停在 7。

起點很多、每個起點要問很多次時，每次都重走一遍太慢。 一個通用的做法叫
\textbf{倍增}（binary lifting）：先算出每個點走 $1$ 步、$2$ 步、$4$ 步、
$8$ 步……會到哪裡：\texttt{up[j][x]} 是從 $x$ 走 $2^j$ 步到的點。 走 $2^j$ 步
等於先走 $2^{j-1}$ 步、再走 $2^{j-1}$ 步，所以每一層都能從上一層算出來。 要走
$k$ 步時，把 $k$ 寫成二進位，例如 $13 = 8 + 4 + 1$，就依序走 $1$ 步、$4$ 步、
$8$ 步。 這和第 1 本第~8~章快速冪把指數拆成二進位是同一個想法。

\begin{code}{cpp}{倍增：從 x 走 k 步}
const int LOG = 60;                    // 2^60 > 10^18
vector<vector<int>> up(LOG, vector<int>(n + 1));
up[0] = f;                             // 走 1 步
for (int j = 1; j < LOG; j++) {
    for (int x = 1; x <= n; x++) {
        up[j][x] = up[j - 1][up[j - 1][x]];   // 走 2^(j-1) 步兩次
    }
}

int jump(int x, long long k) {
    for (int j = 0; k > 0; j++, k >>= 1) {
        if (k & 1) {                   // k 的第 j 位是 1：走 2^j 步
            x = up[j][x];
        }
    }
    return x;
}
\end{code}

建表是 $O(n \log K)$，每次詢問 $O(\log K)$，$K$ 是 $k$ 的上限。 它不需要知道圈
在哪裡，尾巴和圈一視同仁。 代價是記憶體：$n = 2 \times 10^5$、$K = 10^{18}$ 時，
表有 $60 \times 2 \times 10^5$ 個 \texttt{int}，大約 $48$ MB，先確認題目的記憶
體限制。

\subsection{記不下的時候}

第 1 本 4.6 節最後說，狀態的種類太多、記不下每一步時，要另外想辦法。 函數圖上有一個只
用幾個變數的方法，叫\textbf{龜兔賽跑}（Floyd's cycle detection）：烏龜一次走一
步，兔子一次走兩步，從同一個起點出發。 兩者都進入圈之後，兔子每一步追近一格，
一定會追上；追上時兔子比烏龜多走的步數剛好是整數圈。 這時把烏龜放回起點，兩個
都改成一次一步，它們會剛好在圈的入口相遇，走了幾步就是尾巴長。 時間仍是
$O(\mu + \lambda)$（$\mu$、$\lambda$ 是尾巴長和圈長），記憶體只要幾個變數。
細節可以看本章的延伸閱讀；實作練習的「記不下的驗證碼紀錄」就是這種情況。

\subsection{例題：兩台電子骰子何時同號}
\label{sec:ex-dice}

\textbf{題意}（\akiref{when-do-the-two-dice-agree}{兩台電子骰子何時同號}）

\begin{itemize}
\item 第一台骰子顯示 $0$ 到 $m_1 - 1$，照對照表 $f_1$ 換數字：這一秒是 $x$，下
一秒是 $f_1(x)$。 第二台顯示 $0$ 到 $m_2 - 1$，照 $f_2$ 換。
\item 第 $0$ 秒分別顯示 $x_0$、$y_0$，之後兩台每秒同時換一次。
\item 求最小的 $t \ge 0$ 使兩台顯示的數字相同；永遠不會相同就輸出 $-1$。
\item $m_1, m_2 \le 2000$。
\end{itemize}

例如 $f_1 = (1, 0, 0, 3)$、$x_0 = 2$，$f_2 = (1, 2, 4, 0, 1)$、$y_0 = 1$。 兩台
依序顯示：

\begin{center}
\begin{tabular}{c|ccccccc}
\toprule
秒 $t$ & 0 & 1 & 2 & 3 & 4 & 5 & 6 \\
\midrule
第一台 & 2 & 0 & 1 & 0 & 1 & 0 & 1 \\
第二台 & 1 & 2 & 4 & 1 & 2 & 4 & 1 \\
\bottomrule
\end{tabular}
\end{center}

第 $6$ 秒都是 $1$，答案是 $6$。

\textbf{想法}

一秒一秒模擬，相同就輸出，這部分很直接。 難的是答案是 $-1$ 的時候：要模擬到什
麼時候才能確定永遠不會相同？

每一台骰子自己就是一張函數圖：數字是點，對照表是箭頭。 第一台從 $x_0$ 出發，走
出一個 $\rho$；第二台也是。 但只看一台不夠：上表的第一台在第 $3$ 秒就回到第 $1$
秒的 $0$，開始繞圈；第二台在第 $3$ 秒也回到第 $0$ 秒的 $1$。 兩台都在繞圈了，
同號卻發生在第 $6$ 秒。 兩個圈的長度是 $2$ 和 $3$，兩台合起來要 $6$ 秒才會一起
回到同一個組合。

所以狀態要是兩台合起來的\textbf{一對數} $(x, y)$。 下一秒的 $(x, y)$ 只由這一秒
的 $(x, y)$ 決定，所以這些數對也構成一張函數圖，有 $m_1 m_2$ 個點，最多
$4 \times 10^6$ 個。 從 $(x_0, y_0)$ 出發同樣走出一個 $\rho$：最多 $m_1 m_2$ 步
之內，一定會走到出現過的數對。 從那一刻起，之後的每一秒都和上次一模一樣。 如果
在第一次重複之前都沒有同號，之後也不會有，答案是 $-1$。

做法：開一個 $m_1 m_2$ 格的陣列，記哪些數對出現過，數對 $(x, y)$ 編成
$x \cdot m_2 + y$。 每一秒先檢查是否相同，再檢查是否出現過，再往下走一步。

同樣的模型也出現在別的地方：兩條固定班表的公車路線，問最早哪一天兩班車會停在
同一站；兩個照固定規則切換狀態的號誌，問什麼時候會同時亮紅燈。 合起來看成一對
狀態，就回到函數圖。

\textbf{程式}

\begin{example}[label=ex:dice]{cpp}{(dice.cpp) 兩台電子骰子何時同號}
#include <iostream>
#include <vector>
using namespace std;

int main() {
    ios::sync_with_stdio(false);
    cin.tie(nullptr);
    int m1, x, m2, y;
    cin >> m1 >> x;
    vector<int> f1(m1);
    for (int &v : f1) {
        cin >> v;
    }
    cin >> m2 >> y;
    vector<int> f2(m2);
    for (int &v : f2) {
        cin >> v;
    }
    // 狀態是一對 (x, y)，編成 x * m2 + y；最多 m1 * m2 種
    vector<bool> seen((long long)m1 * m2, false);
    long long t = 0;
    while (true) {
        if (x == y) {                      // 先比相等：第 0 秒就可能相同
            cout << t << '\n';
            break;
        }
        long long id = (long long)x * m2 + y;
        if (seen[id]) {                    // 這一對出現過：之後只會重演
            cout << -1 << '\n';
            break;
        }
        seen[id] = true;
        x = f1[x];
        y = f2[y];
        t++;
    }
    return 0;
}
\end{example}

\iosample{
4 2 \\
1 0 0 3 \\
5 1 \\
1 2 4 0 1 \\
}{
6 \\
}

\begin{itemize}
\item \texttt{vector<bool>} 每格只佔一個位元，$4 \times 10^6$ 格大約 $0.5$ MB。
用 \texttt{vector<int>} 也只要 $16$ MB，都放得下。
\item 數對的編號最大是 $m_1 m_2 - 1$，不會超過 \texttt{int}；寫成
\texttt{long long} 是習慣，$m$ 再大一點時不必回頭改。
\item 迴圈最多跑 $m_1 m_2 + 1$ 輪就會停，因為每一輪都標記一個新的數對，數對只有
$m_1 m_2$ 種。
\end{itemize}

\textbf{複雜度}

最多 $m_1 m_2 \le 4 \times 10^6$ 步，每步 $O(1)$。 最大的測試資料在本書測試的
環境下約 $0.02$ 秒。

\subsection{常見錯誤}

\begin{itemize}
\item \textbf{只看一台骰子}：用一台的數字當狀態，它重複就輸出 $-1$。 範例的第
一台在第 $3$ 秒重複，這樣寫會印 \texttt{-1}。
\item \textbf{憑感覺決定跑幾秒}：例如「每台最多 $m$ 種數字，跑 $m_1 + m_2$ 秒一
定夠」。 輸入
\begin{code}{text}{只跑 m1 + m2 秒會錯的輸入}
3 0
1 2 0
6 3
0 1 3 4 5 2
\end{code}
第一台在 $0, 1, 2$ 上繞長度 $3$ 的圈，第二台在 $3, 4, 5, 2$ 上繞長度 $4$ 的圈，
只有數字 $2$ 是共同的，第一次同號在第 $11$ 秒。 只跑 $m_1 + m_2 = 9$ 秒會印
\texttt{-1}。
\item \textbf{只等第 $0$ 秒的數對再出現}：起點可能在尾巴上，之後再也不會出現。
例如第一台 $f_1 = (1, 2, 0, 0)$ 從 $3$ 出發，$3$ 只出現在第 $0$ 秒。 答案是 $-1$
時程式就停不下來。
\item \textbf{先走一步再比較}：漏掉第 $0$ 秒。 $x_0 = y_0$ 時答案是 $0$，這樣寫
會從第 $1$ 秒開始比，印出別的數字或 $-1$。
\end{itemize}

\section{本章常見錯誤}
\label{sec:dsu-pitfalls}

前面各節的錯誤，歸納起來是下面幾類。

\begin{center}
\begin{tabular}{p{0.24\linewidth}p{0.5\linewidth}p{0.16\linewidth}}
\toprule
錯誤 & 例子 & 見 \\
\midrule
資料沒存在根上 & 查詢時沒有先 \texttt{find}；已經同一群還再合併一次，人數重複加 &
第~\ref{sec:ex-club}~節 \\
樹太高 & 兩個改法都沒做；只做路徑壓縮時用遞迴 \texttt{find} &
第~\ref{sec:dsu}~節 \\
有向圖只記去過沒有 & 碰到已經結束的點就當成環 & 第~\ref{sec:cycle-directed}~節 \\
不連通的圖只處理一塊 & 只從點 1 開始找環或上色 & 第~\ref{sec:cycle-directed}、
\ref{sec:ex-cabinets}~節 \\
衝突就停 & 群組沒走完，剩下的被當成新群組 & 第~\ref{sec:ex-cabinets}~節 \\
狀態不完整 & 兩台骰子只看一台；以為一定會回到起點 & 第~\ref{sec:ex-dice}~節 \\
遞迴太深 & 長鏈上遞迴的 DFS 上色、找環 & 第~\ref{sec:cycle-undirected}、
\ref{sec:ex-cabinets}~節 \\
\bottomrule
\end{tabular}
\end{center}

\section{本章練習}
\label{sec:ch7-exercises}

\begin{drills}[label=drills:ch7]
\item 用範例~\ref{ex:repair-roads} 的並查集（按大小合併，一樣大時 \texttt{a} 接到
\texttt{b} 底下，\texttt{find} 做路徑壓縮），$n = 6$，依序執行
\texttt{unite(1, 2)}、\texttt{unite(3, 4)}、\texttt{unite(5, 6)}、
\texttt{unite(1, 3)}、\texttt{unite(5, 1)}、\texttt{find(1)}。 每一步之後的
\texttt{parent[1..6]} 是什麼？

\item 一張無向圖有 $n$ 個點、$m$ 條邊，會依序\textbf{拆掉}所有的邊，每拆一條就
要回答目前有幾個連通元件。 並查集不能拆邊，怎麼用並查集解這題？（提示：從最後
一刻往回看。）

\item 一張無向圖有 $n$ 個點、$m$ 條邊、$c$ 個連通元件，沒有重邊和自環。 用並查
集依序加入這 $m$ 條邊，\texttt{unite} 會回傳幾次 \texttt{false}？ 這個數和「最
少要拿掉幾條邊，圖上才不會有環」有什麼關係？

\item 下面是有人寫的有向圖找環，\texttt{seen} 一開始全是 \texttt{false}，從每
個還沒去過的點呼叫 \texttt{dfs}。 它錯在哪裡？ 給一組會讓它答錯的輸入，並說明
怎麼改。
\begin{code}{cpp}{練習 4}
bool dfs(int u) {
    seen[u] = true;
    for (int v : adj[u]) {
        if (seen[v] || dfs(v)) {
            return true;
        }
    }
    return false;
}
\end{code}

\item 函數圖 $n = 8$，$f = (2, 3, 1, 3, 4, 7, 8, 6)$。
\begin{enumerate}
\item 畫出這張圖，寫出範例~\ref{ex:functional} 會印出的 \texttt{tail} 和
\texttt{cycle}。
\item 從點 5 出發走 $10^{18}$ 步，停在哪裡？
\item 倍增的 \texttt{up[2][5]} 是多少？
\end{enumerate}
\end{drills}

\section{實作練習}
\label{sec:ch7-practice}

下面是 Aki-Judge 上和本章相關的題目，依小節分組，難度的四階和前面各章相同。

\begin{center}
\begin{tabular}{llc}
\toprule
小節 & 題目 & 難度 \\
\midrule
並查集 & \akiref{club-mergers-and-total-points}{社團合併與總積分} & 暖身 \\
 & \akiref{two-dorm-assignment-records}{兩棟宿舍的分配紀錄} & 實戰 \\
 & \akiref{locker-block-assignment}{置物櫃區段分配} & 挑戰 \\
\midrule
找環 & \akiref{server-room-clusters-with-loops}{有繞圈線路的機房群} & 基本 \\
 & \akiref{the-first-contradictory-rule}{第一條自相矛盾的排程規則} & 實戰 \\
\midrule
二分圖 & \akiref{cabinets-for-incompatible-chemicals}{不相容藥品分櫃} & 暖身 \\
 & \akiref{same-team-or-not}{同隊與不同隊的說法} & 基本 \\
 & \akiref{the-robot-that-never-stops}{不停站的導覽機器人} & 實戰 \\
 & \akiref{balancing-two-stages}{兩個舞台的人數平衡} & 挑戰 \\
\midrule
函數圖與週期 & \akiref{the-code-word-on-day-k}{暗號字串的第 K 天} & 暖身 \\
 & \akiref{when-do-the-two-dice-agree}{兩台電子骰子何時同號} & 基本 \\
 & \akiref{the-shuttle-on-a-command-loop}{照指令表繞圈的接駁車} & 基本 \\
 & \akiref{piggy-bank-range-totals}{存錢筒的區間總額} & 實戰 \\
 & \akiref{the-k-th-water-stop}{第 K 次補水站} & 挑戰 \\
 & \akiref{the-code-log-too-long-to-keep}{記不下的驗證碼紀錄} & 挑戰 \\
\bottomrule
\end{tabular}
\end{center}

本章的三題例題也在 Aki-Judge 上，讀完各節之後可以先自己寫一次再對照。 「兩棟
宿舍的分配紀錄」和「同隊與不同隊的說法」都是同隊、不同隊的關係，可以用第
~\ref{sec:dsu-parity}~小節預告的帶權並查集，也可以先建圖再上色，兩種都試試看。

\section{讀完你應該會}
\label{sec:ch7-checklist}

\begin{itemize}
\item 看到「邊一條一條加進來、隨時問同不同群」的題目，能寫出按大小合併加路徑壓
縮的並查集，並把人數、總和這類資料存在根上。
\item 能用並查集、數點和邊、DFS 三種方法判斷無向圖有沒有環，並依題目要的是
「第幾條邊開始有環」「哪幾群有環」還是「環上有哪些點」來選。
\item 能用三種狀態的 DFS 找出有向圖的一個環，說得出為什麼只記「去過沒有」會誤
判。
\item 能用 BFS 上色判斷二分圖，處理不連通的圖，並用奇環說明一張圖為什麼不能
分成兩邊。
\item 看到每個點只有一個「下一個」的題目，能畫出函數圖，一次算出每個點的尾巴長
和圈長，並用週期或倍增求出第 $k$ 步的位置。
\end{itemize}

\textbf{延伸閱讀}（書末會整理完整的參考資料）

\begin{itemize}
\item Antti Laaksonen，\textit{Competitive Programmer's Handbook}，12.3 節與第
15、16 章：找環、二分圖、並查集、函數圖（successor graph）的倍增和龜兔賽跑。
\item Steven Halim 等，\textit{Competitive Programming 4}，Book 1，2.4.2 節與
4.2.7、4.2.8 節：並查集的另一種寫法（按高度合併），以及二分圖和有向圖找環的
整理。
\item Thomas H. Cormen 等，\textit{Introduction to Algorithms}，第 4 版，第 19
章：並查集的完整分析，包含 $\alpha(n)$ 從哪裡來。
\end{itemize}

\addcontentsline{toc}{section}{練習解答}
\section*{練習解答}

\begin{enumerate}
\item 依序是：\texttt{unite(1, 2)} 後 \texttt{2 2 3 4 5 6}；\texttt{unite(3, 4)}
後 \texttt{2 2 4 4 5 6}；\texttt{unite(5, 6)} 後 \texttt{2 2 4 4 6 6}；
\texttt{unite(1, 3)}：兩個根 2、4 都是兩人，2 接到 4，得到
\texttt{2 4 4 4 6 6}（\texttt{find(1)} 走 1、2 時，1 的 \texttt{parent} 本來就
是根 2，壓縮不改變什麼）；\texttt{unite(5, 1)}：\texttt{find(5)} 是 6（兩人），
\texttt{find(1)} 走 1、2、4，順便把 1 接到 4；6 比較小，接到 4，得到
\texttt{4 4 4 4 6 4}；最後的 \texttt{find(1)} 不改變陣列。

\item 倒過來看：最後一刻所有邊都拆光了，有 $n$ 個元件。 從最後一條被拆的邊開
始，依拆除的相反順序一條一條\textbf{加回去}，用並查集維護元件數。 加回第 $i$ 條
邊之前的元件數，就是「拆掉第 $i$ 條邊之後」的答案。 存起來，最後依原本的順序輸出。

\item $m - (n - c)$ 次。 每次回傳 \texttt{true} 元件就少一個，從 $n$ 個變成 $c$ 個，
剛好 $n - c$ 次，其他都是 \texttt{false}。 這也是最少要拿掉的邊數：每個元件留下
一棵 $c_i - 1$ 條邊的樹，共留 $n - c$ 條，其他都要拿掉。

\item 它把「已經去過」都當成環，沒有分辨 $v$ 是還在目前的路上，還是早就結束了。
輸入 \texttt{3 3}、\texttt{1 2}、\texttt{1 3}、\texttt{3 2}（圖~\ref{fig:directed-cycle}
左邊）沒有環，它從 1 走 2，回來再走 3，看到 2 已經去過，就回傳 \texttt{true}。
改法是把 \texttt{seen} 換成三種狀態：進入時設 $1$，迴圈結束、回傳
\texttt{false} 之前設 $2$；判斷條件改成 \texttt{state[v] == 1} 才算環，
\texttt{state[v] == 0} 才往下走。

\item
\begin{enumerate}
\item 圈有兩個：$1 \to 2 \to 3 \to 1$（長 $3$）和 $6 \to 7 \to 8 \to 6$（長 $3$）。
4 指向 3，5 指向 4。 輸出
\texttt{tail: 0 0 0 1 2 0 0 0} 和 \texttt{cycle: 3 3 3 3 3 3 3 3}。
\item 尾巴長 $2$、圈長 $3$，第 $10^{18}$ 步和第 $2 + (10^{18} - 2) \bmod 3 = 4$ 步
相同（$10^{18} \bmod 3 = 1$）：$5 \to 4 \to 3 \to 1 \to 2$，停在 2。
\item 從 5 走 $4$ 步，也是 2。
\end{enumerate}
\end{enumerate}

\chapter{加權圖的最短路徑}
\label{ch:sp}

% 本章的圖共用的樣式（點、已定案的點、剛處理的點、灰掉的點、邊、最短路徑樹的邊、
% 有向邊、強調的有向邊、權重、距離標籤）
\tikzset{
    spv/.style={circle, draw=black!60, thick, minimum size=0.68cm, inner sep=0pt, font=\small},
    spvdone/.style={spv, draw=memStroke, fill=memFill},
    spvnow/.style={spv, draw=optStroke, fill=optFill},
    spvgray/.style={spv, draw=black!35, text=black!45, dashed},
    spe/.style={thick, draw=black!55},
    spetree/.style={line width=2.4pt, draw=memStroke},
    spd/.style={-{Stealth[length=2.2mm]}, thick, draw=black!60, shorten >=1pt},
    spdhl/.style={-{Stealth[length=2.2mm]}, very thick, draw=optStroke, shorten >=1pt},
    spdgray/.style={-{Stealth[length=2.2mm]}, thick, draw=black!35, dashed, shorten >=1pt},
    spw/.style={font=\scriptsize, fill=white, inner sep=1pt},
    spdist/.style={font=\small\bfseries, text=coreStroke},
}

% 章首的外送地圖。#1：「點/距離/樣式/距離標籤的位置」的清單；#2：要畫粗的邊（a/b）。
\newcommand{\spmapgraph}[2]{%
    \foreach \v/\x/\y in {1/0/1, 2/1.5/2, 3/1.5/0, 4/3/2, 5/3/0, 6/4.5/1} {
        \coordinate (q\v) at (\x,\y);
    }
    \foreach \v/\d/\s/\p in {#1} {
        \node[\s, label={[spdist]\p:\d}] (n\v) at (q\v) {\v};
    }
    \foreach \a/\b in {#2} {
        \draw[spetree] (n\a) -- (n\b);
    }
    \foreach \a/\b/\w in {1/2/7, 1/3/2, 2/3/3, 2/4/2, 3/5/7, 4/5/1, 4/6/6, 5/6/2} {
        \draw[spe] (n\a) -- node[spw] {\w} (n\b);
    }
}

晚餐時段，便當店接到一張外送單。 下圖是附近的路口和道路，店在路口 $1$，客
人家在路口 $6$，道路都是雙向的，邊上的數字是騎車通過這條路要幾分鐘。 從店
到客人家最快要幾分鐘？ 往下讀之前，先在圖上找一條最快的路線，並且數一數它
經過幾條路。

\begin{center}
\begin{tikzpicture}[x=1.25cm, y=1.05cm]
\spmapgraph{1//spv/left, 2//spv/above, 3//spv/below, 4//spv/above, 5//spv/below, 6//spv/right}{}
\end{tikzpicture}
\diagcaption{外送的地圖。 店在路口 $1$，客人家在路口 $6$，數字是騎車要幾分鐘。}
\label{fig:sp-map}
\end{center}

從 $1$ 到 $6$ 最少要經過三條路，三條路的走法有兩種：$1 \to 2 \to 4 \to 6$
要 $7 + 2 + 6 = 15$ 分鐘，$1 \to 3 \to 5 \to 6$ 要 $2 + 7 + 2 = 11$ 分鐘。
最快的走法是 $1 \to 3 \to 2 \to 4 \to 5 \to 6$，經過五條路，只要
$2 + 3 + 2 + 1 + 2 = 10$ 分鐘。 《C++ 演算法解題 1：APCS 實作題核心》（以下簡
稱第 1 本）第~11~章的 BFS 數的是「經過幾條邊」，在這張圖上會找到三條路的走
法，答案是錯的。 邊上有數字、要比的是數字總和的圖，就是這一章的主題。

這一章先講最常用的 \textbf{Dijkstra 演算法}和怎麼印出路線，再講幾種建圖方
法：只有 $0$、$1$ 兩種代價的 0/1 BFS、虛擬起點、反向圖、狀態圖。 權重可以是負
數時改用 \textbf{Bellman--Ford 演算法}，它也能找出負環；要算每兩點之間的距離，
用 \textbf{Floyd--Warshall 演算法}。 最後一節整理怎麼從資料範圍選演算法。

這一章用到第 1 本第~11~章的 BFS、反向圖和多源 BFS，第 1 本第~12~章的狀態
BFS，第 1 本 6.5 節的 \texttt{priority\_queue} 和「先不刪」的寫法，以及本書
第~6~章在 DAG 上算最短路。 語法上需要 \texttt{pair} 和結構化綁定（銜接書第~8~
章）、「小的在上」的 \texttt{priority\_queue}（銜接書 8.6 節）、
\texttt{deque}（銜接書 8.5 節）。

\section{BFS 為什麼不夠}
\label{sec:sp-why}

先把名詞說清楚。 每條邊上有一個數字的圖叫\textbf{加權圖}（weighted graph），
那個數字叫邊的\textbf{權重}（weight），可以是時間、距離、費用。 一條路徑的
\textbf{長度}是路上每條邊的權重加起來，不是邊的條數。 從 $s$ 到 $v$ 的所有路
徑裡長度最小的叫\textbf{最短路徑}（shortest path），它的長度叫 $s$ 到 $v$ 的
\textbf{最短距離}，記作 $\delta(v)$；走不到時 $\delta(v) = \infty$。 從一個
起點求出到每個點的最短距離，叫\textbf{單源最短路徑}（single-source shortest
paths）。 這一章前七節都在處理這個問題，第~\ref{sec:sp-fw}~節才處理「每兩點
之間」。

存圖仍然用第 1 本第~11~章的鄰接串列，只是每一格多存權重：

\begin{code}{cpp}{有權重的鄰接串列}
vector<vector<pair<int, long long>>> adj(n + 1);   // adj[u] 裡放 (v, w)
adj[u].push_back({v, w});                          // 一條 u -> v、權重 w 的邊
\end{code}

\subsection{BFS 錯在哪裡}

第 1 本第~11~章說明 BFS 求出最短距離時，用到一件事：點被第一次看到的時候，
就是從離起點最近的那一層看到的，之後再看到它只會更遠。 這件事靠的是「每條邊
都算一步」。 在圖~\ref{fig:sp-map} 上，BFS 從 $4$ 或 $5$ 第一次看到 $6$，
這時已經是三條邊了；可是經過五條邊的那條路反而比較短。 邊的長度不同時，「先
看到」和「比較近」就不是同一件事。

有一個笨方法可以讓 BFS 重新變對：把權重 $w$ 的邊拆成 $w$ 條長度 $1$ 的小
段，中間插 $w - 1$ 個假的點，再做 BFS。 這樣每一步都是一分鐘，BFS 一層一層往
外，就是時間一分鐘一分鐘往後。 在圖~\ref{fig:sp-map} 上，真正的路口被走到的
順序是 $1$（第 $0$ 分鐘）、$3$（第 $2$ 分鐘）、$2$（第 $5$ 分鐘）、$4$（第
$7$ 分鐘）、$5$（第 $8$ 分鐘）、$6$（第 $10$ 分鐘），剛好是最短距離由小到大。

權重到 $10^9$ 時，假點多到存不下，這個方法不能真的用。 但它說明了要做的事：
照最短距離\textbf{由小到大}一個一個確定真正的點。 中間那些假點只是在「等時間
過去」，可以跳過：每次直接找出下一個會被走到的真正的點。 下一節的 Dijkstra
演算法做的就是這件事。

\subsection{鬆弛}

找「下一個會被走到的點」需要一個暫時的答案。 每個點 $v$ 記一個
\texttt{dist[v]}，表示「目前找到的、從起點到 $v$ 最短的走法」的長度，一開始起
點是 $0$，其他都是無限大。 已經知道到 $u$ 可以花 \texttt{dist[u]}，而且有一
條 $u \to v$、權重 $w$ 的邊，那麼到 $v$ 最多只要 \texttt{dist[u] + w}：

\begin{code}{cpp}{鬆弛一條邊}
if (dist[u] + w < dist[v]) {
    dist[v] = dist[u] + w;
}
\end{code}

這個動作叫\textbf{鬆弛}（relaxation）。 本章的每一個演算法都由鬆弛組成，差別
只在\textbf{用什麼順序}鬆弛哪些邊。 \texttt{dist[v]} 永遠是某一條真正的走法
的長度，所以它只會從上面往下逼近 $\delta(v)$，不會比 $\delta(v)$ 小。

鬆弛背後是最短路徑的一個性質：\textbf{最短路徑的一段，也是那兩點之間的最短路
徑}。 如果 $s \to \cdots \to u \to v$ 是到 $v$ 的最短路徑，前面 $s$ 到 $u$ 那
段一定也是到 $u$ 的最短路徑；否則把它換成更短的那條，到 $v$ 的路也跟著變短。
所以
\[
\delta(v) = \min_{\text{每條邊 } u \to v} \bigl(\delta(u) + w(u, v)\bigr),
\]
這就是第 1 本第~9~章的轉移式，「到 $v$ 最少花多少」只依賴前一站的答案。 圖沒
有環時，照本書第~6~章的拓樸順序一個點一個點算就好；有環時沒有事先排好的順
序。 第 1 本第~11~章說過，BFS 是一邊算一邊排，照答案本身由小到大填表。
Dijkstra 演算法也是這樣，只是權重不同時，要用優先佇列找出下一個最小的。

\section{Dijkstra 演算法}
\label{sec:sp-dijkstra}

\textbf{Dijkstra 演算法}（Dijkstra's algorithm）以提出者 Edsger Dijkstra 命
名。 它要求每條邊的權重都\textbf{不是負數}，步驟如下：

\begin{enumerate}
\item \texttt{dist[s]} 設成 $0$，其他點設成無限大，把 $(0, s)$ 放進一個「小的
在上」的優先佇列。
\item 優先佇列不空時，拿出最小的一筆 $(d, u)$。 如果 $d$ 不等於目前的
\texttt{dist[u]}，這筆是過期的紀錄，丟掉；否則 $u$ 的最短距離就是 $d$，從此不
再改變，接著把 $u$ 的每一條邊 $u \to v$ 鬆弛一次，成功時把
$(\texttt{dist[v]}, v)$ 放進優先佇列。
\item 優先佇列空了就結束。 \texttt{dist} 仍是無限大的點就是走不到的點。
\end{enumerate}

下面的程式讀入圖~\ref{fig:sp-map}，印出到每個點的最短距離，以及到終點的一條
最短路徑。 印路徑用到的 \texttt{par} 陣列在第~\ref{sec:sp-path}~小節說明。

\begin{example}[label=ex:sp-dijkstra]{cpp}{(dijkstra.cpp) Dijkstra 演算法}
#include <algorithm>
#include <functional>
#include <iostream>
#include <queue>
#include <utility>
#include <vector>
using namespace std;

const long long INF = 1e18;

int main() {
    int n, m, s, t;
    cin >> n >> m >> s >> t;
    vector<vector<pair<int, long long>>> adj(n + 1);   // (鄰點, 權重)
    for (int i = 0; i < m; i++) {
        int u, v;
        long long w;
        cin >> u >> v >> w;
        adj[u].push_back({v, w});                      // 雙向的道路
        adj[v].push_back({u, w});
    }
    vector<long long> dist(n + 1, INF);
    vector<int> par(n + 1, 0);                         // 最短路上的前一站
    // 小的在上：(目前的距離, 點)
    priority_queue<pair<long long, int>, vector<pair<long long, int>>,
                   greater<pair<long long, int>>> pq;
    dist[s] = 0;
    pq.push({0, s});
    while (!pq.empty()) {
        auto [d, u] = pq.top();
        pq.pop();
        if (d != dist[u]) {                            // 過期的紀錄，丟掉
            continue;
        }
        for (auto [v, w] : adj[u]) {
            if (d + w < dist[v]) {                     // 鬆弛
                dist[v] = d + w;
                par[v] = u;
                pq.push({dist[v], v});
            }
        }
    }
    for (int v = 1; v <= n; v++) {
        cout << (dist[v] == INF ? -1 : dist[v]) << (v < n ? ' ' : '\n');
    }
    if (dist[t] == INF) {
        cout << "no path\n";
        return 0;
    }
    vector<int> path;
    for (int v = t; v != s; v = par[v]) {              // 從終點往回走
        path.push_back(v);
    }
    path.push_back(s);
    reverse(path.begin(), path.end());
    for (size_t i = 0; i < path.size(); i++) {
        cout << path[i] << (i + 1 < path.size() ? ' ' : '\n');
    }
    return 0;
}
\end{example}

\iosample{
6 8 1 6 \\
1 2 7 \\
1 3 2 \\
2 3 3 \\
2 4 2 \\
3 5 7 \\
4 5 1 \\
4 6 6 \\
5 6 2 \\
}{
0 5 2 7 8 10 \\
1 3 2 4 5 6 \\
}

\begin{itemize}
\item \texttt{greater<pair<long long, int>>} 讓優先佇列變成小的在上（銜接書
8.6 節），\texttt{pair} 先比距離、距離相同再比點的編號。 沒有寫這個參數的話，
拿出來的是最大的，見本節的常見錯誤。
\item \texttt{auto [d, u] = pq.top();} 把最上面那筆\textbf{複製}出來再
\texttt{pop}。 寫成參考 \texttt{auto \&[d, u]} 的話，\texttt{pop} 之後
\texttt{d}、\texttt{u} 會變成別筆紀錄的值。
\item 距離可能到 $10^{14}$，要用 \texttt{long long}；無限大用 $10^{18}$，加上
一條邊也不會溢位。
\end{itemize}

\subsection{在圖上跑一次}

圖~\ref{fig:sp-trace} 的四張圖是程式在圖~\ref{fig:sp-map} 上的過程。 橘色的點已經定案，粉紅色的
點是這一步剛處理完的點，藍色數字是目前的 \texttt{dist}。

\begin{figure}[tbp]
\centering
\begin{tikzpicture}[x=0.95cm, y=0.82cm]
\begin{scope}
\spmapgraph{1/0/spvnow/left, 2/7/spv/above, 3/2/spv/below, 4/$\infty$/spv/above, 5/$\infty$/spv/below, 6/$\infty$/spv/right}{}
\node[font=\small] at (2.25,-1.3) {(a) 處理完 $1$};
\end{scope}
\begin{scope}[xshift=6.8cm]
\spmapgraph{1/0/spvdone/left, 2/5/spv/above, 3/2/spvnow/below, 4/$\infty$/spv/above, 5/9/spv/below, 6/$\infty$/spv/right}{}
\node[font=\small] at (2.25,-1.3) {(b) 處理完 $3$};
\end{scope}
\begin{scope}[yshift=-4.3cm]
\spmapgraph{1/0/spvdone/left, 2/5/spvdone/above, 3/2/spvdone/below, 4/7/spvnow/above, 5/8/spv/below, 6/13/spv/right}{}
\node[font=\small] at (2.25,-1.3) {(c) 處理完 $2$、$4$};
\end{scope}
\begin{scope}[xshift=6.8cm, yshift=-4.3cm]
\spmapgraph{1/0/spvdone/left, 2/5/spvdone/above, 3/2/spvdone/below, 4/7/spvdone/above, 5/8/spvdone/below, 6/10/spvdone/right}{1/3, 3/2, 2/4, 4/5, 5/6}
\node[font=\small] at (2.25,-1.3) {(d) 結束，粗線是 \texttt{par}};
\end{scope}
\end{tikzpicture}
\diagcaption{Dijkstra 演算法在外送地圖上的過程。}
\label{fig:sp-trace}
\end{figure}

每一次從優先佇列拿出一筆紀錄，對應下表的一列。 「堆裡的紀錄」依大小列出，實
際在堆積裡不是排好的（第 1 本 6.5 節）。

\begin{center}
\begin{tabular}{cll}
\toprule
拿出 & 做了什麼 & 之後堆裡的紀錄 \\
\midrule
$(0, 1)$ & 處理 $1$：$2$ 變 $7$，$3$ 變 $2$ & $(2,3)\ (7,2)$ \\
$(2, 3)$ & 處理 $3$：$2$ 從 $7$ 變 $5$，$5$ 變 $9$ & $(5,2)\ (7,2)\ (9,5)$ \\
$(5, 2)$ & 處理 $2$：$4$ 變 $7$ & $(7,2)\ (7,4)\ (9,5)$ \\
$(7, 2)$ & 過期（\texttt{dist[2]} 已經是 $5$），丟掉 & $(7,4)\ (9,5)$ \\
$(7, 4)$ & 處理 $4$：$5$ 從 $9$ 變 $8$，$6$ 變 $13$ & $(8,5)\ (9,5)\ (13,6)$ \\
$(8, 5)$ & 處理 $5$：$6$ 從 $13$ 變 $10$ & $(9,5)\ (10,6)\ (13,6)$ \\
$(9, 5)$ & 過期，丟掉 & $(10,6)\ (13,6)$ \\
$(10, 6)$ & 處理 $6$：沒有邊能鬆弛 & $(13,6)$ \\
$(13, 6)$ & 過期，丟掉 & （空） \\
\bottomrule
\end{tabular}
\end{center}

處理的順序是 $1, 3, 2, 4, 5, 6$，最短距離 $0, 2, 5, 7, 8, 10$ 由小到大，和上一
節「把邊拆成小段再 BFS」走到真正路口的順序一樣。

\subsection{過期的紀錄}

點 $2$ 先被 $1$ 更新成 $7$，放了一筆 $(7, 2)$；後來又被 $3$ 更新成 $5$，再放
一筆 $(5, 2)$。 \texttt{priority\_queue} 不能改裡面的值，也不能刪中間的元
素，所以舊的那筆就留在裡面。 這和第 1 本 6.5 節「刪掉堆積裡任意一個」的做法一
樣：\textbf{先不刪}，等它浮到最上面時再檢查，和 \texttt{dist} 對不上就丟掉。

檢查也可以改用一個 \texttt{done} 陣列：拿出 $u$ 時，\texttt{done[u]} 已經是
\texttt{true} 就跳過，否則設成 \texttt{true} 再處理。 兩種寫法在權重非負時效
果相同。 本章用比較 \texttt{d != dist[u]} 的寫法，因為少一個陣列，而且 $d$ 本
來就要拿來算新的距離。

\subsection{為什麼拿出來就定案}
\label{sec:sp-dijkstra-proof}

直覺是：權重沒有負的，路只會越走越長。 優先佇列最上面的 $(d, u)$ 是所有還沒
處理的紀錄裡最小的；任何還沒處理的點，最好也是 $d$ 或更大，從它們繞過來再走到
$u$，只會更長或一樣長，不可能比 $d$ 短。 完整的說明在下面的框裡，第一次讀可以
跳過。

\begin{whybox}[Dijkstra 拿出來的距離就是最短距離]
對處理的次數做歸納：假設之前處理過的點，\texttt{dist} 都已經等於最短距離。 現
在拿出 $(d, u)$，$d = \texttt{dist[u]}$，要說明 $d = \delta(u)$。

\texttt{dist[u]} 是某一條真正走法的長度，所以 $d \ge \delta(u)$。 假設 $d >
\delta(u)$，取一條長度 $\delta(u)$ 的最短路徑 $P$，從起點沿著 $P$ 往 $u$ 走。
起點已經處理過，$u$ 還沒有，所以途中一定有第一個「還沒處理」的點，叫它 $y$，
它在 $P$ 上的前一站 $x$ 已經處理過。

處理 $x$ 時鬆弛過 $x \to y$，所以 $\texttt{dist[y]} \le \delta(x) + w(x, y)$。
右邊是 $P$ 從起點走到 $y$ 那一段的長度，而 $P$ 從 $y$ 走到 $u$ 的剩下那段，每
條邊都不是負的，所以這段的長度不超過整條 $P$：
\[
\texttt{dist[y]} \le \delta(x) + w(x, y) \le \delta(u) < d.
\]
$y$ 還沒處理，\texttt{dist[y]} 最後一次被改的時候放了一筆
$(\texttt{dist[y]}, y)$，這筆還在優先佇列裡，而且比 $(d, u)$ 小，和「拿出來
的是最小的」矛盾。 所以 $d = \delta(u)$。

證明裡用到權重非負的地方只有一處：「$P$ 剩下那段的長度 $\ge 0$」。 有負權重
時這一步不成立，第~\ref{sec:sp-negative}~小節會看到實際出錯的例子。
\end{whybox}

\subsection{複雜度}

每個點只會被處理一次（之後的紀錄都會因為過期被丟掉），處理 $u$ 時把
\texttt{adj[u]} 看一遍，所以全部加起來看了 $m$ 條有向邊（雙向的道路算兩條）。
每看一條邊最多放一筆紀錄，優先佇列最多 $m + 1$ 筆，每次放入、拿出是
$O(\log m)$。 $m$ 不超過 $n^2$，$\log m \le 2 \log n$，所以總共是
$O((n + m) \log n)$。 $n = 10^5$、$m = 2 \times 10^5$ 時大約幾百萬次堆積操作，
很快。

只要一個終點的距離時，可以在拿出終點的那一刻就結束，因為拿出來就定案了。 第~
\ref{sec:ex-cart}~小節的例題就是這樣寫的。

\subsection{還原路徑}
\label{sec:sp-path}

要印出路線，就在每次鬆弛成功時記下「是從哪一點過來的」：\texttt{par[v] = u}。
\texttt{dist[v]} 最後一次被改的時候，\texttt{par[v]} 也跟著改，所以結束時
\texttt{par[v]} 是 $v$ 在某一條最短路徑上的前一站。 從終點沿著 \texttt{par}
一路往回走到起點，再把順序倒過來，就是一條最短路徑。 這和第 1 本第~9~章沿著
\texttt{last} 往回印出一種貼法是同一個做法。

圖~\ref{fig:sp-trace}(d) 的粗線是全部的 \texttt{par} 邊。 起點以外每個點剛好一
條，而且 \texttt{par[v]} 比 $v$ 早處理，往回走不會繞圈，所以粗線是一棵以起點為根
的樹，叫\textbf{最短路徑樹}（shortest-path tree）。 從起點順著它走到任何一點，
都是一條最短路徑。

\subsection{負權重會出錯}
\label{sec:sp-negative}

下面這張有向圖有一條權重 $-4$ 的邊，起點是 $1$。

\begin{center}
\begin{tikzpicture}[x=1.6cm, y=1.1cm]
\node[spv] (1) at (0,0) {1};
\node[spv] (2) at (1.5,0) {2};
\node[spv] (3) at (0.75,-1.3) {3};
\node[spv] (4) at (3,0) {4};
\draw[spd] (1) -- node[spw] {2} (2);
\draw[spd] (1) -- node[spw] {5} (3);
\draw[spdhl] (3) -- node[spw] {$-4$} (2);
\draw[spd] (2) -- node[spw] {1} (4);
\end{tikzpicture}
\diagcaption{有一條負權重邊的圖。 從 $1$ 到 $2$ 最短是 $1 \to 3 \to 2$，長度 $1$。}
\label{fig:sp-neg}
\end{center}

Dijkstra 處理完 $1$ 之後，\texttt{dist[2]} $= 2$、\texttt{dist[3]} $= 5$，接著
拿出 $(2, 2)$，認定 $2$ 的最短距離是 $2$，並用它把 $4$ 設成 $3$。 真正的答案
是 $\delta(2) = 1$、$\delta(4) = 2$：先繞到 $3$ 比較遠，但 $3 \to 2$ 會扣回來。
上面證明裡「剩下那段的長度 $\ge 0$」就是在這裡失效的。

用 \texttt{done} 陣列的寫法，處理 $3$ 時 \texttt{dist[2]} 會被改成 $1$，但 $2$
已經處理過、不會再處理，\texttt{dist[4]} 就停在錯的 $3$。 本章用比較
\texttt{dist} 的寫法，改成 $1$ 時會再放一筆 $(1, 2)$，$2$ 會被處理第二次，這張
圖碰巧得到正確的答案。 可是這樣
「處理過還要再改」的次數沒有上限的保證，有些圖會讓它慢到指數時間。 有負權重
時，正確的做法是第~\ref{sec:sp-bf}~節的 Bellman--Ford；圖沒有環時，用本書第~6~
章的拓樸順序，權重是負數也沒關係。

\subsection{常見錯誤}

\begin{itemize}
\item \textbf{優先佇列的方向反了}：少寫 \texttt{greater<...>}，拿出來的是距離
最大的。 範例的輸入仍然印出正確的答案，因為過期的檢查和鬆弛會一直修正，但順序
亂了，同一個點會被處理很多次。 在一張大約 $5 \times 10^4$ 個點、$5 \times 10^4$
條路的圖上，範例程式在本書測試的環境下要 $0.02$ 秒，這樣改要約 $6.5$ 秒。 如果
程式在拿出終點時就提早結束（第~\ref{sec:ex-cart}~小節的寫法），答案還會直接錯
掉：搬運車的範例 1 會印出 \texttt{15}。
\item \textbf{第一次看到就定案}：沿用 BFS 的習慣，\texttt{dist[v]} 還是無限
大才更新，更新過就不再改。 範例的輸入會印出 \texttt{0 7 2 9 9 15} 和
\texttt{1 2 4 6}：$2$ 被 $1$ 先看到，定成 $7$，後來 $3$ 帶來的 $5$ 被忽略。 要
定案的時機是\textbf{拿出來}的時候，不是放進去的時候。
\item \textbf{無限大太小}：用 $10^9$ 或 \texttt{0x3f3f3f3f}（約 $1.06 \times
10^9$）當無限大，題目的距離到 $10^{14}$ 時，真的距離比「無限大」還大，會被當成
走不到。 \texttt{dist} 用 \texttt{int} 時，相加直接溢位。
\item \textbf{有負權重還用 Dijkstra}：見上一小節。 先看題目的權重範圍有沒有
負數。
\end{itemize}

\section{0/1 BFS}
\label{sec:sp-01bfs}

第 1 本 12.4 節最後預告過一種權重只有 $0$ 和 $1$ 的圖，例如格子地圖上「最少要
拆幾面牆」。 這種圖用 Dijkstra 當然可以，但有一個更簡單、少一個 $\log$ 的做
法。

換一個例子。 社區的巷子都是單行道，夜間清運車要從車庫（路口 $1$）開到處理場
（路口 $n$）。 順著方向開不用申請；逆向開一段巷子，要事先申請一張通行證。 最
少要申請幾張？ 把每條單行道 $u \to v$ 拆成兩條邊：$u \to v$ 代價 $0$，
$v \to u$ 代價 $1$，問題就變成從 $1$ 到 $n$ 的最短距離。

\begin{center}
\begin{tikzpicture}[x=1.5cm, y=1.1cm]
\node[spv] (1) at (0,1) {1};
\node[spv] (2) at (1.5,2) {2};
\node[spv] (3) at (3,2) {3};
\node[spv] (4) at (2.25,1) {4};
\node[spv] (5) at (1.5,0) {5};
\node[spv] (6) at (4,1) {6};
\draw[spd] (1) -- (2);
\draw[spd] (2) -- (3);
\draw[spd] (6) -- (3);
\draw[spd] (5) -- (1);
\draw[spd] (6) -- (5);
\draw[spd] (4) -- (2);
\draw[spd] (4) -- (6);
\end{tikzpicture}
\diagcaption{七條單行道。 從 $1$ 到 $6$，順著箭頭走不用申請，逆著走一段要一張。}
\label{fig:sp-oneway}
\end{center}

圖~\ref{fig:sp-oneway} 裡，經過的巷子最少的是 $1 \to 5 \to 6$，兩段都是逆向，
要申請兩張。 $1 \to 2 \to 3$ 都是順向，最後一段逆著 $6 \to 3$ 開到 $6$，只要
一張。

做法是把 BFS 的佇列換成兩端都能放的 \texttt{deque}：從前端拿出 $u$，鬆弛成功
時，代價 $0$ 的放到\textbf{前端}，代價 $1$ 的放到\textbf{尾端}。 這叫
\textbf{0/1 BFS}（0-1 BFS）。

\begin{example}[label=ex:sp-zeroone]{cpp}{(zeroone.cpp) 0/1 BFS：最少的逆向通行證}
#include <deque>
#include <iostream>
#include <utility>
#include <vector>
using namespace std;

int main() {
    int n, m;
    cin >> n >> m;
    vector<vector<pair<int, int>>> adj(n + 1);   // (鄰點, 代價 0 或 1)
    for (int i = 0; i < m; i++) {
        int u, v;
        cin >> u >> v;                           // 單行道 u -> v
        adj[u].push_back({v, 0});                // 順向：不用申請
        adj[v].push_back({u, 1});                // 逆向：一張申請單
    }
    const int INF = 1e9;
    vector<int> dist(n + 1, INF);
    deque<int> dq;
    dist[1] = 0;
    dq.push_back(1);
    while (!dq.empty()) {
        int u = dq.front();
        dq.pop_front();
        for (auto [v, w] : adj[u]) {
            if (dist[u] + w < dist[v]) {
                dist[v] = dist[u] + w;
                if (w == 0) {
                    dq.push_front(v);            // 代價 0：插到最前面
                } else {
                    dq.push_back(v);             // 代價 1：排到最後面
                }
            }
        }
    }
    cout << (dist[n] == INF ? -1 : dist[n]) << '\n';
    return 0;
}
\end{example}

\iosample{
6 7 \\
1 2 \\
2 3 \\
6 3 \\
5 1 \\
6 5 \\
4 2 \\
4 6 \\
}{
1 \\
}

\texttt{deque} 的變化如下，括號裡是那個點當時的 \texttt{dist}：

\begin{center}
\begin{tabular}{lll}
\toprule
拿出 & 鬆弛成功 & 之後的 \texttt{deque}（前 $\to$ 後） \\
\midrule
 & & $1(0)$ \\
$1$ & $2$ 變 $0$ 放前端，$5$ 變 $1$ 放尾端 & $2(0)\ 5(1)$ \\
$2$ & $3$ 變 $0$ 放前端，$4$ 變 $1$ 放尾端 & $3(0)\ 5(1)\ 4(1)$ \\
$3$ & $6$ 變 $1$ 放尾端 & $5(1)\ 4(1)\ 6(1)$ \\
$5$ & 沒有（$5 \to 6$ 逆向要 $2$，不比 $1$ 好） & $4(1)\ 6(1)$ \\
$4$ & 沒有 & $6(1)$ \\
$6$ & 沒有 & （空） \\
\bottomrule
\end{tabular}
\end{center}

為什麼這樣是對的：\texttt{deque} 裡的 \texttt{dist} 由前往後不會變小，而且頭
尾最多差 $1$。 拿出前端的 $u$（值 $d$）時，裡面的值都在 $d$ 和 $d + 1$ 之間；
放到前端的是 $d$，放到尾端的是 $d + 1$，放完之後仍然成立。 這和第 1 本第~11~
章說明 BFS 時用的是同一件事。 所以點被拿出來的順序，\texttt{dist} 由小到大，
和 Dijkstra 拿出來的順序一樣，正確的理由也一樣。

同一個點可能進 \texttt{deque} 兩次：先以 $d + 1$ 放到尾端，之後又被別的點以 $d$
放到前端。 第二次拿出它時，\texttt{dist} 已經是最終的值，鬆弛都不會成功，不影
響答案。 一個點的 \texttt{dist} 最多被改兩次，所以總共是 $O(n + m)$。

\subsection{常見錯誤}

\begin{itemize}
\item \textbf{當成普通的 BFS}：用 \texttt{queue}，第一次看到就定案。 範例的
輸入會印出 \texttt{2}：$6$ 第一次被看到，是從 $5$（\texttt{dist} $1$）逆向過
去，定成 $2$；後來從 $3$ 順著過去只要 $1$，已經改不了了。
\item \textbf{代價不只 $0$、$1$}：有代價 $2$ 的邊時，放到尾端會破壞「頭尾最多
差 $1$」。 這時改用 Dijkstra。
\end{itemize}

\section{多個起點與虛擬起點}
\label{sec:sp-multi}

第 1 本第~11~章的多源 BFS，把所有起點一起放進佇列、距離都設成 $0$，算出每個
點離\textbf{最近的}起點多遠。 加權圖也一樣：把所有起點的 \texttt{dist} 設成
$0$、一起放進優先佇列，做一次 Dijkstra。

起點各自有不同的「出發成本」時，例如第 $i$ 個起點要先等 $p_i$ 分鐘才能出發，
可以想像多一個\textbf{虛擬起點} $0$，從 $0$ 拉一條權重 $p_i$ 的邊到第 $i$ 個起
點，再從 $0$ 做 Dijkstra。 任何從 $0$ 出發的路徑，第一步一定是這些邊之一，所以
到 $v$ 的最短距離就是 $\min_i \bigl(p_i + \text{起點 } i \text{ 到 } v \text{
的最短距離}\bigr)$。 程式裡不用真的加點和邊：把第 $i$ 個起點的 \texttt{dist}
設成 $p_i$（同一點有好幾個起點時取最小），一起放進優先佇列，就是從 $0$ 處理完
之後的狀態。

\subsection{記住是哪一個起點}
\label{sec:sp-label}

有時候不只要距離，還要知道最近的是\textbf{哪一個}起點，距離相同時取編號小的。
例如三間消防隊同時出動，每個村落要知道哪一隊最先到。 做法是把每個點的
\texttt{dist} 從一個數換成一對數 $(\text{距離}, \text{起點編號})$，用
\texttt{pair} 的大小比較：先比距離，相同再比編號。 沿著一條權重 $w$ 的邊走，
標籤從 $(d, i)$ 變成 $(d + w, i)$。

Dijkstra 的證明（第~\ref{sec:sp-dijkstra-proof}~小節）只用到兩件事：沿著路徑往
前走，標籤不會變小；兩個標籤 $L_1 \le L_2$ 各自接上同一條邊之後，仍然是
$L_1 + w \le L_2 + w$。 $(d, i)$ 加上 $w$ 只動到距離、編號不變，這兩件事都對，
所以優先佇列裡放 $(\text{距離}, \text{編號}, \text{點})$ 照樣做 Dijkstra。 C++
的 \texttt{tuple} 就是三個欄位的 \texttt{pair}，比較也是從第一個依序比，
\texttt{auto [d, id, u]} 一樣可以拆開。

要注意兩個地方。 一是鬆弛的條件要比整個標籤：只比距離的話，先到的起點以距離 $6$
佔住某個點之後，編號更小、距離也是 $6$ 的起點就改不過來。 二是過期的檢查也要
比整個標籤，不然一筆距離相同、編號較大的舊紀錄會被當成有效的。 同一個點有好幾
個起點時，初值取標籤最小的那一個。

同樣的模型也出現在別的地方：每戶人家離哪一間醫院最近；好幾個倉庫一起出貨、各
自的備貨時間不同，每個門市由哪一個倉庫供貨最快。 實作練習的
\akiref{fire-crews-getting-ready}{消防分隊的整備時間}就是這種題目。

\section{反向圖：每個點到同一個終點}
\label{sec:sp-reverse}

Dijkstra 算的是「從一個起點到每個點」。 有向圖的題目常常反過來問：每個人從自
己的位置走到同一個集合點，各要多久？ 對每個人各做一次 Dijkstra 太慢。 第 1 本
第~11~章用過反向圖：把每條邊 $u \to v$ 反過來存成 $v \to u$，在反向圖上從終點
$t$ 出發做一次，得到的就是每個點\textbf{走到} $t$ 的最短距離。 權重照抄，
Dijkstra 本身不用改。

實作時把 Dijkstra 寫成函式 \texttt{dijkstra(g, s)}，回傳整個 \texttt{dist}；讀邊時
同時建 \texttt{adj} 和反向的 \texttt{radj}，再各呼叫一次，得到 \texttt{ds}（從 $s$
出發）和 \texttt{dt}（走到 $t$）。

有了 \texttt{ds} 和 \texttt{dt}，很多問題變成查表：

\begin{itemize}
\item \textbf{一定要經過某個點 $x$}：從 $s$ 經過 $x$ 再到 $t$，最短是
\texttt{ds[x] + dt[x]}。 前後兩段互不影響，各自取最短就好。 能經過的點有好幾
個（例如任何一間有賣某樣東西的店）時，對它們取最小值。
\item \textbf{哪些邊在某條最短路徑上}：邊 $u \to v$（權重 $w$）在某一條 $s$ 到
$t$ 的最短路徑上，剛好是 \texttt{ds[u] + w + dt[v]} 等於 \texttt{ds[t]} 的時
候。
\end{itemize}

例如一張有向圖有七條邊：$1 \to 2$（$3$）、$1 \to 3$（$1$）、$3 \to 2$（$1$）、
$2 \to 5$（$4$）、$3 \to 4$（$5$）、$4 \to 5$（$1$）、$5 \to 3$（$2$），起點
$1$、終點 $5$。 兩次 Dijkstra 的結果：

\begin{center}
\begin{tabular}{cccccc}
\toprule
點 $x$ & 1 & 2 & 3 & 4 & 5 \\
\midrule
\texttt{ds[x]}（從 $1$ 出發） & 0 & 2 & 1 & 6 & 6 \\
\texttt{dt[x]}（走到 $5$） & 6 & 4 & 5 & 1 & 0 \\
\texttt{ds[x] + dt[x]} & 6 & 6 & 6 & 7 & 6 \\
\bottomrule
\end{tabular}
\end{center}

最短距離是 $6$。 一定要經過 $4$ 的話最短是 $7$。 七條邊裡，滿足
\texttt{ds[u] + w + dt[v]} $= 6$ 的是 $1 \to 3$、$3 \to 2$、$2 \to 5$，它們組成
唯一的一條最短路徑 $1 \to 3 \to 2 \to 5$。 注意 \texttt{dt} 一定要在反向圖上
算：從 $5$ 出發在原圖上走得到 $3$，但那是「$5$ 到 $3$」，不是「$3$ 到 $5$」。

\section{狀態圖}
\label{sec:sp-state}

第 1 本第~12~章的狀態 BFS 說過：同一個位置，如果手上的東西、面向、剩下的量不同，
之後能做的事就不同，這時一個點要代表「位置加上這些資訊」。 加權圖也一樣，只是
狀態之間的邊有不同的權重，BFS 換成 Dijkstra。 狀態圖通常不用真的建出來：把狀態
編成一個整數當作 \texttt{dist} 的索引，處理到它的時候，當場算出它能走到哪些狀態、
各花多少。

常見的附加資訊有三種。

\textbf{用了幾次。} 例如手上有 $k$ 張半價券，每條路最多用一張。 狀態是
$(\text{位置}, \text{已經用了幾張})$，共 $n(k + 1)$ 個。 從 $(u, j)$ 走一條權重
$w$ 的路，可以付全額到 $(v, j)$，或在 $j < k$ 時用一張券、付 $w / 2$ 到
$(v, j + 1)$。 畫出來像疊起來的 $k + 1$ 層地圖，券只能讓你往下一層走，所以這種
圖也叫\textbf{分層圖}（layered graph）。 答案是終點在每一層的距離取最小。

\begin{center}
\begin{tikzpicture}[x=1.7cm, y=1.5cm]
\node[font=\small, anchor=east] at (-0.4,1.4) {第 0 層（還沒用券）};
\node[font=\small, anchor=east] at (-0.4,0) {第 1 層（用過一張）};
\node[spv] (a0) at (0,1.4) {1};
\node[spv] (b0) at (1.5,1.4) {2};
\node[spv] (c0) at (3,1.4) {3};
\node[spv] (a1) at (0,0) {1};
\node[spv] (b1) at (1.5,0) {2};
\node[spvdone] (c1) at (3,0) {3};
\draw[spe] (a0) -- node[spw] {8} (b0);
\draw[spe] (b0) -- node[spw] {6} (c0);
\draw[spe] (a0) to[bend left=30] node[spw] {18} (c0);
\draw[spe] (a1) -- node[spw] {8} (b1);
\draw[spe] (b1) -- node[spw] {6} (c1);
\draw[spe] (a1) to[bend right=30] node[spw] {18} (c1);
\draw[spd, dashed] (a0) -- node[spw, pos=0.3] {4} (b1);
\draw[spd, dashed] (b0) -- node[spw, pos=0.3] {3} (c1);
\draw[spdhl, dashed] (a0) -- node[spw, pos=0.3] {9} (c1);
\end{tikzpicture}
\diagcaption{三個點、一張半價券的分層圖。 虛線是用券的邊，只能往下走（往左走的用券
邊省略沒畫）。 不用券最
短是 $1 \to 2 \to 3$ 的 $14$；用一張券，最好是把券用在 $18$ 那條路上，只要 $9$。}
\label{fig:sp-layer}
\end{center}

圖~\ref{fig:sp-layer} 也說明了為什麼不能「先找最短路徑，再把最貴的那段打半
價」：最短路徑 $1 \to 2 \to 3$ 打完折是 $8 + 3 = 11$，比直接把券用在 $18$ 那條
路上的 $9$ 還貴。 用券改變了哪一條路最好，所以要把券數放進狀態裡一起搜。

\textbf{朝向。} 車子、機器人轉彎要花時間時，同一格面向不同，之後要付的轉彎代
價就不同，狀態要包含朝向。 下一小節的例題就是這種。

\textbf{時間。} 有些邊什麼時候能走和時刻有關，例如路口要等綠燈、平台退潮才露出
來。 如果\textbf{早到不吃虧}，也就是可以在原地等，早一點到某個點，能做的事不會
比晚到少，那麼每個點只要記最早到達的時刻就好，不用把時間放進狀態。 鬆弛時算出
「從 $u$ 最早什麼時候能出發」，再加上路程：
\[
\texttt{dist[v]} \leftarrow \min\bigl(\texttt{dist[v]},\ \text{最早能出發的時刻}(\texttt{dist[u]}) + w\bigr).
\]
最早能出發的時刻不會早於到達的時刻，而且到得越早，能出發的時刻也越早（或一樣），
所以 Dijkstra 的證明照樣成立。 反過來，如果規定到了不能等、必須馬上走，早到就可
能比晚到糟，這時就要把時間放進狀態，或另外想辦法。

不管哪一種，動手前先估狀態數和邊數（第 1 本 12.4 節的方法）。 Dijkstra 的時間
大約是 $(\text{狀態數} + \text{邊數}) \times \log(\text{狀態數})$。

\subsection{例題：轉彎很慢的搬運車}
\label{sec:ex-cart}

\textbf{題意}（\akiref{the-cart-that-turns-slowly}{轉彎很慢的搬運車}）

\begin{itemize}
\item 倉庫是 $R \times C$ 的方格。 \texttt{\#} 是貨架，不能進入；數字
\texttt{1} 到 \texttt{9} 是地面，數字是駛入這一格要花的秒數；\texttt{S} 是起
點、\texttt{T} 是終點，駛入這兩格都算 $1$ 秒。
\item 搬運車停在 \texttt{S}，一開始面向東。 每次可以往目前的朝向前進一格，或原
地左轉、右轉 $90^\circ$，轉一次花 $K$ 秒；掉頭就是轉兩次。
\item 開進 \texttt{T} 就算送達，面向哪裡都可以。 求最少的秒數，到不了輸出
$-1$。
\item $R, C \le 500$，$0 \le K \le 10^4$。
\end{itemize}

範例 1 的地圖如下，$K = 3$。

\begin{center}
\begin{tikzpicture}[x=0.9cm, y=0.9cm]
\foreach \t [count=\i from 0] in {2, 3, 1, T} {
    \node[cell] at (\i,1) {\t};
}
\foreach \t [count=\i from 0] in {S, 4, 2} {
    \node[cell] at (\i,0) {\t};
}
\node[cell, fill=black!35] at (3,0) {\#};
\draw[-{Stealth[length=2.2mm]}, very thick, draw=memStroke] (-0.3,0.15) -- (-0.3,1.3) -- (1.8,1.3);
\draw[-{Stealth[length=2.2mm]}, very thick, draw=optStroke, dashed] (-0.1,-0.3) -- (2.3,-0.3) -- (2.3,0.85);
\node[font=\small, anchor=west] at (4,1) {甲：先北再東，到 \texttt{1} 是第 $12$ 秒、面向東};
\node[font=\small, anchor=west] at (4,0) {乙：先東再北，到 \texttt{1} 是第 $10$ 秒、面向北};
\end{tikzpicture}
\diagcaption{範例 1 的兩種走法。 \texttt{T} 只能從左邊那格往東開進去。}
\label{fig:sp-cart}
\end{center}

走法甲：左轉（$3$）、駛入 \texttt{2}（$5$）、右轉（$8$）、駛入 \texttt{3}、
\texttt{1}（$12$），再往東駛入 \texttt{T}，共 $13$ 秒。 走法乙：駛入 \texttt{4}、
\texttt{2}（$6$）、左轉（$9$）、駛入上面的 \texttt{1}（$10$），這時車頭朝北，
還要右轉（$13$）才能開進 \texttt{T}（$14$）。 乙比較早到 \texttt{T} 左邊那格，最
後反而比較慢。 答案是 $13$。

\textbf{想法}

每一格只記一個最小秒數的話，\texttt{T} 左邊那格會記下乙的 $10$ 秒、面向北，甲的
$12$ 秒、面向東被丟掉，最後算出 $14$。 慢一點但方向對的走法，之後可能比較快，所
以狀態要是 $(\text{列}, \text{行}, \text{朝向})$，共 $4RC$ 個。 每個狀態有三條
往外的邊：

\begin{itemize}
\item 右轉、左轉：到同一格的另外兩個朝向，權重 $K$。
\item 前進：到前方那一格、同一個朝向，權重是駛入那一格的秒數。 出界或是貨架就
沒有這條邊。
\end{itemize}

權重有 $K$ 也有 $1$ 到 $9$，不是全部一樣，也不只有 $0$、$1$，所以用 Dijkstra。
掉頭不用另外處理，連續轉兩次就是 $2K$。 終點的四個朝向都算送達，拿出任何一個
\texttt{T} 的狀態，那就是答案，可以馬上結束。

\begin{center}
\begin{tikzpicture}[x=1.15cm, y=1.15cm]
\node[spv] (aE) at (1,0) {東};
\node[spv] (aS) at (0,-1) {南};
\node[spv] (aW) at (-1,0) {西};
\node[spv] (aN) at (0,1) {北};
\draw[spe, <->, >={Stealth[length=2mm]}] (aE) -- node[spw] {$K$} (aS);
\draw[spe, <->, >={Stealth[length=2mm]}] (aS) -- node[spw] {$K$} (aW);
\draw[spe, <->, >={Stealth[length=2mm]}] (aW) -- node[spw] {$K$} (aN);
\draw[spe, <->, >={Stealth[length=2mm]}] (aN) -- node[spw] {$K$} (aE);
\node[spv] (bE) at (5,0) {東};
\node[spv] (bS) at (4,-1) {南};
\node[spv] (bW) at (3,0) {西};
\node[spv] (bN) at (4,1) {北};
\draw[spdhl] (aE) .. controls (2.2,-2.6) and (5.9,-2.6) .. node[spw, pos=0.5] {駛入 $B$ 的秒數} (bE);
\draw[black!40, rounded corners] (-1.6,-1.5) rectangle (1.6,1.5);
\draw[black!40, rounded corners] (2.4,-1.5) rectangle (5.6,1.5);
\node[font=\small] at (0,1.8) {格子 $A$};
\node[font=\small] at (4,1.8) {格子 $B$（在 $A$ 的東邊）};
\end{tikzpicture}
\diagcaption{一格拆成四個狀態。 同一格的相鄰朝向互相連 $K$ 秒的邊；只有面向東的
狀態能前進到東邊那格，而且到達後仍然面向東。}
\label{fig:sp-cart-states}
\end{center}

同樣的模型也出現在別的地方：堆高機、掃地機器人轉向比直走慢；地圖導航算路線時，
左轉要等號誌，比直行和右轉多花時間。

\textbf{程式}

\begin{example}[label=ex:sp-cart]{cpp}{(cart.cpp) 轉彎很慢的搬運車}
#include <functional>
#include <iostream>
#include <queue>
#include <string>
#include <utility>
#include <vector>
using namespace std;

const long long INF = 1e18;
const int dr[4] = {0, 1, 0, -1};   // 0 東、1 南、2 西、3 北（順時針）
const int dc[4] = {1, 0, -1, 0};

int main() {
    ios::sync_with_stdio(false);
    cin.tie(nullptr);
    int R, C;
    long long K;
    cin >> R >> C >> K;
    vector<string> g(R);
    for (string &row : g) {
        cin >> row;
    }
    // 狀態 (r, c, d) 編成一個整數 (r * C + c) * 4 + d
    vector<long long> dist(R * C * 4, INF);
    priority_queue<pair<long long, int>, vector<pair<long long, int>>,
                   greater<pair<long long, int>>> pq;
    for (int r = 0; r < R; r++) {
        for (int c = 0; c < C; c++) {
            if (g[r][c] == 'S') {
                int s = (r * C + c) * 4 + 0;          // 一開始面向東
                dist[s] = 0;
                pq.push({0, s});
            }
        }
    }
    while (!pq.empty()) {
        auto [d, st] = pq.top();
        pq.pop();
        if (d != dist[st]) {
            continue;
        }
        int dir = st % 4, r = st / 4 / C, c = st / 4 % C;
        if (g[r][c] == 'T') {                       // 第一次拿出 T 就是答案
            cout << d << '\n';
            return 0;
        }
        // 兩種轉向：右轉 dir + 1、左轉 dir + 3（都取除以 4 的餘數）
        for (int nd : {(dir + 1) % 4, (dir + 3) % 4}) {
            int ns = (r * C + c) * 4 + nd;
            if (d + K < dist[ns]) {
                dist[ns] = d + K;
                pq.push({dist[ns], ns});
            }
        }
        // 往前一格
        int nr = r + dr[dir], nc = c + dc[dir];
        if (nr < 0 || nr >= R || nc < 0 || nc >= C || g[nr][nc] == '#') {
            continue;
        }
        char ch = g[nr][nc];
        long long cost = (ch >= '1' && ch <= '9') ? ch - '0' : 1;
        int ns = (nr * C + nc) * 4 + dir;
        if (d + cost < dist[ns]) {
            dist[ns] = d + cost;
            pq.push({dist[ns], ns});
        }
    }
    cout << -1 << '\n';
    return 0;
}
\end{example}

\iosample{
2 4 3 \\
231T \\
S42\# \\
}{
13 \\
}

\begin{itemize}
\item 朝向照順時針編號，右轉是加 $1$、左轉是加 $3$（等於減 $1$），再取除以
$4$ 的餘數，不用寫四種情況。
\item 狀態編成 \texttt{(r * C + c) * 4 + d}，拆回來是 \texttt{d = st \% 4}、
\texttt{r = st / 4 / C}、\texttt{c = st / 4 \% C}。 最多 $4 \times 500 \times
500 = 10^6$ 個狀態，\texttt{dist} 約 8 MB。
\item $K = 0$ 時轉向的邊權重是 $0$。 Dijkstra 允許權重 $0$，不用特別處理。
\end{itemize}

\textbf{複雜度}

$4RC$ 個狀態、每個最多三條邊，$O(RC \log(RC))$，大約 $10^6 \times 20$。 最大
的測資在本書測試的環境下約 $0.12$ 秒。

\subsection{常見錯誤}

\begin{itemize}
\item \textbf{只用格子當狀態}：每格記最小秒數和當時的朝向，下一步依朝向的差補
上轉彎的代價。 範例 1 會印出 \texttt{14}，就是上面走法乙的結果。
\item \textbf{起點四個朝向都從 $0$ 開始}：題目規定一開始面向東。 範例 1 會印
出 \texttt{10}：直接面向北出發，省掉了第一次轉彎。
\item \textbf{無限大設成 $10^9$}：測資裡有每走一格都要轉彎的蛇行地圖，答案超
過 $1.6 \times 10^9$。
\item \textbf{每個動作都算一步，用 BFS}：前進和轉彎的代價不同，BFS 求的是動作
數最少，不是秒數最少。
\end{itemize}

\section{Bellman--Ford 與負環}
\label{sec:sp-bf}

權重可以是負數的情況不少：走某段路會拿到回饋點數，遊戲裡經過某個房間會撿到金幣，
或是題目給的是「$x_b - x_a \le d$」這種差距的限制。 第~\ref{sec:sp-negative}~小
節看過，這時 Dijkstra 的「拿出來就定案」不成立。

\textbf{Bellman--Ford 演算法}（Bellman--Ford algorithm）不管順序，直接把\textbf{每
一條邊}都鬆弛一次，當作一輪，重複 $n - 1$ 輪。 邊用一個串列存起來就好，不需要
鄰接串列：

\begin{code}{cpp}{Bellman--Ford 的主迴圈}
struct Edge {
    int u, v;
    long long w;
};
// dist[s] = 0，其他是 INF
for (int round = 1; round < n; round++) {
    for (const Edge &e : edges) {
        if (dist[e.u] != INF && dist[e.u] + e.w < dist[e.v]) {
            dist[e.v] = dist[e.u] + e.w;
        }
    }
}
\end{code}

\texttt{dist[e.u] != INF} 這個檢查不能省。 少了它，走不到的點的「無限大」加上
負的權重，會變成一個比無限大小一點的數，看起來像走得到。

拿圖~\ref{fig:sp-neg} 來跑，故意把邊排成 $2 \to 4$、$3 \to 2$、$1 \to 2$、
$1 \to 3$ 的順序（$n = 4$，跑三輪），表裡是每個點的 \texttt{dist}：

\begin{center}
\begin{tabular}{lccccl}
\toprule
 & 點 1 & 點 2 & 點 3 & 點 4 & 這一輪發生的事 \\
\midrule
開始 & 0 & $\infty$ & $\infty$ & $\infty$ & \\
第 1 輪 & 0 & 2 & 5 & $\infty$ & 前兩條邊的起點還是 $\infty$，跳過 \\
第 2 輪 & 0 & 1 & 5 & 3 & $2 \to 4$ 先用舊的 $2$，之後 $3 \to 2$ 才改成 $1$ \\
第 3 輪 & 0 & 1 & 5 & 2 & $2 \to 4$ 用新的 $1$ \\
\bottomrule
\end{tabular}
\end{center}

邊的順序不好時，好消息要好幾輪才傳得過去。 最短路徑 $1 \to 3 \to 2 \to 4$ 有三
條邊，剛好用了三輪。 下面的框說明為什麼 $n - 1$ 輪一定夠。

\begin{whybox}[$n - 1$ 輪就夠]
先證明：\textbf{第 $i$ 輪結束時，\texttt{dist[v]} 不超過「從起點到 $v$、最多用
$i$ 條邊的走法」裡最短的長度。} 對 $i$ 歸納。 $i = 0$ 時只有起點是 $0$，對。 任
取一條最多 $i$ 條邊、走到 $v$ 的走法，它的最後一條邊是 $u \to v$，前面那段最多
$i - 1$ 條邊，依歸納，第 $i - 1$ 輪結束時 \texttt{dist[u]} 不超過那段的長度。
\texttt{dist} 只會變小，第 $i$ 輪鬆弛 $u \to v$ 時，\texttt{dist[v]} 就不超過
整條走法的長度。

如果從起點走得到的地方沒有負環（下一小節），最短的走法可以不重複經過任何點：
重複經過表示中間繞了一圈，這一圈的長度不是負的，拿掉不會變長。 不重複的走法最
多 $n - 1$ 條邊，所以第 $n - 1$ 輪結束時 \texttt{dist[v]} $\le \delta(v)$；而
\texttt{dist[v]} 永遠是某條真正走法的長度，不會小於 $\delta(v)$，兩邊夾起來就相
等。

注意第 $i$ 輪結束時的 \texttt{dist[v]} 可能比「最多 $i$ 條邊」的答案還小，因為
同一輪裡前面改過的值，後面的邊馬上就用上了（上表第 2 輪的 $3 \to 2$ 就是在
$2 \to 4$ 之後才改）。 題目限制邊數時要小心，見本章練習第~\ref{drill:sp-klimit}~
題。
\end{whybox}

每一輪看 $m$ 條邊，共 $n - 1$ 輪，複雜度 $O(nm)$。 實際上常常不到 $n - 1$ 輪就
不再有任何變化；某一輪什麼都沒改，之後每一輪也都一樣，可以直接停下來。

\subsection{負環和受影響的點}
\label{sec:sp-negcycle}

一圈走回原點、權重總和是負數的環，叫\textbf{負環}（negative cycle）。 從起點走
得到一個負環時，繞一圈就少一點，繞越多圈越少，環上的點沒有最短距離，可以任意
小；從環上走得到的每個點也一樣，先繞很多圈再走過去就好。 這些點的答案是
$-\infty$。

怎麼找出它們：跑完 $n - 1$ 輪之後，再把每條邊檢查一次。 沒有負環的話，上面的框說
\texttt{dist} 已經是最短距離，不可能再鬆弛；所以這時還能鬆弛的邊 $u \to v$，它
的 $v$ 一定受負環影響。 把這些 $v$ 標起來，再沿著邊往下游做一次 BFS 或 DFS，走到
的點全部標成 $-\infty$。

只標「第 $n$ 輪還能鬆弛的點」不夠，一定要往下游擴散。 下游的點可能另外有一條便
宜很多的路，負環每圈只少一點點，第 $n$ 輪時還追不上，那條邊就不會被鬆弛；但繞
夠多圈之後一定追得上。 本節的例題會看到實際的輸入。

\begin{whybox}[標起來再擴散，剛好是受影響的點]
\textbf{標到的點都受影響。} 假設 $u \to v$ 在第 $n$ 輪還能鬆弛，但 $v$ 不受任何
負環影響。 那麼 $u$ 也不受影響（否則 $v$ 在 $u$ 的下游，也會受影響），上面的框
說 $n - 1$ 輪後 \texttt{dist[u]} $= \delta(u)$、\texttt{dist[v]} $= \delta(v)$，
而 $\delta(v) \le \delta(u) + w$，不能再鬆弛，矛盾。 從受影響的點往下游走到的點
當然也受影響。

\textbf{受影響的點都標得到。} 只要說明每個走得到的負環上，至少有一條邊在第 $n$
輪還能鬆弛。 假設某個負環 $x_1 \to x_2 \to \cdots \to x_k \to x_1$ 上每條邊都不
能鬆弛，也就是每條邊都有 $\texttt{dist}[x_{j+1}] \le \texttt{dist}[x_j] +
w(x_j, x_{j+1})$。 這 $k$ 個不等式加起來，左右兩邊的 \texttt{dist} 剛好互相消
掉，剩下 $0 \le$ 環上權重的總和，和負環矛盾。 所以每個負環上都有點被標到，環上
其他點和下游都會在擴散時走到。
\end{whybox}

從起點走不到的負環不影響任何答案。 程式裡 \texttt{dist[e.u] != INF} 的檢查讓這
些環永遠不會被鬆弛，也就不會被標到。 網路上常見的另一種寫法是把所有
\texttt{dist} 一開始都設成 $0$，用來檢查「整張圖有沒有負環」，它會把走不到的負
環也找出來，在這一節的問題裡是錯的。

\subsection{例題：多入口地城的最低花費}
\label{sec:ex-dungeon}

\textbf{題意}（\akiref{cheapest-walk-in-a-multi-entrance-dungeon}{多入口地城的最低花費}）

\begin{itemize}
\item 地城有 $n$ 個房間、$m$ 條單向通道，走過第 $j$ 條通道要付 $w_j$ 枚金幣，
$w_j$ 是負數時代表撿到金幣。 每次經過都重新計算。 可能有通道繞回同一個房間。
\item 有 $k$ 個入口，第 $i$ 個在房間 $a_i$，入場費 $e_i$。 選一個入口進場，沿通
道走任意多步（可以不走），在任何房間都可以停下。
\item 對每個房間，求停在那裡的最小總花費：沒有入口走得到就輸出 \texttt{X}；可以
任意小就輸出 \texttt{-INF}；否則輸出最小值。
\item $n \le 5000$，$m \le 10^4$，$e_i \le 10^9$，$|w_j| \le 10^9$。
\end{itemize}

範例 2 只有一個入口：房間 $1$，入場費 $5$。

\begin{center}
\begin{tikzpicture}[x=1.45cm, y=1.1cm]
\node[spv] (1) at (0,0) {1};
\node[spvnow] (2) at (1.5,0.8) {2};
\node[spvnow] (3) at (3,0.8) {3};
\node[spvnow] (4) at (4.5,0.8) {4};
\node[spv] (5) at (1.5,-0.9) {5};
\node[spvgray] (6) at (3.4,-0.9) {6};
\node[spvgray] (7) at (4.9,-0.9) {7};
\draw[spd] (1) -- node[spw] {3} (2);
\draw[spdhl] (2) to[bend left=25] node[spw] {$-4$} (3);
\draw[spdhl] (3) to[bend left=25] node[spw] {2} (2);
\draw[spd] (3) -- node[spw] {1} (4);
\draw[spd] (1) -- node[spw] {2} (5);
\draw[spdgray] (6) to[bend left=25] node[spw] {$-1$} (7);
\draw[spdgray] (7) to[bend left=25] node[spw] {$-1$} (6);
\draw[spdgray] (7) to[bend left=30] node[spw] {0} (5);
\node[font=\small, anchor=east] at (-0.4,0) {入口，入場費 $5$};
\end{tikzpicture}
\diagcaption{範例 2。 $2 \to 3 \to 2$ 繞一圈是 $-2$，粉紅色的點都是 \texttt{-INF}；
$6$、$7$ 的負環沒有入口走得到。}
\label{fig:sp-dungeon}
\end{center}

房間 $2$、$3$ 在負環上，房間 $4$ 不在環上，但可以先繞很多圈再走 $3 \to 4$，三
個都是 \texttt{-INF}。 房間 $5$ 只能走 $1 \to 5$，花費 $7$。 $6 \leftrightarrow
7$ 也是負環，還有一條 $7 \to 5$，但沒有入口走得到 $6$、$7$，它們是 \texttt{X}，
也不會把房間 $5$ 拉成 \texttt{-INF}。

\textbf{想法}

最直接的做法是每個入口各跑一次 Bellman--Ford，$k \cdot n \cdot m$ 最多 $5000
\times 5000 \times 10^4$，太慢。 和第~\ref{sec:sp-multi}~節一樣用虛擬起點：從
虛擬起點連一條權重 $e_i$ 的邊到每個入口。 程式裡就是把 \texttt{dist[a\_i]} 設
成 $e_i$，其他設成無限大，再照常跑。 這裡權重有負的，所以用 Bellman--Ford，不
用 Dijkstra。

跑完 $n - 1$ 輪（加上虛擬起點是 $n + 1$ 個點，要 $n$ 輪；但虛擬起點的邊在一開始
就鬆弛完了，等於已經跑過一輪，所以 $n - 1$ 輪仍然夠），再用上一小節的方法找出
\texttt{-INF} 的房間。 \texttt{dist} 仍是無限大的房間輸出 \texttt{X}。

範例 2 跑第 1 輪後，房間 $1$ 到 $5$ 的 \texttt{dist} 是 $5, 6, 4, 5, 7$，第 2 輪
是 $5, 4, 2, 3, 7$，第 6 輪是 $5, -4, -6, -5, 7$；$6$、$7$ 一直是無限大。
房間 $2$、$3$、$4$ 每一輪都少 $2$，永遠停不下來。 跑完 $6$ 輪之後檢查，只有
$2 \to 3$ 還能鬆弛（$-4 - 4 < -6$），標出房間 $3$；再往下游擴散，走到 $2$ 和
$4$。 只看第 $7$ 輪的話，連環上的房間 $2$ 都會漏掉。

\textbf{程式}

\begin{example}[label=ex:sp-dungeon]{cpp}{(dungeon.cpp) 多入口地城的最低花費}
#include <iostream>
#include <queue>
#include <vector>
using namespace std;

const long long INF = 1e18;

struct Edge {
    int u, v;
    long long w;
};

int main() {
    ios::sync_with_stdio(false);
    cin.tie(nullptr);
    int n, m, k;
    cin >> n >> m >> k;
    vector<long long> dist(n + 1, INF);
    for (int i = 0; i < k; i++) {
        int a;
        long long e;
        cin >> a >> e;
        dist[a] = e;                    // 虛擬起點連到入口，權重是入場費
    }
    vector<Edge> edges(m);
    vector<vector<int>> adj(n + 1);     // 只用來往下游擴散
    for (Edge &e : edges) {
        cin >> e.u >> e.v >> e.w;
        adj[e.u].push_back(e.v);
    }
    // 第 1 到 n - 1 輪：每輪把每條邊鬆弛一次
    for (int round = 1; round < n; round++) {
        bool changed = false;
        for (const Edge &e : edges) {
            if (dist[e.u] != INF && dist[e.u] + e.w < dist[e.v]) {
                dist[e.v] = dist[e.u] + e.w;
                changed = true;
            }
        }
        if (!changed) {                 // 這一輪沒有變化，之後也不會變
            break;
        }
    }
    // 第 n 輪還能鬆弛的邊，終點受負環影響；再沿邊擴散到下游
    vector<bool> neg(n + 1, false);
    queue<int> q;
    for (const Edge &e : edges) {
        if (dist[e.u] != INF && dist[e.u] + e.w < dist[e.v] && !neg[e.v]) {
            neg[e.v] = true;
            q.push(e.v);
        }
    }
    while (!q.empty()) {
        int u = q.front();
        q.pop();
        for (int v : adj[u]) {
            if (!neg[v]) {
                neg[v] = true;
                q.push(v);
            }
        }
    }
    for (int v = 1; v <= n; v++) {
        if (dist[v] == INF) {
            cout << "X\n";
        } else if (neg[v]) {
            cout << "-INF\n";
        } else {
            cout << dist[v] << '\n';
        }
    }
    return 0;
}
\end{example}

\iosample{
7 8 1 \\
1 5 \\
1 2 3 \\
2 3 -4 \\
3 2 2 \\
3 4 1 \\
1 5 2 \\
6 7 -1 \\
7 6 -1 \\
7 5 0 \\
}{
5 \\
-INF \\
-INF \\
-INF \\
7 \\
X \\
X \\
}

\begin{itemize}
\item 擴散用的 \texttt{adj} 只存鄰點，不用權重。 \texttt{neg} 同時當作 BFS 的
「看過了」，每個點最多放進佇列一次。
\item 有負環時，環上的 \texttt{dist} 在 $n - 1$ 輪裡會一直往下掉。 每一輪每條邊最
多鬆弛一次，所以 \texttt{dist} 對應的走法最多 $(n - 1)m$ 條邊，最負大約
$5000 \times 10^4 \times 10^9 = 5 \times 10^{16}$，\texttt{long long} 裝得下。
\item 範例 1 的房間 $1$ 是入口，入場費 $4$，但從房間 $4$ 進場走 $4 \to 5 \to 1$
只要 $1$。 入口的答案不一定是自己的入場費，所以 \texttt{dist[a]} 只是初值，之
後照樣可以被鬆弛。
\end{itemize}

\textbf{複雜度}

Bellman--Ford $O(nm)$，最多 $5000 \times 10^4 = 5 \times 10^7$ 次鬆弛；擴散
$O(n + m)$。 最大的測資在本書測試的環境下約 $0.14$ 秒。

\subsection{常見錯誤}

\begin{itemize}
\item \textbf{沒有往下游擴散}：只把第 $n$ 輪能鬆弛的點標成 \texttt{-INF}。 輸入
\begin{code}{text}{只標第 n 輪會錯的輸入}
4 5 1
1 0
1 2 0
2 3 0
3 2 -1
3 4 100
1 4 0
\end{code}
正確的輸出是 \texttt{0}、\texttt{-INF}、\texttt{-INF}、\texttt{-INF}，這樣寫會
印出 \texttt{0}、\texttt{-3}、\texttt{-INF}、\texttt{0}。 負環 $2 \to 3 \to 2$
每圈只少 $1$，房間 $4$ 另外有一條 $0$ 的路，第 $n$ 輪時 $3 \to 4$ 還差得遠；連
環上的房間 $2$ 都沒被標到，因為第 $n$ 輪剛好只有 $3$ 被改。
\item \textbf{沒檢查無限大}：拿掉 \texttt{dist[e.u] != INF}。 範例 2 會印出
\texttt{5} 和六個 \texttt{-INF}：走不到的 $6 \leftrightarrow 7$ 從
$10^{18}$ 開始往下掉，被當成走得到的負環，還連帶把房間 $5$ 拉成
\texttt{-INF}。
\item \textbf{用 Dijkstra}：有負權重，答案會錯（第~\ref{sec:sp-negative}~小節）；
有負環時更會一直繞下去。
\item \textbf{每個入口各跑一次}：$k$ 到 $n$ 時太慢。
\end{itemize}

\subsection{SPFA}
\label{sec:sp-spfa}

Bellman--Ford 每一輪都把所有邊看一遍，其實只有「上一輪 \texttt{dist} 變過的點」
往外的邊才可能鬆弛成功。 用一個佇列放這些點：拿出 $u$，鬆弛它的每條邊，成功時
如果 $v$ 不在佇列裡就把它放進去。 這個寫法常被叫做 \textbf{SPFA}（shortest path
faster algorithm），在很多隨機的圖上比 Bellman--Ford 快很多。

但它的最壞情況仍然是 $O(nm)$，而且這個最壞情況很容易構造出來：讓同一個點被反覆改
小、反覆放回佇列。 本章有幾題的測資特別放了這種圖。 例如把第~\ref{sec:ex-cart}~
小節的搬運車改用 SPFA，其中一筆測資有 $10^6$ 個狀態，SPFA 從佇列拿出了超過 $6
\times 10^8$ 次，在本書測試的環境下約 $2.4$ 秒；Dijkstra 只要 $0.01$ 秒。

所以：權重沒有負的，用 Dijkstra，不要用 SPFA。 有負權重時，SPFA 和
Bellman--Ford 的保證一樣是 $O(nm)$，只是平常比較快；拿它找負環還要另外記每個
點進出佇列的次數。 寫 Bellman--Ford 並加上「這一輪沒有變化就停」，通常就夠了。

\section{Floyd--Warshall 演算法}
\label{sec:sp-fw}

前面都是從一個起點出發。 題目要問很多對點之間的距離，而且點不多（幾百個）時，
可以一次算出每兩點之間的最短距離，叫\textbf{全點對最短路徑}（all-pairs shortest
paths）。 \textbf{Floyd--Warshall 演算法}（Floyd--Warshall algorithm）只有三層
迴圈，權重可以是負數（不能有負環）。

它是一個動態規劃。 把點編號 $1$ 到 $n$。 一條路徑除了起點和終點以外經過的點，叫
它的\textbf{中間點}。 定義
\[
d_k[i][j] = \text{從 } i \text{ 到 } j\text{、中間點都在 } \{1, 2, \ldots, k\} \text{ 裡的最短長度}.
\]
$d_0$ 不允許任何中間點，就是直接連的邊（有好幾條取最短，$d_0[i][i] = 0$，沒有邊是
無限大）；$d_n$ 沒有限制，就是答案。 從 $d_{k-1}$ 算 $d_k$ 時，多了點 $k$ 可以
用，一條中間點在 $\{1, \ldots, k\}$ 裡的最短路徑只有兩種：

\begin{itemize}
\item 沒經過 $k$：長度是 $d_{k-1}[i][j]$。
\item 經過 $k$：沒有負環時，最短路徑不必重複經過點，$k$ 只出現一次。 從 $i$ 到 $k$
那段、從 $k$ 到 $j$ 那段，中間點都在 $\{1, \ldots, k - 1\}$ 裡，長度是
$d_{k-1}[i][k] + d_{k-1}[k][j]$。
\end{itemize}

\begin{center}
\begin{tikzpicture}[x=1.5cm, y=1.1cm]
\node[spv] (i) at (0,0) {$i$};
\node[spvnow] (k) at (2,1.1) {$k$};
\node[spv] (j) at (4,0) {$j$};
\draw[spd, dashed] (i) to[bend left=15] node[font=\scriptsize, above left] {$d_{k-1}[i][k]$} (k);
\draw[spd, dashed] (k) to[bend left=15] node[font=\scriptsize, above right] {$d_{k-1}[k][j]$} (j);
\draw[spd, dashed] (i) to[bend right=12] node[font=\scriptsize, below] {$d_{k-1}[i][j]$} (j);
\end{tikzpicture}
\diagcaption{虛線是中間點都在 $\{1, \ldots, k - 1\}$ 裡的最短路徑。 $d_k[i][j]$
取下面那條和上面兩段相加的較小者。}
\label{fig:sp-fw}
\end{center}

所以 $d_k[i][j] = \min\bigl(d_{k-1}[i][j],\ d_{k-1}[i][k] + d_{k-1}[k][j]\bigr)$。
三維的表要 $n^3$ 格，但第 $k$ 層只用到第 $k - 1$ 層，而且第 $k$ 輪裡
$d[i][k]$、$d[k][j]$ 這兩種格子不會被改（拿 $d[i][k]$ 來說，候選值是
$d[i][k] + d[k][k] = d[i][k]$，因為沒有負環時 $d[k][k] = 0$），所以直接在同一張
二維表上更新就好：

\begin{code}{cpp}{Floyd--Warshall}
for (int k = 1; k <= n; k++) {          // 第 k 輪：開放 k 當中間點
    for (int i = 1; i <= n; i++) {
        for (int j = 1; j <= n; j++) {
            if (d[i][k] != INF && d[k][j] != INF) {
                d[i][j] = min(d[i][j], d[i][k] + d[k][j]);
            }
        }
    }
}
\end{code}

\subsection{$k$ 一定在最外層}

三層迴圈的順序不能換。 $k$ 是動態規劃的「階段」：第 $k$ 輪要用到第 $k - 1$ 輪
\textbf{整張表}都算完的值。 把 $k$ 放到最內層，變成「對每一對 $(i, j)$，試遍所有
$k$」，每一格只算一次，用的是當時表上的值，有些還沒算好。

例如有向圖只有三條邊 $1 \to 4$、$4 \to 3$、$3 \to 2$，權重都是 $1$。 照
$i, j, k$ 的順序，算 $d[1][2]$ 的時候 $d[1][3]$ 還沒算（要等 $j = 3$ 才算），試
$k = 3$ 時是 $\infty + 1$；試 $k = 4$ 時 $d[4][2]$ 也還是 $\infty$。 $d[1][2]$
就停在 $\infty$，正確答案是 $3$。 $k$ 放最外層時，第 $3$ 輪先把 $d[4][2]$ 算成
$2$，第 $4$ 輪再用它得到 $d[1][2] = 1 + 2 = 3$。

\subsection{使用時要注意的事}

\begin{itemize}
\item \textbf{複雜度}：$O(n^3)$ 時間、$O(n^2)$ 空間。 $n = 500$ 時是
$1.25 \times 10^8$ 次很簡單的運算，大約零點幾秒；$n$ 到幾千就不能用了。
\item \textbf{無限大}：$d[i][k] + d[k][j]$ 兩個都是無限大時會溢位。 先檢查，或
選一個「兩個相加也不會溢位」的無限大，例如 $10^{18}$（$2 \times 10^{18}$ 仍在
\texttt{long long} 的範圍內）。
\item \textbf{負環}：跑完之後 $d[i][i] < 0$，表示 $i$ 在某個負環上。 有負環時表
裡的值可能變得非常負，甚至溢位，題目可能有負環時要另外處理。
\item \textbf{印出路徑}：另外存 \texttt{nxt[i][j]}，表示從 $i$ 走到 $j$ 的第一步。
$d[i][j]$ 被 $k$ 改小時，令 \texttt{nxt[i][j] = nxt[i][k]}。
\end{itemize}

\subsection{例題：天橋候選方案的效益}
\label{sec:ex-footbridge}

\textbf{題意}（\akiref{evaluating-footbridge-proposals}{天橋候選方案的效益}）

\begin{itemize}
\item 市區有 $n$ 個路口、$m$ 條雙向道路，路網不一定連通。
\item 有 $p$ 個方案，第 $k$ 個是在 $u_k$、$v_k$ 之間加一條長度 $w_k$ 的雙向天
橋。 每個方案\textbf{各自獨立}評估：只加那一座，其他方案的天橋不存在。
\item 對每個方案，求有幾個無序點對 $\{i, j\}$ 的最短距離\textbf{嚴格}變短。 原
本走不到、加了之後走得到，也算變短。
\item $n \le 200$，$m \le 2 \times 10^4$，$p \le 2000$，長度都在 $1$ 到 $10^6$
之間。
\end{itemize}

範例 1 的路網是一條鏈 $1 - 2 - 3 - 4$（長度 $3$、$4$、$5$），路口 $5$ 沒有路。
方案 1 在 $4$、$1$ 之間加長度 $6$ 的天橋：$\{1, 4\}$ 從 $12$ 變成 $6$；
$\{2, 4\}$ 走 $2 \to 1 \to 4$ 是 $9$，和原本一樣，不算；答案是 $1$。 方案 2 把
$5$ 接到 $4$，$5$ 和其他四個路口原本都走不到，答案是 $4$。

\textbf{想法}

最直接的做法：每個方案把天橋加進去，重跑一次 Floyd--Warshall，再和原本的表比較。
每次 $n^3 = 8 \times 10^6$，乘上 $p = 2000$ 是 $1.6 \times 10^{10}$，太慢。

觀察：長度都是正的，一條最短路徑不會經過同一座天橋兩次（兩次之間是一圈，拿掉更
短）。 所以加了 $u$--$v$ 之後，$i$ 到 $j$ 的新最短距離只有三種可能：不走天橋、
走 $u \to v$ 方向、走 $v \to u$ 方向，
\[
d'[i][j] = \min\bigl(d[i][j],\ d[i][u] + w + d[v][j],\ d[i][v] + w + d[u][j]\bigr),
\]
其中 $d$ 是原本的表，天橋前後兩段都是原本路網裡的最短路徑。 Floyd--Warshall 只
要做一次，每個方案再把 $n(n - 1)/2$ 對點檢查一遍，$O(n^2)$。

這題的重點是把距離表當成算好一次、之後一直查的資料。

\textbf{程式}

\begin{example}[label=ex:sp-footbridge]{cpp}{(footbridge.cpp) 天橋候選方案的效益}
#include <algorithm>
#include <iostream>
#include <vector>
using namespace std;

const long long INF = 1e18;

int main() {
    ios::sync_with_stdio(false);
    cin.tie(nullptr);
    int n, m, p;
    cin >> n >> m >> p;
    vector<vector<long long>> d(n + 1, vector<long long>(n + 1, INF));
    for (int i = 1; i <= n; i++) {
        d[i][i] = 0;
    }
    for (int i = 0; i < m; i++) {
        int a, b;
        long long c;
        cin >> a >> b >> c;
        d[a][b] = min(d[a][b], c);           // 重複的道路取最短
        d[b][a] = min(d[b][a], c);
    }
    for (int k = 1; k <= n; k++) {           // k 一定在最外層
        for (int i = 1; i <= n; i++) {
            if (d[i][k] == INF) {
                continue;
            }
            for (int j = 1; j <= n; j++) {
                if (d[k][j] != INF && d[i][k] + d[k][j] < d[i][j]) {
                    d[i][j] = d[i][k] + d[k][j];
                }
            }
        }
    }
    for (int q = 0; q < p; q++) {
        int u, v;
        long long w;
        cin >> u >> v >> w;
        long long count = 0;
        for (int i = 1; i <= n; i++) {
            for (int j = i + 1; j <= n; j++) {
                long long viaUV = INF, viaVU = INF;   // 走天橋 u→v 或 v→u
                if (d[i][u] != INF && d[v][j] != INF) {
                    viaUV = d[i][u] + w + d[v][j];
                }
                if (d[i][v] != INF && d[u][j] != INF) {
                    viaVU = d[i][v] + w + d[u][j];
                }
                if (min(viaUV, viaVU) < d[i][j]) {
                    count++;
                }
            }
        }
        cout << count << '\n';
    }
    return 0;
}
\end{example}

\iosample{
5 3 4 \\
1 2 3 \\
2 3 4 \\
3 4 5 \\
4 1 6 \\
4 5 1 \\
2 3 4 \\
1 3 2 \\
}{
1 \\
4 \\
0 \\
2 \\
}

\begin{itemize}
\item \texttt{d[i][k] == INF} 時整列都不可能更新，先跳過，順便省一點時間。
\item 每個方案只讀 \texttt{d}、不改 \texttt{d}，方案之間互不影響。
\item 兩個方向都要試。 範例的方案 1 寫成 \texttt{4 1}，對 $\{1, 4\}$ 來說，有用
的是從 $1$ 往 $4$ 走天橋。
\end{itemize}

\textbf{複雜度}

Floyd--Warshall $O(n^3) = 8 \times 10^6$，加上方案 $O(pn^2)$，大約 $2000 \times
2 \times 10^4 = 4 \times 10^7$ 對。 最大的測資在本書測試的環境下約 $0.05$ 秒。

\subsection{常見錯誤}

\begin{itemize}
\item \textbf{只試一個方向}：只算 \texttt{d[i][u] + w + d[v][j]}。 範例會印出
\texttt{0 4 0 2}（每行一個），方案 1 漏掉了 $\{1, 4\}$。
\item \textbf{相等也算變短}：條件寫成 \texttt{<=}。 範例會印出 \texttt{2 4 4 2}：
方案 3 的天橋和原本的路一樣長，什麼都沒變短，卻算出 $4$ 對。
\item \textbf{把天橋真的加進表裡}：之後的方案會用到前面方案的天橋。
\item \textbf{無限大太小}：最長的真實距離約 $199 \times 10^6 \approx 2 \times
10^8$，無限大設成 $10^8$ 會把遠的點對當成走不到。 用 \texttt{INT\_MAX} 又會在
相加時溢位。
\end{itemize}

\section{怎麼選演算法}
\label{sec:sp-choose}

看到最短路徑的題目，先回答三個問題：要從幾個起點出發？ 權重的範圍是什麼？ $n$、
$m$ 多大？ 下表的「大約能處理」是以 1 秒、每秒幾億次簡單運算估的（銜接書第~10~
章 10.1 節的表）。

\begin{center}
\begin{tabular}{p{0.27\linewidth}p{0.24\linewidth}p{0.17\linewidth}p{0.2\linewidth}}
\toprule
圖的樣子 & 方法 & 時間 & 大約能處理 \\
\midrule
每條邊都一樣長 & BFS（第 1 本第~11~章） & $O(n + m)$ & $n + m$ 到 $10^7$ \\
權重只有 $0$、$1$ & 0/1 BFS & $O(n + m)$ & $n + m$ 到 $10^7$ \\
權重非負 & Dijkstra & $O((n + m) \log n)$ & $n + m$ 到 $10^6$ \\
沒有環（DAG），權重任意 & 拓樸順序 DP（第~6~章） & $O(n + m)$ & $n + m$ 到 $10^7$ \\
有負權重 & Bellman--Ford & $O(nm)$ & $nm$ 到 $10^8$ \\
每兩點之間 & Floyd--Warshall & $O(n^3)$ & $n$ 到 $500$ \\
每兩點之間，權重非負、$n$ 大 & 每個點各做一次 Dijkstra & $O(n(n + m) \log n)$ & 看 $n$ 和 $m$ \\
\bottomrule
\end{tabular}
\end{center}

幾個補充：

\begin{itemize}
\item \textbf{先看能不能用比較簡單的}。 權重都一樣時不要用 Dijkstra，DAG 上不要用
Bellman--Ford。
\item \textbf{很密的圖}：$m$ 接近 $n^2$ 時，Dijkstra 可以不用優先佇列，每次把
$n$ 個點掃一遍找最小的，$O(n^2)$，比 $O(n^2 \log n)$ 快一點。
\item \textbf{規則不只是「加總」}：有券、有朝向、要等紅綠燈，先想清楚狀態要包含
什麼（第~\ref{sec:sp-state}~節），再按權重選上面的方法。
\item \textbf{估最大距離}：最短路徑最多 $n - 1$ 條邊，最大距離大約是 $(n - 1)
\times$ 最大權重。 $n = 10^5$、權重 $10^9$ 時是 $10^{14}$，一定要
\texttt{long long}；無限大要比它大，而且加上一條邊不能溢位。
\end{itemize}

\section{本章常見錯誤}
\label{sec:sp-pitfalls}

前面各節的錯誤歸納成六類，寫最短路徑的題目時，交出去之前逐項檢查一次。

\begin{center}
\begin{tabular}{p{0.19\linewidth}p{0.47\linewidth}p{0.22\linewidth}}
\toprule
錯誤 & 例子 & 見 \\
\midrule
方法和權重不合 & 加權圖用 BFS；有負權重用 Dijkstra；權重非負卻用 SPFA &
第~\ref{sec:sp-why}、\ref{sec:sp-negative}、\ref{sec:sp-spfa}~節 \\
定案的時機 & 第一次看到就定案；0/1 BFS 當成普通 BFS & 第~\ref{sec:sp-dijkstra}、
\ref{sec:sp-01bfs}~節 \\
狀態或標籤不夠 & 搬運車只記格子；記起點編號時只比距離 & 第~\ref{sec:ex-cart}、
\ref{sec:sp-label}~節 \\
無限大與溢位 & 無限大設成 $10^9$；\texttt{INF + INF}；走不到的點被負權重拉低 &
第~\ref{sec:sp-dijkstra}、\ref{sec:ex-dungeon}、\ref{sec:sp-fw}~節 \\
負環的範圍 & 沒有往下游擴散；所有 \texttt{dist} 從 $0$ 開始 &
第~\ref{sec:sp-negcycle}、\ref{sec:ex-dungeon}~節 \\
迴圈的順序 & Floyd--Warshall 的 $k$ 不在最外層 & 第~\ref{sec:sp-fw}~節 \\
\bottomrule
\end{tabular}
\end{center}

\section{本章練習}
\label{sec:ch8-exercises}

\begin{drills}[label=drills:ch8]
\item 有向圖有八條邊：$1 \to 2$（$4$）、$1 \to 3$（$1$）、$3 \to 2$（$2$）、
$2 \to 4$（$1$）、$3 \to 4$（$6$）、$4 \to 5$（$3$）、$2 \to 5$（$7$）、
$3 \to 5$（$9$）。 從 $1$ 出發，照範例~\ref{ex:sp-dijkstra} 的寫法跑
Dijkstra。
\begin{enumerate}
\item 依序寫出從優先佇列拿出的每一筆紀錄，標出哪幾筆是過期的。
\item 寫出最後的 \texttt{dist} 和 \texttt{par}，以及到 $5$ 的最短路徑。
\end{enumerate}

\item 有人想讓 Dijkstra 處理負權重：先找出最小的權重 $-c$，把每條邊都加上 $c$，
權重就都不是負的了，跑完再換算回來。 給一個反例，說明這樣做為什麼不對。

\item 把範例~\ref{ex:sp-dijkstra} 改成同時算出到每個點的最短路徑\textbf{有幾
條}（權重都是正的）。
\begin{enumerate}
\item 寫出鬆弛時 \texttt{cnt[v]} 要怎麼更新。
\item 如果允許權重是 $0$，你的寫法還對嗎？ 給一個例子。
\end{enumerate}

\item \label{drill:sp-klimit} 題目要「最多經過 $K$ 條邊」的最短路徑，有人直接
跑 $K$ 輪 Bellman--Ford。
\begin{enumerate}
\item 給一個例子，說明跑完第 $1$ 輪，\texttt{dist} 可能已經用了不只一條邊。
\item 怎麼改才能讓第 $i$ 輪結束時，\texttt{dist} 剛好是「最多 $i$ 條邊」的答案？
\end{enumerate}

\item Floyd--Warshall 一開始把 $d[i][i]$ 設成無限大（不是 $0$），其他照舊，權重都
是正的。 跑完之後 $d[i][i]$ 代表什麼？

\item 下面每種情況，選一個最適合的方法，並估計時間。
\begin{enumerate}
\item $n = 2 \times 10^5$、$m = 5 \times 10^5$，權重 $1$ 到 $10^9$，問一個起點到
所有點。
\item $1000 \times 1000$ 的格子，走進空地代價 $0$、走進草叢代價 $1$。
\item $n = 300$，權重可能是負的但保證沒有負環，有 $10^5$ 筆「$a$ 到 $b$ 多遠」的
詢問。
\item $n = 2 \times 10^5$ 的 DAG，權重可能是負的，問一個起點到所有點。
\item $n = 3000$、$m = 6000$，權重可能是負的，問從起點走得到的地方有沒有負環。
\end{enumerate}
\end{drills}

\section{實作練習}
\label{sec:ch8-practice}

下面是 Aki-Judge 上和本章相關的題目，依小節分組，本章的三題例題也列在裡面。 難
度分四階：暖身直接套用本章的寫法；基本要自己建圖或加一點變化；實戰要先把題目轉
換成本章的形狀；挑戰會和其他章的技巧合在一起。

\begin{center}
\begin{tabular}{llc}
\toprule
小節 & 題目 & 難度 \\
\midrule
Dijkstra 與多個起點 & \akiref{earliest-parcel-delivery-through-sorting-hubs}{包裹轉運的最早送達時間} & 暖身 \\
 & \akiref{fire-crews-getting-ready}{消防分隊的整備時間} & 暖身 \\
 & \akiref{platforms-above-the-tide}{退潮才露出的平台} & 基本 \\
 & \akiref{staggered-traffic-lights}{錯開的紅綠燈} & 基本 \\
\midrule
反向圖 & \akiref{costume-room-before-the-hall}{先拿戲服再去禮堂} & 暖身 \\
 & \akiref{one-errand-on-the-way-home}{放學順路買東西} & 基本 \\
 & \akiref{check-in-before-the-auditorium}{先報到再進禮堂} & 實戰 \\
\midrule
狀態圖與建圖 & \akiref{the-cart-that-turns-slowly}{轉彎很慢的搬運車} & 基本（例題） \\
 & \akiref{drone-delivery-with-battery-swaps}{無人機送藥與換電池} & 實戰 \\
 & \akiref{ten-minutes-flat}{一律十分鐘的接駁車} & 實戰 \\
 & \akiref{the-gentlest-wheelchair-route}{坡度最緩的輪椅路線} & 挑戰 \\
\midrule
Bellman--Ford 與負環 & \akiref{cheapest-bike-trips-within-k-legs}{K 段內的最低點數與最少段數} & 暖身 \\
 & \akiref{cheapest-walk-in-a-multi-entrance-dungeon}{多入口地城的最低花費} & 基本（例題） \\
 & \akiref{rehearsal-start-windows}{校慶彩排的可排時段} & 實戰 \\
\midrule
Floyd--Warshall & \akiref{patrol-tasks-on-one-way-walkways}{單行步道的巡邏任務} & 暖身 \\
 & \akiref{evaluating-footbridge-proposals}{天橋候選方案的效益} & 基本（例題） \\
 & \akiref{flood-alert-threshold-on-a-typhoon-day}{颱風天的轉運站警戒門檻} & 實戰 \\
\bottomrule
\end{tabular}
\end{center}

\section{讀完你應該會}
\label{sec:ch8-checklist}

\begin{itemize}
\item 看到邊上有數字的圖，能說明 BFS 為什麼會錯，並依權重的範圍選 BFS、0/1 BFS、
Dijkstra、Bellman--Ford 或 Floyd--Warshall。
\item 能寫出用 \texttt{priority\_queue} 的 Dijkstra，處理過期的紀錄，用
\texttt{par} 印出一條最短路徑，並說得出為什麼有負權重時會錯。
\item 能用虛擬起點處理好幾個起點、用反向圖處理「每個點到同一個終點」，並用
\texttt{ds[x] + dt[x]} 回答「一定要經過某點」的問題。
\item 能把「用了幾張券、車頭朝哪裡」放進狀態，估出狀態數，在狀態圖上跑
Dijkstra。
\item 能用 Bellman--Ford 找出受負環影響的點，用 Floyd--Warshall 算出每兩點的距
離，並說明 $k$ 為什麼要在最外層。
\end{itemize}

\textbf{延伸閱讀}（書末會整理完整的參考資料）

\begin{itemize}
\item Antti Laaksonen，\textit{Competitive Programmer's Handbook}，第 13 章：
Bellman--Ford、Dijkstra、Floyd--Warshall 各一節，篇幅短，有手算的過程。
\item Steven Halim 等，\textit{Competitive Programming 4}，Book 1，4.4、4.5 節：
多源、反向圖、0/1 權重等變形整理得很完整，也比較了 Dijkstra 的兩種寫法和 SPFA。
\item Thomas H. Cormen 等，\textit{Introduction to Algorithms}，第 4 版，第 22
章（Bellman--Ford、Dijkstra 的完整證明，22.4 節用最短路徑解差距限制）與 23.2 節
（Floyd--Warshall）。
\item 劉汝佳，《算法競賽入門經典》第 2 版，11.3 節（最短路問題）。
\end{itemize}

\addcontentsline{toc}{section}{練習解答}
\section*{練習解答}

\begin{enumerate}
\item
\begin{enumerate}
\item 拿出的順序：$(0, 1)$、$(1, 3)$、$(3, 2)$、$(4, 2)$ 過期、$(4, 4)$、
$(7, 4)$ 過期、$(7, 5)$、$(10, 5)$ 過期。 處理 $1$ 時 $2$ 變 $4$、$3$ 變 $1$；
處理 $3$ 時 $2$ 改成 $3$、$4$ 變 $7$、$5$ 變 $10$；處理 $2$ 時 $4$ 改成 $4$（$5$
是 $3 + 7 = 10$，沒有更小）；處理 $4$ 時 $5$ 改成 $7$。 $(4, 2)$ 和 $(4, 4)$ 同
距離，\texttt{pair} 再比點的編號，所以 $(4, 2)$ 先出來。
\item \texttt{dist} $= 0, 3, 1, 4, 7$；\texttt{par[2]} $= 3$、\texttt{par[3]}
$= 1$、\texttt{par[4]} $= 2$、\texttt{par[5]} $= 4$。 到 $5$ 的最短路徑是
$1 \to 3 \to 2 \to 4 \to 5$，長度 $1 + 2 + 1 + 3 = 7$。
\end{enumerate}

\item 每條邊加 $c$，一條有 $\ell$ 條邊的路徑就多了 $\ell c$，邊多的路徑被罰得比較
重，比較的結果會改變。 例如 $1 \to 2$（$-2$）、$2 \to 3$（$-2$）、$1 \to 3$
（$-3$），從 $1$ 到 $3$ 最短是 $1 \to 2 \to 3$ 的 $-4$。 每條加 $3$ 之後變成
$1$、$1$、$0$，最短變成直接走 $1 \to 3$ 的 $0$，換算回來是 $-3$，錯了。

\item
\begin{enumerate}
\item 鬆弛時，\texttt{d + w < dist[v]} 表示找到更短的，舊的條數作廢：
\texttt{cnt[v] = cnt[u]}；\texttt{d + w == dist[v]} 表示一樣短的另一批路：
\texttt{cnt[v] += cnt[u]}。 \texttt{cnt[s] = 1}。 權重都是正的時候，$v$ 的每個
「最短路徑上的前一站」$u$ 都滿足 $\delta(u) < \delta(v)$，一定比 $v$ 早處理，處
理 $u$ 時 \texttt{cnt[u]} 已經是最後的值。 條數可能很大，題目通常會要求取餘數。
\item 不一定對。 例如 $1 \to 2$（$1$）、$1 \to 3$（$1$）、$3 \to 2$（$0$）、
$2 \to 4$（$1$）。 到 $4$ 的最短路徑有 $1 \to 2 \to 4$ 和 $1 \to 3 \to 2 \to 4$
兩條，長度都是 $2$。 但 $(1, 2)$ 比 $(1, 3)$ 先拿出來，處理 $2$ 時
\texttt{cnt[2]} 還是 $1$，傳給 $4$ 的是 $1$；之後處理 $3$ 才把 \texttt{cnt[2]}
加成 $2$，已經來不及，印出 $1$。 權重 $0$ 讓「前一站」和自己距離相同，處理順序
就不保證了。
\end{enumerate}

\item
\begin{enumerate}
\item 邊的順序是 $1 \to 2$（$1$）、$2 \to 3$（$1$）。 第 1 輪先把 \texttt{dist[2]}
改成 $1$，接著 $2 \to 3$ 馬上用到它，\texttt{dist[3]} $= 2$，已經用了兩條邊。
\item 每一輪開始前把 \texttt{dist} 複製一份 \texttt{prev}，鬆弛時一律用
\texttt{prev[e.u] + e.w} 去更新 \texttt{dist[e.v]}。 這樣第 $i$ 輪只會在第
$i - 1$ 輪的結果後面多接一條邊，跑 $K$ 輪就是「最多 $K$ 條邊」的答案。 因為限制
了邊數，有負環也不會無限下降。
\end{enumerate}

\item $d[i][i]$ 是\textbf{經過 $i$ 的最短的環}的長度：從 $i$ 出發、至少走一條邊、
回到 $i$ 的最短長度，沒有這種環就仍是無限大。 推導和正文一樣，只是 $d_0[i][i]$ 改
成 $i$ 到 $i$ 的自環長度（沒有就是無限大），「什麼都不走」的 $0$ 不算，之後 $d[i][i]$
只能由「$i$ 走到 $k$ 再走回 $i$」得到。 權重都是正的，這個環不會被算成 $0$。

\item
\begin{enumerate}
\item Dijkstra，大約 $7 \times 10^5 \times 18 \approx 10^7$ 次堆積操作。
\item 0/1 BFS，$10^6$ 個點、約 $4 \times 10^6$ 條邊，$O(n + m)$。
\item Floyd--Warshall，$300^3 = 2.7 \times 10^7$；之後每筆詢問查表 $O(1)$。
\item 本書第~6~章：先排拓樸順序，再照順序鬆弛，$O(n + m)$。
\item Bellman--Ford 跑 $n$ 輪，$3000 \times 6000 = 1.8 \times 10^7$ 次鬆弛，第 $n$
輪還能鬆弛就有走得到的負環。
\end{enumerate}
\end{enumerate}

\chapter{最小生成樹}
\label{ch:mst}

% 本章的圖共用的樣式（點、強調的點、邊、選中的邊、要拿掉的邊、權重標籤）
\tikzset{
    mstv/.style={circle, draw=black!60, thick, minimum size=0.7cm, inner sep=0pt, font=\small},
    mstvhl/.style={mstv, draw=memStroke, fill=memFill},
    mstvblue/.style={mstv, draw=coreStroke, fill=coreFill},
    mste/.style={thick, draw=black!35},
    mstehl/.style={line width=1.6pt, draw=memStroke},
    mstenew/.style={line width=1.6pt, draw=memStroke, dashed},
    mstebad/.style={line width=1.6pt, draw=optStroke, dashed},
    mstw/.style={font=\scriptsize, text=black!70, fill=white, inner sep=1pt},
    mstnote/.style={font=\scriptsize, text=coreStroke},
}

工業區裡有六座廠房，管委會要鋪光纖把它們連成一個網路。 下圖是可以施工的九條
路線，邊上的數字是造價（萬元）。 訊號可以經過別的廠房轉送，所以不需要每兩座
之間都直接拉線，只要任兩座廠房之間有一條由光纖接起來的路就好。 最少要花多少
錢？ 往下讀之前，先在圖上挑一次，算出你的總造價。

\begin{center}
\begin{tikzpicture}[x=1.5cm, y=1.3cm]
\node[mstv] (1) at (0,1) {1};
\node[mstv] (2) at (1.6,2) {2};
\node[mstv] (6) at (1.6,0) {6};
\node[mstv] (3) at (3.4,2) {3};
\node[mstv] (5) at (3.4,0) {5};
\node[mstv] (4) at (5,1) {4};
\draw[mste] (1) -- node[mstw, above left] {4} (2);
\draw[mste] (1) -- node[mstw, below left] {2} (6);
\draw[mste] (2) -- node[mstw, left] {5} (6);
\draw[mste] (2) -- node[mstw, above] {6} (3);
\draw[mste] (6) -- node[mstw, below] {7} (5);
\draw[mste] (3) -- node[mstw, right] {3} (5);
\draw[mste] (3) -- node[mstw, above right] {8} (4);
\draw[mste] (5) -- node[mstw, below right] {9} (4);
\draw[mste] (6) -- node[mstw, pos=0.35] {10} (3);
\end{tikzpicture}
\diagcaption{六座廠房和九條可以施工的光纖路線，邊上的數字是造價。}
\label{fig:mst-factories}
\end{center}

六座廠房至少要五條線才連得起來，而且用五條剛好連通時，線路不會繞成圈：多出一
個圈，就表示圈上有一條線拿掉之後大家仍然相通，那條線是白花的錢。 所以答案一定
是一棵連到全部六個點的樹。 從九條裡挑五條有 $126$ 種挑法，其中 $55$ 種剛好連
成一棵樹，最便宜的那一種花 $23$ 萬元。 你挑的是 $23$ 嗎？

這一章講怎麼不靠枚舉，直接找出這棵最便宜的樹：\textbf{Kruskal 演算法}把邊由便宜到
貴排好，能接就接；\textbf{Prim 演算法}從一個點長出一棵樹，每次接上最便宜的往外的
邊。 兩者都是貪心，理由是同一個交換論證。 之後再看三種變形：總和要最大、圖要自己
建、路上最窄的一段要盡量寬（瓶頸問題）。

這一章用到：本書第~7~章的並查集（直接沿用那一章的 \texttt{find} 和
\texttt{unite}，不再重講）；第~8~章的 Dijkstra 演算法和
\texttt{priority\_queue}；《C++ 演算法解題 1：APCS 實作題核心》（以下簡稱第 1 本）
第~7~章的貪心和交換論證；第 1 本第~11~章的邊列表、鄰接串列和 BFS；第 1 本第~13~章
「$n$ 個點的樹剛好有 $n - 1$ 條邊」。

\section{生成樹}
\label{sec:mst-tree}

這一章的圖都是無向圖，每條邊有一個\textbf{權重}（weight），可能代表造價、距離或
時間。 同一對點之間可以有好幾條邊（重邊），權重不一定相同。

一張連通的圖有 $n$ 個點。 從它的邊裡挑出一些，如果挑出來的邊能讓全部 $n$ 個點
彼此相通，而且沒有環，這些邊就叫做這張圖的一棵\textbf{生成樹}（spanning tree）。
「生成」指的是它碰到了每一個點。 第 1 本第~13~章說過，$n$ 個點連通又沒有環，剛好
是 $n - 1$ 條邊，所以生成樹一定是 $n - 1$ 條邊；反過來，挑了 $n - 1$ 條邊又讓全
部的點相通，一定沒有環。

生成樹的權重是它的邊的權重總和。 所有生成樹裡權重最小的，叫\textbf{最小生成樹}
（minimum spanning tree，常簡寫成 MST）。 章首的問題就是求最小生成樹，答案是
$23$。 最小生成樹可能不只一棵，例如權重相同的邊可以互換，但它們的總權重都一樣。

題目通常不會寫「求最小生成樹」，而是寫「選一些路線，讓任兩個地點都能到達，總造價
最小」或「剛好選 $n - 1$ 條，讓全部連通」。 權重都是正的時候，兩種說法是同一件事。
權重可以是負的時候就不一樣：第一種說法裡負的邊選越多越省，答案可能有環；第二種限
定了 $n - 1$ 條，仍然是生成樹。 這一章的例題和題單都屬於第二種，或權重都不是負的。

\subsection{和最短路徑不一樣}
\label{sec:mst-vs-sp}

第~8~章的 Dijkstra 從起點 $s$ 出發，替每個點找出最短的路。 每個點記下最短路上的
前一個點，這 $n - 1$ 條「前一步」的邊合起來也是一棵生成樹，叫\textbf{最短路徑樹}
（shortest-path tree）。 它和最小生成樹看起來很像，都是一棵樹、都在挑便宜的邊，
但兩者照顧的東西不同：

\begin{center}
\begin{tikzpicture}[x=1.3cm, y=1.1cm]
\begin{scope}
\node[font=\small] at (1,2.8) {最短路徑樹：總權重 $18$};
\node[mstvhl] (s) at (0,1) {$s$};
\node[mstv] (a) at (2,2) {$a$};
\node[mstv] (b) at (2,1) {$b$};
\node[mstv] (c) at (2,0) {$c$};
\draw[mstehl] (s) -- node[mstw, above left] {5} (a);
\draw[mstehl] (s) -- node[mstw, above, pos=0.6] {6} (b);
\draw[mstehl] (s) -- node[mstw, below left] {7} (c);
\draw[mste] (a) -- node[mstw, right] {2} (b);
\draw[mste] (b) -- node[mstw, right] {2} (c);
\node[mstnote, right=0.25cm] at (a.east) {距離 5};
\node[mstnote, right=0.25cm] at (b.east) {距離 6};
\node[mstnote, right=0.25cm] at (c.east) {距離 7};
\end{scope}
\begin{scope}[xshift=6cm]
\node[font=\small] at (1,2.8) {最小生成樹：總權重 $9$};
\node[mstvhl] (s) at (0,1) {$s$};
\node[mstv] (a) at (2,2) {$a$};
\node[mstv] (b) at (2,1) {$b$};
\node[mstv] (c) at (2,0) {$c$};
\draw[mstehl] (s) -- node[mstw, above left] {5} (a);
\draw[mste] (s) -- node[mstw, above, pos=0.6] {6} (b);
\draw[mste] (s) -- node[mstw, below left] {7} (c);
\draw[mstehl] (a) -- node[mstw, right] {2} (b);
\draw[mstehl] (b) -- node[mstw, right] {2} (c);
\node[mstnote, right=0.25cm] at (a.east) {距離 5};
\node[mstnote, right=0.25cm] at (b.east) {距離 7};
\node[mstnote, right=0.25cm] at (c.east) {距離 9};
\end{scope}
\end{tikzpicture}
\diagcaption{同一張圖上，從 $s$ 出發的最短路徑樹（左）和最小生成樹（右）。 粗線是
樹上的邊，藍字是沿著樹從 $s$ 走過去的距離。}
\label{fig:mst-vs-spt}
\end{center}

左邊的最短路徑樹讓\textbf{每個點到 $s$ 的距離}都最短，總權重 $18$；右邊的最小生成
樹讓\textbf{邊的總和}最小，只要 $9$，代價是 $c$ 沿著樹走離 $s$ 有 $9$，直接走只要
$7$。 想像 $s$ 是機房：要省鋪線的錢選右邊，要每台機器到機房的延遲最小選左邊。 另
外，最小生成樹沒有起點，最短路徑樹換一個起點就整棵不同。

暴力解是枚舉所有 $n - 1$ 條邊的組合，用第~7~章的並查集檢查有沒有環，留下權重最小
的。 章首的圖只有 $\binom{9}{5} = 126$ 種組合；$m = 30$、$n = 16$ 時已經大約
$1.6 \times 10^8$ 種，$m$ 到 $2 \times 10^5$ 完全不可能。 它寫起來不難，留著當對拍用。

\section{Kruskal 演算法}
\label{sec:kruskal}

最便宜的線應該盡量用。 照這個直覺，先試兩條看起來合理的規則，都拿章首的圖來試。

\begin{itemize}
\item \textbf{規則 A：直接拿最便宜的 $n - 1$ 條。} 最便宜的五條是 $2, 3, 4, 5, 6$，
也就是 1--6、3--5、1--2、2--6、2--3。 其中 1--2、2--6、1--6 繞成一個三角形，
廠房 4 卻一條線都沒接到。
\item \textbf{規則 B：每座廠房各自選一條連出去最便宜的線。} 1 選 1--6、2 選
1--2、3 和 5 都選 3--5、4 選 3--4、6 選 1--6，合起來是四條線，分成
$\{1, 2, 6\}$ 和 $\{3, 4, 5\}$ 兩塊，接不起來。
\end{itemize}

規則 A 的問題出在三角形：2--6 這條線接上的兩座廠房，早就經由 1 相通了，它沒有
幫上任何忙，還佔掉了一個名額。 把這一點修好，就是 Kruskal 演算法：

\begin{itemize}
\item 把所有的邊依權重由小到大排好。
\item 依序看每一條邊 $(u, v)$：如果 $u$ 和 $v$ 目前還不相通，就選它；如果早就相
通，選它只會繞成環，跳過。
\item 選滿 $n - 1$ 條就完成了。
\end{itemize}

「目前相不相通」正是第~7~章並查集在做的事：\texttt{unite(u, v)} 回傳
\texttt{true} 表示兩群原本分開、現在合併了，這條邊就要選；回傳 \texttt{false}
表示早就同一群，跳過。

\subsection{在章首的圖上跑一次}

九條邊排序之後依序處理，每一列是一條邊：

\begin{center}
\begin{tabular}{cclll}
\toprule
邊 & 權重 & 兩端 & 結果 & 處理完的群 \\
\midrule
1--6 & 2 & 不同群 & 選 & $\{1, 6\}$ $\{2\}$ $\{3\}$ $\{4\}$ $\{5\}$ \\
3--5 & 3 & 不同群 & 選 & $\{1, 6\}$ $\{2\}$ $\{3, 5\}$ $\{4\}$ \\
1--2 & 4 & 不同群 & 選 & $\{1, 2, 6\}$ $\{3, 5\}$ $\{4\}$ \\
2--6 & 5 & 同一群 & 跳過 & 不變 \\
2--3 & 6 & 不同群 & 選 & $\{1, 2, 3, 5, 6\}$ $\{4\}$ \\
6--5 & 7 & 同一群 & 跳過 & 不變 \\
3--4 & 8 & 不同群 & 選 & $\{1, 2, 3, 4, 5, 6\}$ \\
5--4 & 9 & 同一群 & 跳過 & 不變 \\
6--3 & 10 & 同一群 & 跳過 & 不變 \\
\bottomrule
\end{tabular}
\end{center}

選到 1--6、3--5、1--2、2--3、3--4 五條，總和 $2 + 3 + 4 + 6 + 8 = 23$，畫出來就是
下一節圖~\ref{fig:mst-cut} 的右半邊。 注意 6--5 雖然只要 $7$，比 3--4 的 $8$ 便宜，
卻被跳過了：它要接的兩座廠房早就經由 6--1--2--3--5 相通。

\subsection{程式}

輸入第一行是 $n$、$m$，接著 $m$ 行「$u$ $v$ $w$」。 邊存成邊列表（第 1 本第~11~
章），因為 Kruskal 只需要把邊排序，用不到鄰接串列。 程式印出選到的每一條邊和總
和；圖不連通時印出 \texttt{disconnected}。

\begin{example}[label=ex:mst-kruskal]{cpp}{(kruskal.cpp) Kruskal 演算法}
#include <algorithm>
#include <iostream>
#include <utility>
#include <vector>
using namespace std;

struct Edge {
    int u, v;
    long long w;
};

vector<int> parent, sz;

int find(int x) {
    return parent[x] == x ? x : parent[x] = find(parent[x]);
}

bool unite(int a, int b) {
    a = find(a);
    b = find(b);
    if (a == b) {
        return false;
    }
    if (sz[a] > sz[b]) {
        swap(a, b);
    }
    parent[a] = b;
    sz[b] += sz[a];
    return true;
}

int main() {
    int n, m;
    cin >> n >> m;
    vector<Edge> edges(m);
    for (Edge &e : edges) {
        cin >> e.u >> e.v >> e.w;
    }
    sort(edges.begin(), edges.end(),
         [](const Edge &a, const Edge &b) { return a.w < b.w; });
    parent.resize(n + 1);
    sz.assign(n + 1, 1);
    for (int i = 1; i <= n; i++) {
        parent[i] = i;
    }
    long long total = 0;
    int used = 0;
    for (const Edge &e : edges) {
        if (unite(e.u, e.v)) {         // 兩端原本不連通：選
            total += e.w;
            used++;
            cout << e.u << ' ' << e.v << ' ' << e.w << '\n';
        }                              // 否則會繞成環：跳過
    }
    if (used == n - 1) {
        cout << "total " << total << '\n';
    } else {
        cout << "disconnected\n";
    }
    return 0;
}
\end{example}

\iosample{
6 9 \\
1 2 4 \\
1 6 2 \\
2 6 5 \\
2 3 6 \\
6 5 7 \\
3 5 3 \\
3 4 8 \\
5 4 9 \\
6 3 10 \\
}{
1 6 2 \\
3 5 3 \\
1 2 4 \\
2 3 6 \\
3 4 8 \\
total 23 \\
}

\begin{itemize}
\item \texttt{find} 和 \texttt{unite} 和第~7~章修路的程式（\texttt{roads.cpp}）一字
不差：按大小合併加路徑壓縮。 這一章每一支程式都用同一份。
\item 排序用 lambda 只比權重（銜接書第~9~章），權重相同的邊誰先都可以。
\item \texttt{total} 用 \texttt{long long}。 $n$ 到 $2 \times 10^5$、權重到
$10^9$ 時，總和可以到 $2 \times 10^{14}$。
\item \texttt{used} 數選了幾條。 等於 $n - 1$ 時其實可以直接 \texttt{break}，後面
的邊都會被跳過；不寫也一樣對，只是多跑幾次 \texttt{find}。
\end{itemize}

\textbf{複雜度}：排序 $O(m \log m)$，之後 $m$ 次 \texttt{unite}，每次幾乎是常數
時間（第~7~章）。 整體是 $O(m \log m)$，時間幾乎都花在排序上。

\subsection{圖不連通的時候}
\label{sec:mst-forest}

如果圖本身不連通，就沒有生成樹。 Kruskal 照樣可以跑完：每一塊連通元件各自連成
一棵樹，合起來是一片森林，叫做\textbf{生成森林}（spanning forest）；Kruskal 得到
的是總權重最小的那一片。 圖有 $c$ 個連通元件時，森林有 $n - c$ 條邊，所以只要看
選了幾條：少於 $n - 1$ 條就表示不連通。

題目對不連通的要求各有不同：輸出 $-1$、輸出森林的總權重，或問最少還要加幾條線（答
案是 $c - 1$）。 \texttt{used} 和 $c = n - \texttt{used}$ 都在手上，照題目要的輸出就
好。

\subsection{什麼時候用、什麼時候不能用}

\begin{itemize}
\item 適合：「選一些邊讓所有點連通，總權重最小」或「剛好選 $n - 1$ 條讓全部連
通」，邊直接給在輸入裡，$m$ 到 $2 \times 10^5$ 左右。
\item 不能用：\textbf{有向圖}（例如沿著單行道從總部到每個地點），Kruskal 會選出走
不通的方向，這是另一個問題，超出本書範圍。
\item 不能用：\textbf{只要連通其中幾個點}，其他點可經過可不經過。 這叫 Steiner 樹
問題，目前沒有已知的有效率解法，資料量通常很小，只能枚舉。
\end{itemize}

\section{為什麼 Kruskal 是對的}
\label{sec:mst-cut}

第 1 本第~7~章看過很多「看起來合理、其實會錯」的貪心，例如郵票先貼大面額。
Kruskal 每一步只看眼前最便宜的邊，選了就不回頭，也需要一個理由。 這一節用第 1 本
第~7~章的交換論證證明它。 先介紹一個工具。

\subsection{切割}

把圖上的點分成兩邊，一邊叫 $S$，另一邊是其他的點，這樣的分法叫一個\textbf{切割}
（cut）。 一條邊的兩端如果分在不同邊，就說它\textbf{跨過}這個切割。 下圖把章首的
圖切成 $S = \{1, 2, 6\}$ 和 $\{3, 4, 5\}$，跨過去的邊有 2--3、6--5、6--3 三條。

任何一棵生成樹都必須至少用一條跨過的邊，否則 $S$ 裡的點到不了另一邊。 直覺上，
要過去，就從跨過的邊裡挑最便宜的。 這個直覺是對的：

\begin{center}
\fbox{\parbox{0.88\linewidth}{\textbf{切割性質}（cut property）：任取一個切割，跨過
它的邊裡權重最小的那一條 $e$，一定包含在某一棵最小生成樹裡。}}
\end{center}

證明用交換論證。 任取一棵最小生成樹 $T$。 如果 $T$ 裡有 $e$，就證完了。 如果沒
有，設 $e$ 連接 $u$ 和 $v$，$u$ 在 $S$ 裡、$v$ 不在。

\begin{itemize}
\item \textbf{找一條可以換掉的邊。} $T$ 是生成樹，$u$ 到 $v$ 在 $T$ 上有一條路。
這條路從 $S$ 出發、結束在 $S$ 外面，中途一定有某一條邊 $f$ 跨過切割。
\item \textbf{換完仍然合法。} 把 $T$ 加上 $e$、拿掉 $f$，得到 $T'$。 $T'$ 還是
$n - 1$ 條邊。 原本經過 $f$ 的路，都可以改走「$T$ 上那條路剩下的部分，加上
$e$」繞過去，所以 $T'$ 仍然讓全部的點相通，是一棵生成樹。
\item \textbf{換完沒有變差。} $e$ 是跨過切割的邊裡最小的，$f$ 也跨過切割，所以
$e$ 的權重 $\le$ $f$ 的權重，$T'$ 的總權重 $\le$ $T$ 的總權重。 $T$ 已經是最小，
所以 $T'$ 也是最小生成樹，而且包含 $e$。
\end{itemize}

下圖是一個實際的例子。 $T$ 是一棵沒有用到 2--3 的生成樹，總權重 $24$。 加上
2--3 之後，2--3--5--6--1--2 繞成一個環，環上另一條跨過切割的邊是 6--5，把它拿掉，
總權重變成 $24 + 6 - 7 = 23$。

\begin{center}
\begin{tikzpicture}[x=1.05cm, y=1.1cm]
\begin{scope}
\fill[coreFill, rounded corners=8pt] (-0.5,-0.45) rectangle (2.1,2.45);
\node[font=\small, coreStroke] at (0.2,2.7) {$S$};
\node[font=\small] at (2.5,3.1) {$T$，總權重 $24$；加上 $e$};
\node[mstv] (1) at (0,1) {1};
\node[mstv] (2) at (1.6,2) {2};
\node[mstv] (6) at (1.6,0) {6};
\node[mstv] (3) at (3.4,2) {3};
\node[mstv] (5) at (3.4,0) {5};
\node[mstv] (4) at (5,1) {4};
\draw[mstehl] (1) -- node[mstw, above left] {4} (2);
\draw[mstehl] (1) -- node[mstw, below left] {2} (6);
\draw[mstebad] (6) -- node[mstw, below] {$f$: 7} (5);
\draw[mstehl] (3) -- node[mstw, right] {3} (5);
\draw[mstehl] (3) -- node[mstw, above right] {8} (4);
\draw[mstenew] (2) -- node[mstw, above] {$e$: 6} (3);
\end{scope}
\begin{scope}[xshift=7cm]
\fill[coreFill, rounded corners=8pt] (-0.5,-0.45) rectangle (2.1,2.45);
\node[font=\small, coreStroke] at (0.2,2.7) {$S$};
\node[font=\small] at (2.5,3.1) {$T'$，總權重 $23$};
\node[mstv] (1) at (0,1) {1};
\node[mstv] (2) at (1.6,2) {2};
\node[mstv] (6) at (1.6,0) {6};
\node[mstv] (3) at (3.4,2) {3};
\node[mstv] (5) at (3.4,0) {5};
\node[mstv] (4) at (5,1) {4};
\draw[mstehl] (1) -- node[mstw, above left] {4} (2);
\draw[mstehl] (1) -- node[mstw, below left] {2} (6);
\draw[mstehl] (3) -- node[mstw, right] {3} (5);
\draw[mstehl] (3) -- node[mstw, above right] {8} (4);
\draw[mstehl] (2) -- node[mstw, above] {6} (3);
\end{scope}
\end{tikzpicture}
\diagcaption{切割 $S = \{1, 2, 6\}$（藍底）。 左：生成樹 $T$ 加上跨過切割最便宜的
邊 $e$（虛線），形成一個環，環上另一條跨過切割的邊是 $f$（紅色虛線）。 右：拿掉
$f$ 之後的 $T'$。}
\label{fig:mst-cut}
\end{center}

\subsection{Kruskal 每一步都選了切割上最便宜的邊}

現在把切割性質套到 Kruskal 上。 要證明的是下面這句話，對 Kruskal 的每一步都成立：

\begin{center}
目前選到的邊，全部包含在某一棵最小生成樹裡。
\end{center}

一開始一條邊都還沒選，當然成立。 假設選了一些邊之後還成立，設這些邊包含在最小
生成樹 $T$ 裡。 Kruskal 接著選了邊 $e = (u, v)$，取 $S$ 是 $u$ 目前所在的那一群
（並查集裡的那一群），$v$ 不在 $S$ 裡。 這個切割有兩個特點：

\begin{itemize}
\item 已經選的邊都不跨過它。 已選的邊兩端一定在同一群，而 $S$ 剛好是一整群。
\item $e$ 是跨過它的邊裡最便宜的。 理由在下面的框裡。
\end{itemize}

接著照切割性質的證明做交換：$T$ 加上 $e$ 形成的環上，有另一條跨過切割的邊 $f$，
換掉它得到 $T'$。 因為已選的邊都不跨過切割，被換掉的 $f$ 不會是已選的邊，所以
$T'$ 仍然包含所有已選的邊，再加上 $e$。 這句話在選了 $e$ 之後仍然成立。

一直到選滿 $n - 1$ 條，選到的邊包含在某一棵最小生成樹裡，而它們本身就有
$n - 1$ 條，所以它們就是那一棵最小生成樹。 這和第 1 本第~7~章交換論證的三步對得
起來：任取一個最佳解；把和貪心不一樣的地方換成貪心的選擇，仍然合法、沒有變差；
每換一次，最佳解和貪心相同的邊就多一條，所以會停。

\begin{whybox}[$e$ 是跨過 $S$ 最便宜的邊]
看任何一條比 $e$ 便宜、而且跨過 $S$ 的邊 $g = (x, y)$。 Kruskal 依權重由小到大處
理，$g$ 比 $e$ 先處理過。 處理 $g$ 的時候只有兩種可能：

\begin{itemize}
\item $x$、$y$ 不同群：Kruskal 會選 $g$，兩群合併，$x$、$y$ 從此在同一群。
\item $x$、$y$ 已經同群：群只會合併、不會拆開，從此也一直在同一群。
\end{itemize}

兩種情況下，現在 $x$ 和 $y$ 都在同一群，不可能一個在 $S$ 裡、一個在外面，和「$g$
跨過 $S$」矛盾。 所以跨過 $S$ 的邊沒有比 $e$ 便宜的。 和 $e$ 一樣便宜的邊可能有，
這不影響證明：切割性質只要求 $e$ 不比其他跨過的邊貴。 這也說明了權重相同的邊誰
先處理都可以。
\end{whybox}

\subsection{環性質}

切割性質還有一個反過來的版本，第~\ref{sec:mst-build}~節建圖時會用到：

\begin{center}
\fbox{\parbox{0.88\linewidth}{\textbf{環性質}（cycle property）：如果一條邊 $h$ 在某
個環上，而且比環上其他每一條邊都\textbf{嚴格}重，$h$ 不會出現在任何一棵最小生成
樹裡。}}
\end{center}

理由也是交換：假設某棵最小生成樹 $T$ 用了 $h$。 把 $h$ 拿掉，$T$ 分成兩塊。 那個
環從 $h$ 的一端繞回另一端，路上一定有另一條邊 $g$ 連接這兩塊。 用 $g$ 換掉 $h$ 得
到的仍然是生成樹，而 $g$ 比 $h$ 輕，總權重嚴格變小，和 $T$ 最小矛盾。

在章首的圖上，6--3（權重 $10$）在環 6--3--5--6 上，比 3--5（$3$）和 6--5（$7$）都
重，所以不管怎麼選，最小生成樹都不會用它。 Kruskal 跳過的每一條邊，其實都是這種
「環上最重」的邊：它被跳過時兩端已經經由更便宜的邊相通，那條路加上它就是一個環。

如果所有的邊權重都不相同，用環性質可以證明最小生成樹只有一棵（本章練習第~2~
題）。 有相同權重時可能有好幾棵，但總權重都一樣。

\section{Prim 演算法}
\label{sec:prim}

Kruskal 一次看整張圖，同時長出好幾小塊，最後才接成一棵。 Prim 演算法換一個方向：
從一個點開始，只維護\textbf{一棵}樹，每一步把「一端在樹上、另一端還不在」的邊裡
最便宜的那條接上，樹就多一個點。 接了 $n - 1$ 次，全部的點都在樹上。

它為什麼對？ 每一步的 $S$ 取「目前在樹上的點」，要接的邊正是跨過這個切割最便宜
的那一條，已經選的邊都在 $S$ 裡面、不跨過切割，所以上一節的論證整段可以照搬。

\subsection{和 Dijkstra 長得很像}

要很快找出「往外最便宜的邊」，就把往外的邊放進 \texttt{priority\_queue}。 寫出
來和第~8~章的 Dijkstra 幾乎一樣，只差在堆裡放的數字：從 $u$ 往 $v$ 的邊，Dijkstra
放 $\mathrm{dist}[u] + w$（經過 $u$ 到 $v$ 的距離），Prim 只放 $w$（把 $v$ 接到樹上
的費用）。 只差這一點，結果就是圖~\ref{fig:mst-vs-spt} 左右兩棵不同的樹。 寫 Prim 時
最常見的錯誤，就是順手把 Dijkstra 的 \texttt{dist[u] + w} 抄了過來。

章首的圖從點 1 開始，堆裡的元素寫成（費用，點）：

\begin{center}
\begin{tabular}{clll}
\toprule
步 & 從堆裡拿出 & 結果 & 新放進堆的 \\
\midrule
1 & (0, 1) & 1 上樹 & (4, 2)、(2, 6) \\
2 & (2, 6) & 6 上樹，費用 2 & (5, 2)、(7, 5)、(10, 3) \\
3 & (4, 2) & 2 上樹，費用 4 & (6, 3) \\
4 & (5, 2) & 2 已經在樹上，丟掉 & \\
5 & (6, 3) & 3 上樹，費用 6 & (3, 5)、(8, 4) \\
6 & (3, 5) & 5 上樹，費用 3 & (9, 4) \\
7 & (7, 5) & 5 已經在樹上，丟掉 & \\
8 & (8, 4) & 4 上樹，費用 8 & \\
\bottomrule
\end{tabular}
\end{center}

之後堆裡剩下的 (9, 4)、(10, 3) 拿出來都會丟掉。 費用加起來 $2 + 4 + 6 + 3 + 8 =
23$，選到的邊和 Kruskal 相同，只是順序不一樣：Kruskal 照權重由小到大選，Prim 照
點上樹的順序選，第 5、6 步先花 $6$ 再花 $3$。 費用 $3$ 的 3--5 一直到 3 上樹之後
才「往外」，在那之前它兩端都不在樹上，Prim 看不到它。

\begin{code}{cpp}{Prim 演算法的主迴圈（用 priority\_queue）}
// adj[u]：(鄰居, 權重) 的鄰接串列，點 1 到 n
vector<bool> inTree(n + 1, false);
priority_queue<pair<long long, int>, vector<pair<long long, int>>,
               greater<pair<long long, int>>> pq;   // (費用, 點)，小的在頂端
pq.push({0, 1});                        // 從點 1 開始，接進來不用錢
long long total = 0;
int count = 0;
while (!pq.empty()) {
    auto [w, u] = pq.top();
    pq.pop();
    if (inTree[u]) {                    // 已經在樹裡：過期的紀錄
        continue;
    }
    inTree[u] = true;
    total += w;
    count++;
    for (auto [v, wv] : adj[u]) {
        if (!inTree[v]) {
            pq.push({wv, v});
        }
    }
}
// count < n 表示圖不連通
\end{code}

\begin{itemize}
\item 同一個點可能被放進堆好幾次（表裡的點 2 就有兩筆），只有第一次拿出來的那筆
最便宜，之後的都是過期的紀錄，用 \texttt{inTree} 檢查後丟掉。 這和第~8~章
Dijkstra 丟掉過期紀錄的「先不刪」寫法相同。
\item \texttt{inTree[u] = true} 一定要寫在\textbf{拿出來}的時候，不能寫在放進去的
時候。 放進去時只表示「有一條路可以接到 $u$」，不代表那是最便宜的一條（本章練習
第~4~題）。
\item Prim 只會長出點 1 所在的那一棵樹。 圖不連通時，\texttt{count} 會小於 $n$。
要的是生成森林的話，對每一個還不在樹上的點再開始一次。
\end{itemize}

\textbf{複雜度}：每條邊最多讓兩筆紀錄進堆，堆的操作共 $O(m)$ 次，每次
$O(\log m)$，整體 $O(m \log m)$，和 Kruskal 一樣。

\subsection{稠密圖：不用堆的 Prim}
\label{sec:prim-dense}

有一類題目，任兩個點之間都可以連線，費用由兩點的位置算出來，例如「$n$ 根旗桿，
任兩根都能拉線，費用是兩根的曼哈頓距離 $|x_i - x_j| + |y_i - y_j|$」。 這種圖有
$n(n - 1)/2$ 條邊，邊數是點數平方的量級，叫\textbf{稠密圖}（dense graph）；邊數和點
數差不多的叫\textbf{稀疏圖}（sparse graph）。

$n = 5000$ 時有 $12497500$ 條邊。 Kruskal 要先把這一千兩百多萬條邊全部存下來再排
序；用堆的 Prim 也會把同樣多的紀錄塞進堆裡。 Prim 有一個更簡單的寫法：替每個還
不在樹上的點記一個數 \texttt{best[v]}，表示「把 $v$ 接到目前這棵樹上最少要多少
錢」。 每一步：

\begin{itemize}
\item 掃過所有不在樹上的點，挑 \texttt{best} 最小的 $u$，讓它上樹。
\item 掃過所有不在樹上的點 $v$，算出 $u$ 到 $v$ 的費用 $d$，如果 $d$ 比
\texttt{best[v]} 小，就更新。 $u$ 剛上樹，所以「接到樹上」多了一個選擇。
\end{itemize}

每一步是兩次 $O(n)$ 的掃描，做 $n$ 步，總共 $O(n^2)$。 邊的費用需要時才算，完全
不用存邊。

\begin{code}{cpp}{稠密圖的 Prim，$O(n^2)$：點 $0$ 到 $n - 1$，座標存在 x、y}
vector<long long> best(n, INF);         // 接到目前這棵樹的最低費用
vector<bool> inTree(n, false);
best[0] = 0;
long long total = 0;
for (int round = 0; round < n; round++) {
    int u = -1;
    for (int i = 0; i < n; i++) {       // 樹外 best 最小的點
        if (!inTree[i] && (u == -1 || best[i] < best[u])) {
            u = i;
        }
    }
    inTree[u] = true;
    total += best[u];
    for (int v = 0; v < n; v++) {       // u 上樹之後，更新其他點
        if (!inTree[v]) {
            long long d = llabs(x[u] - x[v]) + llabs(y[u] - y[v]);
            best[v] = min(best[v], d);
        }
    }
}
\end{code}

例如五根旗桿在 $(0, 0)$、$(2, 2)$、$(3, 10)$、$(5, 2)$、$(7, 0)$，最小生成樹是
$(2,2)$--$(5,2)$（$3$）、$(0,0)$--$(2,2)$（$4$）、$(5,2)$--$(7,0)$（$4$）、
$(2,2)$--$(3,10)$（$9$），總和 $20$。

在本書測試的環境下，$n = 5000$ 個隨機的點，這個寫法（補上讀入和輸出）約 $0.04$ 秒，記憶體約
$3$ MB。 把全部 $12497500$ 條邊存下來跑 Kruskal，得到同樣的答案，但要約
$0.72$ 秒和約 $199$ MB，記憶體已經超過很多題目的限制。 輸入直接給一個
$n \times n$ 的鄰接矩陣（第 1 本第~11~章）時，也是同樣的情況，用這個寫法最省事。

\subsection{選哪一個}

Kruskal 和用堆的 Prim 都是 $O(m \log m)$，前者用邊列表，後者用鄰接串列。 大部分題
目用 Kruskal 就好，它最短，也最容易延伸到後面幾節的變形。 看到「任兩點都可以連」而
且 $n$ 只有幾千，或輸入是鄰接矩陣，改用 $O(n^2)$ 的 Prim。 如果 $n$ 到 $10^5$ 又任
兩點都能連，兩種都來不及，要先像第~\ref{sec:mst-build}~節那樣，證明大部分的邊用不
到。

\section{最大生成樹}
\label{sec:max-st}

反過來，要總權重\textbf{最大}的生成樹，叫\textbf{最大生成樹}（maximum spanning
tree）。 做法只改一個地方：Kruskal 的邊改成由\textbf{大到小}排序，Prim 改成每次
接最貴的往外的邊。 章首的圖的最大生成樹總和是 $37$。

為什麼這樣就對？ 把每條邊的權重變號，$w$ 變成 $-w$。 每一棵生成樹的總權重也跟著
變號，原本最大的變成最小的，所以「原圖的最大生成樹」就是「變號之後的最小生成
樹」。 而對變號後的圖由小到大跑 Kruskal，正好就是對原圖由大到小跑。 這個論證沒
有用到權重是正的，切割性質的證明也沒有，所以\textbf{權重是負數時 Kruskal 照樣對}，
最大生成樹和最小生成樹都一樣。 要小心的只是題目的說法（第~\ref{sec:mst-tree}~
節）：限定剛好 $n - 1$ 條時，負的邊該選還是要選。

\subsection{例題：學伴配對的最高默契總分}
\label{sec:ex-buddy}

\textbf{題意}（\akiref{best-total-buddy-compatibility}{學伴配對的最高默契總分}）

\begin{itemize}
\item 有 $n$ 位新生，學生會列出 $m$ 組候選配對，第 $i$ 組是新生 $u_i$ 和 $v_i$，
默契分數 $w_i$，可能是負的。 同一對新生可能出現在好幾組裡，分數不一定相同，每
一組都可以單獨選或不選。
\item 要從中挑出\textbf{剛好 $n - 1$ 組}，讓所有新生沿著挑出來的配對串在一起。
\item 輸出挑出來的分數總和最大是多少，負的分數照算。
\item $n \le 10^5$，$m \le 2 \times 10^5$，$|w_i| \le 10^9$。 保證全部選上時所有
新生會串在一起。
\end{itemize}

例如 $4$ 位新生，配對是 $(1, 2)$ 分數 $3$、$(1, 3)$ 分數 $1$、$(2, 3)$ 分數 $-2$、
$(3, 4)$ 分數 $-7$。 新生 4 只出現在 $(3, 4)$ 裡，不選它就接不上，所以一定要選；
剩下兩組選 $3$ 和 $1$，總和 $3 + 1 - 7 = -3$。

\textbf{想法}

暴力解是枚舉 $m$ 組裡選哪 $n - 1$ 組，$\binom{m}{n - 1}$ 種，只在 $m$ 很小時可行。

「剛好 $n - 1$ 組，全部串在一起」就是生成樹的定義，要總分最大，就是最大生成樹。
照上面說的，把配對依分數由大到小排序跑 Kruskal：會讓兩群新生串起來的就選，兩人
已經串在一起的就跳過。 負分的配對不用特別處理：輪到它時如果還有兩群沒串起來，
就表示比它好的配對都接不起這兩群，只能選它。

同一對新生的重複配對也不用特別處理。 分數高的那一組先被看到，如果選了，同一對
後面的配對兩端已經相通，自然會跳過。

同樣的模型也出現在別的地方：幾個倉庫之間要排固定的運輸路線，每條路線一年的淨
收益有正有負，主管規定剛好開 $n - 1$ 條、而且每個倉庫都要接上，要讓總收益最大。

\textbf{程式}

\begin{example}[label=ex:mst-buddy]{cpp}{(buddy.cpp) 學伴配對的最高默契總分}
#include <algorithm>
#include <iostream>
#include <utility>
#include <vector>
using namespace std;

struct Edge {
    int u, v;
    long long w;
};

vector<int> parent, sz;

int find(int x) {
    return parent[x] == x ? x : parent[x] = find(parent[x]);
}

bool unite(int a, int b) {
    a = find(a);
    b = find(b);
    if (a == b) {
        return false;
    }
    if (sz[a] > sz[b]) {
        swap(a, b);
    }
    parent[a] = b;
    sz[b] += sz[a];
    return true;
}

int main() {
    ios::sync_with_stdio(false);
    cin.tie(nullptr);
    int n, m;
    cin >> n >> m;
    vector<Edge> edges(m);
    for (Edge &e : edges) {
        cin >> e.u >> e.v >> e.w;
    }
    sort(edges.begin(), edges.end(),
         [](const Edge &a, const Edge &b) { return a.w > b.w; });  // 大的先
    parent.resize(n + 1);
    sz.assign(n + 1, 1);
    for (int i = 1; i <= n; i++) {
        parent[i] = i;
    }
    long long total = 0;
    for (const Edge &e : edges) {
        if (unite(e.u, e.v)) {
            total += e.w;               // 負分也照加
        }
    }
    cout << total << '\n';
    return 0;
}
\end{example}

\iosample{
5 7 \\
1 2 4 \\
2 3 6 \\
1 3 5 \\
3 4 -2 \\
4 5 3 \\
2 5 1 \\
1 2 2 \\
}{
15 \\
}

\begin{itemize}
\item 和範例~\ref{ex:mst-kruskal} 只差在比較函式的 \texttt{>}，以及不用印出選到的
邊。 題目保證全部選上時會串在一起，所以不用檢查 \texttt{used}。
\item 範例裡依序選了 $6$、$5$、$3$、$1$，跳過第一組 $(1, 2)$ 的 $4$（1、2 已經經
由 3 串在一起）和最後一組 $(1, 2)$ 的 $2$；$-2$ 的 $(3, 4)$ 輪到時，全部的人已經
串在一起，也跳過。 總和 $6 + 5 + 3 + 1 = 15$。
\end{itemize}

\textbf{複雜度}：$O(m \log m)$。 最大的測試資料在本書測試的環境下約 $0.02$ 秒。

\subsection{常見錯誤}

\begin{itemize}
\item \textbf{照抄最小生成樹的排序}：比較函式寫成 \texttt{a.w < b.w}，求出來的是
最小生成樹。 範例 1 印出 \texttt{4}，上面 $4$ 位新生的例子印出 \texttt{-8}。
\item \textbf{只加正分}：寫成 \texttt{if (e.w > 0) total += e.w;}。 範例 1 剛好
沒選到負分，仍然印出正確的 \texttt{15}；$4$ 位新生的例子印出 \texttt{4}，正確答案
是 \texttt{-3}。 「剛好 $n - 1$ 組」時，負分的配對是被迫要選的。
\item \textbf{總和用 \texttt{int}}：$n = 10^5$ 位新生排成一串、每組都是 $10^9$ 分，
正確答案是 \texttt{99999000000000}。 \texttt{total} 改成 \texttt{int}，在本書測試的
環境下印出 \texttt{-723552768}，加了 \texttt{-Wall -Wextra} 也沒有警告。
\end{itemize}

\section{自己建出圖來}
\label{sec:mst-build}

前面的圖都直接寫在輸入裡。 很多題目要自己決定點是什麼、邊是什麼。 第
~\ref{sec:prim-dense}~節的旗桿就是一例：點是旗桿，任兩根之間都有一條邊。 有時候
邊的數量多到存不下，這時要用環性質證明大部分的邊一定用不到，只留下少數可能有用
的邊。

\subsection{例題：只能拉直線的路燈電纜}
\label{sec:ex-lamps}

\textbf{題意}（\akiref{straight-line-street-lamp-cables}{只能拉直線的路燈電纜}）

\begin{itemize}
\item 校園裡有 $n$ 根路燈，第 $i$ 根在格子點 $(x_i, y_i)$，位置都不同。
\item 兩根路燈在同一條水平線（$y$ 相同）或同一條垂直線（$x$ 相同）上時，可以在
它們之間拉一條直線電纜，費用是兩根的距離。 電纜只在兩端接上路燈，中途經過的路
燈不算接上，兩條電纜交叉也不會相連。
\item 求讓所有路燈都相通的最少總費用；做不到輸出 $-1$。
\item $n \le 2 \times 10^5$，座標的絕對值到 $10^9$。
\end{itemize}

\textbf{想法}

點是路燈，同一行或同一列的任兩根之間有一條邊，答案是這張圖的最小生成樹，圖不連
通時輸出 $-1$。

直接照這個定義建邊：同一列有 $k$ 根路燈，就有 $k(k - 1)/2$ 條邊。 全部 $2 \times
10^5$ 根都在同一列時，大約是 $2 \times 10^{10}$ 條，存不下也排不完。

看同一列上三根依序排列的路燈 $a$、$b$、$c$：

\begin{center}
\begin{tikzpicture}[x=0.9cm, y=1cm]
\draw[black!30] (-0.5,0) -- (7.5,0);
\node[mstv] (a) at (0,0) {$a$};
\node[mstv] (b) at (3,0) {$b$};
\node[mstv] (c) at (7,0) {$c$};
\draw[mstehl] (a) -- node[mstw, below=2pt] {$3$} (b);
\draw[mstehl] (b) -- node[mstw, below=2pt] {$4$} (c);
\draw[mstebad] (a) to[bend left=35] node[mstw, above] {$7 = 3 + 4$} (c);
\end{tikzpicture}
\diagcaption{同一列的三根路燈。 $a$--$c$ 的費用等於 $a$--$b$ 加 $b$--$c$，是三角形
上最貴的邊。}
\label{fig:mst-lamps-line}
\end{center}

$a$--$c$ 的費用等於 $a$--$b$ 加上 $b$--$c$。 三根路燈位置不同，兩段都是正的，所以
$a$--$c$ 比另外兩條都\textbf{嚴格}貴。 $a$、$b$、$c$ 三條邊圍成一個環，照環性質，
$a$--$c$ 不會出現在任何最小生成樹裡，可以直接刪掉。 同樣的道理，同一列上不相鄰
的兩根之間的邊，中間至少夾著一根路燈，都可以刪掉。

刪掉的邊都不在任何一棵最小生成樹裡，所以刪完之後，最小生成樹和原本一樣，連不
連通也不會變。 每一列只要留下「相鄰兩根」之間的邊，每一行也一樣。 把路燈依 $(x, y)$ 排序，
排好之後 $x$ 相同的會排在一起，而且依 $y$ 由小到大，相鄰兩根就是同一行上相鄰的
路燈；再依 $(y, x)$ 排一次，得到同一列上相鄰的路燈。 兩次加起來最多
$2(n - 1)$ 條邊，再跑 Kruskal。

題目說電纜中途經過的路燈不算接上，其實不影響答案：用得到的邊都是相鄰兩根之間的
邊，中途本來就沒有別的路燈。

下圖是範例 1 的七根路燈。 留下來的邊共七條，最小生成樹用了其中六條，總和
$18$。

\begin{center}
\begin{tikzpicture}[x=0.95cm, y=0.95cm]
\foreach \gx in {1,...,6} {\draw[black!12] (\gx,0.5) -- (\gx,5.5); \node[font=\scriptsize, text=black!50] at (\gx,0.2) {\gx};}
\foreach \gy in {1,...,5} {\draw[black!12] (0.5,\gy) -- (6.5,\gy); \node[font=\scriptsize, text=black!50] at (0.2,\gy) {\gy};}
\node[mstv, minimum size=0.55cm] (L1) at (4,5) {1};
\node[mstv, minimum size=0.55cm] (L2) at (1,1) {2};
\node[mstv, minimum size=0.55cm] (L3) at (6,3) {3};
\node[mstv, minimum size=0.55cm] (L4) at (4,1) {4};
\node[mstv, minimum size=0.55cm] (L5) at (2,3) {5};
\node[mstv, minimum size=0.55cm] (L6) at (1,5) {6};
\node[mstv, minimum size=0.55cm] (L7) at (6,1) {7};
\draw[mstehl] (L2) -- node[mstw, below] {3} (L4);
\draw[mstehl] (L4) -- node[mstw, below] {2} (L7);
\draw[mstehl] (L7) -- node[mstw, right] {2} (L3);
\draw[mstehl] (L6) -- node[mstw, above] {3} (L1);
\draw[mstehl] (L5) -- node[mstw, above] {4} (L3);
\draw[mstehl] (L2) -- node[mstw, left] {4} (L6);
\draw[mste, dashed] (L4) -- node[mstw, right] {4} (L1);
\end{tikzpicture}
\diagcaption{範例 1：只留相鄰路燈之間的七條邊，粗線是一棵最小生成樹。 虛線的
4--1 和粗線的 2--6 一樣是 $4$，兩條選哪一條都可以。}
\label{fig:mst-lamps}
\end{center}

\textbf{程式}

\begin{example}[label=ex:mst-lamps]{cpp}{(lamps.cpp) 只能拉直線的路燈電纜}
#include <algorithm>
#include <iostream>
#include <utility>
#include <vector>
using namespace std;

struct Edge {
    int u, v;
    long long w;
};

vector<int> parent, sz;

int find(int x) {
    return parent[x] == x ? x : parent[x] = find(parent[x]);
}

bool unite(int a, int b) {
    a = find(a);
    b = find(b);
    if (a == b) {
        return false;
    }
    if (sz[a] > sz[b]) {
        swap(a, b);
    }
    parent[a] = b;
    sz[b] += sz[a];
    return true;
}

int main() {
    ios::sync_with_stdio(false);
    cin.tie(nullptr);
    int n;
    cin >> n;
    vector<long long> x(n), y(n);
    for (int i = 0; i < n; i++) {
        cin >> x[i] >> y[i];
    }
    vector<Edge> edges;
    vector<int> id(n);
    for (int i = 0; i < n; i++) {
        id[i] = i;
    }
    // 同一行（x 相同）：依 y 排好，只接相鄰的兩根
    sort(id.begin(), id.end(), [&](int a, int b) {
        return x[a] != x[b] ? x[a] < x[b] : y[a] < y[b];
    });
    for (int k = 0; k + 1 < n; k++) {
        int a = id[k], b = id[k + 1];
        if (x[a] == x[b]) {
            edges.push_back({a, b, y[b] - y[a]});
        }
    }
    // 同一列（y 相同）：依 x 排好，只接相鄰的兩根
    sort(id.begin(), id.end(), [&](int a, int b) {
        return y[a] != y[b] ? y[a] < y[b] : x[a] < x[b];
    });
    for (int k = 0; k + 1 < n; k++) {
        int a = id[k], b = id[k + 1];
        if (y[a] == y[b]) {
            edges.push_back({a, b, x[b] - x[a]});
        }
    }
    sort(edges.begin(), edges.end(),
         [](const Edge &a, const Edge &b) { return a.w < b.w; });
    parent.resize(n);
    sz.assign(n, 1);
    for (int i = 0; i < n; i++) {
        parent[i] = i;
    }
    long long total = 0;
    int used = 0;
    for (const Edge &e : edges) {
        if (unite(e.u, e.v)) {
            total += e.w;
            used++;
        }
    }
    cout << (used == n - 1 ? total : -1) << '\n';
    return 0;
}
\end{example}

\iosample{
7 \\
4 5 \\
1 1 \\
6 3 \\
4 1 \\
2 3 \\
1 5 \\
6 1 \\
}{
18 \\
}

\begin{itemize}
\item \texttt{id} 存路燈的編號，排序時用 \texttt{[\&]} 取用外面的 \texttt{x}、
\texttt{y}（銜接書第~9~章），不用把座標和編號綁成 \texttt{struct}。
\item 排好之後 $y$ 由小到大，\texttt{y[b] - y[a]} 一定是正的，就是距離，不用取絕
對值。
\item $n = 1$ 時沒有邊，\texttt{used} 是 $0$，剛好等於 $n - 1$，印出 \texttt{0}，
符合題目的要求。
\end{itemize}

\textbf{複雜度}：三次排序，都是 $O(n \log n)$。 最大的測試資料在本書測試的環境下
約 $0.05$ 秒。

\subsection{常見錯誤}

\begin{itemize}
\item \textbf{只建同一列的邊}：漏掉同一行的那一段。 範例 1 只用同一列的邊，會分成
$\{2, 4, 7\}$、$\{3, 5\}$、$\{1, 6\}$ 三群，程式印出 \texttt{-1}。
\item \textbf{同一列所有的點對都建邊}：答案是對的，但全部在同一列時邊數是
$n(n - 1)/2$，記憶體和時間都不夠。
\item \textbf{忘了檢查連通}：直接印出 \texttt{total}。 範例 2 的路燈分成兩群，正確
答案是 \texttt{-1}，這樣寫會印出森林的總和 \texttt{19}。
\item \textbf{距離或總和用 \texttt{int}}：座標到 $\pm 10^9$，一條電纜的費用就可能到
$2 \times 10^9$，已經超過 \texttt{int}。
\end{itemize}

\subsection{什麼時候用這一招}

能這樣刪邊，是因為「遠的費用剛好等於近的加起來」。 費用和距離有關、連線方向又受限
時，先想想有沒有一條邊一定是某個環上最重的。 如果每一對點各自給一個價格，就沒有
這種捷徑。 另一種建圖的情況是點不是題目裡最顯眼的東西，例如本章實作練習的「倉庫
軌道的轉轍器」。

\section{瓶頸路徑}
\label{sec:bottleneck}

貨車從 $u$ 開到 $v$，路上每座橋有一個限重，貨車的總重不能超過任何一座橋的限重。
這時路線的好壞，只看路上\textbf{最小}的那個限重，叫這條路的\textbf{瓶頸}
（bottleneck）。 要載最多的貨，就要在所有 $u$ 到 $v$ 的路線裡，找瓶頸最大的那一
條。 反過來也常見：每段路有一個坡度，要讓路上最陡的那一段盡量平緩。

這兩種問題和路線的總長無關，第~8~章的最短路徑幫不上忙。 它們的答案都在生成樹
上：

\begin{center}
\fbox{\parbox{0.88\linewidth}{在\textbf{最小}生成樹上，$u$ 到 $v$ 那條唯一的路，路上
最大的邊，是所有 $u$ 到 $v$ 的路線裡最小的。 在\textbf{最大}生成樹上，$u$ 到 $v$
那條路上最小的邊，是所有路線裡最大的。}}
\end{center}

用章首的圖看一個例子。 從 6 到 5 有一條直接的線，權重 $7$。 最小生成樹上的路是
6--1--2--3--5，經過 $2, 4, 6, 3$，最大是 $6$，比直接走的 $7$ 還小。 繞遠路，最大
的一段反而比較小，這正是瓶頸問題和最短路不同的地方。

為什麼？ 用最小生成樹的版本說明，看 Kruskal 是在哪一刻把 $u$ 和 $v$ 接通的。 設
那一刻選的邊是 $e$，權重 $c$。

\begin{itemize}
\item \textbf{樹上的路做得到 $c$。} 那一刻並查集裡 $u$ 和 $v$ 剛好接通，連接它們
的是到目前為止選過的邊，每一條都不超過 $c$，其中有一條是 $e$。 所以樹上的路最大
的邊剛好是 $c$。
\item \textbf{沒有路線能比 $c$ 更好。} 假設有一條 $u$ 到 $v$ 的路線，每一段都
\textbf{嚴格}小於 $c$。 這些邊都排在 $e$ 前面，處理完它們之後，路線上相鄰的點兩
兩同群，所以 $u$ 和 $v$ 早就同群了，Kruskal 不會在 $e$ 這一刻才接通它們，矛盾。
\end{itemize}

最大生成樹的版本把「小」和「大」對調，理由相同。

\subsection{離線處理：只用不超過 $x$ 的邊}
\label{sec:mst-offline}

上面的論證也說明了另一件事：Kruskal 處理完所有權重 $\le x$ 的邊時，並查集的每一
群，正好是「只走權重 $\le x$ 的邊能互相到達」的點。 這給了一種處理大量詢問的方
法。

問題：有 $q$ 筆詢問，第 $i$ 筆問「只走權重 $\le x_i$ 的邊，$a_i$ 和 $b_i$ 能不能
互相到達」。 每筆詢問各做一次 BFS 是 $O(q(n + m))$，$q$、$n$、$m$ 都到
$2 \times 10^5$ 時來不及。

做法是先把全部詢問讀進來，依 $x$ 由小到大排序，再跟著詢問把邊一條一條加進並查
集。 這種「先讀完所有詢問，照自己方便的順序回答，最後再照原本的順序輸出」的做法，
叫\textbf{離線}（offline）處理。 第~7~章練習第~2~題把刪邊倒過來變成加邊，也是離
線處理。

\begin{code}{cpp}{只用權重不超過 x 的邊，a 和 b 連不連通}
sort(edges.begin(), edges.end(),
     [](const Edge &e, const Edge &f) { return e.w < f.w; });
vector<int> order(q);
for (int i = 0; i < q; i++) {
    order[i] = i;
}
sort(order.begin(), order.end(), [&](int i, int j) { return x[i] < x[j]; });
int k = 0;                              // 已經加進並查集的邊數
for (int i : order) {                   // 詢問由 x 小到大處理
    while (k < m && edges[k].w <= x[i]) {
        unite(edges[k].u, edges[k].v);
        k++;
    }
    ans[i] = (find(a[i]) == find(b[i]));
}
\end{code}

\texttt{k} 只往前走，所有詢問加起來最多加 $m$ 條邊，加上兩次排序，整體是
$O(m \log m + q \log q)$。 答案存在 \texttt{ans[i]}，最後依 $i = 0, 1, \ldots$
輸出，順序就和輸入相同。 在章首的圖上問「6 和 5，只用 $\le 5$ 的邊」答案是不能，
「只用 $\le 6$ 的邊」就可以了，和上面算出的瓶頸 $6$ 一致。

同一個寫法可以順手回答「$a$ 所在的那一群有幾個點」（第~7~章的 \texttt{sz}）、
「群裡的總和」這類問題，只要那個量能存在根上、合併時加起來。 下面的例題不是這種
量：推車越窄走得越遠，裝得下的器材卻越少，答案沒辦法隨著合併累加，要換一個想
法。

\subsection{例題：推車寬度與器材配送}
\label{sec:ex-carts}

\textbf{題意}（\akiref{cart-width-and-equipment-deliveries}{推車寬度與器材配送}）

\begin{itemize}
\item 校舍有 $n$ 個地點，地點 1 是器材室。 有 $m$ 條雙向走廊，第 $i$ 條連接 $u_i$
和 $v_i$，寬度 $w_i$。
\item 地點 $2$ 到 $n$ 各要送一件器材，送到地點 $v$ 的器材寬度是 $r_v$。
\item 寬度 $x$ 的推車只能走寬度 $\ge x$ 的走廊，只裝得下寬度 $\le x$ 的器材。 它能
「負責」地點 $v$，是指裝得下送到 $v$ 的器材，而且從器材室出發走得到 $v$。
\item 有 $q$ 台推車，對每一台回答它能負責幾個地點。
\item $n$、$m$、$q$ 都到 $2 \times 10^5$，寬度都到 $10^9$。 圖不一定連通。
\end{itemize}

\textbf{想法}

最直接的做法：每台推車各做一次 BFS，只走夠寬的走廊，再數走得到、又裝得下的地
點。 一次 $O(n + m)$，$q$ 台是 $O(q(n + m))$，大約 $8 \times 10^{10}$，來不及。

換個角度，對每個地點各自想。 寬度 $x$ 的推車能負責 $v$，要同時滿足：

\begin{itemize}
\item 裝得下：$r_v \le x$。
\item 走得到：有一條從器材室到 $v$ 的路線，每條走廊都 $\ge x$，也就是這條路線的瓶
頸 $\ge x$。 把所有路線的瓶頸取最大，記成 $\mathrm{cap}_v$，條件就是
$x \le \mathrm{cap}_v$。
\end{itemize}

所以地點 $v$ 對應一個區間 $[r_v, \mathrm{cap}_v]$，寬度落在這個區間裡的推車都能負
責它。 每台推車的答案就是「有幾個區間包含 $x$」。

\textbf{第一步：算出每個 $\mathrm{cap}_v$。} 這正是瓶頸路徑：在最大生成森林上，器
材室到 $v$ 的路上最窄的走廊。 把走廊由寬到窄跑一次 Kruskal 建出森林，再從地點 1
做一次 BFS，沿路取最小值：$\mathrm{cap}_1$ 設成無限大，走到鄰居 $v$ 時
$\mathrm{cap}_v = \min(\mathrm{cap}_u, w)$。 BFS 走不到的地點和器材室不相通，任何
推車都負責不了。

\textbf{第二步：數有幾個區間包含 $x$。} 先丟掉 $r_v > \mathrm{cap}_v$ 的地點，它們
的區間是空的。 剩下的區間都滿足 $r_v \le \mathrm{cap}_v$，包含 $x$ 等於
「$r_v \le x$」而且「不是 $\mathrm{cap}_v < x$」。 如果 $\mathrm{cap}_v < x$，因為
$r_v \le \mathrm{cap}_v$，一定也有 $r_v \le x$，所以

\[
\text{答案} = (\,r_v \le x \text{ 的個數}\,) - (\,\mathrm{cap}_v < x \text{ 的個數}\,).
\]

把所有 $r_v$ 排序、所有 $\mathrm{cap}_v$ 排序，每台推車用 \texttt{upper\_bound} 和
\texttt{lower\_bound}（銜接書第~9~章）各查一次。

範例 1 的圖畫在下面。 走廊由寬到窄處理時，6--7（$9$）、3--4（$6$）、1--2（$5$）、
2--3（$4$）、4--5（$3$）都被選進森林，1--3（$2$）和 2--5（$1$）被跳過。

\begin{center}
\begin{tikzpicture}[x=1.5cm, y=1.25cm]
\node[mstvhl] (1) at (0,1) {1};
\node[mstv] (2) at (1.6,2) {2};
\node[mstv] (3) at (1.6,0) {3};
\node[mstv] (4) at (3.2,0) {4};
\node[mstv] (5) at (3.2,2) {5};
\node[mstv] (6) at (4.8,2) {6};
\node[mstv] (7) at (4.8,0) {7};
\draw[mstehl] (1) -- node[mstw, above left] {5} (2);
\draw[mste] (1) -- node[mstw, below left] {2} (3);
\draw[mstehl] (2) -- node[mstw, left] {4} (3);
\draw[mstehl] (3) -- node[mstw, below] {6} (4);
\draw[mste] (2) -- node[mstw, above] {1} (5);
\draw[mstehl] (4) -- node[mstw, right] {3} (5);
\draw[mstehl] (6) -- node[mstw, right] {9} (7);
\node[mstnote, above=1pt] at (2.north) {cap 5};
\node[mstnote, below=1pt] at (3.south) {cap 4};
\node[mstnote, below=1pt] at (4.south) {cap 4};
\node[mstnote, above=1pt] at (5.north) {cap 3};
\node[mstnote, above=1pt] at (6.north) {到不了};
\node[mstnote, below=1pt] at (7.south) {到不了};
\end{tikzpicture}
\diagcaption{範例 1 的走廊（數字是寬度）。 粗線是最大生成森林，藍字是從器材室 1
沿森林走過去的 $\mathrm{cap}$。}
\label{fig:mst-carts}
\end{center}

地點 3 值得多看一眼：直接從器材室過去的走廊只有 $2$ 寬，繞 1--2--3 最窄也有 $4$
寬。 器材寬度 $r$ 依序是 $3, 1, 6, 2, 1, 1$，得到的區間是地點 2 的 $[3, 5]$、地點 3
的 $[1, 4]$、地點 5 的 $[2, 3]$；地點 4 的 $[6, 4]$ 是空的，地點 6、7 到不了，都丟掉。

寬度 $3$ 的推車落在三個區間裡，答案 $3$；寬度 $5$ 的只落在 $[3, 5]$，答案 $1$；
寬度 $6$ 和 $9$ 的一個都沒有。 用公式算寬度 $4$：$r_v \le 4$ 的有 $3, 1, 2$ 三個，
$\mathrm{cap}_v < 4$ 的有 $3$ 一個，答案 $3 - 1 = 2$。

\textbf{程式}

\begin{example}[label=ex:mst-carts]{cpp}{(carts.cpp) 推車寬度與器材配送}
#include <algorithm>
#include <iostream>
#include <queue>
#include <utility>
#include <vector>
using namespace std;

struct Edge {
    int u, v;
    long long w;
};

vector<int> parent, sz;

int find(int x) {
    return parent[x] == x ? x : parent[x] = find(parent[x]);
}

bool unite(int a, int b) {
    a = find(a);
    b = find(b);
    if (a == b) {
        return false;
    }
    if (sz[a] > sz[b]) {
        swap(a, b);
    }
    parent[a] = b;
    sz[b] += sz[a];
    return true;
}

int main() {
    ios::sync_with_stdio(false);
    cin.tie(nullptr);
    int n, m, q;
    cin >> n >> m >> q;
    vector<long long> r(n + 1, 0);
    for (int v = 2; v <= n; v++) {
        cin >> r[v];
    }
    vector<Edge> edges(m);
    for (Edge &e : edges) {
        cin >> e.u >> e.v >> e.w;
    }
    // 第一步：最大生成森林（寬的走廊先接）
    sort(edges.begin(), edges.end(),
         [](const Edge &a, const Edge &b) { return a.w > b.w; });
    parent.resize(n + 1);
    sz.assign(n + 1, 1);
    for (int i = 1; i <= n; i++) {
        parent[i] = i;
    }
    vector<vector<pair<int, long long>>> tree(n + 1);
    for (const Edge &e : edges) {
        if (unite(e.u, e.v)) {
            tree[e.u].push_back({e.v, e.w});
            tree[e.v].push_back({e.u, e.w});
        }
    }
    // 第二步：從器材室在森林上 BFS，cap[v] = 路上最窄的走廊
    const long long INF = 2e9;
    vector<long long> cap(n + 1, -1);   // -1：到不了
    cap[1] = INF;
    queue<int> bfs;
    bfs.push(1);
    while (!bfs.empty()) {
        int u = bfs.front();
        bfs.pop();
        for (auto [v, w] : tree[u]) {
            if (cap[v] == -1) {
                cap[v] = min(cap[u], w);
                bfs.push(v);
            }
        }
    }
    // 第三步：每個地點是一個區間 [r_v, cap_v]，空的不要
    vector<long long> lo, hi;
    for (int v = 2; v <= n; v++) {
        if (cap[v] != -1 && r[v] <= cap[v]) {
            lo.push_back(r[v]);
            hi.push_back(cap[v]);
        }
    }
    sort(lo.begin(), lo.end());
    sort(hi.begin(), hi.end());
    for (int j = 0; j < q; j++) {
        long long x;
        cin >> x;
        long long startOk = upper_bound(lo.begin(), lo.end(), x) - lo.begin();
        long long endBad = lower_bound(hi.begin(), hi.end(), x) - hi.begin();
        cout << startOk - endBad << '\n';   // r_v <= x 的，扣掉 cap_v < x 的
    }
    return 0;
}
\end{example}

\iosample{
7 7 6 \\
3 1 6 2 1 1 \\
1 2 5 \\
1 3 2 \\
2 3 4 \\
3 4 6 \\
2 5 1 \\
4 5 3 \\
6 7 9 \\
3 1 4 5 6 9 \\
}{
3 \\
1 \\
2 \\
1 \\
0 \\
0 \\
}

\begin{itemize}
\item \texttt{tree} 只存 Kruskal 選中的邊，是一片森林，在它上面 BFS，每個點的上一
步只有一種走法，\texttt{cap} 就是那條唯一的路上的最小值。
\item 用 BFS，不用遞迴的 DFS。 森林可能是一條 $2 \times 10^5$ 個點的長鏈，遞迴太深
（第 1 本 11.5 節）。
\item \texttt{upper\_bound(..., x)} 找第一個 $> x$ 的位置，前面的個數就是
$r_v \le x$ 的個數；\texttt{lower\_bound(..., x)} 找第一個 $\ge x$ 的位置，前面的
個數是 $\mathrm{cap}_v < x$ 的個數。 兩邊都用對，寬度剛好相等的情況才會算進去。
\end{itemize}

\textbf{複雜度}：排序走廊 $O(m \log m)$，BFS $O(n)$，排序區間端點 $O(n \log n)$，每台
推車 $O(\log n)$。 最大的測試資料在本書測試的環境下約 $0.06$ 秒。

\subsection{常見錯誤}

\begin{itemize}
\item \textbf{在原本的圖上 BFS}：以為從器材室 BFS 走到的路就是最寬的路，直接沿 BFS
取最小值。 BFS 走的是走廊數最少的路，地點 3 會從寬度 $2$ 的那條直接過去。 範例 1
印出 \texttt{1 1 1 1 0 0}（每個數一行，這裡寫成一行），正確是
\texttt{3 1 2 1 0 0}。
\item \textbf{沒有丟掉空的區間}：公式的推導用到了 $r_v \le \mathrm{cap}_v$。 地點 4
的區間是空的，留著它，寬度 $5$ 的推車會被多扣一次，範例 1 印出
\texttt{3 1 2 0 0 0}。
\item \textbf{到不了的地點當成無限寬}：\texttt{cap} 的初值設成無限大，地點 6、7 也
被算進去，範例 1 印出 \texttt{5 3 4 3 2 2}。
\end{itemize}

\section{本章常見錯誤}
\label{sec:mst-pitfalls}

前面各節的錯誤，歸納起來是下面幾類。

\begin{center}
\begin{tabular}{p{0.22\linewidth}p{0.52\linewidth}p{0.16\linewidth}}
\toprule
錯誤 & 例子 & 見 \\
\midrule
和最短路徑搞混 & 用 Dijkstra 的最短路徑樹當最小生成樹；Prim 的堆裡放
\texttt{dist[u] + w}；以為樹上的路就是最短路 & 第~\ref{sec:mst-vs-sp}、
\ref{sec:prim}~節 \\
排序方向 & 最大生成樹照抄由小到大；瓶頸問題用錯了最小或最大生成樹 &
第~\ref{sec:ex-buddy}、\ref{sec:bottleneck}~節 \\
沒檢查連通 & 不連通時直接輸出總和；Prim 只長出點 1 那一棵 &
第~\ref{sec:mst-forest}、\ref{sec:ex-lamps}~節 \\
負權與題目說法 & 剛好 $n - 1$ 條時只加正的、和 $0$ 取較大值 &
第~\ref{sec:ex-buddy}~節 \\
邊太多 & 同一列所有點對都建邊；稠密圖存下全部的邊再跑 Kruskal &
第~\ref{sec:prim-dense}、\ref{sec:ex-lamps}~節 \\
\texttt{int} 溢位 & 總和、座標相減用 \texttt{int} & 第~\ref{sec:ex-buddy}、
\ref{sec:ex-lamps}~節 \\
\bottomrule
\end{tabular}
\end{center}

\section{本章練習}
\label{sec:ch9-exercises}

\begin{drills}[label=drills:ch9]
\item 一張圖有 $5$ 個點、$7$ 條邊：1--2 權重 $3$、1--3 權重 $1$、2--3 權重 $3$、
2--4 權重 $2$、3--4 權重 $5$、3--5 權重 $4$、4--5 權重 $4$。
\begin{enumerate}
\item 照 Kruskal 演算法，依序寫出每條邊選或不選，以及最小生成樹的總權重。
\item 權重相同的邊換個順序處理，選到的邊會不會變？ 總權重呢？
\item 從點 1 開始跑 Prim，寫出點上樹的順序。
\end{enumerate}

\item 判斷下面每一句話對不對。 對的說明理由，錯的給一個反例。
\begin{enumerate}
\item 權重最小的邊，一定在每一棵最小生成樹裡。
\item 權重最大的邊，一定不在最小生成樹裡。
\item 所有邊的權重都不相同時，最小生成樹只有一棵。
\item 每條邊的權重都加上 $10$，最小生成樹的邊不會變；從 $s$ 到 $t$ 的最短路徑也不
會變。
\end{enumerate}

\item 章首的圖（圖~\ref{fig:mst-factories}）多了一條新的路線 6--4，造價 $1$。 新的
最小生成樹總和是多少？ 只用原本最小生成樹的五條邊和這條新邊，能不能算出答案？
說明理由。


\item 有人把第~\ref{sec:prim}~節 Prim 的主迴圈改成放進堆的時候就標記，像 BFS 一樣：
\begin{code}{cpp}{練習 6}
inTree[1] = true;
pq.push({0, 1});
while (!pq.empty()) {
    auto [w, u] = pq.top();
    pq.pop();
    total += w;
    count++;
    for (auto [v, wv] : adj[u]) {
        if (!inTree[v]) {
            inTree[v] = true;
            pq.push({wv, v});
        }
    }
}
\end{code}
給一組會讓它算錯的輸入，說明錯在哪裡。
\end{drills}

\section{實作練習}
\label{sec:ch9-practice}

下面是 Aki-Judge 上和本章相關的題目，依小節分組，難度的四階和前面各章相同。

\begin{center}
\begin{tabular}{llc}
\toprule
小節 & 題目 & 難度 \\
\midrule
Kruskal 與建圖 & \akiref{straight-line-street-lamp-cables}{只能拉直線的路燈電纜} & 基本 \\
 & \akiref{track-switches-in-the-warehouse}{倉庫軌道的轉轍器} & 基本 \\
\midrule
最大生成樹 & \akiref{best-total-buddy-compatibility}{學伴配對的最高默契總分} & 暖身 \\
 & \akiref{farm-road-upgrade-standards}{農路升級的標準} & 實戰 \\
\midrule
瓶頸路徑 & \akiref{the-weakest-link-in-the-class-network}{班級聯絡網的最弱一環} & 暖身 \\
 & \akiref{cart-width-and-equipment-deliveries}{推車寬度與器材配送} & 基本 \\
 & \akiref{power-lines-for-typhoon-season}{颱風季的電線規劃} & 實戰 \\
\midrule
最小生成樹的性質 & \akiref{most-bike-lanes-at-each-intersection}{路口最多接幾條車道} & 實戰 \\
 & \akiref{sites-for-a-new-substation}{新變電所的選址提案} & 挑戰 \\
\bottomrule
\end{tabular}
\end{center}

本章的三題例題也在 Aki-Judge 上，讀完各節之後可以先自己寫一次再對照。 「倉庫軌道的
轉轍器」先想清楚什麼是點，點的數量不是交叉處的個數。 「路口最多接幾條車道」要考慮
所有的最小生成樹，把同一個權重的邊當成一組來看。 「新變電所的選址提案」和本章練習
第~3~題有關，但要很快查出樹上一條路的最大邊，需要第~7~章函數圖用過的倍增想法。

\section{讀完你應該會}
\label{sec:ch9-checklist}

\begin{itemize}
\item 看到「選一些邊讓全部連通、總和最小」或「剛好 $n - 1$ 條」的題目，能認出最小
生成樹，並和「每個點到起點最近」的最短路徑問題分開。
\item 能寫出 Kruskal（排序加並查集），並處理不連通、負權重和最大生成樹。
\item 能用切割性質和交換論證說明 Kruskal 和 Prim 為什麼對，並用環性質判斷哪些邊一
定用不到。
\item 能寫出用堆的 Prim，以及稠密圖上 $O(n^2)$ 的 Prim，並依輸入的形式選一個。
\item 看到「路上最窄（最陡）的一段」這類問題，能用最大（最小）生成樹求出瓶頸，或
把詢問依門檻排序，用並查集離線回答。
\end{itemize}

\textbf{延伸閱讀}（書末會整理完整的參考資料）

\begin{itemize}
\item Antti Laaksonen，\textit{Competitive Programmer's Handbook}，第 15 章：
Kruskal、Prim 和並查集，篇幅短，例子清楚。
\item Thomas H. Cormen 等，\textit{Introduction to Algorithms}，第 4 版，第 21 章：
用「目前選的邊包含在某棵最小生成樹裡」這個迴圈不變量，一次證明 Kruskal 和 Prim；
章末問題有次佳生成樹和瓶頸生成樹。
\item Steven Halim 等，\textit{Competitive Programming 4}，Book 1，4.3 節：整理了最
大生成樹、指定邊、最小生成森林、minimax 路徑、次佳生成樹等變形。
\end{itemize}

\addcontentsline{toc}{section}{練習解答}
\section*{練習解答}

\begin{enumerate}
\item
\begin{enumerate}
\item 排序後是 1--3（$1$）、2--4（$2$）、1--2（$3$）、2--3（$3$）、3--5（$4$）、
4--5（$4$）、3--4（$5$）。 選 1--3、選 2--4、選 1--2（把 $\{1, 3\}$ 和
$\{2, 4\}$ 接起來）、跳過 2--3、選 3--5、跳過 4--5，這時已經選滿四條，3--4 也會跳
過。 總權重 $1 + 2 + 3 + 4 = 10$。
\item 可能會變：2--3 排在 1--2 前面時選 2--3、跳過 1--2；4--5 排在 3--5 前面時選
4--5。 總權重一定還是 $10$，兩種排法都是合法的 Kruskal。
\item 1 先上樹；往外最便宜的是 1--3，3 上樹；接著 1--2 和 2--3 都是 $3$，2 上樹；
接著 2--4（$2$）最便宜，4 上樹；最後 3--5 和 4--5 都是 $4$，5 上樹。 順序是
1、3、2、4、5。
\end{enumerate}

\item
\begin{enumerate}
\item 錯。 三個點兩兩相連、三條邊權重都是 $1$ 時，任兩條都是最小生成樹，每一條邊
都有某一棵最小生成樹沒用到它。 如果權重最小的邊只有一條，這句話就對了：某棵最
小生成樹沒用到它的話，把它加進去會形成環，拿掉環上另一條（比它重的）邊，總權重
嚴格變小，矛盾。
\item 錯。 如果某個點只有一條邊連出去，那條邊一定要選，不管它多重。 例如兩個點
之間只有一條權重 $100$ 的邊。
\item 對。 設 $T_1 \ne T_2$ 都是最小生成樹，$e$ 是只出現在其中一棵的邊裡最輕的，不妨
設在 $T_1$。 $e$ 加進 $T_2$ 形成的環上有一條不在 $T_1$ 的邊 $f$，$f$ 比 $e$ 重，
$T_2$ 用 $e$ 換掉 $f$ 會嚴格變小，矛盾。
\item 前半句對：每一棵生成樹都剛好 $n - 1$ 條邊，總權重都多了 $10(n - 1)$，誰最
小不會變。 後半句錯：路徑的邊數不一定相同。 例如 $s$--$t$ 權重 $25$，另一條路
$s$--$a$--$t$ 兩段各 $5$。 原本走 $s$--$a$--$t$ 是 $10$；各加 $10$ 之後，直接走
是 $35$，繞路是 $40$，最短路變成直接走。
\end{enumerate}

\item 總和 $16$。 舊樹上 6 到 4 的路是 6--1--2--3--4，加上新邊形成的環上最重的是
3--4（$8$），換掉它得到 $23 - 8 + 1 = 16$。 只看舊樹的邊加新邊就夠了：不在舊樹上的
邊 $g$，被 Kruskal 跳過時兩端已經經由更便宜的舊樹邊相通，$g$ 是那個環上最重的（這
張圖權重都不同）；加了新邊這個環還在，照環性質 $g$ 仍然用不到。


\item 輸入 \texttt{3 3}、\texttt{1 2 10}、\texttt{1 3 1}、\texttt{3 2 1}，正確答案是
$2$（1--3 和 3--2）。 從 1 開始，2 和 3 一放進堆就都被標記；拿出 3 時，2 已經標記
過，3--2 這條便宜的邊不會被放進堆；之後拿出 2，用的是 1--2 的 $10$，算出 $11$。
放進堆時只知道「有一條邊可以接到 $v$」，之後可能出現更便宜的邊，所以要等拿出來時
才能確定。 BFS 可以放進去就標記，是因為 BFS 的每條邊都一樣長，先發現的一定最近。
\end{enumerate}

\chapter{樹形 DP}
\label{ch:treedp}

% 本章的圖共用的樣式，和第 1 本第 13 章同名、同定義（那一章不在這本書裡，所以這裡再定義一次）
\tikzset{
    tv/.style={circle, draw=black!60, thick, minimum size=0.72cm, inner sep=0pt, font=\small},
    tvhl/.style={tv, draw=memStroke, fill=memFill},
    tvalt/.style={tv, draw=optStroke, fill=optFill},
    te/.style={thick, draw=black!60},
    tehl/.style={very thick, draw=memStroke},
    tlbl/.style={font=\scriptsize, text=black!60},
}

週末市集在一片廣場上舉辦，廣場的走道連成一棵樹，每個路口可以擺一個攤位，主辦
單位估過每個路口擺攤的營收。 規定只有一條：相鄰的兩個路口不能都擺，不然兩攤會
搶同一批客人，走道也會擠成一團。 擺哪些路口，營收加起來最多？ 往下讀之前，先想
一個挑法，例如「營收高的先擺，擺了就把相鄰的路口劃掉」，或是「從入口算起，一層
擺、一層不擺」，再想想它會不會出錯。

走道網是一棵樹。 《C++ 演算法解題 1：APCS 實作題核心》（以下簡稱第 1 本）第~13~章已經在樹上由下往上算過子樹大小、子樹和、高度，每
個點的答案都由子節點的答案組成；章末預告過，這就是把第 1 本第~9~章的動態規劃搬到樹
上。 這一章接著往下做：每個點像第 1 本第~9~章「選或不選」那樣存好幾個狀
態，再用子節點的狀態合併出自己的。

這一章依序講四件事。 先說明為什麼子樹之間可以分開算，以及合併的寫法；接著把
「選或不選」的兩個狀態用到樹上，解掉上面的擺攤問題；再來是樹上的路徑，以第 1 本第~13
~章的直徑為例，改用動態規劃的算法；最後是\textbf{換根}（rerooting），一次算出
「每個點當根」時的答案，例如每個點到其他所有點的距離和。 章末預告一種狀態是一
整張表的樹形 DP。

這一章不教新的語法，也不重講第 1 本第~13~章的存法和走訪：從根做一次 BFS 得到
\texttt{order} 和 \texttt{parent}，倒著看 \texttt{order} 由下往上算（第 1 本第~13~章
13.3 節），這一章每支程式都直接沿用。

\section{每個點的狀態由子節點合併而來}
\label{sec:tdp-merge}

第 1 本第~13~章的子樹大小是 $\mathit{sz}(v) = 1 + $ 子節點的 $\mathit{sz}$ 加起來。 用第 1 本第
~9~章的四件事來看它：

\begin{itemize}
\item \textbf{狀態}：$\mathit{sz}(v)$ 是「以 $v$ 為根的子樹有幾個點」。 樹形 DP
的狀態幾乎都長這樣：只看 $v$ 的子樹，把題目的問題在這棵子樹上問一次。 常常還要
多記一點 $v$ 自己的資訊，例如 $v$ 選了沒有、塗了什麼顏色。
\item \textbf{轉移}：用 $v$ 自己的值，加上每個子節點的狀態，合併出 $v$ 的狀態。
\item \textbf{初始值}：只有 $v$ 自己、還沒合併任何子節點時的值。 葉子沒有子節點，
它的狀態就是初始值。
\item \textbf{計算順序}：子節點要先算完，也就是倒著看 BFS 順序。
\end{itemize}

答案通常是根的狀態；路徑之類的問題，答案可能落在任何一個點上，要在途中取最大
值，第~\ref{sec:tdp-path}~節會遇到。

\subsection{為什麼子樹可以分開算}

樹形 DP 能成立，靠的是樹的一個性質：$v$ 的各個子樹之間沒有邊。 拿掉 $v$，它
的每個子樹各自成為一塊，兩個子樹裡的點要互相到達，一定得經過 $v$（第 1 本第~13~章：
樹上兩點只有一條路徑）。

\begin{center}
\begin{tikzpicture}[x=1cm, y=1cm]
\node[tvhl] (v) at (3,2.4) {$v$};
\foreach \x/\c in {0.6/1, 3/2, 5.4/3} {
    \node[tv] (c\c) at (\x,1.2) {$c_\c$};
    \draw[te] (v) -- (c\c);
    \draw[black!50] (c\c.south) -- (\x-0.7,-0.4) -- (\x+0.7,-0.4) -- cycle;
}
\node[tlbl] at (1.8,-0.1) {沒有邊};
\node[tlbl] at (4.2,-0.1) {沒有邊};
\end{tikzpicture}
\diagcaption{$v$ 的三個子樹。 子樹和子樹之間沒有邊，它們只透過 $v$ 相連。}
\label{fig:tdp-indep}
\end{center}

所以只要 $v$ 的狀態定了（例如「$v$ 擺攤」），每個子樹怎麼選，只受 $v$ 影響，
不受其他子樹影響。 每個子樹各自挑最好的，再把結果合起來，就是 $v$ 在這個狀態
下最好的結果。 有環的圖沒有這個性質：兩個「子樹」之間可能另有一條邊，一邊的選
擇會牽動另一邊，就不能分開算。

\subsection{一次合併一個子節點}

第 1 本第~13~章由下往上的迴圈，是每個點把自己的值「交給」父節點一次。 樹形 DP 照用
這個寫法，只是交的東西換成幾個狀態，合併的方式由題目決定：

\begin{code}{cpp}{樹形 DP 的骨架}
// 先把每個點的狀態設成初始值（只有自己）
for (int i = n - 1; i >= 1; i--) {      // 倒著看，跳過根 order[0]
    int v = order[i], p = parent[v];
    // 輪到 v 時，v 的子節點都已經交過，v 的狀態是最終的
    // 把 v 的狀態合併進 p 的狀態
}
\end{code}

一次併一個子節點，要求合併的方式可以一個一個累積：加總、取最大、取最小都可以，
先併哪一個子節點都一樣。 第~\ref{sec:tdp-preview}~節的狀態是一整張表，合併就
沒有這麼便宜，要另外想辦法。

\section{選或不選}
\label{sec:tdp-mis}

回到章首的市集。 每個路口的營收 $w_v$ 標在點旁邊，相鄰的兩個路口不能都擺。 在一張圖上挑一些點、任兩個挑到的點都不相鄰，這樣的點集合叫
\textbf{獨立集}（independent set）；章首的問題是挑出營收總和最大的獨立集，叫
\textbf{最大權獨立集}。 下圖是一個有八個路口的市集，入口是 $1$：

\begin{center}
\begin{tikzpicture}[x=1.1cm, y=1.0cm]
\node[tvhl] (1) at (3,3) {1};
\node[tv] (2) at (1.5,2) {2};
\node[tv] (3) at (4.5,2) {3};
\node[tvhl] (4) at (0.8,1) {4};
\node[tvhl] (5) at (2.2,1) {5};
\node[tv] (6) at (4.5,1) {6};
\node[tvhl] (7) at (3.8,0) {7};
\node[tvhl] (8) at (5.2,0) {8};
\foreach \a/\b in {1/2, 1/3, 2/4, 2/5, 3/6, 6/7, 6/8} {\draw[te] (\a) -- (\b);}
\foreach \v/\w in {1/6, 3/2, 6/5, 8/4} {\node[tlbl, right] at (\v.east) {$w = \w$};}
\foreach \v/\w in {2/5, 4/3, 7/2} {\node[tlbl, left] at (\v.west) {$w = \w$};}
\node[tlbl, below] at (5.south) {$w = 4$};
\end{tikzpicture}
\diagcaption{八個路口的走道網，旁邊是擺攤的營收（萬元）。 橘色的 $1, 4, 5, 7, 8$
是最好的擺法，營收 $19$。}
\label{fig:tdp-mis}
\end{center}

\subsection{兩種直覺的挑法都會錯}

最直接的做法是暴力解：每個路口擺或不擺，共 $2^n$ 種，挑出合法的，取營收最大。
$n = 30$ 就超過十億種，只能拿來對拍。

章首的兩種挑法，在這張圖上都不是最好的：

\begin{itemize}
\item \textbf{營收高的先擺}：先擺 $1$（$6$），劃掉 $2$、$3$；再擺 $6$（$5$），
劃掉 $7$、$8$；剩下 $4$、$5$ 都能擺。 總共 $6 + 5 + 3 + 4 = 18$。
\item \textbf{一層擺、一層不擺}：深度是偶數的 $1, 4, 5, 6$ 共 $18$，奇數的
$2, 3, 7, 8$ 共 $13$，取大的 $18$。
\end{itemize}

最好的是 $1, 4, 5, 7, 8$，營收 $19$：$6$ 的營收雖然高，不擺它、改擺它底下的
$7$、$8$ 反而多一萬。 一個路口該不該擺，要看它底下的路口能拿多少，只看眼前的營
收判斷不了。

\subsection{兩個狀態}

先試一個狀態：$\mathit{best}(v)$ 是「$v$ 的子樹裡，營收最多是多少」。 想用子節
點的 $\mathit{best}$ 算 $v$ 時會卡住：$v$ 能不能擺，要看子節點有沒有擺，而
$\mathit{best}(c)$ 把「$c$ 有擺」和「$c$ 沒擺」兩種情況混在一起了。 這和第 1 本第~9
~章招牌燈的第一步一樣，缺的資訊放進狀態：

\begin{itemize}
\item $\mathit{in}(v)$：$v$ \textbf{擺}攤時，$v$ 的子樹裡最大的營收。
\item $\mathit{out}(v)$：$v$ \textbf{不擺}時，$v$ 的子樹裡最大的營收。
\end{itemize}

轉移看 $v$ 自己。 $v$ 擺的話，每個子節點都不能擺；$v$ 不擺的話，每個子節點擺不
擺都可以，各自挑比較好的那一種。 由第~\ref{sec:tdp-merge}~節，各個子樹可以分開
挑，所以
\[
\mathit{in}(v) = w_v + \sum_{c} \mathit{out}(c),\qquad
\mathit{out}(v) = \sum_{c} \max\bigl(\mathit{in}(c), \mathit{out}(c)\bigr),
\]
其中 $c$ 跑遍 $v$ 的所有子節點。 初始值是只有自己：$\mathit{in}(v) = w_v$、
$\mathit{out}(v) = 0$，葉子就停在這裡。 答案是 $\max(\mathit{in}(r),
\mathit{out}(r))$，$r$ 是根。

和第 1 本第~9~章對照：一排燈是每盞燈只有一個「下一盞」，樹是每個點有好幾個子節點，轉
移從「看前一格」變成「把所有子節點加起來」，其他完全一樣。

\subsection{由上往下找出擺在哪裡}

只要營收的話，算到根就結束了。 想知道擺在哪些路口，再順著 BFS 順序從根往下決定一次：
根選 $\mathit{in}$ 和 $\mathit{out}$ 裡大的那一種；之後每個點，如果父節點擺了，
它一定不擺；父節點沒擺，它就選自己 $\mathit{in}$ 和 $\mathit{out}$ 裡大的那一
種。 這和 $\mathit{out}(v)$ 的式子「每個子節點各自挑比較好的」是同一件事，只是倒
過來做。 父節點的決定要先知道，所以這一趟是由上往下。

\begin{example}[label=ex:tdp-mis]{cpp}{(market.cpp) 市集要擺哪些路口}
#include <algorithm>
#include <iostream>
#include <vector>
using namespace std;

int main() {
    ios::sync_with_stdio(false);
    cin.tie(nullptr);
    int n;
    cin >> n;
    vector<long long> w(n + 1);
    for (int v = 1; v <= n; v++) {
        cin >> w[v];
    }
    vector<vector<int>> adj(n + 1);
    for (int i = 0; i < n - 1; i++) {
        int u, v;
        cin >> u >> v;
        adj[u].push_back(v);
        adj[v].push_back(u);
    }

    // 以 1 為根：BFS 順序與父節點（第 1 本第 13 章）
    vector<int> parent(n + 1, 0);
    vector<int> order = {1};
    for (int i = 0; i < (int)order.size(); i++) {
        int u = order[i];
        for (int v : adj[u]) {
            if (v != parent[u]) {
                parent[v] = u;
                order.push_back(v);
            }
        }
    }

    // in[v]：v 擺攤，子樹裡最大營收；out[v]：v 不擺，子樹裡最大營收
    vector<long long> in(n + 1), out(n + 1, 0);
    for (int v = 1; v <= n; v++) {
        in[v] = w[v];                       // 還沒合併任何子節點
    }
    for (int i = n - 1; i >= 1; i--) {      // 倒著看，跳過根
        int v = order[i], p = parent[v];
        in[p] += out[v];                    // p 擺：v 一定不擺
        out[p] += max(in[v], out[v]);       // p 不擺：v 選比較好的
    }
    cout << max(in[1], out[1]) << '\n';

    // 順著 BFS 順序，由上往下決定擺哪些
    vector<bool> go(n + 1, false);
    for (int v : order) {
        if (v == 1) {
            go[v] = in[v] > out[v];
        } else if (go[parent[v]]) {
            go[v] = false;                  // 父節點擺了，自己不能擺
        } else {
            go[v] = in[v] > out[v];
        }
    }
    bool first = true;
    for (int v = 1; v <= n; v++) {
        if (go[v]) {
            cout << (first ? "" : " ") << v;
            first = false;
        }
    }
    cout << '\n';
    return 0;
}
\end{example}

\iosample{
8 \\
6 5 2 3 4 5 2 4 \\
1 2 \\
1 3 \\
2 4 \\
2 5 \\
3 6 \\
6 7 \\
6 8 \\
}{
19 \\
1 4 5 7 8 \\
}

BFS 順序是 $1, 2, 3, 4, 5, 6, 7, 8$。 倒著看的迴圈，每一步之後父節點的兩個值：

\begin{center}
\begin{tabular}{cccccc}
\toprule
\texttt{i} & $v$ & $p$ & 交上去的 & $\mathit{in}(p)$ & $\mathit{out}(p)$ \\
\midrule
7 & 8 & 6 & $\mathit{out} = 0$，$\max = 4$ & 5 & 4 \\
6 & 7 & 6 & $\mathit{out} = 0$，$\max = 2$ & 5 & 6 \\
5 & 6 & 3 & $\mathit{out} = 6$，$\max = 6$ & 8 & 6 \\
4 & 5 & 2 & $\mathit{out} = 0$，$\max = 4$ & 5 & 4 \\
3 & 4 & 2 & $\mathit{out} = 0$，$\max = 3$ & 5 & 7 \\
2 & 3 & 1 & $\mathit{out} = 6$，$\max = 8$ & 12 & 8 \\
1 & 2 & 1 & $\mathit{out} = 7$，$\max = 7$ & 19 & 15 \\
\bottomrule
\end{tabular}
\end{center}

$6$ 的兩個值是 $\mathit{in} = 5$、$\mathit{out} = 6$：不擺比較好，這就是「營收高
的先擺」錯的地方。 由上往下時，$1$ 擺（$19 > 15$），所以 $2$、$3$ 不擺；$2$ 沒
擺，$4$、$5$ 各自比較，都是擺比較好；$3$ 沒擺，$6$ 選 $\mathit{out}$，不擺；$7$、
$8$ 都擺。

\begin{itemize}
\item \texttt{in} 和 \texttt{out} 一開始就放好初始值，每個子節點交一次就加一項，
和第 1 本第~13~章累加 \texttt{sz} 的寫法一樣。
\item \texttt{in[v] > out[v]} 相等時選不擺，營收一樣，只是擺法可能和別的寫
法不同。 題目如果要求印出特定的一種（例如攤位最少），相等時怎麼選要照題目的規定。
\item 營收到 $10^9$、路口到 $2 \times 10^5$ 個時，總和會超過 \texttt{int}，所以用
\texttt{long long}。
\end{itemize}

兩趟都是 $O(n)$。 同樣的模型也出現在別的地方：樹狀的通訊線路上挑幾站裝中繼
器，相鄰兩站同時裝會互相干擾；一條樹狀的配送路線上挑幾個點設取貨櫃，相鄰的兩點
設了會互相分走客人。

\subsection{常見錯誤}

\begin{itemize}
\item \textbf{$v$ 不擺，就以為子節點一定要擺}：把 \texttt{out[p] += max(in[v],
out[v]);} 寫成 \texttt{out[p] += in[v];}。 範例的輸入第一行印出 \texttt{18}。
「不能都擺」只限制了兩個都擺，沒有規定至少要擺一個；$v$ 不擺時，子節點不擺也可
能比較好，例如 $6$。
\item \textbf{由上往下時沒看父節點}：每個點都只比較自己的 \texttt{in} 和
\texttt{out}。 範例的第二行印出 \texttt{1 3 4 5 7 8}，$1$ 和 $3$ 相鄰卻都擺了。
$\mathit{in}(3) = 8$ 比 $\mathit{out}(3) = 6$ 大，是說「只看 $3$ 的子樹的話，$3$
擺比較好」；可是 $1$ 已經擺了，$3$ 根本沒有選擇。
\end{itemize}

\subsection{例題：班服顏色的搭配}
\label{sec:tdp-ex-shirts}

\textbf{題意}（\akiref{class-shirt-colors}{班服顏色的搭配}）

\begin{itemize}
\item $n$ 個班之間有 $n - 1$ 組合作關係，形成一棵樹。 每個班的班服選白色或紅色，
第 $i$ 班選白色的成本是 $a_i$，紅色是 $b_i$。
\item 一組合作關係 $(u, v, c)$ 兩端的顏色不同時，要多付協調費 $c$。
\item 求班服成本加上協調費的最小總和。 $n \le 2 \times 10^5$，$a_i, b_i, c$ 都
到 $10^9$。
\end{itemize}

例如三個班，$1$ 班白 $5$、紅 $4$，$2$、$3$ 班都是白 $1$、紅 $9$，$1$ 班和另外
兩班各有一組協調費 $10$ 的合作關係。 每班各自選便宜的顏色是紅、白、白，班服
$6$ 元，協調費 $20$ 元，共 $26$。 讓 $1$ 班也穿白色，班服多 $1$ 元，協調費全省
下來，共 $7$ 元，這才是最小的。

\textbf{想法}

暴力解是枚舉 $2^n$ 種塗法，$n \le 20$ 時可以用來對拍。 上面的例子已經說明了每
班各自選便宜的會錯：一個班的顏色會影響和鄰居之間要不要付協調費。

這題的形狀和市集一樣，每個點有兩種選擇，只是限制從「相鄰的不能都選」變成「相鄰
的不一樣要付錢」。 狀態照樣分兩個：

\begin{itemize}
\item $\mathit{white}(v)$：$v$ 穿白色時，$v$ 的子樹裡最少要花多少（子樹裡的班服，
加上子樹裡合作關係的協調費）。
\item $\mathit{red}(v)$：$v$ 穿紅色時，同上。
\end{itemize}

$v$ 和子節點 $c$ 之間那組合作關係的協調費記作 $f_c$。 $v$ 穿白色時，子節點 $c$
穿白色，花 $\mathit{white}(c)$；穿紅色，花 $\mathit{red}(c) + f_c$。 兩者取小的，
每個子節點各自決定：
\[
\mathit{white}(v) = a_v + \sum_c \min\bigl(\mathit{white}(c),\ \mathit{red}(c) + f_c\bigr),
\]
$\mathit{red}(v)$ 對稱，把 $a_v$ 換成 $b_v$、白和紅對調。 答案是根的兩個值取小
的。

每組合作關係只在「子節點交給父節點」那一刻算一次。 定根時把協調費存在子節點那
一端（\texttt{fee[v]} 是 $v$ 和父節點之間的協調費），合併時就不會重複算，也不會
漏算。

用題目的範例 1 走一次。 以 $1$ 為根，$1$ 的子節點是 $2$、$3$，$2$ 的子節點是 $4$：

\begin{center}
\begin{tabular}{ccccccc}
\toprule
$v$ & $p$ & \texttt{fee[v]} & 交上去的白 & 交上去的紅 & $\mathit{white}(p)$ & $\mathit{red}(p)$ \\
\midrule
 & & & & & \multicolumn{2}{c}{初始：$1$：$(3, 5)$，$2$：$(4, 1)$} \\
4 & 2 & 1 & $\min(7, 2 + 1) = 3$ & $\min(2, 7 + 1) = 2$ & 7 & 3 \\
3 & 1 & 2 & $\min(2, 6 + 2) = 2$ & $\min(6, 2 + 2) = 4$ & 5 & 9 \\
2 & 1 & 3 & $\min(7, 3 + 3) = 6$ & $\min(3, 7 + 3) = 3$ & 11 & 12 \\
\bottomrule
\end{tabular}
\end{center}

答案是 $\min(11, 12) = 11$。 第一列：$4$ 班自己穿白要 $7$，穿紅只要 $2$；$2$ 班穿
白色時，讓 $4$ 班穿紅色再付 $1$ 元協調費，比較划算。

同樣的模型也出現在別的地方：一棟大樓的每個樓層選兩家網路業者之一，相連的樓層
用不同業者要加裝轉接設備；一個專案的各個模組選兩種框架之一，互相呼叫的模組框架
不同就要多寫轉換層。 「每個點二選一，相連的不一樣要付代價」，在樹上都可以這樣
做。

\textbf{程式}

\begin{example}[label=ex:tdp-shirts]{cpp}{(shirts.cpp) 班服顏色的搭配}
#include <algorithm>
#include <iostream>
#include <utility>
#include <vector>
using namespace std;

int main() {
    ios::sync_with_stdio(false);
    cin.tie(nullptr);
    int n;
    cin >> n;
    // white[v]、red[v]：v 穿白、穿紅時，v 的子樹裡的最小成本
    vector<long long> white(n + 1), red(n + 1);
    for (int v = 1; v <= n; v++) {
        cin >> white[v] >> red[v];          // 先放自己的班服成本
    }
    vector<vector<pair<int, long long>>> adj(n + 1);   // （鄰居, 協調費）
    for (int i = 0; i < n - 1; i++) {
        int u, v;
        long long c;
        cin >> u >> v >> c;
        adj[u].push_back({v, c});
        adj[v].push_back({u, c});
    }

    // 以 1 為根；fee[v] 是 v 和父節點之間那組合作關係的協調費
    vector<int> parent(n + 1, 0);
    vector<long long> fee(n + 1, 0);
    vector<int> order = {1};
    for (int i = 0; i < (int)order.size(); i++) {
        int u = order[i];
        for (auto [v, c] : adj[u]) {
            if (v != parent[u]) {
                parent[v] = u;
                fee[v] = c;
                order.push_back(v);
            }
        }
    }

    for (int i = n - 1; i >= 1; i--) {
        int v = order[i], p = parent[v];
        // 同色不用付，異色付一次協調費
        white[p] += min(white[v], red[v] + fee[v]);
        red[p] += min(red[v], white[v] + fee[v]);
    }
    cout << min(white[1], red[1]) << '\n';
    return 0;
}
\end{example}

\iosample{
4 \\
3 5 \\
4 1 \\
2 6 \\
7 2 \\
1 2 3 \\
1 3 2 \\
2 4 1 \\
}{
11 \\
}

\begin{itemize}
\item \texttt{white} 和 \texttt{red} 讀進來時就是自己的班服成本，也就是初始值，
不用另外開陣列存 $a_i$、$b_i$。
\item $n = 2 \times 10^5$、每項成本 $10^9$ 時總和約 $2 \times 10^{14}$，要用
\texttt{long long}。
\end{itemize}

\textbf{複雜度}

一次 BFS、一次倒著合併，都是 $O(n)$。 最大的測試資料在本書測試的環境下，CPU 時
間不到 $0.1$ 秒。

\subsection{常見錯誤}

\begin{itemize}
\item \textbf{每班各自選便宜的}：先替每班選便宜的顏色，再把兩端不同色的協調費加
上去。 題目的範例 2（就是上面三個班的例子）印出 \texttt{26}，正確答案是
\texttt{7}。 範例 1 剛好是對的，只用範例 1 測試看不出來。
\item \textbf{協調費拿錯一條}：合併時寫成 \texttt{fee[p]}，用到的是 $p$ 和它的父
節點之間的協調費。 範例 1 印出 \texttt{8}，範例 2 印出 \texttt{6}。 根的
\texttt{fee} 是 $0$，所以和根相連的合作關係都變成免費。
\item \textbf{用 \texttt{int}}：把程式裡的 \texttt{long long} 全改成
\texttt{int}。 輸入 $n = 2 \times 10^5$ 的一條鏈，每班兩種顏色都是 $10^9$、協調
費都是 $1$，正確答案是 \texttt{200000000000000}，在本書測試的環境下印出
\texttt{552894464}，加了 \texttt{-Wall -Wextra} 也沒有警告。
\end{itemize}

\subsection{什麼時候用、什麼時候不能用}

\begin{itemize}
\item 每個點有幾種選擇（選或不選、幾種顏色），而限制只發生在\textbf{相鄰}的兩點
之間：把「$v$ 選了哪一種」放進狀態，每種選擇各記一個值。 選擇有 $k$ 種時，每個子
節點要看 $k \times k$ 種組合，$k$ 小的時候沒有問題。
\item 限制牽涉到距離 $2$ 以上的點，例如「祖父和孫子不能都選」，狀態要多記一層。
\item 同樣的問題放在一般的圖上（有環），子樹之間不再互相獨立，這套做法就不能用。
一般圖上的最大獨立集沒有已知的快速算法，只能靠範圍小時的暴力解。
\end{itemize}

\begin{apcsnote}
實作題常常把測試資料分成幾個子題，第一個子題的 $n$ 很小（例如 $n \le 20$）。 樹
形 DP 的題目，$2^n$ 的暴力枚舉通常就能拿下這個子題，而且它正好是對拍用的程式。
先寫暴力解交一次，再寫樹形 DP，兩支程式在小的隨機樹上比對。
\end{apcsnote}

\section{樹上的路徑}
\label{sec:tdp-path}

第 1 本第~13~章用兩次走訪求樹直徑，條件是邊長不能是負的。 這一節換一個想法：用樹形
DP 求最長的路徑。 它比兩次走訪多寫幾行，但邊長是負的也能用，題目對路徑多加條件
時也比較容易改。

\subsection{每條路徑都有一個最高點}

以 $1$ 為根。 樹上任何一條路徑，經過的點裡深度最小的只有一個，叫它這條路徑的
\textbf{最高點}。 從最高點 $v$ 出發，路徑往兩邊走，兩邊都只能往下走（往上就會碰
到比 $v$ 更高的點）；兩邊走進的是 $v$ 的兩個\textbf{不同}的子節點，也可能某一邊
是空的，路徑就停在 $v$。

所以「以 $v$ 為最高點的最長路徑」等於：從 $v$ 往下、經過不同子節點的兩條鏈，最長
和第二長的加起來。 每個點記兩個數：

\begin{itemize}
\item $\mathit{top}_1(v)$：從 $v$ 往下走的鏈，最長的有多長。
\item $\mathit{top}_2(v)$：和最長的那條走進\textbf{不同}子節點的鏈裡，最長的有多長。
\end{itemize}

子節點 $c$ 交給 $v$ 的，是「從 $v$ 經過 $c$ 往下走的最長鏈」：$\mathit{top}_1(c)$
加上 $v$--$c$ 這條邊的長度。 $v$ 收到之後，同時更新最大和第二大。 初始值 $\mathit{top}_1 = \mathit{top}_2 = 0$，代表「停在 $v$，不往
下走」。 輪到 $v$ 交給父節點之前，它的子節點都交過了，這時用
$\mathit{top}_1(v) + \mathit{top}_2(v)$ 更新答案。 所有點的這個值取最大，就是最長
的路徑，因為每條路徑都在它的最高點被算到。

\subsection{分數最高的一段步道}

用一個邊長可以是負的問題試試看。 山上的步道網是一棵樹，每一段步道有一個分數：
風景好的是正的，正在施工、又吵又髒的是負的。 要挑一段路線（不重複經過同一點），
讓經過的步道分數總和最大；一段都不走（停在一個點）算 $0$ 分。

\begin{center}
\begin{tikzpicture}[x=1.15cm, y=1.0cm]
\node[tv] (1) at (2,3) {1};
\node[tv] (2) at (0.8,2) {2};
\node[tvalt] (3) at (3.2,2) {3};
\node[tv] (6) at (0.8,1) {6};
\node[tvhl] (4) at (2.6,1) {4};
\node[tvhl] (5) at (3.8,1) {5};
\node[tvhl] (7) at (3.8,0) {7};
\draw[te] (1) -- node[tlbl, above left] {5} (2);
\draw[te] (2) -- node[tlbl, left] {3} (6);
\draw[te] (1) -- node[tlbl, above right] {$-10$} (3);
\draw[tehl] (3) -- node[tlbl, left] {7} (4);
\draw[tehl] (3) -- node[tlbl, right] {6} (5);
\draw[tehl] (5) -- node[tlbl, right] {2} (7);
\end{tikzpicture}
\diagcaption{邊上的數字是步道的分數。 最好的路線是 $4 - 3 - 5 - 7$，$15$ 分，它的最
高點是 $3$（粉色）。}
\label{fig:tdp-trail}
\end{center}

第 1 本第~13~章的兩次走訪在這裡會錯：從 $1$ 出發，分數最高的終點是 $6$（$5 + 3 = 8$）；
從 $6$ 出發，最高也只有 $8$。 $-10$ 那一段把右邊的好路線擋住了，「先走到最遠的
點」不再保證是某條最好路線的一端。 樹形 DP 不受影響：負的鏈交上去時比 $0$ 小，
不會被選進 $\mathit{top}_1$、$\mathit{top}_2$，等於在那裡停下來。

\begin{code}{cpp}{分數最高的一段步道（由下往上的部分）}
// 以 1 為根做 BFS 時，把 v 到父節點那一段的分數存在 score[v]，
// 寫法和 shirts.cpp 的 fee 一樣
// top1[v]、top2[v]：從 v 往下、走進不同子節點的鏈，最高與第二高的分數
// 初始值 0 代表「停在 v 不往下走」
vector<long long> top1(n + 1, 0), top2(n + 1, 0);
long long best = 0;
for (int i = n - 1; i >= 0; i--) {
    int v = order[i];
    best = max(best, top1[v] + top2[v]);    // 以 v 為最高點的最好路線
    if (i == 0) {
        break;                              // 根沒有父節點
    }
    int p = parent[v];
    long long d = top1[v] + score[v];       // 從 p 經過 v 往下的最好的鏈
    if (d > top1[p]) {
        top2[p] = top1[p];
        top1[p] = d;
    } else if (d > top2[p]) {
        top2[p] = d;
    }
}
\end{code}

圖~\ref{fig:tdp-trail} 的輸入是 $n = 7$ 和六行「$u$ $v$ 分數」：\texttt{1 2 5}、
\texttt{2 6 3}、\texttt{1 3 -10}、\texttt{3 4 7}、\texttt{3 5 6}、\texttt{5 7 2}，
答案是 $15$。 BFS 順序是 $1, 2, 3, 6, 4, 5, 7$，倒著看：

\begin{center}
\begin{tabular}{cccccc}
\toprule
$v$ & $(\mathit{top}_1, \mathit{top}_2)$ & 用來更新 \texttt{best} & 交給 & 交上去的 \texttt{d} & 結果 \\
\midrule
7 & $(0, 0)$ & 0 & 5 & $0 + 2 = 2$ & $\mathit{top}_1(5) = 2$ \\
5 & $(2, 0)$ & 2 & 3 & $2 + 6 = 8$ & $\mathit{top}_1(3) = 8$ \\
4 & $(0, 0)$ & 0 & 3 & $0 + 7 = 7$ & $\mathit{top}_2(3) = 7$ \\
6 & $(0, 0)$ & 0 & 2 & $0 + 3 = 3$ & $\mathit{top}_1(2) = 3$ \\
3 & $(8, 7)$ & \textbf{15} & 1 & $8 - 10 = -2$ & 比 $0$ 小，不採用 \\
2 & $(3, 0)$ & 3 & 1 & $3 + 5 = 8$ & $\mathit{top}_1(1) = 8$ \\
1 & $(8, 0)$ & 8 & -- & -- & -- \\
\bottomrule
\end{tabular}
\end{center}

\begin{itemize}
\item 迴圈這次跑到 \texttt{i = 0}：根也要用 $\mathit{top}_1 + \mathit{top}_2$ 更
新答案，只是不用再往上交。
\item 更新最大和第二大時，新的值比最大的還大，原本最大的要先降成第二大，順序不能
寫反。
\item 所有分數都是 $1$ 時，算的就是第 1 本第~13~章的樹直徑（邊數）。
\end{itemize}

\subsection{常見錯誤}

\begin{itemize}
\item \textbf{只記最大值，算兩次}：答案寫成 \texttt{2 * top1[v]}，想用「最長的鏈
走兩次」代替兩條鏈。 上面的輸入印出 \texttt{16}：它把 $3 - 5 - 7$ 這條鏈當成兩條，可是
兩條鏈要走進不同的子節點。
\item \textbf{第一名換人時，沒把舊的降成第二名}：\texttt{d > top1[p]} 的時候只寫
\texttt{top1[p] = d;}。 上面的輸入剛好先交 $8$ 再交 $7$，答案還是 \texttt{15}；把
輸入的 \texttt{3 4 7} 和 \texttt{3 5 6} 兩行對調，$3$ 會先收到 $7$ 再收到 $8$，
$7$ 就被丟掉，印出 \texttt{8}。 結果取決於子節點交上來的順序，範例過了不代表對。
\end{itemize}

\subsection{例題：翻過山頂的步道}
\label{sec:tdp-ex-summit}

\textbf{題意}（\akiref{over-the-summit}{翻過山頂的步道}）

\begin{itemize}
\item 山上有 $n$ 個休息點，由 $n - 1$ 段步道連成一棵樹，休息點 $i$ 的海拔是
$h_i$，不同的點海拔可能相同。
\item 一條翻山路線不重複經過同一點，海拔先嚴格上升到某個山頂，再嚴格下降。 只上
升、只下降、只有一個點，都算。
\item 求翻山路線最多經過幾個休息點。 $n \le 2 \times 10^5$，$h_i \le 10^9$。
\end{itemize}

下圖是題目的範例 1，括號裡是海拔。 最長的翻山路線是 $4 \to 2 \to 1 \to 3 \to 6$，
海拔 $2, 7, 5, 4, 1$，山頂是 $2$，經過 $5$ 個點。 $3$ 和 $7$ 一樣高，不能接在一
起。

\begin{center}
\begin{tikzpicture}[x=1.15cm, y=1.0cm]
\node[tvhl] (1) at (3,3) {1};
\node[tvalt] (2) at (1.5,2) {2};
\node[tvhl] (3) at (4.5,2) {3};
\node[tvhl] (4) at (0.8,1) {4};
\node[tv] (5) at (2.2,1) {5};
\node[tvhl] (6) at (3.8,1) {6};
\node[tv] (7) at (5.2,1) {7};
\draw[tehl] (1) -- (2);
\draw[tehl] (1) -- (3);
\draw[tehl] (2) -- (4);
\draw[te] (2) -- (5);
\draw[tehl] (3) -- (6);
\draw[te] (3) -- (7);
\foreach \v/\h in {1/5, 3/4} {\node[tlbl, right] at (\v.east) {(\h)};}
\node[tlbl, left] at (2.west) {(7)};
\foreach \v/\h in {4/2, 5/6, 6/1, 7/4} {\node[tlbl, below] at (\v.south) {(\h)};}
\end{tikzpicture}
\diagcaption{範例 1。 粉色的 $2$ 是山頂，橘色的路線從 $2$ 往兩邊下降：一邊到 $4$，
另一邊經過父節點 $1$ 再到 $3$、$6$。}
\label{fig:tdp-summit}
\end{center}

\textbf{想法}

最直接的做法：從每個點出發往外走，帶著「還在上升」或「已經開始下降」的狀態，試
遍所有走法。 一條長長的山脊上，每個起點都可能走 $O(n)$ 步，總共 $O(n^2)$，$n =
2 \times 10^5$ 時來不及，但可以拿來對拍。

這題很像上一小節：一條翻山路線也有一個特別的點，就是山頂。 把路線從山頂切開，兩
邊都是「從山頂出發、海拔一路嚴格下降」的鏈，而且走進山頂的兩個不同的鄰居。 所以
答案是：每個點當山頂，往下的鏈取最長和第二長，加上山頂自己。

最自然的寫法是照上一小節以 $1$ 為根，$\mathit{top}_1$、$\mathit{top}_2$ 只收「比
自己低的子節點」交上來的鏈。 範例 1 會得到 $3$：山頂 $2$ 往下的兩條鏈，一條是
$2 \to 4$，另一條要先往\textbf{父節點} $1$ 走，可是以 $1$ 為根時，$1$ 不是 $2$ 的
子節點，這條鏈不會被交給 $2$。 上一小節的「往下」是離根越來越遠；這一題的「往下」
是海拔越來越低，方向由海拔決定，和根是誰無關。

所以換一個狀態，不分父子：$\mathit{down}(v)$ 是從 $v$ 出發、海拔一路嚴格下降的
鏈，最多經過幾個點（含 $v$）。 它只走向比 $v$ 低的鄰居：
\[
\mathit{down}(v) = 1 + \max\{\mathit{down}(u) : u \text{ 是 } v \text{ 的鄰居，} h_u < h_v\},
\]
沒有比 $v$ 低的鄰居時是 $1$。 計算順序：算 $v$ 時，比它低的鄰居要先算好，依海拔
由低到高排序就行了。 把每條步道想成從高處指向低處的箭頭，這就是第~6~章有向無環
圖的拓樸順序；海拔相同的兩點之間沒有箭頭，誰先誰後都可以。

以 $v$ 為山頂的最長路線是 $1 + $ 比 $v$ 低的鄰居中 $\mathit{down}$ 最大的加上第二
大的。 兩條鏈為什麼不會碰在一起？ 拿掉 $v$ 之後，樹斷成好幾塊，每個鄰居各在一
塊。 一條從 $v$ 往下的鏈，走進鄰居 $u$ 之後海拔都比 $v$ 低，不會再回到 $v$，所以
整條都留在 $u$ 那一塊。 兩條鏈走進不同的鄰居，就在不同的塊裡，沒有共同的點。

範例 1 依海拔由低到高計算：

\begin{center}
\begin{tabular}{ccccc}
\toprule
$v$（海拔） & 比 $v$ 低的鄰居 & 它們的 $\mathit{down}$ & $\mathit{down}(v)$ & 當山頂 \\
\midrule
6 (1) & 無 & -- & 1 & 1 \\
4 (2) & 無 & -- & 1 & 1 \\
3 (4) & 6 & 1 & 2 & 2 \\
7 (4) & 無（3 一樣高） & -- & 1 & 1 \\
1 (5) & 3 & 2 & 3 & 3 \\
5 (6) & 無 & -- & 1 & 1 \\
2 (7) & 1、4、5 & 3、1、1 & 4 & $1 + 3 + 1 = \mathbf{5}$ \\
\bottomrule
\end{tabular}
\end{center}

最後一列就是圖~\ref{fig:tdp-summit}：山頂 $2$ 往 $1$ 那邊的鏈 $1 \to 3 \to 6$ 有
$3$ 個點，往 $4$ 那邊有 $1$ 個點。

這一題和上一小節是同一個想法，換了一種「最高點」：上一小節的最高點是離根最近的
點，這一題是海拔最高的點。 只要每條要找的路徑都能在某個唯一的點切成兩條鏈，就
能在那個點把最好的兩條接起來。 同樣的模型也出現在別的地方：自行車路線想找一段先
爬坡、再下坡的練習路線；滑雪場的纜車一路往上、滑道一路往下，問一趟最多經過幾個
站。

\textbf{程式}

\begin{example}[label=ex:tdp-summit]{cpp}{(summit.cpp) 翻過山頂的步道}
#include <algorithm>
#include <iostream>
#include <vector>
using namespace std;

int main() {
    ios::sync_with_stdio(false);
    cin.tie(nullptr);
    int n;
    cin >> n;
    vector<long long> h(n + 1);
    for (int v = 1; v <= n; v++) {
        cin >> h[v];
    }
    vector<vector<int>> adj(n + 1);
    for (int i = 0; i < n - 1; i++) {
        int a, b;
        cin >> a >> b;
        adj[a].push_back(b);
        adj[b].push_back(a);
    }

    // 依海拔由低到高排出計算順序
    vector<int> byHeight(n);
    for (int i = 0; i < n; i++) {
        byHeight[i] = i + 1;
    }
    sort(byHeight.begin(), byHeight.end(),
         [&](int x, int y) { return h[x] < h[y]; });

    // down[v]：從 v 出發一路嚴格下降，最多經過幾個休息點（含 v）
    vector<int> down(n + 1, 1);
    int best = 1;
    for (int v : byHeight) {
        int top1 = 0, top2 = 0;             // 比 v 低的鄰居裡，最大與第二大的 down
        for (int u : adj[v]) {
            if (h[u] >= h[v]) {
                continue;                   // 不是往下走
            }
            if (down[u] > top1) {
                top2 = top1;
                top1 = down[u];
            } else if (down[u] > top2) {
                top2 = down[u];
            }
        }
        down[v] = 1 + top1;
        best = max(best, 1 + top1 + top2);  // v 當山頂
    }
    cout << best << '\n';
    return 0;
}
\end{example}

\iosample{
7 \\
5 7 4 2 6 1 4 \\
1 2 \\
1 3 \\
2 4 \\
2 5 \\
3 6 \\
3 7 \\
}{
5 \\
}

\begin{itemize}
\item 這支程式沒有定根，也沒有 BFS：計算順序來自海拔，不是來自樹的上下。 它仍然
是樹形 DP，因為「兩條鏈不相交」用到的是樹的性質。
\item 排序的 lambda 用 \texttt{[\&]} 捕捉 \texttt{h}（銜接書第~9~章）。 海拔相同
的點誰先都可以，所以不用比編號。
\item 每個點看一遍自己的鄰居，全部加起來是 $2(n - 1)$ 次。
\end{itemize}

\textbf{複雜度}

排序 $O(n \log n)$，之後 $O(n)$。 最大的測試資料在本書測試的環境下，CPU 時間不到
$0.1$ 秒。

\subsection{常見錯誤}

\begin{itemize}
\item \textbf{以 $1$ 為根、只看子節點}：照上一小節的寫法，只收比自己低的子節點交
上來的鏈。 範例 1 印出 \texttt{3}，原因見「想法」。 山頂在葉子附近、往根的方向
一路下降時，這種寫法一定會漏。
\item \textbf{海拔相同也當成能走}：把 \texttt{h[u] >= h[v]} 寫成
\texttt{h[u] > h[v]}。 題目的範例 2（一條直線，前兩個點海拔都是 $4$）正確答案
是 \texttt{4}，在本書測試的環境下印出 \texttt{6}，比真正能走的還長。 兩個一樣高
的點會互相把對方算進自己的 $\mathit{down}$，等於在兩點之間來回走。
\item \textbf{只記最大值，算兩次}：寫成 \texttt{1 + 2 * top1}，範例 1 印出
\texttt{7}。 和上一小節一樣，兩條鏈要走進不同的鄰居。
\end{itemize}

\subsection{什麼時候用、什麼時候不能用}

\begin{itemize}
\item 題目問「樹上最好的一條路徑」，而且每條路徑都能在一個唯一的點切成兩條往下的
鏈：在那個點取最好和第二好的鏈接起來。 邊長有負的、要對每個子樹各求一次、或要
數最長路徑有幾條時，兩次走訪不夠用，就用這個做法。
\item 要求「從每個點出發的最長路徑」，一次由下往上只看得到往下的鏈，還要知道往父
節點那邊能走多遠，這是下一節換根的工作（本章練習第~\ref{drill:tdp-ecc}~題）。
\end{itemize}

\section{換根 DP}
\label{sec:tdp-reroot}

到目前為止，答案都只要算一次：以某個點為根，由下往上算到根就結束。 有些題目要
問「每個點」：

\begin{quote}
社區有 $n$ 戶人家，由 $n - 1$ 條巷子連成一棵樹，每條巷子走一次算一步。 過年
前，每一戶都要把一盒年菜分別送到其他每一戶，一戶一趟，回程不算。 對每一戶 $v$，
求出 $v$ 送完所有年菜一共要走幾步，也就是 $v$ 到其他每一戶的步數總和。
\end{quote}

最直接的做法是每一戶各做一次 BFS，加總距離，$O(n^2)$。 $n = 2 \times 10^5$ 時
是 $4 \times 10^{10}$ 步，來不及。

\subsection{先算一個根，再換到旁邊}

先只算一戶。 以 $1$ 為根，第 1 本第~13~章的子樹大小 $\mathit{sz}(v)$ 已經會算；再記
$\mathit{sub}(v)$：$v$ 到它子樹裡每一戶的步數總和。 從 $v$ 走到子節點 $c$ 的子
樹裡的每一戶，都是先走一步到 $c$，再走 $c$ 到那一戶的步數，所以
\[
\mathit{sub}(v) = \sum_c \bigl(\mathit{sub}(c) + \mathit{sz}(c)\bigr).
\]
倒著看 BFS 順序算完，$\mathit{sub}(1)$ 就是 $1$ 的答案。

接下來是關鍵：知道 $p$ 的答案，它的子節點 $c$ 的答案差多少？ 拿掉
$p$--$c$ 這條巷子，樹分成兩塊：$c$ 的子樹有 $\mathit{sz}(c)$ 戶，其他 $n -
\mathit{sz}(c)$ 戶在 $p$ 那一邊。 送年菜的人從 $p$ 換成 $c$，到 $c$ 那一塊的每
戶都少走一步，到 $p$ 那一塊的每戶都多走一步：
\[
\mathit{ans}(c) = \mathit{ans}(p) - \mathit{sz}(c) + \bigl(n - \mathit{sz}(c)\bigr).
\]

\begin{center}
\begin{tikzpicture}[x=1.15cm, y=1.0cm]
\node[tv] (1) at (2,3) {1};
\node[tvalt] (2) at (1,2) {2};
\node[tv] (3) at (3,2) {3};
\node[tv] (4) at (0.4,1) {4};
\node[tvhl] (5) at (1.6,1) {5};
\node[tv] (7) at (3,1) {7};
\node[tvhl] (6) at (1.6,0) {6};
\foreach \a/\b in {1/2, 1/3, 2/4, 3/7, 5/6} {\draw[te] (\a) -- (\b);}
\draw[tehl, dashed] (2) -- (5);
\node[tlbl, left] at (2.west) {$\mathit{ans}(2) = 10$};
\node[tlbl, right] at (6.east) {$5$、$6$ 各少走一步};
\end{tikzpicture}
\diagcaption{送年菜的人從 $2$（粉色）換成 $5$。 到 $5$ 的子樹（橘色，$2$ 戶）每戶
少走一步，到其他 $5$ 戶每戶多走一步，$\mathit{ans}(5) = 10 - 2 + 5 = 13$。}
\label{fig:tdp-reroot}
\end{center}

這個式子要先知道父節點的答案，所以第二趟是\textbf{由上往下}，順著 BFS 順序算。
兩趟加起來，每個點的答案都在 $O(n)$ 內算完。 這種「先以一個點為根算好，再把根一
步一步換到鄰居」的做法叫換根 DP。

\begin{example}[label=ex:tdp-distsum]{cpp}{(new\_year.cpp) 每一戶到其他戶的步數總和}
#include <iostream>
#include <vector>
using namespace std;

int main() {
    ios::sync_with_stdio(false);
    cin.tie(nullptr);
    int n;
    cin >> n;
    vector<vector<int>> adj(n + 1);
    for (int i = 0; i < n - 1; i++) {
        int u, v;
        cin >> u >> v;
        adj[u].push_back(v);
        adj[v].push_back(u);
    }

    vector<int> parent(n + 1, 0);
    vector<int> order = {1};
    for (int i = 0; i < (int)order.size(); i++) {
        int u = order[i];
        for (int v : adj[u]) {
            if (v != parent[u]) {
                parent[v] = u;
                order.push_back(v);
            }
        }
    }

    // 第一趟，由下往上：sz[v] 子樹大小，sub[v] 是 v 到子樹裡每戶的步數總和
    vector<long long> sz(n + 1, 1), sub(n + 1, 0);
    for (int i = n - 1; i >= 1; i--) {
        int v = order[i], p = parent[v];
        sz[p] += sz[v];
        sub[p] += sub[v] + sz[v];           // 從 p 出發，每戶都多 v 那 1 步
    }

    // 第二趟，由上往下：送年菜的人從 p 換成 v
    vector<long long> ans(n + 1);
    ans[1] = sub[1];
    for (int i = 1; i < n; i++) {
        int v = order[i], p = parent[v];
        ans[v] = ans[p] - sz[v] + (n - sz[v]);
    }
    for (int v = 1; v <= n; v++) {
        cout << ans[v] << (v == n ? '\n' : ' ');
    }
    return 0;
}
\end{example}

\iosample{
7 \\
1 2 \\
1 3 \\
2 4 \\
2 5 \\
5 6 \\
3 7 \\
}{
11 10 14 15 13 18 19 \\
}

範例就是圖~\ref{fig:tdp-reroot} 的樹，BFS 順序是 $1, 2, 3, 4, 5, 7, 6$。 兩趟算
出來的值：

\begin{center}
\begin{tabular}{c|ccccccc}
\toprule
$v$ & 1 & 2 & 3 & 4 & 5 & 6 & 7 \\
\midrule
$\mathit{sz}(v)$ & 7 & 4 & 2 & 1 & 2 & 1 & 1 \\
$\mathit{sub}(v)$ & 11 & 4 & 1 & 0 & 1 & 0 & 0 \\
$\mathit{ans}(v)$ & 11 & $11 - 4 + 3 = 10$ & $11 - 2 + 5 = 14$ & 15 & 13 & 18 & 19 \\
\bottomrule
\end{tabular}
\end{center}

$2$ 位在比較中間，送完只要 $10$ 步；$7$ 在邊上，要 $19$ 步。

\begin{itemize}
\item 第二趟的 \texttt{i} 從 $1$ 開始往後走，跳過根。 \texttt{order} 裡父節點一定
排在前面，所以輪到 $v$ 時 \texttt{ans[p]} 已經算好。
\item \texttt{sz} 也開成 \texttt{long long}，\texttt{n - sz[v]} 就直接用
\texttt{long long} 算。 一條 $2 \times 10^5$ 個點的鏈，端點的答案是
$n(n-1)/2 \approx 2 \times 10^{10}$，超過 \texttt{int}。
\end{itemize}

\subsection{一般的換根：父節點那一側}

距離和的換根式子很短，是因為「$c$ 那一塊少一步、其他多一步」剛好能直接寫出來。
一般的題目沒有這麼巧，要換一個角度想。

以 $v$ 為根時，$v$ 的鄰居分成兩種：以 $1$ 為根時的子節點，和父節點 $p$。 子節點
那幾塊，就是第一趟算好的子樹，記作 $\mathit{down}(c)$。 父節點那一塊，是「整棵
樹拿掉 $v$ 的子樹」，從 $v$ 看過去，它也是一棵掛在 $v$ 下面的子樹，根是 $p$。
把它的值記作 $\mathit{up}(v)$。 有了這兩種，$v$ 當根的答案就是用 $v$ 自己的值，
合併所有 $\mathit{down}(c)$ 和 $\mathit{up}(v)$，和第一趟的合併方式一樣。

\begin{center}
\begin{tikzpicture}[x=1cm, y=1cm]
\draw[dashed, optStroke, thick] (3,3.35) ellipse (2.4cm and 0.85cm);
\node[tv] (p) at (3,3.2) {$p$};
\node[tv] (s1) at (1.7,3.4) {};
\node[tv] (s2) at (4.3,3.4) {};
\draw[te] (p) -- (s1);
\draw[te] (p) -- (s2);
\node[tlbl, right] at (5.6,3.55) {$\mathit{up}(v)$：拿掉 $v$ 的子樹後};
\node[tlbl, right] at (5.6,3.15) {剩下的部分，根是 $p$};
\node[tvhl] (v) at (3,1.9) {$v$};
\draw[te] (p) -- (v);
\foreach \x/\c in {1.8/1, 4.2/2} {
    \node[tv] (c\c) at (\x,0.9) {$c_\c$};
    \draw[te] (v) -- (c\c);
    \draw[black!50] (c\c.south) -- (\x-0.6,-0.3) -- (\x+0.6,-0.3) -- cycle;
    \node[tlbl] at (\x,-0.55) {$\mathit{down}(c_\c)$};
}
\end{tikzpicture}
\diagcaption{以 $v$ 為根時，$v$ 的每個鄰居各帶著一塊：子節點那幾塊是第一趟的
$\mathit{down}$，父節點那一塊是 $\mathit{up}(v)$。}
\label{fig:tdp-up}
\end{center}

$\mathit{up}(v)$ 怎麼算？ 那一塊的根是 $p$，$p$ 在那一塊裡的鄰居是：$p$ 的父節點
那一側（$\mathit{up}(p)$），以及 $p$ 除了 $v$ 以外的子節點。 所以
$\mathit{up}(v)$ 就是「$p$ 自己的值，合併 $\mathit{up}(p)$ 和 $p$ 除了 $v$ 以外
所有子節點的 $\mathit{down}$」。 它要用到 $\mathit{up}(p)$，所以由上往下算。

難的地方是「除了 $v$ 以外」。 $p$ 可能有很多子節點，每個子節點都重新合併一次其他
兄弟，星狀的樹就變成 $O(n^2)$。 要看合併方式：

\begin{itemize}
\item 合併是加總：先算好所有子節點的總和，扣掉 $v$ 的那一項就好，$O(1)$。
\item 合併是取最大：最大值沒辦法「扣掉」。 改記最大和第二大，以及最大的那一個來
自哪個子節點；如果 $v$ 正好是最大的那一個，就用第二大（本章練習第
~\ref{drill:tdp-ecc}~題）。
\end{itemize}

\subsection{例題：選在哪裡發電}
\label{sec:tdp-ex-power}

\textbf{題意}（\akiref{where-to-build-the-power-station}{選在哪裡發電}）

\begin{itemize}
\item 山上的村子有 $n$ 個用電地點，用 $n - 1$ 條電線連成一棵樹。 地點 $i$ 一天自己
發 $a_i$ 度電、要用 $b_i$ 度電。
\item 總站蓋在地點 $r$ 時，可以從山下引進 $X$ 度電到 $r$。 每條電線只能從離總站近
的一端送到遠的一端；每個地點可以把手上的電自己用、往外送，或丟掉。
\item 對每一個 $r$，求出讓每個地點都用得到 $b_i$ 度電的最小 $X$。
\item $n \le 2 \times 10^5$，$a_i, b_i \le 10^9$。
\end{itemize}

例如四個地點連成一條線 $1 - 2 - 3 - 4$，發電量是 $3, 0, 4, 0$，用電量是 $0, 2, 0,
3$。 總站蓋在 $4$ 時，$3$ 多的 $4$ 度送不回 $4$，$4$ 要的 $3$ 度只能全部引進，答
案是 $3$。 蓋在 $3$ 時，$3$ 要送 $3$ 度給 $4$、$2$ 度給 $2$（$1$ 多的電送不到 $2$），
自己發的 $4$ 度不夠，要引進 $1$ 度。 四個答案依序是 $0, 2, 1, 3$。

\textbf{想法}

先固定總站 $r$，以 $r$ 為根。 記 $d_v = b_v - a_v$，是 $v$ 自己缺多少電（負的表
示有多）。 定義 $\mathit{need}(v)$：$v$ 的子樹整塊向 $v$ 的父節點要多少電。 $v$
要送給每個子節點 $c$ 剛好 $\mathit{need}(c)$ 度（子節點之間不能互相支援，也不能往
回送），自己還缺 $d_v$，所以 $v$ 一共缺 $d_v + \sum_c \mathit{need}(c)$。 這個
數是負的時候，多出來的電送不上去，只能丟掉，向上要的是 $0$：
\[
\mathit{need}(v) = \max\Bigl(0,\ d_v + \sum_c \mathit{need}(c)\Bigr).
\]
總站自己也照這個式子算，$X = \mathit{need}(r)$。

這個 $\max(0, \cdot)$ 是整題的重點。 少了它，某個子樹多的電就會被拿去補別的子樹，
可是電只能往外送，做不到。 所以「全部用電減全部發電」不是答案：題目的範例 1 用電
和發電一樣多，五個答案裡卻只有一個是 $0$。

對一個 $r$ 這是一次樹形 DP，$O(n)$。 每個 $r$ 各做一次是 $O(n^2)$，只能過小的子
題，也可以拿來對拍。

換根。 以 $1$ 為根做第一趟，算出 $\mathit{down}(v)$（就是以 $1$ 為根時的
$\mathit{need}(v)$），並記下 $\mathit{kids}(v)$：$v$ 所有子節點的
$\mathit{down}$ 加起來。 第二趟由上往下，照上一小節的做法算 $\mathit{up}(v)$：
那一塊的根是 $p$，$p$ 收到的是 $\mathit{up}(p)$ 和除了 $v$ 以外的子節點，所以
\[
\mathit{up}(v) = \max\Bigl(0,\ d_p + \bigl(\mathit{kids}(p) - \mathit{down}(v)\bigr) + \mathit{up}(p)\Bigr),
\]
根的 $\mathit{up}$ 是 $0$。 總站蓋在 $v$ 的答案是
\[
\mathit{ans}(v) = \max\bigl(0,\ d_v + \mathit{kids}(v) + \mathit{up}(v)\bigr).
\]

這裡能直接「扣掉 $v$ 的那一項」，是因為 $\max(0, \cdot)$ 是在加總\textbf{之後}才
做的。 $\mathit{kids}(p)$ 是一些數單純相加，拿掉其中一個 $\mathit{down}(v)$，剩
下的就是其他子節點的總和，一點誤差都沒有；扣完再取 $\max(0, \cdot)$，就是那一塊
的值。

用題目的範例 1 走一次。 電線是 $1$--$2$、$1$--$3$、$3$--$4$、$3$--$5$，$d =
(-1, 4, 0, -5, 2)$，BFS 順序是 $1, 2, 3, 4, 5$：

\begin{center}
\begin{tabular}{c|ccccc}
\toprule
$v$ & 1 & 2 & 3 & 4 & 5 \\
\midrule
$d_v$ & $-1$ & 4 & 0 & $-5$ & 2 \\
$\mathit{down}(v)$ & 5 & 4 & 2 & 0 & 2 \\
$\mathit{kids}(v)$ & 6 & 0 & 2 & 0 & 0 \\
$\mathit{up}(v)$ & 0 & 1 & 3 & 5 & 3 \\
$\mathit{ans}(v)$ & 5 & 5 & 5 & 0 & 5 \\
\bottomrule
\end{tabular}
\end{center}

例如 $\mathit{up}(4) = \max(0,\ d_3 + \mathit{kids}(3) - \mathit{down}(4) +
\mathit{up}(3)) = 0 + 2 - 0 + 3 = 5$，$\mathit{ans}(4) = \max(0,\ -5 + 0 + 5) = 0$。
$\mathit{down}(4) = 0$：$4$ 多出 $5$ 度，以 $1$ 為根時送不上去。 換到以 $4$ 為根，
$4$ 那 $5$ 度剛好送給 $\mathit{up}(4) = 5$ 的那一塊，答案變成 $0$。

同樣的模型也出現在別的地方：水只能從水塔往下游流，要決定水塔蓋在哪裡、每個候選
位置要補多少水；物資只能從倉庫往外配送、不能退回，比較每個倉庫地點要額外進多少
貨。

\textbf{程式}

\begin{example}[label=ex:tdp-power]{cpp}{(power.cpp) 選在哪裡發電}
#include <algorithm>
#include <iostream>
#include <vector>
using namespace std;

int main() {
    ios::sync_with_stdio(false);
    cin.tie(nullptr);
    int n;
    cin >> n;
    vector<long long> d(n + 1);             // 用電減發電，負的表示有多的電
    for (int v = 1; v <= n; v++) {
        long long a, b;
        cin >> a >> b;
        d[v] = b - a;
    }
    vector<vector<int>> adj(n + 1);
    for (int i = 0; i < n - 1; i++) {
        int u, v;
        cin >> u >> v;
        adj[u].push_back(v);
        adj[v].push_back(u);
    }

    vector<int> parent(n + 1, 0);
    vector<int> order = {1};
    for (int i = 0; i < (int)order.size(); i++) {
        int u = order[i];
        for (int v : adj[u]) {
            if (v != parent[u]) {
                parent[v] = u;
                order.push_back(v);
            }
        }
    }

    // 第一趟，由下往上（以 1 為根）
    // kids[v]：v 所有子節點的 down 加總；down[v]：v 的子樹向父節點要的電
    vector<long long> kids(n + 1, 0), down(n + 1);
    for (int i = n - 1; i >= 0; i--) {
        int v = order[i];
        down[v] = max(0LL, d[v] + kids[v]);
        if (i > 0) {
            kids[parent[v]] += down[v];
        }
    }

    // 第二趟，由上往下
    // up[v]：把父節點那一側整塊看成 v 的子樹時，它向 v 要的電
    vector<long long> up(n + 1, 0), ans(n + 1);
    for (int i = 0; i < n; i++) {
        int v = order[i];
        if (i > 0) {
            int p = parent[v];
            up[v] = max(0LL, d[p] + kids[p] - down[v] + up[p]);
        }
        ans[v] = max(0LL, d[v] + kids[v] + up[v]);
    }
    for (int v = 1; v <= n; v++) {
        cout << ans[v] << (v == n ? '\n' : ' ');
    }
    return 0;
}
\end{example}

\iosample{
5 \\
2 1 \\
0 4 \\
1 1 \\
5 0 \\
0 2 \\
1 2 \\
1 3 \\
3 4 \\
5 3 \\
}{
5 5 5 0 5 \\
}

\begin{itemize}
\item 第一趟跑到 \texttt{i = 0}，根的 \texttt{down} 也算出來，但第二趟沒有用到它：
根的答案直接由 \texttt{kids[1]} 和 \texttt{up[1] = 0} 算。
\item \texttt{max(0LL, ...)} 的 \texttt{0LL} 是 \texttt{long long} 的 $0$，和第 1 本第~7
~章說過的一樣，\texttt{max} 的兩個引數型別要相同。
\item 答案最大約 $2 \times 10^5 \times 10^9 = 2 \times 10^{14}$，全部用
\texttt{long long}。
\end{itemize}

\textbf{複雜度}

兩趟各 $O(n)$。 最大的測試資料在本書測試的環境下，CPU 時間不到 $0.1$ 秒。

\subsection{常見錯誤}

\begin{itemize}
\item \textbf{$\mathit{up}$ 沒有取 $\max(0, \cdot)$}：寫成 \texttt{up[v] = d[p] +
kids[p] - down[v] + up[p];}。 範例 1 剛好不受影響，題目的範例 2（就是上面四個點的
直線）正確答案是 \texttt{0 2 1 3}，這樣寫印出 \texttt{0 0 0 0}。 以 $2$ 為總站時，
$\mathit{up}(2) = -3 + 0 + 0 = -3$，等於讓 $1$ 多的電倒流回 $2$。
\item \textbf{用總量相減}：輸出 $\max(0, \sum b_i - \sum a_i)$。 範例 1 每個答案都
印成 \texttt{0}，範例 2 也全是 \texttt{0}。
\item \textbf{每個子節點重算兄弟}：算 $\mathit{up}(v)$ 時用迴圈把 $p$ 的其他子節點
重新加一遍，不用 \texttt{kids[p] - down[v]}。 結果是對的，但一個點有 $2 \times
10^5$ 個子節點的星狀樹上，迴圈一共跑大約 $n^2 = 4 \times 10^{10}$ 次，遠超過時間
限制。 這題的測試資料裡有這樣的樹。
\end{itemize}

\begin{apcsnote}
換根題的小子題，通常是「每個點各做一次」的 $O(n^2)$ 就能過的範圍。 先把這支寫好，
它就是對拍用的程式；換根的式子寫錯時，常常只有某些點的答案不對，和暴力解逐點比
對，最容易找到是哪一個點、哪一條式子出錯。
\end{apcsnote}

\subsection{什麼時候用、什麼時候不能用}

\begin{itemize}
\item 題目要求「每個點當根（中心、總站、集合地點）時」的答案，而且固定一個根時是
一次樹形 DP：先想單一根的式子，再想怎麼把根換到鄰居。
\item 合併是加總（或任何可以「扣掉一項」的運算）時，$\mathit{up}$ 可以 $O(1)$ 算；
合併是取最大、最小時，記最好和第二好。
\item 合併方式既不能扣、也不能只記前兩名（例如要排序所有子節點的值），換根就沒
有這麼簡單，要用前綴、後綴的技巧，留到《C++ 演算法解題 3》。
\end{itemize}

\section{預告：依子樹大小合併}
\label{sec:tdp-preview}

到這裡為止，每個點的狀態都是幾個數。 有一類題目，狀態是一整張表。 例如一棵技能
樹，要學一個技能，得先學它的上層技能；每個技能有一個價值，最多只能學 $k$ 個，價
值總和最多是多少？ 狀態要記「在 $v$ 的子樹裡學幾個」：$f_v[j]$ 是 $v$ 的子樹裡學
$j$ 個技能時的最大價值。 $j \ge 1$ 時 $v$ 自己一定有學；$f_v[0] = 0$，表示整個子
樹都不學。 把子節點 $c$ 的表併進 $v$，是第~1~章
背包的合併：$v$ 這邊用掉 $a$ 個、$c$ 的子樹用掉 $b$ 個，加起來是 $a + b$ 個。

\begin{code}{cpp}{兩張表的合併（s 是 v 目前併過的點數）}
// f：v 和已經併進來的子樹，長度 s + 1；g：子節點 c 的表，長度 sz[c] + 1
vector<long long> h(s + sz[c] + 1, NEG);   // NEG 表示做不到
h[0] = 0;                                  // 整個子樹都不學
for (int a = 1; a <= s; a++) {             // a >= 1：v 自己有學，c 才能學
    for (int b = 0; b <= sz[c]; b++) {
        h[a + b] = max(h[a + b], f[a] + g[b]);
    }
}
f = h;
s += sz[c];
\end{code}

一次合併的雙層迴圈跑 $s \times \mathit{sz}(c)$ 次左右。 如果每張表都開成 $k + 1$
格、每次都跑滿，一次合併是 $k^2$，$n$ 次就是 $n k^2$；$n = k = 3000$ 時是 $2.7
\times 10^{10}$，來不及。

關鍵是表的長度只開到\textbf{子樹大小}：$v$ 的子樹只有 $\mathit{sz}(v)$ 個點，
不可能學超過 $\mathit{sz}(v)$ 個。 這樣一來，整棵樹所有合併加起來不超過大約
$n^2$ 次。 理由是把 $s \times \mathit{sz}(c)$ 看成「已經併進來的點」和「$c$ 的
子樹裡的點」兩兩配對的對數；兩個點只會在它們最低的共同祖先那裡、被合併的那一次
配成一對，所以總數不超過 $n(n-1)/2$ 對。 $n = 3000$ 時是 $4.5 \times 10^6$，很
快。

這個做法和第~1~章的背包合在一起，寫起來要注意的細節比較多：表的長度怎麼限制、
「做不到」怎麼表示、合併的順序。 這一章只說到這裡，實作練習的「修剪預算」用得到
同樣的想法，它的表的長度是用時間的上限來限制的。

\section{本章常見錯誤}
\label{sec:tdp-pitfalls}

\begin{center}
\begin{tabular}{p{0.22\linewidth}p{0.46\linewidth}p{0.2\linewidth}}
\toprule
錯誤 & 例子 & 見 \\
\midrule
狀態不夠 & 只記「子樹最多幾分」，不記 $v$ 選了沒有；路徑只記最長的鏈，答案寫成兩倍 &
第~\ref{sec:tdp-mis}、\ref{sec:tdp-path}~節 \\
轉移多加了限制 & $v$ 不選就要求子節點一定選 & 第~\ref{sec:tdp-mis}~節 \\
同一條邊算錯次數 & 協調費用了父節點那一條；兩條鏈走進同一個子節點 &
第~\ref{sec:tdp-ex-shirts}、\ref{sec:tdp-path}~節 \\
「往下」的方向弄錯 & 翻山路線只看子節點，漏掉經過父節點的下坡 &
第~\ref{sec:tdp-ex-summit}~節 \\
截斷放錯位置 & 換根的 $\mathit{up}$ 沒有取 $\max(0, \cdot)$；用總量相減 &
第~\ref{sec:tdp-ex-power}~節 \\
由上往下那一趟漏看父節點 & 找出擺法時沒看父節點有沒有擺；換根時順序寫反 &
第~\ref{sec:tdp-mis}、\ref{sec:tdp-reroot}~節 \\
太慢 & 換根時每個子節點重算兄弟；表的長度沒有用子樹大小限制 &
第~\ref{sec:tdp-ex-power}、\ref{sec:tdp-preview}~節 \\
\bottomrule
\end{tabular}
\end{center}

\section{本章練習}
\label{sec:ch10-exercises}

\begin{drills}[label=drills:ch10]
\item \label{drill:tdp-trace} 一棵樹有 $7$ 個點，邊是 $1$--$2$、$1$--$3$、$1$--$4$、
$2$--$5$、$2$--$6$、$4$--$7$，營收依序是 $w = (1, 3, 2, 1, 2, 2, 5)$。 以 $1$ 為
根，寫出每個點的 $\mathit{in}$、$\mathit{out}$，最大的營收，以及範例
~\ref{ex:tdp-mis} 會印出的擺法。

\item \label{drill:tdp-cover} 大樓的走廊連成一棵樹，每條走廊的兩端至少要有一端
站著一位保全，在點 $v$ 站保全的成本是 $c_v$。 用兩個狀態寫出求最小總成本的轉移式
和初始值。

\item 用第~\ref{sec:tdp-path}~節的 $\mathit{top}_1$、$\mathit{top}_2$，在第 1 本第~13~章的族譜（圖 13.1，根是 $4$，每條邊長度 $1$）上算出每個點的
$\mathit{top}_1$、$\mathit{top}_2$，確認最大的 $\mathit{top}_1 +
\mathit{top}_2$ 等於第 1 本第~13~章算出的直徑 $5$。

\item 有人把範例~\ref{ex:tdp-distsum} 第二趟的迴圈寫成
\texttt{for (int i = n - 1; i >= 1; i--)}。 範例的輸入會印出什麼？ 為什麼前三個
數還是對的？

\item \label{drill:tdp-ecc} 對樹上每個點 $v$，求從 $v$ 出發最遠能走多遠（邊長都
是 $1$）。 以 $1$ 為根，第一趟算出 $\mathit{top}_1(v)$、$\mathit{top}_2(v)$，並
記下 $\mathit{top}_1(v)$ 是從哪一個子節點來的。 寫出第二趟 $\mathit{up}(v)$ 的式
子：從 $v$ 先走到父節點 $p$、再繼續走，最遠能走多遠。 說明為什麼這裡不能像距離
和那樣「扣掉 $v$ 的那一項」。

\item 第~\ref{sec:tdp-preview}~節的合併，一次的雙層迴圈大約跑 $s \times
\mathit{sz}(c)$ 次。 一條 $n$ 個點的鏈（以一端為根），所有合併加起來大約幾次？ 一
棵星狀的樹（根接著 $n - 1$ 個葉子）呢？
\end{drills}

\section{實作練習}
\label{sec:ch10-practice}

下面是 Aki-Judge 上和本章相關的題目，依小節分組。 難度分四階：暖身是直接套用本
章的寫法；基本要自己決定狀態；實戰的題敘比較長，要先把題目轉換成本章的形狀；挑
戰會和其他章的技巧合在一起。

\begin{center}
\begin{tabular}{llc}
\toprule
小節 & 題目 & 難度 \\
\midrule
子節點合併 & \akiref{squads-that-beat-the-level-above}{贏過上一層的小隊} & 暖身 \\
 & \akiref{are-the-squads-evenly-staffed}{小隊人力是否平均} & 暖身 \\
 & \akiref{hot-water-arrival-spread}{熱水到得夠整齊的閥門} & 暖身 \\
 & \akiref{the-descending-token}{往下走的棋子} & 基本 \\
選或不選 & \akiref{non-adjacent-water-stations}{不相鄰的補水站} & 實戰 \\
\midrule
樹上的路徑 & \akiref{longest-contact-chain-in-every-branch}{每個分部內部的最長聯絡鏈} & 基本 \\
 & \akiref{counting-the-farthest-campsite-pairs}{相隔最遠的營位有幾組} & 實戰 \\
\midrule
換根 & \akiref{where-to-put-the-broadcast-server}{廣播主機該放哪一間教室} & 基本 \\
 & \akiref{levels-after-a-branch-spins-off}{分部獨立後的層級數} & 挑戰 \\
\midrule
依子樹大小合併 & \akiref{pruning-budget}{修剪預算} & 挑戰 \\
\bottomrule
\end{tabular}
\end{center}

本章的三題例題也在 Aki-Judge 上，讀完各節之後可以先自己寫一次再對照。 這一章的
題單還在擴充，之後會再補題。

\section{讀完你應該會}
\label{sec:ch10-checklist}

\begin{itemize}
\item 看到「每個點選或不選、相鄰的有限制」的樹上問題，能把 $v$ 的選擇放進狀態，寫
出用子節點合併的轉移，並倒著看 BFS 順序算完，不用遞迴。
\item 算出最佳值之後，能順著 BFS 順序由上往下，找出一組達到最佳值的選法。
\item 求樹上最好的路徑時，能找出每條路徑唯一的「切點」，在那裡用最大和第二大接起
兩條鏈，並說得出什麼時候不能用兩次走訪。
\item 題目要求每個點當根的答案時，能先寫單一根的樹形 DP，再寫出由上往下的
$\mathit{up}$ 式子，並判斷「扣掉一個子節點」能不能直接做。
\item 估得出狀態是一整張表時的合併次數，知道表的長度要用子樹大小限制。
\end{itemize}

\textbf{延伸閱讀}（書末會整理完整的參考資料）

\begin{itemize}
\item Antti Laaksonen，\textit{Competitive Programmer's Handbook}，第 14 章 14.2、
14.3 節：直徑的動態規劃算法，以及每個點出發的最長路徑（就是本章練習第
~\ref{drill:tdp-ecc}~題），篇幅很短。
\item 劉汝佳，《算法競賽入門經典》第 2 版，9.4.2 節（PDF p.434 起）：樹的最大獨立
集、重心、最長路徑，用中文寫，後面有幾題例題。
\item 吳邦一，《AP325：從 APCS 實作題檢測三級到五級》，第 8 章 8.3 節：樹上的例題
和習題很多，其中不少是由下往上的樹形 DP，適合要準備 APCS 的讀者多練。
\end{itemize}

\addcontentsline{toc}{section}{練習解答}
\section*{練習解答}

\begin{enumerate}
\item BFS 順序是 $1, 2, 3, 4, 5, 6, 7$。 葉子 $3, 5, 6, 7$ 的 $(\mathit{in},
\mathit{out})$ 是 $(2, 0)$、$(2, 0)$、$(2, 0)$、$(5, 0)$。 $2$：$\mathit{in} = 3
+ 0 + 0 = 3$，$\mathit{out} = 2 + 2 = 4$。 $4$：$\mathit{in} = 1 + 0 = 1$，
$\mathit{out} = 5$。 $1$：$\mathit{in} = 1 + 4 + 0 + 5 = 10$，$\mathit{out} =
\max(3, 4) + \max(2, 0) + \max(1, 5) = 11$。 最大是 $11$。 由上往下：$1$ 不去
（$10 < 11$）；$2$ 選 $\mathit{out}$，不去；$3$ 去；$4$ 選 $\mathit{out}$，不去；
$5$、$6$ 的父節點 $2$ 沒去，各自去；$7$ 去。 程式印出 \texttt{11} 和
\texttt{3 5 6 7}。

\item $\mathit{on}(v)$：$v$ 站保全時，$v$ 的子樹裡每條走廊都有人顧的最小成本；
$\mathit{off}(v)$：$v$ 不站時的最小成本。 $v$ 不站，$v$ 和子節點之間的走廊只能靠
子節點顧，所以子節點一定要站；$v$ 站了，子節點站不站都可以：
\[
\mathit{on}(v) = c_v + \sum_c \min\bigl(\mathit{on}(c), \mathit{off}(c)\bigr),\qquad
\mathit{off}(v) = \sum_c \mathit{on}(c).
\]
初始值 $\mathit{on}(v) = c_v$、$\mathit{off}(v) = 0$，答案是根的兩個值取小的。 第
~\ref{sec:tdp-mis}~節常見錯誤第一項那條錯的式子，換到這一題就是對的。

\item 每個點的 $(\mathit{top}_1, \mathit{top}_2)$：葉子 $1, 3, 6, 8, 9$ 都是
$(0, 0)$；$5$ 收到 $6$、$8$ 各 $1$，是 $(1, 1)$；$7$ 收到 $3$ 的 $1$，是 $(1, 0)$；
$2$ 收到 $1$ 的 $1$ 和 $5$ 的 $1 + 1 = 2$，是 $(2, 1)$；$4$ 收到 $2$ 的 $3$、$7$
的 $2$、$9$ 的 $1$，是 $(3, 2)$。 最大的和是 $4$ 的 $3 + 2 = 5$，路徑是 $6 - 5 - 2
- 4 - 7 - 3$（或把 $6$ 換成 $8$），和第 1 本第~13~章一樣。

\item 印出 \texttt{11 10 14 5 3 5 5}。 倒著看時先處理 BFS 順序後面的點，它們父節
點的 \texttt{ans} 還是 $0$，例如 $6$ 算成 $0 - 1 + 6 = 5$。 $2$、$3$ 的父節點是根，
\texttt{ans[1]} 在迴圈之前就算好了，所以它們是對的。

\item 第一趟之外再記 $\mathit{from}(v)$：$\mathit{top}_1(v)$ 是經過哪個子節點來的
（沒有子節點時記 $0$）。 第二趟由上往下，$v$ 的父節點是 $p$：
\[
\mathit{up}(v) = 1 + \max\bigl(\mathit{up}(p),\ \text{$\mathit{from}(p) = v$ 時用
$\mathit{top}_2(p)$，否則用 $\mathit{top}_1(p)$}\bigr),
\]
根的 $\mathit{up}$ 是 $0$。 從 $v$ 出發最遠是 $\max(\mathit{top}_1(v),
\mathit{up}(v))$。 $p$ 最長的鏈如果正好經過 $v$，就不能再走回去，要改用第二長的。 只知道最大值，
拿掉其中一個之後的最大值算不出來，所以不能「扣掉」，要多記第二大。

\item 鏈：每個點只有一個子節點，合併時 $s = 1$，$\mathit{sz}(c)$ 從 $1$ 到 $n - 1$，
總共約 $n^2 / 2$ 次。 星狀：根依序併入 $n - 1$ 個葉子，第 $j$ 次 $s = j$、
$\mathit{sz}(c) = 1$，也約 $n^2 / 2$ 次。 兩種都在正文的上限 $n(n-1)/2$ 附近。
\end{enumerate}

\chapter{綜合練習}
\label{ch:capstone}

% 本章的圖共用的樣式（點、兩種顏色的點、邊、選中的邊、邊上的標籤）
\tikzset{
    capv/.style={circle, draw=black!60, thick, minimum size=0.7cm, inner sep=0pt, font=\small},
    capva/.style={capv, draw=memStroke, fill=memFill},
    capvb/.style={capv, draw=coreStroke, fill=coreFill},
    cape/.style={thick, draw=black!45},
    capehl/.style={line width=1.8pt, draw=memStroke},
    capw/.style={font=\scriptsize, text=black!75, fill=white, inner sep=1pt},
    capq/.style={rectangle, rounded corners=3pt, draw=black, thick, fill=black!4,
        minimum width=3.6cm, minimum height=0.9cm, align=center, font=\sffamily\small},
}

學校園遊會的場地上有六個攤位、一個正門和一個舞台，中間的走道各要走幾分鐘，畫
成一張地圖。 總務處、音控組和巡堂老師各拿著這張地圖來問問題：

\begin{itemize}
\item 總務處要在攤位、正門和舞台之間拉延長線，每一對之間能拉的線長度不同。 電
可以經過別的攤位轉過去，要讓每個地方都有電，線的總長最少是多少？
\item 音控組要從正門把音箱推到舞台，最快要幾分鐘？
\item 巡堂老師從正門出發，每個攤位都要去一次，最後回到正門，最快要幾分鐘？
\end{itemize}

三個問題用的是同一張地圖，看起來都在問「最少」。 往下讀之前，先想想每一題分別
要用本書哪一章的方法。

第一題是第~9~章的最小生成樹：要的是一組把大家連起來的線，沒有起點也沒有終點。
第二題是第~8~章的 Dijkstra：要的是一條從正門到舞台的路。 第三題的點只有八個，
每個攤位都要去而且順序自己排，是第~4~章的旅行推銷員問題；地圖上兩個攤位之間不
一定有直接的走道，所以要先用第~8~章算出每兩個地方之間最快要幾分鐘，再做位元
DP。 同樣的地圖、同樣的「最少」，差別藏在「要連起來」「要一條路」「每個都要
去」這幾個字裡。

《C++ 演算法解題 1：APCS 實作題核心》（以下簡稱第 1 本）第~14~章整理過讀題的方
法：一份自問的問題清單、從資料範圍和題目問法找線索、一套固定的解題流程。 那一
章的內容在這裡都適用，不再重講。 本書前十章每一章教一類技巧，例題也放在那一章
裡，讀的時候已經知道要用什麼；這一章不教新的技巧，練的是在動態規劃的各種狀態、
貪心和圖的各種演算法之間做選擇。

第~\ref{sec:cap-questions}~節是針對本書內容的自問清單，第~\ref{sec:cap-clues}~
節把資料範圍、題目問法和狀態的選法整理成表，第~\ref{sec:cap-flow}~節補充解題流
程。 接著三題例題都要組合兩章以上的方法，每題都從判斷開始，包括先想錯、再修正的
過程。 最後是一份不標章節的總題單。

\section{拿到題目時先問自己的問題}
\label{sec:cap-questions}

第 1 本 14.1 節的清單問的是「看懂了沒」「可以多慢」「怎麼變快」。 下面的清單接
在它後面，專門用來在本書的十章之間做選擇。 卡住的時候回來看看漏了哪一題。

\textbf{第一組：選擇會不會互相牽制？}

\begin{enumerate}
\item 每一步的「最好」能不能只看眼前？ 寫一條貪心規則，馬上用三五筆的小例子找
反例（第 1 本第~7~章）。 找不到反例、又說得出交換論證，才用貪心。
\item 前面的選擇會改變後面能做什麼嗎？ 會的話，「做到一半的樣子」要記下來，這
就是動態規劃的狀態。
\end{enumerate}

\textbf{第二組：如果是動態規劃，狀態要記什麼？}

\begin{enumerate}[resume]
\item 暴力解在枚舉什麼：挑哪些、排什麼順序、怎麼切、每一位填什麼？
\item 「之後還能怎麼做」只取決於什麼？ 用一句話寫出來，例如還剩多少容量、做過
哪幾件、現在在哪裡、這一段的左右端、這一棵子樹。 這句話就是狀態。
\item 狀態有幾個、每個狀態的轉移要看幾個？ 兩數相乘，和時間限制比；再算一次
表格要幾個位元組。
\end{enumerate}

\textbf{第三組：能不能畫成圖？}

\begin{enumerate}[resume]
\item 東西之間的關係是「誰要在誰之前」「誰和誰不能同組」「從哪裡可以走到哪裡」
嗎？ 點是什麼、邊是什麼、有沒有方向、有沒有權重？
\item 要的是一條路，還是一組把大家連起來的邊？ 路上的數字要加起來，還是只看最
大（最小）的那一段？
\item 狀態之間的依賴會不會繞回來？ 不會就照順序填表；會的話，把狀態當成點、轉
移當成邊，改用最短路徑。
\end{enumerate}

\textbf{第四組：題目還藏了什麼？}

\begin{enumerate}[resume]
\item 有沒有一個特別小的量（十幾、二十）或特別大的量（$10^{18}$）？ 前者常常
是位元 DP，後者常常是數位 DP 或對答案二分搜。
\item 詢問很多，但都在問同一份資料嗎？ 和詢問無關的部分能不能先算好，詢問時只
查表？
\end{enumerate}

\section{從題目線索判斷技巧}
\label{sec:cap-clues}

\subsection{資料範圍}

第 1 本 14.2 節的範圍表列的是第 1 本各章的方法，下表補上本書各章的方法。 估計
的基準是時間限制一秒、每秒大約 $10^8$ 次簡單運算（第 1 本第~1~章）。

\begin{center}
\begin{tabular}{p{0.38\linewidth}p{0.54\linewidth}}
\toprule
範圍 & 先想到的方法 \\
\midrule
某一個量不超過 $20$，其他量可以很大 & 位元 DP，$2^n \cdot n$ 或 $2^n \cdot n^2$；
列舉子集合的子集合是 $3^n$，要 $n \le 15$ 左右；每個集合只是要彙總它所有子集合
的值時，用子集合加總（SOS DP），$2^n \cdot n$（第~4~章） \\
數值到 $10^{18}$，問「有幾個」「加起來多少」 & 數位 DP（第~5~章）；問「第 $K$
個」時加上沿著位數往下走或二分搜 \\
$n$ 幾百，子問題是一段一段的 & 區間 DP，$O(n^3)$（第~3~章） \\
$n \le 500$，問每兩點之間 & Floyd--Warshall，$O(n^3)$（第~8~章） \\
兩串各幾千 & $O(nm)$ 的表格（第~2~章），放不下時滾動 \\
物品數乘上容量在 $10^8$ 左右以內 & 背包（第~1~章） \\
\end{tabular}

% 表格分成兩段，讓它可以跨頁，不在上一頁留下大片空白
\begin{tabular}{p{0.38\linewidth}p{0.54\linewidth}}
點數、邊數到 $2 \times 10^5$ 的圖或樹 & 拓樸排序、並查集、Dijkstra、Kruskal、樹形
DP，都在 $O((n + m) \log n)$ 以內（第~6~到第~10~章） \\
樹的點數幾千，限制「最多選 $k$ 個點」 & 樹上背包，表的長度只開到子樹大小，合起來
$O(n^2)$（第~10~章） \\
$nm$ 在 $10^8$ 以內，而且有負權重 & Bellman--Ford（第~8~章） \\
\bottomrule
\end{tabular}
\end{center}

前兩列要分清楚：位元 DP 看的是\textbf{個數}很小，數位 DP 看的是\textbf{數值}很大
但位數很少。 還有一種常見的組合：題目同時給一個很小的 $k$ 和一張很大的圖，例如
$10^5$ 個路口裡只有十幾個地方非去不可。 這時先在大圖上跑十幾次 Dijkstra，把那
十幾個地方之間的距離算好，再對它們做位元 DP。

\subsection{題目的問法和資料的形狀}

\begin{center}
\begin{tabular}{p{0.58\linewidth}p{0.34\linewidth}}
\toprule
題目的問法或資料的形狀 & 先試 \\
\midrule
有一個總量不能超過（或要剛好），每樣東西各佔一些，挑一些出來 & 背包（第~1~章） \\
「照原本的順序」「可以跳過」；兩串要互相對應、配對不能交叉；改幾次才相同 & 子序
列 DP（第~2~章） \\
操作只發生在相鄰的東西之間；回文、括號、不交叉的配對 & 區間 DP（第~3~章） \\
要排順序、走遍每個點、分成幾組，而且某個量很小 & 位元 DP（第~4~章） \\
每個集合都要問「它的所有子集合」加起來多少、最大多少 & 子集合加總（SOS DP，第~4~章） \\
\end{tabular}

\begin{tabular}{p{0.58\linewidth}p{0.34\linewidth}}
「$L$ 到 $R$ 之間有幾個數」，條件可以一位一位檢查 & 數位 DP（第~5~章） \\
「誰要在誰之前」，有向、不會繞成圈；最長、幾種走法、最早完成 & 拓樸順序和 DAG
上的 DP（第~6~章） \\
工作有先後、可以同時做；最早幾天做完、哪些工作一延誤就會拖到全部 & 關鍵路徑（第~6~章） \\
邊一條一條加進來，隨時問連不連通 & 並查集（第~7~章） \\
「同組」「不同組」的限制，一次給完 & 二分圖上色（第~7~章） \\
「同組」「不同組」的限制一條一條給，每條都要馬上知道和前面矛盾不矛盾 & 帶權並查集
（第~7~章） \\
每個點剛好一條往外的箭頭，要走很多步 & 函數圖（第~7~章） \\
從某處到某處，路上的數字加起來最小，可以繞來繞去 & 加權最短路徑（第~8~章） \\
選一些邊讓大家都連通、總和最小；路上「最窄」「最陡」的那一段要最好 & 最小生成樹、
瓶頸路徑（第~9~章） \\
要總和第二小的生成樹（至少有一條邊和最小的那棵不同） & 次小生成樹（第~9~章） \\
$n$ 個點、$n - 1$ 條邊、全部連通；每個點的答案由子節點組成；「每個點當根」 & 樹形
DP、換根（第~10~章） \\
樹上挑一些點，挑了下層就要挑上層（或剪掉一個點就剪掉整棵子樹），有個數或預算限制
& 樹上背包（第~10~章） \\
\bottomrule
\end{tabular}
\end{center}

\subsection{動態規劃的狀態怎麼選}
\label{sec:cap-state}

決定要用動態規劃之後，最難的一步是狀態。 下表把本書各章的狀態放在一起，左欄是
「子問題長什麼樣」，也就是第~\ref{sec:cap-questions}~節第 4 題那句話的樣子。

\begin{center}
\begin{tabular}{p{0.4\linewidth}p{0.34\linewidth}p{0.16\linewidth}}
\toprule
子問題長什麼樣 & 狀態 & 本書 \\
\midrule
從左邊長出來，還要知道用掉多少 & 前 $i$ 個、用了多少容量 & 第~1~章 \\
挑出來的東西要接成一串 & 以第 $i$ 個結尾；兩串時是前 $i$ 個和前 $j$ 個 & 第~2~章 \\
拆一次之後，剩下的是中間一段 & 區間 $[l, r]$ & 第~3~章 \\
之後怎麼做只看「哪些做過了」 & 集合 $S$；需要時加上最後一個 & 第~4~章 \\
一位一位填數字 & 第幾位、是否貼著上界、再記一點點 & 第~5~章 \\
輸入本身是一張沒有環的圖 & 圖上的點，照拓樸順序算 & 第~6~章 \\
輸入是一棵樹 & 子樹的根，加上它自己的選擇 & 第~10~章 \\
輸入是一棵樹，還要分配一個總量 & 子樹的根、子樹裡用掉多少 & 第~10~章 10.5 節 \\
狀態之間會繞回來 & 狀態當成點，用最短路徑 & 第~8~章 8.6 節 \\
\bottomrule
\end{tabular}
\end{center}

同一題常常可以用兩種狀態，選狀態數比較少的那一種。 第~1~章 1.7 節「容量太大」
的情況是一個例子：容量到 $10^9$ 但價值總和很小，就把價值當成表格的一維。 狀態
寫出來之後，再問一次：知道這些，之後的決定真的就夠了嗎？ 本章的例題三會看到，
少記一樣東西，答案就會錯。

\subsection{動態規劃、貪心，還是圖}
\label{sec:cap-decide}

下圖把第一組和第三組的問題接成一條路。 它不是保證，只是一個先後順序：先試最簡
單的，不行再往下走。

\begin{center}
\begin{tikzpicture}[node distance=0.6cm and 0.9cm]
\node[process, minimum width=4.2cm] (brute) {寫出暴力解：\\每一步有哪些選擇};
\node[capq, below=of brute] (q1) {最好的一步\\不看後面就能決定？};
\node[process, right=1.2cm of q1, minimum width=3.6cm] (greedy) {貪心\\先找反例、寫交換論證};
\node[process, below=of q1, minimum width=4.2cm] (state) {把「做到一半的樣子」\\寫成狀態};
\node[capq, below=of state] (q2) {狀態之間的依賴\\會繞回來嗎？};
\node[process, right=1.2cm of q2, minimum width=3.6cm] (dp) {照順序填表\\（第~1~到第~6~、10~章）};
\node[process, below=of q2, minimum width=4.2cm] (sp) {狀態當成點、轉移當成邊\\最短路徑（第~8~章）};
\draw[flowarrow] (brute) -- (q1);
\draw[flowarrow] (q1) -- node[above, font=\scriptsize] {能} (greedy);
\draw[flowarrow] (q1) -- node[right, font=\scriptsize] {不能} (state);
\draw[flowarrow] (state) -- (q2);
\draw[flowarrow] (q2) -- node[above, font=\scriptsize] {不會} (dp);
\draw[flowarrow] (q2) -- node[right, font=\scriptsize] {會} (sp);
\end{tikzpicture}
\diagcaption{從暴力解到演算法的先後順序。 圖要選的是一組邊（最小生成樹）時不走這
條路，見下一小節。}
\label{fig:cap-decide}
\end{center}

\textbf{貪心和動態規劃。} 本書好幾章的開頭都先試過一條貪心規則：第~1~章「比值大
的先拿」、第~3~章「每次合併最少的相鄰兩疊」、第~4~章「每一關挑最熟練的人」、
第~10~章「營收高的先擺」，四條都有反例。 它們錯的原因相同：眼前的選擇會改變之
後能選什麼，而規則沒有看到。 反過來，第 1 本第~7~章的區間排程能用貪心，是因為交
換論證證明了「結束最早的那一場」換進任何最佳解都不吃虧。 所以判斷的順序是：先找
反例，找不到再試著寫交換論證；兩件事都卡住時，就用動態規劃，它比較慢，但不需要
猜。

\textbf{動態規劃和最短路徑。} 第 1 本第~9~章 9.2 節說過，每個動態規劃背後都有一
張 DAG：狀態是點，轉移是邊，計算順序就是拓樸順序（本書第~6~章）。 如果畫出來的
圖有環，例如狀態是「人在哪個路口」，而路口之間可以來回走，就排不出計算順序，表
格填不下去。 這時把同一張圖交給最短路徑演算法：權重都是 $1$ 用 BFS，非負用
Dijkstra，有負的用 Bellman--Ford。 第~8~章 8.6 節的搬運車就是這樣：狀態是「格子
加朝向」，轉彎和前進都有代價，狀態之間會繞圈，所以用 Dijkstra 而不是填表。

\subsection{藏起來的最短路徑和生成樹}
\label{sec:cap-hidden}

有些題目讀起來和圖無關，畫出來才發現是最短路徑或生成樹。 下面幾種本書都出現過：

\begin{itemize}
\item \textbf{等待和時刻表。} 「最早什麼時候到」，可以在原地等，而且早到不會
吃虧：每個點只要記最早到達的時刻，就是一次 Dijkstra，只是走一條邊的時間要看出發
的時刻（第~8~章實作練習的「錯開的紅綠燈」「退潮才露出的平台」）。
\item \textbf{不等式。} 一堆「$x_b \ge x_a + d$」的要求，每一條看成一條帶權重的
邊，問每個 $x$ 的最小值、最大值或有沒有解，就是最長路、最短路和找正（負）環，
用 Bellman--Ford（第~8~章實作練習的「校慶彩排的可排時段」）。
\item \textbf{多一層資訊。} 「最多用一張折價券」「最多轉乘 $k$ 次」，把每個點拆
成「用了幾張」「轉了幾次」的好幾份，就是第~8~章 8.6 節的狀態圖。
\item \textbf{最弱的一環。} 「路上最窄的走廊要盡量寬」「最陡的一段要盡量平
緩」，比的是路上最大（最小）的那一段，不是總和。 這是第~9~章的瓶頸路徑，由小到大加邊，用並查集看兩點什麼時候連通就好。
\item \textbf{全部連起來。} 「讓任兩個都能互通」「每一個都要接上」，而且接法是
選一些兩兩之間的連線，就是生成樹，不需要指定起點。 只有一部分點非接不可時，就不
是生成樹了（第~9~章 9.2 節）。
\end{itemize}

\subsection{線索不是口訣}
\label{sec:cap-traps}

上面的表是用來縮小範圍的。 背成「看到某個字就用某個演算法」會套錯，下面幾個例
子的字面線索都指向一個方法，但多了一個條件，那個方法就不能用：

\begin{itemize}
\item \textbf{「總時間不超過 $T$，挑一些攤位，分數最高」}，但每攤花的時間要看
前一站在哪裡。 字面上是背包，可是背包要求每樣東西各有固定的重量。 這是例題三。
\item \textbf{「最陡的坡盡量平緩，一樣平緩時走最短的」}。 字面上是一次 Dijkstra，
把兩個數一起比；但這樣比會錯，要拆成瓶頸和最短路兩步。 這是例題二。
\item \textbf{「分成兩隊，人數盡量平均」}，但有些人不能同隊。 字面上是子集合和，
可是每個人不能自由選隊，要先用二分圖看出誰被綁在一起。 這是例題一。
\item \textbf{「選一些邊，讓每個村子到鄉公所都盡量近」}。 字面上是「選邊、連通、
最小」，像最小生成樹；可是最小生成樹只管總長，不管每個點到鄉公所多遠。 這要的
是從鄉公所出發的最短路徑樹（第~9~章 9.1 節）。
\end{itemize}

每個例子的處理方式都一樣：表格選出候選之後，回頭確認這個方法需要的前提，題目都
滿足嗎？

\section{解題流程：本書要多注意的地方}
\label{sec:cap-flow}

第 1 本 14.3 節的七個步驟（讀題、手算範例、想暴力解、找瓶頸、換技巧、寫程式、
對拍）照用。 本書的題目在其中三步要多花一點力氣：

\begin{itemize}
\item \textbf{想暴力解。} 動態規劃和圖的題目，暴力解通常是「試遍所有挑法、排法
或路線」，只能跑 $n \le 10$ 左右。 這樣就夠了，它的用途是對拍。 題目的第一個子
題常常就是這個大小。
\item \textbf{換技巧。} 拆成兩件事分開決定：先決定\textbf{模型}（狀態是什麼；點
和邊是什麼），再決定\textbf{演算法}（照順序填表、Dijkstra、Kruskal\ldots）。
模型錯了，演算法再快也沒有用；三題例題裡想錯的地方，都是模型。
\item \textbf{寫程式。} 題目由好幾個階段組成時（先建圖、再上色、再做背包），每
寫完一個階段就把中間結果印出來，用範例對一次，再寫下一個階段。 全部寫完才測，
錯了會不知道是哪一段。
\end{itemize}

\section{例題一：兩個舞台的人數平衡}
\label{sec:cap-stages}

\textbf{題意}（\akiref{balancing-two-stages}{兩個舞台的人數平衡}）

\begin{itemize}
\item 晚會有 $n$ 組表演，每組排到舞台 A 或舞台 B。
\item 有 $m$ 對表演準備了相同的曲目，同一對不能排在同一個舞台。 同一對可能被重
複列出。
\item 舞台 A 有 $x$ 組時，兩邊的差是 $|x - (n - x)|$。 求這個差的最小值 $d$，以及
有幾種排法剛好達到 $d$，除以 $10^9 + 7$ 的餘數。 只要有一組排的舞台不同就算不
同的排法，所以整個對調也算另一種。
\item 怎麼排都會違反限制時，只輸出 \texttt{-1}。
\item $n \le 5000$，$m \le 2 \times 10^5$。
\end{itemize}

例如 $n = 7$，限制是 1--2、2--3、4--5、5--6。 把限制畫成圖（下圖），1、2、3 連
成一群，4、5、6 一群，7 自己一群。 $n$ 是奇數，差至少是 $1$；舞台 A 排 1、3、5，
舞台 B 排 2、4、6、7 就是 $3$ 對 $4$。 差是 $1$ 的排法一共有 $6$ 種，答案是
\texttt{1 6}。

\begin{center}
\begin{tikzpicture}[x=1.2cm, y=1cm]
\node[capva] (1) at (0,0) {1};
\node[capvb] (2) at (1,0) {2};
\node[capva] (3) at (2,0) {3};
\node[capva] (4) at (3.6,0) {4};
\node[capvb] (5) at (4.6,0) {5};
\node[capva] (6) at (5.6,0) {6};
\node[capva] (7) at (7.2,0) {7};
\draw[cape] (1) -- (2) -- (3);
\draw[cape] (4) -- (5) -- (6);
\node[font=\scriptsize, text=black!70] at (1,-0.75) {$(2, 1)$};
\node[font=\scriptsize, text=black!70] at (4.6,-0.75) {$(2, 1)$};
\node[font=\scriptsize, text=black!70] at (7.2,-0.75) {$(1, 0)$};
\end{tikzpicture}
\diagcaption{範例一的限制畫成圖，每一群從編號最小的那組開始上兩種顏色。 下方
括號裡是這一群橘色、藍色各有幾組。 每一群只能整群決定「橘色去 A」或「藍色去
A」。}
\label{fig:cap-stages}
\end{center}

\textbf{判斷}

「不能排在同一個舞台」，而且只有兩個舞台：這是第~7~章的二分圖，用 BFS
上兩種顏色，碰到相鄰兩點同色就是奇環，輸出 \texttt{-1}。 可是上色只回答了「能
不能排」。 一群裡只要決定一組在哪一邊，其他組就全部跟著定了，所以每一群剛好有
兩種排法：橘色的去 A，或藍色的去 A。 範例的三群分別讓 A 多出 $2$ 或 $1$ 組、
$2$ 或 $1$ 組、$1$ 或 $0$ 組。

接下來是「兩邊盡量平均」。 字面上像第~1~章 1.5 節的子集合和：把東西分成兩
堆，總和盡量接近。

\textbf{第一次想錯：貪心。} 一群一群處理，每次把人多的那一邊放到目前人少的舞
台。 試範例二：$n = 6$，限制是 3--4、3--5、3--6。 群是 $\{1\}$、$\{2\}$ 和
$\{3, 4, 5, 6\}$，最後一群一邊 $1$ 組（3）、一邊 $3$ 組（4、5、6）。 照編號
處理，1 去 A、2 去 B，兩邊一樣多；最後一群人多的那邊隨便放到一邊，就變成 $4$
對 $2$，差是 $2$。 正確答案是 $0$：1、2、3 在一邊，4、5、6 在另一邊。 先排大
的群也不行，這和「分兩堆」的貪心一樣有反例：五群各讓兩邊差 $3$、$3$、$2$、
$2$、$2$ 組時，由大到小輪流放到少的一邊，最後差 $2$，可是 $3 + 3 = 2 + 2 + 2$，
差可以是 $0$。 前面的選擇會影響後面怎麼湊，這是第一組問題的第 2 題，要動態規
劃。

\textbf{換成背包。} 把每一群看成一組東西，裡面有兩個選項：讓舞台 A 多 $a$ 組，
或多 $b$ 組，每組\textbf{恰好選一個}。 這是第~1~章 1.7 節「每組恰好選一個」的
背包，容量是舞台 A 的組數，最多 $n$。 題目還要排法數，所以把取最大值換成加法，
像第~1~章 1.5 節數方法數那樣：

\begin{itemize}
\item \textbf{狀態}：$\mathit{ways}[x]$ 是處理完前幾群之後，舞台 A 剛好 $x$ 組的
排法數。
\item \textbf{轉移}：處理一群 $(a, b)$ 時，開一個新的陣列，
$\mathit{next}[x + a] \mathrel{+}= \mathit{ways}[x]$，
$\mathit{next}[x + b] \mathrel{+}= \mathit{ways}[x]$。
\item \textbf{初始值}：還沒處理任何一群時，$\mathit{ways}[0] = 1$。
\item \textbf{答案}：全部處理完，在做得到的 $x$ 裡找 $|2x - n|$ 最小的；
$x$ 和 $n - x$ 的差一樣，兩個都要加起來。
\end{itemize}

範例一的三群依序是 $(2, 1)$、$(2, 1)$、$(1, 0)$，表格的變化如下（空白是
$0$）：

\begin{center}
\begin{tabular}{lcccccccc}
\toprule
處理完 & $x = 0$ & $1$ & $2$ & $3$ & $4$ & $5$ & $6$ & $7$ \\
\midrule
還沒開始 & 1 & & & & & & & \\
第一群 $(2, 1)$ & & 1 & 1 & & & & & \\
第二群 $(2, 1)$ & & & 1 & 2 & 1 & & & \\
第三群 $(1, 0)$ & & & 1 & 3 & 3 & 1 & & \\
\bottomrule
\end{tabular}
\end{center}

$x = 3$ 和 $x = 4$ 的差都是 $1$，排法數 $3 + 3 = 6$。

\textbf{程式}

\begin{example}[label=ex:cap-stages]{cpp}{(stages.cpp) 兩個舞台的人數平衡}
#include <cstdlib>
#include <iostream>
#include <queue>
#include <utility>
#include <vector>
using namespace std;

const long long MOD = 1000000007;

int main() {
    ios::sync_with_stdio(false);
    cin.tie(nullptr);
    int n, m;
    cin >> n >> m;
    vector<vector<int>> adj(n + 1);
    for (int i = 0; i < m; i++) {
        int u, v;
        cin >> u >> v;
        adj[u].push_back(v);
        adj[v].push_back(u);
    }
    // 第一步：每一群用 BFS 上兩色，記下兩色各有幾組
    vector<int> color(n + 1, -1);
    vector<pair<int, int>> group;     // 每一群兩種顏色各幾組
    for (int s = 1; s <= n; s++) {
        if (color[s] != -1) {
            continue;
        }
        int cnt[2] = {0, 0};
        queue<int> q;
        color[s] = 0;
        q.push(s);
        while (!q.empty()) {
            int u = q.front();
            q.pop();
            cnt[color[u]]++;
            for (int v : adj[u]) {
                if (color[v] == -1) {
                    color[v] = 1 - color[u];
                    q.push(v);
                } else if (color[v] == color[u]) {
                    cout << -1 << '\n';          // 奇環，分不成兩邊
                    return 0;
                }
            }
        }
        group.push_back({cnt[0], cnt[1]});
    }
    // 第二步：每群恰好選一邊去舞台 A 的背包
    // ways[x]：舞台 A 剛好 x 組的排法數（模 MOD）
    // can[x]：舞台 A 剛好 x 組做不做得到
    vector<long long> ways(n + 1, 0);
    vector<char> can(n + 1, 0);
    ways[0] = 1;
    can[0] = 1;
    int total = 0;                    // 已處理的群加起來幾組
    for (auto [a, b] : group) {
        vector<long long> nways(n + 1, 0);
        vector<char> ncan(n + 1, 0);
        for (int x = 0; x <= total; x++) {
            if (!can[x]) {
                continue;
            }
            // 這一群顏色 0 的去 A，或顏色 1 的去 A
            nways[x + a] = (nways[x + a] + ways[x]) % MOD;
            nways[x + b] = (nways[x + b] + ways[x]) % MOD;
            ncan[x + a] = ncan[x + b] = 1;
        }
        ways = nways;
        can = ncan;
        total += a + b;
    }
    int d = n + 1;
    for (int x = 0; x <= n; x++) {
        if (can[x]) {
            d = min(d, abs(2 * x - n));
        }
    }
    long long count = 0;
    for (int x = 0; x <= n; x++) {    // A 比較多、比較少都要算
        if (can[x] && abs(2 * x - n) == d) {
            count = (count + ways[x]) % MOD;
        }
    }
    cout << d << ' ' << count << '\n';
    return 0;
}
\end{example}

\iosample{
7 4 \\
1 2 \\
2 3 \\
4 5 \\
5 6 \\
}{
1 6 \\
}

\begin{itemize}
\item 「做不做得到」用另一個陣列 \texttt{can} 記，不用 \texttt{ways[x] != 0}
判斷。 排法數是取過餘數的，真正的排法數不是 $0$，餘數卻可能剛好是 $0$。 這種情
況很少見，測試資料可能測不到，但它是錯的。
\item 一群的兩種顏色一樣多時（$a = b$），兩行都加到同一格，那一格加了兩次。 這
是對的：橘色去 A 和藍色去 A 是兩種不同的排法，只是人數相同。
\item 奇環一出現就可以結束：不管其他群怎麼排，這一群已經排不了。
\item \texttt{x} 只跑到 \texttt{total}，也就是已經處理的群加起來的組數；更大的
\texttt{x} 一定做不到。 \texttt{x + a}、\texttt{x + b} 也不會超過 $n$。
\end{itemize}

\textbf{對拍用的暴力解}

題目的第一個子題是 $n \le 15$。 這個大小可以試遍 $2^n$ 種排法，一個一個檢查限
制，正好拿來對拍：

\begin{example}[label=ex:cap-stages-brute]{cpp}{(stages\_brute.cpp) 兩個舞台的暴力解}
#include <cstdlib>
#include <iostream>
#include <vector>
using namespace std;

int main() {
    int n, m;
    cin >> n >> m;
    vector<int> u(m), v(m);
    for (int i = 0; i < m; i++) {
        cin >> u[i] >> v[i];
        u[i]--;
        v[i]--;
    }
    int best = n + 1;
    long long count = 0;
    // mask 的第 i 位是 1：第 i + 1 組排在舞台 A
    for (int mask = 0; mask < (1 << n); mask++) {
        bool ok = true;
        for (int i = 0; i < m && ok; i++) {
            if (((mask >> u[i]) & 1) == ((mask >> v[i]) & 1)) {
                ok = false;
            }
        }
        if (!ok) {
            continue;
        }
        int d = abs(2 * __builtin_popcount(mask) - n);
        if (d < best) {
            best = d;
            count = 0;
        }
        if (d == best) {
            count++;
        }
    }
    if (best > n) {
        cout << -1 << '\n';
    } else {
        cout << best << ' ' << count << '\n';
    }
    return 0;
}
\end{example}

\iosample{
6 3 \\
3 4 \\
3 5 \\
3 6 \\
}{
0 2 \\
}

題目其他兩個子題也是提示。 「群的數量不超過 $18$」是在說：每群只有兩種選擇，
試遍 $2^{18}$ 種也來得及，這一步就已經看出「以群為單位」。 「$m = 0$」時每組自
己一群，舞台 A 有 $x$ 組的排法數是組合數 $\binom{n}{x}$。 讀完子題再決定做法，
常常能少走一段冤枉路。

\textbf{複雜度}

BFS 是 $O(n + m)$。 背包最多 $n$ 群，每群掃最多 $n + 1$ 格，最壞約
$5000 \times 5000 = 2.5 \times 10^7$ 次。 最大的測試資料在本書測試的環境下約
$0.02$ 秒。

\subsection{常見錯誤}

\begin{itemize}
\item \textbf{貪心地放到人少的一邊}：範例二印出差 \texttt{2}，正確是 \texttt{0}。
\item \textbf{兩色一樣多時只加一次}：以為兩種排法人數一樣就是同一種。 輸入
\texttt{2 1}、\texttt{1 2} 時，正確答案是 \texttt{0 2}（1 在 A 或 2 在 A），這樣
寫印出 \texttt{0 1}。
\item \textbf{只算 A 比較少的那一邊}：答案只加 $x \le n / 2$ 的那一格，範例一印
出 \texttt{1 3}。 題目說整個對調也算另一種，$x = 3$ 和 $x = 4$ 都要算。
\item \textbf{沒有檢查奇環}：上色時碰到同色的鄰居沒有處理。 輸入 \texttt{3 3}、
\texttt{1 2}、\texttt{2 3}、\texttt{3 1} 時應該印出 \texttt{-1}，這樣寫印出
\texttt{1 2}。
\end{itemize}

\section{例題二：坡度最緩的輪椅路線}
\label{sec:cap-slope}

\textbf{題意}（\akiref{the-gentlest-wheelchair-route}{坡度最緩的輪椅路線}）

\begin{itemize}
\item $n$ 個地點、$m$ 條雙向步道，第 $i$ 條連接 $u_i$ 和 $v_i$，長度 $l_i$，坡
度等級 $g_i$。 同一對地點之間可能有好幾條步道。
\item 一條路線的「最陡坡度」是路上坡度等級的最大值，「總長度」是長度的總和。
\item 從 $s$ 到 $t$，先讓最陡坡度最小，記為 $G$；最陡坡度等於 $G$ 的路線裡，再
選總長度最短的。 輸出 $G$ 和這個總長度；到不了 $t$ 時輸出 \texttt{-1}。
\item $n \le 10^5$，$m \le 2 \times 10^5$，$l_i, g_i \le 10^9$。
\end{itemize}

下圖是題目的範例一，$s = 1$、$t = 5$，邊上寫著「長度／坡度」。 走 $1 \to 4 \to 5$
最短，可是坡度是 $5$。 最陡坡度可以壓到 $3$，有兩條：$1 \to 3 \to 5$ 長 $5$，
$1 \to 2 \to 3 \to 5$ 長 $11$。 答案是 \texttt{3 5}。

\begin{center}
\begin{tikzpicture}[x=1.5cm, y=1.1cm]
\node[capv] (1) at (0,1) {1};
\node[capv] (2) at (1.5,2.2) {2};
\node[capv] (3) at (3,1) {3};
\node[capv] (4) at (2.25,-0.4) {4};
\node[capv] (5) at (4.5,1) {5};
\draw[cape] (1) -- node[capw, above left] {4/1} (2);
\draw[cape] (2) -- node[capw, above right] {4/1} (3);
\draw[capehl] (1) -- node[capw] {2/2} (3);
\draw[capehl] (3) -- node[capw] {3/3} (5);
\draw[cape] (1) -- node[capw, below left] {1/5} (4);
\draw[cape] (4) -- node[capw, below right] {1/5} (5);
\end{tikzpicture}
\diagcaption{範例一，邊上是「長度／坡度」。 橘色是答案的路線 $1 \to 3 \to 5$。
走到 $3$ 時，從 $2$ 繞過來比較平緩，可是比較長。}
\label{fig:cap-slope}
\end{center}

\textbf{判斷}

「從 $s$ 到 $t$」「最短」，邊有權重：第~8~章的 Dijkstra。 可是這題有兩個標準，
而且第一個標準「最陡的一段最小」比的不是總和，是第~\ref{sec:cap-hidden}~小節說
的「最弱的一環」。

\textbf{第一次想錯：一次 Dijkstra，兩個數一起比。} 把每個點的標籤改成一對數
（最陡坡度，總長度），比較時先比坡度、一樣再比長度；走一條邊時坡度取最大值、長
度相加。 範例一會印出 \texttt{3 11}。 看點 $3$：從 $1 \to 2 \to 3$ 來是
$(1, 8)$，從 $1 \to 3$ 來是 $(2, 2)$。 照這個比法 $(1, 8)$ 比較好，點 $3$ 只留下
它。 可是接上 $3 \to 5$ 這條坡度 $3$ 的步道之後，兩個都變成坡度 $3$，長度是
$11$ 和 $5$，被丟掉的那一個反而比較好。

錯在哪裡？ Dijkstra 每個點只留一個最好的值，靠的是：到 $u$ 比較好的走法，接上
同一條邊之後還是比較好（至少不會變差）。 第~8~章 8.2 節的證明默默用到這件事。
總長度有這個性質：兩個數加上同一個數，大小順序不變。 （坡度，長度）這一對沒有：
取最大值會把坡度的差別抹掉，抹掉之後改比長度，而長度的大小可能和原本相反。 所以
一個點只留一個「最好」的標籤是不夠的。

\textbf{第二次想錯：最小生成樹上的路。} 第~9~章的瓶頸路徑證明過，最小生成樹上 $s$
到 $t$ 的那條路，最陡的一段就是所有路線裡最小的 $G$。 那就在最小生成樹上找這條
路，順便算長度？ 範例一的 Kruskal 依坡度先選 1--2、2--3（坡度 $1$），1--3 會造
成環所以不選，再選 3--5，樹上的路是 $1 \to 2 \to 3 \to 5$，長 $11$，又印出
\texttt{3 11}。 樹上的路保證坡度最小，但它只是坡度最小的路線之一，生成樹不管長
度。

\textbf{修正：拆成兩步。} 兩次想錯的地方都是把兩個標準混在一起。 分開來看：

\begin{enumerate}
\item \textbf{先求 $G$。} 步道依坡度由小到大，用並查集一條一條接上（第~9~章瓶頸路徑的離線做法），接完某一條之後 $s$ 和 $t$ 第一次連通，這一條的坡度就是 $G$。 在
它之前接上的步道都連不通 $s$ 和 $t$，所以沒有最陡坡度比 $G$ 小的路線；接上它之
後，有一條全部坡度都 $\le G$ 的路線。 整個過程沒有連通，就輸出 \texttt{-1}。
\item \textbf{再求長度。} 最陡坡度剛好是 $G$ 的路線，就是「只用坡度 $\le G$ 的步
道」的路線：這種路線的最陡坡度 $\le G$，又不可能比 $G$ 小。 所以把坡度 $\le G$
的步道\textbf{全部}留下（不只是生成樹上的），在上面做一次普通的 Dijkstra。
\end{enumerate}

範例一的第一步接完 1--2、2--3、1--3、3--5 時 $1$ 和 $5$ 連通，$G = 3$。 第二步
只用這四條步道，最短的是 $1 \to 3 \to 5$，長 $5$。

\textbf{程式}

\begin{example}[label=ex:cap-slope]{cpp}{(slope.cpp) 坡度最緩的輪椅路線}
#include <algorithm>
#include <functional>
#include <iostream>
#include <queue>
#include <utility>
#include <vector>
using namespace std;

struct Edge {
    int u, v;
    long long len, g;
};

const long long INF = 1e18;
vector<int> parent, sz;

int find(int x) {
    return parent[x] == x ? x : parent[x] = find(parent[x]);
}

void unite(int a, int b) {
    a = find(a);
    b = find(b);
    if (a == b) {
        return;
    }
    if (sz[a] > sz[b]) {
        swap(a, b);
    }
    parent[a] = b;
    sz[b] += sz[a];
}

int main() {
    ios::sync_with_stdio(false);
    cin.tie(nullptr);
    int n, m, s, t;
    cin >> n >> m >> s >> t;
    vector<Edge> edges(m);
    for (Edge &e : edges) {
        cin >> e.u >> e.v >> e.len >> e.g;
    }
    // 第一步：坡度由小到大加邊，s 和 t 一連通，最後加的坡度就是 G
    vector<Edge> byG = edges;
    sort(byG.begin(), byG.end(),
         [](const Edge &a, const Edge &b) { return a.g < b.g; });
    parent.resize(n + 1);
    sz.assign(n + 1, 1);
    for (int v = 1; v <= n; v++) {
        parent[v] = v;
    }
    long long G = -1;
    for (const Edge &e : byG) {
        unite(e.u, e.v);
        if (find(s) == find(t)) {
            G = e.g;
            break;
        }
    }
    if (G == -1) {
        cout << -1 << '\n';
        return 0;
    }
    // 第二步：只用坡度不超過 G 的步道，做 Dijkstra
    vector<vector<pair<int, long long>>> adj(n + 1);
    for (const Edge &e : edges) {
        if (e.g <= G) {
            adj[e.u].push_back({e.v, e.len});
            adj[e.v].push_back({e.u, e.len});
        }
    }
    vector<long long> dist(n + 1, INF);
    priority_queue<pair<long long, int>, vector<pair<long long, int>>,
                   greater<pair<long long, int>>> pq;
    dist[s] = 0;
    pq.push({0, s});
    while (!pq.empty()) {
        auto [d, u] = pq.top();
        pq.pop();
        if (d != dist[u]) {
            continue;
        }
        for (auto [v, w] : adj[u]) {
            if (d + w < dist[v]) {
                dist[v] = d + w;
                pq.push({dist[v], v});
            }
        }
    }
    cout << G << ' ' << dist[t] << '\n';
    return 0;
}
\end{example}

\iosample{
5 6 1 5 \\
1 2 4 1 \\
2 3 4 1 \\
1 3 2 2 \\
3 5 3 3 \\
1 4 1 5 \\
4 5 1 5 \\
}{
3 5 \\
}

\begin{itemize}
\item 並查集和 Dijkstra 都是前面各章的寫法，沒有改動；這題的工作全在決定「先做
哪一個、把什麼交給下一步」。
\item 第一步走到 $G$ 就停，不必把所有步道接完。 第二步一定走得到 $t$，因為第一
步已經確認過了。
\item 總長度最多約 $10^5 \times 10^9 = 10^{14}$，\texttt{dist} 要用
\texttt{long long}。
\end{itemize}

\textbf{複雜度}

排序 $O(m \log m)$，並查集幾乎是線性，Dijkstra 是 $O((n + m) \log n)$。 最大的
測試資料在本書測試的環境下約 $0.05$ 秒。

\subsection{常見錯誤}

\begin{itemize}
\item \textbf{（坡度，長度）一起比的 Dijkstra}：範例一印出 \texttt{3 11}。 題目的
完整測試資料裡，這樣寫錯了 $11$ 筆。
\item \textbf{只在最小生成樹上找路}：範例一同樣印出 \texttt{3 11}。 生成樹只保
證最陡的一段，長度要回到所有坡度 $\le G$ 的步道上算。
\item \textbf{第二步用「坡度 $< G$」}：等於 $G$ 的步道被拿掉，範例一的 3--5 不見
了，$5$ 走不到，\texttt{dist[t]} 是無限大。 不等號要和 $G$ 的定義一致。
\end{itemize}

\section{例題三：社團博覽會的體驗時間}
\label{sec:cap-expo}

\textbf{題意}（\akiref{booth-hopping-at-the-club-fair}{社團博覽會的體驗時間}）

\begin{itemize}
\item 有 $n$ 個體驗攤位，正門是地點 $0$。 完成攤位 $i$ 的體驗得 $p_i$ 分。
\item $t_{i,j}$ 是從地點 $i$ 直接到地點 $j$ 的分鐘數，已經包含在 $j$ 體驗的時間；
$t_{i,0}$ 是走回正門的時間。 表不一定對稱，也不保證繞路比較慢。 每一步都要直接
走過去，不能經過別的攤位或正門轉過去。
\item 有 $q$ 位新生，第 $k$ 位有 $T_k$ 分鐘。 每人各自從正門出發，去若干個不同的
攤位，最後回到正門，總時間不超過 $T_k$；也可以一攤都不去。 對每個人求最高的分數。
\item $n \le 16$，$q \le 2 \times 10^5$，$t_{i,j} \le 10^6$，$T_k \le 10^9$。
\end{itemize}

下面用題目的範例二說明：三個攤位，分數 $8$、$5$、$5$，時間表如下（第 $i$ 列第
$j$ 行是 $t_{i,j}$）。 攤位 1 來回 $3 + 3 = 6$ 分鐘，得 $8$ 分；攤位 2、3 彼此只
要 $2$ 分鐘，$0 \to 2 \to 3 \to 0$ 要 $4 + 2 + 4 = 10$ 分鐘，得 $10$ 分。 有 $10$
分鐘的人選後者，有 $6$ 到 $9$ 分鐘的人只能去攤位 1。

\begin{center}
\begin{tabular}{c|cccc}
\toprule
從 \textbackslash{} 到 & 0 & 1 & 2 & 3 \\
\midrule
0 & 0 & 3 & 4 & 4 \\
1 & 3 & 0 & 6 & 6 \\
2 & 4 & 6 & 0 & 2 \\
3 & 4 & 6 & 2 & 0 \\
\bottomrule
\end{tabular}
\end{center}

\textbf{判斷}

「時間不超過 $T$，挑一些攤位，每攤最多一次，分數最高」：字面上是第~1~章的 0/1
背包。 檢查背包的前提：每樣東西要有\textbf{固定}的重量。 這裡去一個攤位花多少
時間，要看前一站在哪裡、下一站去哪裡，同樣是攤位 2，接在攤位 3 後面只要 $2$ 分
鐘，從正門過去要 $4$ 分鐘。 沒有固定的重量，背包不能用。

時間和順序有關，就要排順序。 第四組問題的第 9 題：$n \le 16$ 特別小，指向第~4~
章。 排順序時，「之後還要多少時間」只看兩件事：已經去過哪些攤位、現在在哪一攤。
這正是第~4~章 4.3 節旅行推銷員的狀態 $(S, \mathit{last})$。 有些位元 DP 的題目
要先用第~8~章算出地點之間的最短距離，這一題不用，也不能：題目規定每一步要直接
走，經過別的攤位就等於去了那一攤。

\textbf{第一次想錯：每個人各做一次。} 對每位新生，做一次旅行推銷員的 DP，看哪
些集合走得完。 一次 DP 是 $2^{16} \times 16 \times 16 \approx 1.7 \times 10^7$，
乘上 $q = 2 \times 10^5$ 位新生，完全來不及。

\textbf{找瓶頸。} 這是第四組問題的第 10 題。 DP 裡沒有用到 $T$：一個集合走完要
多少時間，和誰來走無關，$T$ 只在最後比一下。 所以 DP 只做一次，對每個集合 $S$
算出
\[
\mathit{cost}(S) = \min_{\mathit{last} \in S}\bigl(\mathit{dp}[S][\mathit{last}]
+ t_{\mathit{last}, 0}\bigr),
\]
也就是去過剛好 $S$ 這些攤位、再回到正門的最少時間，以及 $S$ 的分數總和
$\mathit{score}(S)$。 空集合的時間和分數都是 $0$。

每個詢問變成：$\mathit{cost}(S) \le T$ 的集合裡，最大的 $\mathit{score}(S)$。
$2^{16}$ 個集合每個人都掃一遍，是 $2 \times 10^5 \times 65536 \approx 1.3 \times
10^{10}$，還是太慢。 把 $2^n$ 個 $(\mathit{cost}, \mathit{score})$ 依時間排序，
算分數的前綴最大值，詢問時用 \texttt{upper\_bound} 找到最後一個時間 $\le T$ 的
位置，讀那裡的前綴最大值（第 1 本第~2、3~章）。

範例二排序之後是這樣：

\begin{center}
\begin{tabular}{lcccccccc}
\toprule
集合 & 空 & $\{1\}$ & $\{2\}$ & $\{3\}$ & $\{2, 3\}$ & $\{1, 2\}$ & $\{1, 3\}$ & $\{1, 2, 3\}$ \\
\midrule
時間 & 0 & 6 & 8 & 8 & 10 & 13 & 13 & 15 \\
分數 & 0 & 8 & 5 & 5 & 10 & 13 & 13 & 18 \\
前綴最大值 & 0 & 8 & 8 & 8 & 10 & 13 & 13 & 18 \\
\bottomrule
\end{tabular}
\end{center}

$T = 10$ 時最後一個時間 $\le 10$ 的是 $\{2, 3\}$，答案 $10$；$T = 9$ 落在
$\{3\}$，答案 $8$；$T = 5$ 只有空集合，答案 $0$。

\textbf{程式}

\begin{example}[label=ex:cap-expo]{cpp}{(expo.cpp) 社團博覽會的體驗時間}
#include <algorithm>
#include <iostream>
#include <utility>
#include <vector>
using namespace std;

const long long INF = 1e18;

int main() {
    ios::sync_with_stdio(false);
    cin.tie(nullptr);
    int n, q;
    cin >> n >> q;
    vector<long long> p(n);
    for (long long &x : p) {
        cin >> x;
    }
    vector<vector<long long>> t(n + 1, vector<long long>(n + 1));
    for (auto &row : t) {
        for (long long &x : row) {
            cin >> x;
        }
    }
    // 攤位在程式裡編號 0 到 n - 1，地點編號是攤位編號加 1
    // 正門是地點 0
    int full = (1 << n) - 1;
    // dp[mask][last]：去過的攤位剛好是 mask、現在在 last 的最少時間
    vector<vector<long long>> dp(1 << n, vector<long long>(n, INF));
    for (int i = 0; i < n; i++) {
        dp[1 << i][i] = t[0][i + 1];
    }
    for (int mask = 1; mask <= full; mask++) {
        for (int last = 0; last < n; last++) {
            long long cur = dp[mask][last];
            if (cur == INF) {
                continue;
            }
            for (int j = 0; j < n; j++) {
                if ((mask >> j) & 1) {
                    continue;
                }
                long long cost = cur + t[last + 1][j + 1];
                long long &target = dp[mask | (1 << j)][j];
                target = min(target, cost);
            }
        }
    }
    // 每個集合：走完再回正門的最少時間、分數總和
    // plan[mask]：(時間, 分數)
    vector<pair<long long, long long>> plan(1 << n);
    plan[0] = {0, 0};
    for (int mask = 1; mask <= full; mask++) {
        long long best = INF;
        for (int last = 0; last < n; last++) {
            if (dp[mask][last] < INF) {
                best = min(best, dp[mask][last] + t[last + 1][0]);
            }
        }
        int low = __builtin_ctz(mask);
        long long score = plan[mask & (mask - 1)].second + p[low];
        plan[mask] = {best, score};
    }
    // 依時間排序；maxScore[k]：前 k + 1 個集合裡最高的分數
    sort(plan.begin(), plan.end());
    vector<long long> times(plan.size()), maxScore(plan.size());
    for (size_t k = 0; k < plan.size(); k++) {
        times[k] = plan[k].first;
        maxScore[k] = plan[k].second;
        if (k > 0) {
            maxScore[k] = max(maxScore[k], maxScore[k - 1]);
        }
    }
    for (int i = 0; i < q; i++) {
        long long T;
        cin >> T;
        // 時間 <= T 的集合是前面 k 個，空集合一定在裡面
        auto it = upper_bound(times.begin(), times.end(), T);
        int k = it - times.begin();
        cout << maxScore[k - 1] << '\n';
    }
    return 0;
}
\end{example}

\iosample{
3 4 \\
8 5 5 \\
0 3 4 4 \\
3 0 6 6 \\
4 6 0 2 \\
4 6 2 0 \\
10 6 9 5 \\
}{
10 \\
8 \\
8 \\
0 \\
}

\begin{itemize}
\item DP 的部分和第~4~章 4.4 節的回收推車幾乎一樣，只是那一題只要「全部去過」
的答案，這一題每個集合都要。 \texttt{cost(S)} 的定義是「剛好去過 $S$」，所以每
個集合都要各自加上回正門的時間，不能只看 \texttt{full}。
\item 分數總和用第~4~章的寫法，從拿掉最低位的集合加一項得到。 \texttt{plan[mask
\& (mask - 1)]} 在這時候還沒被排序打亂，因為排序在迴圈後面。
\item 空集合 \texttt{plan[0]} 的時間是 $0$，題目保證 $T \ge 1$，所以
\texttt{upper\_bound} 的結果至少是 $1$，\texttt{k - 1} 不會是 $-1$。
\end{itemize}

\textbf{複雜度}

DP 是 $2^{16} \times 16 \times 16 \approx 1.7 \times 10^7$ 次轉移，表格是
$2^{16} \times 16$ 個 \texttt{long long}，約 $8$ MB。 排序 $2^{16}$ 筆，每個詢問
一次二分搜，約 $16$ 步。 最大的測試資料在本書測試的環境下約 $0.04$ 秒。

\subsection{常見錯誤}

\begin{itemize}
\item \textbf{忘了回程}：\texttt{cost} 只取 \texttt{dp[mask][last]}，沒有加上回
正門的時間。 範例一的第一位新生有 $20$ 分鐘，正確答案是 \texttt{13}，這樣寫印出
\texttt{23}：四攤依序走到最後一攤剛好 $20$ 分鐘，可是走不回來。
\item \textbf{先跑 Floyd--Warshall}：習慣性地先把時間表換成最短時間。 題目規定
不能經過別的攤位轉過去，這樣做會讓「路過」的攤位不算數。 輸入
\texttt{2 1}、\texttt{1 5}、\texttt{0 1 10}、\texttt{1 0 1}、\texttt{10 1 0}、
\texttt{4}（六行）時正確答案是 \texttt{1}，這樣寫印出 \texttt{6}。 題目的兩個範
例都測不出來，完整測試資料裡錯了 $4$ 筆。
\item \textbf{沒有放空集合}：\texttt{plan} 只放非空的集合。 範例二的 $T = 5$ 時
\texttt{k} 是 $0$，讀到 \texttt{maxScore[-1]}，這是未定義的行為。
\item \textbf{每個詢問各做一次 DP}：答案對，但 $q$ 一大就超過時間限制。
\end{itemize}

回頭看三題的判斷：例題一先用圖看出誰被綁在一起，再用背包；例題二先用生成樹的想法
求瓶頸，再用最短路徑；例題三先排除背包，用位元 DP 對每個集合算一次，再用排序和二
分搜回答詢問。 每一題用的都是學過的東西，難在看出要拆成哪幾步、每一步交給哪一章。

\section{考場策略}
\label{sec:cap-exam}

第 1 本 14.8 節整理了 APCS 實作題官方公告的作答與計分規則，以及從規則推出來的時
間分配。 那些建議在這裡照用；下面補充本書這類題目要多注意的幾件事。

\begin{itemize}
\item \textbf{子題是提示。} 例題一的第二個子題「群的數量不超過 $18$」已經點出
「以群為單位」。 動手之前，把每個子題的限制讀一次，問自己這個限制讓什麼做法變得
可行。
\item \textbf{先寫能拿到小子題的暴力解。} 動態規劃和圖的題目，暴力解常常就是第
1 本第~4~章的枚舉。 它先拿到部分分，等一下也是對拍的工具。
\item \textbf{寫表格之前先算記憶體。} 例題三是 $2^{16} \times 16$ 個
\texttt{long long}，$8$ MB；同樣的表格遇到 $n = 20$ 就是 $160$ MB。 每格的值用
\texttt{int} 放得下時，記憶體就少一半。
\end{itemize}

\begin{apcsnote}
官方公告的題本種類、檢測時間、計分方式和上傳規則，整理在第 1 本 14.8 節，依的是
2026 年 11 月的官網公告。 規則會調整，報名前請再看一次官網，以官網和考試時題目
上寫的為準。
\end{apcsnote}

\section{本章常見錯誤}
\label{sec:cap-pitfalls}

這一章的錯誤大多出在「選模型」，演算法本身都寫對了。

\begin{center}
\begin{tabular}{p{0.22\linewidth}p{0.46\linewidth}p{0.22\linewidth}}
\toprule
錯誤 & 例子 & 該問的問題 \\
\midrule
套用前沒確認前提 & 例題三當成背包；例題二以為 Dijkstra 能同時比兩個數 & 這個方法需要
什麼？ \\
貪心沒先找反例 & 例題一放到人少的一邊 & 第一組第 1 題 \\
狀態少記一樣 & 例題二每個點只留一個標籤；例題三忘了回程 & 第二組第 4 題 \\
兩個標準混在一起 & 例題二在生成樹上找路 & 能不能拆成兩步？ \\
多做了題目不准的事 & 例題三先跑 Floyd--Warshall & 回頭讀題 \\
算重或漏算 & 例題一兩色一樣多只加一次、只算一邊 & 用暴力解對拍 \\
\bottomrule
\end{tabular}
\end{center}

最後兩列都是範例測不出來的錯誤。 例題三的 Floyd--Warshall 兩個範例都會過，例題一
「只加一次」要遇到兩色一樣多的群才會錯。 第 1 本第~1~章的對拍腳本，在這一章的每
一題都值得跑一次。

\section{本章練習}
\label{sec:cap-exercises}

\begin{drills}[label=drills:cap]
\item 下面每一題只給大意和範圍。 寫出你第一個想試的方法、它的複雜度，以及它在
本書的哪一章。
\begin{enumerate}
\item $n \le 10^5$ 個路口、$m \le 2 \times 10^5$ 條單向道路，各有通行時間，求消
防隊到每個路口最快要多久。
\item $n \le 15$ 個景點，給定兩兩之間的交通時間，從飯店出發每個景點去一次再回飯
店，最短要多久。
\item 求 $1$ 到 $N$ 之中，各位數字和是 $7$ 的倍數的整數有幾個，$N \le 10^{18}$。
\item $n \le 300$ 堆石頭排成一列，每次合併相鄰兩堆，花費是兩堆的總和，合成一堆的
最小花費。
\item 公司的組織圖是一棵樹，每人有一個開心值，直屬上司和下屬不能同時參加聚餐，
參加的人開心值總和最大。
\item $n \le 2 \times 10^5$ 件工作和它們的先後關係，每件要花一些天，沒有先後關
係的可以同時做，最早第幾天全部做完。
\item $n \le 10^5$ 座小島、$m \le 2 \times 10^5$ 座可以蓋的橋各有造價，要讓所有
小島互通的最小總造價。
\item 兩份長度 $5000$ 的打卡紀錄，找出兩份共有、順序相同的最長一串。
\end{enumerate}

\item 下面每一題都有一個看起來很像的線索。 說明照線索直接套為什麼不對，以及應
該怎麼做。
\begin{enumerate}
\item 從 $1$ 號站坐到 $n$ 號站，路上最貴的那一段票價要盡量便宜。 （線索：從某處
到某處、最小。）
\item 颱風過後，鄉公所只能先搶修 $n - 1$ 條道路，條件是每個村子到鄉公所的最短
距離都要和道路全通時一樣。 要修哪幾條？ （線索：選 $n - 1$ 條邊、連通。）
\item $n \le 2 \times 10^5$ 個不超過 $10^9$ 的正整數，挑一些出來，總和是 $3$ 的
倍數，總和最大是多少。 （線索：挑一些、總和。）
\item 一張有向圖，從 $s$ 出發可以重複走，每個點的分數只在第一次到達時拿到，最多
拿幾分。 （線索：有向圖上的最大值。）
\end{enumerate}

\item 照例題一的程式，寫出下面這組輸入每一群的 $(a, b)$、每處理完一群之後的
\texttt{ways}，以及輸出。
\begin{code}{text}{練習 3 的輸入}
5 3
1 2
1 3
4 5
\end{code}

\item 例題二的延伸。
\begin{enumerate}
\item 把規則倒過來：先讓總長度最短，總長度一樣時再讓最陡坡度最小。 這時「一次
Dijkstra，標籤是（總長度，最陡坡度），依序比較」對不對？ 說明理由。
\item 第一步不用並查集：把所有不同的坡度排好，對 $G$ 二分搜，每次用 BFS 檢查只
走坡度 $\le G$ 的步道能不能從 $s$ 到 $t$。 這樣的複雜度是多少？
\end{enumerate}

\item 例題三的延伸。
\begin{enumerate}
\item 題目改成逛完最後一攤就從那裡離開，不必回正門。 程式要改哪裡？
\item 如果只有一位新生（$q = 1$），還需要排序和前綴最大值嗎？
\end{enumerate}

\item 估計。
\begin{enumerate}
\item 例題一的 $m = 0$、$n = 5000$ 時，背包的內層迴圈一共跑幾次？ 如果
\texttt{x} 每次都從 $0$ 跑到 $n$ 呢？
\item 例題三如果 $n = 20$，\texttt{dp} 表要多少記憶體、多少次轉移？ 還來得及嗎？
\end{enumerate}
\end{drills}

\section{總題單}
\label{sec:cap-practice}

下面是 Aki-Judge 上的題目，從本書各章挑出來，另外加了幾題沒有出現在前面題單裡的。
順序是打亂的，也不標章節：每一題先照第~\ref{sec:cap-questions}~節的問題判斷要用
什麼方法，寫完再看本章最後的解答，那裡只寫用到的技巧和對應的章。 前面已經寫過的
題目，這次從判斷開始重做一次。

\begin{center}
\begin{tabular}{rlc}
\toprule
 & 題目 & 難度 \\
\midrule
1 & \akiref{campus-fiber-upgrade}{校園光纖補線} & 實戰 \\
2 & \akiref{multiples-both-ways}{正讀倒讀都是 K 的倍數} & 實戰 \\
3 & \akiref{the-3d-printer-that-needs-to-cool-down}{需要冷卻的 3D 印表機} & 基本 \\
4 & \akiref{the-robot-that-never-stops}{不停站的導覽機器人} & 實戰 \\
5 & \akiref{squared-scores-for-palindromic-pieces}{回文切段的平方得分} & 實戰 \\
6 & \akiref{one-errand-on-the-way-home}{放學順路買東西} & 基本 \\
7 & \akiref{pruning-budget}{修剪預算} & 挑戰 \\
8 & \akiref{how-light-can-the-heaviest-box-be}{最重那箱可以多輕} & 實戰 \\
9 & \akiref{threshold-skill-tree}{門檻式技能樹} & 實戰 \\
10 & \akiref{drone-payload-balance}{無人機載重配貨} & 基本 \\
11 & \akiref{power-lines-for-typhoon-season}{颱風季的電線規劃} & 實戰 \\
12 & \akiref{songs-that-might-be-kept}{哪些歌曲可能被保留} & 實戰 \\
13 & \akiref{errands-before-school-starts}{開學前的採買路線} & 挑戰 \\
14 & \akiref{the-descending-token}{往下走的棋子} & 基本 \\
15 & \akiref{rehearsal-start-windows}{校慶彩排的可排時段} & 實戰 \\
16 & \akiref{how-far-the-digit-stickers-go}{數字貼紙夠貼到幾號} & 實戰 \\
17 & \akiref{piggy-bank-range-totals}{存錢筒的區間總額} & 實戰 \\
18 & \akiref{rivals-and-partners-in-the-tug-of-war}{有恩怨也有搭檔的分組} & 挑戰 \\
19 & \akiref{orienteering-passphrase}{定向越野的通關密語} & 實戰 \\
20 & \akiref{ten-minutes-flat}{一律十分鐘的接駁車} & 實戰 \\
21 & \akiref{how-late-can-you-arrive}{最晚幾點到夜市} & 挑戰 \\
22 & \akiref{non-adjacent-water-stations}{不相鄰的補水站} & 實戰 \\
\bottomrule
\end{tabular}
\end{center}

本章的三題例題也在 Aki-Judge 上。 讀完題意之後先蓋住後面，自己走一次判斷，再和
書上的過程對照。 第 18 題和例題一是同一種組合，可以接著寫。

\section{讀完你應該會}
\label{sec:cap-checklist}

\begin{itemize}
\item 拿到一題陌生的題目，能從資料範圍和題目問法，在本書十章裡列出兩三個候選，並
說出每個候選需要的前提（固定的重量、權重非負、沒有環、相鄰才有限制等），逐一確認。
\item 寫一條貪心規則之前，會先用小例子找反例；找到反例時，能說出眼前的選擇影響了
後面的什麼，並把它寫成動態規劃的狀態。
\item 能用一句話寫出「之後還能怎麼做只取決於什麼」，對照狀態表選出狀態，並估計狀態
數、轉移數和記憶體。
\item 看到「最弱的一環」「全部連起來」「可以等待」「不等式」這類說法，能認出背後的
瓶頸路徑、生成樹或最短路徑。
\item 題目有兩個標準或兩個階段時，能拆成先後兩步，各用一章的方法，並在中間印出結
果檢查。
\end{itemize}

\textbf{延伸閱讀}（書末會整理完整的參考資料）

\begin{itemize}
\item Steven S. Skiena，\textit{The Algorithm Design Manual}，第 3 版，8.7 節
「Design Graphs, Not Algorithms」：作者的看法是，很多題目的關鍵在把圖建對，而不
是發明新的演算法；可以和本章的例題一、例題二對照。
\item Steven Halim 等，\textit{Competitive Programming 4}，Book 2，8.7 節「Problem
Decomposition」：把需要兩種技巧的題目依組合方式分類，例如「對答案二分搜加上別
的」「圖的前處理加上動態規劃」。
\item Thomas H. Cormen 等，\textit{Introduction to Algorithms}，第 4 版，14.3 節和
15.2 節：分別整理動態規劃和貪心各需要什麼性質，可以用來檢查第~\ref{sec:cap-decide}
~小節的判斷順序。
\end{itemize}

\addcontentsline{toc}{section}{練習解答}
\section*{練習解答}

\begin{enumerate}
\item
\begin{enumerate}
\item Dijkstra，$O((n + m) \log n)$。 第~8~章。
\item 旅行推銷員的位元 DP，狀態是「去過哪些、現在在哪」，$O(2^n n^2)$，
$2^{15} \times 225 \approx 7 \times 10^6$。 第~4~章 4.3 節。
\item 數位 DP，狀態是「第幾位、是否貼著上界、目前數字和除以 $7$ 的餘數」。 第~5~章。
\item 區間 DP，看最後一次合併的分界，$O(n^3) = 2.7 \times 10^7$。 第~3~章。
\item 樹形 DP，每個人記「參加」「不參加」兩個狀態。 第~10~章 10.2 節。
\item 照拓樸順序算每件工作最早的完成日，也就是 DAG 上的最長路（關鍵路徑），
$O(n + m)$。 第~6~章 6.5 節。
\item 最小生成樹，Kruskal，$O(m \log m)$。 第~9~章。
\item 最長共同子序列，$O(nm) = 2.5 \times 10^7$。 第~2~章 2.3 節。
\end{enumerate}

\item
\begin{enumerate}
\item 最短路徑把票價加起來，這裡只看最貴的一段，是瓶頸路徑。 依票價由小到大加
邊，用並查集看 $1$ 和 $n$ 什麼時候連通（第~9~章的瓶頸路徑）。 也可以用 Dijkstra，
把「加起來」換成「取最大值」：取最大值不會讓路變得更好，Dijkstra 的論證仍然成立。
\item 「選 $n - 1$ 條邊、連通」像最小生成樹，但最小生成樹只讓總長最小，不保證每
個村子到鄉公所最近。 要的是從鄉公所出發的最短路徑樹：做一次 Dijkstra，每個村子
留下最短路上的前一條邊（第~9~章 9.1 節、第~8~章 8.2 節）。
\item 「挑一些、總和」像背包，但總和可以到 $2 \times 10^{14}$，容量開不出來。 判斷只需要總和除以 $3$ 的餘數，狀態改成「前 $i$ 個數、總和除以 $3$
餘幾」的最大總和，只有 $3$ 種，$O(n)$（第~1~章 1.7 節「容量太大」的想法）。
\item 「有向圖上的最大值」像第~6~章 DAG 上的 DP，但這張圖可能有環，排不出拓樸順序。
同一個環上的點，走到一個就能全部走到，要先把互相走得到的點合併成一個點，變成 DAG
再做（第~6~章 6.4 節最後一小節；合併的方法叫強連通元件，在《C++ 演算法解題 3》）。
\end{enumerate}

\item BFS 從 1 開始，1 是顏色 $0$，2、3 是顏色 $1$，第一群是 $(1, 2)$；4 是顏色 $0$、
5 是顏色 $1$，第二群是 $(1, 1)$。
\begin{center}
\begin{tabular}{lcccccc}
\toprule
處理完 & $x = 0$ & $1$ & $2$ & $3$ & $4$ & $5$ \\
\midrule
還沒開始 & 1 & & & & & \\
第一群 $(1, 2)$ & & 1 & 1 & & & \\
第二群 $(1, 1)$ & & & 2 & 2 & & \\
\bottomrule
\end{tabular}
\end{center}
第二群兩色一樣多，每一格都加了兩次。 $n = 5$，$x = 2$ 和 $x = 3$ 的差都是 $1$，
輸出 \texttt{1 4}。

\item
\begin{enumerate}
\item 對。 設到某個點有兩種走法，標籤 $(L_1, G_1)$ 比 $(L_2, G_2)$ 好，接上同一
條長度 $l$、坡度 $g$ 的步道。 如果 $L_1 < L_2$，接上之後仍然 $L_1 + l < L_2 + l$；
如果 $L_1 = L_2$、$G_1 \le G_2$，接上之後長度相同，而
$\max(G_1, g) \le \max(G_2, g)$。 比較好的接上之後仍然不比較差，而且每走一步長度至
少加 $1$，標籤只會變大，第~8~章 8.2 節的論證可以照用。 原本的順序不行，是因為取最
大值可能把第一個數變成相同，這時第二個數的順序可以和原本相反。
\item 不同的坡度最多 $m$ 種，二分搜約 $\log_2 m \approx 18$ 次，每次 BFS 是
$O(n + m)$，加上排序，共 $O(m \log m + (n + m) \log m)$，和並查集的做法一樣來得及。
\end{enumerate}

\item
\begin{enumerate}
\item 算 \texttt{plan} 時不加回程：\texttt{best} 改成
\texttt{dp[mask][last]} 的最小值。 其他都不用改。
\item 不需要。 只問一次時，把 $2^n$ 個集合掃一遍，在時間 $\le T$ 的集合裡取最高
分，$O(2^n)$，比排序還快。 排序和前綴最大值是為了讓很多次詢問共用同一份整理好的
資料。
\end{enumerate}

\item
\begin{enumerate}
\item $m = 0$ 時每組自己一群，共 $5000$ 群，處理第 $k$ 群時 \texttt{total} 是
$k - 1$，迴圈跑 $k$ 次，總共 $1 + 2 + \cdots + 5000 \approx 1.25 \times 10^7$ 次。
每次都從 $0$ 跑到 $n$ 的話是 $5000 \times 5001 \approx 2.5 \times 10^7$ 次，多一
倍，不過也還來得及。
\item $2^{20} \times 20 \approx 2.1 \times 10^7$ 格，每格 $8$ 位元組是 $160$ MB
（這裡的 MB 和第~4~章一樣是 $2^{20}$ 位元組），改用 \texttt{int}（最大時間
$21 \times 10^6$ 放得下）是 $80$ MB。 轉移是
$2^{20} \times 20 \times 20 \approx 4.2 \times 10^8$ 次，在 $1$ 秒內很勉強。 這
也是第~4~章把「$n$ 不超過 $20$ 左右」當成位元 DP 訊號的原因。
\end{enumerate}
\end{enumerate}

\subsection*{總題單的解答}

只寫用到的技巧和對應的章。 有些題不只一種做法，這裡寫的是其中一種。

\begin{enumerate}
\item 舊光纖先用並查集合併；新光纖依造價排序做一次 Kruskal，記下每接受一條之後
的累計造價；每個預算用二分搜找最多能接受幾條（第~7、9~章，第 1 本第~3~章）。
\item 數位 DP，狀態同時記「$x$ 除以 $K$ 的餘數」和「倒過來的數除以 $K$ 的餘數」，
位數固定時，第幾位在倒過來的數裡乘上 $10$ 的幾次方是確定的（第~5~章）。
\item 冷卻時間併進結束時刻，每個委託佔用 $[s_i, t_i + r_i)$，變成有價值的區間排
程；依結束時刻排序，用二分搜找前面接得上的（第 1 本第~9~章 9.7 節）。 依結束時刻
貪心會錯，因為酬勞不同。
\item 每個站拆成「走了偶數步」「走了奇數步」兩個狀態做 BFS；第 $k$ 分鐘能不能在 $v$，
只要看和 $k$ 同奇偶的最短步數有沒有超過 $k$，因為可以來回走一條步道多花 $2$ 分鐘
（第~7~章二分圖的奇偶想法，第 1 本第~12~章的狀態 BFS）。 站 $s$ 沒有步道時要另外處理。
\item 先用區間的方法算出每一段是不是回文（第~3~章 3.5 節），再做「把前 $i$ 顆切
成好幾段」的一維 DP（第 1 本第~9~章）。
\item 從學校在原圖上、從家在反向圖上各做一次 Dijkstra，每種商品取所有賣它的地點
$v$ 的 $d_s(v) + d_t(v)$ 最小值（第~8~章 8.5 節）。
\item 剪掉一個點就是剪掉整棵子樹。 可以照第~10~章 10.5 節的樹上背包合併，表的維度
改成花掉的時間；但時間不像點數那樣受子樹大小限制，每次合併都跑滿 $T$ 會太慢。 比
較穩的做法是用 DFS 的前序把每棵子樹排成連續的一段，在前序上做背包：留下這個點就
看下一個位置，剪掉就跳過整段子樹，$O(nT)$（第~1~章）。
\item 箱子由輕到重一箱一箱加進 0/1 背包的可達表，記下每個總重第一次變成湊得出來
時，加的是哪一箱；那一箱就是最輕可能的「最重那箱」。 詢問直接查表（第~1~章 1.5
節）。
\item 把 Kahn 演算法的「入度變成 $0$」換成「學會的相關技能數達到門檻」，一天處理
一層（第~6~章 6.2 節）。
\item 「剛好裝滿」的 0/1 背包，算出每個總重的最高運費，在 $L$ 到 $W$ 之間取最大值
（第~1~章 1.3 節）。
\item 和例題二同一種拆法：先依風險做 Kruskal，找出最小的最高風險 $B$；再只用風險
$\le B$ 的電線，依造價做一次 Kruskal（第~9~章）。
\item 正著、倒著各做一次最長共同子序列；位置 $i$ 和 $j$ 的代號相同，而且前面一段
加上後面一段再加 $1$ 等於最長長度時，這兩個位置可能被用到（第~2~章，和第 1 本第
~14~章例題二「從兩端各算一次」同一類想法）。
\item 從家和每一家店各做一次 Dijkstra，得到店與店之間的最短距離；再做位元 DP，狀
態是「辦好了哪幾種事、現在在哪一家店」，轉移時檢查先後條件（第~8~章、第~4~章）。
\item 由下往上算，依深度的奇偶決定這一步是誰走，取子節點的最大值或最小值（第~10~
章）。
\item 差分約束：每條要求變成一條帶權重的邊，用 Bellman--Ford 從下限往上推、從上
限往下壓，各做一次；做完還在變就是不可行（Bellman--Ford 見第~8~章 8.7 節）。
\item 對貼好的號碼數二分搜；檢查時一位一位算出 $S$ 到 $S + k - 1$ 每個數字各用了
幾張（第~5~章，第 1 本第~3~章）。
\item 每天的金額只由前一天決定，是函數圖上的一條路。 找出尾巴和圈，配上前綴和回
答區間總和（第~7~章的函數圖，第 1 本第~5~章）。
\item 搭檔和恩怨用二分圖上色或帶權並查集（第~7~章）分成好幾群，每群只能整群選
一邊。 問題變成子集合和，相同大小的群合併成多重背包做二進位拆分，可達性用
\texttt{bitset} 存，加一件物品只要一次位移和一次 OR（第~7~章、第~1~章 1.6 節）。
\item 狀態是「蓋到第 $k$ 個章、人在哪個檢查點」。 $k$ 每一步加一，就算路線圖有
環，狀態圖也是一層一層的 DAG，逐層推（第~6~章 6.4 節，第~8~章 8.6 節）。
\item 每條接駁車路線加一個虛擬點：站點走到虛擬點的權重是 $r_{j,x} + c_j$，虛擬點
回到站點是 $r_{j,y}$，再做 Dijkstra（第~8~章）。 直接在停靠站之間兩兩連邊會太多。
\item 已經買過的攤位一定是連續的一段，狀態是「買過哪一段、站在左端還是右端」，從
整段往回推「最晚可以站在那裡的時刻」（第~3~章）。
\item 樹形 DP，每個點記三種狀態：自己設站；沒設站但被某個子節點蓋到；沒設站，要
靠父節點（第~10~章 10.2 節）。
\end{enumerate}


% 書末（不編章號）
% 書末：參考資料與延伸閱讀（不編章號）
\chapter*{參考資料與延伸閱讀}
\addcontentsline{toc}{chapter}{參考資料與延伸閱讀}
\markboth{參考資料與延伸閱讀}{參考資料與延伸閱讀}
\label{ch:references}

各章最後的「延伸閱讀」提到的書，完整的書目資料都列在這裡，每本附一兩句它適合
做什麼，以及和本書哪幾章搭配。 括號裡的章節號碼是那本書的章節，以這裡列出的
版本為準；換了版本，章節號碼可能不同。

本書的講解、例子和程式都是自己寫的，參考這些書的地方是題型的分類、講解的先後
和經典的結論。 想看另一種講法、更嚴謹的證明，或想多做一些題目時，可以從下面
挑一本。

\section*{篇幅短、適合快速複習}

\begin{itemize}
\item Antti Laaksonen，\textit{Competitive Programmer's Handbook}，2018 年 7 月
3 日草稿版，作者公開的 PDF。

每個主題只用幾頁、一兩個小例子講完，適合讀完本書某一章後，花十幾分鐘從另一個
角度再看一次。 和本書搭配的地方：第~1、2~章（第 7 章的背包、最長遞增子序列、
編輯距離），第~6~章（第 16 章前半），第~7~章（12.3 節，第 15、16 章），第~8~章
（第 13 章），第~9~章（第 15 章），第~10~章（第 14 章）。

\item Antti Laaksonen，\textit{Guide to Competitive Programming: Learning and
Improving Algorithms Through Contests}，第 3 版，Springer，2024。

和上一本出自同一位作者，章節安排相近，是正式出版的書，內容較新。 本書第~4~%
章引用了它 6.2.4、6.2.5 節：從排列改成子集合的動態規劃，以及一列一列數鋪磚方法
的寫法。
\end{itemize}

\section*{題型分類細、題目多}

\begin{itemize}
\item Steven Halim、Felix Halim、Suhendry Effendy，\textit{Competitive
Programming 4}，Book 1 與 Book 2，2020（作者自行出版）。

把題目依技巧分得很細，每一節後面附一長串分好類的練習題，適合打算參加程式競
賽、想大量練習的讀者。 和本書搭配的地方：第~1~章（Book 1 的 3.5.2 節），第~2~章（Book 1
的 3.5.2 節，Book 2 的 6.3 節），第~4~章（Book 1 的 3.5.2 節，Book 2 第 8 章），第~5~章
（Book 2 的 6.3.3 節），第~6~章（Book 1 的 4.2.6 節），第~7~章（Book 1 的
2.4.2、4.2.7、4.2.8 節），第~8~章（Book 1 的 4.4、4.5 節），第~9~章（Book 1
的 4.3 節）。

\item 劉汝佳，《算法競賽入門經典》，第 2 版，清華大學出版社，2014。

簡體中文。 例題多，程式寫得精簡，適合想用中文多看幾種題型的讀者。 和本書搭配
的地方：第~1~章（9.3.2 節），第~3~章（9.4.1 節），第~5~章（10.4 節例題
10-22），第~6~章（6.4.3、9.2 節），第~8~章（11.3 節），第~10~章（9.4.2 節）。

\item 吳邦一，《AP325：從 APCS 實作題檢測三級到五級》，v1.5，2024 年 9 月，作者
公開的講義。

以 APCS 實作題的程度為目標，每個主題附很多例題和習題，還有測試資料，適合要準
備 APCS 的讀者多練。 它的第 6 章（動態規劃）、第 7 章（基本圖論演算法）、第 8
章（樹上演算法）和本書的範圍重疊最多；各章引用的是第~3~章（6.2.4 節）和第~10~%
章（8.3 節）。
\end{itemize}

\section*{嚴謹的證明與演算法設計}

\begin{itemize}
\item Thomas H. Cormen、Charles E. Leiserson、Ronald L. Rivest、Clifford Stein，
\textit{Introduction to Algorithms}，第 4 版，The MIT Press，2022。

大學演算法課常用的教科書，每個演算法都有完整的正確性證明和分析。 想知道本書
「為什麼一定對」框裡省略的細節，或要修演算法課的讀者，可以對照著讀。 和本書搭
配的地方：第~2~章（14.4 節），第~3~章（14.2 節），第~6~章（20.4、22.2 節），
第~7~章（第 19 章），第~8~章（第 22 章、23.2 節），第~9~章（第 21 章）。

\item Steven S. Skiena，\textit{The Algorithm Design Manual}，第 3 版，Springer，
2020。

重點放在怎麼想出演算法、怎麼把實際的問題對應到已知的問題，後半是一份常見問題
的目錄。 適合想練習建模的大學生和自學的人。 和本書搭配的地方：第~2~章（10.2
節），第~4~章（10.9 節）。
\end{itemize}

\section*{本系列的其他書}

\begin{itemize}
\item 《C++ 演算法解題 1：APCS 實作題核心》：本書的前一本，各章提到的「第 1 本」。
排序、二分搜、前綴和、堆疊與佇列、貪心、一維與格子的動態規劃、BFS 與 DFS、樹的
基本做法都在那裡。
\item 《C++ 程式設計：給學過 C 的人》：各章提到的「銜接書」，本書用到的 C++ 語法
和 STL 都在那裡。
\end{itemize}


\end{document}
